\documentclass[12pt,oneside,openany]{book}
\usepackage[T1]{fontenc}
\usepackage[utf8]{inputenc}
\usepackage{amsmath,amssymb,amsthm,mathtools,mathrsfs}
\usepackage{booktabs,array}
\usepackage[numbers,sort&compress]{natbib}
\usepackage{environ}
\PassOptionsToPackage{hypertexnames=false}{hyperref}
\makeatletter

\usepackage{lmodern}
\usepackage[sc]{mathpazo}
\usepackage{bm}
\linespread{1.1}
\usepackage{microtype}
\microtypesetup{stretch=15,shrink=15}

\usepackage{geometry}
\geometry{letterpaper,left=1in,right=1in,top=0.75in,bottom=0.75in,includeheadfoot,
  headheight=15pt,headsep=16pt,footskip=26pt}
\setlength{\parindent}{1.5em}
\setlength{\parskip}{0pt}
\setlength{\emergencystretch}{3em}
\tolerance=1000
\clubpenalty=3000 \widowpenalty=3000 \displaywidowpenalty=1500
\raggedbottom
\allowdisplaybreaks[2]

\usepackage{enumitem}
\setlist{topsep=4pt plus 2pt minus 1pt,itemsep=2pt plus 1pt,parsep=0pt}
\setlist[itemize]{leftmargin=1.5em}
\setlist[enumerate]{leftmargin=2em}

\usepackage[newparttoc]{titlesec}
\usepackage{titletoc,fancyhdr}

\newcommand{\trilogy@spacedsc}[1]{{\scshape\textls[90]{\MakeLowercase{#1}}}}
\newcommand{\trilogy@headpar}{\rightskip=0pt plus 0.5\hsize\relax
  \parfillskip=0pt plus 0.5\hsize\relax\hyphenpenalty=10000\relax\exhyphenpenalty=10000\relax}

\titleformat{\part}[display]
  {\normalfont\centering}
  {\large\trilogy@spacedsc{\partname\ \thepart}}
  {22pt}
  {\fontsize{22}{28}\selectfont}
\assignpagestyle{\part}{empty}

\titleformat{\chapter}[display]
  {\normalfont\trilogy@headpar}
  {\trilogy@spacedsc{\chaptertitlename\ \thechapter}}
  {12pt}
  {\fontsize{20}{25}\selectfont}
\titleformat{name=\chapter,numberless}[display]
  {\normalfont\trilogy@headpar}
  {}
  {0pt}
  {\fontsize{20}{25}\selectfont\trilogy@starmark}
\titlespacing*{\chapter}{0pt}{12pt}{30pt}
\newcommand{\trilogy@starmark}[1]{#1\markboth{#1}{#1}}

\titleformat{\section}
  {\normalfont\fontsize{13.5}{17}\selectfont\bfseries\trilogy@headpar}{\thesection}{0.75em}{}
\titlespacing*{\section}{0pt}{22pt plus 4pt minus 2pt}{9pt plus 2pt}
\titleformat{\subsection}
  {\normalfont\normalsize\bfseries\trilogy@headpar}{\thesubsection}{0.75em}{}
\titlespacing*{\subsection}{0pt}{15pt plus 3pt minus 2pt}{6pt plus 1pt}
\titleformat{\paragraph}[runin]
  {\normalfont\normalsize\bfseries}{}{0pt}{}
\titlespacing*{\paragraph}{0pt}{9pt plus 2pt minus 1pt}{0.6em}

\setcounter{tocdepth}{1}
\setcounter{secnumdepth}{1}

\newcommand{\shorthead}[2]{\expandafter\def\csname trilogy@sh@\detokenize{#1}\endcsname{#2}}
\newcommand{\trilogy@lookup}[1]{%
  \ifcsname trilogy@sh@\detokenize{#1}\endcsname
    \expandafter\let\expandafter\trilogy@headtext\csname trilogy@sh@\detokenize{#1}\endcsname
  \else
    \def\trilogy@headtext{#1}%
  \fi}
\newcommand{\trilogy@headfont}{\normalfont\small\itshape}
\DeclareRobustCommand{\trilogy@hd}[2]{%
  \ifx\relax#1\relax\else{\upshape#1}\hspace{0.8em}\fi#2}
\newsavebox{\trilogy@chapbox}
\newsavebox{\trilogy@secbox}
\newsavebox{\trilogy@cachedchapbox}
\newsavebox{\trilogy@cachedsecbox}
\newif\iftrilogy@headtoolong
\newif\iftrilogy@cachedheadtoolong
\let\trilogy@cachedheadkey\relax
\newcommand{\trilogy@computeheads}{%
  \trilogy@headtoolongfalse
  \sbox\trilogy@chapbox{\trilogy@headfont\nouppercase{\leftmark}}%
  \sbox\trilogy@secbox{}%
  \protected@edef\trilogy@l{\leftmark}%
  \protected@edef\trilogy@r{\rightmark}%
  \ifx\trilogy@r\@empty\else\ifx\trilogy@r\trilogy@l\else
    \sbox\trilogy@secbox{\trilogy@headfont\nouppercase{\rightmark}}%
  \fi\fi
  \dimen@=\wd\trilogy@chapbox\advance\dimen@ 4em\relax
  \ifdim\dimen@>\headwidth
    \sbox\trilogy@chapbox{\trilogy@headfont\@chapapp\ \thechapter}%
    \trilogy@headtoolongtrue
  \fi
  \dimen@=\wd\trilogy@chapbox\advance\dimen@\wd\trilogy@secbox\advance\dimen@ 8em\relax
  \ifdim\dimen@>\headwidth
    \sbox\trilogy@secbox{\trilogy@headfont
      \expandafter\let\csname trilogy@hd \endcsname\trilogy@hdnum\nouppercase{\rightmark}}%
  \fi}
\newcommand{\trilogy@heads}[1]{%
  \protected@edef\trilogy@headkey{%
    {\leftmark}{\rightmark}{\the\headwidth}%
    {\fontname\font}{\the\dimexpr1em\relax}{\thechapter}{\@chapapp}}%
  \ifx\trilogy@headkey\trilogy@cachedheadkey\else
    \trilogy@computeheads
    \global\setbox\trilogy@cachedchapbox=\copy\trilogy@chapbox
    \global\setbox\trilogy@cachedsecbox=\copy\trilogy@secbox
    \global\let\trilogy@cachedheadkey\trilogy@headkey
    \iftrilogy@headtoolong
      \global\trilogy@cachedheadtoolongtrue
    \else
      \global\trilogy@cachedheadtoolongfalse
    \fi
  \fi
  \setbox\trilogy@chapbox=\copy\trilogy@cachedchapbox
  \setbox\trilogy@secbox=\copy\trilogy@cachedsecbox
  \iftrilogy@cachedheadtoolong
    \ifx\relax#1\relax\else\PackageWarningNoLine{trilogy-style}{Running head too long on
      page \thepage; give the chapter a \string\shorthead}\fi
  \fi}
\newcommand{\trilogy@hdnum}[2]{\ifx\relax#1\relax\else\S\,{\upshape#1}\fi}

\pagestyle{fancy}
\fancyhf{}
\fancyhead[L]{\trilogy@heads{warn}\usebox\trilogy@chapbox}
\fancyhead[R]{\trilogy@heads{}%
  \ifdim\wd\trilogy@secbox>\z@\usebox\trilogy@secbox\hspace{2em}\fi
  \normalfont\small\thepage}
\renewcommand{\headrulewidth}{0pt}
\renewcommand{\footrulewidth}{0pt}
\fancypagestyle{plain}{\fancyhf{}\fancyfoot[C]{\normalfont\small\thepage}%
  \renewcommand{\headrulewidth}{0pt}\renewcommand{\footrulewidth}{0pt}}
\renewcommand{\chaptermark}[1]{\trilogy@lookup{#1}%
  \markboth{\trilogy@hd{\ifnum\c@secnumdepth>\m@ne\if@mainmatter\thechapter\fi\fi}{\trilogy@headtext}}{}}
\renewcommand{\sectionmark}[1]{\trilogy@lookup{#1}%
  \markright{\trilogy@hd{\ifnum\c@secnumdepth>\z@\thesection\fi}{\trilogy@headtext}}}

\contentsmargin{2.2em}
\titlecontents{part}[0pt]
  {\addvspace{16pt plus 2pt}\normalfont}
  {\trilogy@spacedsc{\partname\ \thecontentslabel}\hspace{1em}}
  {}
  {\hfill\contentspage}
  [\addvspace{4pt}]
\titlecontents{chapter}[2.2em]
  {\addvspace{7pt plus 1pt}\normalfont}
  {\contentslabel{2.2em}}
  {\hspace*{-2.2em}}
  {\hfill\contentspage}
\titlecontents{section}[4.8em]
  {\normalfont\small}
  {\contentslabel{2.6em}}
  {\hspace*{-2.6em}}
  {\titlerule*[0.9em]{.}\contentspage}

\@ifpackageloaded{natbib}{%
  \renewcommand{\bibsection}{\chapter*{References}%
    \addcontentsline{toc}{chapter}{References}}%
}{}

\newlength{\trilogyenvskip}
\setlength{\trilogyenvskip}{10pt plus 3pt minus 2pt}
\newtheoremstyle{plain}{\trilogyenvskip}{\trilogyenvskip}
  {\itshape}{}{\bfseries}{.}{0.5em}{}
\newtheoremstyle{definition}{\trilogyenvskip}{\trilogyenvskip}
  {\normalfont}{}{\bfseries}{.}{0.5em}{}
\newtheoremstyle{remark}{\trilogyenvskip}{\trilogyenvskip}
  {\normalfont}{}{\itshape}{.}{0.5em}{}

\usepackage[font=small,labelsep=period,skip=8pt]{caption}

\usepackage{etoolbox}
\patchcmd{\@footnotetext}{\footnotesize}{\small}{}{}

\usepackage[hidelinks,bookmarksnumbered=true]{hyperref}
\usepackage{bookmark}
\hypersetup{pdfauthor={},pdfpagemode=UseNone,pdfdisplaydoctitle=true,
  pdfstartview={FitH \hypercalcbp{\paperheight}}}

\makeatother
%
\usepackage{xurl}%
\usepackage[nameinlink,capitalize,noabbrev]{cleveref}
\usepackage{aliascnt}
\hypersetup{pdftitle={Aging in Spherical Spin Glass Dynamics},pdfauthor={P. M. Aronow and Patrick Lopatto}}
\renewcommand{\chaptername}{Section}
\numberwithin{equation}{chapter}

\theoremstyle{plain}
\newtheorem{theorem}{Theorem}[chapter]
\newaliascnt{proposition}{theorem}
\newtheorem{proposition}[proposition]{Proposition}
\aliascntresetthe{proposition}
\newaliascnt{lemma}{theorem}
\newtheorem{lemma}[lemma]{Lemma}
\aliascntresetthe{lemma}
\newaliascnt{corollary}{theorem}
\newtheorem{corollary}[corollary]{Corollary}
\aliascntresetthe{corollary}
\theoremstyle{definition}
\newaliascnt{definition}{theorem}
\newtheorem{definition}[definition]{Definition}
\aliascntresetthe{definition}
\newaliascnt{example}{theorem}
\newtheorem{example}[example]{Example}
\aliascntresetthe{example}
\theoremstyle{remark}
\newaliascnt{remark}{theorem}
\newtheorem{remark}[remark]{Remark}
\aliascntresetthe{remark}
\makeatletter
\@addtoreset{proposition}{chapter}\@addtoreset{lemma}{chapter}\@addtoreset{corollary}{chapter}
\@addtoreset{definition}{chapter}\@addtoreset{example}{chapter}\@addtoreset{remark}{chapter}
\makeatother
\crefname{theorem}{Theorem}{Theorems}
\crefname{proposition}{Proposition}{Propositions}
\crefname{lemma}{Lemma}{Lemmas}
\crefname{corollary}{Corollary}{Corollaries}
\crefname{definition}{Definition}{Definitions}
\crefname{example}{Example}{Examples}
\crefname{remark}{Remark}{Remarks}
\crefname{chapter}{Section}{Sections}
\crefname{part}{Part}{Parts}
\crefname{section}{Section}{Sections}
\crefname{subsection}{Section}{Sections}
\crefname{table}{Table}{Tables}
\crefname{equation}{}{}
\crefformat{equation}{(#2#1#3)}
\Crefformat{equation}{(#2#1#3)}
\crefrangeformat{equation}{(#3#1#4)--(#5#2#6)}
\Crefrangeformat{equation}{(#3#1#4)--(#5#2#6)}
\crefmultiformat{equation}{(#2#1#3)}{ and~(#2#1#3)}{, (#2#1#3)}{, and~(#2#1#3)}
\Crefmultiformat{equation}{(#2#1#3)}{ and~(#2#1#3)}{, (#2#1#3)}{, and~(#2#1#3)}
\newcommand{\creflastconjunction}{, and~}
\newcommand{\creflastgroupconjunction}{, and~}

\makeatletter
\newcommand{\npp@noteblock}[2]{%
  \ifx\relax#1\relax\else{\centering\trilogy@spacedsc{#1}\par}\vspace{6pt}\fi
  {\centering
   \begin{minipage}{0.84\textwidth}\small\normalfont\setlength{\parindent}{1.5em}%
     \noindent\ignorespaces#2\end{minipage}\par}}
\NewEnviron{authorsnote}{%
  \clearpage\thispagestyle{empty}%
  \vspace*{\stretch{2}}%
  \npp@noteblock{Authors' note}{\BODY}%
  \vspace*{\stretch{3}}%
  \clearpage}
\NewEnviron{abstractpage}{%
  \clearpage\thispagestyle{empty}%
  \vspace*{\stretch{2}}%
  \npp@noteblock{Abstract}{\BODY}%
  \vspace*{\stretch{3}}%
  \clearpage}
\newcommand{\npptitlepage}[2]{%
  \hypersetup{pageanchor=false}%
  \begin{titlepage}
  \centering\normalfont
  \vspace*{1.7in}
  {\large\strut\par}
  \vspace{1.05in}
  {\fontsize{28}{35}\selectfont #1\par}
  \vspace{1.1in}
  {\Large\def\\{\par\vspace{10pt}}#2\par}
  \vfill
  \end{titlepage}%
  \hypersetup{pageanchor=true}}
\makeatother

\newcommand{\E}{\mathbb E}
\newcommand{\1}{\mathbf 1}
\newcommand{\norm}[1]{\lVert#1\rVert}

\clubpenalty=10000
\patchcmd{\thebibliography}{\clubpenalty4000}{\clubpenalty10000}{}{}

\newcommand{\Prob}{\mathbb P}
\newcommand{\R}{\mathbb R}
\newcommand{\N}{\mathbb N}
\newcommand{\dd}{\mathop{}\!\mathrm d}
\newcommand{\ip}[2]{\left\langle#1,#2\right\rangle}
\newcommand{\op}{\mathrm{op}}
\newcommand{\HS}{\mathrm{HS}}
\newcommand{\tr}{\operatorname{tr}}
\newcommand{\Span}{\operatorname{span}}
\newcommand{\Id}{\operatorname{Id}}
\newcommand{\ALG}{\operatorname{ALG}}
\newcommand{\Eth}{E_{\mathrm{th}}}
\newcommand{\qd}{q_{\mathrm d}}
\newcommand{\bd}{\beta_{\mathrm d}}
\DeclareMathOperator{\diver}{div}


\begin{document}
\frontmatter
\npptitlepage{Aging in Spherical Spin Glass Dynamics}{P.\ M.\ Aronow\\Patrick Lopatto}
\begin{authorsnote}
With the exception of this note, this entire document is machine written. We have taken care to supervise this writing to ensure that the proofs meet a minimum standard of readability and that previous literature is appropriately cited. However, the resulting document remains lacking in exposition and overall coherence. We are nonetheless releasing this manuscript on Hexagon to disseminate the result as quickly as possible, because we believe it may be of interest to the mathematics and physics communities. We welcome suggestions concerning mathematical corrections, omitted citations, and attribution of ideas.

We thank Luca Leuzzi for many helpful comments on the theory of spherical spin glasses and his encouragement to explore their dynamics rigorously.

P.L.\ was partially supported by NSF grant DMS-2450004.

\end{authorsnote}
\begin{abstractpage}
We study Langevin dynamics for the pure spherical $p$-spin glass with $p\ge3$, with covariance $Nq^p$ and generator $\Delta_S+\beta\nabla_SH_N\cdot\nabla_S$, started from a uniform point independent of the disorder. We let the dimension tend to infinity on fixed time horizons and study the long-time behavior of the resulting hard-spherical Cugliandolo--Kurchan equations for the correlation $C$ and the response $R$. At every positive temperature, their solution has a unique stationary limit at bounded time differences, which satisfies the fluctuation--dissipation theorem (FDT), and its energy converges. For $\beta\ge\beta_{\mathrm d}$, where $\beta_{\mathrm d}$ is the dynamical transition, we show that the plateau $q$ of the stationary correlation is the larger solution of the marginality relation $\beta^2p(p-1)q^{p-2}(1-q)^2=1$, and we compute the limiting energy. Correlations with every fixed earlier state tend to zero. The limits of large waiting time and large time difference therefore do not commute, which proves aging of the correlation, including at the transition point.

We also prove that the response at a late time, as a function of the time difference, converges to its stationary profile in $L^2$ and uniformly over all earlier times; that FDT in the later time holds uniformly over all earlier times; and that the integral of the squared violation of FDT in the earlier time tends to zero, while for $\beta>\beta_{\mathrm d}$ the positive part $(\partial_sC-R)_+$ of this violation keeps a positive integrated mass. For $\beta\ge\beta_{\mathrm d}$, the stationary correlation approaches its plateau as a power law, with the exponent predicted by mode-coupling theory and a finite positive amplitude. For the cubic model and $\beta\ge\beta_{\mathrm d}$, including the transition point, we prove a lower bound on the total susceptibility and a product inequality for correlations at three times; for $\beta>\beta_{\mathrm d}$, we prove, for each fixed level below the plateau, an explicit asymptotic lower bound on the first passage time to that level multiplied by the entropy production during the passage. The aging proofs for the cubic and quartic models represent the correlation by vectors in a Hilbert space indexed by trees, and general pure interactions are treated with averages over finitely many independent futures and established algorithmic threshold bounds. The full fluctuation--dissipation relation of the aging regime and the time scale of aging remain open.

\end{abstractpage}
\tableofcontents
\mainmatter

\chapter{Introduction}\label{sec:introduction}

\section{Spin glass dynamics and aging}\label{sec:intro-background}

The static theory of mean-field spin glasses describes the Gibbs measure of a random Hamiltonian in high dimension. For the spherical $p$-spin model studied in this paper, the limiting free energy is given by the Crisanti--Sommers formula, a variational problem over probability measures on $[0,1]$ whose minimizer is the Parisi measure \cite{cs1992}. This formula was proved by Talagrand for even models \cite{talagrand2006} and by Chen in general \cite{chen2013}. For the pure model, the Parisi measure is a point mass at zero at high temperature, and it has one step of replica symmetry breaking at every temperature below the static transition \cite[Section~5]{auffingerchen2015}; the two-point support property follows from \cite[Proposition~2.2]{talagrand2006}. These results concern the equilibrium measure. In this paper, we study the dynamics.

The natural dynamics is the Langevin diffusion on the sphere, in which the configuration moves in the direction of increasing Hamiltonian, perturbed by Brownian noise; the inverse temperature sets the strength of the drift relative to the noise. This diffusion is reversible with respect to the Gibbs measure, and for each fixed dimension $N$, it converges to the Gibbs measure as time tends to infinity. At sufficiently low temperature, however, the spectral gap of its generator is exponentially small in $N$, with probability tending to one over the disorder \cite{gj2019}, so the slowest relaxation modes have relaxation times that grow exponentially with $N$. On time scales that do not grow with $N$, the dynamics can stay far from equilibrium. In this paper, we take the limit $N\to\infty$ first, on every fixed time interval, and only then study the long-time behavior of the limiting dynamics. In this order of limits, the Gibbs measure and its Parisi measure need not describe the behavior at long times. The physical theory predicts instead that the system ages.

Aging is observed in disordered materials that do not reach equilibrium on any accessible time scale. After a rapid cooling to low temperature, called a quench, the response of such a material to a small perturbation depends on the time that has elapsed since the quench, and it becomes slower as the sample gets older. In spin glasses, Lundgren, Svedlindh, Nordblad, and Beckman found that magnetic relaxation in CuMn depends on the time spent at low temperature before the measurement \cite{lundgren1983}; the experimental picture is reviewed in \cite{vincent1997}. In a trap model, the system jumps between traps and spends a random time in each of them. Bouchaud introduced weak ergodicity breaking through a trap model whose mean trapping time diverges: the system keeps aging, although the phase space does not split into mutually inaccessible regions \cite[p.~1707]{bouchaud1992}. Here the correlation measures the overlap of the configurations at two times, and the response measures the linear effect of a small force applied at the earlier time; we define both precisely in \cref{sec:intro-model}. In the formulation used for mean-field aging, the correlation with every fixed earlier configuration eventually vanishes, while relaxation becomes slower as the system gets older \cite[arXiv version~1, Section~3]{ck1994}.

Mean-field spin glasses with $p$-spin interactions are a setting in which aging can be derived from a microscopic dynamics. Derrida introduced the random energy model, in which the energies of different configurations are independent, as a limit of a family of $p$-spin models \cite{derrida1980}, and the Ising $p$-spin models themselves were studied by Gross and M\'ezard \cite{grossmezard1984} and by Gardner \cite{gardner1985}. Sompolinsky and Zippelius developed the dynamical mean-field theory of spin glass relaxation \cite{sz1981,sz1982}. This theory replaces the interaction of one coordinate with all the others by a random force and a memory of its own past, both determined self-consistently, and it describes the dynamics through closed equations for the correlation and response functions. Kirkpatrick and Thirumalai connected the dynamics of $p$-spin models with the glass transition of supercooled liquids, which are liquids cooled below their freezing temperature without crystallizing \cite{kt1987a,kt1987b}. For $p\ge3$, the dynamical equations resemble those of the mode-coupling theory of supercooled liquids, a self-consistent theory in which correlation functions satisfy nonlinear equations with memory and, at low temperature, develop a plateau, that is, a positive limit at large time differences (see \cite{goetzesjogren1992} for a review). In these equations, the correlation of the equilibrium dynamics stops decaying to zero below a dynamical transition temperature. This temperature is higher than the static transition temperature, at which the Parisi measure changes. Kirkpatrick, Thirumalai, and Wolynes built on this correspondence a scaling theory of glasses in finite dimension, in which relaxation below the dynamical transition proceeds by thermally activated rearrangements, where the noise carries the system across barriers of the energy landscape \cite{ktw1989}. The spherical $p$-spin model was solved at equilibrium by Crisanti and Sommers \cite{cs1992}. Crisanti, Horner, and Sommers studied its dynamics in \cite{chs1993}, including the dynamical transition above the static one and the power-law relaxation near it. Cugliandolo and Kurchan constructed an asymptotic solution of the nonequilibrium equations at long times \cite{ck1993}. We describe their predictions in \cref{sec:intro-picture}, after we introduce the model and the equations that they analyzed.

\section{The model and the Cugliandolo--Kurchan equations}\label{sec:intro-model}

We fix an integer $p\ge3$ and $\beta>0$, and write $\xi(q)=q^p$. For $N\ge2$, let $S_N=\{x\in\R^N:\norm{x}^2=N\}$. The Hamiltonian $H_N$ is a centered Gaussian process on $\R^N$ with covariance
\[
 \E H_N(x)H_N(y)=N\xi\Bigl(\frac{x\cdot y}{N}\Bigr),
\]
realized as in \eqref{eq:hamiltonian}. This is the pure spherical $p$-spin model in the normalization of the rigorous static theory. Its Gibbs measure at inverse temperature $\beta$ has density proportional to $e^{\beta H_N}$ with respect to the uniform probability measure on $S_N$. Let $X_t$ be the Langevin dynamics on $S_N$ with generator
\[
 \Delta_S+\beta\nabla_SH_N\cdot\nabla_S,
\]
where $\Delta_S$ and $\nabla_S$ are the Laplace--Beltrami operator and gradient on $S_N$; equivalently, $X_t$ solves \eqref{eq:hard-sde}, whose noise has amplitude $\sqrt2$. The Gibbs measure is reversible for this diffusion. The initial point $X_0$ is uniform on $S_N$ and independent of the disorder and of the Brownian motion. Since the uniform measure is the Gibbs measure at $\beta=0$, this initial condition models a quench from infinite temperature. The drift points in the direction of increasing $H_N$, so the limiting energies and the threshold energy in this paper are positive; physics works usually minimize the Hamiltonian, which reverses these signs.

Since the uniform initial law is not invariant for the dynamics, its natural observables depend on two times. The correlation
\[
 C_N(t,s)=\frac{X_t\cdot X_s}N
\]
is the overlap of the configurations at times $s$ and $t$. The response describes how the configuration at a time $t$ reacts to a small force applied at an earlier time $s$. Precisely, for a deterministic force $h(u)$ that is bounded and continuous on each fixed time horizon and projected onto the tangent space of $S_N$, the response is the zero-initial strong solution $V_t$ of the forced linearized equation along $X$, with the drift and noise derivatives evaluated on $X$. It satisfies $V_t=\int_0^tU_N(t,u)h(u)\,\dd u$ for the random response matrix $U_N(t,u)$ (\cref{sec:langevin}). We write
\[
 \bar r_N(t,s)=\frac1N\E\tr U_N(t,s),\qquad 0\le s\le t,
\]
for its normalized trace, averaged over the disorder, the initial point, and the noise. Thus $\bar r_N(t,s)$ is the mean linear response of a coordinate at time $t$ to a unit impulse of force on the same coordinate at time $s$, that is, to a force concentrated near time $s$ with unit time integral. We measure the response to a force added to the drift without a factor $\beta$. The finite-dimensional response has diagonal value $\bar r_N(t,t)=\kappa_N=(N-1)/N$, because the force is projected onto the $(N-1)$-dimensional tangent space, and its limit below has diagonal value one (\cref{rem:fdt-convention}).

As $N\to\infty$, these functions have deterministic limits on fixed time horizons. For every fixed $0\le T<\infty$, the correlation $C_N(t,s)$ converges in probability to a function $C(t,s)$, uniformly on $[0,T]^2$, and the averaged response $\bar r_N(t,s)$ converges to a function $R(t,s)$, uniformly on $0\le s\le t\le T$. Together with a function $\lambda(t)$, which keeps $C(t,t)=1$, the limits form the unique solution of the system of integro-differential equations
\begin{align*}
 \partial_tR(t,s)&=-\lambda(t)R(t,s)
 +\beta^2\int_s^t\xi''(C(t,u))R(t,u)R(u,s)\,\dd u,\\
 \partial_tC(t,s)&=-\lambda(t)C(t,s)
 +\beta^2\int_0^t\xi''(C(t,u))R(t,u)C(u,s)\,\dd u\\
 &\hspace{1em}+\beta^2\int_0^s\xi'(C(t,u))R(s,u)\,\dd u,\\
 \lambda(t)&=1+\beta^2\int_0^t
 \bigl[\xi'(C(t,u))+C(t,u)\xi''(C(t,u))\bigr]R(t,u)\,\dd u,
\end{align*}
for $t\ge s\ge0$, with $C(t,t)=R(t,t)=1$, where $C$ is extended to $s>t$ by symmetry. We call these the Cugliandolo--Kurchan (CK) equations; they are stated with their conventions in \eqref{eq:ck-response}--\eqref{eq:ck-boundary}. Existence and uniqueness of the solution were proved by Dembo, Guionnet, and Mazza \cite[Proposition~1.1]{dgm2007}, and the solution satisfies $0\le C\le1$ and $R\ge0$ (see \cref{prop:ck-wellposed}). Sellke proved fixed-horizon convergence of $C_N$ and of the Brownian observable representing the integrated response \cite[Theorem~4.3]{sellke2024}. The force-response identification in \cref{lem:brown-force-response} equates the annealed mean of its rescaled version $\Gamma_N(t,s)=X_t\cdot B_s/(\sqrt2\,N)$ with $\int_0^s\bar r_N(t,u)\,\dd u$, where $B$ drives the diffusion. He built on the theorem of Ben Arous, Dembo, and Guionnet for the soft-spherical dynamics \cite[Theorem~1.2]{bag2006}, in which the constraint $\norm{X_t}^2=N$ is replaced by a confining potential. We state these results in our normalization in \cref{prop:hard-bridge,lem:brown-force-response}, and we deduce the uniform convergence of $\bar r_N$ in \cref{rem:no-pointwise-response}. Throughout the paper, the \emph{CK solution} is this solution $(C,R,\lambda)$, which describes the Langevin dynamics started from the uniform point.

The CK equations involve only $C$, $R$, and $\lambda$; the disorder no longer appears. The function $\lambda$ is a Lagrange multiplier for the spherical constraint, and the integrals are memory terms. In the physical derivation through dynamical mean-field theory, each coordinate of $X_t$ behaves in the limit like the solution $x_t$ of the linear equation (see \cite[arXiv version~2, Section~9.5]{cugliandolo2004})
\[
 \dd x_t=\Bigl(-\lambda(t)x_t+\beta^2\int_0^t\xi''(C(t,u))R(t,u)x_u\,\dd u+\eta_t\Bigr)\dd t+\sqrt2\,\dd b_t.
\]
Here $x_0$, $b$, and $\eta$ are independent; $x_0$ is a standard Gaussian variable, $b$ is a standard Brownian motion, and $\eta$ is a centered Gaussian process with covariance $\E\eta_t\eta_u=\beta^2\xi'(C(t,u))$. The interaction of one coordinate with all the others produces the memory kernel $\beta^2\xi''(C(t,u))R(t,u)$, which weights the effect of the past values of $x$ on its present drift, and the noise $\eta$, which is called colored because its values at different times are correlated. The CK equations express the self-consistency conditions that $C(t,s)=\E x_tx_s$, that $R(t,s)$ is the response of $x_t$ to a force applied at time $s$, and that $\E x_t^2=1$, which determines $\lambda$. We use this effective process only as a guide. Its variation-of-constants formula writes $C(t,s)$ as the sum of the contributions of the initial condition, of the Brownian noise, and of the colored noise, each propagated by the response. This is the Green covariance identity of \cref{prop:green-covariance}, which we prove directly from the CK equations.

We write
\[
 e_\beta(t)=\beta p\int_0^tC(t,u)^{p-1}R(t,u)\,\dd u,\qquad
 \chi(t,s)=\int_s^tR(t,u)\,\dd u.
\]
The first quantity is the limit of the mean energy $N^{-1}\E H_N(X_t)$, and $\lambda(t)=1+p\beta e_\beta(t)$ (\cref{prop:energy-bridge}). The second is the integrated response, or susceptibility. It is the response at time $t$ to a constant unit force switched on at time $s$. A field $h$ added to the Hamiltonian adds the force $\beta h$ to the drift, so the magnetization per unit field is $\beta\chi(t,s)$ in our conventions. It is the analog of the magnetization measured in aging experiments after a small field is switched on at time $s$, the zero-field-cooled protocol \cite{vincent1997}.

We also use the entropy production of the dynamics. Fix the disorder, and let $\rho_t$ be the law of $X_t$. Since the dynamics is reversible, the relative entropy of $\rho_t$ with respect to the Gibbs measure decreases in time. Its rate of decrease is the relative Fisher information, the integral against $\rho_t$ of the squared norm of the gradient of the logarithm of the density of $\rho_t$ with respect to the Gibbs measure. We call this rate the entropy production at time $t$. The initial relative entropy is at most $\beta^2N/2$ on average over the disorder, so the total entropy production per coordinate is at most $\beta^2/2$ on average (\cref{lem:entropy-budget}). We write $a_N(t)$ for the entropy production at time $t$, divided by $N$ and averaged over the disorder. Following the usual terminology for disorder averages, we also call $a_N(t)$ the normalized annealed Fisher information (\cref{sec:entropy-clusters}).

Other works use other normalizations. Cugliandolo and Kurchan \cite[Equations~(1)--(2)]{ck1993} minimize a Hamiltonian whose couplings are indexed by distinct unordered $p$-tuples and have variance $p!/(2N^{p-1})$. Their limiting covariance convention is $N\xi(x\cdot y/N)/2$, with an inverse temperature $\beta^{\mathrm{CK}}$ in front of the drift and noise with variance rate $2$. At the level of the limiting equations, their convention agrees with ours when $\beta^{\mathrm{CK}}=\sqrt2\,\beta$, with no change of time. In the physical normalization of time, where the drift is the gradient itself and the noise has variance rate $2\theta$ at temperature $\theta$, a rescaling of time gives $\beta^{\mathrm{CK}}=1/\theta$, and hence $\beta=1/(\sqrt2\,\theta)$. The mathematical works \cite{bag2006,dgm2007,sellke2024} use Brownian noise with variance rate one; their drift parameter is $\alpha=\beta/2$, and their time is $v=2t$. \Cref{tab:conventions} in \cref{sec:conventions} collects the dictionary. All formulas in this paper are written in our conventions, and the formulas that we quote from \cite{ck1993,cs1995,dgm2007,sellke2024} have been translated.

\section{The physical picture}\label{sec:intro-picture}

We first recall what happens in equilibrium. Suppose that the diffusion is started from the Gibbs measure instead of the uniform measure. Then the process is stationary, and its correlation and response depend only on the time difference $\tau=t-s$, which we also call the lag: they have the form $c(\tau)$ and $r(\tau)$. They satisfy the fluctuation--dissipation theorem (FDT),
\[
 r(\tau)=-c'(\tau),\qquad \tau>0,
\]
which expresses the response to a small force through the spontaneous fluctuations of the system. FDT is a consequence of reversibility (\cref{rem:fdt-convention}); see Kubo \cite{kubo1957} for the general linear-response theory. Agarwal \cite{agarwal1972} treats response around nonthermal steady states under detailed balance; nonequilibrium diffusion response identities are discussed in \cite{bmw2009}. In our conventions it carries no factor $\beta$. When $C(t,s)=c(t-s)$, it reads $R(t,s)=\partial_sC(t,s)$. Reversibility also gives a spectral representation. By the spectral theorem for the generator, the equilibrium correlation is a mixture of decaying exponentials, $c(\tau)=\int_{[0,\infty)}e^{-x\tau}\,\nu(\dd x)$ for a probability measure $\nu$. By Bernstein's theorem \cite{widder1941}, a continuous function $c$ on $[0,\infty)$ with $c(0)=1$ has this form if and only if it is \emph{completely monotone} on $(0,\infty)$, that is, infinitely differentiable there with
\[
 (-1)^nc^{(n)}(\tau)\ge0,\qquad n\ge0,\ \tau>0.
\]
In particular, a completely monotone function is nonnegative, nonincreasing, and convex. We use the terms FDT and complete monotonicity in this sense throughout.

At high temperature, the dynamics started from the uniform point approaches this equilibrium behavior. Dembo, Guionnet, and Mazza proved that for $\beta$ small enough, the CK solution relaxes exponentially fast and converges to a stationary solution that satisfies FDT \cite[Proposition~1.2 and Theorem~1.3]{dgm2007}.

Below the dynamical transition, the picture is different. Let $s$ be the waiting time, which is the age of the system when a measurement starts, and let $\tau=t-s$ be the lag. Cugliandolo and Kurchan constructed an asymptotic solution of the CK equations at long times \cite{ck1993}. It has the following features; the equation numbers refer to the arXiv version of \cite{ck1993}.
\begin{itemize}
\item \emph{A stationary regime.} When $s\to\infty$ at fixed lag $\tau$, we have $C(s+\tau,s)\to c(\tau)$ and $R(s+\tau,s)\to r(\tau)$, and the limits satisfy FDT. On bounded windows of late times, the correlation and the response behave as in equilibrium.
\item \emph{A positive plateau.} The stationary correlation does not decay to zero. It decreases to a positive limit $q=\lim_{\tau\to\infty}c(\tau)$, called the plateau; in the physics literature, $q$ is also called the dynamical Edwards--Anderson parameter.
\item \emph{Weak ergodicity breaking and an aging regime.} Below the plateau, the correlation decays on time scales that grow with the waiting time $s$, and $C(t,s)\to0$ as $t\to\infty$ for every fixed $s$. The aging regime consists of pairs of late times $s<t$ with $t-s\to\infty$ at which $C(t,s)$ stays below the plateau.
\item \emph{Marginality.} The plateau is fixed by a marginality condition \cite[Equation~(7)]{ck1993}, which says that a stability eigenvalue vanishes; we explain its meaning below. In our conventions, it reads $\beta^2p(p-1)q^{p-2}(1-q)^2=1$.
\item \emph{A constant fluctuation--dissipation ratio.} In the aging regime, FDT fails, and response and correlation are related instead by $R(t,s)\approx x\,\partial_sC(t,s)$, with a constant $x=(p-2)(1-q)/q$ given in the paragraph after \cite[Equation~(12)]{ck1993}; the corresponding response moments are \cite[Equation~(13)]{ck1993}. Equilibrium FDT corresponds to $x=1$, and below the dynamical transition $x<1$. Cugliandolo, Kurchan, and Peliti later interpreted $x$ as the ratio of the temperature of the noise to an effective temperature of the slow motion, which is therefore higher \cite{ckp1997}.
\item \emph{The energy.} The limit of the energy $e_\beta(t)$ as $t\to\infty$ is given by \cite[Equation~(9)]{ck1993}, combined with the relation between response and correlation in the aging regime; in our conventions, it is $\beta[1+(p-2)q^{p-1}-(p-1)q^p]$.
\end{itemize}
The first three features imply that the two iterated limits of the correlation differ:
\[
 \lim_{\tau\to\infty}\lim_{s\to\infty}C(s+\tau,s)=q>0,
 \qquad
 \lim_{s\to\infty}\lim_{\tau\to\infty}C(s+\tau,s)=0.
\]
Taking the limit of large waiting time first describes the relaxation of an arbitrarily old system over a bounded time window. Taking the limit of large lag first asks whether the configuration at a fixed time is eventually forgotten. In equilibrium, both limits equal $\lim_{\tau\to\infty}c(\tau)$. In this model, their difference is a mathematical signature of aging: at large waiting times, the correlation first approaches the plateau, as it would in equilibrium, and it then decays further on a time scale that diverges with the waiting time. Thus stationarity on bounded windows does not describe the relaxation over the whole history. The noncommutation of the limits detects this slower relaxation without determining its time scale.

The dynamical transition is the temperature at which $x$ reaches one \cite{ck1993}. In our conventions, the marginality condition and $x=1$ give the transition parameters
\[
 \qd=\frac{p-2}{p-1},\qquad \bd^2=\frac{(p-1)^{p-1}}{p(p-2)^{p-2}}.
\]
For the cubic model, $\qd=1/2$ and $\bd=\sqrt{4/3}$. For $\beta<\bd$, the physical picture predicts a zero plateau \cite{ck1993,chs1993}. Equivalently, $\bd$ is the smallest $\beta$ for which the equation $\beta^2pq^{p-2}(1-q)=1$ has a solution $q\in(0,1)$, and this solution is $\qd$ (\cref{sec:conventions}). In the physical picture, this is the equation for the plateau of the equilibrium dynamics in the replica symmetric phase, so the equilibrium dynamics stops decorrelating at the inverse temperature $\bd$ \cite{chs1993}.

The picture can be compared with the static theory. Let $\beta_{\mathrm s}$ be the static transition, so that the Parisi measure is the point mass at zero for $\beta\le\beta_{\mathrm s}$ and has two atoms for $\beta>\beta_{\mathrm s}$ \cite{auffingerchen2015}. Then $\bd<\beta_{\mathrm s}$; for $p=3$, we have $\bd\approx1.155$ and $\beta_{\mathrm s}\approx1.207$ (see \cite[Theorem~1]{subag2023} for the equations that determine $\beta_{\mathrm s}$ in our normalization). Since $\xi''(q)=p(p-1)q^{p-2}$, the marginality condition reads $\beta^2\xi''(q)(1-q)^2=1$. Consider a trial measure with one step of replica symmetry breaking, $\theta\delta_0+(1-\theta)\delta_q$ with $\theta\in(0,1)$, in the Crisanti--Sommers functional. In the replica calculation, its stability against perturbations of the overlaps within a block is measured by the replicon eigenvalue, which is a positive multiple of $1-\beta^2\xi''(q)(1-q)^2$. Thus the marginality condition says that the replicon eigenvalue vanishes, so that the solution is marginally stable. If, in addition, the derivative of the functional in $q$ vanishes, then $\theta=x$. For $\beta>\beta_{\mathrm s}$, the Parisi measure is obtained instead by optimizing the functional over both $\theta$ and $q$, and it is strictly stable. Barrat, Franz, and Parisi identified the asymptotic dynamical condition with this static marginal stability \cite[Section~IV, Equation~(25)]{bfp1997}; see also \cite[arXiv version~2, Section~15.1]{cugliandolo2004}.

In this sense, the plateau and the energy of the dynamics are those of marginally stable states rather than those of the equilibrium states that the Parisi measure describes, and the fluctuation--dissipation ratio $x$ is the weight $\theta$ of the corresponding marginal solution. Cugliandolo and Kurchan noted that, in this model, the overlap distribution defined by the dynamics differs from the static one \cite[the paragraph after Equation~(15)]{ck1993}. Franz, M\'ezard, Parisi, and Peliti related the dynamical response to the static Parisi measure under additional hypotheses, and they pointed out that the pure spherical $p$-spin model does not satisfy this relation \cite{fmpp1998}. For $\bd<\beta\le\beta_{\mathrm s}$, the Parisi measure carries no positive overlap, while the predicted dynamical plateau is positive.

The limiting energy also has a landscape interpretation, developed alongside the calculation of Cugliandolo and Kurchan. The Thouless--Anderson--Palmer (TAP) free energy is a function of the local magnetizations whose stable local minima represent candidate metastable states. Crisanti and Sommers counted the solutions of the TAP equations of the spherical model and computed the threshold energy at which these solutions become marginally stable \cite[Equations~(8)--(9)]{cs1995}. At this threshold, the lower edge of the limiting spectrum of the Hessian of the TAP free energy reaches zero \cite[arXiv version~1, p.~4]{ck1993}; see also \cite[arXiv version~2, Section~15.4]{cugliandolo2004}. The energy of these threshold states depends on the temperature, and it agrees with the limiting energy of Cugliandolo and Kurchan \cite[Equation~(14)]{cs1995}. As $\beta\to\infty$, it tends to the threshold energy $\Eth=2\sqrt{(p-1)/p}$ of the landscape, which is strictly smaller than the ground-state energy $\lim_{N\to\infty}N^{-1}\E\max_{S_N}H_N$; for $p=3$, we have $\Eth\approx1.633$, while the ground-state energy is approximately $1.657$ \cite{chensen2017}. Kurchan and Laloux related slow dynamics to the geometry of the landscape near the threshold \cite[Section~V]{kurchanlaloux1996}. Related analyses of metastable states, barriers, and stationary points of the TAP free energy are in \cite{kpv1993,cgp1998}.

Cugliandolo and Kurchan developed the structure of the aging regime further in \cite{ck1994}, where they studied the Sherrington--Kirkpatrick model. They distinguished weak ergodicity breaking from weak long-term memory, which means that the integrated response to a force applied during any fixed time interval $[0,S]$ vanishes as $t\to\infty$ \cite[arXiv version~1, Section~3]{ck1994}. They also observed that the asymptotic equations of the aging regime do not change when time is replaced by an increasing function of time and the response is transformed accordingly, so that these equations do not determine the time scale of aging \cite[arXiv version~1, Section~4]{ck1994}. Further, they studied the relations between correlations at three widely separated times \cite[arXiv version~1, Section~5]{ck1994}. Two important examples of such relations are the dynamical ultrametric relation $C(t,u)=\min\{C(t,s),C(s,u)\}$ for $u<s<t$, analogous to the ultrametric organization of equilibrium overlaps, and a product relation; for spherical $p$-spin models, they discussed a product relation \cite[arXiv version~1, Section~5.1, Equation~(5.1.8)]{ck1994}. Reviews of the two regimes and their response are given in \cite{bckm1998,cugliandolo2004}. The idealized mode-coupling scenario of the glass transition was developed by Bengtzelius, G\"otze, and Sj\"olander and by Leutheusser \cite{bgs1984,leutheusser1984}. Mode-coupling theory also describes how a critical or marginal stationary correlation $c(\tau)$ approaches its plateau: $c(\tau)-q$ decays like $\tau^{-a}$, where $a\in(0,1/2)$ solves $\Gamma(1-a)^2/\Gamma(1-2a)=\lambda_{\mathrm{MCT}}$, $\Gamma$ is Euler's gamma function, and the constant $\lambda_{\mathrm{MCT}}$, called the exponent parameter, is not to be confused with the multiplier $\lambda(t)$. The correspondence between disordered models and mode-coupling theory was worked out by Bouchaud, Cugliandolo, Kurchan, and M\'ezard in \cite{bckm1996}. They gave this exponent relation at equilibrium \cite[arXiv version, Equation~(IV.3)]{bckm1996} and for the stationary profile of the aging solution \cite[arXiv version, Equation~(IV.6)]{bckm1996}, where they attribute the calculation to Cugliandolo and Le Doussal \cite{cugliandololedoussal1996}. For the pure spherical model, the exponent parameter of the aging solution is $x/2=(p-2)(1-q)/(2q)$. This relation is written explicitly by Kim and Latz \cite[the paragraph after Equation~(12)]{kimlatz2001}; see also \cite[arXiv version~2, Equation~(12.36)]{cugliandolo2004}.

The order of limits is essential in this picture. As explained in \cref{sec:intro-background}, the dynamics converges to its Gibbs measure for each finite $N$. In the CK description, the dimension tends to infinity first, on every fixed time interval, and only then do the observation and waiting times tend to infinity. Aging in this sense is a property of a deterministic pair of functions $(C,R)$, and it is not a statement about time scales that grow with $N$. On such time scales, the dynamics can be activated, and the relevant mechanisms are different; we discuss the rigorous results on activated aging in \cref{sec:intro-related}. Neither a theorem about finite-dimensional mixing nor a formal solution of the asymptotic equations determines the behavior of the CK solution in the order of limits used here.

\section{Main results}\label{sec:main-results}

The goal of this paper is to prove the main features of the picture of \cref{sec:intro-picture} for pure $p$-spin models with $p\ge3$ started from a uniform point. The physical calculation imposes a long-time form on the solution and solves reduced asymptotic equations. We prove its conclusions for the CK solution itself, which is the limit of the Langevin dynamics, without assuming that form. At every positive temperature, we prove that the CK solution becomes stationary on bounded windows of lags and satisfies FDT there, that its energy converges, and that its correlation and response with every fixed earlier time tend to zero. For $\beta\ge\bd$, that is, at and below the dynamical transition temperature, we prove that the plateau is marginal and that the limiting energy is the one predicted by Cugliandolo and Kurchan. The two iterated limits of the correlation then differ, which is the sense in which the solution ages. We also prove that the response at a late time to forces applied at all earlier times concentrates at short lags in the mean-square sense, that the violation of FDT vanishes in a mean-square sense, and that, for $\beta>\bd$, its positive part keeps a positive integrated mass. For $\beta\ge\bd$, we prove that the stationary correlation approaches its plateau as a power law with a positive and finite amplitude. For the cubic model, we prove a lower bound on the susceptibility and an inequality between correlations at three times for $\beta\ge\bd$, and, for $\beta>\bd$, a lower bound that relates the time needed to relax below the plateau to the entropy production. We do not identify the time scale of aging, and we do not prove the full fluctuation--dissipation relation of the aging regime.

Let $(C,R,\lambda)$ be the CK solution of \cref{sec:intro-model}. All long-time limits in the following statements concern this deterministic solution. They describe the Langevin dynamics only in the order of limits described in \cref{sec:intro-picture}: first $N\to\infty$ on a fixed time horizon, through \cref{prop:hard-bridge}, and then the long-time limits.

\begin{theorem}[Stationary limit, marginal plateau, and aging]\label{thm:main-aging}
Fix $p\ge3$ and $\beta>0$. There exist functions $c_\beta,r_\beta\colon[0,\infty)\to\R$ and numbers $e_\beta(\infty)$ and $\ell_\beta$ such that $e_\beta(T)\to e_\beta(\infty)$ as $T\to\infty$ and
\begin{equation}
 C(T+t,T+s)\longrightarrow c_\beta(t-s),\qquad
 R(T+t,T+s)\longrightarrow r_\beta(t-s),\qquad
 \lambda(T)\longrightarrow\ell_\beta,
 \label{eq:main-stationary}
\end{equation}
where the first two limits are locally uniform in $t\ge s\ge0$. The function $c_\beta$ is completely monotone on $(0,\infty)$, $r_\beta=-c_\beta'$, and $c_\beta(0)=r_\beta(0)=1$. For every finite $S\ge0$,
\begin{equation}
 \lim_{t\to\infty}\sup_{0\le s\le S}\bigl\{C(t,s)+R(t,s)\bigr\}=0.
 \label{eq:main-forgetting}
\end{equation}
If $\beta\ge\bd$, then the plateau $q_\beta=\lim_{\tau\to\infty}c_\beta(\tau)$ is the larger of the two solutions $q\in(0,1)$ of
\begin{equation}
 \beta^2p(p-1)q^{p-2}(1-q)^2=1,
 \label{eq:main-marginality}
\end{equation}
and $q_\beta\ge\qd$. In this range,
\begin{align}
 e_\beta(\infty)
 &=\beta\bigl[1+(p-2)q_\beta^{p-1}-(p-1)q_\beta^p\bigr],
 \label{eq:main-energy}\\
 \ell_\beta
 &=\frac1{1-q_\beta}+\beta^2p\bigl(1-q_\beta^{p-1}\bigr),
 \label{eq:main-multiplier}
\end{align}
and the correlation limits do not commute:
\begin{equation}
 \lim_{\tau\to\infty}\lim_{s\to\infty}C(s+\tau,s)=q_\beta>0,
 \qquad
 \lim_{s\to\infty}\lim_{\tau\to\infty}C(s+\tau,s)=0.
 \label{eq:main-aging}
\end{equation}
\end{theorem}

\Cref{thm:main-aging} has two parts. The first holds at every temperature. It says that the CK solution becomes stationary at bounded lags: near every late time $T$, the correlation and the response are close to functions of the lag alone, these functions satisfy FDT, and the correlation has the spectral representation of an equilibrium correlation. Thus, on bounded windows of late times, the correlation and the response have the stationary FDT structure of an equilibrium dynamics, even at low temperature, where the dynamics keeps aging. The limits in \eqref{eq:main-stationary} are taken along the whole family $T\to\infty$, so the stationary profile does not depend on a subsequence. By \eqref{eq:main-forgetting}, the solution also forgets every fixed earlier time: its correlation with the configuration at a fixed time $s$, and its response to a force applied at time $s$, tend to zero.

The second part concerns $\beta\ge\bd$, and it identifies the plateau and the energy with those of the asymptotic solution of Cugliandolo and Kurchan. The relation \eqref{eq:main-marginality} is their marginality condition \cite[Equation~(7)]{ck1993}, and \eqref{eq:main-energy} is their energy \cite[Equation~(9)]{ck1993}, combined with the relation between response and correlation in the aging regime stated after \cite[Equation~(12)]{ck1993}; both are written here in the conventions of \cref{sec:conventions}. Since the plateau is positive while the correlation with every fixed earlier time vanishes, the two iterated limits in \eqref{eq:main-aging} differ. We call this noncommutation \emph{correlation aging}. It does not identify the time scale on which $C(s+\tau,s)$ decreases below $q_\beta$. In physical terms, \cref{thm:main-aging} shows that, for $\beta\ge\bd$, the dynamics started from a uniform point does not reach equilibrium: its plateau is the marginal one, and its limiting energy is the energy predicted for the threshold states. For example, if $\bd<\beta<\beta_{\mathrm s}$, the Parisi measure is the point mass at zero and the equilibrium energy per coordinate is $\beta$, while by \eqref{eq:main-energy}, we have $e_\beta(\infty)=\beta-\beta q_\beta^{p-1}[(p-1)q_\beta-(p-2)]<\beta$, since $q_\beta>\qd$.

The theorem includes the transition point $\beta=\bd$, where $q_\beta=\qd$ and $e_\beta(\infty)=\bd$. The value $e_{\bd}(\infty)=\bd$ coincides with the value $e_\beta(\infty)=\beta$ obtained whenever the plateau is zero (\cref{sel:intermediate-range}). Letting $\beta\to\infty$ in \eqref{eq:main-energy}, we find that $e_\beta(\infty)\to\Eth=2\sqrt{(p-1)/p}$ (\cref{rem:low-temperature}). This limit is the threshold energy of the pure $p$-spin landscape \cite[Equation~(2.13)]{abc2013}, and it agrees with the low-temperature results of Sellke for Langevin dynamics on bounded time horizons \cite[Theorems~1.1 and~1.7]{sellke2024}. The limit $\beta\to\infty$ is taken after all other limits, and it says nothing about gradient flow at zero temperature.

We prove \cref{thm:main-aging} at the end of \cref{sec:aging}, by combining \cref{sel:all-temperature,sel:selected-sector,sel:fixed-wait-forgetting,sel:aging} with \cref{prop:stationary-clusters,sel:cluster-complete-monotonicity}. The main difficulty is that the correlation equation involves the whole history: on late windows of bounded length, its memory integrals involve the overlap of late states with the remote past, which a limit on such windows does not control. The selection of the plateau and the forgetting of earlier times rest on two facts, which we state as \cref{prop:threshold-inputs} in \cref{sec:threshold-inputs}. To state them, we write $C(t,s)=\ip{v_t}{v_s}$ for unit vectors $v_t$ in a deterministic Hilbert space whose components are indexed by trees, and we project $v_t$ onto the components whose leaves record the initial condition or noise up to an earlier time, while their internal interaction times may be later (\cref{sec:geometry-core}). The representation comes with an energy functional, the tree energy, which reproduces $e_\beta(t)$ on $v_t$. The first fact says that, for every stationary limit with positive plateau, the tree energy of these projections, after normalization, does not exceed the threshold energy $\Eth$ in the limit. The second fact says that if the whole translated family converges to one stationary profile whose plateau is positive and marginal, then the solution forgets every fixed earlier time. We prove that proposition for $p\in\{3,4\}$ in \cref{sec:tree-ceilings}, using a deterministic inequality in this Hilbert space, and for every $p\ge3$ in \cref{app:general-selection}, using the algorithmic threshold of Huang and Sellke \cite[Theorem~1 and Corollary~1.8]{hs2023}.

To describe what happens for $\beta<\bd$, we use the following elementary facts about the marginality equation. Let
\[
 q_{\mathrm m}=\frac{p-2}{p},\qquad
 \beta_{\mathrm m}^2=\frac{p^{p-1}}{4(p-1)(p-2)^{p-2}}.
\]
The function $q\mapsto p(p-1)q^{p-2}(1-q)^2$ attains its maximum $\beta_{\mathrm m}^{-2}$ on $[0,1]$ at $q_{\mathrm m}$, so the marginality equation $\beta^2p(p-1)q^{p-2}(1-q)^2=1$ has a solution in $(0,1)$ if and only if $\beta\ge\beta_{\mathrm m}$; the value $\beta_{\mathrm m}$ is not a physical transition. For $\beta\ge\beta_{\mathrm m}$, we write $q_+(\beta)$ for its solution in $[q_{\mathrm m},1)$. We have $\beta_{\mathrm m}<\bd$, $q_{\mathrm m}<\qd$, and $q_+(\bd)=\qd$, and for $\beta\ge\bd$ we write $q_\beta=q_+(\beta)$; see \cref{sec:marginality,sec:selected-sector}.

For $\beta<\bd$, \cref{thm:main-aging} does not determine the plateau on the whole interval. For $\beta<\beta_{\mathrm m}$, the plateau is zero. For $\beta_{\mathrm m}\le\beta<\bd$, the proof shows that it is either zero or $q_+(\beta)$, and it proves convergence in both cases (\cref{sel:all-temperature,sel:intermediate-range}). The physical picture predicts the zero plateau there \cite{ck1993,chs1993}. We return to this question in \cref{sec:discussion-below}.

In the remaining statements, when $p$ and $\beta$ are fixed, we abbreviate $c=c_\beta$, $r=r_\beta$, $q=\lim_{\tau\to\infty}c_\beta(\tau)$, and $\ell=\ell_\beta$. All these statements use the stationary profile selected in \cref{thm:main-aging}. The only exception in notation is \cref{thm:main-fisher-cost}, where $c$ denotes a level in $(0,q)$, as in \cref{sec:three-time,resp:fisher-chapter}.

The stationary limit in \cref{thm:main-aging} describes the response $R(t,s)$ at bounded lags $t-s$, and the forgetting statement describes it at bounded earlier times $s$. The next theorem concerns the response at a late time $t$ to forces applied at all earlier times $s\in[0,t]$ simultaneously. We write it as a function of the lag $v=t-s\in[0,t]$, so that $R(t,t-v)$ is the response at time $t$ to an impulse at time $t-v$. We also write
\[
 E(t,s)=\partial_sC(t,s)-R(t,s),\qquad t>s>0,
\]
for the violation of FDT in the earlier time $s$, and $z_+=\max\{z,0\}$ for $z\in\R$. In equilibrium, $E=0$. For $\beta>\bd$, the physical picture predicts $E(t,s)\approx(1-x)\,\partial_sC(t,s)\ge0$ in the aging regime, and this quantity is expected to be small there, because the correlation changes on a diverging time scale.

\begin{theorem}[Response at late times and the violation of FDT]\label{thm:main-response}
Fix $p\ge3$ and $\beta>0$. Then, as $t\to\infty$,
\begin{equation}
\begin{aligned}
 \mathbf1_{[0,t]}(v)R(t,t-v)&\longrightarrow r(v),\\
 \mathbf1_{[0,t]}(v)\,\partial_sC(t,s)\big|_{s=t-v}&\longrightarrow r(v)
\end{aligned}
 \qquad\text{strongly in }L^2(0,\infty),
 \label{eq:main-strong-rows}
\end{equation}
where the functions on the left are set equal to zero for $v>t$. Moreover,
\begin{equation}
 \lim_{t\to\infty}\sup_{0\le s\le t}|R(t,s)-r(t-s)|=0,
 \qquad \chi(t,0)=o(\sqrt t).
 \label{eq:main-uniform-response}
\end{equation}
Let $E(t,s)=\partial_sC(t,s)-R(t,s)$. There exists a nonnegative integrable function $a_F$ on $[0,\infty)$, obtained as an $L^1_{\mathrm{loc}}$ limit of the normalized annealed Fisher informations along one subsequence of dimensions, such that $\int_0^\infty a_F(s)\,\dd s\le\beta^2/2$, $|E(t,s)|^2\le a_F(s)$ for $t>s>0$, and $\sup_{s\ge T}a_F(s)\to0$ as $T\to\infty$. Further,
\begin{equation}
 \lim_{t\to\infty}\int_0^t|E(t,s)|^2\,\dd s=0.
 \label{eq:main-defect-l2}
\end{equation}
If $\beta>\bd$, let
\[
 q=q_\beta,\qquad x=\frac{(p-2)(1-q)}q\in(0,1),\qquad
 m_\infty=q\bigl(1-x^{1/p}\bigr)>0.
\]
Then
\begin{equation}
 \lim_{A\to\infty}\lim_{t\to\infty}
 \frac{p\int_0^{t-A}C(t,u)^{p-1}R(t,u)\,\dd u}
 {C(t,t-A)^p-C(t,0)^p}=x,
 \label{eq:main-weighted-ratio}
\end{equation}
while, for every fixed $S\ge0$,
\begin{equation}
 \lim_{L\to\infty}\liminf_{t\to\infty}
 \int_S^{t-L}E(t,u)_+\,\dd u\ge m_\infty,
 \qquad
 \liminf_{t\to\infty}t\int_0^tE(t,u)_+^2\,\dd u\ge m_\infty^2.
 \label{eq:main-defect-budget}
\end{equation}
\end{theorem}

The first half of \cref{thm:main-response} says that the response and the derivative of the correlation in the earlier time concentrate at short lags. The stationary limit in \cref{thm:main-aging} controls the response only on bounded windows of lags, while \eqref{eq:main-strong-rows} concerns all lags $v\in[0,t]$, an interval whose length tends to infinity. In particular, for large $A$, the squared response integrated over lags larger than $A$ is small at late times, and by \eqref{eq:main-uniform-response}, the response converges uniformly over all earlier times. This does not imply that the response integrated over these lags vanishes: the total susceptibility $\chi(t,0)$ grows more slowly than $\sqrt t$, but we do not prove that it stays bounded. The convergence of the second function in \eqref{eq:main-strong-rows} is a strong form of FDT in the earlier time: the violation $E(t,\cdot)$ tends to zero in $L^2(0,t)$.

The bound $|E|^2\le a_F$ explains why FDT holds in this sense. The function $a_F$ is a limit, along one subsequence of dimensions, of the normalized entropy production $a_N$ of \cref{sec:intro-model}, and its integral is at most $\beta^2/2$. The bound $|E(t,s)|^2\le a_F(s)$ says that the violation of FDT at an earlier time $s$ is controlled by the entropy produced at time $s$. Since $\sup_{s\ge T}a_F(s)\to0$ as $T\to\infty$, the violation is uniformly small when the earlier time $s$ is large. This is the mechanism of Cugliandolo, Dean, and Kurchan \cite{cdk1997}, which we make uniform in the dimension. We call $a_F$ the Fisher majorant. It need not be unique, but we choose one dimension subsequence and one majorant and use them throughout, including in \cref{thm:main-fisher-cost}.

FDT also holds in the later time, uniformly over all earlier times, for every $p\ge3$ and $\beta>0$:
\[
 \lim_{t\to\infty}\sup_{0\le s\le t}|\partial_tC(t,s)+r(t-s)|
 =\lim_{t\to\infty}\sup_{0\le s\le t}|\partial_tC(t,s)+R(t,s)|=0
\]
(\cref{resp:forward-derivative}). On the diagonal, $\partial_tC(t,t)=-1$ denotes the derivative from the causal side with the earlier time held fixed; $c'(0)=-1$ is a right derivative.

The second half of \cref{thm:main-response} shows that the violation of FDT nevertheless carries a positive integrated mass in the aging regime. For every fixed $S\ge0$, the integral of the positive part of $E(t,\cdot)$ over $[S,t-L]$ has lower limit at least $m_\infty$ when first $t\to\infty$ and then $L\to\infty$, as in \eqref{eq:main-defect-budget}, while the $L^2$ norm of $E(t,\cdot)$ tends to zero. These statements hold together because the violation is integrated over the growing interval $[S,t-L]$, whose length is of order $t$; the theorem does not say how the violation is distributed over this interval. By the second inequality in \eqref{eq:main-defect-budget}, the $L^2$ norm of $E_+$ cannot decay faster than $t^{-1/2}$. Cugliandolo, Dean, and Kurchan explained in \cite{cdk1997} why entropy production allows small pointwise violations of FDT to coexist with a nonzero integrated violation over an expanding time interval. In \cref{thm:main-response}, we prove such a coexistence for the CK solution. The ratio in \eqref{eq:main-weighted-ratio} identifies one moment of the slow response with the fluctuation--dissipation ratio $x$ of \cite{ck1993}. Indeed, if the response satisfied $R(t,u)=x\,\partial_uC(t,u)$ for $u\le t-A$, as the constant ratio of \cite{ck1993} predicts for large $A$, then the ratio would equal $x$, since $p\int_0^{t-A}C(t,u)^{p-1}\partial_uC(t,u)\,\dd u=C(t,t-A)^p-C(t,0)^p$. We prove this weighted identity, but not the pointwise relation.

The Fisher majorant is constructed in \cref{thm:fisher-clock}. The convergence in \eqref{eq:main-strong-rows} and the bound $\chi(t,0)=o(\sqrt t)$ are proved in \cref{resp:strong-row,resp:uniform-row,resp:backward-fdt}, the identity \eqref{eq:main-weighted-ratio} in \cref{resp:weighted-ratio}, and the bounds \eqref{eq:main-defect-budget} in \cref{resp:persistent-defect}.

The convergence over all earlier times also shows where the correlation at the plateau comes from. The Green covariance identity of \cref{prop:green-covariance} writes $C(t,s)$ as the sum of three terms: the contribution $R(t,0)R(s,0)$ of the initial condition, a thermal term, which is the contribution of the Brownian noise, and a memory term, which is the contribution of the colored noise of the effective process. For $t\ge s$, the thermal term and its stationary counterpart are
\[
 B(t,s)=2\int_0^sR(t,u)R(s,u)\,\dd u,
 \qquad b_{\rm th}(\tau)=2\int_0^\infty r(\tau+v)r(v)\,\dd v.
\]
For every $p\ge3$ and $\beta>0$, we show in \cref{resp:thermal} that
\begin{equation}
 \lim_{S\to\infty}\sup_{t\ge s\ge S}|B(t,s)-b_{\rm th}(t-s)|=0.
 \label{eq:main-thermal}
\end{equation}
The function $b_{\rm th}$ is completely monotone and tends to zero. Hence the thermal term vanishes on every late pair of times whose lag diverges. The initial term also vanishes, by \eqref{eq:main-forgetting}. Thus any correlation retained on such a pair is carried by the memory term, which comes from the interaction of each coordinate with the others.

For the cubic model, we prove two further predictions of the physical picture in one-sided form. Cugliandolo and Kurchan described the aging regime through the relations between correlations at three times \cite[arXiv version~1, Section~5]{ck1994}. For spherical models, they discussed a relation in which correlations below the plateau compose multiplicatively: for three times that are far apart, $C(t,u)/q$ is the product of $C(t,s)/q$ and $C(s,u)/q$ \cite[arXiv version~1, Section~5.1, Equation~(5.1.8)]{ck1994}.

\begin{theorem}[Three-time inequality and susceptibility in the cubic model]\label{thm:main-geometry}
Let $p=3$ and $\beta\ge\bd$. For every sequence of triples $u_n<s_n<t_n$ with $u_n\to\infty$, $s_n-u_n\to\infty$, and $t_n-s_n\to\infty$,
\begin{equation}
 \liminf_{n\to\infty}
 \bigl\{qC(t_n,u_n)-C(t_n,s_n)C(s_n,u_n)\bigr\}\ge0.
 \label{eq:main-product}
\end{equation}
The same holds for $qC(t_n,s_n)-C(t_n,u_n)C(s_n,u_n)$ and for $qC(s_n,u_n)-C(t_n,s_n)C(t_n,u_n)$. Further, under the same hypotheses,
\begin{equation}
 \liminf_{t\to\infty}\chi(t,0)\ge2(1-q).
 \label{eq:main-susceptibility}
\end{equation}
\end{theorem}

For the cubic model, $\bd=\sqrt{4/3}$ and $q\ge1/2$, with $q=1/2$ at $\beta=\bd$; all three product inequalities and the susceptibility bound include this point. The inequality \eqref{eq:main-product} is a one-sided version of the product relation of Cugliandolo and Kurchan for the CK solution, with no assumption on the time scale of aging. In words, after dividing correlations by $q$, the correlation across two successive long time intervals is asymptotically at least the product of the correlations across the two intervals. For a level $c\in(0,q)$, consider the last time after $s$ at which the correlation with the configuration at time $s$ is at least $c$. The inequality \eqref{eq:main-product} implies that these last times compose: for fixed levels $c_1,c_2\in(0,q)$ and a fixed level below $c_1c_2/q$, and for all sufficiently large $s$, the last time for the lower level is at least the last time for $c_2$, started from the last time for $c_1$ (\cref{cor:geometry-last-crossing}). These statements allow the correlation to return above a level before its last crossing, and \eqref{eq:main-product} does not assert a product equality.

The bound \eqref{eq:main-susceptibility} concerns the total response to a constant force switched on at time zero. If FDT held at all pairs of times, we would have $\chi(t,0)=1-C(t,0)\to1$. Of this amount, the stationary regime contributes $\int_0^\infty r(v)\,\dd v=1-q$, so \eqref{eq:main-susceptibility} forces at least another $1-q$ of integrated response outside every fixed window of lags. In the full physical prediction, including the absence of response mass at correlation zero, the aging regime contributes $xq$ instead of $q$, which gives $\chi(t,0)\to(1-q)+xq=(p-1)(1-q)$; for $p=3$, this is $2(1-q)$, which is smaller than one when $q>1/2$. We prove only the lower bound. Both results come from the geometry of the tree vectors of \cref{sec:geometry-core} whose energies approach the threshold: the product inequality from \cref{thm:geometry-threshold-product}, and the susceptibility bound from a limiting operator associated with such vectors (\cref{thm:geometry-universal-root}). They are proved in \cref{thm:geometry-actual-product,thm:geometry-susceptibility-lower}.

The next result relates relaxation below the plateau to entropy production. For a level $c\in(0,q)$ and a waiting time $s$, we consider the time $\tau_c(s)$ that elapses after $s$ until the correlation with the configuration at time $s$ first reaches the level $c$.

\begin{theorem}[Entropy production and relaxation below the plateau in the cubic model]\label{thm:main-fisher-cost}
Let $p=3$ and $\beta>\bd$, and let $a_F$ be the Fisher majorant of \cref{thm:main-response}. If $s_n\to\infty$, $\tau_n\to\infty$, and
\[
 \tau_n\int_{s_n}^{s_n+\tau_n}a_F(u)\,\dd u\longrightarrow0,
\]
then $C(s_n+\tau_n,s_n)\to q$ and $\chi(s_n+\tau_n,s_n)\to1-q$. For $0<c<q$, the first passage time
\[
 \tau_c(s)=\inf\{\tau\ge0:C(s+\tau,s)\le c\}
\]
is finite, tends to infinity as $s\to\infty$, and satisfies
\begin{equation}
 \liminf_{s\to\infty}\tau_c(s)
 \int_s^{s+\tau_c(s)}a_F(u)\,\dd u
 \ge\min\biggl\{1,
 \biggl[\frac{3\beta^2(2q-1)(q-c)^2}{\ell+15\beta^2}\biggr]^2\biggr\}>0.
 \label{eq:main-passage-cost}
\end{equation}
The same lower bound holds after replacing the integral by $\int_s^\infty a_F(u)\,\dd u$.
\end{theorem}

The first statement concerns windows $[s_n,s_n+\tau_n]$ whose starting times and lengths both tend to infinity. If the length of the window times the integral of $a_F$ over it tends to zero, then at the end of the window, the correlation with its start tends to the plateau, and the susceptibility tends to its stationary value $1-q=\int_0^\infty r(v)\,\dd v$, as if the solution were stationary. Relaxation below the plateau therefore requires entropy production: to reach a fixed level $c<q$, the elapsed time multiplied by the integral of $a_F$ over the elapsed time must stay bounded below, and \eqref{eq:main-passage-cost} gives an explicit lower bound. As in \cref{resp:fisher-chapter}, we call the quantity bounded in \eqref{eq:main-passage-cost} the Fisher cost of the passage. Denoting the right side of \eqref{eq:main-passage-cost} by $\kappa_c$, the bound with $\int_s^\infty a_F$ shows that
\[
 \tau_c(s)\ge\frac{\kappa_c-o(1)}{\int_s^\infty a_F(u)\,\dd u}\qquad(s\to\infty).
\]
Since $a_F$ is integrable, the right side tends to infinity: the time needed to relax below the plateau grows at least as fast as the reciprocal of the remaining integral $\int_s^\infty a_F$ of the chosen Fisher majorant. This is a quantitative form of aging, but it does not give a growth law for $\tau_c(s)$, because we do not know how fast $\int_s^\infty a_F$ decays. The strict condition $\beta>\bd$ is essential to this positive bound: its factor $2q-1$ vanishes at the transition. The proof, in \cref{resp:fisher-cost}, combines the backward response equation with the product inequality of \cref{thm:main-geometry}.

Our last main result says that, for $\beta\ge\bd$, the stationary correlation approaches its plateau algebraically: $c(t)-q$ is asymptotic to a positive constant times $t^{-a}$, with the exponent predicted by mode-coupling theory.

\begin{theorem}[Power-law approach to the plateau]\label{thm:main-stationary-tail}
Fix $p\ge3$ and $\beta\ge\bd$. There exists a unique $a\in(0,1/2)$ such that
\begin{equation}
 \frac{\Gamma(1-a)^2}{\Gamma(1-2a)}=\frac{(p-2)(1-q)}{2q},
 \label{eq:main-tail-exponent}
\end{equation}
and there exists $A_\beta\in(0,\infty)$ such that, as $t\to\infty$,
\begin{equation}
 c(t)-q\sim A_\beta t^{-a},\qquad
 (-1)^nc^{(n)}(t)\sim A_\beta(a)_nt^{-a-n}\quad(n\ge1).
 \label{eq:main-stationary-tail}
\end{equation}
Here $(a)_n=a(a+1)\cdots(a+n-1)$, and $n$ is fixed.
\end{theorem}

The relation \eqref{eq:main-tail-exponent} is the exponent relation of mode-coupling theory for the stationary profile of the aging solution, with exponent parameter $x/2$ \cite[arXiv version, Equation~(IV.6)]{bckm1996}, \cite{kimlatz2001}. At the transition, $x=1$ and $a\approx0.395$ for every $p$; as $\beta\to\infty$, the right side of \eqref{eq:main-tail-exponent} tends to zero and $a$ tends to $1/2$. The critical power law, with an amplitude that depends on the dynamics at finite time differences, is part of the classical asymptotic analysis of mode-coupling theory \cite[Section~4, Equation~(6), pp.~880--881]{goetze1998}. In \cref{thm:stationary-critical-amplitude}, we prove it for the stationary profile selected by the CK solution, together with a positive finite amplitude and the asymptotics of every fixed derivative. In particular, the response $r=-c'$ decays like $aA_\beta t^{-a-1}$, and the approach to the plateau has no single exponential relaxation time. These exponents describe how the stationary profile approaches its plateau after the limit of large waiting time; they are not exponents of the aging regime. None of the geometric or Fisher-cost arguments uses this power law.

\section{Relation to previous work}\label{sec:intro-related}

We now compare our results and methods with previous work.

\paragraph{The physical predictions and their later refinements.}
The predictions tested by our results come from the asymptotic solution of Cugliandolo and Kurchan \cite{ck1993}, described in \cref{sec:intro-picture}. Their derivation assumes that the stationary and aging regimes separate, and it imposes an asymptotic form on the correlation and response in the aging regime. The time scale of the aging regime is not determined by this calculation, because the asymptotic equations admit reparametrizations of time \cite[arXiv version~1, Section~4]{ck1994}. In \cref{thm:main-aging}, we prove the stationary FDT, the marginal plateau, the energy, and the forgetting of fixed earlier times for the CK solution itself, without assuming that the regimes separate. \Cref{thm:main-response} identifies one weighted moment of the slow response with the ratio $x$, and \cref{thm:main-geometry} gives one side of the product relation between three times. The full slow fluctuation--dissipation relation and the time scale of aging remain open. The effective-temperature interpretation of $x$ was developed by Cugliandolo, Kurchan, and Peliti \cite{ckp1997}. Franz, M\'ezard, Parisi, and Peliti related the dynamical response to the equilibrium overlap distribution under additional hypotheses \cite{fmpp1998}, and they pointed out that the unperturbed spherical $p$-spin model does not satisfy this relation. We do not use equilibrium overlaps. Altieri, Biroli, and Cammarota developed a formulation of aging within dynamical mean-field theory in \cite{altieri2020}, and Chamon, Kennett, Castillo, and Cugliandolo developed the reparametrization invariance of the slow dynamics in \cite{chamon2002}. Ghimenti, Berthier, Kurchan, and van Wijland related accelerated glass dynamics to time-reparametrization softness, through numerical comparisons and an asymptotic analysis of a nonreciprocally coupled spherical $p$-spin model \cite[pp.~4--5 and SI Appendix, Section~7]{ghimenti2026}.

The picture has also been tested numerically and refined. Kim and Latz integrated the CK equations numerically over long times and compared the result with the predicted fluctuation--dissipation relation and with the crossover to the aging regime \cite{kimlatz2001}. Andreanov and Lef\`evre proposed a form of the time scale of aging by matching the stationary and aging regimes within a scaling expansion \cite{andreanovlefevre2006}. Folena, Franz, and Ricci-Tersenghi found that in zero-temperature quenches of mixed spherical models the long-time dynamics can depend on the initial condition and can reach energies below the threshold \cite[Sections~V--VII]{folena2020}. Folena and Zamponi tested weak ergodicity breaking numerically in mean-field spin glasses over long times, mainly in mixed models \cite{folenazamponi2023}, and Lang, Sachdev, and Diehl developed a numerical renormalization method that reaches longer integration times \cite{lang2025}. Citro and Ricci-Tersenghi obtained numerical evidence for strong ergodicity breaking, in which the correlation with the initial configuration stays positive at long times, in the spherical $3+4$ model quenched from near the dynamical transition to zero temperature \cite{citro2025}. The studies of mixed models and of other initial conditions indicate that the selection of the threshold in \cref{thm:main-aging} depends on our hypotheses.

\paragraph{Rigorous limits on fixed time horizons and the CK equations.}
The rigorous derivation of mean-field dynamical equations began with the large deviation methods on path space of Ben Arous and Guionnet for Langevin dynamics of spin glasses: a large deviation principle for asymmetric interactions \cite{bg1995} and a large deviation upper bound for symmetric interactions \cite{bg1997}, in regimes of short times or high temperature. Grunwald proved an analogous asymmetric large deviation principle and symmetric upper bound for Glauber dynamics, the spin-flip dynamics of Ising spin glasses, without restrictions on the fixed time horizon or temperature \cite{grunwald1996}. Guionnet obtained averaged propagation of chaos, that is, asymptotic independence of finitely many coordinates, and a quenched single-spin limit expressed as convergence in law over the disorder \cite{guionnet1997}. For spherical models, Ben Arous, Dembo, and Guionnet proved in \cite[Theorem~1.2]{bag2006} that the empirical correlation and integrated response of the soft-spherical Langevin dynamics converge on fixed time horizons to the solution of the CK equations. Dembo, Guionnet, and Mazza took the limit of these deterministic equations as the soft constraint becomes hard \cite[Proposition~1.1]{dgm2007}. In the same paper, they proved exponential relaxation and convergence to a stationary FDT solution at sufficiently high temperature \cite[Proposition~1.2 and Theorem~1.3]{dgm2007}, and they classified the solutions of the scalar memory equation by the largest admissible plateau and a stability inequality \cite[Proposition~1.4]{dgm2007}. They also showed that the CK solution has stationary FDT limits at every temperature, under the asymptotic conditions (6.5) and (6.6) of their paper \cite[Proposition~6.1]{dgm2007}. Sellke proved the convergence theorem for the hard-spherical diffusion itself \cite[Theorem~4.3]{sellke2024}, by coupling it with soft-spherical diffusions \cite[Lemmas~3.1 and~3.2]{sellke2024}. Guionnet and Mazza studied the long-time behavior of nonlinear Kraichnan equations, which are integral equations for response kernels with a prescribed covariance, in \cite{guionnetmazza2005}, using positive noncrossing expansions: series with nonnegative terms indexed by pairings of ordered points whose arcs do not cross. Dembo and Subag derived dynamical equations on fixed time horizons for initial conditions in prescribed bands near critical points \cite{ds2020} and for Gibbs distributed initial conditions \cite{ds2025}. The latter result includes the initialization at inverse temperature zero, which is the uniform start, together with the limit of a hard spherical constraint \cite[Theorem~1.8(a) and Proposition~1.10]{ds2025}. They also characterize stationary FDT dynamics and relaxation within equilibrium bands \cite[Proposition~1.14 and Corollary~1.18]{ds2025}. The equilibrium relaxation conclusions concern Gibbs starts at the ambient temperature, whereas our long-time selection follows a quench from the uniform law to inverse temperature $\beta>0$. Dembo, Lubetzky, and Zeitouni proved universality on fixed time horizons for Langevin-like dynamics with independent asymmetric interactions \cite[Theorem~1.1]{dlz2021}. For diffusions with linear random matrix interactions, Dembo and Gheissari proved universality of averaged observables for both symmetric and asymmetric couplings, using a stochastic Taylor expansion \cite[Theorem~1 and Corollary~3]{dembogheissari2021}. Surveys of this program are \cite{ba2002} and \cite[pp.~117--121]{guionnet2007}.

Rigorous dynamical mean-field theory has recently been extended to first-order methods with random data \cite{ccm2021} and to stochastic gradient descent \cite{gerbelot2024}, and its equations have been derived by a dynamical cavity method \cite{dandi2026}. Long-time results are known for Langevin dynamics in spiked matrix models, which are random matrices with a low-rank signal \cite{lss2023}. For Langevin diffusions with adaptive drift, Fan, Ko, Loureiro, Lu, and Shen proved convergence of the linear response on fixed horizons \cite{fan2025a}, and long-time results under an approximate time-translation invariance hypothesis, which they verify under a log-Sobolev inequality that is uniform in the dimension \cite[Theorems~2.5 and~2.9]{fan2025b}. For the orthogonally invariant Sherrington--Kirkpatrick model, Fan, Misiakiewicz, Wang, and Wen proved stationary FDT limits under rapid mixing on time scales that do not depend on the dimension and suitable confining curvature \cite[Lemma~4.1]{fan2026}. The models and hypotheses of these works differ from ours.

We take the results of \cite{bag2006,dgm2007,sellke2024} as input, through \cref{prop:ck-wellposed,prop:hard-bridge}, and prove long-time statements for the CK solution at every temperature. The main difference from \cite[Proposition~6.1]{dgm2007} is that our stationary limits (\cref{prop:stationary-clusters}) require no asymptotic hypothesis on the solution: entropy production supplies the needed control.

\paragraph{Aging for the two-spin model and on other time scales.}
Ben Arous, Dembo, and Guionnet proved aging for the soft-spherical model with $p=2$ in the order of limits used here \cite{bag2001}. Their analysis uses the spectral decomposition of the interaction matrix, which reduces the dynamics to coupled Ornstein--Uhlenbeck processes. This linear structure is absent for $p\ge3$, where a positive stationary plateau must satisfy the nonlinear marginality condition \eqref{eq:main-marginality}.

A different rigorous theory of aging concerns time scales that grow with the dimension, on which the dynamics is activated and is trapped in deep states for long times. Ben Arous, Bovier, and Gayrard proved aging for the Glauber dynamics of the random energy model below the critical temperature \cite{bbg2003a,bbg2003b}. Ben Arous and \v{C}ern\'y showed that the arcsine law of a stable subordinator is a universal aging scheme for trap models: under their hypotheses, the probability that the process occupies the same trap at two late times converges to a function of their ratio given by the generalized arcsine law \cite{bc2008}. For random-hopping-time dynamics, in which the jumps do not depend on the energy and the holding times do, Ben Arous, Bovier, and \v{C}ern\'y proved that Ising $p$-spin models with $p\ge3$ age like the random energy model on intermediate exponential time scales \cite{bbc2008}, and Ben Arous and G\"un obtained extremal aging, which is governed by the longest holding times, on subexponential time scales \cite{bg2012}. Bovier and Gayrard proved convergence of clock processes in random environments, and aging for random-hopping-time dynamics of the $p$-spin SK model \cite{boviergayrard2013}. For the Metropolis dynamics of the random energy model, \v{C}ern\'y and Wassmer proved convergence of the clock process and explained why aging of the two-time correlation does not follow directly from it \cite[Theorem~1.1 and Remark~1.2(3)]{cw2017}, and Gayrard proved two-time aging \cite{gayrard2019}. These results let the observation times grow with $N$ and describe the dynamics through clock processes, which record the total time spent in the states visited by the dynamics; this is the trap mechanism of \cite{bouchaud1992}. In the CK limit, the dimension tends to infinity first on fixed horizons, activation is invisible, and the aging of \cref{thm:main-aging} is a property of the deterministic CK solution. The differences between these orders of limits are discussed in the surveys \cite{ba2002,bc2006} and \cite[pp.~120--121]{guionnet2007}.

\paragraph{Landscape and algorithmic thresholds, shattering, and metastability.}
In physics, the threshold energy $\Eth=2\sqrt{(p-1)/p}$ is the energy of the marginal TAP states at zero temperature \cite[Equation~(9)]{cs1995}. Auffinger, Ben Arous, and \v{C}ern\'y used random matrix theory to compute the complexity of the critical points of the pure spherical $p$-spin Hamiltonian, that is, the exponential growth rate of their number. Up to the sign convention, $\Eth$ is the energy $E_\infty$ of their analysis \cite[Equation~(2.13)]{abc2013}. Auffinger and Ben Arous extended these computations to mixed models \cite{ab2013}. In pure models, Subag proved by a second moment computation that, for each fixed energy strictly between the threshold and the ground-state energy, the number of critical points with energy above it, divided by its expectation, converges to one in $L^2$, in our sign convention \cite[Theorem~1 and Corollary~2]{subag2017}. In optimization, $\Eth$ is the algorithmic threshold of pure models, that is, the largest energy that algorithms of a natural stable class can reach. Subag constructed an algorithm for spherical models, based on the Hessian, that reaches the energy $\int_0^1\sqrt{\xi''(u)}\,\dd u$, which equals $\Eth$ for $\xi(u)=u^p$ \cite{subag2021}. Montanari \cite{montanari2019} and El Alaoui, Montanari, and Sellke \cite{eams2021} constructed message passing algorithms for Ising models. These algorithms iteratively use the gradient of the Hamiltonian, and with probability tending to one, they reach energies arbitrarily close to the ground-state energy under a no-overlap-gap assumption on the Parisi order parameter: in \cite[Assumption~1]{montanari2019}, that the support of the Parisi measure is an interval $[0,q_*(\beta)]$ for all sufficiently large $\beta$, and in \cite[arXiv version~1, Assumption~2]{eams2021}, that the zero-temperature Parisi minimizer is strictly increasing on $[0,1)$. Huang and Sellke proved a matching hardness result. Consider algorithms that map the Gaussian disorder coefficients to points of the ball of radius $\sqrt N$ in $\R^N$, with Lipschitz constants, in the Euclidean metrics, bounded independently of $N$. Then the probability that the energy per coordinate of the output exceeds the algorithmic threshold by a fixed positive amount tends to zero as $N\to\infty$. They proved this for even models in \cite{hs2021}, and for multi-species spherical models, including odd pure models and external fields, in \cite[Theorem~1 and Corollary~1.8]{hs2023}. Sellke proved that at low temperature, Langevin dynamics reaches energies close to $\Eth$ on bounded time horizons \cite[Theorem~1.1]{sellke2024}. He also proved that, from an initial point independent of the disorder, it stays below $\Eth$ on each fixed time horizon with probability tending to one as $N\to\infty$ \cite[Theorem~1.7]{sellke2024}. Our energy formula agrees with these results when the zero-temperature limit is taken after the long-time limits (\cref{rem:low-temperature}).

For every stationary limit with positive plateau, we use $\Eth$ as an upper bound for the limiting energy of normalized prefix centers, which are the projections of late tree vectors described in \cref{sec:main-results} (\cref{prop:threshold-inputs}). For $p\in\{3,4\}$, we prove this bound with a deterministic inequality for unit vectors in a Hilbert space graded by ordered trees (\cref{thm:geometry-tree-ceiling,thm:geometry-quartic-tree-ceiling}). The inequality has a sharp constant, and the deficit from this constant controls the squared norm of the components of a vector at every fixed height, where the height of a tree is the largest distance from its root to a leaf. Branching constructions appear in the algorithms of Subag \cite[Remark~6]{subag2021} and in the overlap geometry of Huang and Sellke \cite[Remarks~1.2 and~1.3]{hs2021}; here we retain quantitative height defects inside the causal covariance realization. For general $p$, we use the theorem of Huang and Sellke \cite{hs2023} together with the coupling of Sellke \cite[Lemmas~3.1 and~3.2]{sellke2024}.

Sellke's work on the marginal stability of approximate ground states \cite{sellke2025} concerns a different part of the landscape, and we do not relate the stationary response to a Hessian. The finite-dimensional dynamics has also been studied directly. Gheissari and Jagannath estimated the spectral gap of spherical spin glass dynamics \cite{gj2019}, and Ben Arous, Gheissari, and Jagannath bounded the evolution of the energy and the gradient by differential inclusions \cite{bgj2020}. Ben Arous and Jagannath separated shattering, a decomposition into exponentially many well-separated and individually subdominant regions whose union has the full limiting free energy, from metastability, which concerns long residence times near a region before escape \cite{bj2024}. El Alaoui, Montanari, and Sellke proved shattering of the Gibbs measure of pure spherical models in a range of temperatures, and showed that Langevin dynamics started from the Gibbs measure keeps a large correlation with its initial point on exponentially long time scales \cite[Corollary~2.9]{shattering2025}. This plateau concerns equilibrium initial data and time scales that grow with $N$. The plateau in \cref{thm:main-aging} concerns a uniform initial point and the limit in which the dimension tends to infinity first.

\paragraph{Methods.}
We describe the origin of the main tools of the proofs; \cref{sec:proof-overview} explains how they fit together.

\emph{Entropy production and FDT.} Agarwal \cite{agarwal1972} treats response around nonthermal steady states under detailed balance; nonequilibrium diffusion response identities are given in \cite{bmw2009}. Cugliandolo, Dean, and Kurchan bounded violations of FDT by the entropy production in \cite[arXiv version~2, Equations~(12)--(14) and~(18)]{cdk1997}. We use the Bochner formula for a Laplacian with a gradient drift \cite[Proposition~3]{bakryemery1985} and the calculus of reversible diffusions \cite{bgl2014}. In \cref{sec:entropy-clusters}, we establish dimension-uniform bounds for the projected hard-spherical response and pass them to a single Fisher majorant for the CK solution, which controls the violation of FDT at all pairs of times $t>s>0$ and gives stationary FDT limits.

\emph{Effective processes and the Green identity.} The covariance identity of \cref{prop:green-covariance} has the variation-of-constants interpretation of the linear effective single-spin equation with white and colored Gaussian noise \cite[arXiv version~2, Section~9.5, Equation~(9.56), p.~55]{cugliandolo2004}. We prove the identity directly from the CK equations in \cref{sec:ck-structure}.

\emph{Expansions with nonnegative terms.} Positive noncrossing expansions of the response equation appear in the work of Guionnet and Mazza \cite[Section~2.1]{guionnetmazza2005} and of Dembo, Guionnet, and Mazza \cite[Section~2]{dgm2007}. In \cref{sec:geometry-core}, we expand the covariance instead. We organize the expansion by ordered trees whose leaves record either the initial condition or a time of the thermal noise. This realizes $C(t,s)$ as the inner product of unit vectors in one deterministic Hilbert space, and the projection onto trees with early leaves plays the role of conditioning on the early noise.

\emph{Independent futures.} Trajectories with a common past and independent future noise were used by Cugliandolo and Dean \cite[arXiv version~1, Section~II]{cugliandolodean1995} and by Barrat, Burioni, and M\'ezard, who classified aging by the overlaps of such trajectories \cite[Section~2]{bbm1996}. The overlap of two independent continuations of lengths $a$ and $b$ equals the correlation at the reflected lag $a+b$ in equilibrium \cite[Section~2, Equation~(9)]{bbm1996}. Franz and Virasoro proposed averaging cloned futures to estimate the magnetization of the region of configurations that the dynamics explores on intermediate time scales \cite[arXiv version~3, Section~III]{franzvirasoro2000}, and noise-averaged magnetizations and their relation to the threshold are central to the dynamical TAP approach of Biroli \cite[arXiv version~2, Section~III.B]{biroli1999}. Time-averaged magnetizations whose normalized squared norm tends to $q$ also appear in \cite[arXiv version~1, Section~5.1, Equation~(5.1.9)]{ck1994}. Our prefix centers are deterministic counterparts of these averages. In \cref{sel:prefix-reflection}, we prove the reflected formula for the stationary limits of the CK solution: the older memory cancels exactly in a closed integral equation for the times after the earlier time, before any limit is taken. For general $p$, \cref{app:general-selection} uses finitely many branched future trajectories of the diffusion.

\emph{Complete monotonicity.} The reflected kernel $(a,b)\mapsto c(a+b)$ is positive semidefinite, and together with the boundedness and continuity of $c$, this implies that $c$ is completely monotone, by the Bernstein--Widder theorem and its semigroup form \cite{widder1941,bcr1984}. Franosch and Voigtmann proved complete monotonicity for Brownian mode-coupling equations of mixtures by a different route, an iteration argument \cite[arXiv version~1, Section~III]{fv2002}.

\emph{Gaussian conditioning.} In \cref{app:general-selection}, we prove forgetting for general $p$ by revealing the disorder in stages, through the gradients of the Hamiltonian at finitely many points chosen from the history, and by applying the external-field threshold of \cite[Theorem~1 and Corollary~1.8]{hs2023} at each stage. An Euler approximation of an earlier state lies in the span of the initial point, finitely many Brownian increments, and the computed gradients, and we condition on the Gaussian coefficients needed to determine the gradients on this span. Gaussian conditioning is classical \cite{janson1997}. Adaptive Gaussian conditioning is used by Bayati and Montanari for approximate message passing \cite[arXiv version~4, Section~3.8]{bayatimontanari2011} and by Bolthausen for the iterative construction of solutions of the TAP equations \cite[Section~3]{bolthausen2014}. El Alaoui, Montanari, and Sellke conditioned the disorder tensor on adaptive gradient observations in their analysis of message passing for spin glasses \cite[arXiv version~1, Appendix~A]{eams2021}, and Huang and Sellke extended this conditioning to multi-species spherical models \cite[Appendix~A]{hs2024amp}. We adapt the method to a fixed finite history of the Langevin dynamics.

\emph{Mode-coupling and renewal tools.} The scalar memory equation of a stationary profile belongs to the class studied in mode-coupling theory \cite[Sections~3.1--3.2]{goetzesjogren1992} and \cite[Section~4.2.1, Equations~(4.15)--(4.17), pp.~192--193]{goetze2009}. The structural theory of these equations is developed in \cite{gs1995} and \cite{fv2002}, including the characterization of the plateau \cite[arXiv version~1, Section~IV]{fv2002}, and Sj\"ogren interpreted them through regenerative processes, which restart at random times \cite{sjogren2003}. Our proof of \cref{thm:main-stationary-tail} in \cref{sec:stationary-relaxation} combines comparison principles, a weighted Schur inequality for the convolution in the scalar equation, a factorization of gamma variables, which turns a quadratic balance into a renewal equation, the key renewal theorem \cite{feller1971}, and the monotone density theorem \cite[Section~1.7.3]{bgt1987}.

\emph{The three-time inequality.} The product relation for three times is part of the physical analysis of Cugliandolo and Kurchan \cite[arXiv version~1, Section~5.1, Equation~(5.1.8)]{ck1994}. We prove a uniform one-sided inequality for the CK solution in \cref{thm:geometry-actual-product}, from the constrained geometry of tree vectors whose energies approach the threshold (\cref{thm:geometry-threshold-product}). We use it for the bound on the Fisher cost in \cref{thm:main-fisher-cost}.

\paragraph{Concurrent work.}
Several recent works study related questions with different models, initial conditions, or hypotheses. Dembo and Subag analyzed dynamics started from Gibbs measures, including limits on fixed time horizons and the stationary FDT dynamics and equilibrium band relaxation described above \cite{ds2025}. Their theorem includes the initialization at inverse temperature $\beta_0=0$, hence the uniform start \cite[Theorem~1.8(a) and Proposition~1.10]{ds2025}; for initialization from a low-temperature Gibbs measure, \cite[Theorem~1.8(b)]{ds2025} assumes additional conditions on the mixture and on the equilibrium Parisi measure (genericity and strict one-step replica symmetry breaking, in the sense of \cite[Section~1]{ds2025}). The works \cite{fan2025a,fan2025b,fan2026} on dynamical mean-field equations, described above, concern other models; their long-time results use approximate time-translation invariance or mixing hypotheses. Lang, Citro, Leuzzi, and Ricci-Tersenghi computed the aging phase diagram and the asymptotic energies of mixed spherical models at zero temperature, within a formal asymptotic analysis \cite{lang2026}. Cipolloni and Cugliandolo studied coupled spherical $p$-spin systems, including their aging regimes and fluctuation--dissipation relations \cite{cipolloni2026}. Our contribution relative to these works is the set of long-time conclusions for the uniformly initialized CK solution of pure models: the selection of the plateau and the energy, the forgetting of fixed earlier times, and the results on the response.

\section{Scope and open problems}\label{sec:intro-scope}

Our conclusions concern the CK solution started from a uniform point, not an arbitrary solution of a reduced asymptotic system, and they apply to Langevin dynamics only through the limit on fixed time horizons of \cref{prop:hard-bridge}. For every $p\ge3$, the stationary limits, the convergence of the energy, the forgetting of fixed earlier times, and the convergence of the response over all earlier times hold at every positive temperature. The identification of the plateau and the energy, correlation aging, and the stationary power law hold for $\beta\ge\bd$, and the weighted ratio and the persistent violation of FDT hold for $\beta>\bd$. For the cubic model, the three-time inequality and the susceptibility bound hold for $\beta\ge\bd$, while the positive Fisher cost requires $\beta>\bd$.

Several parts of the physical picture remain open, and we formulate them in \cref{sec:discussion}:
\begin{itemize}
\item the full slow fluctuation--dissipation relation, of which \eqref{eq:main-weighted-ratio} is one weighted moment (\cref{sec:discussion-fdr});
\item an upper bound matching \eqref{eq:main-susceptibility} for the susceptibility (\cref{sec:discussion-susceptibility});
\item the time scale of aging, which the three-time inequality and the Fisher cost constrain but do not determine (\cref{sec:discussion-clock});
\item the selection of the plateau for $\beta_{\mathrm m}\le\beta<\bd$ (\cref{sec:discussion-below});
\item the behavior on time scales that grow with the dimension, and at zero temperature (\cref{sec:discussion-limits}).
\end{itemize}

\section{Organization of the paper}\label{sec:intro-organization}

\Cref{sec:proof-overview} gives an overview of the proofs. They have three main ingredients. From entropy production, we obtain stationary FDT limits on bounded late windows. The representation of the correlation by vectors in a Hilbert space indexed by trees separates the contributions of the earlier and the later noise. An energy bound for the earlier contributions selects the marginal plateau and shows that fixed earlier configurations are forgotten. The rest of the paper has five parts.

\Cref{part:foundations} contains the foundations. \Cref{sec:model} states the model, the conventions, the CK equations, and the published inputs, including the limit theorem on fixed time horizons. In \cref{sec:ck-structure}, we prove the energy identity, the Green covariance identity, and estimates that are uniform in time. In \cref{sec:entropy-clusters}, we use entropy production to construct the Fisher majorant and to show that all limits of the CK solution on late time windows are stationary and satisfy FDT.

\Cref{part:selection} proves \cref{thm:main-aging}. \Cref{sec:scalar} derives the scalar equation of a stationary limit and its constraints. \Cref{sec:geometry-core} realizes $C$ as the Gram kernel of unit vectors in a Hilbert space graded by trees, and proves the reflection identity and complete monotonicity. \Cref{sec:stationary-selection} states the two threshold inputs (\cref{prop:threshold-inputs}) and deduces from them the selection of the plateau, the forgetting of earlier times, and correlation aging. \Cref{sec:tree-ceilings} proves the threshold inputs for $p\in\{3,4\}$ with deterministic inequalities in the tree space, and \cref{app:general-selection} proves them for every $p\ge3$.

\Cref{part:response} proves \cref{thm:main-response}. In \cref{resp:chapter}, we prove the convergence of the response and of the derivatives of the correlation over all earlier times, and the convergence of the thermal term. In \cref{resp:integrated-chapter}, we prove the weighted ratio and the persistent violation of FDT.

\Cref{part:geometry} concerns the cubic model and proves \cref{thm:main-geometry,thm:main-fisher-cost}. \Cref{sec:geometry-results} studies tree vectors near the threshold, \cref{sec:three-time} transfers their geometry to the CK correlations and proves the susceptibility bound, and \cref{resp:fisher-chapter} proves the bound on the Fisher cost.

\Cref{part:relaxation} contains \cref{sec:stationary-relaxation}, which proves \cref{thm:main-stationary-tail}, and \cref{sec:discussion}, which discusses the open problems. \Cref{sec:stationary-relaxation} can be read once \cref{part:selection} is available.

\Cref{tab:main-results-map} lists, for each main theorem, where it is proved and the main results that the proof uses.

\begin{table}[ht]
\centering
\small
\renewcommand{\arraystretch}{1.3}
\begin{tabular}{@{}>{\raggedright\arraybackslash}p{0.15\textwidth}>{\raggedright\arraybackslash}p{0.27\textwidth}>{\raggedright\arraybackslash}p{0.5\textwidth}@{}}
\toprule
Result & Proved in & Main inputs\\
\midrule
\Cref{thm:main-aging} & \Cref{sec:stationary-selection}; the proof is assembled at the end of \cref{sec:aging} & \Cref{prop:hard-bridge}; \cref{prop:stationary-clusters}; \cref{prop:scalar-clusters}; \cref{sel:prefix-reflection,sel:cluster-complete-monotonicity}; \cref{prop:threshold-inputs}, proved in \cref{sec:tree-ceilings} for $p=3,4$ and in \cref{app:general-selection} for all $p$\\
\Cref{thm:main-response} & \Cref{thm:fisher-clock}; \cref{resp:strong-row,resp:uniform-row,resp:backward-fdt}; \cref{resp:weighted-ratio,resp:persistent-defect} & \Cref{thm:main-aging}; \cref{prop:green-covariance}; \cref{sel:prefix-reflection}; \cref{thm:geometry-tree-ceiling} ($p=3$) or \cref{resp:appendix-p} (all $p$); \cref{resp:transport}\\
\eqref{eq:main-thermal} & \Cref{resp:thermal} & \Cref{sel:cluster-complete-monotonicity}; \cref{resp:strong-row}\\
\Cref{thm:main-geometry} & \Cref{thm:geometry-actual-product,thm:geometry-susceptibility-lower} & \Cref{thm:geometry-tree-ceiling}; \cref{sel:threshold-centers}; \cref{thm:geometry-universal-root,thm:geometry-threshold-product}; \cref{resp:strong-row,resp:uniform-row,resp:forward-derivative}\\
\Cref{thm:main-fisher-cost} & \Cref{resp:fisher-cost} & \Cref{prop:backward-response}; \cref{thm:fisher-clock}; \cref{resp:uniform-row,resp:forward-derivative}; \cref{thm:geometry-actual-product}\\
\Cref{thm:main-stationary-tail} & \Cref{thm:stationary-critical-amplitude} & \Cref{prop:scalar-clusters}; \cref{sel:cluster-complete-monotonicity}; \cref{sel:selected-sector}\\
\bottomrule
\end{tabular}
\caption{The main results, the sections and results in which they are proved, and the main inputs of their proofs.}
\label{tab:main-results-map}
\end{table}

\chapter{Overview of the proofs}\label{sec:proof-overview}

In this section, we explain the main steps of the proofs and the ideas behind them. We describe each step informally and point to the precise statements in later sections. In \cref{sec:overview-difficulty}, we explain what makes the long-time behavior of the CK solution hard to analyze. \Cref{sec:overview-clusters,sec:overview-scalar,sec:overview-tree,sec:overview-selection,sec:overview-inputs} describe the proof of \cref{thm:main-aging}, and \cref{sec:overview-response,sec:overview-geometry,sec:overview-power-law} describe the proofs of the other main theorems. Throughout, $p\ge3$ and $\beta>0$ are fixed, $m=p-1$, and $(C,R,\lambda)$ is the CK solution of \cref{sec:intro-model}. We illustrate several formulas in the cubic model, where $m=2$, $\bd=\sqrt{4/3}$, and $\qd=1/2$.

\section{The difficulty}\label{sec:overview-difficulty}

After the dimension limit, which we take first on every fixed time horizon (\cref{prop:hard-bridge}), we must study the deterministic solution $(C,R,\lambda)$ of the integro-differential equations \eqref{eq:ck-response}--\eqref{eq:ck-multiplier} as both of its times tend to infinity. The aging statement \eqref{eq:main-aging} combines two assertions of different kinds. The first concerns bounded windows of lags at late times, where we must show that the solution becomes stationary and that its stationary correlation decreases to a positive plateau $q_\beta$. The second concerns the whole history: we must show that $C(t,s)\to0$ as $t\to\infty$ for every fixed $s$. Note that a positive stationary plateau implies neither the second assertion nor the noncommutation of the two limits in \eqref{eq:main-aging}. The stationary profile records $C(T+t,T+s)$ only for bounded $t-s$, and it carries no information about $C(t,s)$ for fixed $s$.

Cugliandolo and Kurchan derived the plateau, the energy, and the loss of memory from an asymptotic solution of the CK equations \cite{ck1993}. Their derivation separates a stationary regime at bounded lags from an aging regime below the plateau, in which it imposes an asymptotic relation between response and correlation; within this ansatz, the plateau is fixed by the marginality condition. A proof cannot start from these assumptions. We do not know in advance that the two regimes separate, on which time scale the aging regime lives, or how response and correlation are related there, and our results determine neither the time scale nor the full relation.

The difficulty is visible in the equations. If we translate both observation times to late times and keep their difference bounded, the memory integral of the response equation \eqref{eq:ck-response} involves only the times between them. The memory integrals of the correlation equation \eqref{eq:ck-correlation}, however, run over $[0,t]$ and $[0,s]$, and on a late window they involve the overlap of late states with the remote past, which a limit on late windows does not control.

We make no assumption on this older memory. There are three main ingredients. The first is entropy production. We prove that the total entropy produced by the diffusion relative to its Gibbs measure is at most $\beta^2N/2$ on average over the disorder, and deduce that every limit of the solution on late windows is stationary and satisfies FDT (\cref{sec:overview-clusters}). The closed response equation then constrains the plateau, but it does not select it (\cref{sec:overview-scalar}). The second ingredient is a realization of $C$ as the Gram kernel of unit vectors $v_t$ in a Hilbert space graded by trees, whose leaves record the times of the thermal noise. Projecting a late state onto the trees with early leaves plays the role of conditioning on the early noise, and the Gram kernel of these projections satisfies an equation from which the older memory cancels exactly (\cref{sec:overview-tree}). The third ingredient is an energy ceiling for normalized projections of late states. Together with the scalar constraints, it forces every positive plateau to be marginal, and when it is saturated, it forces the projections to forget every fixed earlier state (\cref{sec:overview-selection,sec:overview-inputs}). For $p\in\{3,4\}$, the ceiling is a deterministic inequality in the tree space; for general $p$, we use the algorithmic threshold of Huang and Sellke \cite[Theorem~1 and Corollary~1.8]{hs2023}.

\section{From the diffusion to stationary clusters}\label{sec:overview-clusters}

\paragraph{Inputs.}
The bridge of \cref{prop:hard-bridge} is Sellke's convergence theorem for hard-spherical Langevin dynamics \cite[Theorem~4.3]{sellke2024}: on every fixed time horizon, empirical correlations and a Brownian observable converge, and by \cref{lem:brown-force-response}, this observable is the integrated response to an additive force. In finite dimension, the normalized trace response has diagonal $\kappa_N=(N-1)/N$. Its annealed mean converges to $R$ uniformly on fixed causal triangles, by the compact-time modulus of the tangent flow and uniqueness of the limit identified by the primitives (\cref{rem:no-pointwise-response}); this is stronger than the convergence of primitives that the published bridge gives. Integrating by parts in the Gaussian disorder, we find that $e_\beta(t)$ is the limit of the mean energy and that $\lambda=1+p\beta e_\beta$ (\cref{prop:energy-bridge}). Our basic structural tool is the Green covariance identity of \cref{prop:green-covariance},
\begin{equation}\label{eq:overview-green}
\begin{aligned}
 C(t,s)={}&R(t,0)R(s,0)+2\int_0^{t\wedge s}R(t,u)R(s,u)\,\dd u\\
 &+\beta^2p\int_0^t\int_0^sR(t,u)R(s,w)C(u,w)^{p-1}\,\dd w\,\dd u ,
\end{aligned}
\end{equation}
whose three terms come from the initial condition, the thermal noise, and the colored memory, each propagated by the response. It has the variation-of-constants interpretation of the linear effective single-spin equation with white and colored Gaussian noise \cite[arXiv version~2, Section~9.5, equation~(9.56), p.~55]{cugliandolo2004}. We prove the identity directly from the CK equations, whose derivation and rigorous fixed-horizon limits are discussed in \cite{ck1993,bag2006,ds2020}. We also show that $1\le\lambda\le\Lambda$ and $0\le R\le C$, and that the translated kernels $C_T(a,b)=C(T+a,T+b)$, $R_T(a,b)=R(T+a,T+b)$, and $\lambda_T(a)=\lambda(T+a)$ are precompact as $T\to\infty$ (\cref{sec:ck-structure}).

\paragraph{Entropy production.}
Fix the disorder, let $\rho_t$ be the density of $X_t$ with respect to the uniform measure $\omega_N$ on $S_N$, and recall that the diffusion is reversible with respect to the Gibbs density $g_N\propto e^{\beta H_N}$. The relative entropy of $\rho_t$ with respect to $g_N$ decreases at the rate $I_H(t)$, the relative Fisher information. Since the initial law is uniform, the initial relative entropy is $\log Z_{N,\beta}-\beta\int H_N\,\dd\omega_N$, whose mean is at most $\log\E Z_{N,\beta}=\beta^2N/2$ by Jensen's inequality. With $a_N(t)=N^{-1}\E I_H(t)$, we arrive at the budget
\begin{equation}\label{eq:overview-budget}
 \int_0^\infty a_N(t)\,\dd t\le\frac{\beta^2}2
\end{equation}
of \cref{lem:entropy-budget}. The link with FDT is the exact identity of \cref{lem:finite-dimensional-fdt}. If $c_H$ and $r_H$ are the correlation and the response at fixed disorder, and $m_\tau(x)$ is the conditional mean of $X_{s+\tau}$ given $X_s=x$, then for $t>s$ we have
\[
 r_H(t,s)-\partial_sc_H(t,s)
 =\frac1N\int\rho_s\,\Bigl\langle m_{t-s},\nabla_S\log\frac{\rho_s}{g_N}\Bigr\rangle\,\dd\omega_N .
\]
At $s=0$, the derivative in the earlier time is taken from the right. The right side vanishes in equilibrium, and by the Cauchy--Schwarz inequality its square is at most $I_H(s)/N$. After averaging over the disorder, we find that $|\bar r_N(t,s)-\partial_s\bar c_N(t,s)|^2\le a_N(s)$: the violation of FDT at the earlier time $s$ is controlled by the entropy production at time $s$. This is the mechanism of Cugliandolo, Dean, and Kurchan \cite[arXiv version~2, equations~(12)--(13), p.~3, and~(18), p.~4]{cdk1997}.

\paragraph{A single Fisher majorant.}
We need this bound for the CK solution at every pair of times. However, the budget \eqref{eq:overview-budget} bounds $a_N$ only in $L^1$, and the integrability of a limit of $a_N$ would not make it small at every late time. To overcome this, we turn to the Bochner formula \cite[Proposition~3, equation~(4b), p.~187]{bakryemery1985}, which shows that $e^{-Kt}I_H(t)$ is nonincreasing on an event of overwhelming probability for the disorder, with $K$ depending only on $p$ and $\beta$. We apply Helly's selection theorem to the Fisher informations restricted to this disorder event (compare the selection proof in \cite[Theorem~25.9, p.~336]{billingsley1995}), and use the landscape moments and entropy budget to show that the disorder outside this event contributes an $L^1$ mass that vanishes as $N\to\infty$. This gives one subsequence $(N_j)$, which we keep for the rest of the paper, along which the full normalized annealed Fisher information $a_{N_j}$ converges in $L^1_{\mathrm{loc}}$ to a function $a_F$, the Fisher majorant. We prove in \cref{thm:fisher-clock} that
\begin{equation}\label{eq:overview-fisher}
\begin{gathered}
 \int_0^\infty a_F(u)\,\dd u\le\frac{\beta^2}2,\qquad
 \lim_{T\to\infty}\sup_{s\ge T}a_F(s)=0,\\
 |\partial_sC(t,s)-R(t,s)|^2\le a_F(s),\qquad t>s>0,
\end{gathered}
\end{equation}
where the decay of $\sup_{s\ge T}a_F(s)$ follows from integrability and the monotonicity of $e^{-Kt}a_F(t)$. We choose the right-continuous representative of this limit and use the same $a_F$ below, including for the physical relaxation cost. One function thus dominates the backward defect $E(t,s)=\partial_sC(t,s)-R(t,s)$ on the whole triangle. In the same way, using a second current identity and the gradient bound for the semigroup, we show that a correlation at fixed lag changes little when both times are shifted. For $0\le\tau\le L$ and $0\le w_1\le w_2$, the difference $|C(w_2+\tau,w_2)-C(w_1+\tau,w_1)|$ is at most $B_L\int_{w_1}^{w_2}\sqrt{a_F(u)}\,\dd u$, where $B_L$ depends only on $p$, $\beta$, and $L$ (\cref{lem:waiting-clock-current}); compare \cite[arXiv version~2, equation~(15), p.~3]{cdk1997}.

\paragraph{Stationary clusters.}
As $a_F$ is integrable, the right sides of these bounds vanish on late windows of bounded length. Combining this with the precompactness of the translated kernels, we prove in \cref{prop:stationary-clusters} that every limit of $(C_{T_n},R_{T_n},\lambda_{T_n})$ along a sequence $T_n\to\infty$ has the form $c(|a-b|)$, $r(a-b)$, and $\ell$, where $c(0)=r(0)=1$, $0\le c\le1$, $r=-c'\ge0$, and the constant $\ell$ equals $1+p\beta e_\infty$ for the limit $e_\infty$ of the energies. Stationarity comes from the flatness bound, and FDT from integrating the Fisher bound over a bounded window. As the memory integral of the response equation runs only between the two observation times, we also obtain the closed stationary response equation
\begin{equation}\label{eq:overview-stationary-response}
 r'(\tau)=-\ell r(\tau)+\beta^2p(p-1)\int_0^\tau c(\tau-u)^{p-2}r(\tau-u)r(u)\,\dd u ,
\end{equation}
in which the coefficient $\beta^2p(p-1)$ is $6\beta^2$ for $p=3$. We call $(c,r,\ell)$ a stationary cluster and $q=\lim_{\tau\to\infty}c(\tau)$ its plateau. Note that \cref{prop:stationary-clusters} determines neither $q$ nor $\ell$, that different sequences $T_n$ may a priori lead to different clusters, and that it says nothing about $C(t,s)$ for fixed $s$.

\section{The scalar equation and its constraints}\label{sec:overview-scalar}

Integrating \eqref{eq:overview-stationary-response} once, we obtain the scalar memory equation of \cref{prop:scalar-clusters},
\begin{equation}\label{eq:overview-scalar}
 c'(t)=-1+\int_0^t\phi\bigl(c(t-u)\bigr)r(u)\,\dd u,\qquad
 \phi(y)=\ell-\beta^2p+\beta^2py^{p-1},
\end{equation}
which is the equation classified by Dembo, Guionnet, and Mazza \cite[arXiv version~1, Proposition~1.4, p.~5, and Section~5, pp.~17--20]{dgm2007}. We derive its consequences directly, because \cref{sec:stationary-selection} needs the constants and a strict inequality in exactly the following form. As $t\to\infty$, the integral tends to $\phi(q)(1-q)$, and the limit of $c'(t)$ must vanish because $c$ is bounded. This gives the plateau identity
\[
 \phi(q)(1-q)=1,\qquad\text{that is,}\qquad
 \ell=L(q)=\frac1{1-q}+\beta^2p\bigl(1-q^{p-1}\bigr).
\]
Next, $\phi$ is nondecreasing and $c$ is nonincreasing. Replacing $c(t-u)$ by $c(t)$ in \eqref{eq:overview-scalar}, we find that $c'(t)\ge F(c(t))$, where $F(y)=\phi(y)(1-y)-1$. The memory term of the stationary response equation is nonnegative, so $r(t)\ge e^{-\ell t}>0$. Then $c$ decreases strictly from $1$ to $q$, every $y\in(q,1)$ equals $c(t_y)$ for some $t_y$, and $F(y)\le c'(t_y)<0$. We have obtained the strict sign condition
\[
 F(q)=0,\qquad F(y)<0\quad(q<y<1).
\]
In particular, $F'(q)\le0$, which is the submarginality inequality $\beta^2p(p-1)q^{p-2}(1-q)^2\le1$. In the cubic model, these read $\ell=(1-q)^{-1}+3\beta^2(1-q^2)$, $F(y)=(\ell-3\beta^2+3\beta^2y^2)(1-y)-1$, and $6\beta^2q(1-q)^2\le1$. Finally, \cref{lem:stationary-uniqueness} states that, for a prescribed multiplier, the stationary system has at most one solution.

These are constraints on a possible cluster, not a selection of its plateau. If $q=0$, then $\ell=1+\beta^2p$ and the inequality says nothing. If $q>0$, then $\ell$ is determined by $q$, and the plateau is submarginal. To prove equality in the marginality relation, we need an inequality in the opposite direction. It will come from an upper bound on the energy of the part of a late state that was already determined at an earlier time.

\section{The causal tree realization and reflection}\label{sec:overview-tree}

\paragraph{The tree space.}
In the Green identity \eqref{eq:overview-green}, the first two terms form the inner product of $R(t,0)e_0+\sqrt2\,\mathbf1_{[0,t]}R(t,\cdot)$ and $R(s,0)e_0+\sqrt2\,\mathbf1_{[0,s]}R(s,\cdot)$ in $\R e_0\oplus L^2(\R_+)$, where a unit vector $e_0$ represents the initial condition. If $C(u,w)=\ip{v_u}{v_w}$, then $C(u,w)^m=\ip{v_u^{\otimes m}}{v_w^{\otimes m}}$, and the third term is an inner product as well. We therefore look for unit vectors such that
\begin{equation}\label{eq:overview-realization}
 v_t=R(t,0)e_0+\sqrt2\,\mathbf1_{[0,t]}R(t,\cdot)
 +\beta\sqrt p\int_0^tR(t,u)\,\mathcal I\bigl(v_u^{\otimes m}\bigr)\,\dd u,
\end{equation}
where $\mathcal I$ is an isometry from $m$-fold tensor powers into a space orthogonal to the first two terms. Iterating \eqref{eq:overview-realization}, we expand $v_t$ into terms indexed by finite rooted trees in which every vertex has either no children or $m$ ordered children (\cref{sec:tree-space}). Each leaf carries the initial vector or the time of a noise increment, and each internal vertex records one use of the memory integral. For example, when $p=3$, the component of $v_t$ indexed by the tree with two leaves, evaluated at two noise times $y_1,y_2\le t$, is
\[
 w_t(y_1,y_2)=2\sqrt3\,\beta\int_{y_1\vee y_2}^tR(t,u)R(u,y_1)R(u,y_2)\,\dd u .
\]
Noise received at times $y_1$ and $y_2$ is transported to a time $u$ by the response, combined there by the cubic interaction, and transported to time $t$; at a leaf that carries the initial vector, a factor $\sqrt2R(u,y_i)$ becomes $R(u,0)$ (\cref{ex:tree-first-heights}). Accordingly, $\mathcal H$ is an orthogonal sum, over trees, of tensor powers of $\R e_0\oplus L^2(\R_+)$, and $\mathcal I$ grafts $m$ trees onto a new root. We write $\Pi_j$ for the projection onto trees of height $j$, and $P_T$ for the prefix projection, which keeps each component only where all of its leaf times are at most $T$. As the leaves of a graft are the leaves of its subtrees, $P_T$ commutes with grafting.

\paragraph{The realization.}
In \cref{prop:geometry-tree-realization}, we construct a norm-continuous curve of unit vectors $v_t\in\mathcal H$ that solves \eqref{eq:overview-realization} and satisfies
\[
 \ip{v_t}{v_s}=C(t,s),\qquad P_tv_t=v_t,\qquad
 h(v_t)=e_\beta(t),\quad\text{where }h(v)=\sqrt p\,\ip{v}{\mathcal I(v^{\otimes m})}.
\]
We call $h$ the tree energy. Every component of $v_t$ is a nonnegative function of the leaf labels. The construction is a positive Picard iteration: the Gram kernels of the truncations at height $n$ increase with $n$, because $R$ and $C$ are nonnegative, and their limit solves the Green equation, whose solution with values in $[-1,1]$ is unique. Positive noncrossing expansions of the response equation appear in \cite[Section~2.1]{guionnetmazza2005} and \cite[Section~2]{dgm2007}; here we expand the covariance.

\paragraph{Prefix centers.}
For $T,a\ge0$, we call $M_T(a)=P_Tv_{T+a}$ a prefix center. It retains the components whose initial and thermal-noise leaf labels lie at or before $T$; their internal memory integrals may still extend beyond $T$. For $s\le T$, we have $P_Tv_s=v_s$, and therefore $\ip{v_{T+a}}{v_s}=\ip{M_T(a)}{v_s}$; that is, the correlation of a late state with any state up to time $T$ is carried by its prefix center. At first sight, the Gram kernel of the prefix centers is as hard to compute as $C$, since the centers are built from the whole older history. However, the older memory cancels exactly. Let $K_T(t,s)=\ip{P_Tv_t}{P_Tv_s}$. Applying $P_T$ to \eqref{eq:overview-realization}, we see that $K_T$ satisfies a Green identity in which the noise integral stops at $T$. If $u\le T$, then $P_Tv_u=v_u$, and therefore $K_T(u,w)=C(u,w)$. When we subtract the identity for $K_T$ from \eqref{eq:overview-green}, the initial terms, the noise integrals over $[0,T]$, and every block of the memory integral with $u\le T$ or $w\le T$ cancel. For $t,s\ge T$, only the future rectangle remains:
\[
\begin{aligned}
 C(t,s)-K_T(t,s)={}&2\int_T^{t\wedge s}R(t,u)R(s,u)\,\dd u\\
 &+\beta^2p\int_T^t\int_T^sR(t,u)R(s,w)\bigl\{C(u,w)^m-K_T(u,w)^m\bigr\}\,\dd w\,\dd u .
\end{aligned}
\]
In the variables $a=t-T$ and $b=s-T$, this is the closed equation \eqref{sel:closed-prefix-rectangle} for $Q_T(a,b)=\ip{M_T(a)}{M_T(b)}$ on a fixed rectangle. It holds for every $T$, before any limit is taken, and it uses no decay of the older memory.

\paragraph{Reflection.}
Let $T_n\to\infty$ be a sequence along which the translated kernels converge to a stationary cluster $(c,r,\ell)$. The coefficients of the closed equation converge, and by a Gronwall argument, every solution of the limiting equation with values in $[0,1]$ is the limit of $Q_{T_n}$. Comparing double Laplace transforms, and using only the stationary response equation, we check that $c(a+b)$ is such a solution. This proves the reflection identity of \cref{sel:prefix-reflection}: for every $A<\infty$,
\begin{equation}\label{eq:overview-reflection}
 \lim_{n\to\infty}\sup_{0\le a,b\le A}\bigl|\ip{M_{T_n}(a)}{M_{T_n}(b)}-c(a+b)\bigr|=0 .
\end{equation}
In particular, $\norm{M_{T_n}(L)}^2\to d_L=c(2L)$. We compute the energy of a center by the same kind of exact subtraction. When we subtract $h(v_t)=e_\beta(t)$ from $h(P_Tv_t)$, only an integral over the last $L$ units of time remains \eqref{sel:exact-prefix-energy-subtraction}. Passing to the limit, we find that
\begin{equation}\label{eq:overview-center-energy}
 h\bigl(M_{T_n}(L)\bigr)\longrightarrow
 h_L=e_\infty+\beta p\int_0^Lr(\tau)\bigl\{c(2L-\tau)^{p-1}-c(\tau)^{p-1}\bigr\}\,\dd\tau .
\end{equation}
These limits hold for every stationary cluster, whatever its plateau.

We may think of $M_T(a)$ as the average of the state at time $T+a$ over independent continuations of the thermal noise after time $T$. For a stationary reversible Markov process, the conditional means of two such continuations of lengths $a$ and $b$ have overlap $c(a+b)$ (\cref{rem:tree-independent-futures}); this reflected formula appears in the work of Barrat, Burioni, and M\'ezard \cite[Section~2, Equation~(9)]{bbm1996}. Averages over independent futures and time-averaged magnetizations \cite{cugliandolodean1995,franzvirasoro2000,biroli1999,ck1994} are discussed in \cref{sec:intro-related}. In our proof, the reflected overlap comes from the exact cancellation above.

\paragraph{Complete monotonicity.}
By \eqref{eq:overview-reflection}, the Hankel kernel $(a,b)\mapsto c(a+b)$ is a limit of Gram kernels, and it is positive semidefinite. Using a short semigroup argument in the spirit of the Bernstein--Widder theorem \cite{widder1941,bcr1984}, we deduce the representation of \cref{sel:cluster-complete-monotonicity},
\[
 c(t)=\int_{[0,\infty)}e^{-xt}\,\nu(\dd x),\qquad \nu(\{0\})=q,
\]
for a probability measure $\nu$. In particular, $c$ and $r$ are completely monotone. Franosch and Voigtmann proved complete monotonicity for Brownian mode-coupling equations of mixtures by a different route, an iteration argument \cite[arXiv version~1, Section~III]{fv2002}. Finally, after time $T$ no new noise leaf survives the projection $P_T$. By \cref{lem:geometry-projected-drift}, the curve $t\mapsto P_Tv_t$ is then differentiable on $(T,\infty)$, and its derivative is the image under $P_T$ of an explicit polynomial function of $v_t$. We use this in \cref{resp:chapter} to bound $\partial_tC$ far from the diagonal.

\section{Selection of the plateau from two threshold inputs}\label{sec:overview-selection}

\paragraph{The threshold inputs.}
Besides the fact that the correlation of a late state with an earlier state is carried exactly by the center, \cref{sec:stationary-selection} takes as input two facts about prefix centers, which we state as \cref{prop:threshold-inputs}. Recall that $h$ is homogeneous of degree $p$, so that $h_L/d_L^{p/2}$ is the limiting energy of the normalized center $M_{T_n}(L)/\norm{M_{T_n}(L)}$. Let $\Eth=2\sqrt{(p-1)/p}$ be the threshold energy of the pure $p$-spin landscape \cite[Equation~(2.13)]{abc2013}. Informally, \cref{prop:threshold-inputs} states the following.
\begin{enumerate}[label=(\alph*)]
\item For every stationary cluster with positive plateau, the limiting center energies satisfy $h_L\le\Eth\,d_L^{p/2}$ for every $L\ge0$.
\item If the translated kernels converge along every sequence $T\to\infty$ to one stationary profile, and if its plateau is positive and marginal, then the solution forgets every fixed earlier state: $\sup_{0\le s\le S}C(t,s)\to0$ as $t\to\infty$, for every $S<\infty$.
\end{enumerate}
We prove \cref{thm:main-aging} from these inputs in \cref{sec:stationary-selection}, and we prove the inputs in \cref{sec:tree-ceilings,app:general-selection}.

\paragraph{An algebraic identity.}
Since $c$ decreases to $q$, we have $d_L\to q$. Using dominated convergence in \eqref{eq:overview-center-energy} and $p\int_0^\infty c^{p-1}r=1-q^p$, we find that
\[
 h_L\longrightarrow h=e_\infty-\beta\bigl\{1-q^p-pq^{p-1}(1-q)\bigr\}.
\]
If $q>0$, inserting the plateau identity $\ell=L(q)$ and $e_\infty=(\ell-1)/(p\beta)$ into this formula, we obtain the identity of \cref{lem:center-energy},
\begin{equation}\label{eq:overview-amgm}
 \frac h{q^{p/2}}=\sqrt{\frac{p-1}p}\Bigl(z+\frac1z\Bigr)\ge\Eth,\qquad
 z=\beta\sqrt{p(p-1)}\,q^{(p-2)/2}(1-q).
\end{equation}
Equality holds if and only if $z=1$, that is, if and only if $\beta^2p(p-1)q^{p-2}(1-q)^2=1$. For $p=3$, we have $h=e_\infty-\beta\{1-q^3-3q^2(1-q)\}$, $z=\beta\sqrt{6q}\,(1-q)$, and $\Eth=2\sqrt{2/3}$. Note that this identity is algebra about a stationary cluster, and that it uses no threshold input.

\paragraph{Marginality.}
Let $(c,r,\ell)$ be a stationary cluster with $q>0$. Letting $L\to\infty$ in part~\ref{item:threshold-ceiling} of \cref{prop:threshold-inputs}, we find that $h\le\Eth q^{p/2}$, while \eqref{eq:overview-amgm} gives $h\ge\Eth q^{p/2}$. Equality must hold, and the plateau is marginal. This is the inequality in the opposite direction that the scalar constraints were missing. For $\beta>\beta_{\mathrm m}$, the marginality equation has two roots in $(0,1)$, and the strict sign condition excludes the smaller one. Indeed, at a marginal plateau we have $F(q)=F'(q)=0$, and if $q<q_{\mathrm m}=(p-2)/p$, then $F''(q)>0$, so that $F$ would be positive just above $q$. Hence $q=q_+(\beta)$, the root in $[q_{\mathrm m},1)$ (\cref{sel:positive-marginality}).

\paragraph{Convergence at every temperature.}
It remains to show that the cluster does not depend on the sequence $T_n$. By marginality and the plateau identity, the multiplier of every cluster is either $1+\beta^2p$, when $q=0$, or $L(q_+(\beta))$; it lies in a set with at most two elements. The function $\lambda$ is continuous and bounded. By the intermediate value theorem, every point between $\liminf_{t\to\infty}\lambda(t)$ and $\limsup_{t\to\infty}\lambda(t)$ is the limit of $\lambda(T_n)$ along some sequence $T_n\to\infty$, and by compactness, each such limit is the multiplier of a cluster. An interval contained in a set with at most two elements is a single point, so $\lambda$ converges. Once the multiplier is fixed, \cref{lem:stationary-uniqueness} identifies the profile, and precompactness upgrades subsequential convergence to convergence of the whole translated family (\cref{sel:all-temperature}). This holds for every $\beta>0$, and the energy converges as well.

\paragraph{The selected plateau.}
For $\beta\ge\bd$, we exclude the zero plateau with the same sign condition. If $q=0$, then $\ell=1+\beta^2p$ and $F(y)=y\{\beta^2py^{p-2}(1-y)-1\}$, and we find that
\[
 F(\qd)=\qd\Bigl(\frac{\beta^2}{\bd^2}-1\Bigr)\ge0,
\]
which contradicts the sign condition on $(0,1)$. For $p=3$, this reads $F(1/2)=(3\beta^2/4-1)/2$. At $\beta=\bd$, we have $F(\qd)=0$, and the contradiction comes from the strictness of the sign condition; the endpoint needs no continuity argument in $\beta$. We conclude that $q=q_\beta=q_+(\beta)\ge\qd$, and the energy formula \eqref{eq:main-energy} follows from $\ell=L(q_\beta)$ and marginality (\cref{sel:selected-sector}). For $\beta<\beta_{\mathrm m}$, only the zero plateau is possible. For $\beta_{\mathrm m}\le\beta<\bd$, we prove convergence, but we do not decide between the two alternatives (\cref{sel:intermediate-range}).

\paragraph{Forgetting and aging.}
If $q>0$, the convergence of the whole translated family and marginality are the hypotheses of part~\ref{item:threshold-forgetting} of \cref{prop:threshold-inputs}, which implies that $\sup_{s\le S}C(t,s)\to0$. If $q=0$, we need no threshold input. For $s\le t-L$, we have $C(t,s)=\ip{M_{t-L}(L)}{v_s}\le\norm{M_{t-L}(L)}$ by the Cauchy--Schwarz inequality. The right side tends to $c(2L)^{1/2}$ as $t\to\infty$, and $c(2L)\to0$ as $L\to\infty$. Using the domination $R\le C$, we transfer forgetting to the response (\cref{sel:fixed-wait-forgetting}). For $\beta\ge\bd$, the ordered limits \eqref{eq:main-aging} follow: the first inner limit is $c(\tau)$, which tends to $q_\beta>0$, and the second vanishes by forgetting (\cref{sel:aging}).

\paragraph{The order of the argument.}
But there is an apparent circularity here, since the threshold inputs are proved in later sections, and part~\ref{item:threshold-forgetting} assumes a marginal plateau, which is what part~\ref{item:threshold-ceiling} is used to prove. The way out is the order of deduction. Part~\ref{item:threshold-ceiling} concerns an arbitrary stationary cluster with positive plateau, and its proofs in \cref{sec:tree-ceilings,app:general-selection} use only \cref{sec:model,sec:ck-structure,sec:entropy-clusters,sec:scalar,sec:geometry-core}. We deduce marginality from part~\ref{item:threshold-ceiling} and the identity \eqref{eq:overview-amgm}, which uses no threshold input. Convergence at every temperature and marginality are then the hypotheses of part~\ref{item:threshold-forgetting}, from which we deduce forgetting.

\section{Why the threshold inputs hold}\label{sec:overview-inputs}

\subsection{The cubic tree ceiling}
For $p=3$, grafting is binary, the tree energy is $h=\sqrt3f$ with $f(v)=\ip{v}{\mathcal I(v\otimes v)}$, and $\Eth=\sqrt3F_*$ with $F_*=2\sqrt2/3$. We prove in \cref{thm:geometry-tree-ceiling} that $|f(v)|<F_*$ for every unit vector $v$, that $F_*$ is the supremum, and that, with $b_j=\norm{\Pi_jv}^2$,
\begin{equation}\label{eq:overview-height-defect}
 \sum_{j\ge0}b_j^2\le4\sqrt2\,\bigl(F_*-f(v)\bigr).
\end{equation}
The proof is a heightwise Cauchy--Schwarz inequality followed by a Riemann sum. Let $B_j=b_0+\dots+b_j$, with $B_{-1}=B_{-2}=0$. Since a graft is one level taller than its taller subtree, the component of $\mathcal I(v\otimes v)$ at height $j$ is built from the components of $v$ below height $j$, and its squared norm is $B_{j-1}^2-B_{j-2}^2=b_{j-1}s_j^2$, where $s_j=(B_{j-1}+B_{j-2})^{1/2}$ (\cref{lem:tree-graft-heights}). Using the Cauchy--Schwarz inequality at each height, and then $2\sqrt{b_jb_{j-1}}\le b_j+b_{j-1}$, we find that
\[
 f(v)=\sum_{j\ge1}\ip{\Pi_jv}{\Pi_j\mathcal I(v\otimes v)}
 \le\sum_{j\ge1}\sqrt{b_jb_{j-1}}\,s_j
 \le\frac12\sum_{j\ge1}(b_j+b_{j-1})s_j .
\]
We have $s_j\le\sqrt{2B_{j-1}}$ and $s_{j+1}\le\sqrt{2B_j}$, and the last sum is therefore at most the trapezoid sum of the concave function $\sqrt{2y}$ over the partition $0\le B_0\le B_1\le\dots$ of $[0,1]$. This trapezoid sum is at most
\[
 \int_0^1\sqrt{2y}\,\dd y=\frac{2\sqrt2}3=F_* ,
\]
and keeping track of the gaps in these steps gives \eqref{eq:overview-height-defect}. The supremum is approached by unit vectors with mass $1/(n+1)$ at each of the heights $0,\dots,n$, whose component at each height $j\ge1$ is a positive multiple of the graft of the lower components. For these vectors, the first inequality above is an equality, and $f(v)$ is the midpoint Riemann sum of the same integral over the first $n$ of the $n+1$ intervals of the partition; the omitted interval contributes at most $\sqrt2/(n+1)$.

The defect bound says more than $f(v)<F_*$. Together with the Cauchy--Schwarz inequality, it implies that $\norm{\Pi_{\le H}v}^2\le(4\sqrt2(H+1)(F_*-f(v)))^{1/2}$ for every $H\ge0$, where $\Pi_{\le H}=\Pi_0+\dots+\Pi_H$: a unit vector close to the ceiling has little mass at every fixed set of heights. The proof also retains the gaps of the two inequalities above, an angular defect and a Young defect \eqref{eq:geometry-angular-young-defects}, which we use in \cref{sec:geometry-results}. For $p=4$, a difference of cubes replaces the difference of squares, and the ceiling is $f(v)=\ip{v}{\mathcal I(v^{\otimes3})}<\sqrt3/2$ for unit vectors $v$, with $h=2f$ and $\Eth=\sqrt3$ (\cref{thm:geometry-quartic-tree-ceiling}). To prove part~\ref{item:threshold-ceiling} for $p\in\{3,4\}$, we apply the ceiling to the normalized centers $M_{T_n}(L)/\norm{M_{T_n}(L)}$ and let $n\to\infty$ using \cref{sel:prefix-reflection}. The ceilings are elementary inequalities in the tree space, and their proofs use neither the landscape results of \cite{abc2013} nor the algorithmic results of \cite{hs2021,hs2023}.

\subsection{Height escape and forgetting}
Under the hypotheses of part~\ref{item:threshold-forgetting}, the limits of \cref{sel:prefix-reflection} hold as $T\to\infty$ without passing to a sequence, and the normalized centers saturate the threshold. Indeed, the deficit $\delta_L=\Eth-h_L/d_L^{p/2}$ is nonnegative by part~\ref{item:threshold-ceiling}, and $h_L/d_L^{p/2}$ tends to $h/q^{p/2}$, which equals $\Eth$ by the equality case of \eqref{eq:overview-amgm}. Then $\delta_L\to0$, and by the defect bound, the mass of the centers leaves every fixed set of heights (\cref{sel:threshold-centers}):
\[
 \limsup_{t\to\infty}\norm{\Pi_{\le H}P_{t-L}v_t}^2\le d_L\bigl\{\kappa_p(H+1)\delta_L\bigr\}^{1/2},
\]
with $\kappa_3=4\sqrt{2/3}$ and $\kappa_4=\sqrt3$. An earlier state behaves in the opposite way. The curve $s\mapsto v_s$ is norm-continuous, so the set $\{v_s:s\le S\}$ is compact, and $\norm{(\Id-\Pi_{\le H})v_s}\to0$ as $H\to\infty$, uniformly in $s\le S$. For $s\le S\le t-L$, we have $C(t,s)=\ip{P_{t-L}v_t}{v_s}$, and splitting $v_s$ at height $H$, we find that
\[
 |C(t,s)|\le\norm{\Pi_{\le H}P_{t-L}v_t}+\norm{(\Id-\Pi_{\le H})v_s} .
\]
We first let $t\to\infty$ with $L$ and $H$ fixed, then let $L\to\infty$, and finally let $H\to\infty$. This proves part~\ref{item:threshold-forgetting} for $p\in\{3,4\}$.

\subsection{General pure models}
For general $p$, we have no deterministic ceiling in the tree space. In \cref{app:general-selection}, we prove both inputs for every $p\ge3$ by an argument in finite dimension, which does not use the tree ceilings.

\paragraph{Forks.}
In finite dimension, the analog of a prefix center is a mean over independent futures. At a late time $T$, we continue the trajectory by $K$ independent Brownian futures of length $L$, which we call a fork. We let $m_{N,T}$ be the average of their endpoints, the center of the fork. With a current identity of the kind used in \cref{sec:entropy-clusters}, we show that the overlap of two branches of lengths $a$ and $d$ concentrates near $c(a+d)$ along $T_n$; this is the reflection identity in finite dimension (\cref{sel:fork-moments}). We compute the energy of the center exactly by a Gaussian integration by parts, and subtracting the energy of a single branch removes the common past in the dimension limit. Unlike the cancellation of \cref{sec:geometry-core}, this one holds only in the limit, because on a fixed time horizon the overlaps of all branches with a state of the common past concentrate around the same value. As $N=N_j\to\infty$ and then $T=T_n\to\infty$, we find that (\cref{sel:finite-forks})
\[
 \frac{\norm{m_{N,T}}^2}N\longrightarrow Q_{K,L}=\frac1K+\Bigl(1-\frac1K\Bigr)c(2L),\qquad
 \frac{\E H_N(m_{N,T})}N\longrightarrow h_{K,L},
\]
where $h_{K,L}$ is explicit. Using the hard--soft approximation of Sellke \cite[Lemmas~3.1 and~3.2]{sellke2024}, we approximate the normalized center in mean square by Lipschitz functions of the disorder, with constants that do not depend on $N$. The algorithmic threshold of Huang and Sellke \cite[Theorem~1 and Corollary~1.8]{hs2023} applies to these approximations, and it implies that $h_{K,L}\le\Eth Q_{K,L}^{p/2}$; this follows the argument by which Sellke showed that the energy of the dynamics itself stays below the threshold \cite[Section~3]{sellke2024}. Letting $K\to\infty$, we have $Q_{K,L}\to d_L$ and $h_{K,L}\to h_L$, which proves part~\ref{item:threshold-ceiling}. Every horizon, number of branches, and approximation parameter is fixed while the dimension tends to infinity, and the limits are taken in the order
\[
 N_j\to\infty,\qquad n\to\infty,\qquad \varepsilon\downarrow0,\qquad K\to\infty,\qquad L\to\infty,\qquad M\to\infty
\]
of \eqref{sel:aging-limit-order}, where $\varepsilon$ is the accuracy of the Lipschitz approximation and $M$ is the number of steps of the Euler--Maruyama scheme below.

\paragraph{Adaptive disclosure.}
Part~\ref{item:threshold-forgetting} needs more. Under marginality, the normalized fork centers saturate the threshold as $K\to\infty$ and then $L\to\infty$ (\cref{lem:fork-parameter-limits}). If a fork ends at time $t$, the mean overlap of its center with the earlier state $X_s$ is the mean of $X_t\cdot X_s/N$, which tends to $C(t,s)$. We must therefore show that a Lipschitz algorithm close to the threshold is almost orthogonal to $X_s$. For directions independent of the disorder, this follows from the external-field version of the threshold \cite[Corollary~1.8]{hs2023}, which by convex duality implies the overlap penalty of \cref{sel:absolute-field-cost}. Let $\zeta_N$ be an auxiliary seed independent of the disorder tensor $J$, let $w=w(\zeta_N)$ satisfy $\norm w=\sqrt N$, and let $\mathcal A_N(J,\zeta_N)$ take values in the ball of radius $\sqrt N$. If, for almost every value of the seed, $J\mapsto\mathcal A_N(J,\zeta_N)$ is $L_{\mathrm{Lip}}$-Lipschitz, with a deterministic constant $L_{\mathrm{Lip}}$ that does not depend on $N$, then
\[
 \limsup_{N\to\infty}\Bigl\{\E\frac{H_N(\mathcal A_N)}N+\bd\,\E\Bigl|\frac{\ip w{\mathcal A_N}}N\Bigr|^p\Bigr\}\le\Eth .
\]
However, $X_s$ depends on the disorder. To overcome this, we disclose the disorder in stages (\cref{sel:adaptive-disclosure}). At each stage, we reveal the part of the tensor that determines the next gradient query. Conditionally on what has been revealed, the rest of the tensor is standard Gaussian on a smaller space, and the overlap penalty applies to the new query direction. By induction over finitely many stages, we show that an algorithm that saturates the threshold is asymptotically orthogonal to every span obtained in this way from a seed-measurable initial span of fixed dimension by a fixed number of gradient queries (\cref{sel:adaptive-orthogonality}). In the proof of \cref{sel:general-fixed-wait}, we apply this only on the fixed fork product law, whose seed consists of $X_0$, the full trunk Wiener path, and the $K$ full future Wiener paths, jointly independent of $J$. For its Euler--Maruyama scheme with $M$ steps, the initial span is generated by $X_0$ and the $M$ trunk Brownian increments, and the scheme adds at most $M$ gradient queries. The resulting span contains a vector approximating this fork's trunk state $X_s$, with dimension-limit normalized mean-square error at most $Cs/M$, where $C$ depends only on $p$, $\beta$, and $s$. We take the limits in \eqref{sel:aging-limit-order}, using $\varepsilon=1/(m+1)$, and let $M\to\infty$ last. This proves that $C(t,s)\to0$ for every fixed $s$ (\cref{sel:general-fixed-wait}), uniformly in $s\le S$ by the modulus of continuity of $C$. Adaptive Gaussian conditioning of this kind goes back to Bolthausen \cite[Section~3]{bolthausen2014} and to Bayati and Montanari \cite[arXiv version~4, Section~3.8]{bayatimontanari2011}. Conditioning of the disorder tensor on adaptive gradient observations was developed for message passing in spin glasses in \cite[arXiv version~1, Appendix~A]{eams2021} and extended to multi-species spherical models in \cite[Appendix~A]{hs2024amp}; we combine it with the external-field threshold.

\section{Response localization and persistent violation}\label{sec:overview-response}

\paragraph{The remote response.}
\Cref{thm:main-aging} controls the response only on bounded windows of lags, while a row $s\mapsto R(t,s)$ has length $t$. Convergence on every bounded window does not exclude a small contribution spread over the remote part $[0,t-L]$ of the row. In \cref{resp:chapter}, we first prove that the remote part carries no squared mass:
\begin{equation}\label{eq:overview-remote}
 \lim_{L\to\infty}\limsup_{t\to\infty}\int_0^{t-L}R(t,u)^2\,\dd u=0 .
\end{equation}
Again we use the prefix centers. The component of $v_{T+L}$ at height zero is the leaf vector $R(T+L,0)e_0+\sqrt2\,\mathbf1_{[0,T+L]}R(T+L,\cdot)$, and $P_T$ restricts its noise part to $[0,T]$. This yields
\[
 2\int_0^TR(T+L,u)^2\,\dd u\le\norm{\Pi_0M_T(L)}^2\le\norm{M_T(L)}^2\longrightarrow c(2L)\qquad(T\to\infty),
\]
which proves \eqref{eq:overview-remote} when $q=0$. If $q>0$, then $c(2L)\to q$, and we need to show that the normalized centers have vanishing mass at height zero. For $p=3$, this follows from the term $j=0$ of \eqref{eq:overview-height-defect}, because the normalized centers saturate the threshold. For general $p$, we test the center of a fork against the Gaussian vector $G_N=\int_0^Tk(u)\,\dd B_u$, built from the Brownian path with a weight $k$ proportional to $R(T+L,\cdot)$. A weighted form of the Brownian response identity shows that the mean overlap of $m_{N,T}$ with $G_N$ tends to $(2\int_0^TR(T+L,u)^2\,\dd u)^{1/2}$. Since $G_N$ is independent of the disorder, the overlap penalty forces this overlap to be small, with no adaptive disclosure (\cref{prop:remote-response-general}).

\paragraph{Whole rows.}
Together with convergence on recent strips, \eqref{eq:overview-remote} implies that the lag rows $v\mapsto\mathbf1_{[0,t]}(v)R(t,t-v)$ converge to $r$ in $L^2(0,\infty)$, and by the Cauchy--Schwarz inequality, we have $\chi(t,0)=o(\sqrt t)$ (\cref{resp:strong-row}). By the backward response equation, we have $\partial_sR\le\Lambda R$, so that a large remote value $R(t,s)$ would force squared mass on $[s-1,s]$; this upgrades the convergence to the uniform norm (\cref{resp:uniform-row}). For the earlier derivative, we differentiate the Green identity in its earlier time. On a bounded interval of earlier times, forgetting makes every term vanish as $t\to\infty$, except a colored term $\int_0^tR(t,u)C(u,s)^{p-1}\,\dd u$, which we control by positivity and a modulus of continuity uniform in $t$: the Green identity bounds $C(t,s)$ from below by a weighted integral of this term over earlier times, with positive weights; forgetting makes that integral vanish, and the uniform modulus in \eqref{resp:colored-modulus} then forces the term to vanish pointwise. This shows that $E(t,s)\to0$ for every fixed $s>0$. Since $|E(t,s)|^2\le a_F(s)$ with an integrable $a_F$, dominated convergence over the whole row then shows that $\int_0^t|E(t,s)|^2\,\dd s\to0$ (\cref{resp:backward-fdt}). This is where we need one majorant for the whole triangle. For every $p\ge3$, the projected drift and the reflection identity also bound $\partial_tC$ far from the diagonal (\cref{resp:forward-derivative}). More precisely, for every $L>0$,
\[
 \limsup_{t\to\infty}\sup_{0\le s\le t-L}|\partial_tC(t,s)|
 \le\sqrt{c''(2L)}\le\frac{\sqrt{1-q}}{eL},
\]
and, using convergence on bounded lag windows, we obtain
\[
 \lim_{t\to\infty}\sup_{0\le s\le t}|\partial_tC(t,s)+r(t-s)|
 =\lim_{t\to\infty}\sup_{0\le s\le t}|\partial_tC(t,s)+R(t,s)|=0.
\]
The diagonal derivative is taken from the causal side and equals $-1$, while $c'(0)=-1$ is a right derivative. These conclusions about forward FDT hold at every positive temperature. For every $p\ge3$, we also deduce \eqref{eq:main-thermal} from complete-row convergence (\cref{resp:thermal}).

\paragraph{Persistent integrated violation.}
These statements about squared norms are compatible with a persistent violation of FDT. Let $\beta>\bd$, and recall that $x=(p-2)(1-q)/q\in(0,1)$. The energy identity fixes one weighted integral of the response over the whole row,
\[
 p\int_0^tC(t,u)^{p-1}R(t,u)\,\dd u=\frac{e_\beta(t)}\beta\longrightarrow1-d_*,
\]
where $d_*=1-e_\beta(\infty)/\beta=(1-x)q^p$ by the energy formula \eqref{eq:main-energy}. A recent strip of length $A$ contributes $p\int_0^Ac^{p-1}r=1-c(A)^p$, as it would for an equilibrium profile. The rest of the row carries the weighted response $c(A)^p-d_*$, while $C(t,\cdot)^p$ rises across it from $C(t,0)^p\to0$ to $C(t,t-A)^p\to c(A)^p$. It follows that the ratio in \eqref{eq:main-weighted-ratio} tends to $1-d_*/c(A)^p$, and then to $x$ as $A\to\infty$ (\cref{resp:weighted-ratio}). In the same way, we find the exact identity
\[
 p\int_0^tC(t,u)^{p-1}E(t,u)\,\dd u=1-C(t,0)^p-\frac{e_\beta(t)}\beta\longrightarrow d_*>0 .
\]
To remove the weight $C^{p-1}$, we prove the transport inequality of \cref{resp:transport}: if $z\ge0$ is absolutely continuous on $[u,v]$, $r_0\ge0$, and $J=\int_u^vz^{p-1}r_0$, then $\int_u^v(z'-r_0)_+\ge[z(v)-(z(u)^p+pJ)^{1/p}]_+$. Applying this on $[S,t-L]$ with $z=C(t,\cdot)$ and $r_0=R(t,\cdot)$, we find that, in the limit $t\to\infty$, the positive part of the defect has integral at least $m_L=c(L)-[c(L)^p-d_*]^{1/p}$ there. As $L\to\infty$, $m_L$ increases to $m_\infty=q(1-x^{1/p})$ (\cref{resp:persistent-defect}). This lower bound holds on an interval whose length is of order $t$, and by the Cauchy--Schwarz inequality, we have $\liminf_{t\to\infty}t\int_0^tE(t,u)_+^2\,\dd u\ge m_\infty^2$. This is the coexistence of small pointwise violation and integrated violation over an expanding interval described by Cugliandolo, Dean, and Kurchan \cite{cdk1997}.

\section{Three-time geometry, susceptibility, and the Fisher cost}\label{sec:overview-geometry}

In the cubic model, we use more of the tree geometry than the energy ceiling. Throughout this section, $p=3$.

\paragraph{Saturated tree vectors.}
In \cref{sec:geometry-results}, which is deterministic, we study unit vectors $v\in\mathcal H$ with a small deficit $\delta(v)=F_*-f(v)$. Using the retained defects, we bound weighted sums of the variations of the height amplitudes $\norm{\Pi_jv}$ and of the angular errors between each component $\Pi_jv$ and the graft $\Pi_j\mathcal I(v\otimes v)$ of the lower components. We deduce that the root operator $\mathcal I^*v$, viewed as a Hilbert--Schmidt operator on $\mathcal H$, is then close to an operator whose coefficients in the normalized height directions depend only on the height masses (\cref{lem:geometry-coherent-root}). In the coordinate $x\in(0,1)$ given by the cumulative mass, these operators converge to one universal operator. More precisely, if $\norm{v_n}=1$ and $f(v_n)\to F_*$, there exist isometric embeddings $V_n$ of $\mathcal H$ into $L^2(0,1)\oplus\mathcal K$, where $\mathcal K$ is a fixed auxiliary Hilbert space, such that $V_nv_n=\mathbf1$ and
\[
 V_n\mathcal I^*v_nV_n^*\longrightarrow\mathcal T\oplus0\quad\text{in Hilbert--Schmidt norm},
\]
where $\mathcal T$ is the integral operator with kernel $T(x,y)=(2\max(x,y))^{-1/2}$ (\cref{thm:geometry-universal-root}). The operator $\mathcal T$ is positive and trace class, with $\tr\mathcal T=\sqrt2$ and $\ip{\mathbf1}{\mathcal T\mathbf1}=F_*$. For two vectors with nonnegative inner products at every height, the same defects show that little of their overlap is carried by heights at which a smoothed cosine between their prefixes is far from one (\cref{lem:geometry-pair-deficit}). For three such vectors, the prefixes that carry the overlaps of two of the pairs are nested, which leads to the product inequality of \cref{thm:geometry-threshold-product},
\[
 \ip{v_1}{v_2}\ge\ip{v_1}{v_3}\ip{v_2}{v_3}-\omega\Bigl(\max_i\delta(v_i)\Bigr),\qquad
 \lim_{\delta\downarrow0}\omega(\delta)=0 .
\]

\paragraph{Three times.}
Let $\beta\ge\bd$, including $\beta=\bd$, where $q=1/2$. To apply these results to the CK solution, we use the normalized centers $U_{s,D}=M_{s,D}/\norm{M_{s,D}}$ of a fixed smoothing age $D$, where $M_{s,D}=P_{s-D}v_s$. Their components are nonnegative, and by \cref{sel:threshold-centers} their deficits are small when $D$ is large. The obstacle is to relate their overlaps to the correlations. A center is close to the average $D^{-1}\int_{s-D}^sv_u\,\dd u$ of the trajectory over the preceding window (\cref{lem:geometry-center-average}), and the overlap of two such averages is close to $C(t,s)$ when the derivatives of $C$ are small far from the diagonal, which is what \cref{resp:chapter} provides. We find that $\ip{U_{t,D}}{U_{s,D}}\approx C(t,s)/q$ for late pairs whose distance is large compared with $D$ (\cref{prop:geometry-center-gram}), and the product inequality becomes \cref{thm:geometry-actual-product}: for late triples $u<s<t$ with diverging gaps,
\[
 qC(t,u)\ge C(t,s)C(s,u)-o(1).
\]
This is a one-sided version of the spherical product relation of Cugliandolo and Kurchan \cite[arXiv version~1, Section~5.1, Equation~(5.1.8)]{ck1994}, which we prove for the actual CK solution without any assumption on the time scale of aging. It implies composition bounds for the last times at which the correlation crosses given levels (\cref{cor:geometry-last-crossing}).

\paragraph{The susceptibility.}
For the bound on $\chi(t,0)$, we use the root operator of a center. For $p=3$, the root frame \eqref{eq:geometry-prefix-root-frame} reads
\[
 \mathcal I^*P_Tv_t=\beta\sqrt3\int_0^tR(t,u)\,|P_Tv_u\rangle\langle P_Tv_u|\,\dd u,
\]
so that its trace is a weighted response mass. We take $t_n\to\infty$ along which $\chi(t_n,0)$ tends to its lower limit, ages $L_n\to\infty$, and the normalized centers $U_n$ of $M_{T_n}(L_n)$, where $T_n=t_n-L_n$; these centers approach the ceiling. We split their root operators into an old part $A_n$, the integral over $u\le T_n$, and a recent part $B_n$. On recent windows of fixed length, the projected states $P_{T_n}v_u$ are close to $\sqrt q\,U_n$ by reflection, and in the coordinates of the universal operator, $B_n$ converges to $\beta\sqrt3\sqrt q\,(1-q)|\mathbf1\rangle\langle\mathbf1|$, which equals $2^{-1/2}|\mathbf1\rangle\langle\mathbf1|$ by marginality. By \cref{thm:geometry-universal-root}, the old part then converges in Hilbert--Schmidt norm to $\mathcal T-2^{-1/2}|\mathbf1\rangle\langle\mathbf1|$, a positive operator with trace $2^{-1/2}$.

However, convergence in Hilbert--Schmidt norm does not give convergence of the trace. To overcome this, we rely on the lower semicontinuity of the trace on positive operators, after replacing each old term $|v_u\rangle\langle v_u|$ by $|M_{u,D}\rangle\langle M_{u,D}|$, whose norm converges. By causality, the errors on time strips at distance at least $D$ from each other are orthogonal (\cref{lem:geometry-disjoint-strips}), and their total is controlled by the squared remote response, which vanishes by \eqref{eq:overview-remote}. The trace of the smoothed old part is asymptotically $\beta\sqrt3\,q^{-1/2}c(2D)$ times the old response mass $\chi(t_n,0)-(1-q)$, and comparing it with $2^{-1/2}$, we find that
\[
 \liminf_{t\to\infty}\chi(t,0)\ge(1-q)+\frac{q(1-q)}{c(2D)}\longrightarrow2(1-q)\qquad(D\to\infty),
\]
which is \cref{thm:geometry-susceptibility-lower}. This holds for every $\beta\ge\bd$, including the endpoint, but we have no matching upper bound.

\paragraph{A Fisher cost.}
In \cref{resp:fisher-chapter}, we combine the same Fisher majorant $a_F$ from \eqref{eq:overview-fisher} with the product inequality. Let $\beta>\bd$, so that $q>1/2$, and consider windows $[s_n,t_n]$ with $s_n\to\infty$, $t_n-s_n\to\infty$, and $C(t_n,s_n)\to z\in(0,q)$, such that
\[
 \varepsilon_n^2=(t_n-s_n)\int_{s_n}^{t_n}a_F(u)\,\dd u\longrightarrow\varepsilon^2\le1 .
\]
By the Cauchy--Schwarz inequality and $|E(w,u)|^2\le a_F(u)$, the integrated backward defect on every subinterval of the window is at most $\varepsilon_n$; in particular, the response mass of the window is $1-C(t_n,s_n)$ up to $\varepsilon_n$. Integrating the backward response equation \eqref{eq:backward-response} over the window, we find that $1-\ell(1-C(t_n,s_n))+3\beta^2I_n$ is at most of order $\varepsilon_n$ in the limit, where
\[
 I_n=\int_{s_n}^{t_n}\bigl[1-C(u,s_n)^2\bigr]R(t_n,u)\,\dd u .
\]
In the bulk of the window, we have $C(t_n,u)C(u,s_n)\le qC(t_n,s_n)+o(1)$ by the product inequality. As the row $u\mapsto C(t_n,u)$ has derivative $R(t_n,u)+E(t_n,u)$, a clipped primitive turns this into an explicit lower bound for $I_n$, and we compute the recent strip from stationary convergence. With the calibration $6\beta^2q(1-q)^2=1$ and $\ell=(1-q)^{-1}+3\beta^2(1-q^2)$, the two estimates are compatible only if
\[
 3\beta^2(2q-1)(q-z)^2\le(\ell+9\beta^2)\varepsilon+6\beta^2\varepsilon^2
\]
(\cref{resp:fisher-obstruction}), and the left side is positive because $q>1/2$. This strict margin vanishes at $\beta=\bd$, which is why the positive Fisher cost is restricted to $\beta>\bd$, even though the product inequality includes the endpoint. For $\varepsilon=0$, this is impossible, which excludes every limit $z\in(0,q)$ of $C(t_n,s_n)$ along windows on which $\tau\int a_F$ tends to zero. Limits above $q$ are excluded by the remote envelope \eqref{eq:geometry-remote-envelope}. We also exclude a zero limit: the first crossing of a fixed level $z'\in(0,q)$ occurs inside the window, at a lag that tends to infinity, and the crossing windows have no larger Fisher cost, so that the obstruction applies to them. Every such window therefore ends at the plateau. At the first passage below a level $c\in(0,q)$, the inequality bounds $\varepsilon$ from below, which gives \eqref{eq:main-passage-cost} (\cref{resp:fisher-cost}).

\section{The stationary power law}\label{sec:overview-power-law}

For every $p\ge3$, let $\beta\ge\bd$. In \cref{sec:stationary-relaxation}, we study the scalar equation \eqref{eq:overview-scalar} for the selected profile, using only this equation with its plateau identity, marginality, and complete monotonicity. The function $u(z)=[c((1-q)z)-q]/(1-q)$ satisfies (\cref{lem:sc-normalized-equation})
\[
 u'+u+M(u)*u'=0,\qquad u(0)=1,
\]
where $*$ is the convolution on $[0,\infty)$ and $M$ is a polynomial with nonnegative coefficients such that $M(y)=y+\vartheta y^2+O(y^3)$, with $\vartheta=(p-2)(1-q)/(2q)$. For $p=3$, we have $M(y)=y+\vartheta y^2$. The linear coefficient of $M$ is exactly one, by marginality. For a pure power $w(t)=t^{-a}$, the beta integral gives
\[
 (w*w)'=\frac{\Gamma(1-a)^2}{\Gamma(1-2a)}\,w^2 ,
\]
so that the exponent relation \eqref{eq:main-tail-exponent} says that the two quadratic terms of the equation balance for $w$. This is the classical exponent relation of mode-coupling theory \cite[Section~4, p.~881]{goetze1998}, in its version for the aging state \cite[Equation~(IV.6)]{bckm1996} and \cite{kimlatz2001}. As the balance is quadratic, it fixes the exponent but not the amplitude. The proof of an amplitude has three further steps.

\paragraph{Comparison.}
Using a comparison principle (\cref{lem:sc-comparison}), we turn the sign of the residual of an explicit profile into an inequality between that profile and $u$. For the shifted power $g(t)=(1+t)^{-a}$, the choice of $a$ makes the quadratic part of the residual positive but smaller than $g^3$. The sign of the residual is then decided by cubic terms. We fix it by correcting $g$ with a small multiple of $g^2$ (\cref{lem:sc-shifted,lem:sc-critical-residuals}). Our comparison profiles keep $u$ itself up to a fixed time and continue it by a rescaled corrected power, which absorbs the memory of the initial segment. We obtain $A_-(1+t)^{-a}\le u(t)\le A_+(1+t)^{-a}$ with constants $0<A_-\le A_+<\infty$ (\cref{thm:stationary-critical-endpoint}). These bounds determine the exponent, but they still allow $t^au(t)$ to oscillate between $A_-$ and $A_+$ on a logarithmic time scale.

\paragraph{A logarithmic Dirichlet form.}
The convolution square $u*u$ satisfies the quadratic balance $(u*u)'=\vartheta u^2+\mathcal E$, where the error $\mathcal E$ is integrable against the weight $t^{-(1-2a)}$ (\cref{lem:stationary-critical-balance-error}). We integrate the balance against this weight and compare the result with the Schur test for a homogeneous kernel, keeping its exact remainder. We find that the function $D(x)=e^{ax}u(e^x)$ has finite energy for an explicit Dirichlet form on the line (\cref{lem:stationary-critical-dirichlet}), which implies that $D(x+y)-D(x)\to0$ as $x\to\infty$, for every fixed $y$. This still allows $D$ to vary slowly without converging.

\paragraph{A renewal equation.}
We take the Laplace transform of the balance at the point $e^{-x}$. Since the product of independent gamma and beta variables with suitable parameters has a gamma distribution (\cref{lem:sc-gamma-beta}), the transformed balance becomes a linear renewal equation for a smoothed version of $D^2$. Its source is integrable: the variance part of the source is controlled by the Dirichlet form, and the rest by $\mathcal E$. Using the key renewal theorem \cite[Chapter~XI, Section~1]{feller1971}, with a renewal density that we control near zero and at infinity (\cref{lem:log-beta-renewal-density}), we show that the smoothed function converges. Its difference from $D^2$ tends to zero by the second step, so that $D$ converges, and the bounds of the first step make the limit positive and finite. This proves that $c(t)-q\sim A_\beta t^{-a}$. Finally, complete monotonicity and the monotone density theorem \cite[Section~1.7.3]{bgt1987} give the asymptotics of every fixed derivative, without differentiating an asymptotic relation (\cref{thm:stationary-critical-amplitude}). The tools are classical, and the structural theory of these memory equations is developed in \cite{gs1995,fv2002,sjogren2003}; we apply them to the stationary profile selected by the CK solution.

\section{Dependencies between the sections}\label{sec:overview-dependencies}

\Cref{tab:main-results-map} lists, for each main theorem, the section in which it is proved and the main results that its proof uses. We add three remarks on the order of reading. First, for \cref{thm:main-aging}, a reader interested in $p\in\{3,4\}$ may read \cref{sec:tree-ceilings} and skip \cref{app:general-selection}, and a reader interested in general $p$ may skip \cref{sec:tree-ceilings}, since \cref{app:general-selection} does not use the tree ceilings. For \cref{thm:main-response} with $p\ne3$ and a positive plateau, \cref{resp:appendix-p} uses the forks and the overlap penalty of \cref{app:general-selection}. Second, \cref{sec:geometry-results} uses only the tree space and the cubic ceiling with its retained defects. Third, \cref{part:relaxation} uses only \cref{prop:stationary-clusters,prop:scalar-clusters,sel:selected-sector,sel:cluster-complete-monotonicity}, and the power law is not used in the proofs of the other main theorems; it serves only to verify the power-law hypotheses in the last assertions of \cref{resp:remote-earlier,resp:thermal}.


\part{The Cugliandolo--Kurchan system and its stationary limits}\label{part:foundations}
\chapter{The model and the hard-spherical Cugliandolo--Kurchan equations}\label{sec:model}

In this section, we fix the model, the dynamics, and our conventions, and state the hard-spherical Cugliandolo--Kurchan (CK) equations together with the published results that connect them with the Langevin dynamics. We prove three statements. \Cref{prop:hard-bridge} transfers Sellke's convergence theorem for the hard-spherical diffusion on fixed time horizons \cite[arXiv version~5, Theorem~4.3, p.~23]{sellke2024} to our normalization and to a uniform initial point. \Cref{lem:brown-force-response} identifies the Brownian observable in that theorem with the integrated response to an additive force. \Cref{lem:landscape-moments} collects bounds on the Hamiltonian and its derivatives that are uniform in the dimension. This section also uses the variational-response estimates from Steps~1 and~2 of the proof of \cref{prop:energy-bridge}, which depend only on \cref{lem:landscape-moments}: \cref{rem:fdt-convention} uses the actual $L^2$ derivative and the quenched moment estimate \eqref{eq:tangent-flow-bound}, and the proof of \cref{lem:brown-force-response} uses \eqref{eq:finite-lag-flow-moments}.

These results are the finite-dimensional input of the paper, and we use them throughout \cref{sec:ck-structure}. The landscape bounds enter again in the entropy estimates of \cref{sec:entropy-clusters} and in the branched trajectories of \cref{app:general-selection}, and we use the bridge wherever we obtain a statement about the CK solution as a limit of the diffusion, in \cref{sec:entropy-clusters,app:general-selection,resp:chapter}.

Throughout the paper, $p\ge3$ and $\beta>0$ are fixed, and $\xi(q)=q^p$.

\section{The Langevin dynamics}\label{sec:langevin}

For $N\ge2$, let $S_N=\{x\in\R^N:\norm{x}^2=N\}$, and let $\omega_N$ be the uniform probability measure on $S_N$. The disorder is an ordered coefficient tensor $J=(J_{i_1\cdots i_p})_{i_1,\dots,i_p=1}^N$ whose $N^p$ entries are independent standard Gaussian random variables. The Hamiltonian is
\begin{equation}
 H_N(x)=N^{-(p-1)/2}\sum_{i_1,\ldots,i_p=1}^N
 J_{i_1\cdots i_p}x_{i_1}\cdots x_{i_p},\qquad x\in\R^N.
 \label{eq:hamiltonian}
\end{equation}
Since the entries of $J$ are independent and the tensor is not symmetrized, we have
\begin{equation}\label{eq:hamiltonian-covariance}
 \E H_N(x)H_N(y)
 =N^{-(p-1)}\sum_{i_1,\ldots,i_p=1}^N\prod_{j=1}^px_{i_j}y_{i_j}
 =N^{-(p-1)}(x\cdot y)^p
 =N\xi\Bigl(\frac{x\cdot y}{N}\Bigr)
\end{equation}
for all $x,y\in\R^N$. No factor $p!$ appears, because the coefficients are ordered.

We write $\nabla H_N$ and $\nabla^2H_N$ for the Euclidean gradient and Hessian. For $x\in\R^N$, let $P_x=I-xx^{\mathsf T}/N$. For $x\in S_N$, the matrix $P_x$ is the orthogonal projection onto the tangent space $T_xS_N=\{v\in\R^N:v\cdot x=0\}$. We write $\nabla_S$, $\nabla_S^2$, and $\Delta_S$ for the gradient, the Riemannian Hessian, and the Laplace--Beltrami operator of $S_N$, with the metric induced from $\R^N$. Note that $\nabla_SH_N(x)=P_x\nabla H_N(x)$. We also write
\[
 \kappa_N=\frac{N-1}{N}.
\]

We consider the diffusion
\begin{equation}
 \dd X_t=\beta P_{X_t}\nabla H_N(X_t)\,\dd t
 +\sqrt2P_{X_t}\,\dd B_t-\kappa_NX_t\,\dd t,
 \label{eq:hard-sde}
\end{equation}
where $B$ is a standard Brownian motion in $\R^N$. The initial point $X_0$ has law $\omega_N$, and $X_0$, $J$, and $B$ are independent. The equation has a unique strong solution up to its explosion time, because its coefficients are polynomials. The solution stays on $S_N$. By It\^o's formula, $Z_t=\norm{X_t}^2-N$ satisfies a linear equation $\dd Z_t=Z_t\,\dd A_t$ for a continuous semimartingale $A$, because $x\cdot P_xv=-(\norm x^2-N)(x\cdot v)/N$ and $2\tr(P_x^2)-2\kappa_N\norm x^2$ vanishes on $S_N$; since $Z_0=0$, we have $Z\equiv0$, and there is no explosion.

For a smooth function $F$ on $\R^N$ and $x\in S_N$, It\^o's formula shows that the generator of \eqref{eq:hard-sde} is
\begin{equation}\label{eq:generator}
 \mathcal L_NF(x)=\tr\bigl(P_x\nabla^2F(x)\bigr)-\kappa_N\,x\cdot\nabla F(x)
 +\beta\,\nabla_SH_N(x)\cdot\nabla F(x).
\end{equation}
The first two terms are $\Delta_SF(x)$, by the classical formula for the Laplacian of a sphere of radius $\sqrt N$, and the last one is $\beta\,\nabla_SH_N\cdot\nabla_SF$. This shows that $\mathcal L_N=\Delta_S+\beta\nabla_SH_N\cdot\nabla_S$. Equivalently, $\mathcal L_NF=g_N^{-1}\operatorname{div}_S(g_N\nabla_SF)$ with
\[
 g_N=Z_{N,\beta}^{-1}e^{\beta H_N},\qquad Z_{N,\beta}=\int e^{\beta H_N}\,\dd\omega_N,
\]
so that the Gibbs measure $g_N\,\dd\omega_N$ is reversible. The sign convention is gradient ascent, in the sense that the drift pushes $X_t$ toward large values of $H_N$, and the limiting energies of this paper are positive. Because the noise amplitude is $\sqrt2$, the diffusion part of $\mathcal L_N$ is the Laplacian itself.

We write
\[
 C_N(t,s)=\frac{X_t\cdot X_s}N,\qquad s,t\ge0,
\]
for the empirical correlation.

\paragraph{Response.}
We measure the response to an additive tangent force in the drift, without a factor $\beta$. Write
\[
 b(x)=\beta P_x\nabla H_N(x)-\kappa_Nx,\qquad
 \sigma_j(x)=\sqrt2P_xe_j,\quad j=1,\dots,N.
\]
On each finite horizon, the response to a bounded continuous deterministic force $h$ is the zero-initial-state strong solution of the forced linearized equation
\begin{equation}\label{eq:perturbed-sde}
 \dd V_t=Db(X_t)[V_t]\,\dd t
 +P_{X_t}h(t)\,\dd t
 +\sum_{j=1}^ND\sigma_j(X_t)[V_t]\,\dd B^j_t,
 \qquad V_0=0.
\end{equation}
It is constructed on each finite horizon in the space of predictable square-integrable processes. More generally, the predictable $L^2$ construction provides an affine variational strong solution for arbitrary predictable square-integrable drift and noise sources. We state the response-kernel representation below for bounded continuous deterministic forces and bounded continuous adapted random directions.
For each fixed waiting time $u\ge0$ and finite horizon $H\ge u$, the response matrix $U_N(t,u)$, $u\le t\le H$, is the continuous adapted matrix whose columns solve the homogeneous variational equation along the actual trajectory $X$:
\begin{equation}\label{eq:response-matrix}
\begin{aligned}
 \dd\bigl(U_N(t,u)v\bigr)
 &=Db(X_t)[U_N(t,u)v]\,\dd t
 +\sum_{j=1}^ND\sigma_j(X_t)[U_N(t,u)v]\,\dd B^j_t,
 \\
 U_N(u,u)v&=P_{X_u}v,
\end{aligned}
\end{equation}
for $v\in\R^N$. The equation is constructed from the $N$ coordinate columns; linearity and pathwise uniqueness identify every direction with their linear combination. The choices on different finite horizons agree on their common time interval, on one event for all vectors. The matrix acts as zero in the normal direction $X_u$ and maps every vector into $T_{X_t}S_N$. For each fixed $H>u$, its terminal values are also the Fr\'echet derivatives into $L^2(\Prob_B;\R^N)$ of the nonlinear solution restarted at $X_u+P_{X_u}v$, evaluated at $v=0$. For this ambient restart, choose smooth compactly supported extensions agreeing with the displayed coefficients on a neighborhood of the sphere; the variational equation along $X$ uses those physical coefficients. This definition needs no simultaneous samplewise smoothness in all starting points. For classical background on smooth spatial stochastic flows, see \cite[Section~1, pp.~1022--1023]{kunita1987icm} and \cite[Section~4.6, Theorem~4.6.5, p.~173]{kunita1990}.

For every bounded continuous deterministic force $h$, the forced variational equation has the Bochner integral representation in $L^2(\Prob_B;\R^N)$,
\begin{equation}\label{eq:response-variation}
 V_t=\int_0^tU_N(t,u)h(u)\,\dd u.
\end{equation}
To justify the representation on a fixed horizon, integrate the causal homogeneous response in its forcing time and apply stochastic Fubini. This gives the forced variational equation, and pathwise uniqueness identifies the result with $V$. The same causal-source argument applies to bounded continuous adapted random directions, including the disorder directions used in \cref{prop:energy-bridge}. In words, $U_N(t,u)$ is the homogeneous linearized response to a tangent-projected impulse at time $u$. We define its annealed normalized trace by
\begin{equation}\label{eq:annealed-response}
 \bar r_N(t,u)=\frac1N\E\tr U_N(t,u),\qquad 0\le u\le t.
\end{equation}
This expectation is finite by the fourth-moment bound \eqref{eq:finite-lag-flow-moments} below. The finite-$N$ diagonal is $\bar r_N(t,t)=\kappa_N$, because $U_N(t,t)=P_{X_t}$ and $\tr P_x=N-1$ on $S_N$; the limiting response will have diagonal value one.

\begin{remark}[The equilibrium fluctuation--dissipation relation in this convention]\label{rem:fdt-convention}
In this convention, stationary equilibrium FDT reads $r=-c'$, with no factor of $\beta$; at lag zero, the derivative is taken from the right. To see this, fix $J$ and let $P^H_\tau=e^{\tau\mathcal L_N}$. For the Gibbs initial law, the correlation and response at lag $\tau$ are
\[
 c_{N,G}(\tau)=\frac1N\sum_{i=1}^N\int x_iP^H_\tau x_i\,g_N\,\dd\omega_N,\qquad
 r_{N,G}(\tau)=\frac1N\sum_{i=1}^N\int\nabla_Sx_i\cdot\nabla_SP^H_\tau x_i\,g_N\,\dd\omega_N.
\]
The second formula is the trace of the mean of $U_N$ for the Gibbs initial law. The actual tangent-projected restart is Fr\'echet differentiable into $L^2$, and the bounded linear expectation operator commutes with that derivative. The Markov restart property identifies the expected response with the expected derivative of $P^H_{t-u}\operatorname{id}$ at the waiting state. Gibbs stationarity then gives the displayed integral, whose integrand is $\sum_i\nabla_Sx_i\cdot\nabla_SP^H_{t-u}x_i$ because $\nabla_Sx_i=P_xe_i$. The fixed-disorder moment estimate \eqref{eq:tangent-flow-bound} ensures integrability. Since $\mathcal L_N$ is symmetric with respect to $g_N\,\dd\omega_N$, we have
\[
 c_{N,G}'(\tau)=\frac1N\sum_{i=1}^N\int x_i\,\mathcal L_NP^H_\tau x_i\,g_N\,\dd\omega_N
 =-r_{N,G}(\tau).
\]
A force written as $\beta P_xh$ would have response $\beta r_{N,G}$, and the relation would read $r=-\beta c'$ in that convention.
\end{remark}

\section{The hard-spherical CK system}\label{sec:hard-ck}

By a solution of the hard-spherical CK system, we mean a triple $(C,R,\lambda)$ of continuous functions, with $C$ defined and symmetric on $[0,\infty)^2$, $R$ defined on the causal triangle $\{(t,s):0\le s\le t\}$, and $\lambda$ defined on $[0,\infty)$, such that, for every $s\ge0$, the functions $t\mapsto C(t,s)$ and $t\mapsto R(t,s)$ are continuously differentiable on $[s,\infty)$ and satisfy, for $t\ge s$,
\begin{align}
 \partial_tR(t,s)&=-\lambda(t)R(t,s)
 +\beta^2\int_s^t\xi''(C(t,u))R(t,u)R(u,s)\,\dd u,
 \label{eq:ck-response}\\
 \partial_tC(t,s)&=-\lambda(t)C(t,s)
 +\beta^2\int_0^t\xi''(C(t,u))R(t,u)C(u,s)\,\dd u\notag\\
 &\hspace{1em}+\beta^2\int_0^s\xi'(C(t,u))R(s,u)\,\dd u,
 \label{eq:ck-correlation}\\
 \lambda(t)&=1+\beta^2\int_0^t
 \bigl[\xi'(C(t,u))+C(t,u)\xi''(C(t,u))\bigr]R(t,u)\,\dd u\notag\\
 &=1+\beta^2p^2\int_0^tC(t,u)^{p-1}R(t,u)\,\dd u,
 \label{eq:ck-multiplier}
\end{align}
with the boundary data
\begin{equation}\label{eq:ck-boundary}
 C(t,t)=R(t,t)=1,\qquad t\ge0.
\end{equation}
At $t=s$, the derivatives are one-sided. The second line of \eqref{eq:ck-multiplier} holds because $\xi'(q)+q\xi''(q)=pq^{p-1}+p(p-1)q^{p-1}=p^2q^{p-1}$.

The correlation equation \eqref{eq:ck-correlation} describes $C$ only on the causal side $t\ge s$, and the diagonal requires a word. At $s=t$, the symmetry $C(u,t)=C(t,u)$ and \eqref{eq:ck-multiplier} turn \eqref{eq:ck-correlation} into
\begin{equation}\label{eq:ck-diagonal-slope}
 \partial_tC(t,s)\big|_{s=t}=-\lambda(t)+\bigl(\lambda(t)-1\bigr)=-1.
\end{equation}
The derivative in \eqref{eq:ck-diagonal-slope} holds the earlier argument fixed and approaches the diagonal from the causal side. The function $t\mapsto C(t,t)$ is constant; its derivative is zero. For the diffusion, It\^o's formula gives $\dd(X_t\cdot X_t)/N=2X_t\cdot\dd X_t/N+2\kappa_N\,\dd t$, where the drift part of $2X_t\cdot\dd X_t/N$ is $-2\kappa_N\,\dd t$. Thus the inward drift and the quadratic variation balance exactly at finite $N$, and their coefficients tend to $-2$ and $+2$ as $N\to\infty$. The multiplier $\lambda$ enforces the limiting balance and keeps $C(t,t)=1$; this is the meaning of the term $1$ in \eqref{eq:ck-multiplier}.

We write
\begin{equation}\label{eq:energy-chi-def}
 e_\beta(t)=\beta p\int_0^tC(t,u)^{p-1}R(t,u)\,\dd u,\qquad
 \chi(t,s)=\int_s^tR(t,u)\,\dd u.
\end{equation}
By \eqref{eq:ck-multiplier}, we have $\lambda(t)=1+p\beta e_\beta(t)$, and \cref{prop:energy-bridge} shows that $e_\beta(t)$ is the limit of the mean energy $N^{-1}\E H_N(X_t)$. The quantity $\chi(t,s)$ is the integrated response, or susceptibility.

We use the following published facts throughout the paper.

\begin{proposition}[Well-posedness and positivity]\label{prop:ck-wellposed}
The hard-spherical CK system \eqref{eq:ck-response}--\eqref{eq:ck-boundary} has a unique solution $(C,R,\lambda)$. It satisfies
\begin{equation}
0\le C(t,s)\le1\quad(s,t\ge0),\qquad R(t,s)\ge0\quad(t\ge s\ge0),
\label{eq:ck-physical-bounds}
\end{equation}
and $C$ is a positive-semidefinite kernel: for all $n\ge1$, times $t_1,\dots,t_n\ge0$, and $a\in\R^n$, we have $\sum_{i,j=1}^na_ia_jC(t_i,t_j)\ge0$.
\end{proposition}

Existence and uniqueness, the bounds \eqref{eq:ck-physical-bounds}, and the positive semidefiniteness of $C$ are proved by Dembo, Guionnet, and Mazza \cite[arXiv version~1, Proposition~1.1, p.~4, and Lemma~2.1, p.~6]{dgm2007} for the hard-spherical equations in the noise-one convention recalled in \cref{sec:conventions}. Their hypotheses are met by the polynomial $\xi(q)=q^p$, the initial value $C(0,0)=1$, and any fixed temperature; their Proposition~1.1 has no restriction to high temperature. Existence and uniqueness in the noise-one convention are also recorded in \cite[arXiv version~5, Proposition~4.2, p.~23]{sellke2024}. The time change \eqref{eq:sellke-timechange} below, which is a bijection between solutions of the two systems, transfers all of these properties to \eqref{eq:ck-response}--\eqref{eq:ck-boundary}. The positive semidefiniteness of $C$ can also be read off from \cref{prop:hard-bridge}. For fixed times $t_i$ and coefficients $a_i$, we have
\[
 \sum_{i,j=1}^na_ia_jC_N(t_i,t_j)=\frac1N\Bigl\|\sum_{i=1}^na_iX_{t_i}\Bigr\|^2\ge0,
\]
and the left side converges in probability to $\sum_{i,j}a_ia_jC(t_i,t_j)$.

We record two consequences that are used constantly. Because $C$ and $R$ are continuous, each of them is bounded on every compact set of its domain. The right sides of \eqref{eq:ck-response} and \eqref{eq:ck-correlation} are continuous on the causal triangle $\{0\le s\le t\}$, and so are the partial derivatives $\partial_tR$ and $\partial_tC$.

\section{Comparison of conventions}\label{sec:conventions}

The factors of two in the CK equations depend on the noise amplitude, on the normalization of time, temperature, and covariance, and on how an external force enters the drift. We compare our convention with the two conventions of the works on which we rely.

\paragraph{The noise-one convention.}
The works of Ben Arous, Dembo, and Guionnet \cite{bag2006}, of Dembo, Guionnet, and Mazza \cite{dgm2007}, and of Sellke \cite{sellke2024} use the noise coefficient one. On $S_N$, the generator is
\[
 \tfrac12\Delta_S+\alpha\nabla_SH_N\cdot\nabla_S,
\]
with the same covariance $N\xi(x\cdot y/N)$. Here $\alpha>0$ is the drift coefficient, which Sellke denotes by $\beta$; his Gibbs measure is proportional to $e^{2\alpha H_N}$ \cite[arXiv version~5, Equation~(1.2) and the following paragraph, p.~2]{sellke2024}. In \cite{dgm2007}, the dynamics descends a Hamiltonian, and reversing its sign gives the law of $H_N$. In \cite[Equations~(1.2), (1.6), and~(1.7)]{bag2006}, specialized to $\nu(q)=q^p$, the displayed Hamiltonian $H^{\mathrm{BAG}}$ has the law of $-2H_N$, and the force in the drift is $-\nabla H^{\mathrm{BAG}}/2$, which has the law of $\nabla H_N$.
If $x_v$ is the noise-one diffusion with $\alpha=\beta/2$, then $X_t=x_{2t}$ is our diffusion; this is Step~1 of the proof of \cref{prop:hard-bridge}. Under this time change, the multiplier of the noise-one system is multiplied by two, its memory coefficient $\alpha^2$ is multiplied by four, and its noise diagonal one becomes two. This produces exactly \eqref{eq:ck-response}--\eqref{eq:ck-multiplier}. We do not rescale the covariance $N\xi$.

\paragraph{The convention of Cugliandolo and Kurchan.}
In \cite[Equations~(1)--(4)]{ck1993}, the couplings $J_{i_1\cdots i_p}$, $i_1<\dots<i_p$, have variance $p!/(2N^{p-1})$, so that the Hamiltonian $H^{\mathrm{CK}}$ has covariance $N(x\cdot y/N)^p/2$ up to terms that do not affect the limiting equations, and the energy is minimized. The Langevin equation is $\partial_ts=-\beta^{\mathrm{CK}}\nabla H^{\mathrm{CK}}(s)-z(t)s+\eta$, with white noise of variance $2$, and the response has diagonal value one. At the level of the limiting CK equations, the covariance of $-\sqrt2H^{\mathrm{CK}}$ has the normalization of the covariance of $H_N$, and the identification is
\[
 H_N\leftrightarrow-\sqrt2\,H^{\mathrm{CK}},\qquad
 \beta^{\mathrm{CK}}=\sqrt2\,\beta,
\]
with no change of time. With these substitutions, the causal forms of equations~(3)--(4) of \cite{ck1993}, for $t>s$, coincide with \eqref{eq:ck-response}--\eqref{eq:ck-correlation},\footnote{In arXiv:cond-mat/9303036v1, the factor $C^{p-1}(t,t')$ in the first memory integral of equation~(3) should read $C^{p-1}(t,t'')$.} their multiplier $z(t)=1-p\beta^{\mathrm{CK}}E(t)$ is $\lambda(t)$, and their energy per spin is $E(t)=-e_\beta(t)/\sqrt2$. Their equation~(7) for the threshold plateau is the marginality relation $\beta^2p(p-1)q^{p-2}(1-q)^2=1$.

\Cref{tab:conventions} collects the dictionary.

\begin{table}[ht]
\centering
\small
\renewcommand{\arraystretch}{1.25}
\begin{tabular}{@{}lccc@{}}
\toprule
 & this paper & noise one \cite{dgm2007,sellke2024} & CK \cite{ck1993}\\
\midrule
time & $t$ & $v=2t$ & $t$\\
covariance of $H$ & $Nq^p$ & $Nq^p$ & $Nq^p/2$\\
temperature parameter & $\beta$ & $\alpha=\beta/2$ & $\beta^{\mathrm{CK}}=\sqrt2\,\beta$\\
diffusion part of generator & $\Delta_S$ & $\Delta_S/2$ & $\Delta_S$\\
multiplier & $\lambda(t)$ & $\mu(v)=\lambda(v/2)/2$ & $z(t)=\lambda(t)$\\
constant term of multiplier & $1$ & $1/2$ & $1$\\
memory coefficients & $\beta^2\xi''$, $\beta^2\xi'$ & $\alpha^2\xi''$, $\alpha^2\xi'$ & $(\beta^{\mathrm{CK}})^2\xi''/2$, $(\beta^{\mathrm{CK}})^2\xi'/2$\\
noise diagonal & $2$ & $1$ & $2$\\
mean energy & $e_\beta(t)$ & $e_\beta(v/2)$ & $E(t)=-e_\beta(t)/\sqrt2$\\
equilibrium FDT & $r=-c'$ & $r=-2c'$ & $r=-c'$\\
transition, $p=3$ & $\bd^2=4/3$ & $\alpha_{\mathrm d}^2=1/3$ & $(\beta^{\mathrm{CK}}_{\mathrm d})^2=8/3$\\
transition, general $p$ & $\bd$ & $\alpha_{\mathrm d}=\bd/2$ & $\beta^{\mathrm{CK}}_{\mathrm d}=\sqrt2\,\bd$\\
plateau at transition & $\qd$ & $\qd$ & $\qd$\\
\bottomrule
\end{tabular}
\caption{The dictionary between the conventions. The CK column refers to the normalization of the limiting CK equations. In each column, $r$ and $c$ denote the stationary response and correlation as functions of the lag measured in that column's time, and $\bd^2=(p-1)^{p-1}/(p(p-2)^{p-2})$ and $\qd=(p-2)/(p-1)$. In the convention of \cite{ck1993}, the transition temperature for $p=3$ is $(3/8)^{1/2}$.}
\label{tab:conventions}
\end{table}

In our convention, the predicted dynamical transition $(\bd,\qd)$ is the point of the marginality curve $\beta^2p(p-1)q^{p-2}(1-q)^2=1$ at which $x=(p-2)(1-q)/q$ equals one; see the text after equation~(12) and the concluding discussion in \cite{ck1993}. The condition $x=1$ gives $\qd=(p-2)/(p-1)$. Substituting this value into the marginality relation, we obtain $\bd^2=(p-1)^{p-1}/(p(p-2)^{p-2})$. Equivalently, $\bd^2=1/(p\,g(\qd))$ with $g(q)=q^{p-2}(1-q)$. Indeed, we have $g'(q)=q^{p-3}[(p-2)-(p-1)q]$, and $g$ is maximal on $[0,1]$ at $\qd$, with $g(\qd)=(p-2)^{p-2}/(p-1)^{p-1}$.

\section{The bridge on fixed time horizons}\label{sec:published-inputs}

The CK equations describe the diffusion \eqref{eq:hard-sde} only after we send the dimension to infinity on a fixed time horizon. Three published results combine to give this limit.
Ben Arous, Dembo, and Guionnet proved the limit for the soft-spherical diffusion, including the empirical correlation and the Brownian observable for the integrated response \cite[arXiv version~1, Equation~(1.9), p.~4, and Theorem~1.2, p.~5]{bag2006}. Dembo, Guionnet, and Mazza proved that the deterministic soft-spherical equations converge to the hard-spherical ones \cite[Proposition~1.1 and its proof]{dgm2007}. The convergence of the deterministic equations is not by itself a statement about the hard-spherical diffusion. Sellke proved the convergence theorem for the hard-spherical diffusion itself \cite[arXiv version~5, Theorem~4.3, p.~23]{sellke2024}, by coupling it with soft-spherical diffusions \cite[arXiv version~5, Lemma~3.1, pp.~16--18]{sellke2024}. We use Sellke's theorem for the correlation and for the Brownian observable. The next proposition states it in our normalization.

\begin{proposition}[The bridge on fixed time horizons]\label{prop:hard-bridge}
Let $(C,R,\lambda)$ be the solution of the hard-spherical CK system of \cref{prop:ck-wellposed}, and let
\[
 C_N(t,s)=\frac{X_t\cdot X_s}N,\qquad
 \Gamma_N(t,s)=\frac{X_t\cdot B_s}{\sqrt2\,N}.
\]
For every fixed $0\le T<\infty$, we have, uniformly on $0\le s\le t\le T$,
\begin{equation}
 C_N(t,s)\longrightarrow C(t,s),\qquad
 \Gamma_N(t,s)\longrightarrow\int_0^sR(t,u)\,\dd u
 \quad\text{in probability}.
\label{eq:hard-bridge}
\end{equation}
That is, for every $\varepsilon>0$,
\begin{equation}\label{eq:hard-bridge-precise}
 \lim_{N\to\infty}\Prob\Bigl(\sup_{0\le s\le t\le T}\Bigl\{|C_N(t,s)-C(t,s)|+\Bigl|\Gamma_N(t,s)-\int_0^sR(t,u)\,\dd u\Bigr|\Bigr\}>\varepsilon\Bigr)=0.
\end{equation}
Both convergences also hold after taking expectations, uniformly on the same triangle. The limiting response satisfies $R(t,t)=1$.
\end{proposition}

\begin{proof}
The proof has five steps. In Step~1, we write our diffusion as a time change of Sellke's diffusion. In Steps~2 and~3, we check that the same time change maps Sellke's CK kernels to the solution of \eqref{eq:ck-response}--\eqref{eq:ck-boundary} and his Brownian observable to $\Gamma_N$; this proves \eqref{eq:hard-bridge} for a deterministic initial point. In Step~4, we use a rotation to extend the convergence to the uniform initial point. In Step~5, we prove the convergence of expectations by a uniform integrability bound.

\paragraph{Step 1: the time change.}
Sellke's diffusion \cite[arXiv version~5, Equation~(1.2), p.~2]{sellke2024} has noise coefficient one and drift coefficient $\alpha$ (which Sellke denotes by $\beta$):
\begin{equation}\label{eq:sellke-sde}
 \dd x_v=\alpha P_{x_v}\nabla H_N(x_v)\,\dd v+P_{x_v}\,\dd W_v-\frac{\kappa_N}2x_v\,\dd v,
\end{equation}
where $W$ is a standard Brownian motion in $\R^N$ and $H_N$ is the Hamiltonian \eqref{eq:hamiltonian}. We take $\alpha=\beta/2$. Given our $X_0$, $J$, and $B$, let
\[
 W_v=\sqrt2\,B_{v/2},\qquad x_v=X_{v/2},\qquad v\ge0.
\]
Then $W$ is a standard Brownian motion with respect to the filtration $(\mathcal F_{v/2})_{v\ge0}$, where $(\mathcal F_t)$ is the filtration of $X_0$, $J$, and $B$. For a continuous adapted integrand $G$, we have $\int_0^{v/2}G_r\,\dd r=\frac12\int_0^vG_{w/2}\,\dd w$ and $\sqrt2\int_0^{v/2}G_r\,\dd B_r=\int_0^vG_{w/2}\,\dd W_w$. Inserting these two identities into the integral form of \eqref{eq:hard-sde} at time $v/2$, we obtain
\[
 x_v=x_0+\int_0^v\Bigl[\frac\beta2P_{x_w}\nabla H_N(x_w)-\frac{\kappa_N}2x_w\Bigr]\dd w
 +\int_0^vP_{x_w}\,\dd W_w,
\]
which is \eqref{eq:sellke-sde} with $\alpha=\beta/2$. Conversely, $X_t=x_{2t}$ and $B_t=W_{2t}/\sqrt2$. The disorder and the initial point are unchanged, and they remain independent of $W$.

\paragraph{Step 2: the CK kernels.}
Let $C^{\mathrm S}$, $R^{\mathrm S}$, and $\mu^{\mathrm S}$ be the hard-spherical CK kernels and multiplier of \cite[arXiv version~5, Definition~4, pp.~22--23]{sellke2024} with parameter $\alpha$; by \cite[arXiv version~5, Proposition~4.2, p.~23]{sellke2024}, which rests on \cite[arXiv version~1, Theorem~1.2, p.~5]{bag2006} and the proof of \cite[Proposition~1.1]{dgm2007}, they are the unique continuous, locally bounded solution of the following system. In our notation for $\xi$, and with the later time first as in \cite{sellke2024}, they solve, for $v\ge w\ge0$,
\begin{align*}
 \partial_vR^{\mathrm S}(v,w)&=-\mu^{\mathrm S}(v)R^{\mathrm S}(v,w)
 +\alpha^2\int_w^v\xi''(C^{\mathrm S}(v,y))R^{\mathrm S}(v,y)R^{\mathrm S}(y,w)\,\dd y,\\
 \partial_vC^{\mathrm S}(v,w)&=-\mu^{\mathrm S}(v)C^{\mathrm S}(v,w)
 +\alpha^2\int_0^v\xi''(C^{\mathrm S}(v,y))R^{\mathrm S}(v,y)C^{\mathrm S}(y,w)\,\dd y\\
 &\hspace{1em}+\alpha^2\int_0^w\xi'(C^{\mathrm S}(v,y))R^{\mathrm S}(w,y)\,\dd y,\\
 \mu^{\mathrm S}(v)&=\frac12+\alpha^2p^2\int_0^vC^{\mathrm S}(v,y)^{p-1}R^{\mathrm S}(v,y)\,\dd y,
\end{align*}
with $C^{\mathrm S}(v,v)=R^{\mathrm S}(v,v)=1$. We claim that
\begin{equation}
 C(t,s)=C^{\mathrm S}(2t,2s),\quad
 R(t,s)=R^{\mathrm S}(2t,2s),\quad
 \lambda(t)=2\mu^{\mathrm S}(2t).
\label{eq:sellke-timechange}
\end{equation}
Let $\tilde C$, $\tilde R$, and $\tilde\lambda$ denote the right sides of \eqref{eq:sellke-timechange}. By the uniqueness in \cref{prop:ck-wellposed}, it suffices to show that $(\tilde C,\tilde R,\tilde\lambda)$ solves \eqref{eq:ck-response}--\eqref{eq:ck-boundary}. The boundary data are immediate. For the multiplier, the substitution $y=2u$ gives
\[
 \tilde\lambda(t)=1+2\alpha^2p^2\int_0^{2t}C^{\mathrm S}(2t,y)^{p-1}R^{\mathrm S}(2t,y)\,\dd y
 =1+4\alpha^2p^2\int_0^t\tilde C(t,u)^{p-1}\tilde R(t,u)\,\dd u,
\]
which is \eqref{eq:ck-multiplier} since $4\alpha^2=\beta^2$. For the response, we have $\partial_t\tilde R(t,s)=2(\partial_vR^{\mathrm S})(2t,2s)$ by the chain rule, and the substitution $y=2u$ in the memory integral yields
\[
 \partial_t\tilde R(t,s)=-2\mu^{\mathrm S}(2t)\tilde R(t,s)
 +4\alpha^2\int_s^t\xi''(\tilde C(t,u))\tilde R(t,u)\tilde R(u,s)\,\dd u.
\]
This is \eqref{eq:ck-response}. Each memory term acquires the coefficient $4\alpha^2=\beta^2$, with one factor of two from differentiating in the new time and the other from changing the integration variable. The same computation, applied to the two memory integrals of the correlation equation, yields \eqref{eq:ck-correlation}. This proves \eqref{eq:sellke-timechange}.

\paragraph{Step 3: the observables and a deterministic initial point.}
Fix $x_0\in S_N$, and consider for the moment the diffusion started from the deterministic point $X_0=x_0$. Sellke's theorem \cite[arXiv version~5, Theorem~4.3, p.~23]{sellke2024} concerns the diffusion \eqref{eq:sellke-sde} started from a deterministic point of $S_N$, with the Hamiltonian \eqref{eq:hamiltonian}. Applied on the horizon $[0,2T]$, it states that, almost surely,
\[
 \frac{x_v\cdot x_w}N\longrightarrow C^{\mathrm S}(v,w),\qquad
 \frac{x_v\cdot W_w}N\longrightarrow\int_0^wR^{\mathrm S}(v,y)\,\dd y,
\]
uniformly on $0\le w\le v\le2T$, for any coupling of the diffusions in different dimensions. In particular, these convergences hold in probability, which is all we use. The hypotheses of the theorem are satisfied, since the Hamiltonian is the pure $p$-spin Hamiltonian with ordered independent standard Gaussian coefficients, the initial point is deterministic, and the horizon $2T$ and the parameter $\alpha=\beta/2$ are fixed as $N\to\infty$. By Step~1, we have $C_N(t,s)=x_{2t}\cdot x_{2s}/N$ and
\[
 \Gamma_N(t,s)=\frac{X_t\cdot B_s}{\sqrt2\,N}=\frac12\cdot\frac{x_{2t}\cdot W_{2s}}N.
\]
By \eqref{eq:sellke-timechange} and the substitution $y=2u$, the limit of the latter is
\[
 \frac12\int_0^{2s}R^{\mathrm S}(2t,y)\,\dd y=\int_0^sR^{\mathrm S}(2t,2u)\,\dd u=\int_0^sR(t,u)\,\dd u.
\]
The triangle $0\le w\le v\le2T$ corresponds to $0\le s\le t\le T$. This proves \eqref{eq:hard-bridge-precise} for the diffusion started from $x_0$.

\paragraph{Step 4: the uniform initial point.}
For $x_0\in S_N$, let $\delta_N(x_0,\varepsilon)$ be the probability in \eqref{eq:hard-bridge-precise} for the diffusion started from $x_0$. We claim that $\delta_N(x_0,\varepsilon)$ does not depend on $x_0$. Let $e_1$ be the first coordinate vector, let $b=x_0/\sqrt N$, and let $O=O_{x_0}$ be the Householder reflection
\[
 O=I-\frac{2(e_1-b)(e_1-b)^{\mathsf T}}{\norm{e_1-b}^2}\quad(b\ne e_1),\qquad O=I\quad(b=e_1).
\]
It is orthogonal, depends measurably on $x_0$, and satisfies $O(\sqrt N e_1)=x_0$. Let $X$ be the diffusion started from $x_0$, and set $Y_t=O^{\mathsf T}X_t$, $B'_t=O^{\mathsf T}B_t$, and $J'=(O^{\mathsf T})^{\otimes p}J$, that is,
\[
 J'_{j_1\cdots j_p}=\sum_{i_1,\ldots,i_p=1}^NJ_{i_1\cdots i_p}O_{i_1j_1}\cdots O_{i_pj_p}.
\]
The matrix $O^{\otimes p}$ is orthogonal on $\R^{N^p}$, so $J'$ is again a tensor of independent standard Gaussian entries. Also, $B'$ is a standard Brownian motion, independent of $J'$. Let $H'_N$ be the Hamiltonian with coefficients $J'$; then $H'_N(y)=H_N(Oy)$. Using $P_{Oy}=OP_yO^{\mathsf T}$ and $\nabla H'_N(y)=O^{\mathsf T}\nabla H_N(Oy)$, and applying $O^{\mathsf T}$ to \eqref{eq:hard-sde}, we obtain
\[
 \dd Y_t=\beta P_{Y_t}\nabla H'_N(Y_t)\,\dd t+\sqrt2P_{Y_t}\,\dd B'_t-\kappa_NY_t\,\dd t,\qquad Y_0=\sqrt Ne_1.
\]
We conclude that $(Y,J',B')$ has the law of the diffusion started from $\sqrt Ne_1$, together with its disorder and its noise. The matrix $O$ is orthogonal, so that $Y_t\cdot Y_s=X_t\cdot X_s$ and $Y_t\cdot B'_s=X_t\cdot B_s$; in other words, $C_N$ and $\Gamma_N$ take the same values for both triples. This implies that $\delta_N(x_0,\varepsilon)=\delta_N(\sqrt Ne_1,\varepsilon)$.

Now let $X_0$ be uniform and independent of $(J,B)$. The strong solution is a measurable function of $(X_0,J,B)$. By independence and Fubini's theorem, the probability in \eqref{eq:hard-bridge-precise} is
\[
 \int_{S_N}\delta_N(x_0,\varepsilon)\,\omega_N(\dd x_0)=\delta_N(\sqrt Ne_1,\varepsilon),
\]
which tends to zero by Step~3. This proves \eqref{eq:hard-bridge-precise}, that is, \eqref{eq:hard-bridge}.

\paragraph{Step 5: expectations.}
Since $|X_t\cdot B_s|\le\sqrt N\norm{B_s}$, we have
\[
 \sup_{0\le s\le t\le T}|\Gamma_N(t,s)|\le\frac{\sup_{s\le T}\norm{B_s}}{\sqrt{2N}}.
\]
By Doob's maximal inequality, $\E\sup_{s\le T}\norm{B_s}^2\le4\E\norm{B_T}^2=4NT$, and therefore
\[
 \E\sup_{0\le s\le t\le T}|\Gamma_N(t,s)|^2\le2T
\]
for every $N$. The function $(t,s)\mapsto\int_0^sR(t,u)\,\dd u$ is continuous on the triangle, and we let $\rho_T$ be its maximum there. Then the random variables $\sup_{0\le s\le t\le T}|\Gamma_N(t,s)-\int_0^sR(t,u)\,\dd u|$ have second moments at most $2(2T+\rho_T^2)$, and in particular they are uniformly integrable. They tend to zero in probability, and uniform integrability implies that their expectations tend to zero. The same holds for $\sup_{0\le s\le t\le T}|C_N(t,s)-C(t,s)|$, which is bounded by $2$. We have shown that
\begin{equation}\label{eq:bridge-l1}
 \lim_{N\to\infty}\E\sup_{0\le s\le t\le T}\Bigl(|C_N(t,s)-C(t,s)|^2
 +\Bigl|\Gamma_N(t,s)-\int_0^sR(t,u)\,\dd u\Bigr|\Bigr)=0,
\end{equation}
where for the square we used $|C_N-C|^2\le2|C_N-C|$. Because $|\E Z|\le\E|Z|$, this implies the convergence of the expectations, uniformly on the triangle. Finally, $R(t,t)=1$ is part of the boundary data \eqref{eq:ck-boundary}.
\end{proof}

\section{The response as a Brownian primitive}\label{sec:brownian-primitive}

\Cref{prop:hard-bridge} concerns a Brownian observable. The next lemma identifies its mean with the integrated response to an additive force, which is how the Brownian observable is used as a response in \cite[arXiv version~1, Equation~(1.9), p.~4, and Theorem~1.2, p.~5]{bag2006} and in \cite[arXiv version~5, Theorem~4.3, p.~23]{sellke2024}. The identification is the classical connection between linear response and Gaussian integration by parts. The original Wiener translation formulas of Cameron and Martin \cite[Theorems~1--2, pp.~386--387]{cameronmartin1944} assume a shift whose derivative has bounded variation and use Wiener paths with variance $t/2$; for a modern deterministic-shift formulation, see \cite[Exercise~4.29, p.~27, and Theorem~4.44, p.~31]{hairer2025spde}. We use the backward It\^o equation for the actual conditional mean, which fixes the tangent projection and the normalization of the noise directly.

\begin{lemma}[Brownian primitive of the averaged force response]\label{lem:brown-force-response}
Let $U_N(t,u)$ be the pathwise response matrix \eqref{eq:response-matrix}, and let $\bar r_N(t,u)=N^{-1}\E\tr U_N(t,u)$ as in \eqref{eq:annealed-response}. Then, for $0\le s\le t$,
\begin{equation}
 \int_0^s\bar r_N(t,u)\,\dd u
 =\E\Gamma_N(t,s)
 =\frac{\E[X_t\cdot B_s]}{\sqrt2\,N}.
\label{eq:brown-force-response}
\end{equation}
The same identity holds with the disorder and the initial point fixed, and with the expectation taken over the Brownian motion only. At equal times, $\bar r_N(t,t)=(N-1)/N$.
\end{lemma}

\begin{proof}
Fix the disorder $J$ and the initial point $x_0$, and write $\E_B$ for expectation over $B$. Put $m_\tau=P^H_\tau\operatorname{id}$, the mean of the actual solution started from a sphere point, and
\[
 F_u=m_{t-u}(X_u),\qquad 0\le u\le t.
\]
The mean has the spatial and time regularity needed for the backward Kolmogorov equation. Applying It\^o's formula along the actual sphere trajectory, the drift vanishes and
\[
 \dd F_u=\sqrt2\,Dm_{t-u}(X_u)P_{X_u}\,\dd B_u.
\]
The predictable integrand is square-integrable on this fixed horizon, $F_0=m_t(x_0)$ is deterministic, and $F_t=X_t$. The actual $L^2$ derivative of the projected restart, followed by the bounded linear expectation operator, and the Markov restart identity give
\[
 \E_BU_N(t,u)=\E_B[Dm_{t-u}(X_u)P_{X_u}]
\]
at every fixed deterministic $u$. These are statements about the constructed solution and its variational equation, rather than assumed response identities.

For $0\le s\le t$, the It\^o product formula, or the covariance identity for square-integrable stochastic integrals, yields
\[
 \E_B[X_t^iB_s^i]
 =\sqrt2\int_0^s\E_B[Dm_{t-u}(X_u)P_{X_u}]_{ii}\,\dd u.
\]
The deterministic initial mean contributes zero because $\E_BB_s^i=0$. Equivalently, the coordinate form is
\begin{equation}\label{eq:cameron-martin-derivative}
 \int_0^s\E_B[U_N(t,u)]_{ii}\,\dd u
 =\frac1{\sqrt2}\E_B[X_t^iB_s^i].
\end{equation}
Summing over $i$ and dividing by $N$, we obtain the quenched form of \eqref{eq:brown-force-response}:
\begin{equation}\label{eq:quenched-brown-force-response}
 \int_0^s\frac1N\E_B\tr U_N(t,u)\,\dd u
 =\frac{\E_B[X_t\cdot B_s]}{\sqrt2\,N}.
\end{equation}
The right side is integrable over the disorder and initial point, since its absolute value is bounded by $\sqrt{s/2}$. The normalized trace on the left has a fourth moment bounded uniformly on each fixed causal window by \eqref{eq:finite-lag-flow-moments}. Fubini's theorem therefore permits averaging over the independent disorder and initial point, giving \eqref{eq:brown-force-response}. Finally, $\bar r_N(t,t)=(N-1)/N$ because $U_N(t,t)=P_{X_t}$ and $\tr P_{X_t}=N-1$.
\end{proof}

Combining \cref{prop:hard-bridge,lem:brown-force-response}, we find that, for every $0\le T<\infty$,
\begin{equation}\label{eq:trace-response-primitive}
 \lim_{N\to\infty}\sup_{0\le s\le t\le T}\Bigl|\int_0^s\bar r_N(t,u)\,\dd u-\int_0^sR(t,u)\,\dd u\Bigr|=0.
\end{equation}
In particular, the susceptibility $\chi(t,s)$ of \eqref{eq:energy-chi-def} is the limit of the integrated response $\int_s^t\bar r_N(t,u)\,\dd u=\E[X_t\cdot(B_t-B_s)]/(\sqrt2N)$, uniformly on $0\le s\le t\le T$.

\begin{remark}\label{rem:no-pointwise-response}
\Cref{prop:hard-bridge,lem:brown-force-response} give the convergence \eqref{eq:trace-response-primitive} of the primitives of the annealed response. Primitive convergence alone does not imply pointwise convergence of $\bar r_N$. The published convergence theorem supplies this statement about primitives, which suffices for the integrated uses in this paper. Where a response is integrated against a continuous weight, as in the proof of \cref{prop:energy-bridge}, we may pass to the limit through the primitives, using the bound on the total variation of $\bar r_N(t,u)\,\dd u$ supplied by \eqref{eq:finite-lag-flow-moments}.

The annealed means satisfy the stronger conclusion
\[
 \sup_{0\le s\le t\le T}|\bar r_N(t,s)-R(t,s)|\longrightarrow0
 \qquad(N\to\infty),\qquad 0\le T<\infty.
\]
Indeed, the projected variational equation, its time-increment estimates, and the uniform exponential moments of the landscape give a common modulus of continuity for the annealed traces on each fixed causal time window. Together with \eqref{eq:finite-lag-flow-moments}, this gives precompactness by the Arzel\`a--Ascoli theorem. Every uniformly convergent subsequence has the same primitives as $R$ by \eqref{eq:trace-response-primitive}; continuity and uniqueness of a function from its primitives identify its limit with $R$. Precompactness then gives convergence along the whole sequence. This conclusion concerns the deterministic annealed normalized traces; we assert no convergence or concentration of the random response matrices $U_N$.
\end{remark}

\section{Landscape bounds}\label{sec:landscape}

For later estimates, we define the normalized multilinear norm of the coefficient tensor,
\begin{equation}
Q_N=\frac1{\sqrt N}
\sup_{\norm{v_1}=\cdots=\norm{v_p}=1}
\left|\sum_{i_1,\ldots,i_p=1}^NJ_{i_1\cdots i_p}
               v_1^{i_1}\cdots v_p^{i_p}\right|,
\label{eq:tensor-norm}
\end{equation}
where the supremum is over unit vectors $v_1,\dots,v_p\in\R^N$. Closely related uniform tensor estimates are in \cite[arXiv version~1, Appendix~B, Lemma~B.1 and Equations~(B.5)--(B.7), pp.~27--29]{bag2006}. The next lemma shows that $Q_N$ is of order one, with uniform moments and an exponentially small probability of being large, and that it controls the Hamiltonian and its first two derivatives on $S_N$.

\begin{lemma}[Uniform landscape moments]\label{lem:landscape-moments}
For every fixed $0\le k<\infty$, we have $\sup_N\E Q_N^k<\infty$, and for every fixed $b<\infty$, we have $\sup_N\E e^{bQ_N}<\infty$. There exist constants $A_p,c_p>0$, depending only on $p$, such that $\Prob(Q_N>A_p)\le e^{-c_pN}$ for all $N\ge2$. Moreover, on $S_N$,
\begin{equation}
\begin{gathered}
 \frac{|H_N|}{N}\le Q_N,\qquad
 \frac{\norm{\nabla H_N}}{\sqrt N}\le pQ_N,\qquad
 \norm{\nabla^2H_N}_{\op}\le p(p-1)Q_N,\\
 \frac{|\Delta_SH_N|}{N}\le p^2Q_N,\qquad
 \norm{\nabla^2_SH_N}_{\op}\le p^2Q_N.
\end{gathered}
\label{eq:landscape-derivative-bounds}
\end{equation}
\end{lemma}

Because $\nabla_SH_N=P_x\nabla H_N$ and $\norm{P_x}_{\op}\le1$, the gradient bound implies $\norm{\nabla_SH_N}\le pQ_N\sqrt N$ on $S_N$.

\begin{proof}
The proof has three steps. We first use a net argument to reduce $Q_N$ to finitely many Gaussian variables, then deduce a tail bound, and finally prove the derivative bounds. We write $J(v_1,\dots,v_p)=\sum_{i_1,\ldots,i_p}J_{i_1\cdots i_p}v_1^{i_1}\cdots v_p^{i_p}$ for the multilinear form of $J$, and
\[
 M_N=\sup_{\norm{v_1}=\cdots=\norm{v_p}=1}|J(v_1,\dots,v_p)|=\sqrt N\,Q_N.
\]

\paragraph{Step 1: a net and the telescoping bound.}
Let $\varepsilon=1/(2p)$, and let $\mathcal N$ be a maximal $\varepsilon$-separated subset of the unit sphere of $\R^N$. By maximality, every unit vector is within distance $\varepsilon$ of $\mathcal N$. The balls of radius $\varepsilon/2$ around the points of $\mathcal N$ are disjoint and contained in the ball of radius $1+\varepsilon/2$. Comparing volumes gives
\[
 |\mathcal N|\le\Bigl(\frac{1+\varepsilon/2}{\varepsilon/2}\Bigr)^N=\Bigl(1+\frac2\varepsilon\Bigr)^N=(1+4p)^N.
\]
Let $M_N^{\mathrm{net}}$ be the maximum of $|J(w_1,\dots,w_p)|$ over $w_1,\dots,w_p\in\mathcal N$. Given unit vectors $v_1,\dots,v_p$, choose $w_k\in\mathcal N$ such that $\norm{v_k-w_k}\le\varepsilon$. By multilinearity, we have the telescoping identity
\[
 J(v_1,\dots,v_p)-J(w_1,\dots,w_p)
 =\sum_{k=1}^pJ(w_1,\dots,w_{k-1},v_k-w_k,v_{k+1},\dots,v_p).
\]
In the $k$th term, all arguments except $v_k-w_k$ are unit vectors, and by homogeneity its absolute value is at most $\norm{v_k-w_k}M_N\le\varepsilon M_N$. Summing over $k$, we have $|J(v_1,\dots,v_p)|\le M_N^{\mathrm{net}}+p\varepsilon M_N$. Taking the supremum over $v_1,\dots,v_p$, and using that $M_N$ is finite, we conclude that
\[
 M_N\le(1-p\varepsilon)^{-1}M_N^{\mathrm{net}}=2M_N^{\mathrm{net}}.
\]

\paragraph{Step 2: the tail bound and the moments.}
For fixed unit vectors $w_1,\dots,w_p$, the variable $J(w_1,\dots,w_p)$ is a centered Gaussian with variance $\sum_{i_1,\ldots,i_p}(w_1^{i_1}\cdots w_p^{i_p})^2=\prod_k\norm{w_k}^2=1$, because the entries of $J$ are independent. By the Gaussian tail bound, we have $\Prob(|J(w_1,\dots,w_p)|>y)\le2e^{-y^2/2}$ for $y\ge0$. The product net $\mathcal N^p$ has at most $(1+4p)^{pN}$ elements. By Step~1 and a union bound, for $x\ge0$,
\[
 \Prob(Q_N>x)\le\Prob\Bigl(M_N^{\mathrm{net}}>\frac{x\sqrt N}2\Bigr)
 \le2\exp\Bigl(N\kappa_p-\frac{Nx^2}8\Bigr),\qquad \kappa_p=p\log(1+4p).
\]
Let $A_p=4\sqrt{\kappa_p}$. For $x\ge A_p$, the inequality $x^2/16\ge\kappa_p$ gives
\begin{equation}\label{eq:tensor-tail}
 \Prob(Q_N>x)\le2e^{-Nx^2/16}\le2e^{-x^2/8},\qquad x\ge A_p,\ N\ge2.
\end{equation}
At $x=A_p$, the first bound in \eqref{eq:tensor-tail} is $2e^{-N\kappa_p}\le e^{-c_pN}$ with $c_p=\kappa_p-\frac12\log2>0$, because $N\ge2$. For $k>0$ and $b\ge0$, we write the moments as integrals of the tail and use \eqref{eq:tensor-tail} beyond $A_p$:
\[
 \E Q_N^k\le A_p^k+\int_{A_p}^\infty kx^{k-1}\,2e^{-x^2/8}\,\dd x,\qquad
 \E e^{bQ_N}\le e^{bA_p}+\int_{A_p}^\infty be^{bx}\,2e^{-x^2/8}\,\dd x.
\]
Both right sides are finite and do not depend on $N$. The case $k=0$ is immediate, and for $b<0$, we have $e^{bQ_N}\le1$. This proves the assertions about the moments and the tail.

\paragraph{Step 3: the derivative bounds.}
We write a point of $S_N$ as $x=\sqrt Nv$ with $\norm v=1$. Then $H_N(x)=N^{-(p-1)/2}J(x,\dots,x)=\sqrt N\,J(v,\dots,v)$ and $|H_N(x)|\le\sqrt NM_N=NQ_N$. Each derivative of $H_N$ replaces one copy of $x$ by a test direction. For a unit vector $w$, we have
\[
 \nabla H_N(x)\cdot w=N^{-(p-1)/2}\sum_{k=1}^pJ(x,\dots,x,\underbrace{w}_{k},x,\dots,x),
\]
a sum of $p$ terms, each of absolute value at most $N^{-(p-1)/2}N^{(p-1)/2}M_N=M_N$. Summing these $p$ bounds gives $\norm{\nabla H_N(x)}\le pM_N=p\sqrt NQ_N$. Similarly, $\nabla^2H_N(x)(w,w')$ is $N^{-(p-1)/2}$ times a sum of $p(p-1)$ terms in which two of the slots carry the unit vectors $w,w'$ and the others carry $x$. Each term is at most $N^{(p-2)/2}M_N$ in absolute value, and the bound $|\nabla^2H_N(x)(w,w')|\le p(p-1)N^{-1/2}M_N=p(p-1)Q_N$ follows. This proves the first three inequalities in \eqref{eq:landscape-derivative-bounds}.

For the intrinsic Hessian, let $w\in T_xS_N$ be a unit vector, and let $\gamma(r)=x\cos(r/\sqrt N)+\sqrt Nw\sin(r/\sqrt N)$ be the great circle through $x$ in the direction $w$. It is a unit-speed geodesic of $S_N$ with $\gamma''(0)=-x/N$, and we have
\[
 \nabla^2_SH_N(x)(w,w)=\frac{\dd^2}{\dd r^2}H_N(\gamma(r))\Big|_{r=0}
 =\nabla^2H_N(x)(w,w)-\frac{x\cdot\nabla H_N(x)}N.
\]
The polynomial $H_N$ is homogeneous of degree $p$, and Euler's identity reads $x\cdot\nabla H_N(x)=pH_N(x)$. By polarization, we obtain, as quadratic forms on $T_xS_N$,
\begin{equation}\label{eq:intrinsic-hessian}
 \nabla^2_SH_N(x)=P_x\nabla^2H_N(x)P_x-\frac{pH_N(x)}NP_x.
\end{equation}
By the bounds already proved, $\norm{\nabla^2_SH_N}_{\op}\le p(p-1)Q_N+pQ_N=p^2Q_N$. Finally, $\Delta_SH_N(x)$ is the trace of $\nabla^2_SH_N(x)$ on the $(N-1)$-dimensional space $T_xS_N$, which gives $|\Delta_SH_N(x)|\le(N-1)p^2Q_N\le Np^2Q_N$. This proves the last two inequalities.
\end{proof}

\chapter{Structural identities and uniform estimates}\label{sec:ck-structure}

In this section, we prove the structural identities of the CK solution and the estimates on it that hold uniformly in time. \Cref{prop:energy-bridge} shows that the CK integral $e_\beta(t)$ is the limit of the mean energy of the diffusion, and that $\lambda(t)=1+p\beta e_\beta(t)$. \Cref{prop:green-covariance}, the Green covariance identity, expresses the correlation as an explicit quadratic functional of the response. We realize it by unit vectors $Z_t$ in a Hilbert space. From these, we deduce an energy--response tradeoff and the bound $1\le\lambda\le\Lambda$ (\cref{prop:uniform-bounds}), Lipschitz and H\"older bounds in time (\cref{prop:global-regularity}), an equation for the response in its earlier time (\cref{prop:backward-response}), and the domination $R\le C$ together with the compactness of the CK solution along late time windows (\cref{prop:response-compactness}). All constants are uniform over the half-line $[0,\infty)$.

We use the results of \cref{sec:model}. \Cref{prop:energy-bridge} uses \cref{prop:hard-bridge,lem:brown-force-response,lem:landscape-moments}, and \cref{prop:global-regularity} uses \cref{prop:hard-bridge,lem:landscape-moments,prop:energy-bridge}. All other proofs in this section are deterministic, and they use only the CK equations \eqref{eq:ck-response}--\eqref{eq:ck-boundary} and the published facts of \cref{prop:ck-wellposed}. The results of this section are used throughout the rest of the paper. The bounds and the compactness enter the construction of stationary limits in \cref{sec:entropy-clusters}. The Green covariance identity is the starting point of the tree realization in \cref{sec:geometry-core} and of the response estimates in \cref{resp:chapter}. The energy identity identifies the energy of a stationary cluster in \cref{sec:entropy-clusters,sec:stationary-selection}. We use the backward response equation in \cref{resp:chapter,resp:fisher-chapter}, and the domination $R\le C$ transfers the forgetting of correlations to the response in \cref{sec:stationary-selection}.

The Green covariance identity is the variation-of-constants formula for the linear effective process of dynamical mean-field theory, in which a single coordinate is driven by its own response kernel, by white noise, and by a Gaussian noise with covariance $\beta^2\xi'(C)$. The single-coordinate retarded equation and its colored Gaussian noise are described in \cite[arXiv version~2, Section~9.5, Equation~(9.56), p.~55; see also Equations~(9.13)--(9.14), p.~45]{cugliandolo2004}. The dynamical equations behind this picture are those of Cugliandolo and Kurchan \cite[arXiv version~1, Equations~(2)--(4), pp.~2--3]{ck1993}, established for the soft-spherical model in \cite[arXiv version~1, Theorem~1.2, p.~5]{bag2006}; for a conditional Gaussian-field description in a related setting, see \cite[arXiv version~2, Remark~1.4, Equations~(1.23)--(1.24), p.~5]{ds2020}. These works provide the background. We prove the identity in the exact form \eqref{eq:green-covariance} directly from the CK equations, so that the later uses of the complete memory are transparent. The other estimates of this section are elementary, and we include their proofs for completeness.

Throughout this section, $(C,R,\lambda)$ is the solution of the hard-spherical CK system given by \cref{prop:ck-wellposed}, and $e_\beta$ is defined by \eqref{eq:energy-chi-def}.

\section{The energy identity}\label{sec:energy-identity}

\begin{proposition}[Energy identity on fixed horizons]\label{prop:energy-bridge}
For every $T<\infty$, the annealed normalized energy converges uniformly on $[0,T]$, and its limit is
\begin{equation}
e_\beta(t)=\lim_{N\to\infty}\frac1N\E H_N(X_t)
=\beta p\int_0^t C(t,u)^{p-1}R(t,u)\,\dd u,
\qquad
\lambda(t)=1+p\beta e_\beta(t).
\label{eq:energy-identity}
\end{equation}
\end{proposition}

In the physics literature, the energy is expressed in the same way through the correlation--response integral of \cite[Equation~(5)]{ck1993}. The sentence after that equation prints the energy with the sign opposite to their equation~(9); the dictionary of \cref{sec:conventions}, with $E(t)=-e_\beta(t)/\sqrt2$, agrees with the sign in equation~(9). The proof below justifies the hard-sphere form through the primitives of the response and a rank-one correction in finite dimension. It is a Gaussian integration by parts in the disorder. Differentiating the trajectory with respect to one Gaussian coefficient produces the pathwise response matrix $U_N(t,u)$ of \eqref{eq:response-matrix}, and differentiating the covariance \eqref{eq:hamiltonian-covariance} splits the result into a trace term and a rank-one correction. The trace term carries the correlation weight $C_N(t,u)^{p-1}$ against the response; in the limit, it becomes the CK integral. Two estimates make this rigorous. The first is a fourth-moment bound for the response columns and normalized trace that is uniform in the dimension and in the waiting time $u$, and the second is the convergence \eqref{eq:trace-response-primitive} of the primitives of the annealed response. The rank-one correction is controlled by the operator norm, whose second moment is at most a constant times $N$.

\begin{proof}
The proof has four steps. In Step~1, we bound the drift derivative and the noise derivative trace along the sphere. In Step~2, a stopped fourth-moment estimate gives uniform bounds on the response columns and normalized trace. In Step~3, we integrate by parts in the disorder and obtain the exact identity \eqref{eq:hard-energy-ibp}. In Step~4, we pass to the limit $N\to\infty$.

Recall the variational definition \eqref{eq:response-matrix} and the normalized trace \eqref{eq:annealed-response}. The physical coefficients are
\[
 b(x)=\beta P_x\nabla H_N(x)-\kappa_Nx,\qquad
 \sigma_j(x)=\sqrt2P_xe_j,\quad j=1,\dots,N.
\]

\paragraph{Step 1: bounds for the variational coefficients.}
The landscape derivative bounds give $\norm{Db(x)[v]}\le A_N\norm v$ on $S_N$, with $A_N=1+\beta p(p+1)Q_N$.
For the noise, $D\sigma_j(x)[v]=-\sqrt2N^{-1}(vx_j+xv_j)$. Since $\norm x^2=N$, the inequality $\norm{a+b}^2\le2\norm a^2+2\norm b^2$ implies
\begin{equation}\label{eq:tangent-flow-lognorm}
 \sum_{j=1}^N\norm{D\sigma_j(x)[v]}^2\le B_N\norm v^2,
 \qquad B_N=\frac8N.
\end{equation}
These estimates apply to all ambient vectors, so no samplewise diffeomorphism or invertibility assertion is needed.

\paragraph{Step 2: column, trace and operator moments.}
Fix $u\ge0$ and a deterministic $v\in\R^N$. The true variational solution starts from $P_{X_u}v$, which is $\mathcal F_u$-measurable and has norm at most $\norm v$. Apply It\^o's formula to the fourth power of its norm, stopped at the first exit from a ball. The bounds in Step~1, Gronwall's inequality, and Fatou's lemma when removing the stopping give, conditionally on the disorder,
\begin{equation}\label{eq:tangent-flow-bound}
 \E_B\norm{U_N(t,u)v}^4
 \le\norm v^4\exp\bigl\{(4A_N+2+6B_N)(t-u)\bigr\},
 \qquad 0\le u\le t.
\end{equation}
For a $\mathcal F_u$-measurable initial direction of norm at most a deterministic $a$, the same estimate holds with $a^4$ in place of $\norm v^4$. For $N\ge1$, the rate in \eqref{eq:tangent-flow-bound} is at most $54+4\beta p(p+1)Q_N$. Consequently, for $T<\infty$, the constant
\[
 K_T=e^{54T}\sup_N\E e^{4\beta p(p+1)TQ_N}<\infty,
 \qquad \Theta_T=\sqrt{1+K_T},
\]
is independent of the dimension and waiting time. Its finiteness follows from \cref{lem:landscape-moments}. Put $\tau_N(t,u)=N^{-1}\tr U_N(t,u)$. Jensen's inequality for the average of the diagonal entries, $|e_i\cdot U_Ne_i|\le\norm{U_Ne_i}$, and the elementary bound $\norm{U_N}_{\op}^2\le\sum_{i=1}^N\norm{U_Ne_i}^2$ show that
\begin{equation}\label{eq:finite-lag-flow-moments}
\begin{aligned}
 \sup_N\sup_{u\ge0,\,0\le t-u\le T}\sup_{1\le i\le N}
 \E\norm{U_N(t,u)e_i}^4&\le K_T,\\
 \sup_N\sup_{u\ge0,\,0\le t-u\le T}
 \E|\tau_N(t,u)|^4&\le K_T,\\
 \E\norm{U_N(t,u)}_{\op}^2&\le N\Theta_T^2,
 \qquad 0\le t-u\le T.
\end{aligned}
\end{equation}
In particular, $\E|\tau_N(t,u)|^2\le\Theta_T^2$, $\E|\tau_N(t,u)|\le\Theta_T$, and $\E\norm{U_N(t,u)}_{\op}\le\sqrt N\Theta_T$. The bounds depend only on the lag window $T$, even though the waiting state depends on the disorder and the past. Jointly measurable response representatives and their product-time integrability follow from the actual projected-restart derivative and these column moments.

\paragraph{Step 3: Gaussian integration by parts.}
For a multi-index $\nu=(i_1,\dots,i_p)\in\{1,\dots,N\}^p$, let $f_\nu(x)=N^{-(p-1)/2}x_{i_1}\cdots x_{i_p}$, so that $H_N=\sum_\nu J_\nu f_\nu$. On $S_N$, we have $|x_i|\le\sqrt N$, so $|f_\nu|\le\sqrt N$ and $\norm{\nabla f_\nu}\le p$. Changing the coefficient $J_\nu$ by $\varepsilon$ adds $\varepsilon\beta P_x\nabla f_\nu(x)$ to the drift of \eqref{eq:hard-sde}, which is a tangent force $\varepsilon\beta P_{X_u}\nabla f_\nu(X_u)$ at time $u$. As in \eqref{eq:response-variation}, the variation-of-constants formula yields
\begin{equation}\label{eq:disorder-variation}
 \partial_{J_\nu}X_t=\beta\int_0^tU_N(t,u)\nabla f_\nu(X_u)\,\dd u,
\end{equation}
where the derivative is taken in $L^2$ for the actual terminal solution. The first-order expansion of $f_\nu$ on the compact sphere gives
\[
\begin{aligned}
 \partial_{J_\nu}\E_{X_0,B}f_\nu(X_t)
 &=\E_{X_0,B}[\nabla f_\nu(X_t)\cdot\partial_{J_\nu}X_t],\\
 \left|\partial_{J_\nu}\E_{X_0,B}f_\nu(X_t)\right|
 &\le\beta p^2\int_0^t\E_{X_0,B}\norm{U_N(t,u)}_{\op}\,\dd u.
\end{aligned}
\]
Step~2 makes the right side integrable under the disorder law at every fixed $N$. As $X_0$ and $B$ are independent of $J$, we apply Gaussian integration by parts to $J\mapsto\E_{X_0,B}f_\nu(X_t)$, conditionally on the other disorder coefficients. Summing over the $N^p$ multi-indices and using Fubini's theorem, we obtain
\[
 \E H_N(X_t)=\sum_\nu\E[J_\nu f_\nu(X_t)]
 =\beta\int_0^t\E\tr\Bigl[U_N(t,u)\sum_\nu\nabla f_\nu(X_u)\nabla f_\nu(X_t)^{\mathsf T}\Bigr]\dd u.
\]
To compute the sum over $\nu$, we differentiate the identity $\sum_\nu f_\nu(x)f_\nu(y)=N^{-(p-1)}(x\cdot y)^p$, the middle expression in \eqref{eq:hamiltonian-covariance}, once in $x_j$ and once in $y_i$. With $c=x\cdot y/N$, this gives
\[
 \sum_\nu\partial_if_\nu(y)\,\partial_jf_\nu(x)=pc^{p-1}\delta_{ij}+\frac{p(p-1)}Nc^{p-2}x_iy_j.
\]
Taking $x=X_t$ and $y=X_u$, so that $c=C_N(t,u)$, we have
\[
 \sum_\nu\nabla f_\nu(X_u)\nabla f_\nu(X_t)^{\mathsf T}
 =pC_N(t,u)^{p-1}I+\frac{p(p-1)}NC_N(t,u)^{p-2}X_tX_u^{\mathsf T}.
\]
Inserting this into the previous display, using $\tr(U_NX_tX_u^{\mathsf T})=X_u^{\mathsf T}U_NX_t$, and dividing by $N$, we obtain
\begin{align}
\frac1N\E H_N(X_t)
={}&\beta p\int_0^t
 \E\Bigl[C_N(t,u)^{p-1}\,\frac{\tr U_N(t,u)}N\Bigr]\dd u
 \nonumber\\
&+\frac{\beta p(p-1)}{N^2}\int_0^t
 \E\bigl[C_N(t,u)^{p-2}X_u^{\mathsf T}U_N(t,u)X_t\bigr]\dd u.
\label{eq:hard-energy-ibp}
\end{align}
The second line is a rank-one correction. Using $|C_N|\le1$, $|X_u^{\mathsf T}U_N(t,u)X_t|\le N\norm{U_N(t,u)}_{\op}$, and Step~2, its absolute value is at most $\beta p(p-1)T\Theta_T/\sqrt N$ for $0\le t\le T$. This tends to zero uniformly in $t\le T$.

\paragraph{Step 4: the limit $N\to\infty$.}
Fix $T<\infty$. In this step, only the dimension tends to infinity, with the horizon $T$ fixed. We first replace the random weight $C_N(t,u)^{p-1}$ in the first line of \eqref{eq:hard-energy-ibp} by $C(t,u)^{p-1}$. Using that $x\mapsto x^{p-1}$ is $(p-1)$-Lipschitz on $[-1,1]$, the bound $\E|\tau_N(t,u)|^2\le\Theta_T^2$, and the Cauchy--Schwarz inequality, we find that the error is at most
\[
\begin{aligned}
 &\beta p(p-1)\int_0^t\E\bigl[|C_N(t,u)-C(t,u)|\,|\tau_N(t,u)|\bigr]\dd u\\
 &\qquad\le\beta p(p-1)T\,\Theta_T\Bigl(\E\sup_{0\le u\le t\le T}|C_N(t,u)-C(t,u)|^2\Bigr)^{1/2},
\end{aligned}
\]
which tends to zero uniformly in $t\le T$ by \eqref{eq:bridge-l1}. After this replacement, and by Fubini's theorem, the first line of \eqref{eq:hard-energy-ibp} becomes $\beta p\int_0^tw(t,u)\bar r_N(t,u)\,\dd u$, where
\[
 w(t,u)=C(t,u)^{p-1},\qquad 0\le u\le t\le T.
\]

It remains to show that
\begin{equation}\label{eq:weighted-response-limit}
 \lim_{N\to\infty}\sup_{0\le t\le T}\Bigl|\int_0^tw(t,u)\bigl(\bar r_N(t,u)-R(t,u)\bigr)\,\dd u\Bigr|=0.
\end{equation}
The convergence of primitives \eqref{eq:trace-response-primitive} does not imply this by itself, because $\bar r_N(t,\cdot)$ may change sign; we also need a bound on the total variation of the signed measure $\bar r_N(t,u)\,\dd u$. Because $|\bar r_N(t,u)|\le\E|\tau_N(t,u)|\le\Theta_T$, we have
\[
 \int_0^t|\bar r_N(t,u)|\,\dd u\le T\Theta_T,\qquad 0\le t\le T,
\]
for every $N$. Let $r_T$ be the maximum of the continuous function $R$ on the triangle. Let
\[
 \eta_N=\sup_{0\le s\le t\le T}\Bigl|\int_0^s\bigl(\bar r_N(t,u)-R(t,u)\bigr)\,\dd u\Bigr|,
\]
which tends to zero by \eqref{eq:trace-response-primitive}. Fix $\delta>0$ and a partition $0=\tau_0<\tau_1<\dots<\tau_m=T$ of mesh at most $\delta$. Let $w_\delta(t,u)=w(t,\tau_{k-1})$ for $u\in[\tau_{k-1},\tau_k)$, and let $\omega_T(\delta)$ be the modulus of continuity of $w$ on the compact triangle $\{0\le u\le t\le T\}$, so that $\sup|w-w_\delta|\le\omega_T(\delta)$ and $\omega_T(\delta)\to0$ as $\delta\to0$. For fixed $t\le T$, the integral $\int_0^tw_\delta(t,u)(\bar r_N-R)(t,u)\,\dd u$ is a sum of at most $m$ terms $w(t,\tau_{k-1})\int_{\tau_{k-1}\wedge t}^{\tau_k\wedge t}(\bar r_N-R)(t,u)\,\dd u$. Each integral, including the one over a possibly truncated last interval, is a difference of two primitives evaluated at points of $[0,t]$, and $|w|\le1$. This gives
\[
 \sup_{0\le t\le T}\Bigl|\int_0^tw(t,u)\bigl(\bar r_N-R\bigr)(t,u)\,\dd u\Bigr|
 \le\omega_T(\delta)\,T\,(\Theta_T+r_T)+2m\eta_N,
\]
where the first term bounds the integral of $(w-w_\delta)(\bar r_N-R)$ by the total variation bounds above. We first let $N\to\infty$ with the partition fixed, which removes the term $2m\eta_N$, and then let $\delta\to0$. This proves \eqref{eq:weighted-response-limit}.

Combining Step~3, the replacement of the weight, and \eqref{eq:weighted-response-limit}, we conclude that
\[
 \lim_{N\to\infty}\sup_{0\le t\le T}\Bigl|\frac1N\E H_N(X_t)-\beta p\int_0^tC(t,u)^{p-1}R(t,u)\,\dd u\Bigr|=0,
\]
which is the first assertion of \eqref{eq:energy-identity}. The identity $\lambda(t)=1+p\beta e_\beta(t)$ follows from \eqref{eq:ck-multiplier}.
\end{proof}

\section{The Green covariance identity}\label{sec:green-covariance}

We now regard the CK solution as fixed and use its coefficients to define a linear problem. Let
\begin{equation}
\Sigma(t,u)=\beta^2p(p-1)C(t,u)^{p-2}R(t,u),\qquad 0\le u\le t,
\label{eq:self-energy-kernel}
\end{equation}
which is continuous and, by \eqref{eq:ck-physical-bounds}, nonnegative. Let $\mathcal D$ be the linear Volterra operator
\[
\mathcal D f(t)=f'(t)+\lambda(t)f(t)-\int_0^t\Sigma(t,u)f(u)\,\dd u.
\]
Note that $\beta^2\xi''(C)R=\Sigma$ and $\beta^2\xi'(C)=\beta^2pC^{p-1}$. The CK equations take the following form. We extend $R(\cdot,s)$ by zero on $[0,s)$. Then, for fixed $s$, \eqref{eq:ck-response} says that $\mathcal DR(\cdot,s)(t)=0$ for $t>s$, with $R(s,s)=1$. That is, the response is the causal Green kernel of $\mathcal D$, with unit diagonal. Also, \eqref{eq:ck-correlation} says that
\begin{equation}\label{eq:ck-correlation-linear}
 \mathcal DC(\cdot,s)(t)=\beta^2p\int_0^sC(t,u)^{p-1}R(s,u)\,\dd u,\qquad t>s.
\end{equation}
Although the coefficients $\lambda$ and $\Sigma$ come from the nonlinear CK system, they are fixed continuous functions in what follows.

\begin{proposition}[Green covariance identity]\label{prop:green-covariance}
For every $t,s\ge0$,
\begin{align}
C(t,s)={}&R(t,0)R(s,0)
 +2\int_0^{t\wedge s}R(t,u)R(s,u)\,\dd u
 \nonumber\\
&+\beta^2p\int_0^t\int_0^s
       R(t,u)R(s,v)C(u,v)^{p-1}\,\dd v\,\dd u.
\label{eq:green-covariance}
\end{align}
In particular,
\begin{equation}
C(t,0)=R(t,0),\qquad
2\int_0^tR(t,u)^2\,\dd u\le1.
\label{eq:initial-response-row-bound}
\end{equation}
\end{proposition}

The three terms of \eqref{eq:green-covariance} are the contributions of the initial condition, of the white noise, whose covariance is $2$ in our convention, and of the colored noise with covariance $\beta^2pC^{p-1}=\beta^2\xi'(C)$, each propagated by the response.

\begin{proof}
Let $Q(t,s)$ be the right side of \eqref{eq:green-covariance}. It is continuous and symmetric. The proof has four steps. We verify that $Q$ satisfies the same linear equation as $C$ off the diagonal, and the same equation as $C$ along the diagonal; a Gronwall argument then shows that $Q=C$. Finally, we deduce \eqref{eq:initial-response-row-bound}. Throughout, all integrands are continuous and bounded on bounded sets, and whenever we exchange the order of integration below, we apply Fubini's theorem on a bounded region.

\paragraph{Step 1: the off-diagonal equation.}
We write $Q=Q_1+Q_2+Q_3$ for the three terms:
\[
 Q_1(t,s)=R(t,0)R(s,0),\qquad
 Q_2(t,s)=2\int_0^{t\wedge s}R(t,u)R(s,u)\,\dd u,
\]
\[
 Q_3(t,s)=\beta^2p\int_0^t\int_0^sR(t,u)R(s,v)K(u,v)\,\dd v\,\dd u,\qquad K(u,v)=C(u,v)^{p-1}.
\]
Fix $s$ and let $t>s$. By \eqref{eq:ck-response}, we have $\partial_tR(t,0)=-\lambda(t)R(t,0)+\int_0^t\Sigma(t,z)R(z,0)\,\dd z$, and hence
\[
 \partial_tQ_1(t,s)=-\lambda(t)Q_1(t,s)+\int_0^t\Sigma(t,z)Q_1(z,s)\,\dd z.
\]
For $t>s$, we have $Q_2(t,s)=2\int_0^sR(t,u)R(s,u)\,\dd u$. Differentiating under the integral and using \eqref{eq:ck-response} gives
\[
 \partial_tQ_2(t,s)=-\lambda(t)Q_2(t,s)+2\int_0^sR(s,u)\int_u^t\Sigma(t,z)R(z,u)\,\dd z\,\dd u.
\]
The region $\{0\le u\le s,\ u\le z\le t\}$ is $\{0\le z\le t,\ 0\le u\le z\wedge s\}$, and exchanging the integrals yields
\[
 \partial_tQ_2(t,s)=-\lambda(t)Q_2(t,s)+\int_0^t\Sigma(t,z)Q_2(z,s)\,\dd z.
\]
Note that the inner integral produces $Q_2(z,s)$ with its upper limit $z\wedge s$, which accounts correctly for the memory times $z<s$. For the colored term, differentiating in the upper limit $t$ produces a boundary term with $R(t,t)=1$, and differentiating $R(t,u)$ inside the integral produces the memory term after the same exchange of integrals. We obtain
\[
 \partial_tQ_3(t,s)=-\lambda(t)Q_3(t,s)+\int_0^t\Sigma(t,z)Q_3(z,s)\,\dd z+\beta^2p\int_0^sK(t,v)R(s,v)\,\dd v.
\]
Adding the three equations, we conclude that, for $t>s$,
\begin{equation}\label{eq:green-off-diagonal}
 \partial_tQ(t,s)=-\lambda(t)Q(t,s)+\int_0^t\Sigma(t,z)Q(z,s)\,\dd z+\beta^2p\int_0^sC(t,v)^{p-1}R(s,v)\,\dd v.
\end{equation}
By \eqref{eq:ck-correlation-linear}, the correlation $C$ satisfies the same equation. Note that the white-noise integral has no boundary term when $t>s$, because its upper limit $t\wedge s=s$ does not move.

\paragraph{Step 2: the diagonal equation.}
Along the diagonal, the upper limit of the white-noise integral moves. We differentiate $Q(t,t)=Q_1(t,t)+Q_2(t,t)+Q_3(t,t)$, using \eqref{eq:ck-response}, $R(t,t)=1$, and the symmetry of $K$. Both arguments now move, and the computations of Step~1 give
\[
 \frac{\dd}{\dd t}Q_i(t,t)
 =-2\lambda(t)Q_i(t,t)
  +2\int_0^t\Sigma(t,u)Q_i(u,t)\,\dd u+b_i(t),
 \qquad i=1,2,3,
\]
where $b_1(t)=0$, $b_2(t)=2$, and $b_3(t)=2\beta^2p\int_0^tK(t,u)R(t,u)\,\dd u$. Adding these identities, we obtain
\begin{equation}\label{eq:green-diagonal}
\frac{\dd}{\dd t}Q(t,t)
=-2\lambda(t)Q(t,t)
 +2\int_0^t\Sigma(t,u)Q(u,t)\,\dd u
 +2\beta^2p\int_0^t C(t,u)^{p-1}R(t,u)\,\dd u+2.
\end{equation}
The last term $2$ is the boundary term $2R(t,t)^2$ from the moving upper limit of the white-noise integral. The constant function $C(t,t)=1$ satisfies the same equation. Indeed, if we substitute $Q=C$ and $C(t,t)=1$, the right side of \eqref{eq:green-diagonal} becomes
\[
 -2\lambda(t)+2\int_0^t\Sigma(t,u)C(t,u)\,\dd u+2\beta^2p\int_0^tC(t,u)^{p-1}R(t,u)\,\dd u+2,
\]
and since $\Sigma(t,u)C(t,u)=\beta^2p(p-1)C(t,u)^{p-1}R(t,u)$, this vanishes by the multiplier equation \eqref{eq:ck-multiplier}. Both $Q$ and $C$ equal one at $(0,0)$.

\paragraph{Step 3: uniqueness.}
Let $D=Q-C$. It is continuous and symmetric, $D(0,0)=0$, and by Steps~1 and~2 it satisfies the homogeneous equations
\[
\begin{aligned}
 \partial_tD(t,s)&=-\lambda(t)D(t,s)+\int_0^t\Sigma(t,z)D(z,s)\,\dd z,\qquad t>s,\\
 \frac{\dd}{\dd t}D(t,t)&=-2\lambda(t)D(t,t)+2\int_0^t\Sigma(t,z)D(z,t)\,\dd z.
\end{aligned}
\]
Fix $T<\infty$, and let $\bar\lambda_T=\sup_{[0,T]}|\lambda|$, $\bar\Sigma_T=\sup_{0\le u\le t\le T}|\Sigma(t,u)|$, and
\[
 d(t)=\sup_{0\le u,v\le t}|D(u,v)|,\qquad 0\le t\le T,
\]
which is nondecreasing and bounded. Integrating the diagonal equation from $0$ to $s$ and bounding $|D(z,v)|\le d(v)$ for $z\le v$, we find
\[
 |D(s,s)|\le(2\bar\lambda_T+2\bar\Sigma_TT)\int_0^sd(v)\,\dd v.
\]
For $t>s$, we integrate the off-diagonal equation from $s$ to $t$ and find
\[
 |D(t,s)|\le|D(s,s)|+(\bar\lambda_T+\bar\Sigma_TT)\int_s^td(v)\,\dd v,
\]
where we used that for $v\in[s,t]$ and $z\le v$, both arguments of $D(v,s)$ and $D(z,s)$ lie in $[0,v]$. These two bounds and the symmetry of $D$ imply that $|D(t',s')|\le(2\bar\lambda_T+2\bar\Sigma_TT)\int_0^td(v)\,\dd v$ whenever $t',s'\le t$. Taking the supremum over $t'$ and $s'$, we have
\[
 d(t)\le(2\bar\lambda_T+2\bar\Sigma_TT)\int_0^td(v)\,\dd v,\qquad 0\le t\le T.
\]
By Gronwall's inequality, $d\equiv0$ on $[0,T]$. Since the horizon $T$ was arbitrary, we conclude that $Q=C$, which is \eqref{eq:green-covariance}.

\paragraph{Step 4: consequences.}
At $s=0$, the white-noise and colored terms of \eqref{eq:green-covariance} vanish and $R(0,0)=1$, so that $C(t,0)=R(t,0)$. At $s=t$, the left side of \eqref{eq:green-covariance} is one, and the three terms on the right are nonnegative, because $R\ge0$ and $C\ge0$ by \eqref{eq:ck-physical-bounds}. In particular, $2\int_0^tR(t,u)^2\,\dd u\le1$.
\end{proof}

\paragraph{Feature vectors.}
We now realize \eqref{eq:green-covariance} as a Gram identity. By \cref{prop:ck-wellposed}, $C$ is a positive-semidefinite kernel, and the Schur product theorem \cite[Section~4, Satz~VII, p.~14]{schur1911} implies that its entrywise power $K=C^{p-1}$ is positive semidefinite as well. Let $\mathcal K$ be the Hilbert space generated by vectors $V_u$, $u\ge0$, with
\[
 \ip{V_u}{V_v}=C(u,v)^{p-1}.
\]
Explicitly, we take finite formal linear combinations of the symbols $V_u$, define the inner product by bilinearity, take the quotient by the vectors of norm zero, and complete. This is the classical positive-kernel Hilbert-space construction \cite[Part~I, Section~2(4), p.~344]{aronszajn1950}. We call $\mathcal K$ the colored feature space. The vectors $V_u$ have norm one, and
\begin{equation}
\norm{V_u-V_v}^2=2\{1-C(u,v)^{p-1}\},
\label{eq:colored-feature-modulus}
\end{equation}
so $u\mapsto V_u$ is norm-continuous, because $C$ is continuous. We may therefore define the Bochner integral
\[
 A_t=\int_0^tR(t,u)V_u\,\dd u\in\mathcal K
\]
of a continuous function on a compact interval. Since bounded linear functionals commute with the integral, we have
\[
 \ip{A_t}{A_s}=\int_0^t\int_0^sR(t,u)R(s,v)C(u,v)^{p-1}\,\dd v\,\dd u,\qquad
 \ip{V_t}{A_t}=\int_0^tC(t,u)^{p-1}R(t,u)\,\dd u.
\]
In the direct sum $\R\oplus L^2(\R_+)\oplus\mathcal K$, we define
\begin{equation}\label{eq:full-feature}
Z_t=R(t,0)\ \oplus\ \sqrt2\,R(t,\cdot)\mathbf 1_{[0,t]}\ \oplus\ \beta\sqrt p\,A_t.
\end{equation}
Adding the inner products of the three orthogonal components, we see that \eqref{eq:green-covariance} says exactly
\[
 \ip{Z_t}{Z_s}=C(t,s),\qquad \norm{Z_t}=1.
\]
The initial, white-noise, and colored contributions to the correlation are orthogonal components of $Z_t$, and the colored component retains the complete history. \Cref{sec:geometry-core} refines the colored component into a causal realization graded by trees.

\section{The energy--response tradeoff and the multiplier bound}\label{sec:uniform-bounds}

\begin{proposition}[Energy--response tradeoff and uniform multiplier bound]\label{prop:uniform-bounds}
For every $t\ge0$,
\begin{equation}
\frac{e_\beta(t)^2}{p}+C(t,0)^2
                +2\int_0^tR(t,u)^2\,\dd u\le1,
\label{eq:energy-response-tradeoff}
\end{equation}
and
\begin{equation}
1\le\lambda(t)\le
\Lambda,\qquad\Lambda=\frac{1+\sqrt{1+4\beta^2p^3}}2.
\label{eq:uniform-multiplier}
\end{equation}
\end{proposition}

Both estimates are deterministic consequences of the CK equations and of \cref{prop:ck-wellposed}. \Cref{prop:energy-bridge} is needed only to interpret $e_\beta(t)$ as an energy.

\begin{proof}
The proof has three steps. We first prove the tradeoff, then deduce a first bound on the multiplier, and finally improve it to \eqref{eq:uniform-multiplier} using a lower bound on the response.

\paragraph{Step 1: the tradeoff.}
Let $M_t=\int_0^tC(t,u)^{p-1}R(t,u)\,\dd u$. With the notation of the feature vectors above, we have $M_t=\ip{V_t}{A_t}$, and since $\norm{V_t}=1$, the Cauchy--Schwarz inequality yields $M_t^2\le\norm{A_t}^2$. On the diagonal, \eqref{eq:green-covariance} reads $1=R(t,0)^2+2\int_0^tR(t,u)^2\,\dd u+\beta^2p\norm{A_t}^2$. Combining these two facts, we find
\[
 \beta^2pM_t^2+R(t,0)^2+2\int_0^tR(t,u)^2\,\dd u\le1.
\]
Since $e_\beta(t)=\beta pM_t$ and $\lambda(t)-1=\beta^2p^2M_t$, we have $e_\beta(t)^2/p=\beta^2pM_t^2=(\lambda(t)-1)^2/(\beta^2p^3)$. Using $C(t,0)=R(t,0)$ from \eqref{eq:initial-response-row-bound}, we obtain \eqref{eq:energy-response-tradeoff}, and equivalently
\begin{equation}
\frac{(\lambda(t)-1)^2}{\beta^2p^3}
 +R(t,0)^2+2\int_0^tR(t,u)^2\,\dd u\le1.
\label{eq:multiplier-response-tradeoff}
\end{equation}

\paragraph{Step 2: a first bound.}
Note that $M_t\ge0$ and $\lambda(t)\ge1$, because $C\ge0$ and $R\ge0$. By \eqref{eq:multiplier-response-tradeoff}, we also have $\lambda(t)\le1+\beta p^{3/2}$. In particular, $\Lambda_0=\sup_{t\ge0}\lambda(t)$ satisfies $1\le\Lambda_0<\infty$.

\paragraph{Step 3: the improvement.}
Since $\Sigma\ge0$ and $R\ge0$, \eqref{eq:ck-response} implies that $\partial_tR(t,u)\ge-\lambda(t)R(t,u)$ for $t\ge u$. Then $t\mapsto R(t,u)\exp\int_u^t\lambda(v)\,\dd v$ is nondecreasing, and since $R(u,u)=1$, we have
\[
 R(t,u)\ge\exp\Bigl(-\int_u^t\lambda(v)\,\dd v\Bigr)\ge e^{-\Lambda_0(t-u)}.
\]
Squaring this bound and integrating gives
\[
 R(t,0)^2+2\int_0^tR(t,u)^2\,\dd u
 \ge e^{-2\Lambda_0t}+\frac{1-e^{-2\Lambda_0t}}{\Lambda_0}
 \ge\frac1{\Lambda_0},
\]
where in the last inequality we used that the middle expression is a convex combination of $1$ and $1/\Lambda_0$, and that $\Lambda_0\ge1$. Inserting this into \eqref{eq:multiplier-response-tradeoff}, we find $(\lambda(t)-1)^2\le\beta^2p^3(1-\Lambda_0^{-1})$ for every $t$. Taking times $t$ with $\lambda(t)\to\Lambda_0$ (the supremum need not be attained), we obtain
\[
 (\Lambda_0-1)^2\le\beta^2p^3\bigl(1-\Lambda_0^{-1}\bigr)=\beta^2p^3\,\frac{\Lambda_0-1}{\Lambda_0}.
\]
If $\Lambda_0>1$, we divide by $\Lambda_0-1$ and find $\Lambda_0^2-\Lambda_0-\beta^2p^3\le0$, which implies that $\Lambda_0$ is at most the positive root $\Lambda$ of $x^2-x-\beta^2p^3$. If $\Lambda_0=1$, then $\Lambda_0\le\Lambda$ because $\Lambda\ge1$. This proves \eqref{eq:uniform-multiplier}.
\end{proof}

\section{Regularity in time}\label{sec:regularity}

\begin{proposition}[Uniform regularity in time]\label{prop:global-regularity}
There exist finite constants $D_{p,\beta}$ and $L_{p,\beta}$, depending only on $p$ and $\beta$, such that $e_\beta$ and $\lambda$ are $L_{p,\beta}$-Lipschitz on $[0,\infty)$, and, for all $u,v\ge0$ and $h\ge0$,
\begin{equation}
|C(u+h,v)-C(u,v)|\le D_{p,\beta}h+\sqrt{2h},\qquad
|C(u,v+h)-C(u,v)|\le D_{p,\beta}h+\sqrt{2h}.
\label{eq:global-modulus}
\end{equation}
Moreover, with $\Lambda$ as in \eqref{eq:uniform-multiplier}, for all $t\ge s\ge0$,
\begin{equation}\label{eq:diagonal-increment}
 0\le1-C(t,s)\le\Lambda(t-s).
\end{equation}
\end{proposition}

We obtain the estimates \eqref{eq:global-modulus} from the diffusion in finite dimension and pass them to the limit with $u$, $v$, and $h$ fixed; their constants do not depend on the times involved. The bound \eqref{eq:diagonal-increment} is stronger near the diagonal. It implies that any unit vectors $v_t$ in a Hilbert space with $\ip{v_t}{v_s}=C(t,s)$, such as $Z_t$ in \eqref{eq:full-feature}, satisfy
\[
 \norm{v_t-v_s}^2=2\{1-C(t,s)\}\le2\Lambda|t-s|,
\]
so that $t\mapsto v_t$ is $1/2$-H\"older. The modulus \eqref{eq:global-modulus} alone would only yield the exponent $1/4$.

\begin{proof}
Let $\gamma_j=\sup_N\E Q_N^j$ for $j=1,2$, which are finite by \cref{lem:landscape-moments} and depend only on $p$.

\paragraph{Step 1: the energy and the multiplier.}
Let $e_N(t)=N^{-1}\E H_N(X_t)$. By It\^o's formula on $S_N$ and \eqref{eq:generator},
\[
 H_N(X_t)-H_N(X_s)=\int_s^t\bigl[\Delta_SH_N(X_r)+\beta\norm{\nabla_SH_N(X_r)}^2\bigr]\dd r+\sqrt2\int_s^t\nabla_SH_N(X_r)\cdot\dd B_r.
\]
For fixed $J$, the integrand of the stochastic integral is bounded, and the stochastic integral has mean zero; by \eqref{eq:landscape-derivative-bounds}, the integrands are integrable over the disorder as well. By \eqref{eq:landscape-derivative-bounds}, we have $|\Delta_SH_N|/N\le p^2Q_N$ and $\norm{\nabla_SH_N}^2/N\le\norm{\nabla H_N}^2/N\le p^2Q_N^2$ on $S_N$. Taking expectations gives
\[
 |e_N(t)-e_N(s)|\le p^2(\gamma_1+\beta\gamma_2)|t-s|,
\]
with a constant independent of $N$, $s$, and $t$. By \cref{prop:energy-bridge}, we have $e_N(t)\to e_\beta(t)$ for every $t$, and the Lipschitz bound passes to the limit, showing that $e_\beta$ is Lipschitz with constant $L_e=p^2(\gamma_1+\beta\gamma_2)$. By \eqref{eq:energy-identity}, $\lambda=1+p\beta e_\beta$ is Lipschitz with constant $p\beta L_e$. We take $L_{p,\beta}=\max\{1,p\beta\}L_e$.

\paragraph{Step 2: the correlation modulus.}
We decompose the increment of the path into its drift and martingale parts:
\[
 X_{u+h}-X_u=\int_u^{u+h}\bigl[\beta P_{X_r}\nabla H_N(X_r)-\kappa_NX_r\bigr]\dd r+\sqrt2\int_u^{u+h}P_{X_r}\,\dd B_r.
\]
By \eqref{eq:landscape-derivative-bounds}, the drift has norm at most $(1+\beta pQ_N)\sqrt N$. By the It\^o isometry and $\tr(P_x^2)=\tr P_x=N-1$, the martingale part has mean square norm $2(N-1)h\le2Nh$. Using the Cauchy--Schwarz inequality for the martingale part, we find
\[
 \E\frac{\norm{X_{u+h}-X_u}}{\sqrt N}\le(1+\beta p\gamma_1)h+\sqrt{2h}.
\]
By the Cauchy--Schwarz inequality and $\norm{X_v}=\sqrt N$, we also have, pathwise,
\[
 |C_N(u+h,v)-C_N(u,v)|=\frac{|(X_{u+h}-X_u)\cdot X_v|}N\le\frac{\norm{X_{u+h}-X_u}}{\sqrt N},
\]
and no relation between the times $u$ and $v$ is needed. Taking expectations, we deduce that $|\E C_N(u+h,v)-\E C_N(u,v)|\le D_{p,\beta}h+\sqrt{2h}$ with $D_{p,\beta}=1+\beta p\gamma_1$. We now let $N\to\infty$ with $u$, $v$, and $h$ fixed. By \cref{prop:hard-bridge}, applied on the horizon $\max\{u+h,v\}$ and using the symmetry of $C_N$ and $C$, we have $\E C_N\to C$ at both points, which proves the first bound in \eqref{eq:global-modulus}. The second bound follows from the first and the symmetry of $C$.

\paragraph{Step 3: the diagonal increment.}
Let $t>s$. All terms on the right side of \eqref{eq:ck-correlation} other than the first are nonnegative, by \eqref{eq:ck-physical-bounds}. Since $0\le C\le1$ and $\lambda\le\Lambda$ by \cref{prop:uniform-bounds}, we find
\[
 \partial_tC(t,s)\ge-\lambda(t)C(t,s)\ge-\Lambda.
\]
Integrating from $s$ to $t$ and using $C(s,s)=1$, we obtain $C(t,s)\ge1-\Lambda(t-s)$, which is the upper bound in \eqref{eq:diagonal-increment}. The lower bound is $C\le1$.
\end{proof}

\section{The backward response equation}\label{sec:backward-response}

The response equation \eqref{eq:ck-response} is an equation in the later time $t$. The next proposition gives the corresponding equation in the earlier time $s$.

\begin{proposition}[Backward response equation]\label{prop:backward-response}
For $t>s\ge0$, the response satisfies
\begin{equation}
\partial_sR(t,s)=\lambda(s)R(t,s)
                  -\int_s^tR(t,u)\Sigma(u,s)\,\dd u.
\label{eq:backward-response}
\end{equation}
Both partial derivatives $\partial_tR$ and $\partial_sR$ extend continuously from $\{0<s<t\}$ to each compact causal triangle $\{0\le s\le t\le T\}$. At $s=0<t$, $\partial_sR$ is the right derivative; on the positive diagonal, $\partial_tR$ is the right derivative and $\partial_sR$ the left derivative along their causal rows. At $(0,0)$ these symbols denote the continuous extensions of the interior slopes.
\end{proposition}

\begin{proof}
The idea is to write the response as a series of Volterra compositions, to reverse the order of the factors, and to differentiate the resulting formula in its lower endpoint; this needs only the continuity of $\Sigma$, not its derivative. Guionnet and Mazza give a related positive noncrossing expansion for the normalized response in a nonlinear Kraichnan equation \cite[arXiv version~1, Section~2.1 and Lemma~2.1, pp.~7--8]{guionnetmazza2005}. Here we keep $\Sigma$ fixed and use the Neumann series of the resulting linear Volterra equation.
For continuous kernels $k_1$ and $k_2$ on the causal triangle, we define the Volterra composition
\[
 (k_1*k_2)(t,s)=\int_s^tk_1(t,u)k_2(u,s)\,\dd u,\qquad 0\le s\le t,
\]
which is associative by Fubini's theorem on bounded regions. Let
\[
 G(t,s)=\exp\Bigl(-\int_s^t\lambda(v)\,\dd v\Bigr),\qquad 0\le s\le t.
\]
Note that $0<G\le1$, because $\lambda\ge1$ by \cref{prop:uniform-bounds}.

\paragraph{Step 1: the Neumann series.}
By \eqref{eq:ck-response}, $\partial_t[R(t,s)/G(t,s)]=(\Sigma*R)(t,s)/G(t,s)$ for $t\ge s$. Integrating from $s$ to $t$ and using $R(s,s)=1$ gives the variation-of-constants formula $R=G+G*\Sigma*R$. Fix $T$, and let $\bar\Sigma_T=\sup_{0\le u\le t\le T}\Sigma(t,u)$. The kernel $(G*\Sigma)^n*G$ is an integral over $2n$ ordered intermediate times in $[s,t]$ of a product of $n$ factors $\Sigma$ and $n+1$ factors $G$. As $0<G\le1$, it is at most $\bar\Sigma_T^n$ times the volume of the simplex of ordered times:
\[
 \bigl|(G*\Sigma)^n*G\bigr|(t,s)\le\bar\Sigma_T^n\,\frac{(t-s)^{2n}}{(2n)!},\qquad 0\le s\le t\le T.
\]
The Neumann series $\sum_{n\ge0}(G*\Sigma)^n*G$ therefore converges uniformly on the triangle, and its sum $\tilde R$ satisfies $\tilde R=G+G*\Sigma*\tilde R$. The difference $\tilde R-R$ satisfies the homogeneous equation $f=G*\Sigma*f$, and iterating this equation $n$ times shows that $|f|\le\bar\Sigma_T^nT^{2n}\sup|f|/(2n)!$ for every $n$, which forces $f=0$. We have shown that
\[
 R=\sum_{n\ge0}(G*\Sigma)^n*G.
\]

\paragraph{Step 2: reversing the composition.}
By associativity, $(G*\Sigma)^n*G=G*(\Sigma*G)^n$, and both are the composition of the word $G\Sigma G\cdots\Sigma G$. Separating the last two factors $\Sigma*G$ of each term with $n\ge1$, we obtain
\[
 R=G+\sum_{n\ge0}(G*\Sigma)^n*G*\Sigma*G=G+R*\Sigma*G.
\]
Explicitly, for $0\le s\le t$,
\[
 R(t,s)=G(t,s)+\int_s^tR(t,u)\,(\Sigma*G)(u,s)\,\dd u,\qquad
 (\Sigma*G)(u,s)=\int_s^u\Sigma(u,v)G(v,s)\,\dd v.
\]

\paragraph{Step 3: differentiation in the earlier time.}
We have $\partial_sG(v,s)=\lambda(s)G(v,s)$ and $G(s,s)=1$. Differentiating the inner integral in its lower endpoint yields
\[
 \partial_s(\Sigma*G)(u,s)=-\Sigma(u,s)+\lambda(s)(\Sigma*G)(u,s),
\]
which is continuous in $(u,s)$. The outer integral has lower endpoint $s$, and its boundary term vanishes because $(\Sigma*G)(s,s)=0$. By the Leibniz rule, we have
\[
\begin{aligned}
 \partial_sR(t,s)&=\lambda(s)G(t,s)+\int_s^tR(t,u)\bigl[-\Sigma(u,s)+\lambda(s)(\Sigma*G)(u,s)\bigr]\dd u\\
 &=\lambda(s)\bigl[G+R*\Sigma*G\bigr](t,s)-(R*\Sigma)(t,s).
\end{aligned}
\]
Since $G+R*\Sigma*G=R$, this is \eqref{eq:backward-response}. We differentiated only continuous integrands through their moving endpoints.

The right side of \eqref{eq:backward-response} is continuous on the closed triangle $\{0\le s\le t\le T\}$, with value $\lambda(t)$ at $s=t$. The forward derivative given by \eqref{eq:ck-response} is also continuous there, with value $-\lambda(s)$ at $t=s$. This proves the assertions about continuity.
\end{proof}

\section{Response domination and compactness}\label{sec:response-compactness}

For $T\ge0$, we consider the translated kernels
\begin{equation}\label{eq:translated-kernels}
 C_T(a,b)=C(T+a,T+b),\quad
 R_T(a,b)=R(T+a,T+b)\ \ (a\ge b),\quad
 \lambda_T(a)=\lambda(T+a),
\end{equation}
defined for offsets $a,b\ge-T$. On every compact set of offsets, they are defined once $T$ is large enough. We consider correlations on squares of offsets, and responses only on closed causal triangles $\{a\ge b\}$, where their diagonal value is one.

\begin{proposition}[Response domination and bounded-lag compactness]\label{prop:response-compactness}
Let $a_{p,\beta}=\beta^2p(p-1)$. For every $t\ge s\ge0$,
\begin{equation}
 0\le R(t,s)\le C(t,s)\le1,
 \qquad
 |\partial_tR(t,s)|+|\partial_sR(t,s)|
 \le2\{\Lambda+a_{p,\beta}(t-s)\},
 \label{eq:response-derivative-bounds}
\end{equation}
with the boundary interpretations of \cref{prop:backward-response}. The translated triples $(C_T,R_T,\lambda_T)$ are precompact for locally uniform convergence as $T\to\infty$, in the sense that every sequence $T_n\to\infty$ has a subsequence along which $C_{T_n}$, $R_{T_n}$, and $\lambda_{T_n}$ converge uniformly on every compact subset of $\R^2$, of $\{(a,b)\in\R^2:a\ge b\}$, and of $\R$, respectively.
\end{proposition}

\begin{proof}
The proof has three steps. We first prove the domination $R\le C$, then the derivative bounds, and finally the compactness.

\paragraph{Step 1: the domination.}
Fix $s\ge0$, and let $D_s(t)=C(t,s)-R(t,s)$ for $t\ge s$. We split the first memory integral of \eqref{eq:ck-correlation} at $s$. With the notation \eqref{eq:self-energy-kernel}, the correlation satisfies
\[
 \partial_tC(t,s)=-\lambda(t)C(t,s)+\int_s^t\Sigma(t,u)C(u,s)\,\dd u+F_s(t),\qquad t\ge s,
\]
where
\[
 F_s(t)=\int_0^s\Sigma(t,u)C(u,s)\,\dd u+\beta^2p\int_0^s C(t,u)^{p-1}R(s,u)\,\dd u\ge0.
\]
In words, for fixed $s$, the correlation solves the same linear Volterra equation as the response, with the same initial value one at $t=s$ and the additional nonnegative forcing $F_s$, which collects the memory from times before $s$. Subtracting \eqref{eq:ck-response}, we find
\[
 D_s'(t)=-\lambda(t)D_s(t)+\int_s^t\Sigma(t,u)D_s(u)\,\dd u+F_s(t),\qquad D_s(s)=0.
\]
Let $\tilde D_s(t)=\int_s^tR(t,v)F_s(v)\,\dd v$. Differentiating, using $R(t,t)=1$ and \eqref{eq:ck-response}, and exchanging the order of integration over $\{s\le v\le u\le t\}$, we find that $\tilde D_s$ satisfies the same equation with $\tilde D_s(s)=0$. The difference $f=D_s-\tilde D_s$ satisfies $f'=-\lambda f+\int_s^t\Sigma(t,u)f(u)\,\dd u$ and $f(s)=0$. On $[s,T]$, with $\bar f(t)=\sup_{[s,t]}|f|$, we obtain $\bar f(t)\le(\bar\lambda_T+\bar\Sigma_TT)\int_s^t\bar f(v)\,\dd v$, with $\bar\lambda_T$ and $\bar\Sigma_T$ as in the proof of \cref{prop:green-covariance}, and Gronwall's inequality shows that $f=0$. We have shown that
\[
 C(t,s)-R(t,s)=\int_s^tR(t,v)F_s(v)\,\dd v\ge0.
\]
Together with \eqref{eq:ck-physical-bounds}, this implies that $0\le R\le C\le1$.

\paragraph{Step 2: the derivative bounds.}
By Step~1 and \eqref{eq:ck-physical-bounds}, we have $0\le\Sigma(t,u)\le a_{p,\beta}$. In the forward equation \eqref{eq:ck-response}, the memory term is therefore at most $a_{p,\beta}(t-s)$ in absolute value, and the multiplier term is at most $\Lambda$ by \eqref{eq:uniform-multiplier}. This gives $|\partial_tR(t,s)|\le\Lambda+a_{p,\beta}(t-s)$. The same argument applied to the backward equation \eqref{eq:backward-response} gives $|\partial_sR(t,s)|\le\Lambda+a_{p,\beta}(t-s)$. This proves \eqref{eq:response-derivative-bounds}; at $t=s$, we use the continuous extensions of \cref{prop:backward-response}.

\paragraph{Step 3: compactness.}
For $m\in\N$, let $W_m=[-m,m]^2$ and $W_m^{\mathrm{caus}}=\{(a,b)\in W_m:a\ge b\}$. For $T\ge m$, the translated kernels are defined on these sets. On them, we have the following bounds, uniformly in $T\ge m$:
\begin{itemize}
\item $|C_T|\le1$, and $C_T$ has the modulus of continuity \eqref{eq:global-modulus} in each coordinate, by \cref{prop:global-regularity};
\item $0\le R_T\le1$, and $R_T$ is Lipschitz on the convex set $W_m^{\mathrm{caus}}$ with constant $2\{\Lambda+2a_{p,\beta}m\}$, by \eqref{eq:response-derivative-bounds}, since the lag is at most $2m$ there;
\item $1\le\lambda_T\le\Lambda$, and $\lambda_T$ is $L_{p,\beta}$-Lipschitz, by \cref{prop:uniform-bounds,prop:global-regularity}.
\end{itemize}
Let $T_n\to\infty$. By the Arzel\`a--Ascoli theorem, there exists a subsequence along which the three components converge uniformly on $W_1$, $W_1^{\mathrm{caus}}$, and $[-1,1]$. From it, we extract a further subsequence that converges uniformly on $W_2$, $W_2^{\mathrm{caus}}$, and $[-2,2]$, and so on. The diagonal subsequence converges uniformly on $W_m$, $W_m^{\mathrm{caus}}$, and $[-m,m]$ for every $m$. Every compact set of offsets in the respective domains lies in one of these sets for $m$ large, and this proves the precompactness. Offsets may be negative; each compact set lies in the domain once $T_n$ exceeds its distance from the origin. We use the response only inside the causal triangles, and we need no continuity across the diagonal.
\end{proof}

\chapter{Entropy production and stationary limits}\label{sec:entropy-clusters}

In this section, we prove that the CK solution becomes stationary on late time windows and satisfies the fluctuation--dissipation theorem (FDT) there. Our tool is the relative entropy of the diffusion with respect to its Gibbs measure. It decreases at a rate equal to the relative Fisher information. Since the initial law is uniform, the entropy produced over all times is at most $\beta^2N/2$ on average over the disorder (\cref{lem:entropy-budget}). We then show in \cref{lem:finite-dimensional-fdt} that the Fisher information at an earlier time $s$ controls the backward FDT defect $\partial_sC-R$ at that time. In \cref{thm:fisher-clock}, we pass these estimates to the limit along one subsequence of dimensions. We obtain an integrable function $a_F$, the Fisher majorant, whose square root dominates the defect for all $t>s>0$. In \cref{lem:waiting-clock-current}, we establish that the Fisher information also controls the change of a correlation at fixed lag when both observation times are shifted. Combining these estimates with compactness, we prove \cref{prop:stationary-clusters}. It states that every limit point of the translated kernels is a stationary FDT profile with a constant multiplier, and that its response satisfies a closed equation.

We use the following results of \cref{sec:model,sec:ck-structure}: the fixed-horizon bridge (\cref{prop:hard-bridge}), the Brownian primitive of the response (\cref{lem:brown-force-response}), the landscape bounds (\cref{lem:landscape-moments}), the energy identity (\cref{prop:energy-bridge}), the Green covariance identity (\cref{prop:green-covariance}), and the uniform estimates of \cref{prop:uniform-bounds,prop:global-regularity,prop:backward-response,prop:response-compactness}. The stationary clusters of \cref{prop:stationary-clusters} are the starting point of \cref{sec:scalar,sec:geometry-core,sec:stationary-selection}. We use the Fisher majorant again in \cref{app:general-selection,resp:chapter,sec:three-time,resp:fisher-chapter}, and \cref{app:general-selection} reuses the subsequence of dimensions chosen in \cref{thm:fisher-clock}.

The entropy dissipation identity, the Bochner formula, and the gradient bound for the semigroup are standard calculus for reversible diffusions. The Bochner formula for a Laplacian with a gradient drift is recorded in \cite[Proposition~3, equation~(4b), p.~187]{bakryemery1985}, and the diffusion and entropy calculus is presented in \cite[Chapter~3 and Proposition~5.2.2, p.~245]{bgl2014}. We derive the spherical response identity of \cref{lem:finite-dimensional-fdt} directly; for related response identities around nonthermal steady states under detailed balance, see \cite[Sections~2--3, pp.~27--30]{agarwal1972}; nonequilibrium diffusion response identities are given in \cite[equations~(6) and~(11), pp.~010602-2--010602-3]{bmw2009}. Cugliandolo, Dean, and Kurchan described the physical mechanism that links entropy production to violations of FDT, and they derived bounds on these violations from the entropy production \cite[arXiv version~2, equations~(12)--(13), p.~3, and~(18), p.~4]{cdk1997}; their integrated bounds are discussed in the unnumbered passage ``Integrated bounds for purely relaxational systems'' on p.~3 of that arXiv version. For the projected hard-spherical response, we establish dimension-uniform bounds and pass them to a single Fisher majorant for the CK solution. We give the calculations in full because the sign, the spherical projection, and the normalization of the response all matter later.

\paragraph{Standing notation.}
Throughout this section, $p\ge3$ and $\beta>0$ are fixed. For a fixed realization of the disorder $J$, the Hamiltonian $H_N$ is fixed. Recall from \eqref{eq:generator} that the generator of \eqref{eq:hard-sde} is
\[
 \mathcal L_N=\Delta_S+\beta\ip{\nabla_SH_N}{\nabla_S\,\cdot\,}.
\]
We write $P_\tau=P^H_\tau=e^{\tau\mathcal L_N}$ for its semigroup. All differential operators act on functions on $S_N$ and are spherical: $\nabla_S$, $\diver_S$, $\Delta_S$, and $\nabla_S^2$ are the gradient, divergence, Laplacian, and Hessian of $S_N$, as in \cref{sec:langevin}, and $\operatorname{Ric}$ is its Ricci tensor. We write $\Gamma(f,g)=\ip{\nabla_Sf}{\nabla_Sg}$ and $\Gamma(f)=\Gamma(f,f)$. The subscript $H$ marks quantities computed for the fixed Hamiltonian $H_N$. The expectation $\E_J$ averages over the Gaussian tensor, while $\E$ also averages over the initial point and the Brownian motion. Let $\rho_t$ be the density of the law of $X_t$ given $J$ with respect to $\omega_N$. The actual diffusion realizes a strictly positive density on $S_N$, with $\rho_0=1$ because the initial point is uniform and independent of $J$. We use its $C^3$ spatial regularity and the joint continuity, up to time zero, of $\rho_t$, $\nabla_S\rho_t$, its forward-equation slope, and the spatial gradient of that slope. The forward equation holds on $S_N$ at positive times and gives its right derivative at zero. All density, score, and current quantities used below are jointly measurable in the disorder and their time parameters; we establish their integrability over the disorder in the estimates below. We require no positivity or forward equation for an extension off $S_N$. Finally, $Q_N$ is the tensor norm \eqref{eq:tensor-norm}. Derivatives at time zero are always right derivatives; derivatives at positive times are ordinary derivatives.

\section{Relative entropy and an exact fluctuation--dissipation identity}\label{sec:entropy-fdt}

Recall from \cref{sec:langevin} that the Gibbs density $g_N=Z_{N,\beta}^{-1}e^{\beta H_N}$, with $Z_{N,\beta}=\int e^{\beta H_N}\,\dd\omega_N$, is reversible for $\mathcal L_N$. We write
\[
 v^{\mathrm{cur}}_t=\nabla_S\log(\rho_t/g_N)
\]
for the relative score of $\rho_t$; the superscript distinguishes it from the tree vectors $v_t$ of \cref{sec:geometry-core}. The probability current of the diffusion is $j_t=\beta\rho_t\nabla_SH_N-\nabla_S\rho_t=-\rho_tv^{\mathrm{cur}}_t$, which satisfies $\partial_t\rho_t+\diver_Sj_t=0$; we write all formulas below in terms of $v^{\mathrm{cur}}_t$. We define the relative entropy, the relative Fisher information, and its normalized annealed average by
\[
 D_H(t)=\int\rho_t\log(\rho_t/g_N)\,\dd\omega_N,
 \qquad
 I_H(t)=\int\norm{v^{\mathrm{cur}}_t}^2\rho_t\,\dd\omega_N,
 \qquad
 a_N(t)=\frac1N\E_J I_H(t).
\]

\begin{lemma}[Finite normalized entropy budget]\label{lem:entropy-budget}
For every realization of the disorder and every $t\ge0$,
\begin{equation}
 \partial_t\rho_t=\diver_S(\rho_tv^{\mathrm{cur}}_t),\qquad
 D_H'(t)=-I_H(t),\qquad
 \int_0^\infty I_H(t)\,\dd t\le D_H(0).
 \label{eq:quenched-entropy-dissipation}
\end{equation}
The first identity is an equation on $S_N$, and both derivatives at $t=0$ are taken from the right.
Moreover,
\begin{equation}
 \int_0^\infty a_N(t)\,\dd t\le\frac{\beta^2}{2},
 \qquad
 \frac{D_H(0)}N\le2\beta Q_N.
 \label{eq:normalized-entropy-budget}
\end{equation}
\end{lemma}

\begin{proof}
Note that $\rho_tv^{\mathrm{cur}}_t=\nabla_S\rho_t-\beta\rho_t\nabla_SH_N$. Then the forward equation $\partial_t\rho_t=\Delta_S\rho_t-\beta\diver_S(\rho_t\nabla_SH_N)$ of the diffusion is the first identity in \eqref{eq:quenched-entropy-dissipation}. Next, we differentiate the relative entropy. Using $\int\partial_t\rho_t\,\dd\omega_N=0$ and integrating by parts, we find
\[
 D_H'(t)=\int\log(\rho_t/g_N)\,\diver_S(\rho_tv^{\mathrm{cur}}_t)\,\dd\omega_N
 =-\int\rho_t\norm{v^{\mathrm{cur}}_t}^2\,\dd\omega_N=-I_H(t).
\]
At positive times, we may differentiate under the integral because the forward slope is continuous and the density is strictly positive on the compact sphere. Since the entropy and Fisher information are continuous up to zero, the positive-time identity gives $D_H'(0+)=-I_H(0)$. The same argument, applied to the continuous forward slope, gives the right derivative of the density at zero. Integrating from $0$ to $T$, we have $\int_0^TI_H(t)\,\dd t=D_H(0)-D_H(T)\le D_H(0)$, where we used the fact that the relative entropy of two probability densities is nonnegative. Letting $T\to\infty$ proves the third assertion of \eqref{eq:quenched-entropy-dissipation}.

We now bound the initial entropy. Since $\rho_0=1$, we have
\[
 D_H(0)=\log Z_{N,\beta}-\beta\int H_N\,\dd\omega_N.
\]
By \eqref{eq:landscape-derivative-bounds}, we have $|H_N|\le NQ_N$ on $S_N$. As $\omega_N$ is a probability measure, this implies that
\[
 -\beta NQ_N\le\log Z_{N,\beta}\le\beta NQ_N,
 \qquad
 \Bigl|\int H_N\,\dd\omega_N\Bigr|\le NQ_N.
\]
This proves the second inequality in \eqref{eq:normalized-entropy-budget}, and by \cref{lem:landscape-moments}, both $\log Z_{N,\beta}$ and $\int H_N\,\dd\omega_N$ are integrable over the disorder. For each $x\in S_N$, the variable $H_N(x)$ is a centered Gaussian of variance $N\xi(1)=N$. By Fubini's theorem, we have $\E_J\int H_N\,\dd\omega_N=0$ and
\[
 \E_JZ_{N,\beta}=\int\E_Je^{\beta H_N(x)}\,\dd\omega_N(x)=e^{\beta^2N/2}.
\]
Using Jensen's inequality, we obtain
\[
 \E_JD_H(0)=\E_J\log Z_{N,\beta}\le\log\E_JZ_{N,\beta}=\frac{\beta^2N}{2}.
\]
Finally, by Tonelli's theorem and the third assertion of \eqref{eq:quenched-entropy-dissipation},
\[
 \int_0^\infty a_N(t)\,\dd t=\frac1N\E_J\int_0^\infty I_H(t)\,\dd t\le\frac1N\E_JD_H(0)\le\frac{\beta^2}{2},
\]
which is the first inequality in \eqref{eq:normalized-entropy-budget}.
\end{proof}

We next express correlations and responses through the semigroup. Let $m_\tau=P_\tau\operatorname{id}$ be the vector of functions $(P_\tau x_1,\dots,P_\tau x_N)$. By the Markov property, $m_\tau(x)$ is the conditional mean of $X_{s+\tau}$ given $X_s=x$ and $J$. Jensen's inequality and $\norm{X_{s+\tau}}=\sqrt N$ imply that $\norm{m_\tau(x)}\le\sqrt N$. For $t\ge s$, we define the quenched correlation and the normalized impulse response by
\begin{align}
 c_H(t,s)&=\frac1N\int\rho_s(x)\ip{x}{m_{t-s}(x)}\,\dd\omega_N(x),
 \label{eq:quenched-correlation}\\
 r_H(t,s)&=\frac1N\int\rho_s\sum_{i=1}^N\Gamma(x_i,P_{t-s}x_i)\,\dd\omega_N.
 \label{eq:quenched-force-response}
\end{align}
By the Markov property again, $c_H(t,s)=N^{-1}\E[X_t\cdot X_s\mid J]$, and $\E_Jc_H(t,s)=\E C_N(t,s)$. The function $r_H$ is the response to an additive tangent force. Recall that a deterministic force $h(s)\in\R^N$ enters the linearized equation \eqref{eq:perturbed-sde} as $P_{X_s}h(s)$, without a factor of $\beta$. Because $\nabla_Sx_i=P_xe_i$, the actual mean response is expressed through the derivatives of the transition mean. For bounded continuous $h$, \eqref{eq:response-variation} and the expected projected-restart derivative identify the $i$th coordinate of the mean linearized response with
\[
 \int_0^t\int\rho_s\sum_{j=1}^Nh_j(s)\,\Gamma(x_j,P_{t-s}x_i)\,\dd\omega_N\,\dd s.
\]
In other words, $Nr_H(t,s)$ is the trace of the impulse response matrix of the quenched mean. As in \cref{rem:fdt-convention}, the actual tangent-projected restart is Fr\'echet differentiable into $L^2$, and the bounded linear expectation operator commutes with that derivative. It follows that the expected trace of $U_N(t,s)$ is the mean of $\sum_i\Gamma(x_i,P_{t-s}x_i)(X_s)$. The actual waiting-state law is $\rho_s\,\dd\omega_N$, so this is exactly the integral in \eqref{eq:quenched-force-response}. Integrability follows from the fixed-disorder fourth-moment estimate \eqref{eq:tangent-flow-bound}. We set
\[
 \bar c_N=\E_Jc_H,\qquad \bar r_N=\E_Jr_H.
\]
By the discussion above, $\bar r_N$ is the annealed normalized response \eqref{eq:annealed-response}, and $r_H(s,s)=N^{-1}\int\rho_s\operatorname{tr}P_x\,\dd\omega_N=\kappa_N$, where $\kappa_N=(N-1)/N$ as in \cref{sec:langevin}.

\begin{lemma}[Exact nonequilibrium fluctuation--dissipation identity]\label{lem:finite-dimensional-fdt}
For $0\le s<t$,
\begin{equation}
 r_H(t,s)-\partial_sc_H(t,s)
 =\frac1N\int\rho_s\ip{m_{t-s}}{v^{\mathrm{cur}}_s}\,\dd\omega_N,
 \label{eq:quenched-fdt-current}
\end{equation}
where the derivative is taken with the terminal time $t$ fixed, from the right when $s=0$. In particular,
\begin{equation}
 |\bar r_N(t,s)-\partial_s\bar c_N(t,s)|^2\le a_N(s),
 \label{eq:annealed-fdt-bound}
\end{equation}
and, for $0\le s\le t$,
\begin{equation}
 \Bigl|1-\bar c_N(t,s)-\int_s^t\bar r_N(t,u)\,\dd u\Bigr|^2
 \le(t-s)\int_s^ta_N(u)\,\dd u.
 \label{eq:finite-dimensional-integrated-fdt}
\end{equation}
\end{lemma}

\begin{proof}
We begin with an integration by parts formula for the law at time $s$. For smooth functions $f,g$ on $S_N$, we have
\begin{equation}
 \int\rho_sf\mathcal L_Ng\,\dd\omega_N
 =-\int\rho_s\Gamma(f,g)\,\dd\omega_N
 -\int\rho_sf\ip{v^{\mathrm{cur}}_s}{\nabla_Sg}\,\dd\omega_N.
 \label{eq:gibbs-current-integration-by-parts}
\end{equation}
Indeed, integrating $\Delta_Sg$ by parts against $\rho_sf$ gives $-\int\rho_s\Gamma(f,g)-\int f\ip{\nabla_S\rho_s}{\nabla_Sg}$, and the remaining term combines with the drift term of $\mathcal L_N$ because $\nabla_S\rho_s-\beta\rho_s\nabla_SH_N=\rho_sv^{\mathrm{cur}}_s$.

Next, we differentiate \eqref{eq:quenched-correlation} in $s$, using the forward equation $\partial_s\rho_s=\mathcal L_N^*\rho_s$, where $\mathcal L_N^*$ is the adjoint with respect to $\omega_N$, and the backward equation $\partial_sP_{t-s}x_i=-\mathcal L_NP_{t-s}x_i$. This gives
\[
 \partial_sc_H(t,s)=\frac1N\int\rho_s\sum_{i=1}^N\bigl[\mathcal L_N(x_iP_{t-s}x_i)-x_i\mathcal L_NP_{t-s}x_i\bigr]\,\dd\omega_N.
\]
Using the product rule $\mathcal L_N(fg)=f\mathcal L_Ng+g\mathcal L_Nf+2\Gamma(f,g)$, we can write the bracket as $(P_{t-s}x_i)\mathcal L_Nx_i+2\Gamma(x_i,P_{t-s}x_i)$. Applying \eqref{eq:gibbs-current-integration-by-parts} with $f=P_{t-s}x_i$ and $g=x_i$ to the first of these two terms, we obtain
\[
 \partial_sc_H(t,s)=\frac1N\int\rho_s\sum_{i=1}^N\bigl[\Gamma(x_i,P_{t-s}x_i)-(P_{t-s}x_i)\ip{v^{\mathrm{cur}}_s}{\nabla_Sx_i}\bigr]\,\dd\omega_N.
\]
The first term is $r_H(t,s)$. In the second, we have $\sum_i(P_{t-s}x_i)\ip{v^{\mathrm{cur}}_s}{P_xe_i}=\ip{m_{t-s}}{v^{\mathrm{cur}}_s}$, because $v^{\mathrm{cur}}_s$ is tangent and $\nabla_Sx_i=P_xe_i$. This proves \eqref{eq:quenched-fdt-current}, with its sign and coefficient. As a check of the normalization, consider $\beta=0$. Then the uniform law is invariant, $\rho_s=1$, and $v^{\mathrm{cur}}_s=0$. The coordinates satisfy $\Delta_Sx_i=-\kappa_Nx_i$, so that $c_H(t,s)=e^{-\kappa_N(t-s)}$ and $r_H(t,s)=\kappa_Ne^{-\kappa_N(t-s)}=\partial_sc_H(t,s)$.

By the Cauchy--Schwarz inequality and $\norm{m_{t-s}}\le\sqrt N$, the square of the right side of \eqref{eq:quenched-fdt-current} is at most $I_H(s)/N$. To average over the disorder, we need a bound on $\partial_sc_H$ that is integrable over $J$. In the formula for $\partial_sc_H$ before the integration by parts, the first term is $N^{-1}\int\rho_s\ip{m_{t-s}}{\mathcal L_N\operatorname{id}}\,\dd\omega_N$. Using $\mathcal L_N\operatorname{id}=\beta\nabla_SH_N-\kappa_N\operatorname{id}$ and the bound $\norm{\nabla_SH_N}\le pQ_N\sqrt N$ from \eqref{eq:landscape-derivative-bounds}, we see that its absolute value is at most $\kappa_N+\beta pQ_N$. The second term is $2N^{-1}\int\rho_s\operatorname{tr}(D_Sm_{t-s}P_x)\,\dd\omega_N$, and its absolute value is at most $2\norm{D_Sm_{t-s}}_{\op}\le2e^{\beta p^2Q_N(t-s)}$ by the gradient bound \eqref{eq:conditional-mean-gradient} below, whose proof uses only the Bochner formula. We conclude that
\[
 |\partial_sc_H(t,s)|\le1+\beta pQ_N+2e^{\beta p^2Q_N(t-s)},
 \qquad
 |r_H(t,s)|\le e^{\beta p^2Q_N(t-s)}.
\]
These bounds are integrable over the disorder by \cref{lem:landscape-moments}, uniformly for $t-s$ in compact sets, and dominated convergence justifies $\partial_s\bar c_N=\E_J\partial_sc_H$ for $0<s<t$. The quenched and annealed correlations, responses, and currents are continuous on each closed causal row, so the identities for the earlier-time derivatives of the correlations at positive times extend to the right derivative at $s=0$. By Jensen's inequality, we have
\[
 |\bar r_N(t,s)-\partial_s\bar c_N(t,s)|^2\le\E_J|r_H(t,s)-\partial_sc_H(t,s)|^2\le\frac1N\E_JI_H(s)=a_N(s),
\]
which is \eqref{eq:annealed-fdt-bound}.

For \eqref{eq:finite-dimensional-integrated-fdt}, note that $\bar c_N(t,t)=1$ because $\norm{X_t}^2=N$. The function $u\mapsto\bar c_N(t,u)$ is continuous on $[s,t]$, with a bounded derivative on $[s,t)$; in particular, it is Lipschitz and is the integral of its derivative. Integrating \eqref{eq:annealed-fdt-bound} over $u\in[s,t]$, we obtain
\[
 \Bigl|1-\bar c_N(t,s)-\int_s^t\bar r_N(t,u)\,\dd u\Bigr|
 \le\int_s^t\sqrt{a_N(u)}\,\dd u,
\]
and one more application of the Cauchy--Schwarz inequality completes the proof.
\end{proof}

\section{An integrable Fisher majorant}\label{sec:fisher-majorant}

Recall that the landscape cutoff $A_p$ of \cref{lem:landscape-moments} satisfies $\Prob(Q_N>A_p)\le e^{-c_pN}$. We write $G_N=\{Q_N\le A_p\}$ for the corresponding event on the disorder and set
\[
 K=2\beta p^2A_p.
\]
The constant $K$ depends only on $p$ and $\beta$. We extract the subsequence in the following theorem once, before any later waiting-time sequence, stationary cluster, or finite-fork profile is chosen, and we keep the same strictly increasing dimension subsequence $(N_j)$ for the rest of the paper. We take all limits in $N$ at fixed times; the long-time statements concern only the limiting function $a_F$ and the deterministic CK solution.

\begin{theorem}[Fisher control of the backward FDT defect]\label{thm:fisher-clock}
There exist a dimension subsequence $(N_j)$ and a nonnegative, locally bounded, right-continuous function $a_F$ on $[0,\infty)$ such that
\begin{equation}
 a_{N_j}\longrightarrow a_F\quad\text{in }L^1_{\mathrm{loc}}([0,\infty)),
 \qquad \int_0^\infty a_F(u)\,\dd u\le\frac{\beta^2}{2}.
 \label{eq:fisher-density-budget}
\end{equation}
The function $e^{-Kt}a_F(t)$ is nonincreasing. In particular, for $s\ge\varepsilon>0$,
\begin{equation}
 a_F(s)\le\frac{K}{1-e^{-K\varepsilon}}\int_{s-\varepsilon}^sa_F(u)\,\dd u,
 \qquad \sup_{s\ge T}a_F(s)\longrightarrow0\quad(T\to\infty).
 \label{eq:fisher-density-tail}
\end{equation}
For every $t>s>0$, the earlier-time derivative $\partial_2C(t,s)=\partial_sC(t,s)$ exists, is continuous, and satisfies
\begin{equation}
 |\partial_2C(t,s)-R(t,s)|^2\le a_F(s).
 \label{eq:fisher-global-fdt}
\end{equation}
Every row $C(t,\cdot\,)$ is absolutely continuous on $[0,t]$, and for $t\ge v\ge u\ge0$,
\begin{equation}
 \Bigl|C(t,v)-C(t,u)-\int_u^vR(t,s)\,\dd s\Bigr|
 \le\int_u^v\sqrt{a_F(s)}\,\dd s.
 \label{eq:fisher-global-row}
\end{equation}
Finally, $e_\beta$ is locally absolutely continuous, and
\begin{equation}
 |e_\beta'(s)|^2\le p^2A_p^2a_F(s)\quad\text{for almost every }s\ge0.
 \label{eq:fisher-energy-density}
\end{equation}
We write
\begin{equation}
 \mu(\dd s)=a_F(s)\,\dd s,
 \qquad \mu([0,\infty))\le\beta^2/2.
 \label{eq:entropy-measure}
\end{equation}
\end{theorem}

\begin{proof}
The proof has six steps. The first three construct $a_F$. In Step~1, we use the Bochner formula to show that, on the event $G_N$, the Fisher information grows at most like $e^{Kt}$. Then the average of the Fisher information over $G_N$, multiplied by $e^{-Kt}$, is a bounded nonincreasing function of $t$, and we use Helly's selection theorem to extract a subsequence along which these functions converge (Step~2). In Step~3, we show that the Fisher information on the complement of $G_N$ has total mass tending to zero, which leads to \eqref{eq:fisher-density-budget} and \eqref{eq:fisher-density-tail}. In the last three steps, we transfer the finite-dimensional estimates to the CK solution. In Step~4, we pass the integrated form of \cref{lem:finite-dimensional-fdt} to the limit, which gives \eqref{eq:fisher-global-row}, together with \eqref{eq:fisher-global-fdt} for almost every $s$. In Step~5, we deduce from the Green covariance identity that $\partial_2C$ is continuous, which upgrades the bound to every $s$. In Step~6, we treat the energy in the same way.

\paragraph{Step 1: the Bochner identity and growth of the Fisher information.}
Fix the disorder. Let $f_t=\rho_t/g_N$ and $\psi_t=\log f_t$, and let $\nu_N(\dd x)=g_N(x)\,\omega_N(\dd x)$ be the Gibbs measure. Then $v^{\mathrm{cur}}_t=\nabla_S\psi_t$ and $I_H(t)=\int f_t\Gamma(\psi_t)\,\dd\nu_N$. Recall that the generator $\mathcal L_N$ is symmetric with respect to $\nu_N$, with $\int f\mathcal L_Ng\,\dd\nu_N=-\int\Gamma(f,g)\,\dd\nu_N$. In terms of $f_t$, the forward equation reads $\partial_tf_t=\mathcal L_Nf_t$. As $\mathcal L_N(e^{\psi})=e^{\psi}(\mathcal L_N\psi+\Gamma(\psi))$, we also have $\partial_t\psi_t=\mathcal L_N\psi_t+\Gamma(\psi_t)$. Differentiating $I_H$, we find
\[
 I_H'(t)=\int(\mathcal L_Nf_t)\Gamma(\psi_t)\,\dd\nu_N
 +2\int f_t\Gamma(\psi_t,\mathcal L_N\psi_t)\,\dd\nu_N
 +2\int f_t\Gamma(\psi_t,\Gamma(\psi_t))\,\dd\nu_N.
\]
By symmetry, the first integral equals $\int f_t\mathcal L_N\Gamma(\psi_t)\,\dd\nu_N$. It also equals $-\int\Gamma(f_t,\Gamma(\psi_t))\,\dd\nu_N=-\int f_t\Gamma(\psi_t,\Gamma(\psi_t))\,\dd\nu_N$, because $\nabla_Sf_t=f_t\nabla_S\psi_t$. The third integral is therefore $-2\int f_t\mathcal L_N\Gamma(\psi_t)\,\dd\nu_N$, and we obtain
\[
 I_H'(t)=-\int f_t\mathcal L_N\Gamma(\psi_t)\,\dd\nu_N+2\int f_t\Gamma(\psi_t,\mathcal L_N\psi_t)\,\dd\nu_N
 =-2\int f_t\Gamma_2(\psi_t)\,\dd\nu_N,
\]
where $\Gamma_2(\psi)=\frac12\mathcal L_N\Gamma(\psi)-\Gamma(\psi,\mathcal L_N\psi)$. For $\mathcal L_N=\Delta_S+\beta\ip{\nabla_SH_N}{\nabla_S\,\cdot\,}$, the Bochner formula reads $\Gamma_2(\psi)=\norm{\nabla_S^2\psi}_{\HS}^2+(\operatorname{Ric}-\beta\nabla_S^2H_N)(\nabla_S\psi,\nabla_S\psi)$. Using $f_t\,\dd\nu_N=\rho_t\,\dd\omega_N$, we conclude that
\begin{equation}
 I_H'(t)=-2\int\rho_t\Bigl\{
 \norm{\nabla_S^2\psi_t}_{\HS}^2
 +(\operatorname{Ric}-\beta\nabla_S^2H_N)(\nabla_S\psi_t,\nabla_S\psi_t)\Bigr\}\,\dd\omega_N.
 \label{eq:fisher-bochner-identity}
\end{equation}
This identity is asserted for $t>0$. Its differentiations and integrations by parts use the spatial $C^3$ density, positivity on the compact sphere, and joint continuity of the forward slope and its spatial gradient. In particular, we compute the time derivative of the relative score by differentiating the intrinsic forward equation along the sphere. Continuity of $I_H$ at zero will suffice to include the initial time in the resulting comparison.

For $N\ge2$, the sphere $S_N$ of radius $\sqrt N$ has $\operatorname{Ric}=\frac{N-2}{N}$ times its metric, which is nonnegative, and $\norm{\nabla_S^2H_N}_{\op}\le p^2Q_N$ by \eqref{eq:landscape-derivative-bounds}. Then \eqref{eq:fisher-bochner-identity} implies that $I_H'\le2\beta p^2Q_NI_H$ at positive times. On the event $G_N$, we have $I_H'\le KI_H$, so $e^{-Kt}I_H(t)$ is nonincreasing for $t>0$. Its continuity at zero extends this monotonicity to $[0,\infty)$. We define
\[
 b_N(t)=\frac1N\E_J\bigl[\mathbf 1_{G_N}I_H(t)\bigr],
 \qquad
 \tilde b_N(t)=e^{-Kt}b_N(t).
\]
As an average of nonincreasing functions, $\tilde b_N$ is nonincreasing; no differentiation under $\E_J$ is needed. At time zero, $\nabla_S\psi_0=-\beta\nabla_SH_N$. Differentiating the covariance $N(x\cdot y/N)^p$ in both variables and setting $y=x\in S_N$, we see that $\E_J[\nabla H_N(x)\nabla H_N(x)^{\mathsf T}]=pI+\frac{p(p-1)}{N}xx^{\mathsf T}$. The projection $P_x$ removes the second term, because $P_xx=0$. We obtain $\E_J\norm{\nabla_SH_N(x)}^2=p\operatorname{tr}P_x=p(N-1)$, and by Tonelli's theorem,
\[
 a_N(0)=\frac{\beta^2}{N}\int\E_J\norm{\nabla_SH_N}^2\,\dd\omega_N=\beta^2p\,\frac{N-1}{N}.
\]
We have shown that $\tilde b_N$ is nonincreasing and $0\le\tilde b_N(t)\le\tilde b_N(0)\le a_N(0)\le\beta^2p$ for all $t\ge0$. We also note that the bound $I_H'\le2\beta p^2Q_NI_H$ holds for every disorder. It implies $I_H(t)\le I_H(0)e^{2\beta p^2Q_Nt}\le\beta^2p^2NQ_N^2e^{2\beta p^2Q_Nt}$, and by \cref{lem:landscape-moments}, $a_N$ is finite and bounded on compact intervals, uniformly in $N$.

\paragraph{Step 2: Helly selection.}
We use the Helly diagonal-selection argument; compare the proof of \cite[Theorem~25.9, p.~336]{billingsley1995}.
The sequences $(\tilde b_N(\theta))_N$ are bounded, and a diagonal extraction over the positive rationals gives a subsequence $(N_j)$ along which $\tilde b_{N_j}(\theta)$ converges for every rational $\theta>0$. The limits form a nonincreasing function on the positive rationals with values in $[0,\beta^2p]$. Let $\tilde b(t)$ be the right-continuous function on $[0,\infty)$ defined by the limit of these values as $\theta\downarrow t$ through the rationals. Let $t>0$ be a continuity point of $\tilde b$. For rationals $\theta<t<\theta'$, we have $\tilde b_{N_j}(\theta')\le\tilde b_{N_j}(t)\le\tilde b_{N_j}(\theta)$ by monotonicity. Letting $j\to\infty$ and then $\theta\uparrow t$, $\theta'\downarrow t$, we find that $\tilde b_{N_j}(t)\to\tilde b(t)$. Since a monotone function has at most countably many discontinuities, we conclude that $\tilde b_{N_j}\to\tilde b$ almost everywhere. We define
\[
 a_F(t)=e^{Kt}\tilde b(t).
\]
Then $a_F$ is nonnegative and right-continuous, $a_F(t)\le\beta^2pe^{Kt}$, and $e^{-Kt}a_F(t)=\tilde b(t)$ is nonincreasing. For every $M<\infty$, we have $0\le b_{N_j}(t)\le\beta^2pe^{KM}$ for $t\in[0,M]$, and by dominated convergence, $b_{N_j}\to a_F$ in $L^1([0,M])$. The value of $a_F$ at zero is the right limit; no convergence of $b_N(0)$ to $a_F(0)$ is used. We have shown that
\[
 b_{N_j}\longrightarrow a_F\quad\text{in }L^1_{\mathrm{loc}}([0,\infty)).
\]

\paragraph{Step 3: the Fisher mass on the bad event.}
The difference $a_N-b_N=N^{-1}\E_J[\mathbf 1_{G_N^c}I_H]$ is nonnegative. By the quenched budget in \eqref{eq:quenched-entropy-dissipation} and the bound $D_H(0)\le2\beta NQ_N$ of \cref{lem:entropy-budget}, we have
\[
\begin{aligned}
 \int_0^\infty(a_N-b_N)(t)\,\dd t
 &=\frac1N\E_J\Bigl[\mathbf 1_{G_N^c}\int_0^\infty I_H(t)\,\dd t\Bigr]\\
 &\le2\beta\E_J\bigl[Q_N\mathbf 1_{G_N^c}\bigr]
 \le2\beta\bigl(\E_JQ_N^2\bigr)^{1/2}\Prob(G_N^c)^{1/2}.
\end{aligned}
\]
The right side tends to zero by \cref{lem:landscape-moments}. Together with Step~2, this implies that $a_{N_j}\to a_F$ in $L^1_{\mathrm{loc}}([0,\infty))$. For every $M$, we have $\int_0^Ma_F=\lim_j\int_0^Ma_{N_j}\le\beta^2/2$ by \cref{lem:entropy-budget}, and monotone convergence in $M$ proves \eqref{eq:fisher-density-budget}.

The monotonicity of $e^{-Kt}a_F(t)$ means that $a_F(u)\ge e^{-K(s-u)}a_F(s)$ for $u\le s$. Integrating this inequality over $u\in[s-\varepsilon,s]$, we have
\[
 \int_{s-\varepsilon}^sa_F(u)\,\dd u\ge a_F(s)\int_0^\varepsilon e^{-Kw}\,\dd w=a_F(s)\,\frac{1-e^{-K\varepsilon}}{K},
\]
which is the first part of \eqref{eq:fisher-density-tail}; here $K>0$ because $\beta>0$ and $A_p>0$. Taking $\varepsilon=1$, we have $\sup_{s\ge T}a_F(s)\le K(1-e^{-K})^{-1}\mu([T-1,\infty))$ for $T\ge1$, and the right side tends to zero because $\mu$ is a finite measure. This proves \eqref{eq:fisher-density-tail}. We note that integrability of $a_F$ alone would not give this uniform decay; it comes from the monotonicity of $e^{-Kt}a_F(t)$.

\paragraph{Step 4: passage of the row inequality to the limit.}
Fix $t\ge v\ge u\ge0$. Recall from the proof of \cref{lem:finite-dimensional-fdt} that $\bar c_N(t,\cdot\,)$ is the integral of its derivative on $[u,t]$. Integrating \eqref{eq:annealed-fdt-bound} over $[u,v]$, we have
\[
 \Bigl|\bar c_N(t,v)-\bar c_N(t,u)-\int_u^v\bar r_N(t,s)\,\dd s\Bigr|
 \le\int_u^v\sqrt{a_N(s)}\,\dd s.
\]
We let $N=N_j\to\infty$ in each term. By \cref{prop:hard-bridge}, $\bar c_N(t,v)=\E C_N(t,v)\to C(t,v)$, and similarly at $(t,u)$. The response integral equals $\E[X_t\cdot(B_v-B_u)]/(\sqrt2N)$ by \cref{lem:brown-force-response}; it is a difference of the Brownian observables of \cref{prop:hard-bridge}, whose expectations converge to $\int_u^vR(t,s)\,\dd s$. For the right side, we use the elementary inequality $|\sqrt x-\sqrt y|^2\le|x-y|$ for $x,y\ge0$. On every compact interval $I$, it yields
\[
 \int_I\bigl|\sqrt{a_{N_j}}-\sqrt{a_F}\bigr|^2\le\int_I|a_{N_j}-a_F|\longrightarrow0,
\]
and by the Cauchy--Schwarz inequality, $\sqrt{a_{N_j}}\to\sqrt{a_F}$ in $L^1(I)$. This proves \eqref{eq:fisher-global-row} for every $t\ge v\ge u\ge0$. We fixed the subsequence in Step~2, before choosing the endpoints, and the correlations and the primitives of the response converge along the full sequence of dimensions. No further extraction is involved.

Next, fix $t$ and let $F_t(s)=C(t,s)-\int_0^sR(t,w)\,\dd w$ for $0\le s\le t$. By \eqref{eq:fisher-global-row}, we have, for every finite family of disjoint intervals $(u_i,v_i)\subset[0,t]$,
\[
 \sum_i|F_t(v_i)-F_t(u_i)|\le\int_{\bigcup_i(u_i,v_i)}\sqrt{a_F(s)}\,\dd s.
\]
The right side is small when the total length of the intervals is small, because $\sqrt{a_F}$ is locally bounded. Hence $F_t$ is absolutely continuous. As $s\mapsto\int_0^sR(t,w)\,\dd w$ is continuously differentiable, the row $C(t,\cdot\,)$ is absolutely continuous on $[0,t]$. At every point $s\in(0,t)$ where $F_t$ is differentiable and which is a Lebesgue point of $\sqrt{a_F}$, dividing \eqref{eq:fisher-global-row} by $v-u$ and letting the interval shrink to $s$, we have $|F_t'(s)|\le\sqrt{a_F(s)}$. We have shown that
\[
 |\partial_2C(t,s)-R(t,s)|^2\le a_F(s)\quad\text{for almost every }s\in(0,t).
\]
The exceptional set may depend on $t$; this causes no difficulty, since we carry out Step~5 on each fixed row with the fixed function $a_F$.

\paragraph{Step 5: the continuous representative of the earlier derivative.}
Let $\dot R(s,u)=\partial_1R(s,u)$. Recall from \eqref{eq:ck-response} and \cref{prop:backward-response} that $\dot R$ is continuous on every compact causal triangle, with one-sided values on the diagonal. For $t>s$, the Green identity \eqref{eq:green-covariance} reads
\[
\begin{aligned}
 C(t,s)={}&R(t,0)R(s,0)+2\int_0^sR(t,u)R(s,u)\,\dd u\\
 &+\beta^2p\int_0^tR(t,u)\int_0^sR(s,v)C(u,v)^{p-1}\,\dd v\,\dd u.
\end{aligned}
\]
The variable $s$ enters only through $R(s,0)$, $R(s,u)$, $R(s,v)$, and the upper limits of the integrals over $u$ and $v$. Differentiating in $s$ and using $R(s,s)=1$, we obtain, for $t>s>0$,
\begin{align}
 \partial_2C(t,s)={}&R(t,0)\dot R(s,0)+2R(t,s)
 +2\int_0^sR(t,u)\dot R(s,u)\,\dd u\notag\\
 &+\beta^2p\int_0^tR(t,u)
 \Bigl\{C(u,s)^{p-1}+\int_0^s\dot R(s,v)C(u,v)^{p-1}\,\dd v\Bigr\}\,\dd u.
 \label{eq:green-earlier-derivative}
\end{align}
All integrands are continuous on the fixed compact sets of integration, which justifies differentiation under the integral sign. No derivative of $C(u,v)^{p-1}$ is needed, because $u$ and $v$ are integration variables that do not depend on $s$. The right side of \eqref{eq:green-earlier-derivative} is jointly continuous on the open triangle $t>s>0$. Hence, for each fixed $t$, the row $C(t,\cdot\,)$ is continuously differentiable on $(0,t)$, and its classical derivative agrees with the almost-everywhere derivative of Step~4.

Fix $t>s>0$, and choose points $s_n\downarrow s$ with $s_n<t$ at which the bound of Step~4 holds. Then $|\partial_2C(t,s_n)-R(t,s_n)|^2\le a_F(s_n)$. As $n\to\infty$, the left side converges to $|\partial_2C(t,s)-R(t,s)|^2$ by the continuity of \eqref{eq:green-earlier-derivative} and of $R$, and $a_F(s_n)\to a_F(s)$ because $a_F$ is right-continuous. This proves \eqref{eq:fisher-global-fdt}.

\paragraph{Step 6: the energy.}
Let $e_H(s)=N^{-1}\int\rho_sH_N\,\dd\omega_N$, so that $e_N(s)=\E_Je_H(s)=N^{-1}\E H_N(X_s)$. By the first identity in \eqref{eq:quenched-entropy-dissipation} and an integration by parts, we have
\[
 e_H'(s)=-\frac1N\int\rho_s\ip{v^{\mathrm{cur}}_s}{\nabla_SH_N}\,\dd\omega_N.
\]
Using $\norm{\nabla_SH_N}\le pQ_N\sqrt N$ and the Cauchy--Schwarz inequality, we find that $|e_H'(s)|\le pQ_N\sqrt{I_H(s)/N}$. We split the average over the disorder according to $G_N$. On $G_N$, we have $Q_N\le A_p$, and Jensen's inequality yields $\E_J[\mathbf 1_{G_N}pQ_N\sqrt{I_H(s)/N}]\le pA_p\sqrt{b_N(s)}$. On the complement, by the Cauchy--Schwarz inequality, $\E_J[\mathbf 1_{G_N^c}pQ_N\sqrt{I_H(s)/N}]\le\sqrt{\eta_N(s)}$, where
\[
 \eta_N(s)=\frac{p^2}{N}\E_J\bigl[Q_N^2I_H(s)\mathbf 1_{G_N^c}\bigr].
\]
By \eqref{eq:quenched-entropy-dissipation} and $D_H(0)\le2\beta NQ_N$, the total mass of $\eta_N$ satisfies
\[
 \int_0^\infty\eta_N(s)\,\dd s\le2\beta p^2\E_J\bigl[Q_N^3\mathbf 1_{G_N^c}\bigr]
 \le2\beta p^2\bigl(\E_JQ_N^6\bigr)^{1/2}\Prob(G_N^c)^{1/2}\longrightarrow0.
\]
For fixed $u\le v$, Tonelli's theorem and the two bounds above give
\[
 |e_N(v)-e_N(u)|\le\int_u^v\E_J|e_H'(s)|\,\dd s
 \le pA_p\int_u^v\sqrt{b_N(s)}\,\dd s+\Bigl((v-u)\int_0^\infty\eta_N(s)\,\dd s\Bigr)^{1/2}.
\]
We let $N=N_j\to\infty$. By \cref{prop:energy-bridge}, $e_N(v)\to e_\beta(v)$ and $e_N(u)\to e_\beta(u)$; by Step~2 and the argument of Step~4, $\int_u^v\sqrt{b_{N_j}}\to\int_u^v\sqrt{a_F}$; and the last term tends to zero. We conclude that
\[
 |e_\beta(v)-e_\beta(u)|\le pA_p\int_u^v\sqrt{a_F(s)}\,\dd s,\qquad 0\le u\le v.
\]
As in Step~4, we see that $e_\beta$ is locally absolutely continuous and that $|e_\beta'(s)|\le pA_p\sqrt{a_F(s)}$ at almost every $s$, which is \eqref{eq:fisher-energy-density}. Here we use only the convergence of energies at fixed times, not the convergence of their derivatives.
\end{proof}

We note two features of \cref{thm:fisher-clock}. First, the majorant $a_F$ is produced by one subsequence of dimensions and need not be unique; the inequalities it satisfies concern the unique deterministic CK solution; the bounds on the defect hold at every pair of times, and the bound on the derivative of the energy holds for almost every time. Second, \eqref{eq:fisher-global-fdt} and \eqref{eq:fisher-density-tail} imply that the defect $\partial_2C(t,s)-R(t,s)$ tends to zero uniformly in $t$ as $s\to\infty$. This does not identify the slow fluctuation--dissipation ratio: both $R$ and $\partial_2C$ may tend to zero while $R/\partial_2C$ has a nontrivial limit, and the defect may accumulate over a diverging interval of earlier times. For $\beta>\bd$, \cref{resp:integrated-chapter} shows that a positive amount of defect does accumulate in this way.

The bound \eqref{eq:fisher-global-fdt} also holds at the two ends of each row, with one-sided derivatives. We record this for later use.

\begin{remark}[The ends of a row]\label{rem:fisher-endpoints}
Let $t>0$. The one-sided derivatives $\partial_2C(t,0+)$ and $\partial_2C(t,t-)$ of $s\mapsto C(t,s)$ exist, and
\[
 |\partial_2C(t,0+)-R(t,0)|^2\le a_F(0),
 \qquad
 \partial_2C(t,t-)=R(t,t)=1.
\]
In other words, \eqref{eq:fisher-global-fdt} holds for $0\le s\le t$, with derivatives at $s=0$ and $s=t$ taken from the interior of the row. Indeed, let $\Delta(t,s)$ denote the right side of \eqref{eq:green-earlier-derivative}. It is continuous for $0\le s\le t$, and by the mean value theorem, the one-sided derivatives at the ends of the row exist and equal $\Delta(t,0)$ and $\Delta(t,t)$. The first bound follows by choosing $s_n\downarrow0$ in \eqref{eq:fisher-global-fdt} and using the right continuity of $a_F$ at zero. For the second, we differentiate the diagonal Green identity
\[
 1=C(t,t)=R(t,0)^2+2\int_0^tR(t,u)^2\,\dd u
 +\beta^2p\int_0^t\int_0^tR(t,u)R(t,v)C(u,v)^{p-1}\,\dd v\,\dd u
\]
in $t$. Using $R(t,t)=1$ and the symmetry of $C$, we find that the initial, white-noise, and colored-noise terms have derivatives $2R(t,0)\dot R(t,0)$, $2+4\int_0^tR(t,u)\dot R(t,u)\,\dd u$, and twice the colored-noise term of $\Delta(t,t)$, respectively. Comparing with \eqref{eq:green-earlier-derivative}, we see that the derivative of the right side equals $2\Delta(t,t)-2$; the white-noise boundary term is counted once on the diagonal, while $2\Delta(t,t)$ counts it twice. Since the left side is constant, $\Delta(t,t)=1$.
\end{remark}

\section{Flatness in the waiting time}\label{sec:waiting-time-flatness}

The bound \eqref{eq:fisher-global-fdt} controls the earlier observation time with the terminal time fixed. To compare two correlations with the same lag at different waiting times, we also need to control the effect of the probability current on a conditional-mean observable. A bound of this kind for one-time observables appears in \cite[arXiv version~2, equation~(15), p.~3]{cdk1997}.

\begin{lemma}[Semigroup gradient bound and the waiting-time derivative]\label{lem:waiting-clock-current}
For every $\tau\ge0$ and $x\in S_N$,
\begin{equation}
 \norm{D_Sm_\tau(x)}_{\op}\le e^{\beta p^2Q_N\tau}.
 \label{eq:conditional-mean-gradient}
\end{equation}
Let $\Psi_\tau(x)=N^{-1}\ip{x}{m_\tau(x)}$. Then
\begin{align}
 \norm{\nabla_S\Psi_\tau(x)}
 &\le N^{-1/2}\bigl(1+e^{\beta p^2Q_N\tau}\bigr),
 \label{eq:waiting-observable-gradient}\\
 \frac{\dd}{\dd w}c_H(w+\tau,w)
 &=-\int\rho_w\ip{v^{\mathrm{cur}}_w}{\nabla_S\Psi_\tau}\,\dd\omega_N.
 \label{eq:waiting-clock-current}
\end{align}
The waiting-time identity holds for all $w\ge0$, with a right derivative at $w=0$.
For $0\le\tau\le L$ and $0\le w_1\le w_2$,
\begin{equation}
 |C(w_2+\tau,w_2)-C(w_1+\tau,w_1)|
 \le B_L\int_{w_1}^{w_2}\sqrt{a_F(u)}\,\dd u,
 \qquad B_L=1+e^{\beta p^2A_pL}.
 \label{eq:ck-measure-waiting-shift}
\end{equation}
\end{lemma}

\begin{proof}
We begin with the gradient bound. Let $K_N=\beta p^2Q_N$. Recall from Step~1 of the proof of \cref{thm:fisher-clock} that $\operatorname{Ric}\ge0$ and $\norm{\nabla_S^2H_N}_{\op}\le p^2Q_N$. By the Bochner formula, we have the curvature bound $\Gamma_2(f)\ge-K_N\Gamma(f)$ for smooth $f$. Fix $\tau$, $x$, and a smooth function $f$, and let $\Phi(u)=P_u\Gamma(P_{\tau-u}f)(x)$ for $0\le u\le\tau$. Using $\partial_u\Gamma(P_{\tau-u}f)=-2\Gamma(P_{\tau-u}f,\mathcal L_NP_{\tau-u}f)$, we find
\[
 \Phi'(u)=P_u\bigl[\mathcal L_N\Gamma(P_{\tau-u}f)-2\Gamma(P_{\tau-u}f,\mathcal L_NP_{\tau-u}f)\bigr](x)
 =2P_u\Gamma_2(P_{\tau-u}f)(x)\ge-2K_N\Phi(u),
\]
where we used that $P_u$ preserves order. By Gronwall's inequality, $\Phi(\tau)\ge e^{-2K_N\tau}\Phi(0)$, that is,
\[
 \norm{\nabla_SP_\tau f}^2\le e^{2K_N\tau}P_\tau\norm{\nabla_Sf}^2.
\]
This is the classical gradient bound under a lower bound on the curvature; see \cite[Theorem~3.2.3, p.~144]{bgl2014}.
We apply it to $f_z(x)=\ip{z}{x}$ for a unit vector $z\in\R^N$. Then $\norm{\nabla_Sf_z}=\norm{P_xz}\le1$ and $\nabla_SP_\tau f_z=\nabla_S\ip{z}{m_\tau}=(D_Sm_\tau)^*z$. By the gradient bound, we now have $\norm{(D_Sm_\tau(x))^*z}\le e^{K_N\tau}$. Taking the supremum over unit vectors $z$ proves \eqref{eq:conditional-mean-gradient}. This bound has no factor depending on the dimension; applying the scalar estimate only to the $N$ coordinate functions and summing would lose such a factor.

Next, we have $\nabla_S\Psi_\tau=N^{-1}[P_xm_\tau+(D_Sm_\tau)^*x]$. As $\norm x=\sqrt N$ and $\norm{m_\tau}\le\sqrt N$, \eqref{eq:conditional-mean-gradient} implies \eqref{eq:waiting-observable-gradient}. At fixed lag $\tau$, we have $c_H(w+\tau,w)=\int\rho_w\Psi_\tau\,\dd\omega_N$ by \eqref{eq:quenched-correlation}. Differentiating at $w>0$ and using the first identity in \eqref{eq:quenched-entropy-dissipation} and an integration by parts, we obtain \eqref{eq:waiting-clock-current}. The current is continuous up to zero, and the exact current integral over a waiting-time interval therefore gives the same identity with the right derivative at $w=0$.

It remains to prove \eqref{eq:ck-measure-waiting-shift}. Let $0\le\tau\le L$. On the event $G_N$, we have $Q_N\le A_p$, and by \eqref{eq:waiting-observable-gradient}, \eqref{eq:waiting-clock-current}, and the Cauchy--Schwarz inequality,
\[
 \Bigl|\frac{\dd}{\dd w}c_H(w+\tau,w)\Bigr|\le B_L\sqrt{I_H(w)/N}.
\]
Averaging over $G_N$ and using Jensen's inequality, we obtain $\E_J[\mathbf 1_{G_N}\sqrt{I_H(w)/N}]\le\sqrt{b_N(w)}$, with $b_N$ as in Step~1 of the proof of \cref{thm:fisher-clock}. On the complement of $G_N$, we use only $|c_H|\le1$ at the two endpoints; no control of the gradient is needed there. Combining these bounds, we find
\[
 |\bar c_N(w_2+\tau,w_2)-\bar c_N(w_1+\tau,w_1)|
 \le B_L\int_{w_1}^{w_2}\sqrt{b_N(w)}\,\dd w+2\Prob(G_N^c).
\]
We now let $N=N_j\to\infty$. The correlations converge by \cref{prop:hard-bridge}, $\int_{w_1}^{w_2}\sqrt{b_{N_j}}\to\int_{w_1}^{w_2}\sqrt{a_F}$ by Steps~2 and~4 of the proof of \cref{thm:fisher-clock}, and $\Prob(G_N^c)\le e^{-c_pN}\to0$. This proves \eqref{eq:ck-measure-waiting-shift}.
\end{proof}

We now collect the uniform late estimates that follow from \cref{thm:fisher-clock,lem:waiting-clock-current}. Since $\mu$ is a finite measure, they show that, on late windows of bounded length, the integrated FDT relation holds, and the correlation at fixed lag and the energy do not depend on the waiting time, up to errors that vanish as the window moves to infinity.

\begin{corollary}[Uniform late estimates]\label{cor:uniform-late-estimates}
Fix $L,\Theta\ge0$. For every $T\ge0$,
\begin{equation}
 \sup_{s\ge T,\,0\le\tau\le L}
 \Bigl|1-C(s+\tau,s)-\int_s^{s+\tau}R(s+\tau,u)\,\dd u\Bigr|
 \le\sqrt{L\mu([T,\infty))},
 \label{eq:uniform-late-integrated-fdt}
\end{equation}
and for every $T\ge\Theta$,
\begin{align}
 \sup_{s\ge T,\,|\theta|\le\Theta,\,0\le\tau\le L}
 |C(s+\theta+\tau,s+\theta)-C(s+\tau,s)|
 &\le B_L\sqrt{\Theta\mu([T-\Theta,\infty))},
 \label{eq:uniform-late-clock-stationarity}\\
 \sup_{s\ge T,\,|\theta|\le\Theta}|e_\beta(s+\theta)-e_\beta(s)|
 &\le pA_p\sqrt{\Theta\mu([T-\Theta,\infty))}.
 \label{eq:uniform-late-energy-flatness}
\end{align}
The three right sides tend to zero as $T\to\infty$ with $L$ and $\Theta$ fixed.
\end{corollary}

\begin{proof}
For \eqref{eq:uniform-late-integrated-fdt}, we use \eqref{eq:fisher-global-row} on the row $t=s+\tau$ between $u=s$ and $v=t$, together with $C(t,t)=1$. Its right side is at most $\sqrt{\tau\mu([s,s+\tau])}\le\sqrt{L\mu([T,\infty))}$ by the Cauchy--Schwarz inequality. For \eqref{eq:uniform-late-clock-stationarity}, we apply \eqref{eq:ck-measure-waiting-shift} with $w_1$ and $w_2$ the smaller and the larger of $s$ and $s+\theta$. The interval $[w_1,w_2]$ has length $|\theta|\le\Theta$ and lies in $[T-\Theta,\infty)$, and the Cauchy--Schwarz inequality gives the bound as before. The proof of \eqref{eq:uniform-late-energy-flatness} is the same, using the increment bound $|e_\beta(v)-e_\beta(u)|\le pA_p\int_u^v\sqrt{a_F}$ from Step~6 of the proof of \cref{thm:fisher-clock}. The last assertion holds because $\mu([T,\infty))\to0$ as $T\to\infty$.
\end{proof}

\section{Stationary translated clusters}\label{sec:translated-clusters}

Recall the translated kernels $C_T$, $R_T$, and $\lambda_T$ of \eqref{eq:translated-kernels}, with correlations considered on squares of offsets and responses only on closed causal triangles $\{a\ge b\}$. We say that a triple $(C_\infty,R_\infty,\lambda_\infty)$ is a \emph{translated cluster} if there exists a sequence $T_n\to\infty$ such that $C_{T_n}\to C_\infty$, $R_{T_n}\to R_\infty$, and $\lambda_{T_n}\to\lambda_\infty$ uniformly on compact sets of offsets in their respective domains.

\begin{proposition}[Stationary FDT translated clusters]\label{prop:stationary-clusters}
The translated family $(C_T,R_T,\lambda_T)$ is precompact for locally uniform convergence on compact offset sets. Every translated cluster has the form
\[
 C_\infty(a,b)=c(|a-b|),\qquad
 R_\infty(a,b)=r(a-b)\quad(a\ge b),\qquad
 \lambda_\infty(a)=\ell,
\]
where $\ell$ is a constant and
\begin{equation}
 c(0)=1,\qquad r(0)=1,\qquad c'=-r,\qquad
 0\le c\le1,\quad r\ge0.
 \label{eq:cluster-fdt-properties}
\end{equation}
In particular, $c$ is nonincreasing, its limit $q=\lim_{\tau\to\infty}c(\tau)$ belongs to $[0,1)$, and $\int_0^\infty r(u)\,\dd u=1-q$. Along the same sequence $T_n$, the energy translates $e_\beta(T_n+a)$ converge locally uniformly to a constant $e_\infty$, and $\ell=1+p\beta e_\infty$. Moreover, $r$ is continuously differentiable on $[0,\infty)$ and satisfies the closed stationary response equation
\begin{equation}
 r'(\tau)=-\ell r(\tau)
 +\beta^2p(p-1)\int_0^\tau c(\tau-u)^{p-2}r(\tau-u)r(u)\,\dd u.
 \label{eq:stationary-cluster-response}
\end{equation}
\end{proposition}

We call the triple $(c,r,\ell)$ of \cref{prop:stationary-clusters} a \emph{stationary cluster}, and $q$ its \emph{plateau}. Different sequences $T_n$ may a priori lead to different stationary clusters.

\begin{proof}
The proof has four steps. We first recall compactness. We then show that the correlation of every translated cluster is stationary, and that its response is stationary and satisfies FDT. Finally, we identify the multiplier and pass to the limit in the response equation. The order of limits is the same throughout: we have already taken the dimension limit in \cref{thm:fisher-clock}, the window of offsets is fixed, and only $T_n\to\infty$ varies.

\paragraph{Step 1: compactness.}
The precompactness of the translated family is proved in \cref{prop:response-compactness}. The bounds $0\le C_T\le1$ and $R_T\ge0$, the symmetry of $C_T$, and the diagonal values $C_T(a,a)=R_T(a,a)=1$ pass to every translated cluster.

\paragraph{Step 2: stationarity of the correlation.}
Fix a sequence $T_n\to\infty$ along which the translated family converges to a translated cluster $(C_\infty,R_\infty,\lambda_\infty)$. Let $a\ge b$ and $\theta$ be fixed offsets. We have
\[
 C_{T_n}(a+\theta,b+\theta)-C_{T_n}(a,b)=C(s_n+\theta+\tau,s_n+\theta)-C(s_n+\tau,s_n),
\]
where $s_n=T_n+b$ and $\tau=a-b$.
Once $s_n\ge|\theta|$, we may apply \eqref{eq:uniform-late-clock-stationarity} with $L=\tau$, $\Theta=|\theta|$, and $T=s_n$, and its right side tends to zero because $s_n\to\infty$. We conclude that $C_\infty(a+\theta,b+\theta)=C_\infty(a,b)$. Let $c(\tau)=C_\infty(\tau,0)$ for $\tau\ge0$. Taking $\theta=-b$, we have $C_\infty(a,b)=c(a-b)$ for $a\ge b$, and by symmetry $C_\infty(a,b)=c(|a-b|)$. The function $c$ is continuous, $0\le c\le1$, and $c(0)=1$.

\paragraph{Step 3: the response and FDT.}
Let $a\ge b$. By \eqref{eq:uniform-late-integrated-fdt} with $s=T_n+b$, $\tau=a-b$, and $L=a-b$, we have
\[
 \Bigl|1-C_{T_n}(a,b)-\int_b^aR_{T_n}(a,u)\,\dd u\Bigr|\le\sqrt{(a-b)\,\mu([T_n+b,\infty))}\longrightarrow0.
\]
Since the response converges uniformly on the causal triangle containing the segment $\{(a,u):b\le u\le a\}$, we may pass to the limit inside the integral. This gives
\[
 \int_b^aR_\infty(a,u)\,\dd u=1-c(a-b).
\]
For $a=0$ and $b=-\tau$, this reads $1-c(\tau)=\int_{-\tau}^0R_\infty(0,u)\,\dd u$. As $R_\infty(0,\cdot\,)$ is continuous on $(-\infty,0]$, the function $c$ is continuously differentiable on $[0,\infty)$, with $c'(\tau)=-R_\infty(0,-\tau)$ and a one-sided derivative at $\tau=0$. Differentiating the general identity in $b$ with $a$ fixed, we have $R_\infty(a,b)=-c'(a-b)$ for $a\ge b$. We set $r=-c'$. Then $r$ is continuous and nonnegative, $r(0)=R_\infty(0,0)=1$, and $R_\infty(a,b)=r(a-b)$. This proves that the response is stationary and satisfies FDT; neither property was assumed.

As $r\ge0$ and $c\ge0$, the function $c$ is nonincreasing and bounded below, and the limit $q$ exists. As $r$ is continuous with $r(0)=1$, there exists $\delta>0$ such that $r\ge1/2$ on $[0,\delta]$. Then $q\le c(\delta)\le1-\delta/2<1$. Integrating $c'=-r$ gives $\int_0^\infty r(u)\,\dd u=\lim_{\tau\to\infty}(1-c(\tau))=1-q$.

\paragraph{Step 4: the multiplier and the response equation.}
By \eqref{eq:energy-identity}, we have $e_\beta(t)=(\lambda(t)-1)/(p\beta)$ at every time $t$. Then $e_\beta(T_n+a)\to(\lambda_\infty(a)-1)/(p\beta)$ locally uniformly in $a$. For fixed offsets $a$ and $a'$, \eqref{eq:uniform-late-energy-flatness} with $s=T_n+a'$, $\theta=a-a'$, and $\Theta=|a-a'|$ implies that $e_\beta(T_n+a)-e_\beta(T_n+a')\to0$. The limit of the energy translates is therefore a constant $e_\infty$, and $\lambda_\infty$ is the constant $\ell=1+p\beta e_\infty$.

Finally, we pass to the limit in the response equation. Fix $\tau\ge0$ and let $s=T_n$. Integrating \eqref{eq:ck-response} in its first variable from $s$ to $s+\tau$, using $R(s,s)=1$, and substituting $u=s+u'$ in the memory integral, we obtain
\[
\begin{aligned}
 R(s+\tau,s)-1=\int_0^\tau\Bigl[&-\lambda(s+v)R(s+v,s)
 +\beta^2p(p-1)\int_0^vC(s+v,s+u')^{p-2}\\
 &\quad\times R(s+v,s+u')R(s+u',s)\,\dd u'\Bigr]\dd v.
\end{aligned}
\]
Here we used $\xi''(x)=p(p-1)x^{p-2}$. All arguments of $C$, $R$, and $\lambda$ are translated times with offsets in $[0,\tau]$, in a fixed causal triangle. The memory integral of \eqref{eq:ck-response} runs only over times between the two observation times, and no other part of the history enters. By the locally uniform convergence along $T_n$, together with $C_\infty(a,b)=c(|a-b|)$, $R_\infty(a,b)=r(a-b)$, and $\lambda_\infty=\ell$, we find
\[
 r(\tau)=1+\int_0^\tau\Bigl[-\ell r(v)+\beta^2p(p-1)\int_0^vc(v-u)^{p-2}r(v-u)r(u)\,\dd u\Bigr]\dd v.
\]
The integrand is continuous in $v$, which implies that $r$ is continuously differentiable on $[0,\infty)$ and satisfies \eqref{eq:stationary-cluster-response}, with $r'(0)=-\ell$. This completes the proof.
\end{proof}

\Cref{prop:stationary-clusters} describes every translated cluster, but it does not determine the plateau $q$ or the multiplier $\ell$, it does not show that the whole translated family converges, and it says nothing about the correlation with a fixed early state. We answer these questions in \cref{sec:stationary-selection}, using the scalar equation of \cref{sec:scalar} and the causal covariance representation of \cref{sec:geometry-core}.


\part{Plateau selection, forgetting, and correlation aging}\label{part:selection}
\chapter{The scalar equation of a stationary cluster}\label{sec:scalar}

The CK response equation \eqref{eq:ck-response} is causal, in the sense that its memory integral runs only over times between the earlier and the later observation time. When both times are translated to late times and their lag stays bounded, the equation involves only the translated kernels on a bounded window, and it passes to every translated cluster. This is how we obtained the closed equation \eqref{eq:stationary-cluster-response} in \cref{prop:stationary-clusters}. The correlation equation \eqref{eq:ck-correlation} is different. Its memory integrals run over the whole past, and the stationary limit does not control their older part. In this section, we extract from the closed response equation the information that it carries about a stationary cluster $(c,r,\ell)$.

The main result, \cref{prop:scalar-clusters}, is a scalar memory equation for $c$. It determines the multiplier $\ell$ in terms of the plateau $q$, and it gives a strict sign condition on the interval $(q,1)$, which implies the inequality $\beta^2p(p-1)q^{p-2}(1-q)^2\le1$. We also show that the response equation with a prescribed multiplier has at most one solution (\cref{lem:stationary-uniqueness}). We use only \cref{prop:stationary-clusters}. In \cref{sec:stationary-selection}, the multiplier formula and the sign condition enter the proofs of marginality (\cref{sel:positive-marginality}) and of the selection of the plateau (\cref{sel:selected-sector}), and the uniqueness lemma identifies the stationary profile in \cref{sel:all-temperature}. The scalar equation \eqref{eq:scalar-memory} is also the starting point of \cref{sec:stationary-relaxation}.

The scalar memory equation \eqref{eq:scalar-memory} is the equation classified by Dembo, Guionnet, and Mazza \cite[arXiv version~1, Proposition~1.4, p.~5, and Section~5, pp.~17--20]{dgm2007}. For $b>0$ and a nondecreasing function $\phi\colon[0,1]\to\R$ with Lipschitz continuous derivative and $\sup_{0\le y\le1}\phi(y)(1-y)\ge b$, they characterized the plateau of its solution through the largest admissible root and proved the corresponding stability inequality. The sign condition \eqref{eq:scalar-sign} and the inequality \eqref{eq:cluster-submarginality} already follow from their scalar classification; here we obtain them by a direct comparison of the memory term. These hypotheses hold for the polynomial $\phi$ of \eqref{eq:scalar-memory}, with $b=1$ by \eqref{eq:plateau-multiplier}, but we do not use their result. We derive the equation and its consequences directly from \cref{prop:stationary-clusters}, for two reasons. Our response is normalized differently, and \cref{sec:stationary-selection} needs the constants and the strict inequality in the form given here. We also note that their reduction of the CK solution to the scalar equation at all temperatures \cite[arXiv version~1, Proposition~6.1, pp.~21--22]{dgm2007} assumes the asymptotic conditions (6.5) and (6.6) of that paper. Here, \cref{prop:stationary-clusters} applies to every stationary cluster without such assumptions.

Throughout this section, $p\ge3$ and $\beta>0$ are fixed. Derivatives of stationary profiles at $t=0$ are right derivatives on $[0,\infty)$.

\begin{proposition}[Scalar constraints]\label{prop:scalar-clusters}
Let $(c,r,\ell)$ be a stationary cluster with plateau $q$, as in \cref{prop:stationary-clusters}, and let $B=\beta^2p$ and $\gamma=\ell-B$. Then $r=-c'$ is strictly positive on $[0,\infty)$, and
\begin{equation}
 c'(t)=-1+\int_0^t\phi(c(t-u))r(u)\,\dd u,
 \qquad \phi(y)=\gamma+By^{p-1}.
 \label{eq:scalar-memory}
\end{equation}
The plateau $q=\lim_{t\to\infty}c(t)$ belongs to $[0,1)$ and satisfies
\begin{equation}
 \phi(q)(1-q)=1,\qquad
 \ell=L(q),\qquad\text{where }L(q)=\frac1{1-q}+\beta^2p(1-q^{p-1}).
 \label{eq:plateau-multiplier}
\end{equation}
Moreover, if $F(y)=\phi(y)(1-y)-1$, then
\begin{equation}
 F(q)=0,\qquad F(y)<0\quad(q<y<1).
 \label{eq:scalar-sign}
\end{equation}
In particular,
\begin{equation}
 \beta^2p(p-1)q^{p-2}(1-q)^2\le1.
 \label{eq:cluster-submarginality}
\end{equation}
\end{proposition}

\begin{proof}
We first show that $r$ is positive and integrate the response equation once, which gives \eqref{eq:scalar-memory}. We then let $t\to\infty$ in \eqref{eq:scalar-memory} to identify the multiplier. Finally, we compare the memory integral in \eqref{eq:scalar-memory} with its value at the current correlation, which yields the sign condition.

\paragraph{Step 1: positivity and the scalar equation.}
With $\beta^2p(p-1)=B(p-1)$, the stationary response equation \eqref{eq:stationary-cluster-response} reads
\[
 r'(t)=-\ell r(t)+B(p-1)\int_0^tc(t-u)^{p-2}r(t-u)r(u)\,\dd u.
\]
By variation of constants and $r(0)=1$, we have
\[
 r(t)=e^{-\ell t}+B(p-1)\int_0^te^{-\ell(t-v)}\int_0^vc(v-u)^{p-2}r(v-u)r(u)\,\dd u\,\dd v.
\]
Every factor in the double integral is nonnegative, because $c\ge0$ and $r\ge0$. This implies that $r(t)\ge e^{-\ell t}>0$. Then $c$ is strictly decreasing, and $c(t)>q$ for every $t\ge0$.

Next, we integrate the response equation from $0$ to $t$. Using $r(0)=1$ and $\int_0^tr(u)\,\dd u=1-c(t)$, we have
\[
 r(t)-1=-\ell\bigl(1-c(t)\bigr)+B(p-1)\int_0^t\int_0^vc(v-u)^{p-2}r(v-u)r(u)\,\dd u\,\dd v.
\]
The double integral is over the bounded triangle $\{0\le u\le v\le t\}$, where the integrand is continuous; Fubini's theorem applies there without any assumption on the behavior of $c$ and $r$ at infinity. Exchanging the order of integration and substituting $w=v-u$, we find
\[
\begin{aligned}
 &(p-1)\int_0^t\int_0^vc(v-u)^{p-2}r(v-u)r(u)\,\dd u\,\dd v\\
 &\qquad=\int_0^tr(u)\Bigl[(p-1)\int_0^{t-u}c(w)^{p-2}r(w)\,\dd w\Bigr]\dd u.
\end{aligned}
\]
As $c'=-r$ and $c(0)=1$, the derivative of $c(w)^{p-1}$ is $-(p-1)c(w)^{p-2}r(w)$, and the bracket equals $1-c(t-u)^{p-1}$. Inserting this into the previous identity, we obtain
\[
 r(t)=1-(\ell-B)\bigl(1-c(t)\bigr)-B\int_0^tc(t-u)^{p-1}r(u)\,\dd u.
\]
Finally, we substitute $r(t)=-c'(t)$ on the left and $1-c(t)=\int_0^tr(u)\,\dd u$ on the right. Recalling that $\gamma=\ell-B$, we arrive at
\[
 c'(t)=-1+\int_0^t\bigl(\gamma+Bc(t-u)^{p-1}\bigr)r(u)\,\dd u,
\]
which is \eqref{eq:scalar-memory}.

\paragraph{Step 2: the multiplier.}
By \cref{prop:stationary-clusters}, the limit $q$ exists, $q<1$, and $\int_0^\infty r(u)\,\dd u=1-q$. We claim that the integral in \eqref{eq:scalar-memory} converges to $\phi(q)(1-q)$ as $t\to\infty$. Let $M_\phi=\sup_{[0,1]}|\phi|$, and fix $U>0$. For $t>U$, we split the integral at $U$ and find
\[
\begin{aligned}
 &\Bigl|\int_0^t\phi(c(t-u))r(u)\,\dd u-\phi(q)(1-q)\Bigr|\\
 &\qquad\le\sup_{0\le u\le U}|\phi(c(t-u))-\phi(q)|\int_0^Ur(u)\,\dd u
 +2M_\phi\int_U^\infty r(u)\,\dd u.
\end{aligned}
\]
Here we used $0\le c\le1$ and $\int_0^\infty r=1-q$. For $0\le u\le U$, we have $q\le c(t-u)\le c(t-U)$. Then $c(t-u)\to q$ uniformly in $u\in[0,U]$ as $t\to\infty$, and the first term tends to zero because $\phi$ is continuous. Letting $t\to\infty$ and then $U\to\infty$, we obtain the claim, since $r$ is integrable. By \eqref{eq:scalar-memory}, we have $c'(t)\to\phi(q)(1-q)-1$. If this limit were nonzero, the bounded function $c$ would grow or decrease linearly at infinity, which is impossible. We conclude that $\phi(q)(1-q)=1$. Since $\phi(q)=\ell-B+Bq^{p-1}$, this identity is equivalent to $\ell=L(q)$, which proves \eqref{eq:plateau-multiplier}.

\paragraph{Step 3: the sign condition.}
The function $\phi$ is nondecreasing on $[0,1]$, because $B>0$. As $c$ is nonincreasing, we have $c(t-u)\ge c(t)$ for $0\le u\le t$. Using this and $r\ge0$ in \eqref{eq:scalar-memory}, we find
\[
 c'(t)\ge-1+\phi(c(t))\int_0^tr(u)\,\dd u=-1+\phi(c(t))\bigl(1-c(t)\bigr)=F(c(t)).
\]
By Step~1, $c$ decreases continuously and strictly from $c(0)=1$ to its limit $q$. Every $y\in(q,1)$ is then equal to $c(t_y)$ for some $t_y\in(0,\infty)$, and
\[
 F(y)\le c'(t_y)=-r(t_y)<0.
\]
Together with $F(q)=\phi(q)(1-q)-1=0$, which is \eqref{eq:plateau-multiplier}, this proves \eqref{eq:scalar-sign}.

Since the polynomial $F$ vanishes at $q$ and is negative on $(q,1)$, we have $F'(q)\le0$. Using $\phi(q)=1/(1-q)$, we compute
\[
 F'(q)=B(p-1)q^{p-2}(1-q)-\phi(q)=\frac{\beta^2p(p-1)q^{p-2}(1-q)^2-1}{1-q}.
\]
As $1-q>0$, the inequality $F'(q)\le0$ is \eqref{eq:cluster-submarginality}.
\end{proof}

\Cref{prop:scalar-clusters} gives constraints that every stationary cluster satisfies; it does not determine the plateau. If $q=0$, then \eqref{eq:plateau-multiplier} gives $\ell=L(0)=1+\beta^2p$, and \eqref{eq:cluster-submarginality} carries no information. If $q>0$, the multiplier is determined by $q$, and \eqref{eq:cluster-submarginality} is a one-sided version of the marginality relation $\beta^2p(p-1)q^{p-2}(1-q)^2=1$ of \cref{thm:main-aging}. To select $q$, \cref{sec:stationary-selection} combines these constraints with the upper bound on the energy of prefix centers from \cref{prop:threshold-inputs}.

We now turn to uniqueness. We write $f*g$ for the convolution on $[0,\infty)$, given by $(f*g)(\tau)=\int_0^\tau f(\tau-u)g(u)\,\dd u$. A solution of the system in the next lemma is understood in integrated form: a pair $(c,r)$ of locally bounded Borel measurable functions on $[0,\infty)$ such that, for every $t\ge0$,
\[
 c(t)=1-\int_0^tr(u)\,\dd u,
 \qquad
 r(t)=1+\int_0^t\Bigl[-\ell r(v)+\beta^2p(p-1)\bigl((c^{p-2}r)*r\bigr)(v)\Bigr]\dd v.
\]
The integrated equations are pointwise identities, not merely identities almost everywhere. As the integrands on their right sides are locally integrable, both functions are continuous. Every pair of continuously differentiable functions satisfying the system, with right derivatives at zero, is a solution in this sense; by \cref{prop:stationary-clusters}, this applies to every stationary cluster with multiplier $\ell$.

\begin{lemma}[Uniqueness with prescribed multiplier]\label{lem:stationary-uniqueness}
For a fixed finite $\ell$, the system
\[
 c'=-r,\qquad
 r'=-\ell r+\beta^2p(p-1)\bigl(c^{p-2}r\bigr)*r,
 \qquad c(0)=r(0)=1,
\]
has at most one solution in this locally bounded integrated sense on $[0,\infty)$.
\end{lemma}

\begin{proof}
Let $(c_1,r_1)$ and $(c_2,r_2)$ be two solutions, and fix $T>0$. Let $M\ge1$ be a bound for $|c_1|$, $|c_2|$, $|r_1|$, and $|r_2|$ on $[0,T]$, and define, for $0\le t\le T$,
\[
 D(t)=\sup_{0\le v\le t}\bigl(|c_1(v)-c_2(v)|+|r_1(v)-r_2(v)|\bigr).
\]
We first bound the difference of the memory integrands. Let $0\le u\le v\le T$, and write $a=v-u$ and $b=u$. Telescoping the three factors, we have
\[
\begin{aligned}
 &\bigl|c_1(a)^{p-2}r_1(a)r_1(b)-c_2(a)^{p-2}r_2(a)r_2(b)\bigr|\\
 &\qquad\le|c_1(a)^{p-2}-c_2(a)^{p-2}|\,|r_1(a)r_1(b)|
 +|c_2(a)|^{p-2}|r_1(a)-r_2(a)|\,|r_1(b)|\\
 &\qquad\quad+|c_2(a)^{p-2}r_2(a)|\,|r_1(b)-r_2(b)|\\
 &\qquad\le pM^{p-1}\{D(a)+D(b)\}.
\end{aligned}
\]
In the last line, we used $|x^{p-2}-y^{p-2}|\le(p-2)M^{p-3}|x-y|$ for $|x|,|y|\le M$, and $|c_1-c_2|,|r_1-r_2|\le D$. Since $D\ge0$, we have $\int_0^v\{D(v-u)+D(u)\}\,\dd u=2\int_0^vD(u)\,\dd u\le2\int_0^tD(u)\,\dd u$ for $v\le t$. Subtracting the integrated equations of the two solutions, we obtain, for $0\le t\le T$,
\[
\begin{aligned}
 |c_1(t)-c_2(t)|&\le\int_0^tD(u)\,\dd u,\\
 |r_1(t)-r_2(t)|&\le\bigl(|\ell|+2\beta^2p^2(p-1)M^{p-1}T\bigr)\int_0^tD(u)\,\dd u.
\end{aligned}
\]
The right sides are nondecreasing in $t$. Taking the supremum over $[0,t]$, we find
\[
 D(t)\le K_T\int_0^tD(u)\,\dd u,
 \qquad K_T=1+|\ell|+2\beta^2p^2(p-1)M^{p-1}T,
\]
where $K_T$ depends on $\ell$, $p$, $\beta$, $T$, and the bound $M$. The function $D$ is bounded and nondecreasing on $[0,T]$, and Gronwall's inequality implies that $D=0$ on $[0,T]$. As $T$ was arbitrary, the two solutions coincide.
\end{proof}

\chapter{The causal tree realization and reflection}\label{sec:geometry-core}

In this section, we realize the CK correlation as the Gram kernel of a curve $v_t$ of unit vectors in a deterministic Hilbert space. We use this realization to study the part of a late state that was already determined at an earlier time. The main results are the exact causal realization, \cref{prop:geometry-tree-realization}; the evolution of a projected state, \cref{lem:geometry-projected-drift}; the reflection identity for stationary clusters, \cref{sel:prefix-reflection}; and the spectral representation of every stationary correlation profile, \cref{sel:cluster-complete-monotonicity}.

We use the CK equations of \cref{sec:model}, the physical bounds \eqref{eq:ck-physical-bounds}, the Green covariance identity of \cref{prop:green-covariance}, the energy identity \eqref{eq:energy-identity}, the uniform multiplier bound \eqref{eq:uniform-multiplier}, the diagonal increment bound \eqref{eq:diagonal-increment}, the regularity of the response from \cref{prop:backward-response,prop:response-compactness}, and the stationary clusters of \cref{prop:stationary-clusters}. The results of this section are used in most of the later sections. The limits $d_L$ and $h_L$ of \cref{sel:prefix-reflection} are the objects of \cref{prop:threshold-inputs}, and they enter the proofs of marginality and forgetting in \cref{sec:stationary-selection}. The ceilings of \cref{sec:tree-ceilings} and the geometry of \cref{sec:geometry-results,sec:three-time} rest on the tree space and its height grading. In \cref{resp:chapter}, we use the leaf component of a projected state and \cref{lem:geometry-projected-drift}. The complete monotonicity asserted in \cref{thm:main-aging} is part of \cref{sel:cluster-complete-monotonicity}, which is also used in \cref{resp:chapter,sec:three-time,sec:stationary-relaxation}.

We start from the Green identity \eqref{eq:green-covariance}, which we derived in \cref{prop:green-covariance} by variation of constants from the CK equations of \cite{bag2006,dgm2007}. Positive noncrossing expansions of the associated response equation occur in the work of Guionnet and Mazza \cite[Section~2.1 and Lemma~2.1]{guionnetmazza2005} and of Dembo, Guionnet, and Mazza \cite[Section~2]{dgm2007}. Guionnet and Mazza also describe their full Fock-space interpretation. We build on these positive expansions. The tree realization below is a positive expansion of the covariance, graded by ordered trees so that every term records its height and the times of the thermal noise that it uses; we do not identify this grading with a Wiener chaos decomposition. All constructions take place in one deterministic Hilbert space. The prefix projection $P_T$ defined below plays the role of conditioning on the noise up to time $T$, but we never condition on the finite-dimensional probability space.

\paragraph{Standing assumptions and notation.}
Throughout this section, $p\ge3$ and $\beta>0$ are fixed, $m=p-1$, and $(C,R,\lambda)$ is the uniformly initialized hard CK solution of \cref{sec:model}. All Hilbert spaces and tensor products are real. We write
\[
 \Sigma(t,u)=\beta^2p(p-1)C(t,u)^{p-2}R(t,u),\qquad 0\le u\le t,
\]
for the kernel \eqref{eq:self-energy-kernel}, so that the response equation \eqref{eq:ck-response} reads
\begin{equation}\label{eq:tree-response-equation}
 \partial_tR(t,s)=-\lambda(t)R(t,s)+\int_s^t\Sigma(t,u)R(u,s)\,\dd u,\qquad t>s.
\end{equation}
By \eqref{eq:ck-physical-bounds} and \cref{prop:response-compactness}, we have $0\le R(t,s)\le C(t,s)\le1$ for $t\ge s$, and in particular $\Sigma\ge0$.

\section{The tree Hilbert space}\label{sec:tree-space}

We first explain the purpose of the construction. Recall from \cref{prop:green-covariance} that the correlation is the sum of three terms,
\[
\begin{aligned}
 C(t,s)={}&R(t,0)R(s,0)+2\int_0^{t\wedge s}R(t,u)R(s,u)\,\dd u\\
 &+\beta^2p\int_0^t\int_0^sR(t,u)R(s,w)C(u,w)^{m}\,\dd w\,\dd u.
\end{aligned}
\]
The first term comes from the initial condition, the second from the thermal noise, and the third from the colored memory. The first two terms are the inner product of the vectors $R(t,0)e_0+\sqrt2\,\mathbf 1_{[0,t]}R(t,\cdot)$ and $R(s,0)e_0+\sqrt2\,\mathbf 1_{[0,s]}R(s,\cdot)$ in $\R e_0\oplus L^2(\R_+)$. The second component of such a vector records the noise received at each time $u\le t$, transported to time $t$ by the response. If $C(u,w)=\ip{v_u}{v_w}$ for some vectors $v_u$, then $C(u,w)^m=\ip{v_u^{\otimes m}}{v_w^{\otimes m}}$, and the third term is also an inner product. This suggests looking for vectors $v_t$ with
\[
 v_t=R(t,0)e_0+\sqrt2\,\mathbf 1_{[0,t]}R(t,\cdot)+\beta\sqrt p\int_0^tR(t,u)\mathcal I\bigl(v_u^{\otimes m}\bigr)\,\dd u,
\]
where $\mathcal I$ is an isometry from the $m$-fold tensor power into a space orthogonal to the first two terms. Iterating this equation, we expand $v_t$ as a sum of terms indexed by trees. Each leaf of a tree carries either the initial vector $e_0$ or the time of a noise increment, and each internal vertex records one use of the memory integral. The point of this bookkeeping is that every term of $v_t$ comes with the times of the noise it uses. The terms whose noise times are at most $T$ form the part of the state at time $t$ built from the initial vector and noise received by time $T$; their memory integrals may still extend beyond $T$. We now make this precise.

\paragraph{Trees and grafting.}
Let $\mathcal E=\R e_0\oplus L^2(\R_+)$, where $e_0$ is a unit vector. By a \emph{tree} we mean a finite rooted tree in which every vertex has either no children or exactly $m$ children, and the children of each vertex are ordered. We call vertices without children leaves. We write $|\tau|$ for the number of leaves of a tree $\tau$, whose leaves we list from left to right, and $\bullet$ for the tree with one vertex. For trees $\tau_1,\dots,\tau_m$, we write $[\tau_1,\dots,\tau_m]$ for the tree whose root has the subtrees $\tau_1,\dots,\tau_m$, in this order. We call this operation grafting. Every tree other than $\bullet$ is the graft of a unique ordered $m$-tuple of trees.

We define $\mathcal H$ as the orthogonal direct sum, over all trees $\tau$, of one copy of $\mathcal E^{\otimes|\tau|}$:
\[
 \mathcal H=\bigoplus_\tau\mathcal E^{\otimes|\tau|}.
\]
The summand indexed by $\bullet$ is $\mathcal E$. We call it the one-leaf summand. The space $\mathcal H$ is separable, because there are countably many trees and each summand is separable. Since tensor products distribute over orthogonal direct sums, we have
\[
 \mathcal H^{\otimes m}\cong\bigoplus_{(\tau_1,\dots,\tau_m)}\mathcal E^{\otimes(|\tau_1|+\dots+|\tau_m|)}.
\]
We define the grafting map $\mathcal I\colon\mathcal H^{\otimes m}\to\mathcal H$ by sending the summand indexed by $(\tau_1,\dots,\tau_m)$ identically onto the summand indexed by $[\tau_1,\dots,\tau_m]$; the leaves of $\tau_1$ come first, then those of $\tau_2$, and so on. Since every tree other than $\bullet$ is the graft of a unique $m$-tuple, $\mathcal I$ is an isometry, and its range is the orthogonal complement of the one-leaf summand. Its adjoint satisfies $\mathcal I^*\mathcal I=\Id$, and $\mathcal I^*$ vanishes on the one-leaf summand. We use ordered trees because they make grafting an isometry, with no symmetry factors. When $p=3$, we identify a tensor $x\otimes y\in\mathcal H\otimes\mathcal H$ with the Hilbert--Schmidt operator $z\mapsto\ip yz\,x$ on $\mathcal H$; this is how we use $\mathcal I^*$ in \cref{sec:geometry-results,sec:three-time}.

\paragraph{Leaf labels and positivity.}
We now realize $\mathcal E$ as an $L^2$ space. Let $\mathbb L=\{\star\}\sqcup[0,\infty)$ carry the measure $\delta_\star+\dd u$. Then $\mathcal E=L^2(\mathbb L)$, with $e_0=\mathbf 1_{\{\star\}}$, and $\mathcal E^{\otimes k}=L^2(\mathbb L^k)$. An element of the summand indexed by $\tau$ is then a function of the leaf labels of $\tau$. The label $\star$ marks the initial vector, and a label $u\ge0$ marks a noise time. We say that a vector of $\mathcal H$ is nonnegative if each of its components is a nonnegative function almost everywhere. Tensor products and grafts of nonnegative vectors are nonnegative.

\paragraph{Heights.}
We define the height of $\bullet$ to be zero, and the height of $[\tau_1,\dots,\tau_m]$ to be one plus the largest height of $\tau_1,\dots,\tau_m$. We write $\Pi_j$ for the orthogonal projection of $\mathcal H$ onto the sum of the summands indexed by trees of height $j$, and we set $\Pi_{\le H}=\sum_{j=0}^H\Pi_j$. Note that $\sum_{j\ge0}\Pi_j=\Id$ in the strong sense, because every tree has a finite height.

\paragraph{Prefix projections.}
For $T\ge0$, let $\mathcal E_T=\R e_0\oplus L^2([0,T])$, and let $\varpi_T$ be the orthogonal projection of $\mathcal E$ onto $\mathcal E_T$, that is, $\varpi_T(ae_0+g)=ae_0+\mathbf 1_{[0,T]}g$. We define the prefix projection
\[
 P_T=\bigoplus_\tau\varpi_T^{\otimes|\tau|},
\]
and we write $\mathcal H_T$ for its range, the closed subspace generated by trees whose leaf factors lie in $\mathcal E_T$. In terms of leaf labels, $P_T$ multiplies the component indexed by $\tau$ by the indicator function that all leaf labels lie in $\{\star\}\cup[0,T]$. Since $P_T$, $\Pi_j$, and $\mathcal I$ act summand by summand, and the leaves of a graft are the leaves of its root subtrees, we have
\begin{equation}
 P_TP_S=P_{T\wedge S},\qquad P_T\Pi_j=\Pi_jP_T,\qquad
 P_T\mathcal I=\mathcal I P_T^{\otimes m}.
 \label{eq:geometry-prefix-commutation}
\end{equation}
In words, $P_T$ is a causal projection. A vector built from leaves with noise times at most $T$ belongs to $\mathcal H_T$, and $P_T$ commutes with grafting. The projections $P_T$, $\Id-P_T$, and $\Pi_j$ preserve nonnegativity, since each of them multiplies every component by an indicator function.

\section{The exact causal realization}\label{sec:tree-realization}

\begin{proposition}[Exact causal realization]
\label{prop:geometry-tree-realization}
There exists a norm-continuous curve $t\mapsto v_t\in\mathcal H$, $t\ge0$, with $v_t\in\mathcal H_t$ for every $t$, such that
\begin{align}
 v_t={}&R(t,0)e_0+\sqrt2\,\mathbf 1_{[0,t]}R(t,\cdot)
       +\beta\sqrt p\int_0^tR(t,u)\mathcal I\bigl(v_u^{\otimes m}\bigr)\,\dd u,
       \label{eq:geometry-tree-realization}\\
 \ip{v_t}{v_s}={}&C(t,s),\qquad \norm{v_t}=1,
       \label{eq:geometry-tree-gram}
\end{align}
where the integral is a Bochner integral in $\mathcal H$. For every tree $\tau$, the component of $v_t$ indexed by $\tau$ is a nonnegative function of the leaf labels; in particular, $\ip{\Pi_jv_t}{\Pi_jv_s}\ge0$ for all $j\ge0$ and $t,s\ge0$. Let $h$ be the tree energy, defined by
\begin{equation}
 h(v)=\sqrt p\,\ip{v}{\mathcal I\bigl(v^{\otimes m}\bigr)},\qquad v\in\mathcal H.
 \label{eq:geometry-tree-energy}
\end{equation}
Then $h(v_t)=e_\beta(t)$ and $\lambda(t)=1+p\beta h(v_t)$ for every $t\ge0$. Moreover, for all $T\ge0$ and $t\ge0$,
\begin{equation}
 \mathcal I^*P_Tv_t
 =\beta\sqrt p\int_0^tR(t,u)\bigl(P_Tv_u\bigr)^{\otimes m}\,\dd u.
 \label{eq:geometry-prefix-root-frame}
\end{equation}
\end{proposition}

Before the proof, we write out the first terms of the expansion in the cubic case.

\begin{example}[The first two heights for $p=3$]\label{ex:tree-first-heights}
Let $p=3$, so that grafting is binary. We write $\eta_t$ for the component of $v_t$ at height zero, which is the leaf vector
\[
 \eta_t=R(t,0)e_0+\sqrt2\,\mathbf 1_{[0,t]}R(t,\cdot).
\]
The only tree of height one is $[\bullet,\bullet]$, with two leaves. Iterating \eqref{eq:geometry-tree-realization} once, we see that the component of $v_t$ at height one is
\[
 w_t=\beta\sqrt3\int_0^tR(t,u)\,\mathcal I(\eta_u\otimes\eta_u)\,\dd u.
\]
As a function of two noise times $y_1,y_2\in[0,t]$, it is
\[
 w_t(y_1,y_2)=2\sqrt3\,\beta\int_{y_1\vee y_2}^tR(t,u)R(u,y_1)R(u,y_2)\,\dd u.
\]
Noise received at the times $y_1$ and $y_2$ is transported to a time $u$ by the response, combined there by the cubic interaction, and transported to time $t$. At labels involving $\star$, a factor $\sqrt2R(u,y_i)$ is replaced by $R(u,0)$. The trees of height two are $[[\bullet,\bullet],\bullet]$, $[\bullet,[\bullet,\bullet]]$, and $[[\bullet,\bullet],[\bullet,\bullet]]$, and the corresponding components of $v_t$ are
\begin{gather*}
 \beta\sqrt3\int_0^tR(t,u)\,\mathcal I(w_u\otimes\eta_u)\,\dd u,\qquad
 \beta\sqrt3\int_0^tR(t,u)\,\mathcal I(\eta_u\otimes w_u)\,\dd u,\\
 \beta\sqrt3\int_0^tR(t,u)\,\mathcal I(w_u\otimes w_u)\,\dd u.
\end{gather*}
The prefix projection acts only on the labels. For instance, we see that $P_Tw_t$ keeps the values of $w_t$ at noise times $y_1,y_2\le T$, while the time integral defining it still runs up to $t$. Finally, the Gram kernel of $\eta_t+w_t$ is
\[
 A_0(t,s)+3\beta^2\int_0^t\int_0^sR(t,u)R(s,w)A_0(u,w)^2\,\dd w\,\dd u,
 \qquad A_0(t,s)=\ip{\eta_t}{\eta_s},
\]
which is the first step of the Picard iteration of the Green identity \eqref{eq:green-covariance}. In the proof below, we show that these iterations increase to $C$.
\end{example}

\begin{proof}[Proof of \cref{prop:geometry-tree-realization}]
The proof has five steps. We first construct the component of $v_t$ indexed by each tree through a finite recursion. We then show that the Gram kernels of the truncations at bounded height increase to $C$, using a uniqueness property of the Green equation. This allows us to sum the components. Finally, we compute the projected root frame and the energy.

\paragraph{Step 1: the tree components.}
Let $\eta_t=R(t,0)e_0+\sqrt2\,\mathbf 1_{[0,t]}R(t,\cdot)\in\mathcal E$, and set $w_\bullet(t)=\eta_t$. If $\tau=[\tau_1,\dots,\tau_m]$, we define, in the summand indexed by $\tau$,
\begin{equation}\label{eq:tree-component-recursion}
 w_\tau(t)=\beta\sqrt p\int_0^tR(t,u)\,w_{\tau_1}(u)\otimes\dots\otimes w_{\tau_m}(u)\,\dd u.
\end{equation}
This is a recursion on the number of vertices. The curve $t\mapsto\eta_t$ is norm-continuous, because $R$ is uniformly continuous on compact causal triangles and bounded by one. By induction on the number of vertices, each $w_\tau$ is a norm-continuous curve, and the integral in \eqref{eq:tree-component-recursion} is a Bochner integral of a continuous integrand. Also by induction, $w_\tau(t)$ is nonnegative, because $R\ge0$ and tensor products and integrals of nonnegative vectors are nonnegative, and $w_\tau(t)$ lies in $\mathcal E_t^{\otimes|\tau|}$, because $w_{\tau_i}(u)\in\mathcal E_u^{\otimes|\tau_i|}\subset\mathcal E_t^{\otimes|\tau_i|}$ for $u\le t$.

\paragraph{Step 2: the truncated Gram kernels.}
For $n\ge0$, let $v^{(n)}_t$ be the sum of $w_\tau(t)$ over the finitely many trees $\tau$ of height at most $n$, and let $K_n(t,s)=\ip{v^{(n)}_t}{v^{(n)}_s}$. A tree of height at most $n\ge1$ is either $\bullet$ or the graft of $m$ trees of height at most $n-1$. Summing \eqref{eq:tree-component-recursion} over these trees, we have
\[
 v^{(n)}_t=\eta_t+\beta\sqrt p\int_0^tR(t,u)\,\mathcal I\bigl((v^{(n-1)}_u)^{\otimes m}\bigr)\,\dd u.
\]
Recall that $\eta_t$ is orthogonal to the range of $\mathcal I$ and that $\mathcal I$ is an isometry. Using also $\ip{x^{\otimes m}}{y^{\otimes m}}=\ip xy^m$, we find that $K_0=A_0$ and
\begin{equation}\label{eq:tree-truncated-green}
 K_n(t,s)=A_0(t,s)+\beta^2p\int_0^t\int_0^sR(t,u)R(s,w)K_{n-1}(u,w)^m\,\dd w\,\dd u,\qquad n\ge1,
\end{equation}
where
\[
 A_0(t,s)=\ip{\eta_t}{\eta_s}=R(t,0)R(s,0)+2\int_0^{t\wedge s}R(t,u)R(s,u)\,\dd u.
\]
The squared coefficients are $(\sqrt2)^2=2$ and $(\beta\sqrt p)^2=\beta^2p$, which are the coefficients of the Green identity \eqref{eq:green-covariance}. Note that $K_n\ge0$, because the components are nonnegative. We claim that
\begin{equation}\label{eq:tree-truncated-monotone}
 0\le K_{n-1}\le K_n\le C,\qquad n\ge1.
\end{equation}
Indeed, the Green identity reads $C=A_0+\beta^2p\int\!\!\int RR\,C^m$, and its double integral is nonnegative by \eqref{eq:ck-physical-bounds}. This implies that $K_0=A_0\le C$. If $0\le K_{n-1}\le C$, then \eqref{eq:tree-truncated-green} and $R\ge0$ imply that $K_n\le C$. In the same way, $K_1\ge K_0$, and if $K_{n-1}\ge K_{n-2}\ge0$, then $K_n\ge K_{n-1}$. This proves \eqref{eq:tree-truncated-monotone}. Then $K_n$ increases pointwise to a limit $K$ with $0\le K\le C\le1$, and by monotone convergence in \eqref{eq:tree-truncated-green}, $K$ solves the Green equation
\begin{equation}\label{eq:tree-green-equation}
 K(t,s)=A_0(t,s)+\beta^2p\int_0^t\int_0^sR(t,u)R(s,w)K(u,w)^m\,\dd w\,\dd u.
\end{equation}

\paragraph{Step 3: uniqueness for the Green equation.}
We prove that, for every $A<\infty$, \eqref{eq:tree-green-equation} has at most one solution on $[0,A]^2$ with values in $[-1,1]$. The kernel $C$ is such a solution by \cref{prop:green-covariance}, so uniqueness implies that $K=C$, that is, $K_n$ increases to $C$. Let $K$ and $K'$ be two such solutions and let $\Delta=K-K'$. Using $|x^m-y^m|\le m|x-y|$ for $x,y\in[-1,1]$ and $0\le R\le1$, we have, for $t,s\le A$,
\[
 |\Delta(t,s)|\le\beta^2pm\int_0^t\int_0^s|\Delta(u,w)|\,\dd w\,\dd u.
\]
Let $\rho(a)=\sup_{t,s\le a}|\Delta(t,s)|$, a bounded nondecreasing function. We bound $|\Delta(u,w)|$ by $\rho(\max\{u,w\})$ and divide the square $[0,a]^2$ according to which coordinate is larger. For $a\le A$, this yields
\[
 \rho(a)\le\beta^2pm\int_0^a\int_0^a\rho(\max\{u,w\})\,\dd w\,\dd u
 =2\beta^2pm\int_0^az\rho(z)\,\dd z
 \le2\beta^2pmA\int_0^a\rho(z)\,\dd z.
\]
By Gronwall's inequality, we have $\rho=0$ on $[0,A]$. As $A$ is arbitrary, we conclude that $K=C$.

\paragraph{Step 4: the curve $v_t$.}
Using Step 3, we find that $\sum_\tau\norm{w_\tau(t)}^2=\lim_nK_n(t,t)=C(t,t)=1$. The components lie in orthogonal summands, and the series $v_t=\sum_\tau w_\tau(t)$ therefore converges in $\mathcal H$, with $v^{(n)}_t\to v_t$. Then $\ip{v_t}{v_s}=\lim_nK_n(t,s)=C(t,s)$, which is \eqref{eq:geometry-tree-gram}. By the identity $\norm{v_t-v_s}^2=2\{1-C(t,s)\}$ and the continuity of $C$, the curve $t\mapsto v_t$ is norm-continuous. Each component of $v_t$ is nonnegative and lies in $\mathcal E_t^{\otimes|\tau|}$, which shows that $v_t\in\mathcal H_t$. The inner product $\ip{\Pi_jv_t}{\Pi_jv_s}$ is a sum of inner products of nonnegative functions, and it is nonnegative.

It remains to prove \eqref{eq:geometry-tree-realization}. The integrand $u\mapsto R(t,u)\mathcal I(v_u^{\otimes m})$ is continuous on $[0,t]$, and the integral is a Bochner integral. The orthogonal projection onto the summand indexed by $\tau=[\tau_1,\dots,\tau_m]$ is bounded and commutes with the integral, and it maps $\mathcal I(v_u^{\otimes m})$ to $w_{\tau_1}(u)\otimes\dots\otimes w_{\tau_m}(u)$. Comparing with \eqref{eq:tree-component-recursion}, we see that the component of the right side of \eqref{eq:geometry-tree-realization} indexed by $\tau$ is $w_\tau(t)$. Its one-leaf component is $\eta_t=w_\bullet(t)$. Both sides of \eqref{eq:geometry-tree-realization} therefore have the same components.

\paragraph{Step 5: the root frame and the energy.}
We apply $P_T$ to \eqref{eq:geometry-tree-realization}. The projection $P_T$ is bounded and commutes with grafting by \eqref{eq:geometry-prefix-commutation}, which gives
\[
 P_Tv_t=P_T\eta_t+\beta\sqrt p\int_0^tR(t,u)\,\mathcal I\bigl((P_Tv_u)^{\otimes m}\bigr)\,\dd u.
\]
The vector $P_T\eta_t$ lies in the one-leaf summand, which $\mathcal I^*$ annihilates, and $\mathcal I^*\mathcal I=\Id$. Applying $\mathcal I^*$ to the last display yields \eqref{eq:geometry-prefix-root-frame}. We now let $T=t$. For $u\le t$, we have $v_u\in\mathcal H_u\subset\mathcal H_t$, so $P_tv_u=v_u$. Taking the inner product of \eqref{eq:geometry-prefix-root-frame} with $v_t^{\otimes m}$, we obtain
\[
\begin{aligned}
 h(v_t)&=\sqrt p\,\ip{\mathcal I^*v_t}{v_t^{\otimes m}}
 =p\beta\int_0^tR(t,u)\ip{v_u}{v_t}^m\,\dd u\\
 &=p\beta\int_0^tC(t,u)^{p-1}R(t,u)\,\dd u
 =e_\beta(t),
\end{aligned}
\]
where the last equality is \eqref{eq:energy-identity}. By the same identity, we also have $\lambda(t)=1+p\beta e_\beta(t)=1+p\beta h(v_t)$. This completes the proof.
\end{proof}

\section{The evolution of a projected state}\label{sec:tree-projected-drift}

After time $T$, no new noise leaf survives the projection $P_T$, and the projected curve $P_Tv_t$ becomes differentiable. The next lemma computes its derivative. In \cref{resp:chapter}, we use it to turn reflection identities into bounds on derivatives at remote times (\cref{resp:forward-derivative}). The polynomial $h$ of \eqref{eq:geometry-tree-energy} is Fr\'echet differentiable on $\mathcal H$, and we write $\nabla h(v)$ for its gradient.

\begin{lemma}[Local polynomial drift after a causal projection]
\label{lem:geometry-projected-drift}
Define $\mathcal B\colon\mathcal H\to\mathcal H$ by
\[
 \mathcal B(v)=-v+\beta\nabla h(v)-p\beta h(v)v.
\]
For every fixed $T\ge0$, the curve $t\mapsto P_Tv_t$ is continuously differentiable on $(T,\infty)$, and
\begin{equation}
 \frac{\dd}{\dd t}P_Tv_t=P_T\mathcal B(v_t),\qquad t>T.
 \label{eq:geometry-projected-drift}
\end{equation}
The map $\mathcal B$ is Lipschitz on the closed unit ball of $\mathcal H$. If there exists $K<\infty$ such that $1-C(t,s)\le K|t-s|$ for all $t,s\ge0$, then there exists $K'<\infty$, depending only on $p$, $\beta$, and $K$, such that
\[
 \Bigl\|\frac{\dd}{\dd t}P_Tv_t-\frac{\dd}{\dd s}P_Tv_s\Bigr\|\le K'|t-s|^{1/2}
 \qquad\text{for all }T\ge0\text{ and }t,s>T.
\]
\end{lemma}

\begin{proof}
We differentiate the projected realization, identify the result with the gradient of $h$ along the curve, and prove the Lipschitz bounds. Throughout, we write
\[
 \Phi(u)=\beta\sqrt p\,\mathcal I\bigl(v_u^{\otimes m}\bigr),\qquad
 \Psi(t)=\int_0^tR(t,u)\Phi(u)\,\dd u,
\]
so that $v_t=\eta_t+\Psi(t)$ by \eqref{eq:geometry-tree-realization}, with $\eta_t$ the leaf vector of the proof of \cref{prop:geometry-tree-realization}. The curve $\Phi$ is continuous, with $\norm{\Phi(u)}\le\beta\sqrt p$.

\paragraph{Step 1: differentiation.}
By \cref{prop:backward-response}, the partial derivative $\partial_tR$ extends continuously to every compact causal triangle. For $t>T$, we have
\[
 P_T\eta_t=R(t,0)e_0+\sqrt2\,\mathbf 1_{[0,T]}R(t,\cdot),
\]
and the cutoff of the noise component no longer moves with $t$. On every compact interval of $(T,\infty)$, the difference quotients of $R(\cdot,u)$ converge to $\partial_tR(\cdot,u)$ uniformly in $u\in[0,T]$. It follows that $P_T\eta_t$ is continuously differentiable on $(T,\infty)$, with derivative $\partial_tR(t,0)e_0+\sqrt2\,\mathbf 1_{[0,T]}\partial_tR(t,\cdot)$. Inserting the response equation \eqref{eq:tree-response-equation}, and noting that the noise component of $P_T\eta_z$ is $\sqrt2R(z,y)$ for $y\le z\wedge T$ and zero for other $y$, we have
\begin{equation}\label{eq:tree-leaf-derivative}
 \frac{\dd}{\dd t}P_T\eta_t=-\lambda(t)P_T\eta_t+\int_0^t\Sigma(t,z)P_T\eta_z\,\dd z,\qquad t>T.
\end{equation}
Next, using that $R$ and $\partial_tR$ are continuous on compact causal triangles and that $R(t,t)=1$, we see that the curve $\Psi$ is continuously differentiable on $(0,\infty)$, with
\[
 \Psi'(t)=\Phi(t)+\int_0^t\partial_tR(t,u)\Phi(u)\,\dd u.
\]
We insert \eqref{eq:tree-response-equation} and exchange the order of integration on the triangle $\{0\le u\le z\le t\}$, which is allowed because the integrand is continuous and bounded in norm. We obtain
\begin{equation}\label{eq:tree-memory-derivative}
 \Psi'(t)=\Phi(t)-\lambda(t)\Psi(t)+\int_0^t\Sigma(t,z)\Psi(z)\,\dd z.
\end{equation}
Adding \eqref{eq:tree-leaf-derivative} and the image of \eqref{eq:tree-memory-derivative} under $P_T$, we have, for $t>T$,
\begin{equation}\label{eq:tree-projected-derivative}
 \frac{\dd}{\dd t}P_Tv_t=P_T\Bigl[-\lambda(t)v_t+\Phi(t)+\int_0^t\Sigma(t,z)v_z\,\dd z\Bigr].
\end{equation}

\paragraph{Step 2: the gradient along the curve.}
For $v,w\in\mathcal H$, the derivative of $h$ is
\[
 Dh(v)[w]=\sqrt p\,\ip{w}{\mathcal I(v^{\otimes m})}
 +\sqrt p\sum_{i=1}^m\ip{\mathcal I^*v}{v^{\otimes(i-1)}\otimes w\otimes v^{\otimes(m-i)}}.
\]
Recall from \eqref{eq:geometry-prefix-root-frame} with $T=t$ that $\mathcal I^*v_t=\beta\sqrt p\int_0^tR(t,u)v_u^{\otimes m}\,\dd u$. Each of the $m$ terms of the sum then equals $\beta\sqrt p\int_0^tR(t,u)C(t,u)^{m-1}\ip{v_u}w\,\dd u$, and we find that
\begin{equation}\label{eq:tree-gradient-along-curve}
 \beta\nabla h(v_t)=\Phi(t)+\beta^2p(p-1)\int_0^tC(t,u)^{p-2}R(t,u)v_u\,\dd u
 =\Phi(t)+\int_0^t\Sigma(t,u)v_u\,\dd u.
\end{equation}
Although ordered grafting is not symmetric, the $m$ contractions agree here because the root tensor of $v_t$ is an integral of symmetric tensor powers; this produces the factor $p-1$ in $\Sigma$. We also have $\lambda(t)=1+p\beta h(v_t)$ by \cref{prop:geometry-tree-realization}. Inserting this and \eqref{eq:tree-gradient-along-curve} into \eqref{eq:tree-projected-derivative}, we arrive at \eqref{eq:geometry-projected-drift}. Its right side is continuous in $t$, because $t\mapsto v_t$ and $\mathcal B$ are continuous, and this shows that $P_Tv_t$ is continuously differentiable on $(T,\infty)$.

\paragraph{Step 3: the Lipschitz bounds.}
The multilinear form $(y_0,\dots,y_m)\mapsto\ip{y_0}{\mathcal I(y_1\otimes\dots\otimes y_m)}$ has norm at most one, and we obtain $h$ and $Dh$ by inserting $v$ or $w$ into its $p=m+1$ slots. Changing one slot at a time, we find, for $v,v'$ in the closed unit ball,
\begin{gather*}
 |h(v)|\le\sqrt p,\qquad |h(v)-h(v')|\le p\sqrt p\,\norm{v-v'},\\
 \norm{\nabla h(v)-\nabla h(v')}\le p(p-1)\sqrt p\,\norm{v-v'}.
\end{gather*}
Writing $h(v)v-h(v')v'=h(v)(v-v')+\{h(v)-h(v')\}v'$, we have
\[
\begin{aligned}
 \norm{\mathcal B(v)-\mathcal B(v')}
 &\le\bigl\{1+\beta p(p-1)\sqrt p+p\beta(\sqrt p+p\sqrt p)\bigr\}\norm{v-v'}\\
 &=\bigl(1+2\beta p^2\sqrt p\bigr)\norm{v-v'}.
\end{aligned}
\]
Now suppose that $1-C(t,s)\le K|t-s|$ for all $t,s$. Then $\norm{v_t-v_s}^2=2\{1-C(t,s)\}\le2K|t-s|$. Since $P_T$ is a contraction, \eqref{eq:geometry-projected-drift} implies that, for $t,s>T$,
\[
 \Bigl\|\frac{\dd}{\dd t}P_Tv_t-\frac{\dd}{\dd s}P_Tv_s\Bigr\|
 \le\bigl(1+2\beta p^2\sqrt p\bigr)\sqrt{2K}\,|t-s|^{1/2}.
\]
This is the last assertion, with $K'=(1+2\beta p^2\sqrt p)\sqrt{2K}$.
\end{proof}

\begin{remark}[The diagonal increment bound]\label{rem:geometry-diagonal-increment}
The hypothesis of the last assertion of \cref{lem:geometry-projected-drift} holds with $K=\Lambda$, the constant of \eqref{eq:uniform-multiplier}. Indeed, by \eqref{eq:diagonal-increment} in \cref{prop:global-regularity} and the symmetry of $C$, we have $1-C(t,s)\le\Lambda|t-s|$ for all $t,s\ge0$. The derivatives of $P_Tv_t$ are then $1/2$-H\"older on $(T,\infty)$ with the constant $(1+2\beta p^2\sqrt p)\sqrt{2\Lambda}$, which does not depend on $T$.
\end{remark}

\section{Reflection and complete monotonicity}
\label{sec:selection}

We now study the part of a late state that was already determined at an earlier time. By \cref{prop:stationary-clusters}, entropy production identifies every bounded-lag cluster of the translated kernels as a stationary profile that satisfies the FDT relation. It neither selects the plateau of this profile nor controls the older part of the correlation memory. To study the plateau, we project late states onto the information carried by earlier noise leaves. For $T\ge0$ and $a\ge0$, we define the \emph{prefix center}
\[
 M_T(a)=P_Tv_{T+a}.
\]
It is the part of the state at time $T+a$ built from the initial vector and noise received by time $T$, including memory interactions after $T$. The main observation of this section is that the Gram kernel of the prefix centers satisfies a closed equation on a fixed future rectangle. Every memory term that involves a time before $T$ cancels exactly, before any limit is taken. The stationary FDT relation then identifies the limit of this kernel. This is the reflection identity, which says that the overlap of the centers at offsets $a$ and $b$ converges to $c(a+b)$.

Recall that a stationary cluster $(c,r,\ell)$ is a locally uniform limit of the translated kernels $(C_{T_n},R_{T_n},\lambda_{T_n})$ along a sequence $T_n\to\infty$, as in \cref{prop:stationary-clusters}, and that its energy is $e_\infty=(\ell-1)/(p\beta)$.

\begin{proposition}[Reflection of a stationary cluster]
\label{sel:prefix-reflection}
Fix $p\ge3$ and $\beta>0$. Let $T_n\to\infty$ be a sequence along which the translated kernels converge to a stationary cluster $(c,r,\ell)$, as in \cref{prop:stationary-clusters}, and let $e_\infty=(\ell-1)/(p\beta)$ be its energy. Then, for every finite $A\ge0$,
\begin{equation}
 \lim_{n\to\infty}\sup_{0\le a,b\le A}\bigl|\ip{M_{T_n}(a)}{M_{T_n}(b)}-c(a+b)\bigr|=0.
 \label{sel:prefix-gram}
\end{equation}
Moreover, for every fixed $L\ge0$,
\begin{align}
 &\lim_{n\to\infty}\norm{M_{T_n}(L)}^2=d_L,\qquad d_L=c(2L),
 \label{sel:prefix-norm}\\
 &\lim_{n\to\infty}h\bigl(M_{T_n}(L)\bigr)=h_L,\qquad
 h_L=e_\infty+\beta p\int_0^Lr(\tau)\bigl\{c(2L-\tau)^{p-1}-c(\tau)^{p-1}\bigr\}\,\dd\tau.
 \label{sel:prefix-energy}
\end{align}
\end{proposition}

The plateau $q$ of the cluster may be zero; the proposition covers every $p\ge3$, every $\beta>0$, and every stationary cluster.

\begin{remark}[Independent futures]\label{rem:tree-independent-futures}
The reflected lag $a+b$ has a simple heuristic explanation. Let $(x_t)$ be a stationary reversible Markov process in $\R^N$ with correlation $c(t)=\E[x_0\cdot x_t]/N$. Two continuations from a common state $x_T$, with independent noise after time $T$, have the conditional means $\E[x_{T+a}\mid x_T]$ and $\E[x_{T+b}\mid x_T]$ at times $T+a$ and $T+b$. The Markov semigroup of the process is self-adjoint in the stationary law, so stationarity and the semigroup property give
\[
 \frac1N\E\bigl[\E[x_{T+a}\mid x_T]\cdot\E[x_{T+b}\mid x_T]\bigr]=\frac1N\E[x_0\cdot x_{a+b}]=c(a+b).
\]
Trajectories with a common past and independent future noise were studied by Barrat, Burioni, and M\'ezard, and the formula $c(a+b)$ for the overlap of two such continuations appears in \cite[equation~(9) and Section~4.1]{bbm1996}. Franz and Virasoro proposed weighted averages over cloned futures to improve estimates of quasi-state magnetizations \cite[arXiv version~3, Section~III]{franzvirasoro2000}; see also Biroli \cite[arXiv version~2, Section~III.B]{biroli1999}. The prefix center plays the role of the conditional mean in the tree space. The CK dynamics is not stationary, and \cref{sel:prefix-reflection} justifies the reflected overlap for the stationary clusters of the CK solution, by an exact cancellation in the CK equations.
\end{remark}

\begin{proof}[Proof of \cref{sel:prefix-reflection}]
The proof has four steps. First, for every fixed $T$, we derive a closed equation for the projected Gram kernel on the future rectangle. Second, we pass to the limit in its coefficients and show that the limiting equation is stable, in the sense that any solution with values in $[0,1]$ is the limit of the projected kernels. Third, we verify with Laplace transforms that $c(a+b)$ is such a solution, which proves \eqref{sel:prefix-gram} and \eqref{sel:prefix-norm}. Fourth, we compute the energy of a prefix center by an exact subtraction. Throughout, we write $B=\beta^2p$.

\paragraph{Step 1: the closed equation on the future rectangle.}
Fix $T\ge0$, and let $K_T(t,s)=\ip{P_Tv_t}{P_Tv_s}$ for $t,s\ge0$. We apply $P_T$ to \eqref{eq:geometry-tree-realization}. The one-leaf component $P_T\eta_t$ is orthogonal to the grafted part, $P_T$ commutes with grafting by \eqref{eq:geometry-prefix-commutation}, and $\mathcal I$ is an isometry. As in Step 2 of the proof of \cref{prop:geometry-tree-realization}, we obtain, for $t,s\ge T$,
\begin{equation}\label{eq:tree-projected-green}
\begin{split}
 K_T(t,s)={}&R(t,0)R(s,0)+2\int_0^TR(t,u)R(s,u)\,\dd u\\
 &+B\int_0^t\int_0^sR(t,u)R(s,w)K_T(u,w)^m\,\dd w\,\dd u.
\end{split}
\end{equation}
Here the noise integral stops at $T$, since the noise components of $P_T\eta_t$ and $P_T\eta_s$ vanish beyond $T$. The projection has no effect as soon as one argument lies in the past. Indeed, if $u\le T$, then $v_u\in\mathcal H_T$, which implies that $P_Tv_u=v_u$ and
\[
 K_T(u,w)=\ip{v_u}{P_Tv_w}=\ip{P_Tv_u}{v_w}=\ip{v_u}{v_w}=C(u,w).
\]
The same holds if $w\le T$. We subtract \eqref{eq:tree-projected-green} from the Green identity \eqref{eq:green-covariance}. The initial terms cancel, the noise integrals over $[0,T]$ cancel, and in the memory integrals every block with $u\le T$ or $w\le T$ cancels. Only the integrals over $(T,t]\times(T,s]$ remain, and we have
\[
\begin{aligned}
 C(t,s)-K_T(t,s)={}&2\int_T^{t\wedge s}R(t,u)R(s,u)\,\dd u\\
 &+B\int_T^t\int_T^sR(t,u)R(s,w)\bigl\{C(u,w)^m-K_T(u,w)^m\bigr\}\,\dd w\,\dd u.
\end{aligned}
\]
For $a,b\ge0$, let $Q_T(a,b)=K_T(T+a,T+b)=\ip{M_T(a)}{M_T(b)}$. Shifting all times by $T$, we arrive at
\begin{align}
 Q_T(a,b)={}&C(T+a,T+b)
 -2\int_0^{a\wedge b}R(T+a,T+u)R(T+b,T+u)\,\dd u
 \notag\\
 &-B\int_0^a\int_0^b R(T+a,T+u)R(T+b,T+w)
 \notag\\
 &\hspace{6em}\cdot\bigl\{C(T+u,T+w)^m-Q_T(u,w)^m\bigr\}\,\dd w\,\dd u.
 \label{sel:closed-prefix-rectangle}
\end{align}
This equation holds for every $T$; no limit has been taken, and no decay of the older memory has been used. We also note that
\[
 0\le Q_T(a,b)\le C(T+a,T+b)\le1.
\]
Indeed, the components of $P_Tv_t$ and $(\Id-P_T)v_t$ are nonnegative, which implies $C(t,s)=K_T(t,s)+\ip{(\Id-P_T)v_t}{(\Id-P_T)v_s}\ge K_T(t,s)\ge0$.

\paragraph{Step 2: the limiting equation, uniqueness, and stability.}
We fix $0\le A<\infty$ and write $Q_n=Q_{T_n}$, $C_n(a,b)=C(T_n+a,T_n+b)$, and $R_n(a,u)=R(T_n+a,T_n+u)$. We also abbreviate $\Theta_n(u,w)=R_n(a,u)R_n(b,w)$ and $\theta(u,w)=r(a-u)r(b-w)$, for fixed $a$ and $b$. By \cref{prop:stationary-clusters}, we have $C_n(a,b)\to c(|a-b|)$ uniformly for $a,b\in[0,A]$, and $R_n(a,u)\to r(a-u)$ uniformly for $0\le u\le a\le A$. The limiting form of \eqref{sel:closed-prefix-rectangle} is the equation
\begin{align}
 Q(a,b)={}&c(|a-b|)-2\int_0^{a\wedge b}r(a-u)r(b-u)\,\dd u
 \notag\\
 &-B\int_0^a\int_0^b r(a-u)r(b-w)\bigl\{c(|u-w|)^m-Q(u,w)^m\bigr\}\,\dd w\,\dd u
 \label{sel:limiting-prefix-rectangle}
\end{align}
for a function $Q$ on $[0,A]^2$. We first show that it has at most one bounded solution. Let $Q$ and $Q'$ be two solutions with $|Q|,|Q'|\le\Gamma$, and let $\Delta=Q-Q'$. Using $0\le r\le1$ and $|x^m-y^m|\le m\Gamma^{m-1}|x-y|$ for $|x|,|y|\le\Gamma$, we have
\[
 |\Delta(a,b)|\le Bm\Gamma^{m-1}\int_0^a\int_0^b|\Delta(u,w)|\,\dd w\,\dd u,
\]
and the Gronwall argument of Step 3 in the proof of \cref{prop:geometry-tree-realization} shows that $\Delta=0$.

Next, we prove stability. Let $Q$ be a solution of \eqref{sel:limiting-prefix-rectangle} with values in $[0,1]$. Subtracting \eqref{sel:limiting-prefix-rectangle} from \eqref{sel:closed-prefix-rectangle} at $T=T_n$, we find
\begin{align*}
 Q_n(a,b)-Q(a,b)={}&\bigl\{C_n(a,b)-c(|a-b|)\bigr\}
 -2\int_0^{a\wedge b}\bigl\{\Theta_n(u,u)-\theta(u,u)\bigr\}\,\dd u\\
 &-B\int_0^a\int_0^b\bigl\{\Theta_n(u,w)C_n(u,w)^m-\theta(u,w)c(|u-w|)^m\bigr\}\,\dd w\,\dd u\\
 &+B\int_0^a\int_0^b\bigl\{\Theta_n(u,w)-\theta(u,w)\bigr\}Q(u,w)^m\,\dd w\,\dd u\\
 &+B\int_0^a\int_0^b\Theta_n(u,w)\bigl\{Q_n(u,w)^m-Q(u,w)^m\bigr\}\,\dd w\,\dd u.
\end{align*}
Let $\varepsilon_n$ be the supremum over $a,b\in[0,A]$ of the sum of the absolute values of the first four terms on the right. All kernels take values in $[0,1]$ and converge uniformly, and the intervals and rectangles of integration have measure bounded in terms of $A$ only. Then $\varepsilon_n\to0$. In the last term, we use $0\le\Theta_n\le1$ and $|x^m-y^m|\le m|x-y|$ for $x,y\in[0,1]$, which applies by Step 1. This gives
\[
 |Q_n(a,b)-Q(a,b)|\le\varepsilon_n+Bm\int_0^a\int_0^b|Q_n(u,w)-Q(u,w)|\,\dd w\,\dd u,\qquad a,b\in[0,A].
\]
Let $\rho_n(z)=\sup_{a,b\le z}|Q_n(a,b)-Q(a,b)|$. As in Step 3 of the proof of \cref{prop:geometry-tree-realization}, we have $\rho_n(z)\le\varepsilon_n+2BmA\int_0^z\rho_n(y)\,\dd y$ for $z\le A$, and Gronwall's inequality implies that
\begin{equation}\label{eq:tree-prefix-stability}
 \sup_{0\le a,b\le A}|Q_n(a,b)-Q(a,b)|\le\varepsilon_ne^{2BmA^2}\longrightarrow0.
\end{equation}
It remains to find a solution of \eqref{sel:limiting-prefix-rectangle} with values in $[0,1]$.

\paragraph{Step 3: identification of the solution through Laplace transforms.}
We check that $Q(a,b)=c(a+b)$ solves \eqref{sel:limiting-prefix-rectangle}. Note that $0\le c\le1$. Step 2 then proves \eqref{sel:prefix-gram}, and the case $a=b=L$ proves \eqref{sel:prefix-norm}.

For a bounded function $f$ on $[0,\infty)$, let $D_f(a,b)=f(|a-b|)-f(a+b)$. Note that $c(|u-w|)^m-c(u+w)^m=D_{c^m}(u,w)$. Then the function $c(a+b)$ solves \eqref{sel:limiting-prefix-rectangle} if and only if, for all $a,b\ge0$,
\begin{equation}\label{eq:tree-reflection-difference}
 D_c(a,b)=2\int_0^{a\wedge b}r(a-u)r(b-u)\,\dd u
 +B\int_0^a\int_0^b r(a-u)r(b-w)D_{c^m}(u,w)\,\dd w\,\dd u.
\end{equation}
We prove \eqref{eq:tree-reflection-difference} by comparing double Laplace transforms. For a bounded measurable $f$ on $[0,\infty)$ and $z>0$, we write $\widehat f(z)=\int_0^\infty e^{-zt}f(t)\,\dd t$, and for a bounded measurable $F$ on $[0,\infty)^2$ and $z_1,z_2>0$, we write $\widehat F(z_1,z_2)=\int_0^\infty\int_0^\infty e^{-z_1a-z_2b}F(a,b)\,\dd b\,\dd a$.

We begin with the transform of $D_f$. Splitting $[0,\infty)^2$ into the regions $a\ge b$ and $a<b$, and substituting $a=b+x$ in the first and $b=a+x$ in the second, we find
\[
 \int_0^\infty\int_0^\infty e^{-z_1a-z_2b}f(|a-b|)\,\dd b\,\dd a=\frac{\widehat f(z_1)+\widehat f(z_2)}{z_1+z_2}.
\]
Substituting $y=a+b$ and integrating over $a\in[0,y]$, we have, for $z_1\ne z_2$,
\[
 \int_0^\infty\int_0^\infty e^{-z_1a-z_2b}f(a+b)\,\dd b\,\dd a
 =\int_0^\infty f(y)\,\frac{e^{-z_2y}-e^{-z_1y}}{z_1-z_2}\,\dd y
 =\frac{\widehat f(z_2)-\widehat f(z_1)}{z_1-z_2}.
\]
Subtracting the two identities and putting them over the common denominator $z_1^2-z_2^2$, we find that
\begin{equation}\label{eq:tree-laplace-difference}
 \widehat{D_f}(z_1,z_2)=\frac{2\{z_1\widehat f(z_1)-z_2\widehat f(z_2)\}}{z_1^2-z_2^2},\qquad z_1\ne z_2.
\end{equation}

Next, we compute the transforms of the stationary profiles. Let $\sigma(u)=Bm\,c(u)^{m-1}r(u)=-B(c^m)'(u)$. The stationary response equation \eqref{eq:stationary-cluster-response} reads $r'=-\ell r+\sigma*r$, where $*$ denotes convolution on $[0,\infty)$. The functions $r$ and $\sigma$ are nonnegative and integrable, with $\int_0^\infty r=1-q$ and $\int_0^\infty\sigma=B(1-q^m)$. Then $r'$ is integrable, and an integration by parts gives $z\widehat r(z)-1=-\ell\widehat r(z)+\widehat\sigma(z)\widehat r(z)$. Recall that $r$ is continuous and nonnegative with $r(0)=1$, so that $\widehat r(z)>0$. This yields
\begin{equation}\label{eq:tree-laplace-response}
 \widehat r(z)^{-1}=z+\ell-\widehat\sigma(z).
\end{equation}
In the same way, $c'=-r$ and $(c^m)'=-\sigma/B$, together with $c(0)=1$, give
\begin{equation}\label{eq:tree-laplace-correlation}
 z\widehat c(z)=1-\widehat r(z),\qquad z\widehat{c^m}(z)=1-\widehat\sigma(z)/B.
\end{equation}

We now compute the transform of the right side of \eqref{eq:tree-reflection-difference}. Integrating first over $a,b\ge u$, we find that the transform of the noise term is
\[
 2\int_0^\infty e^{-(z_1+z_2)u}\,\dd u\;\widehat r(z_1)\widehat r(z_2)=\frac{2\widehat r(z_1)\widehat r(z_2)}{z_1+z_2}.
\]
The memory term is a two-dimensional convolution of $r(a)r(b)$ with the bounded function $D_{c^m}$, and its transform is $B\widehat r(z_1)\widehat r(z_2)\widehat{D_{c^m}}(z_1,z_2)$. All these integrals converge absolutely, because $r$ is integrable and $D_{c^m}$ is bounded, which justifies the use of Fubini's theorem. Using \eqref{eq:tree-laplace-difference} with $f=c^m$ and \eqref{eq:tree-laplace-correlation}, we can write
\[
 B\,\widehat{D_{c^m}}(z_1,z_2)=\frac{2\{\widehat\sigma(z_2)-\widehat\sigma(z_1)\}}{z_1^2-z_2^2}.
\]
Adding the transforms of the two terms, we have, for $z_1\ne z_2$,
\begin{align*}
 \frac{2\widehat r(z_1)\widehat r(z_2)}{z_1+z_2}
 +B\widehat r(z_1)\widehat r(z_2)\widehat{D_{c^m}}(z_1,z_2)
 &=\frac{2\widehat r(z_1)\widehat r(z_2)}{z_1^2-z_2^2}\bigl\{z_1-z_2+\widehat\sigma(z_2)-\widehat\sigma(z_1)\bigr\}\\
 &=\frac{2\widehat r(z_1)\widehat r(z_2)}{z_1^2-z_2^2}\bigl\{\widehat r(z_1)^{-1}-\widehat r(z_2)^{-1}\bigr\}\\
 &=\frac{2\{\widehat r(z_2)-\widehat r(z_1)\}}{z_1^2-z_2^2}
 =\widehat{D_c}(z_1,z_2).
\end{align*}
In the second line, we used \eqref{eq:tree-laplace-response} at $z_1$ and at $z_2$. In the last equality, we used \eqref{eq:tree-laplace-difference} with $f=c$ and the first identity in \eqref{eq:tree-laplace-correlation}.

Both sides of \eqref{eq:tree-reflection-difference} are bounded continuous functions on $[0,\infty)^2$. Indeed, the left side is bounded by one, the noise term by $2\int_0^\infty r=2(1-q)$, and the memory term by $B(1-q)^2$. We have shown that their double Laplace transforms agree for $z_1\ne z_2$, and by continuity of the transforms in $(z_1,z_2)$ they agree everywhere on $(0,\infty)^2$. Let $F$ denote either side. For fixed $z_2>0$, the function $a\mapsto\int_0^\infty e^{-z_2b}F(a,b)\,\dd b$ is bounded and continuous, by dominated convergence, and its Laplace transform is $\widehat F(\cdot,z_2)$. By the uniqueness theorem for Laplace transforms of continuous functions, the two such functions coincide for every $a\ge0$. For fixed $a$, the functions $b\mapsto F(a,b)$ are bounded and continuous and, by what we just proved, have the same Laplace transform. They therefore coincide. This proves \eqref{eq:tree-reflection-difference}, and hence \eqref{sel:prefix-gram} and \eqref{sel:prefix-norm}.

\paragraph{Step 4: the energy subtraction.}
By \eqref{eq:geometry-prefix-root-frame}, we have
\[
 h(P_Tv_t)=\sqrt p\,\ip{\mathcal I^*P_Tv_t}{(P_Tv_t)^{\otimes m}}
 =\beta p\int_0^tR(t,u)K_T(t,u)^m\,\dd u.
\]
For $u\le T$, we have $K_T(t,u)=C(t,u)$ by Step 1. We subtract the identity $h(v_t)=e_\beta(t)=\beta p\int_0^tR(t,u)C(t,u)^m\,\dd u$ of \cref{prop:geometry-tree-realization} from the last display. The integrals over $[0,T]$ cancel. At $t=T+L$, using $K_T(T+L,u)=Q_T(L,u-T)=Q_T(u-T,L)$, we obtain
\begin{equation}
 h(P_Tv_{T+L})
 =e_\beta(T+L)+\beta p\int_T^{T+L}R(T+L,u)
 \bigl\{Q_T(u-T,L)^m-C(T+L,u)^m\bigr\}\,\dd u.
 \label{sel:exact-prefix-energy-subtraction}
\end{equation}
Only an integral over a future interval of fixed length $L$ remains. Now let $T=T_n$. By \cref{prop:stationary-clusters}, we have $e_\beta(T_n+L)\to e_\infty$. With the change of variables $\tau=T_n+L-u\in[0,L]$, we have
\begin{gather*}
 R(T_n+L,T_n+L-\tau)\to r(\tau),\qquad C(T_n+L,T_n+L-\tau)\to c(\tau),\\
 Q_n(L-\tau,L)\to c(2L-\tau),
\end{gather*}
uniformly in $\tau\in[0,L]$, by \cref{prop:stationary-clusters} and \eqref{sel:prefix-gram}. All these functions take values in $[0,1]$, and the integral in \eqref{sel:exact-prefix-energy-subtraction} converges to the integral in \eqref{sel:prefix-energy}. This completes the proof.
\end{proof}

By the reflection identity, the Hankel kernel $(a,b)\mapsto c(a+b)$ is a limit of Gram kernels. This leads to the following representation of every stationary profile.

\begin{corollary}[Spectral representation of every stationary cluster]
\label{sel:cluster-complete-monotonicity}
Let $(c,r,\ell)$ be a stationary cluster, as in \cref{prop:stationary-clusters}, with plateau $q$. Then there exists a probability measure $\nu$ on $[0,\infty)$ such that
\begin{equation}
 c(t)=\int_{[0,\infty)}e^{-xt}\,\nu(\dd x)\quad(t\ge0),\qquad
 \nu(\{0\})=q.
 \label{eq:stationary-spectral-measure}
\end{equation}
In particular, $c$ and $r=-c'$ are completely monotone on $(0,\infty)$.
\end{corollary}

\begin{proof}
The representation of a bounded continuous function with positive semidefinite Hankel matrices as a Laplace transform is classical; see Widder \cite[Chapter~IV, Section~12, Theorems~12a--12b, pp.~160--161]{widder1941} for Bernstein's theorem on completely monotone functions, and Berg, Christensen, and Ressel \cite[Chapter~4, Corollary~4.5, p.~114]{bcr1984} for positive definite functions on semigroups. Complete monotonicity for Brownian mode-coupling equations is also classical; Franosch and Voigtmann proved it for their mixture model by an iteration argument \cite[arXiv version~1, Section~III]{fv2002}. Here it follows from the reflection identity, and we include the short semigroup argument to keep the proof self-contained.
The argument has three steps. We show that the Hankel kernel is positive semidefinite, construct a semigroup of shifts, and apply the spectral theorem.

\paragraph{Step 1: positivity.}
Let $T_n\to\infty$ be a sequence along which the translated kernels converge to $(c,r,\ell)$. For $a_1,\dots,a_k\ge0$ and $\alpha_1,\dots,\alpha_k\in\R$, we have
\[
 \sum_{i,j=1}^k\alpha_i\alpha_j\ip{M_{T_n}(a_i)}{M_{T_n}(a_j)}=\Bigl\|\sum_{i=1}^k\alpha_iM_{T_n}(a_i)\Bigr\|^2\ge0.
\]
Letting $n\to\infty$ and using \eqref{sel:prefix-gram}, we obtain
\begin{equation}\label{eq:tree-hankel-positive}
 \sum_{i,j=1}^k\alpha_i\alpha_jc(a_i+a_j)\ge0.
\end{equation}

\paragraph{Step 2: the shift semigroup.}
Let $\mathcal V$ be the real vector space of finite formal linear combinations of symbols $\varphi_a$, $a\ge0$, with the bilinear form determined by $\ip{\varphi_a}{\varphi_b}=c(a+b)$. By \eqref{eq:tree-hankel-positive}, this form is positive semidefinite. Let $\mathcal G$ be the Hilbert space obtained by taking the quotient of $\mathcal V$ by its null vectors and completing. For $w=\sum_i\alpha_i\varphi_{a_i}\in\mathcal V$ and $\sigma\ge0$, let $w_\sigma=\sum_i\alpha_i\varphi_{a_i+\sigma}$ and
\[
 \psi_w(\sigma)=\ip{w_\sigma}{w_\sigma}=\sum_{i,j}\alpha_i\alpha_jc(a_i+a_j+2\sigma).
\]
Since $|c|\le1$, we see that $\psi_w(\sigma)\le(\sum_i|\alpha_i|)^2$. For $n\ge0$, we have $\psi_w((n+1)\sigma)=\ip{w_{n\sigma}}{w_{(n+2)\sigma}}$, and the Cauchy--Schwarz inequality gives
\begin{equation}\label{eq:tree-log-convex}
 \psi_w((n+1)\sigma)^2\le\psi_w(n\sigma)\,\psi_w((n+2)\sigma).
\end{equation}
We claim that $\psi_w(\sigma)\le\psi_w(0)$. If $\psi_w(0)=0$, then \eqref{eq:tree-log-convex} with $n=0$ implies that $\psi_w(\sigma)=0$. Suppose that $\psi_w(0)>0$ and $\psi_w(\sigma)>\psi_w(0)$, and let $\kappa=\psi_w(\sigma)/\psi_w(0)>1$. Using \eqref{eq:tree-log-convex} and induction on $n$, we find that $\psi_w((n+1)\sigma)\ge\kappa\,\psi_w(n\sigma)>0$ for every $n\ge0$. Then $\psi_w(n\sigma)\ge\kappa^n\psi_w(0)\to\infty$, which contradicts the bound $\psi_w\le(\sum_i|\alpha_i|)^2$. This proves the claim.

The claim says that $w\mapsto w_\sigma$ does not increase the norm. In particular, it maps null vectors to null vectors, and it defines a contraction $S_\sigma$ of $\mathcal G$ with $S_\sigma\varphi_a=\varphi_{a+\sigma}$. Each $S_\sigma$ is self-adjoint, because $\ip{S_\sigma\varphi_a}{\varphi_b}=c(a+b+\sigma)=\ip{\varphi_a}{S_\sigma\varphi_b}$. Further, $S_0=\Id$ and $S_\sigma S_{\sigma'}=S_{\sigma+\sigma'}$, so that $S_\sigma=S_{\sigma/2}^*S_{\sigma/2}$ is a positive operator.

\paragraph{Step 3: the spectral representation.}
For an integer $j\ge1$, the operator $S_{1/j}$ is positive and $(S_{1/j})^j=S_1$. By uniqueness of positive $j$-th roots, we have $S_{1/j}=S_1^{1/j}$, and therefore $S_\sigma=S_1^\sigma$ for every rational $\sigma\ge0$. Let $\mu$ be the spectral measure of the positive contraction $S_1$ at the vector $\varphi_0$. It is a measure on $[0,1]$ with total mass $\norm{\varphi_0}^2=c(0)=1$, and for rational $\sigma>0$ we have
\[
 c(\sigma)=\ip{\varphi_0}{S_\sigma\varphi_0}=\int_{(0,1]}y^\sigma\,\mu(\dd y).
\]
Both sides are continuous in $\sigma>0$, the right side by dominated convergence, and the identity extends to every $\sigma>0$. Letting $\sigma\downarrow0$, we see that $\mu((0,1])=c(0)=1$, that is, $\mu(\{0\})=0$. Let $\nu$ be the image of $\mu$ under the map $y\mapsto-\log y$ from $(0,1]$ to $[0,\infty)$. Then $\nu$ is a probability measure, and the identity above becomes the first relation in \eqref{eq:stationary-spectral-measure}, also at $t=0$. By dominated convergence, $c(t)\to\nu(\{0\})$ as $t\to\infty$, which gives $\nu(\{0\})=q$. Finally, for $n\ge0$ and $t>0$, the function $x\mapsto x^ne^{-xt}$ is bounded on $[0,\infty)$, uniformly for $t$ in compact subsets of $(0,\infty)$. Differentiating under the integral, we find
\[
 (-1)^nc^{(n)}(t)=\int_{[0,\infty)}x^ne^{-xt}\,\nu(\dd x)\ge0,\qquad
 (-1)^nr^{(n)}(t)=\int_{[0,\infty)}x^{n+1}e^{-xt}\,\nu(\dd x)\ge0,
\]
and $c$ and $r$ are completely monotone on $(0,\infty)$.
\end{proof}

\chapter{Selection of the plateau and correlation aging}\label{sec:stationary-selection}

In this section, we prove the parts of \cref{thm:main-aging} that concern the plateau, the energy, and the loss of memory. We begin with \cref{sel:positive-marginality}, which states that every positive plateau of a stationary cluster is marginal. In \cref{sel:all-temperature}, we prove that the translated kernels, the multiplier, and the energy converge at every temperature. \Cref{sel:selected-sector} identifies the plateau and the energy for $\beta\ge\bd$, and \cref{sel:fixed-wait-forgetting} shows that the solution forgets every fixed earlier state. \Cref{sel:aging} states the two ordered limits of the correlation. At the end of \cref{sec:aging}, we combine these results into a proof of \cref{thm:main-aging}.

We take as input two facts about prefix centers, which we state as \cref{prop:threshold-inputs} in \cref{sec:threshold-inputs} and prove in \cref{sec:tree-ceilings,app:general-selection}. Besides these two facts, we use the stationary clusters of \cref{prop:stationary-clusters}, the scalar constraints of \cref{prop:scalar-clusters}, the uniqueness result \cref{lem:stationary-uniqueness}, the causal tree realization of \cref{prop:geometry-tree-realization}, the reflection identity of \cref{sel:prefix-reflection}, and the uniform bounds of \cref{sec:ck-structure}. The convergence proved in \cref{sel:all-temperature} and the marginality proved in \cref{sel:positive-marginality} are exactly the hypotheses of part~\ref{item:threshold-forgetting} of \cref{prop:threshold-inputs}. We use the selected profile, the marginality relation, and the forgetting theorem throughout \cref{part:response,part:geometry,part:relaxation}.

\paragraph{Standing assumptions and notation.}
Throughout this section, we fix $p\ge3$ and $\beta>0$, and $(C,R,\lambda)$ denotes the uniformly initialized hard CK solution of \cref{sec:model}. All long-time limits below concern this deterministic solution, for which the dimension limit has already been taken on each fixed time horizon. Recall from \eqref{eq:energy-identity} that the energy $e_\beta(t)$ satisfies $\lambda(t)=1+p\beta e_\beta(t)$. For $T\ge0$, the translated kernels are $C_T(a,b)=C(T+a,T+b)$ for $a,b\ge-T$, $R_T(a,b)=R(T+a,T+b)$ for $-T\le b\le a$, and $\lambda_T(a)=\lambda(T+a)$ for $a\ge-T$.

We fix the Fisher majorant $a_F$ and the dimension subsequence $(N_j)$ of \cref{thm:fisher-clock} once, before choosing any waiting-time sequence or stationary cluster. All comparisons with finite-dimensional forks use this same dimension subsequence. We may pass to subsequences of waiting times to obtain stationary clusters, without changing the dimension extraction. The majorant is an auxiliary choice; the CK solution and the stationary profile identified below are unique.

As in \cref{prop:stationary-clusters}, a \emph{stationary cluster} is a triple $(c,r,\ell)$ obtained as the locally uniform limit of $(C_{T_n},R_{T_n},\lambda_{T_n})$ along some sequence $T_n\to\infty$. \Cref{prop:stationary-clusters} implies that $c(0)=r(0)=1$, $c'=-r$, $0\le c\le1$, and $r\ge0$, that the plateau $q=\lim_{\tau\to\infty}c(\tau)$ lies in $[0,1)$, and that $\int_0^\infty r(u)\,\dd u=1-q$. By \eqref{eq:response-derivative-bounds}, we have $0\le R\le C$, and passing to the limit gives $r\le c$. As in \cref{prop:stationary-clusters,sel:prefix-reflection}, we write $e_\infty=(\ell-1)/(p\beta)$ for the energy of the cluster. With $B=\beta^2p$ and $\gamma=\ell-B$, we use the functions of \cref{prop:scalar-clusters},
\[
 \phi(y)=\gamma+By^{p-1},\qquad F(y)=\phi(y)(1-y)-1.
\]
Recall from that proposition that $F(q)=0$, that $F(y)<0$ for $q<y<1$, and that
\[
 \ell=\frac1{1-q}+\beta^2p\bigl(1-q^{p-1}\bigr).
\]

We also use the tree realization of \cref{sec:geometry-core}. The unit vectors $v_t$ satisfy $\ip{v_t}{v_s}=C(t,s)$, and $v_t$ lies in the closed subspace $\mathcal H_t$ generated by trees whose noise leaves carry times at most $t$. We write $P_T$ for the orthogonal projection onto $\mathcal H_T$. The tree energy is $h(v)=\sqrt p\,\ip{v}{\mathcal I(v^{\otimes m})}$ with $m=p-1$, and $h(v_t)=e_\beta(t)$.

\section{Prefix centers and the threshold inputs}\label{sec:threshold-inputs}

For $T\ge0$ and $a\ge0$, we define the prefix center
\[
 M_T(a)=P_Tv_{T+a}.
\]
Let $T_n\to\infty$ be a sequence along which the translated kernels converge to a stationary cluster $(c,r,\ell)$ with energy $e_\infty$. Recall from \cref{sel:prefix-reflection} that, for every fixed $L\ge0$,
\[
\begin{aligned}
 \norm{M_{T_n}(L)}^2&\longrightarrow d_L=c(2L),\\
 h\bigl(M_{T_n}(L)\bigr)&\longrightarrow
 h_L=e_\infty+\beta p\int_0^Lr(\tau)\bigl\{c(2L-\tau)^{p-1}-c(\tau)^{p-1}\bigr\}\,\dd\tau.
\end{aligned}
\]
Note that $h$ is homogeneous of degree $p$. If $d_L>0$, the energy of the normalized center $M_{T_n}(L)/\norm{M_{T_n}(L)}$ converges to $h_L/d_L^{p/2}$.

We now explain why the prefix centers are the right objects for selecting the plateau. The state $v_{T+L}$ is built from the initial vector and from the noise leaves at times up to $T+L$. The projection $P_T$ keeps the part built from the initial vector and noise received by time $T$, including memory interactions after $T$. In physical terms, we may think of $M_T(L)$ as the average of the state at time $T+L$ over independent continuations of the thermal noise after time $T$. %
This interpretation goes back to the trajectories with a common past and independent futures studied by Cugliandolo and Dean \cite[arXiv version~1, Section~II]{cugliandolodean1995} and by Barrat, Burioni, and M\'ezard \cite[Section~4]{bbm1996}, and to the averages over cloned futures proposed by Franz and Virasoro \cite[arXiv version~3, Section~III]{franzvirasoro2000}. In agreement with this interpretation, the limiting squared norm $d_L=c(2L)$ is the stationary correlation at the reflected lag $2L$, which is the overlap of two continuations of length $L$ from a common past; this reflected formula appears in \cite[Section~2, Equation~(9)]{bbm1996}. Time-averaged magnetizations of a similar kind appear in the analysis of Cugliandolo and Kurchan \cite[arXiv version~1, Section~5.1]{ck1994}, and noise-averaged magnetizations and their relation to the threshold appear in the dynamical TAP approach of Biroli \cite[arXiv version~2, Section~III.B]{biroli1999}. We use three properties of the centers.
\begin{itemize}
\item Every earlier state $v_s$ with $s\le T$ lies in $\mathcal H_T$, so that $\ip{v_{T+L}}{v_s}=\ip{M_T(L)}{v_s}$. The correlation of a late state with a fixed earlier state is carried entirely by the center.
\item For a stationary cluster with $q>0$, the normalized limiting energies of the centers satisfy $h_L/d_L^{p/2}\le\Eth$ for every $L\ge0$. Here $\Eth=2\sqrt{(p-1)/p}$ is the threshold energy of the pure $p$-spin landscape \cite[Equation~(2.13)]{abc2013}.
\item Suppose that the whole translated family converges to one stationary profile whose plateau is positive and marginal. Then the normalized energies of the centers approach $\Eth$ as $L\to\infty$, and this forces the solution to forget every fixed earlier state.
\end{itemize}
The first property is exact. The ceiling in the second and the forgetting in the third are the content of the following proposition, which is the only input that we take from \cref{sec:tree-ceilings,app:general-selection}.

\begin{proposition}[Threshold inputs]\label{prop:threshold-inputs}
Fix $p\ge3$ and $\beta>0$.
\begin{enumerate}[label=(\alph*)]
\item\label{item:threshold-ceiling}
Let $T_n\to\infty$ be a sequence along which the translated kernels converge to a stationary cluster $(c,r,\ell)$ with plateau $q>0$, as in \cref{prop:stationary-clusters}, and let $d_L$ and $h_L$ be the limits in \eqref{sel:prefix-norm} and \eqref{sel:prefix-energy}. Then, for every $L\ge0$,
\begin{equation}\label{eq:threshold-input-ceiling}
 h_L\le \Eth\, d_L^{p/2}.
\end{equation}
\item\label{item:threshold-forgetting}
Suppose that the translated kernels converge along every sequence $T\to\infty$ to one stationary profile $(c,r)$, and that its plateau $q$ is positive and satisfies $\beta^2p(p-1)q^{p-2}(1-q)^2=1$. Then, for every $0\le S<\infty$,
\begin{equation}\label{eq:threshold-input-forgetting}
 \lim_{t\to\infty}\sup_{0\le s\le S}|C(t,s)|=0.
\end{equation}
\end{enumerate}
For $p\in\{3,4\}$, both parts are proved in \cref{sec:tree-ceilings}. For every $p\ge3$, they are proved in \cref{app:general-selection}.
\end{proposition}

For $p\in\{3,4\}$, we prove both parts in \cref{sec:tree-ceilings} using a deterministic inequality on the tree Hilbert space (\cref{thm:geometry-tree-ceiling,thm:geometry-quartic-tree-ceiling}). It states that a unit tree vector has energy strictly below $\Eth$, and that near equality forces its mass to leave every fixed set of tree heights. Applying this inequality to the normalized centers, we obtain part~\ref{item:threshold-ceiling}. For part~\ref{item:threshold-forgetting}, we show that the mass of the centers escapes from every fixed height as they approach the threshold, while an earlier state $v_s$ with $s\le S$ is approximated uniformly by finitely many heights. From the same argument, we also obtain the quantitative estimate \cref{sel:threshold-centers}, which we use again in \cref{resp:chapter,sec:geometry-results,sec:three-time}.

For every $p\ge3$, we prove both parts in \cref{app:general-selection}. There we replace the projection by the mean of finitely many independent futures of the finite-dimensional dynamics from a common past. Using the fixed-horizon approximation of Sellke \cite[Lemmas~3.1--3.2]{sellke2024}, we approximate the normalized mean by globally Lipschitz functions of the disorder coefficients. The algorithmic threshold of Huang and Sellke \cite[Theorem~1 and Corollary~1.8]{hs2023} bounds the energy of these approximations by $\Eth$ in the dimension limit, and we transfer this bound to the mean by sending the approximation error to zero. Part~\ref{item:threshold-ceiling} follows from \cref{sel:finite-forks} by letting the number of futures tend to infinity at fixed $L$. Part~\ref{item:threshold-forgetting} follows from \cref{sel:general-fixed-wait} and the finite-net argument in \cref{sec:general-proof-inputs}. Its proof combines the external-field form of the same threshold with adaptive Gaussian conditioning and an Euler approximation of the earlier state.

The references to later sections do not create a circular argument. The proofs of part~\ref{item:threshold-ceiling} in \cref{sec:tree-ceilings,app:general-selection} apply to an arbitrary stationary cluster with positive plateau, and they use only results of \cref{sec:model,sec:ck-structure,sec:entropy-clusters,sec:scalar,sec:geometry-core}. Part~\ref{item:threshold-forgetting} contains in its hypotheses the two facts that we deduce from part~\ref{item:threshold-ceiling} in this section: convergence of the whole translated family, and marginality of the plateau. Besides these hypotheses, its proofs use the algebraic identity \eqref{sel:center-amgm}, which we prove in \cref{lem:center-energy} without using \cref{prop:threshold-inputs}. In this section, the order of deduction is as follows:
\begin{itemize}
\item part~\ref{item:threshold-ceiling} and \cref{lem:center-energy} imply marginality, \cref{sel:positive-marginality};
\item marginality implies convergence at every temperature, \cref{sel:all-temperature}, and the selection of the plateau, \cref{sel:selected-sector};
\item these results verify the hypotheses of part~\ref{item:threshold-forgetting}, which then gives forgetting, \cref{sel:fixed-wait-forgetting}, and correlation aging, \cref{sel:aging}.
\end{itemize}

\section{Marginality of positive plateaus}\label{sec:marginality}

We set
\begin{equation}
 q_{\mathrm m}=\frac{p-2}{p},\qquad
 \beta_{\mathrm m}^{\,2}=\frac{p^{p-1}}{4(p-1)(p-2)^{p-2}},\qquad
 \Eth=2\sqrt{\frac{p-1}{p}}.
 \label{sel:spinodal-constants}
\end{equation}
To explain the role of $q_{\mathrm m}$ and $\beta_{\mathrm m}$, let $G(q)=p(p-1)q^{p-2}(1-q)^2$ for $q\in[0,1]$, so that the marginality relation $\beta^2p(p-1)q^{p-2}(1-q)^2=1$ reads $\beta^2G(q)=1$. We have
\[
 G'(q)=p(p-1)q^{p-3}(1-q)\bigl\{(p-2)-pq\bigr\}.
\]
The function $G$ vanishes at $0$ and $1$, increases strictly on $[0,q_{\mathrm m}]$, and decreases strictly on $[q_{\mathrm m},1]$. Its maximum is
\[
 G(q_{\mathrm m})=\frac{4(p-1)(p-2)^{p-2}}{p^{p-1}}=\beta_{\mathrm m}^{-2}.
\]
In other words, $\beta_{\mathrm m}$ is the inverse temperature at which $\beta^2G$ first reaches one on $[0,1]$. For $\beta\ge\beta_{\mathrm m}$, we denote by $q_+(\beta)$ the unique root of $\beta^2G(q)=1$ in $[q_{\mathrm m},1)$; at $\beta=\beta_{\mathrm m}$, it is the double root $q_{\mathrm m}$. For $\beta>\beta_{\mathrm m}$, the equation has a second root, which lies in $(0,q_{\mathrm m})$, and for $\beta<\beta_{\mathrm m}$ it has no root in $[0,1]$.

The next lemma computes the limit of the center energies of \cref{sel:prefix-reflection} as $L$ tends to infinity, and compares it with $\Eth$. It is a computation with the scalar constraints of \cref{prop:scalar-clusters} and uses no threshold input. We use it again in \cref{sec:tree-ceilings,app:general-selection}.

\begin{lemma}[Limiting energy of prefix centers]\label{lem:center-energy}
Let $(c,r,\ell)$ be a stationary cluster with plateau $q$, as in \cref{prop:stationary-clusters}, let $e_\infty=(\ell-1)/(p\beta)$, and let $d_L$ and $h_L$ be defined by \eqref{sel:prefix-norm} and \eqref{sel:prefix-energy}. Then $d_L\to q$ and $h_L\to h$ as $L\to\infty$, where
\begin{equation}
 h=e_\infty-\beta\bigl\{1-q^p-pq^{p-1}(1-q)\bigr\}.
 \label{sel:limiting-center-energy}
\end{equation}
If $q>0$, define $z=\beta\sqrt{p(p-1)}\,q^{(p-2)/2}(1-q)>0$. Then
\begin{equation}
 \frac{h}{q^{p/2}}
 =\sqrt{\frac{p-1}{p}}\left(z+\frac1z\right)
 \ge \Eth,
 \label{sel:center-amgm}
\end{equation}
and equality holds in \eqref{sel:center-amgm} if and only if $\beta^2p(p-1)q^{p-2}(1-q)^2=1$.
\end{lemma}

\begin{proof}
As $c$ is nonincreasing with limit $q$, we have $d_L=c(2L)\to q$. For $h_L$, we extend the integrand in \eqref{sel:prefix-energy} by zero outside $[0,L]$ and write the integral as an integral over $[0,\infty)$ of
\[
 g_L(\tau)=
 \begin{cases}
 r(\tau)\bigl\{c(2L-\tau)^{p-1}-c(\tau)^{p-1}\bigr\},&0\le\tau\le L,\\
 0,&\tau>L.
 \end{cases}
\]
For fixed $\tau$, we have $2L-\tau\to\infty$ as $L\to\infty$, and $g_L(\tau)$ converges to $r(\tau)\{q^{p-1}-c(\tau)^{p-1}\}$. Since $0\le c\le1$, we have $|g_L(\tau)|\le r(\tau)$, and $r$ is integrable with $\int_0^\infty r=1-q$. By dominated convergence, we obtain
\[
 \lim_{L\to\infty}h_L
 =e_\infty+\beta pq^{p-1}(1-q)-\beta p\int_0^\infty r(\tau)c(\tau)^{p-1}\,\dd\tau.
\]
Using $r=-c'$, $c(0)=1$, and $c(\tau)\to q$, we find
\[
 \int_0^\infty r(\tau)c(\tau)^{p-1}\,\dd\tau
 =-\frac1p\int_0^\infty\bigl(c^p\bigr)'(\tau)\,\dd\tau
 =\frac{1-q^p}p.
\]
Inserting this into the previous display, we obtain \eqref{sel:limiting-center-energy}.

Now suppose that $q>0$. Using the formula for $\ell$ in \eqref{eq:plateau-multiplier}, we write the energy of the cluster as
\[
 e_\infty=\frac{\ell-1}{p\beta}=\frac{q}{p\beta(1-q)}+\beta\bigl(1-q^{p-1}\bigr).
\]
We insert this into \eqref{sel:limiting-center-energy} and note that $(1-q^{p-1})-(1-q^p)=-q^{p-1}(1-q)$. This gives
\[
 h=\frac{q}{p\beta(1-q)}+\beta(p-1)q^{p-1}(1-q).
\]
By the definition of $z$, we have
\[
 \sqrt{\frac{p-1}p}\,z=\beta(p-1)q^{(p-2)/2}(1-q),
 \qquad
 \sqrt{\frac{p-1}p}\,\frac1z=\frac{q^{-(p-2)/2}}{p\beta(1-q)},
\]
and the sum of the two right sides is $h/q^{p/2}$. This proves the identity in \eqref{sel:center-amgm}. For $z>0$, we have $z+z^{-1}\ge2$, with equality if and only if $z=1$. Recalling that $2\sqrt{(p-1)/p}=\Eth$, we deduce the inequality in \eqref{sel:center-amgm}, and equality holds if and only if $z=1$, that is, if and only if $z^2=\beta^2p(p-1)q^{p-2}(1-q)^2=1$.
\end{proof}

\begin{lemma}[Marginality of every positive cluster plateau]\label{sel:positive-marginality}
Let $(c,r,\ell)$ be a stationary cluster, as in \cref{prop:stationary-clusters}, with plateau $q>0$. Then
\begin{equation}
 \beta^2p(p-1)q^{p-2}(1-q)^2=1,
 \qquad q\ge q_{\mathrm m}.
 \label{sel:positive-plateau-equation}
\end{equation}
\end{lemma}

\begin{proof}
First, part~\ref{item:threshold-ceiling} of \cref{prop:threshold-inputs} gives an upper bound on the limiting center energy that matches the lower bound in \eqref{sel:center-amgm}, and this forces marginality. Second, the strict sign condition of \cref{prop:scalar-clusters} excludes the smaller root.

\paragraph{Step 1: marginality.}
Let $T_n\to\infty$ be a sequence along which the translated kernels converge to $(c,r,\ell)$. Since $q>0$, part~\ref{item:threshold-ceiling} of \cref{prop:threshold-inputs} applies along this sequence, and \eqref{eq:threshold-input-ceiling} states that $h_L\le\Eth d_L^{p/2}$ for every $L\ge0$. Letting $L\to\infty$ and using $d_L\to q$ and $h_L\to h$ from \cref{lem:center-energy}, we obtain
\[
 h\le\Eth\,q^{p/2}.
\]
Together with the lower bound in \eqref{sel:center-amgm}, this forces $h/q^{p/2}=\Eth$. The equality case of \cref{lem:center-energy} then implies the first relation in \eqref{sel:positive-plateau-equation}.

\paragraph{Step 2: the larger root.}
We use the functions $\phi$ and $F$ from the beginning of this section; recall that $F(q)=0$ and $F(y)<0$ for $q<y<1$. Differentiating, we have
\[
 F'(y)=B(p-1)y^{p-2}(1-y)-\phi(y),\qquad
 F''(y)=B(p-1)y^{p-3}\bigl\{(p-2)-py\bigr\}.
\]
Since $\phi(q)=1/(1-q)$ by \eqref{eq:plateau-multiplier}, we find
\[
 F'(q)=\frac{\beta^2p(p-1)q^{p-2}(1-q)^2-1}{1-q}=0,
\]
where we used Step~1 in the last equality. Suppose that $0<q<q_{\mathrm m}$. Then $F''(q)>0$. Since $F$ is a polynomial with $F(q)=F'(q)=0$, Taylor's theorem implies that
\[
 F(q+\varepsilon)=\tfrac12F''(q)\varepsilon^2+O(\varepsilon^3)>0
\]
for all sufficiently small $\varepsilon>0$. This contradicts $F(q+\varepsilon)<0$. We conclude that $q\ge q_{\mathrm m}$.
\end{proof}

The first relation in \eqref{sel:positive-plateau-equation} is the marginality condition of the asymptotic solution of Cugliandolo and Kurchan \cite[Equation~(7)]{ck1993}, written in the conventions of \cref{sec:model}. In view of the discussion after \eqref{sel:spinodal-constants}, the relations \eqref{sel:positive-plateau-equation} say that $\beta\ge\beta_{\mathrm m}$ and $q=q_+(\beta)$. A stationary cluster with positive plateau can occur only for $\beta\ge\beta_{\mathrm m}$, and its plateau is then determined by $\beta$.

\section{Convergence at every fixed temperature}\label{sec:all-temperature}

\begin{proposition}[Convergence at every fixed temperature]\label{sel:all-temperature}
For every $p\ge3$ and $\beta>0$, the multiplier $\lambda(t)$ and the energy $e_\beta(t)$ have limits $\ell_\beta$ and $e_\beta(\infty)$ as $t\to\infty$. All bounded-lag translates of $C$ and $R$ converge locally uniformly to a unique stationary FDT profile $(c,r)$, that is, for every $0\le A<\infty$,
\begin{equation}
 \sup_{0\le s\le t\le A}
 \bigl\{|C(T+t,T+s)-c(t-s)|+|R(T+t,T+s)-r(t-s)|\bigr\}
 \longrightarrow0\qquad(T\to\infty).
 \label{eq:all-temperature-window}
\end{equation}
Its plateau $q$ belongs to the following set:
\begin{equation}
 q\in\begin{cases}
 \{0\}, & 0<\beta<\beta_{\mathrm m},\\
 \{0,q_+(\beta)\}, & \beta\ge\beta_{\mathrm m}.
 \end{cases}
 \label{sel:two-possible-plateaux}
\end{equation}
The positive branch is used only on its domain $\beta\ge\beta_{\mathrm m}$. The limiting multiplier is
\[
 \ell_\beta=\frac1{1-q}+\beta^2p\bigl(1-q^{p-1}\bigr),
\]
as in \eqref{eq:plateau-multiplier}, and $e_\beta(\infty)=(\ell_\beta-1)/(p\beta)$.
\end{proposition}

\begin{proof}
First, the multiplier of every stationary cluster lies in a set with at most two elements. Second, the late limit points of $\lambda$ form an interval, each point of which is the multiplier of a stationary cluster; combined with the first step, this forces $\lambda$ to converge. Third, we identify all stationary clusters using \cref{lem:stationary-uniqueness} and upgrade subsequential convergence to convergence of the whole translated family by precompactness.

\paragraph{Step 1: possible multipliers.}
Let $(c,r,\ell)$ be a stationary cluster with plateau $q$. If $q=0$, then \eqref{eq:plateau-multiplier} gives $\ell=1+\beta^2p$. If $q>0$, then \cref{sel:positive-marginality} and the discussion after it imply that $\beta\ge\beta_{\mathrm m}$ and $q=q_+(\beta)$, and by \eqref{eq:plateau-multiplier} we have $\ell=(1-q_+(\beta))^{-1}+\beta^2p(1-q_+(\beta)^{p-1})$. In both cases, $\ell$ belongs to the set
\[
 \mathcal L_\beta=
 \begin{cases}
 \{1+\beta^2p\}, & \beta<\beta_{\mathrm m},\\[2pt]
 \bigl\{1+\beta^2p,\ (1-q_+(\beta))^{-1}+\beta^2p\bigl(1-q_+(\beta)^{p-1}\bigr)\bigr\}, & \beta\ge\beta_{\mathrm m},
 \end{cases}
\]
which has at most two elements.

\paragraph{Step 2: convergence of the multiplier.}
Recall from \cref{prop:uniform-bounds,prop:global-regularity} that $\lambda$ is continuous on $[0,\infty)$ and satisfies $1\le\lambda\le\Lambda$. Let
\[
 a=\liminf_{t\to\infty}\lambda(t),\qquad b=\limsup_{t\to\infty}\lambda(t),
\]
which are finite. We say that $y$ is a late limit point of $\lambda$ if there exists a sequence $T_n\to\infty$ such that $\lambda(T_n)\to y$. We claim that every $y\in[a,b]$ is a late limit point. For $y\in\{a,b\}$, this holds by the definition of the lower and upper limits. If $a<y<b$, then for every $n$ there exist times $t_1,t_2\ge n$ such that $\lambda(t_1)<y<\lambda(t_2)$. By the intermediate value theorem, there exists $T_n$ between $t_1$ and $t_2$ such that $\lambda(T_n)=y$, and $T_n\ge n$. This proves the claim.

Now let $y\in[a,b]$, and let $T_n\to\infty$ be such that $\lambda(T_n)\to y$. By \cref{prop:stationary-clusters}, a subsequence of the translated family $(C_{T_n},R_{T_n},\lambda_{T_n})$ converges locally uniformly to a stationary cluster $(c,r,\ell)$, and along it $\lambda_{T_n}$ converges to the constant $\ell$. Evaluating at offset zero, we find $\ell=y$, and Step~1 gives $y\in\mathcal L_\beta$. We have shown that
\[
 [a,b]\subseteq\mathcal L_\beta.
\]
An interval contained in a finite set is a single point. Hence $a=b$, and $\lambda(t)$ converges to a limit $\ell_\beta\in\mathcal L_\beta$. By \eqref{eq:energy-identity}, we have $e_\beta(t)=(\lambda(t)-1)/(p\beta)$, and the energy converges to $e_\beta(\infty)=(\ell_\beta-1)/(p\beta)$.

\paragraph{Step 3: identification of the profile.}
Let $(c,r,\ell)$ be any stationary cluster, obtained along a sequence $T_n\to\infty$. Since $\lambda(T_n+a)\to\ell_\beta$ for every offset $a$, we have $\ell=\ell_\beta$. By \eqref{eq:cluster-fdt-properties} and \eqref{eq:stationary-cluster-response}, the pair $(c,r)$ solves
\[
 c'=-r,\qquad
 r'=-\ell_\beta r+\beta^2p(p-1)\bigl(c^{p-2}r\bigr)*r,\qquad
 c(0)=r(0)=1,
\]
and it is bounded, since $0\le r\le c\le1$. \Cref{lem:stationary-uniqueness} states that this system has at most one locally bounded solution. All stationary clusters therefore coincide; we denote their common profile by $(c,r)$. The argument uses only the value $\ell_\beta$ of the multiplier. In particular, it applies when the two elements of $\mathcal L_\beta$ coincide, which happens at $\beta=\bd$. Every stationary cluster also has the same energy $e_\infty=(\ell_\beta-1)/(p\beta)=e_\beta(\infty)$.

Suppose that \eqref{eq:all-temperature-window} fails for some $0\le A<\infty$. Then there exist $\varepsilon>0$ and a sequence $T_n\to\infty$ such that, for every $n$, the supremum in \eqref{eq:all-temperature-window} at $T=T_n$ is at least $\varepsilon$. Using the precompactness in \cref{prop:stationary-clusters}, we extract a subsequence of the translates that converges locally uniformly to a stationary cluster, which is $(c,r,\ell_\beta)$ by the preceding paragraph. Along this subsequence, the supremum over the compact set $\{0\le s\le t\le A\}$ tends to zero, which is a contradiction. This proves \eqref{eq:all-temperature-window}.

Finally, if the plateau $q$ of $c$ is positive, then \cref{sel:positive-marginality} implies that $\beta\ge\beta_{\mathrm m}$ and $q=q_+(\beta)$. This proves \eqref{sel:two-possible-plateaux}, and it shows that the second alternative is absent for $\beta<\beta_{\mathrm m}$. The formula for $\ell_\beta$ is \eqref{eq:plateau-multiplier} for the cluster $(c,r,\ell_\beta)$.
\end{proof}

\section{The dynamical transition and the selected energy}\label{sec:selected-sector}

Recall the dynamical transition parameters
\[
 \qd=\frac{p-2}{p-1},\qquad
 \bd^2=\frac{(p-1)^{p-1}}{p(p-2)^{p-2}}.
\]
We compare them with the constants in \eqref{sel:spinodal-constants}. First, $\qd-q_{\mathrm m}=(p-2)/(p(p-1))>0$. Second, a direct computation shows that $G(\qd)=p(p-2)^{p-2}/(p-1)^{p-1}=\bd^{-2}$. As $G$ decreases strictly on $[q_{\mathrm m},1]$, we find that $\bd^{-2}<G(q_{\mathrm m})=\beta_{\mathrm m}^{-2}$, that is, $\bd>\beta_{\mathrm m}$, and that $q_+(\bd)=\qd$. For $\beta>\bd$, we have $G(q_+(\beta))=\beta^{-2}<G(\qd)$, and then $q_+(\beta)>\qd$. For $\beta\ge\bd$, we write
\[
 q_\beta=q_+(\beta),
\]
so that $q_\beta\ge\qd$, with equality exactly at $\beta=\bd$.

\begin{theorem}[The selected stationary sector]\label{sel:selected-sector}
For every $p\ge3$ and $\beta\ge\bd$, the unique stationary plateau is $q_\beta$. In particular, the limit $\ell_\beta$ of $\lambda(t)$ and the limiting energy satisfy
\begin{align}
 \ell_\beta
 &=\frac1{1-q_\beta}+\beta^2p\bigl(1-q_\beta^{p-1}\bigr),
 \label{sel:selected-multiplier}\\
 e_\beta(\infty)
 &=\beta\bigl\{1+(p-2)q_\beta^{p-1}-(p-1)q_\beta^p\bigr\}.
 \label{sel:selected-energy}
\end{align}
For every $0\le A<\infty$,
\begin{equation}
 \sup_{0\le s\le t\le A}
 \bigl\{|C(T+t,T+s)-c(t-s)|+|R(T+t,T+s)-r(t-s)|\bigr\}
 \longrightarrow0\qquad(T\to\infty),
 \label{sel:stationary-window}
\end{equation}
where $r=-c'$ and $c(\tau)\to q_\beta$.
At the endpoint, $e_{\bd}(\infty)=\bd$ and $\ell_{\bd}=1+p\bd^2$.
\end{theorem}

\begin{proof}
In view of \cref{sel:all-temperature}, it suffices to exclude the alternative $q=0$. Suppose that $q=0$. Then \eqref{eq:plateau-multiplier} gives $\ell_\beta=1+\beta^2p$, which means that $\gamma=\ell_\beta-\beta^2p=1$ and
\[
 F(y)=\bigl(1+\beta^2py^{p-1}\bigr)(1-y)-1
 =y\bigl\{\beta^2py^{p-2}(1-y)-1\bigr\}.
\]
Since $\qd^{\,p-2}(1-\qd)=(p-2)^{p-2}/(p-1)^{p-1}$, we have $\beta^2p\,\qd^{\,p-2}(1-\qd)=\beta^2/\bd^2$ and
\[
 F(\qd)=\qd\left(\frac{\beta^2}{\bd^2}-1\right).
\]
For $\beta\ge\bd$, this is nonnegative. On the other hand, $\qd$ lies in $(0,1)=(q,1)$, so $F(\qd)<0$ by \eqref{eq:scalar-sign}. This contradiction shows that $q>0$.

At $\beta=\bd$, we have $F(\qd)=0$, and the contradiction comes from the strict inequality in \eqref{eq:scalar-sign}. In \cref{prop:scalar-clusters}, that strict inequality follows from the bound $c'(t)\ge F(c(t))$ and from $c'=-r<0$ at every finite time. In this way, the endpoint $\beta=\bd$ is included without any continuity argument in $\beta$.

We have shown that $q>0$, and \cref{sel:all-temperature} then implies that $q=q_+(\beta)=q_\beta$. The formula \eqref{sel:selected-multiplier} is the formula for $\ell_\beta$ in \cref{sel:all-temperature}, \eqref{sel:stationary-window} is \eqref{eq:all-temperature-window}, and the relation $r=-c'$ is part of \eqref{eq:cluster-fdt-properties}. For the energy, we use the marginality relation $\beta^2G(q_\beta)=1$. Multiplying it by $q_\beta/(p\beta(1-q_\beta))$, we obtain
\[
 \frac{q_\beta}{p\beta(1-q_\beta)}=\beta(p-1)q_\beta^{p-1}(1-q_\beta).
\]
Using $e_\beta(\infty)=(\ell_\beta-1)/(p\beta)$ and \eqref{sel:selected-multiplier}, we find
\[
 e_\beta(\infty)
 =\frac{q_\beta}{p\beta(1-q_\beta)}+\beta\bigl(1-q_\beta^{p-1}\bigr)
 =\beta(p-1)q_\beta^{p-1}(1-q_\beta)+\beta\bigl(1-q_\beta^{p-1}\bigr),
\]
and expanding the right side yields \eqref{sel:selected-energy}. At $\beta=\bd$, we have $q_\beta=\qd$ and $(p-1)\qd=p-2$, so that
\[
 (p-2)\qd^{\,p-1}-(p-1)\qd^{\,p}=\qd^{\,p-1}\bigl\{(p-2)-(p-1)\qd\bigr\}=0.
\]
This shows that $e_{\bd}(\infty)=\bd$, and \eqref{eq:energy-identity} gives $\ell_{\bd}=1+p\bd e_{\bd}(\infty)=1+p\bd^2$.
\end{proof}

The plateau and the energy in \cref{sel:selected-sector} are those predicted by Cugliandolo and Kurchan. In the normalization of \cite{ck1993}, the inverse temperature is $\sqrt2\,\beta$ and the energy is $-e/\sqrt2$; with this translation, \eqref{sel:selected-energy} is their Equation~(9), evaluated with the fluctuation--dissipation ratio $x=(p-2)(1-q)/q$ given after their Equation~(12).

\begin{remark}[The subsequent low-temperature limit]\label{rem:low-temperature}
As $\beta\to\infty$, we have $q_\beta\to1$, $\beta(1-q_\beta)\to[p(p-1)]^{-1/2}$, and
\[
 \lim_{\beta\to\infty}e_\beta(\infty)=2\sqrt{\frac{p-1}{p}}=\Eth.
\]
This agrees with the threshold energy of the pure $p$-spin landscape, computed by Crisanti and Sommers \cite[Equation~(9)]{cs1995} in the normalization of \cite{ck1993} and by Auffinger, Ben Arous, and \v{C}ern\'y \cite[Equation~(2.13)]{abc2013} in the present one. It also agrees with the results of Sellke \cite[Theorems~1.1 and~1.7]{sellke2024}, which show that, at low temperature, the energy of Langevin dynamics on a bounded time horizon comes close to $\Eth$. Here $\beta\to\infty$ is taken after the dimension limit and after both long-time limits, and the calculation says nothing about the gradient flow at zero temperature.

To verify these limits, recall that $q_\beta\ge\qd>0$. The marginality relation gives
\[
 (1-q_\beta)^2=\frac1{\beta^2p(p-1)q_\beta^{p-2}}\le\frac1{\beta^2p(p-1)\qd^{\,p-2}}\longrightarrow0,
\]
so $q_\beta\to1$. Taking square roots in the same relation, we find
\[
 \beta(1-q_\beta)=\frac{q_\beta^{-(p-2)/2}}{\sqrt{p(p-1)}}\longrightarrow\frac1{\sqrt{p(p-1)}}.
\]
Let $g(y)=1+(p-2)y^{p-1}-(p-1)y^p$, so that $e_\beta(\infty)=\beta g(q_\beta)$ by \eqref{sel:selected-energy}. We have $g(1)=0$ and
\[
 g'(y)=(p-1)y^{p-2}\bigl\{(p-2)-py\bigr\},\qquad g'(1)=-2(p-1).
\]
Let $\varepsilon=1-q_\beta$. By Taylor's theorem, there exists $\eta\in(q_\beta,1)$ such that
\[
 g(q_\beta)=g(1)-g'(1)\varepsilon+\tfrac12g''(\eta)\varepsilon^2
 =2(p-1)\varepsilon+\tfrac12g''(\eta)\varepsilon^2.
\]
On $[0,1]$, we have $|g''|\le K_p$ with $K_p=(p-1)\{(p-2)^2+p(p-1)\}$, and therefore
\[
 \bigl|e_\beta(\infty)-2(p-1)\beta(1-q_\beta)\bigr|\le\tfrac12K_p\,\beta(1-q_\beta)^2.
\]
As $\beta\to\infty$, the right side equals $\frac12K_p\cdot\beta(1-q_\beta)\cdot(1-q_\beta)$ and tends to zero, while $2(p-1)\beta(1-q_\beta)\to2(p-1)/\sqrt{p(p-1)}=\Eth$.
\end{remark}

\begin{remark}[The unselected intermediate interval]\label{sel:intermediate-range}
For $0<\beta<\beta_{\mathrm m}$, \cref{sel:all-temperature} gives $q=0$; then $\ell_\beta=1+\beta^2p$ and $e_\beta(\infty)=\beta$. For $\beta_{\mathrm m}\le\beta<\bd$, it proves convergence but allows either the zero branch or the positive branch $q_+(\beta)$. In the latter case, writing $q_+=q_+(\beta)$, we have
\[
 e_\beta(\infty)=\beta+
 \frac{q_+\{(p-2)-(p-1)q_+\}}{p\beta(p-1)(1-q_+)^2}>\beta.
\]
Indeed, the marginality relation can be written as $\beta^2pq_+^{p-1}=q_+/((p-1)(1-q_+)^2)$. Combining it with the formula for $\ell_\beta$ in \cref{sel:all-temperature}, we have
\[
 e_\beta(\infty)-\beta=\frac{\ell_\beta-1-\beta^2p}{p\beta}
 =\frac1{p\beta}\left\{\frac{q_+}{1-q_+}-\frac{q_+}{(p-1)(1-q_+)^2}\right\},
\]
and bringing the bracket to a common denominator gives the stated formula. For the sign, note that $\beta<\bd$ implies $G(q_+)=\beta^{-2}>G(\qd)$. As $G$ decreases strictly on $[q_{\mathrm m},1]$, we obtain $q_{\mathrm m}\le q_+<\qd$. In particular, $(p-2)-(p-1)q_+>0$, and the correction is positive. We do not decide which alternative occurs in this interval, and the transition is not classified from below there; see \cref{sec:discussion}.
\end{remark}

\section{Forgetting of fixed waiting times}\label{sec:forgetting}

\begin{theorem}[Forgetting of fixed waiting times]\label{sel:fixed-wait-forgetting}
Let $p\ge3$ and $\beta>0$, and let $(c,r)$ be the stationary profile of \cref{sel:all-temperature}. Then
\begin{equation}
 \lim_{t\to\infty}\sup_{0\le s\le S}
       \bigl\{C(t,s)+R(t,s)\bigr\}=0\qquad(0\le S<\infty).
 \label{sel:fixed-wait-decay}
\end{equation}
If the stationary plateau is zero, the stronger whole-row conclusion holds:
\begin{equation}
 \sup_{0\le s\le t}
 \bigl\{|C(t,s)-c(t-s)|+|R(t,s)-r(t-s)|\bigr\}
 \longrightarrow0\qquad(t\to\infty).
 \label{sel:zero-plateau-whole-row}
\end{equation}
\end{theorem}

\begin{proof}
The proof has four steps. We first reduce \eqref{sel:fixed-wait-decay} to a statement about $C$. When the plateau $q$ is positive, we verify the hypotheses of part~\ref{item:threshold-forgetting} of \cref{prop:threshold-inputs}. When $q=0$, we argue directly: the norm of the prefix center tends to zero when first $t\to\infty$ and then $L\to\infty$. Finally, we prove \eqref{sel:zero-plateau-whole-row} by splitting each row into a recent strip and a remote part.

\paragraph{Step 1: reduction to the correlation.}
Recall from \eqref{eq:response-derivative-bounds} that $0\le R(t,s)\le C(t,s)$ for $t\ge s\ge0$. For $t\ge S$, we have
\[
 \sup_{0\le s\le S}\bigl\{C(t,s)+R(t,s)\bigr\}\le2\sup_{0\le s\le S}C(t,s),
\]
and it suffices to show that $\sup_{0\le s\le S}C(t,s)\to0$ as $t\to\infty$.

\paragraph{Step 2: positive plateau.}
Suppose that $q>0$. Then the translated kernels converge along every sequence $T\to\infty$ to the single profile $(c,r)$ by \cref{sel:all-temperature}, and the plateau satisfies $\beta^2p(p-1)q^{p-2}(1-q)^2=1$ by \cref{sel:positive-marginality}. These are the hypotheses of part~\ref{item:threshold-forgetting} of \cref{prop:threshold-inputs}. Then \eqref{eq:threshold-input-forgetting} implies that $\sup_{0\le s\le S}|C(t,s)|\to0$.

\paragraph{Step 3: zero plateau.}
Suppose that $q=0$. Fix $L\ge0$ and let $t\ge L$. For $0\le s\le t-L$, the vector $v_s$ lies in $\mathcal H_s\subseteq\mathcal H_{t-L}$ by \cref{prop:geometry-tree-realization}, so $P_{t-L}v_s=v_s$. As $P_{t-L}$ is an orthogonal projection, we have
\[
 C(t,s)=\ip{v_t}{P_{t-L}v_s}=\ip{P_{t-L}v_t}{v_s}=\ip{M_{t-L}(L)}{v_s}.
\]
By the Cauchy--Schwarz inequality and $\norm{v_s}=1$, we obtain
\[
 \sup_{0\le s\le t-L}C(t,s)^2\le\norm{M_{t-L}(L)}^2.
\]
We now let $t\to\infty$ with $L$ fixed. By \cref{sel:all-temperature}, the translated kernels converge to the stationary cluster $(c,r,\ell_\beta)$ along every sequence $T_n\to\infty$. Then \cref{sel:prefix-reflection} applies along every such sequence, and \eqref{sel:prefix-norm} gives $\norm{M_{T_n}(L)}^2\to c(2L)$. As the sequence is arbitrary, we have $\norm{M_T(L)}^2\to c(2L)$ as $T\to\infty$. Setting $T=t-L$, we obtain
\begin{equation}
 \limsup_{t\to\infty}\sup_{0\le s\le t-L}C(t,s)^2\le c(2L).
 \label{sel:zero-plateau-remote-bound}
\end{equation}
This bound holds for every plateau, but it is useful only when $q=0$. Fix $0\le S<\infty$. For $t\ge S+L$, we have $[0,S]\subseteq[0,t-L]$, and \eqref{sel:zero-plateau-remote-bound} implies $\limsup_{t\to\infty}\sup_{0\le s\le S}C(t,s)^2\le c(2L)$ for every $L\ge0$. Letting $L\to\infty$ and using $c(2L)\to q=0$, we find that $\sup_{0\le s\le S}C(t,s)\to0$. Together with Steps~1 and~2, this proves \eqref{sel:fixed-wait-decay}.

\paragraph{Step 4: the whole row when $q=0$.}
Fix $L>0$, and let $t\ge L$. We split $[0,t]$ into the recent strip $[t-L,t]$ and the remote part $[0,t-L]$. On the recent strip, we use \eqref{eq:all-temperature-window} with $T=t-L$ and $A=L$. Writing $s=T+u$ with $u\in[0,L]$, we have $C(t,s)=C(T+L,T+u)$ and $c(t-s)=c(L-u)$, and similarly for $R$ and $r$. We obtain
\[
 \sup_{t-L\le s\le t}
 \bigl\{|C(t,s)-c(t-s)|+|R(t,s)-r(t-s)|\bigr\}\longrightarrow0
 \qquad(t\to\infty).
\]
On the remote part, we have $t-s\ge L$. Recall that $0\le R\le C$, that $0\le r\le c$, and that $c$ is nonincreasing. For $0\le s\le t-L$, we find
\[
\begin{aligned}
 |C(t,s)-c(t-s)|+|R(t,s)-r(t-s)|
 &\le C(t,s)+c(t-s)+R(t,s)+r(t-s)\\
 &\le2C(t,s)+2c(L).
\end{aligned}
\]
By \eqref{sel:zero-plateau-remote-bound}, the upper limit as $t\to\infty$ of the supremum of the right side over $0\le s\le t-L$ is at most $2c(2L)^{1/2}+2c(L)$. Combining the two parts, we have
\[
 \limsup_{t\to\infty}\sup_{0\le s\le t}
 \bigl\{|C(t,s)-c(t-s)|+|R(t,s)-r(t-s)|\bigr\}
 \le2c(2L)^{1/2}+2c(L),
\]
where $t\to\infty$ is taken with $L$ fixed. Since $q=0$, the right side tends to zero as $L\to\infty$. This proves \eqref{sel:zero-plateau-whole-row}.
\end{proof}

\section{Correlation aging}\label{sec:aging}

\begin{corollary}[Correlation aging in the ordered limits]\label{sel:aging}
For every $p\ge3$ and $\beta\ge\bd$,
\begin{equation}
 \lim_{\tau\to\infty}\lim_{s\to\infty}C(s+\tau,s)=q_\beta>0,
 \qquad
 \lim_{s\to\infty}\lim_{\tau\to\infty}C(s+\tau,s)=0.
 \label{sel:aging-ordered-limits}
\end{equation}
\end{corollary}

\begin{proof}
Fix $\tau\ge0$. Applying \eqref{sel:stationary-window} with $A=\tau$ and $T=s$ to the pair of offsets $(\tau,0)$, we find that $\lim_{s\to\infty}C(s+\tau,s)=c(\tau)$. As $\tau\to\infty$, we have $c(\tau)\to q_\beta$ by \cref{sel:selected-sector}, and $q_\beta\ge\qd>0$. This gives the first limit. Next, fix $s\ge0$. Applying \eqref{sel:fixed-wait-decay} with $S=s$ and $t=s+\tau$, we find that $C(s+\tau,s)\to0$ as $\tau\to\infty$. The inner limit is zero for every $s$, and so is the outer limit.
\end{proof}

The first inner limit in \eqref{sel:aging-ordered-limits} is the stationary correlation $c(\tau)$, and the second inner limit vanishes at every fixed waiting time. \Cref{sel:aging} concerns only these two ordered limits. It does not identify the time scale on which $C(s+\tau,s)$ decreases below $q_\beta$, and it does not give a fluctuation--dissipation relation in the slow regime. We discuss these questions in \cref{sec:discussion}.

\paragraph{Proof of \cref{thm:main-aging}.}
We collect the results of this section, which rely on \cref{prop:threshold-inputs}, proved in \cref{sec:tree-ceilings,app:general-selection}. Fix $p\ge3$ and $\beta>0$. By \cref{sel:all-temperature}, the multiplier and the energy converge to $\ell_\beta$ and $e_\beta(\infty)$, and the translated kernels converge locally uniformly to the stationary profile $(c_\beta,r_\beta)=(c,r)$. This is \eqref{eq:main-stationary}. In Step~3 of the proof of \cref{sel:all-temperature}, we showed that every stationary cluster equals $(c,r,\ell_\beta)$, so the profile is unique. By \cref{prop:stationary-clusters}, we have $c(0)=r(0)=1$ and $r=-c'$, and $c$ is completely monotone by \cref{sel:cluster-complete-monotonicity}. The forgetting statement \eqref{eq:main-forgetting} is \eqref{sel:fixed-wait-decay}. For $\beta\ge\bd$, \cref{sel:selected-sector} shows that the plateau is $q_\beta\ge\qd$, which is the larger solution of \eqref{eq:main-marginality}, and it gives \eqref{eq:main-energy} and \eqref{eq:main-multiplier}. Finally, \eqref{eq:main-aging} is \eqref{sel:aging-ordered-limits}. This completes the proof of \cref{thm:main-aging}.

\chapter{Tree ceilings and the threshold inputs for cubic and quartic models}\label{sec:tree-ceilings}

In this section, we prove \cref{prop:threshold-inputs} for $p\in\{3,4\}$. That proposition is stated in \cref{sec:threshold-inputs}, and \cref{sec:stationary-selection} deduces from it the selection of the plateau, forgetting, and correlation aging. We prove the proposition for every $p\ge3$ in \cref{app:general-selection}.

Both parts rest on a deterministic inequality in the tree space of \cref{sec:tree-space}. For $p=3$, the cubic tree ceiling, \cref{thm:geometry-tree-ceiling}, states that every unit tree vector has energy strictly below $\Eth$, and that a unit vector with energy close to $\Eth$ has little mass at every fixed finite set of heights. \Cref{thm:geometry-quartic-tree-ceiling} is the analogous statement for $p=4$, and \cref{cor:tree-ceiling-energy} restates both in terms of the tree energy. In \cref{sec:tree-input-ceiling}, we prove part~\ref{item:threshold-ceiling} of \cref{prop:threshold-inputs} by applying a ceiling to normalized prefix centers and passing to the limit with \cref{sel:prefix-reflection}. Under the hypotheses of part~\ref{item:threshold-forgetting}, the normalized centers approach the threshold; this is \cref{sel:threshold-centers}. Their mass then leaves every fixed set of heights, while an earlier state is approximated by finitely many heights, uniformly on a bounded time interval. Combining these two facts, we prove part~\ref{item:threshold-forgetting} in \cref{sec:tree-input-forgetting}.

Besides the tree space and the realization of \cref{prop:geometry-tree-realization}, we use the reflection identity of \cref{sel:prefix-reflection}, the precompactness in \cref{prop:stationary-clusters}, the plateau identity \eqref{eq:plateau-multiplier}, and the limiting center energy of \cref{lem:center-energy}. The last lemma is an algebraic statement about a stationary cluster and does not use \cref{prop:threshold-inputs}. Hence the arguments of this section and of \cref{sec:stationary-selection} are not circular. The ceilings are used again later: the participation bound enters the proof of \cref{resp:strong-row}, and the retained angular and Young defects \eqref{eq:geometry-angular-young-defects} are used in \cref{sec:geometry-results}. \Cref{sel:threshold-centers} is used in \cref{resp:chapter,sec:three-time}.

The value $\Eth=2\sqrt{(p-1)/p}$ is the threshold energy of the pure $p$-spin model. It appears in the complexity of critical points computed by Auffinger, Ben Arous, and \v{C}ern\'y \cite[equation~(2.13)]{abc2013}, where the Hamiltonian is minimized and the threshold is $-\Eth$. It is the threshold for Lipschitz algorithms identified by Huang and Sellke, for even models in \cite{hs2021} and including odd pure models in \cite[Theorem~1 and Corollary~1.8]{hs2023}, and Sellke proved that low-temperature Langevin dynamics approaches it when the temperature is also sent to zero \cite[Theorems~1.1 and~1.7]{sellke2024}. The ceilings below are elementary inequalities in the tree space, and we do not use the cited results in the proofs of this section. They also differ from the branching constructions of the algorithmic literature, such as those of Subag \cite[Remark~6]{subag2021} and of Huang and Sellke \cite[Remarks~1.2--1.3 and Sections~3 and~7]{hs2021}, which concern configurations in finite dimension. Here we retain quantitative height defects inside the causal covariance realization and use their saturation to prove escape from every fixed set of heights.

\paragraph{Standing notation.}
Throughout this section, $p\in\{3,4\}$ and $m=p-1$, and $\mathcal H$, $\mathcal I$, $\Pi_j$, and $P_T$ are the tree space, the grafting isometry, and the height and prefix projections of \cref{sec:tree-space} for this arity. For $v\in\mathcal H$ and $j\ge0$, we write
\[
 b_j=\norm{\Pi_jv}^2,\qquad a_j=\sqrt{b_j},\qquad B_j=\sum_{i\le j}b_i,
\]
with $b_j=B_j=0$ for $j<0$. Recall that $\sum_j\Pi_j=\Id$, so that the masses $b_j$ sum to $\norm v^2$ and $B_j$ increases to $\norm v^2$. The tree energy is $h(v)=\sqrt p\,\ip v{\mathcal I(v^{\otimes m})}$, as in \eqref{eq:geometry-tree-energy}. It is a homogeneous polynomial of degree $p$, that is, $h(\alpha v)=\alpha^ph(v)$ for $\alpha\ge0$.

\section{The cubic tree ceiling}\label{sec:tree-cubic-ceiling}

We begin with a computation that holds for both arities. It expresses the height profile of a graft through the cumulative masses $B_j$.

\begin{lemma}[Heights of a graft]\label{lem:tree-graft-heights}
Let $v\in\mathcal H$. Then $\Pi_0\mathcal I(v^{\otimes m})=0$, and for every $j\ge1$,
\begin{equation}\label{eq:tree-graft-heights}
 \Pi_j\mathcal I(v^{\otimes m})=\Pi_j\mathcal I\bigl((\Pi_{\le j-1}v)^{\otimes m}\bigr),\qquad
 \norm{\Pi_j\mathcal I(v^{\otimes m})}^2=B_{j-1}^m-B_{j-2}^m.
\end{equation}
\end{lemma}

\begin{proof}
The ranges of the projections $\Pi_i$, $i\ge0$, are mutually orthogonal and span $\mathcal H$. It follows that $\mathcal H^{\otimes m}$ is the orthogonal sum of the subspaces $\Pi_{i_1}\mathcal H\otimes\dots\otimes\Pi_{i_m}\mathcal H$ over $(i_1,\dots,i_m)$. Because a graft has height one plus the largest height of its root subtrees, the map $\mathcal I$ sends the subspace indexed by $(i_1,\dots,i_m)$ into the range of $\Pi_j$ with $j=1+\max_ki_k$. In particular, $\Pi_0\mathcal I=0$, and for $j\ge1$,
\[
 \Pi_j\mathcal I(v^{\otimes m})=\sum_{\max_ki_k=j-1}\mathcal I\bigl(\Pi_{i_1}v\otimes\dots\otimes\Pi_{i_m}v\bigr).
\]
Only the components $\Pi_iv$ with $i\le j-1$ appear, which gives the first identity in \eqref{eq:tree-graft-heights}. The terms of the sum are orthogonal, since $\mathcal I$ is an isometry and the subspaces are orthogonal, and the term indexed by $(i_1,\dots,i_m)$ has squared norm $b_{i_1}\cdots b_{i_m}$. Summing over the $m$-tuples with largest entry $j-1$, we obtain
\[
 \norm{\Pi_j\mathcal I(v^{\otimes m})}^2=\Bigl(\sum_{i\le j-1}b_i\Bigr)^m-\Bigl(\sum_{i\le j-2}b_i\Bigr)^m=B_{j-1}^m-B_{j-2}^m,
\]
as desired.
\end{proof}

We now treat the cubic case. The next estimate is the reason that a normalized prefix center selects marginality. It is an inequality in the tree space, and it involves no finite-dimensional critical points.

\begin{theorem}[The cubic tree ceiling and its retained defects]
\label{thm:geometry-tree-ceiling}
Let $p=3$, so that grafting is binary, and let $F_*=2\sqrt2/3$. For a unit vector $v\in\mathcal H$, define
\[
 f(v)=\ip{v}{\mathcal I(v\otimes v)},\qquad \delta=F_*-f(v),
\]
and let $b_j=\norm{\Pi_jv}^2$, $a_j=\sqrt{b_j}$, and $B_j=\sum_{i\le j}b_i$ for $j\ge0$, with $B_{-1}=B_{-2}=0$. Then $|f(v)|<F_*$, and
\begin{equation}
 \sum_{j\ge0}b_j^2\le4\sqrt2\,\delta,
 \qquad
 \norm{\Pi_{\le H}v}^2\le\sqrt{4\sqrt2\,(H+1)\,\delta}\quad(H\ge0).
 \label{eq:geometry-height-defect}
\end{equation}
The supremum of $f$ over unit vectors equals $F_*$, and it is not attained. Further, let $s_j=\sqrt{B_{j-1}+B_{j-2}}$, and for $j\ge1$ let $\theta_j\in[0,\pi]$ be the angle between $\Pi_jv$ and $\Pi_j\mathcal I(v\otimes v)$, with $\theta_j=0$ when one of these two vectors is zero. Then
\begin{equation}
 \sum_{j\ge1}a_ja_{j-1}s_j(1-\cos\theta_j)\le\delta,
 \qquad
 \frac12\sum_{j\ge1}s_j(a_j-a_{j-1})^2\le\delta.
 \label{eq:geometry-angular-young-defects}
\end{equation}
\end{theorem}

\begin{proof}
The proof has three steps. First, we write $f(v)$ as a sum over heights and bound it by two sums $S$ and $Y$, keeping track of the gaps. Second, we compare $Y$ with the integral $\int_0^1\sqrt{2y}\,\dd y=F_*$, again keeping the gap. Third, we construct unit vectors whose values of $f$ approach $F_*$.

\paragraph{Step 1: heightwise Cauchy--Schwarz and Young inequalities.}
By \cref{lem:tree-graft-heights} with $m=2$, we have, for $j\ge1$,
\[
 \norm{\Pi_j\mathcal I(v\otimes v)}^2=B_{j-1}^2-B_{j-2}^2=b_{j-1}(B_{j-1}+B_{j-2})=a_{j-1}^2s_j^2.
\]
Using that the height projections are orthogonal and that $\Pi_0\mathcal I=0$, we find that
\begin{equation}\label{eq:tree-cubic-angles}
 f(v)=\sum_{j\ge1}\ip{\Pi_jv}{\Pi_j\mathcal I(v\otimes v)}=\sum_{j\ge1}a_ja_{j-1}s_j\cos\theta_j.
\end{equation}
When one of the two vectors in the $j$-th inner product is zero, the coefficient $a_ja_{j-1}s_j$, which is the product of their norms, vanishes. The convention on $\theta_j$ therefore does not matter. We define
\[
 S=\sum_{j\ge1}a_ja_{j-1}s_j,\qquad Y=\frac12\sum_{j\ge1}(b_j+b_{j-1})s_j.
\]
All these series converge absolutely, because $s_j\le\sqrt2$ and $\sum_jb_j=1$. Using \eqref{eq:tree-cubic-angles} and the inequality $a_ja_{j-1}\le\frac12(b_j+b_{j-1})$, we see that $|f(v)|\le S\le Y$, with the exact gaps
\begin{equation}\label{eq:tree-cubic-gaps}
 S-f(v)=\sum_{j\ge1}a_ja_{j-1}s_j(1-\cos\theta_j),\qquad
 Y-S=\frac12\sum_{j\ge1}s_j(a_j-a_{j-1})^2.
\end{equation}

\paragraph{Step 2: comparison with an integral.}
Let $g(y)=\sqrt{2y}$, a concave function on $[0,1]$ with $\int_0^1g(y)\,\dd y=F_*$. Reindexing the second half of $Y$ and using $s_1=\sqrt{b_0}$, we have
\[
 Y=\frac12b_0^{3/2}+\frac12\sum_{j\ge1}b_j(s_j+s_{j+1}).
\]
Since $B_j$ increases to one, we may split $\int_0^1g$ at the points $B_0\le B_1\le\dots$. Using $\int_0^{B_0}g=F_*b_0^{3/2}$, we find that
\begin{equation}\label{eq:tree-cubic-integral-gap}
 F_*-Y=\Bigl(F_*-\frac12\Bigr)b_0^{3/2}
 +\sum_{j\ge1}\Bigl[\int_{B_{j-1}}^{B_j}g(y)\,\dd y-\frac{b_j}{2}(s_j+s_{j+1})\Bigr].
\end{equation}
We bound each bracket from below. Because $g$ is concave, the trapezoid rule underestimates its integral, that is,
\[
 \int_{B_{j-1}}^{B_j}g(y)\,\dd y\ge\frac{b_j}2\bigl\{g(B_{j-1})+g(B_j)\bigr\}.
\]
We have $s_j\le g(B_{j-1})$, because $B_{j-2}\le B_{j-1}$. If $b_j>0$, then
\[
 g(B_j)-s_{j+1}=\sqrt{2B_j}-\sqrt{B_j+B_{j-1}}
 =\frac{b_j}{\sqrt{2B_j}+\sqrt{B_j+B_{j-1}}}\ge\frac{b_j}{2\sqrt2},
\]
because the denominator is at most $2\sqrt2$. Inserting these bounds, we see that the bracket with index $j$ is at least $\frac{b_j}2\cdot\frac{b_j}{2\sqrt2}=b_j^2/(4\sqrt2)$, and this also holds when $b_j=0$. For the first term, we use $b_0^{3/2}\ge b_0^2$ and
\[
 F_*-\frac12-\frac1{4\sqrt2}=\frac{13\sqrt2-12}{24}>0.
\]
Combining these bounds with \eqref{eq:tree-cubic-integral-gap}, we obtain
\begin{equation}\label{eq:tree-cubic-participation-gap}
 F_*-Y\ge\frac1{4\sqrt2}\sum_{j\ge0}b_j^2.
\end{equation}

We now collect the gaps. We have the exact decomposition
\[
 \delta=F_*-f(v)=(F_*-Y)+(Y-S)+\bigl(S-f(v)\bigr),
\]
in which all three terms are nonnegative. By \eqref{eq:tree-cubic-participation-gap}, the first term is at least $(4\sqrt2)^{-1}\sum_jb_j^2$, which yields the first inequality in \eqref{eq:geometry-height-defect}. The second and third terms are the two sums in \eqref{eq:geometry-angular-young-defects} by \eqref{eq:tree-cubic-gaps}, which proves \eqref{eq:geometry-angular-young-defects}. Using the Cauchy--Schwarz inequality, we also have
\[
 \norm{\Pi_{\le H}v}^2=\sum_{j=0}^Hb_j\le\Bigl((H+1)\sum_{j\ge0}b_j^2\Bigr)^{1/2}\le\sqrt{4\sqrt2\,(H+1)\,\delta},
\]
which is the second inequality in \eqref{eq:geometry-height-defect}. Finally, some $b_j$ is positive, because $\sum_jb_j=1$, and we have $|f(v)|\le Y\le F_*-(4\sqrt2)^{-1}\sum_jb_j^2<F_*$. In particular, no unit vector attains the value $F_*$.

\paragraph{Step 3: sharpness.}
We construct unit vectors $v^{(n)}$ with $f(v^{(n)})\to F_*$. We give the construction for a general arity $m\in\{2,3\}$, since the proof of \cref{thm:geometry-quartic-tree-ceiling} uses it with $m=3$. Fix an integer $n\ge1$ and let $\rho=(n+1)^{-1}$. Let $w_0=\sqrt\rho\,e_0$, a vector in the one-leaf summand. For $j=1,\dots,n$, we define inductively
\[
 \gamma_j=\Pi_j\mathcal I\bigl(w_{j-1}^{\otimes m}\bigr),\qquad u_j=\sqrt\rho\,\frac{\gamma_j}{\norm{\gamma_j}},\qquad w_j=w_{j-1}+u_j.
\]
We check by induction that this is well defined, and that $w_{j-1}$ has mass $\rho$ at each height $0,\dots,j-1$ and no mass at larger heights. This holds for $w_0$. If it holds for $w_{j-1}$, then \cref{lem:tree-graft-heights} applied to $w_{j-1}$ gives
\[
 \norm{\gamma_j}^2=(j\rho)^m-\bigl((j-1)\rho\bigr)^m>0,
\]
and $u_j$ is well defined. The vector $u_j$ lies in the range of $\Pi_j$ and has squared norm $\rho$, and therefore $w_j$ has mass $\rho$ at each height $0,\dots,j$ and no mass at larger heights. Let $v^{(n)}=w_n$. It is a unit vector, with $b_j=\rho$ for $0\le j\le n$ and $b_j=0$ for $j>n$.

We now show that all the angles vanish. Let $1\le j\le n$. Since $\Pi_{\le j-1}v^{(n)}=w_{j-1}$, \cref{lem:tree-graft-heights} implies that $\Pi_j\mathcal I((v^{(n)})^{\otimes m})=\gamma_j$. The component $\Pi_jv^{(n)}=u_j$ is a positive multiple of $\gamma_j$, which means that $\theta_j=0$. For $j>n$, the component $\Pi_jv^{(n)}$ vanishes. Only the heights $1,\dots,n$ contribute to $f$, and we find that
\begin{equation}\label{eq:tree-sharp-value}
 f\bigl(v^{(n)}\bigr)=\sum_{j=1}^n\ip{u_j}{\gamma_j}=\sum_{j=1}^n\sqrt\rho\,\norm{\gamma_j}
 =\sum_{j=1}^n\sqrt\rho\,\bigl\{(j\rho)^m-\bigl((j-1)\rho\bigr)^m\bigr\}^{1/2}.
\end{equation}
For $m=2$, each summand equals $\rho\sqrt{(2j-1)\rho}=\rho\,g\bigl((j-\frac12)\rho\bigr)$. In other words, $f(v^{(n)})$ is the midpoint Riemann sum for $\int g$ over the first $n$ of the $n+1$ intervals $[(j-1)\rho,j\rho]$ that partition $[0,1]$. Since $g$ is continuous on $[0,1]$ and bounded by $\sqrt2$, the omitted last interval contributes at most $\sqrt2\rho$, and we obtain
\[
 \lim_{n\to\infty}f\bigl(v^{(n)}\bigr)=\int_0^1\sqrt{2y}\,\dd y=F_*.
\]
Together with Step 2, this shows that the supremum of $f$ equals $F_*$ and is not attained.
\end{proof}

\section{The quartic tree ceiling}\label{sec:tree-quartic-ceiling}

For $p=4$, grafting is ternary. The proof follows the cubic one, with a difference of cubes of the cumulative masses in place of a difference of squares. The comparison function becomes linear, and the ceiling calculation telescopes.

\begin{lemma}[The quartic tree ceiling]
\label{thm:geometry-quartic-tree-ceiling}
Let $p=4$, so that grafting has arity three. For a unit vector $v\in\mathcal H$, define $f(v)=\ip{v}{\mathcal I(v^{\otimes3})}$. Then
\begin{equation}\label{eq:tree-quartic-ceiling}
 |f(v)|<\frac{\sqrt3}{2},\qquad
 \sum_{j\ge0}\norm{\Pi_jv}^4\le2\sqrt3\Bigl(\frac{\sqrt3}{2}-f(v)\Bigr).
\end{equation}
The supremum of $f$ over unit vectors equals $\sqrt3/2$, and it is not attained. In particular, let $h(v)=2f(v)$ be the tree energy and $\Eth=\sqrt3$. If $(v_n)$ are unit vectors with $h(v_n)\to\Eth$, then $\norm{\Pi_{\le H}v_n}\to0$ for every $H\ge0$.
\end{lemma}

\begin{proof}
The proof has three steps, parallel to those of \cref{thm:geometry-tree-ceiling}: a heightwise bound, a comparison that keeps the gap, and the sharpness construction.

\paragraph{Step 1: heightwise bound.}
In this proof, we let
\[
 s_j=\sqrt{B_{j-1}^2+B_{j-1}B_{j-2}+B_{j-2}^2},\qquad j\ge1.
\]
By \cref{lem:tree-graft-heights} with $m=3$, we have
\[
 \norm{\Pi_j\mathcal I(v^{\otimes3})}^2=B_{j-1}^3-B_{j-2}^3=b_{j-1}s_j^2=a_{j-1}^2s_j^2.
\]
As in Step 1 of the proof of \cref{thm:geometry-tree-ceiling}, using the Cauchy--Schwarz inequality at each height and the inequality $a_ja_{j-1}\le\frac12(b_j+b_{j-1})$, we have
\[
 |f(v)|\le S\le Y,\qquad S=\sum_{j\ge1}a_ja_{j-1}s_j,\qquad Y=\frac12\sum_{j\ge1}(b_j+b_{j-1})s_j.
\]
Reindexing the second half of $Y$ and using $s_1=B_0=b_0$, we have
\[
 Y=\frac12b_0^2+\frac12\sum_{j\ge1}b_j(s_j+s_{j+1}).
\]

\paragraph{Step 2: the comparison.}
Note that $b_j(B_j+B_{j-1})=B_j^2-B_{j-1}^2$. As $B_j$ increases to one, the sum $\sum_{j\ge0}b_j(B_j+B_{j-1})$ telescopes to one, and we can write
\[
 \frac{\sqrt3}2=\frac{\sqrt3}2\,b_0^2+\frac12\sum_{j\ge1}b_j\bigl(\sqrt3B_{j-1}+\sqrt3B_j\bigr),
\]
and subtracting the expression for $Y$, we arrive at
\begin{equation}\label{eq:tree-quartic-gap}
 \frac{\sqrt3}2-Y=\frac{\sqrt3-1}2\,b_0^2
 +\frac12\sum_{j\ge1}b_j\Bigl[\bigl(\sqrt3B_{j-1}-s_j\bigr)+\bigl(\sqrt3B_j-s_{j+1}\bigr)\Bigr].
\end{equation}
We have $s_j\le\sqrt3B_{j-1}$, because $B_{j-2}\le B_{j-1}$. To bound the second difference, let $X\ge X'\ge0$. Then
\[
 \frac{(2X+X')^2}{3}-\bigl(X^2+XX'+X'^2\bigr)=\frac{(X-X')(X+2X')}{3}\ge0.
\]
This implies that $\sqrt{X^2+XX'+X'^2}\le(2X+X')/\sqrt3$, and therefore
\[
 \sqrt3X-\sqrt{X^2+XX'+X'^2}\ge\sqrt3X-\frac{2X+X'}{\sqrt3}=\frac{X-X'}{\sqrt3}.
\]
Taking $X=B_j$ and $X'=B_{j-1}$, we find that $\sqrt3B_j-s_{j+1}\ge b_j/\sqrt3$, and the summand with index $j\ge1$ in \eqref{eq:tree-quartic-gap} is at least $b_j^2/(2\sqrt3)$. The first term is also at least $b_0^2/(2\sqrt3)$, because
\[
 \frac{\sqrt3-1}2-\frac1{2\sqrt3}=\frac{2-\sqrt3}{2\sqrt3}>0.
\]
We conclude that
\[
 \frac{\sqrt3}2-Y\ge\frac1{2\sqrt3}\sum_{j\ge0}b_j^2.
\]
Since $f(v)\le Y$, this yields the second inequality in \eqref{eq:tree-quartic-ceiling}. Further, $|f(v)|\le Y$ and $\sum_jb_j^2>0$ for a unit vector, which gives $|f(v)|<\sqrt3/2$. This also shows that the value $\sqrt3/2$ is not attained.

\paragraph{Step 3: sharpness and finite heights.}
Let $v^{(n)}$ be the unit vectors constructed in Step 3 of the proof of \cref{thm:geometry-tree-ceiling}, now with $m=3$, and let $\rho=(n+1)^{-1}$. Inserting $m=3$ into \eqref{eq:tree-sharp-value}, we obtain
\[
 f\bigl(v^{(n)}\bigr)=\sum_{j=1}^n\sqrt\rho\,\bigl\{\rho^3\bigl(j^3-(j-1)^3\bigr)\bigr\}^{1/2}
 =\rho^2\sum_{j=1}^n\sqrt{3j^2-3j+1}.
\]
For $j\ge1$, we have $3(j-1)^2\le3j^2-3j+1\le3j^2$, so that
\[
 \frac{\sqrt3}2\cdot\frac{n(n-1)}{(n+1)^2}\le f\bigl(v^{(n)}\bigr)\le\frac{\sqrt3}2\cdot\frac{n}{n+1}.
\]
Both bounds tend to $\sqrt3/2$, which proves that the supremum of $f$ equals $\sqrt3/2$.

Finally, since $h=2f$, the second inequality in \eqref{eq:tree-quartic-ceiling} reads $\sum_jb_j^2\le\sqrt3\{\Eth-h(v)\}$. For $H\ge0$, the Cauchy--Schwarz inequality gives
\[
 \norm{\Pi_{\le H}v}^2=\sum_{j=0}^Hb_j\le\Bigl((H+1)\sum_{j\ge0}b_j^2\Bigr)^{1/2}\le\bigl\{\sqrt3\,(H+1)\bigl(\Eth-h(v)\bigr)\bigr\}^{1/2},
\]
and the right side tends to zero along any sequence of unit vectors with $h(v_n)\to\Eth$.
\end{proof}

For later use, we restate the two ceilings in terms of the tree energy $h$ and the threshold $\Eth$.

\begin{corollary}[The ceilings in terms of the tree energy]\label{cor:tree-ceiling-energy}
Let $p\in\{3,4\}$, and let $\kappa_3=4\sqrt{2/3}$ and $\kappa_4=\sqrt3$. Every unit vector $v\in\mathcal H$ satisfies $|h(v)|<\Eth$ and
\begin{equation}\label{eq:tree-ceiling-energy}
 \sum_{j\ge0}\norm{\Pi_jv}^4\le\kappa_p\bigl(\Eth-h(v)\bigr),\qquad
 \norm{\Pi_{\le H}v}^2\le\bigl\{\kappa_p(H+1)\bigl(\Eth-h(v)\bigr)\bigr\}^{1/2}\quad(H\ge0).
\end{equation}
The supremum of $h$ over unit vectors equals $\Eth$.
\end{corollary}

\begin{proof}
For $p=3$, we have $h=\sqrt3f$ and $\Eth=2\sqrt{2/3}=\sqrt3F_*$, so that $F_*-f(v)=\{\Eth-h(v)\}/\sqrt3$. Then \cref{thm:geometry-tree-ceiling} implies that $|h(v)|<\Eth$ and that $\sum_jb_j^2\le(4\sqrt2/\sqrt3)\{\Eth-h(v)\}=\kappa_3\{\Eth-h(v)\}$. For $p=4$, we have $h=2f$ and $\Eth=\sqrt3$, so that $\sqrt3/2-f(v)=\{\Eth-h(v)\}/2$. Then \cref{thm:geometry-quartic-tree-ceiling} implies that $|h(v)|<\Eth$ and that $\sum_jb_j^2\le\sqrt3\{\Eth-h(v)\}=\kappa_4\{\Eth-h(v)\}$. In both cases, the second inequality in \eqref{eq:tree-ceiling-energy} follows from the first by the Cauchy--Schwarz inequality, as in the two proofs above, and the supremum statement follows from those of \cref{thm:geometry-tree-ceiling,thm:geometry-quartic-tree-ceiling}.
\end{proof}

\section{The energy ceiling for prefix centers}\label{sec:tree-input-ceiling}

We now prove the first threshold input. Recall the prefix centers $M_T(a)=P_Tv_{T+a}$ of \cref{sec:selection}.

\begin{proof}[Proof of part~\ref{item:threshold-ceiling} of \cref{prop:threshold-inputs} for $p\in\{3,4\}$]
Let $T_n\to\infty$ be a sequence along which the translated kernels converge to a stationary cluster $(c,r,\ell)$ with plateau $q>0$, and fix $L\ge0$. We abbreviate $M_n=M_{T_n}(L)$. By \eqref{sel:prefix-norm}, we have $\norm{M_n}^2\to d_L=c(2L)$, and $d_L\ge q>0$ because $c$ is nonincreasing with limit $q$. Then $M_n\ne0$ for all large $n$, and $U_n=M_n/\norm{M_n}$ is a unit vector. \Cref{cor:tree-ceiling-energy} implies that $h(U_n)<\Eth$. Since $h$ is homogeneous of degree $p$, we have
\[
 h(M_n)=\norm{M_n}^p\,h(U_n)<\Eth\norm{M_n}^p.
\]
Letting $n\to\infty$ and using \eqref{sel:prefix-energy} on the left side and \eqref{sel:prefix-norm} on the right side, we obtain $h_L\le\Eth d_L^{p/2}$, which is \eqref{eq:threshold-input-ceiling}.
\end{proof}

The proof uses only the limits of \cref{sel:prefix-reflection} along the given sequence. In particular, it uses neither marginality nor convergence of the whole translated family, which \cref{sec:stationary-selection} deduces from part~\ref{item:threshold-ceiling}.

\section{Threshold-saturating prefix centers}\label{sec:tree-saturated-centers}

Under the hypotheses of part~\ref{item:threshold-forgetting} of \cref{prop:threshold-inputs}, the upper bound of part~\ref{item:threshold-ceiling} becomes an equality in the limit of large offsets. The next proposition states this and converts it into a bound on the mass of the centers at low heights. It is the input to the proof of part~\ref{item:threshold-forgetting} below, and we use it again in \cref{resp:chapter,sec:three-time}.

Suppose that the translated kernels converge along every sequence $T\to\infty$ to one stationary profile $(c,r)$ with plateau $q$. Every stationary cluster is then of the form $(c,r,\ell)$, and \eqref{eq:plateau-multiplier} determines its multiplier as $\ell=(1-q)^{-1}+\beta^2p(1-q^{p-1})$. The quantities $d_L$ and $h_L$ in \eqref{sel:prefix-norm} and \eqref{sel:prefix-energy} then do not depend on the sequence, and the energy $e_\infty=(\ell-1)/(p\beta)$ in \eqref{sel:prefix-energy} is the limit of the energy $e_\beta(t)$; see Step 1 of the proof below.

\begin{proposition}[Threshold-saturating prefix centers]
\label{sel:threshold-centers}
Let $p\in\{3,4\}$ and $\beta>0$. Suppose that the translated kernels converge along every sequence $T\to\infty$ to one stationary profile $(c,r)$, and that its plateau $q$ is positive and satisfies $\beta^2p(p-1)q^{p-2}(1-q)^2=1$. Let $d_L$ and $h_L$ be defined by \eqref{sel:prefix-norm} and \eqref{sel:prefix-energy} for this profile. Then
\begin{equation}
 \delta_L=\Eth-\frac{h_L}{d_L^{p/2}}\ge0\quad(L\ge0),
 \qquad \lim_{L\to\infty}\delta_L=0.
 \label{sel:center-deficit}
\end{equation}
Moreover, for every $L\ge0$ and every integer $H\ge0$,
\begin{equation}
 \limsup_{t\to\infty}
 \norm{\Pi_{\le H}P_{t-L}v_t}^2
 \le d_L\bigl\{\kappa_p(H+1)\delta_L\bigr\}^{1/2},
 \label{sel:low-height-center-bound}
\end{equation}
where $\kappa_3=4\sqrt{2/3}$ and $\kappa_4=\sqrt3$ are the constants of \cref{cor:tree-ceiling-energy}.
\end{proposition}

The normalized prefix center approaches the threshold; the proposition makes no threshold assertion about the state $v_t$ itself.

\begin{proof}
The proof has three steps. We first show that the limits of \cref{sel:prefix-reflection} hold as $T\to\infty$ without passing to a sequence. Then we prove \eqref{sel:center-deficit} from part~\ref{item:threshold-ceiling} of \cref{prop:threshold-inputs} and the identity \eqref{sel:center-amgm}. Finally, we apply \cref{cor:tree-ceiling-energy} to normalized centers.

\paragraph{Step 1: reflection along every sequence.}
Let $T_n\to\infty$ be arbitrary. Using the precompactness in \cref{prop:stationary-clusters}, we see that every subsequence of $(T_n)$ has a further subsequence along which $(C_{T_n},R_{T_n},\lambda_{T_n})$ converges to a stationary cluster. By hypothesis, this cluster is $(c,r,\ell)$, with the multiplier $\ell$ given above. Along the further subsequence, \cref{sel:prefix-reflection} gives the limits \eqref{sel:prefix-norm} and \eqref{sel:prefix-energy}, with values $d_L$ and $h_L$ that do not depend on the subsequence, and $\lambda(T_n)\to\ell$. Since every subsequence has a further subsequence with the same limits, the whole sequence has them. As $(T_n)$ was arbitrary, we conclude that, for every $L\ge0$,
\begin{equation}\label{eq:tree-reflection-all-times}
 \lim_{T\to\infty}\norm{M_T(L)}^2=d_L,\qquad \lim_{T\to\infty}h\bigl(M_T(L)\bigr)=h_L,
\end{equation}
and that $\lambda(T)\to\ell$, so that $e_\beta(T)\to e_\infty$ by \eqref{eq:energy-identity}.

\paragraph{Step 2: the deficits.}
Recall that $c$ is nonincreasing with limit $q>0$. Then $d_L\ge q>0$, and $\delta_L$ is well defined. By part~\ref{item:threshold-ceiling} of \cref{prop:threshold-inputs}, proved in \cref{sec:tree-input-ceiling}, we have $h_L\le\Eth d_L^{p/2}$, that is, $\delta_L\ge0$. To prove that $\delta_L\to0$, recall from \cref{lem:center-energy} that $d_L\to q$ and $h_L\to h$ as $L\to\infty$, where $h$ is given by \eqref{sel:limiting-center-energy}, and that
\[
 \frac h{q^{p/2}}=\sqrt{\frac{p-1}p}\Bigl(z+\frac1z\Bigr),\qquad z=\beta\sqrt{p(p-1)}\,q^{(p-2)/2}(1-q),
\]
by \eqref{sel:center-amgm}. The hypothesis $\beta^2p(p-1)q^{p-2}(1-q)^2=1$ says that $z^2=1$. Since $z>0$, we obtain $z=1$ and $h/q^{p/2}=2\sqrt{(p-1)/p}=\Eth$. We conclude that
\[
 \lim_{L\to\infty}\frac{h_L}{d_L^{p/2}}=\frac h{q^{p/2}}=\Eth,
\]
which proves \eqref{sel:center-deficit}.

\paragraph{Step 3: the mass at low heights.}
Fix $L\ge0$ and $H\ge0$. By \eqref{eq:tree-reflection-all-times}, we have $\norm{M_T(L)}^2\to d_L>0$, and in particular $M_T(L)\ne0$ for all large $T$. We let $U_T=M_T(L)/\norm{M_T(L)}$. Using the homogeneity of $h$ and \eqref{eq:tree-reflection-all-times}, we find that
\[
 h(U_T)=\frac{h(M_T(L))}{\norm{M_T(L)}^p}\longrightarrow\frac{h_L}{d_L^{p/2}}=\Eth-\delta_L\qquad(T\to\infty).
\]
By \cref{cor:tree-ceiling-energy}, we have $\norm{\Pi_{\le H}U_T}^2\le\{\kappa_p(H+1)(\Eth-h(U_T))\}^{1/2}$, and therefore
\[
 \limsup_{T\to\infty}\norm{\Pi_{\le H}U_T}^2\le\bigl\{\kappa_p(H+1)\delta_L\bigr\}^{1/2}.
\]
Since $\norm{\Pi_{\le H}M_T(L)}^2=\norm{M_T(L)}^2\norm{\Pi_{\le H}U_T}^2$ and $\norm{M_T(L)}^2\to d_L$, we see that
\[
 \limsup_{T\to\infty}\norm{\Pi_{\le H}M_T(L)}^2\le d_L\bigl\{\kappa_p(H+1)\delta_L\bigr\}^{1/2}.
\]
Setting $T=t-L$, we have $M_{t-L}(L)=P_{t-L}v_t$, and $t\to\infty$ if and only if $T\to\infty$. This proves \eqref{sel:low-height-center-bound}.
\end{proof}

\section{Forgetting at marginal saturation}\label{sec:tree-input-forgetting}

We now prove the second threshold input. When $s\le t-L$, a fixed earlier state $v_s$ lies in the range of $P_{t-L}$, and its overlap with the late state $v_t$ equals its overlap with the prefix center $P_{t-L}v_t$. \Cref{sel:threshold-centers} shows that the center has almost no mass at low heights once $L$ is large, while the earlier states $v_s$, $s\le S$, have almost all their mass at bounded heights. The overlap is therefore small.

\begin{proof}[Proof of part~\ref{item:threshold-forgetting} of \cref{prop:threshold-inputs} for $p\in\{3,4\}$]
Suppose that the hypotheses of part~\ref{item:threshold-forgetting} hold, and fix $0\le S<\infty$. The proof uses three parameters: the offset $L\ge0$, the height cutoff $H\ge0$, and the time $t$. We first let $t\to\infty$ with $L$ and $H$ fixed, then let $L\to\infty$, and finally let $H\to\infty$.

Let $L,H\ge0$, let $t\ge S+L$, and let $0\le s\le S$. Then $s\le t-L$, and $v_s\in\mathcal H_s\subset\mathcal H_{t-L}$ by \cref{prop:geometry-tree-realization}, that is, $P_{t-L}v_s=v_s$. As $P_{t-L}$ is an orthogonal projection, we have
\[
 C(t,s)=\ip{v_t}{v_s}=\ip{v_t}{P_{t-L}v_s}=\ip{P_{t-L}v_t}{v_s}.
\]
We split $v_s=\Pi_{\le H}v_s+(\Id-\Pi_{\le H})v_s$ and move the first projection to the other side of the inner product. Using $\norm{v_s}=1$ and $\norm{P_{t-L}v_t}\le1$, we find
\begin{equation}\label{eq:tree-forgetting-split}
 |C(t,s)|\le\norm{\Pi_{\le H}P_{t-L}v_t}+\norm{(\Id-\Pi_{\le H})v_s}.
\end{equation}

We first show that the second term is small uniformly in $s\le S$. Let
\[
 \varepsilon_H=\sup_{0\le s\le S}\norm{(\Id-\Pi_{\le H})v_s}.
\]
The map $s\mapsto v_s$ is norm-continuous by \cref{prop:geometry-tree-realization}, and the set $\{v_s:0\le s\le S\}$ is compact. Recall that $\sum_j\Pi_j=\Id$ strongly, which means that $(\Id-\Pi_{\le H})x\to0$ as $H\to\infty$ for every $x\in\mathcal H$. Let $\varepsilon>0$. There exist $s_1,\dots,s_k\in[0,S]$ such that every $v_s$ with $s\le S$ lies within $\varepsilon$ of some $v_{s_i}$, and there exists $H_0$ such that $\norm{(\Id-\Pi_{\le H})v_{s_i}}<\varepsilon$ for all $i$ and $H\ge H_0$. Since $\Id-\Pi_{\le H}$ is a contraction, this implies $\norm{(\Id-\Pi_{\le H})v_s}<2\varepsilon$ for all $s\le S$ and $H\ge H_0$. As $\varepsilon$ is arbitrary, we have shown that $\varepsilon_H\to0$ as $H\to\infty$.

We now take the limits in the stated order. Taking the supremum over $s\le S$ in \eqref{eq:tree-forgetting-split} and letting $t\to\infty$, we obtain from \eqref{sel:low-height-center-bound} that
\[
 \limsup_{t\to\infty}\sup_{0\le s\le S}|C(t,s)|
 \le d_L^{1/2}\bigl\{\kappa_p(H+1)\delta_L\bigr\}^{1/4}+\varepsilon_H.
\]
The left side depends on neither $L$ nor $H$. We let $L\to\infty$ with $H$ fixed. Since $d_L\le1$ and $\delta_L\to0$ by \eqref{sel:center-deficit}, the first term on the right tends to zero. Then we let $H\to\infty$, and $\varepsilon_H\to0$. We conclude that $\lim_{t\to\infty}\sup_{0\le s\le S}|C(t,s)|=0$, which is \eqref{eq:threshold-input-forgetting}.
\end{proof}

This completes the proof of \cref{prop:threshold-inputs} for $p\in\{3,4\}$.

\chapter{The threshold inputs for general pure models}\label{app:general-selection}

In this section, we prove \cref{prop:threshold-inputs} for every $p\ge3$.
For $p\in\{3,4\}$, \cref{sec:tree-ceilings} proves both parts of that
proposition by a deterministic inequality for tree vectors. The argument
given here does not use the tree ceilings and works for every pure model.
Part~\ref{item:threshold-ceiling} of \cref{prop:threshold-inputs} is the
input to the proof of marginality in \cref{sel:positive-marginality}, and
part~\ref{item:threshold-forgetting} is the input to the proof of forgetting
in \cref{sel:fixed-wait-forgetting}. The finite forks of \cref{sec:general-forks}, the overlap penalty of
\cref{sel:absolute-field-cost}, and the saturation identity
\eqref{sel:fork-saturation} are used again in \cref{resp:appendix-p}, where
we extend the localization of the response to every $p\ge3$.

We use the landscape bounds of \cref{lem:landscape-moments}, the fixed-horizon bridge of \cref{prop:hard-bridge}, and the Brownian form of the
averaged response in \cref{lem:brown-force-response}. From
\cref{sec:ck-structure}, we use the energy identity of \cref{prop:energy-bridge} and the bound
\eqref{eq:finite-lag-flow-moments} on the moments of the response, proved there. From
\cref{sec:entropy-clusters}, we use the entropy budget of
\cref{lem:entropy-budget}, the integration by parts formula
\eqref{eq:gibbs-current-integration-by-parts}, the Fisher majorant of
\cref{thm:fisher-clock}, and the stationary clusters of
\cref{prop:stationary-clusters}. We also use the formulas
\eqref{sel:prefix-norm} and \eqref{sel:prefix-energy} for $d_L$ and $h_L$.
The proof of part~\ref{item:threshold-forgetting} uses, in addition, the
identities \eqref{sel:limiting-center-energy} and \eqref{sel:center-amgm}
and the correlation modulus \eqref{eq:global-modulus}.

\paragraph{Ideas of the proof.}
We obtain the energy ceiling in part~\ref{item:threshold-ceiling} from the
algorithmic threshold of Huang and Sellke. At a late time $T$, we continue
the trajectory by $K$ independent Brownian futures of length $L$, which we
call a fork, and average their endpoints. The current identity of
\cref{sec:entropy-clusters} shows that two branches of the fork have
overlap close to $c(a+d)$, where $a$ and $d$ are their lengths. We compute
the energy of the average exactly by Gaussian integration by parts.
Subtracting the energy of a single branch removes the common past in the
dimension limit. Using the hard--soft approximation of Sellke, we approximate the average by
a function of the disorder that is Lipschitz with a constant independent of
the dimension. The algorithmic threshold then bounds its energy. Letting
$K\to\infty$, we obtain part~\ref{item:threshold-ceiling}.

For part~\ref{item:threshold-forgetting}, we add a deterministic field to
the threshold. By convex duality, the field version of the threshold
penalizes the overlap of a near-threshold algorithm with any direction
that is independent of the disorder. However, at positive times, an earlier state $X_s$ depends on the disorder. We therefore disclose the disorder in stages,
along a finite sequence of gradient queries, and apply the penalty to one
new direction at each stage. An Euler--Maruyama approximation places $X_s$
close to the span
of finitely many such queries. Under marginality, the normalized average of
a fork saturates the threshold, and the penalty forces its overlap with
$X_s$ to vanish. Up to normalization, this overlap is $C(t,s)$.

\paragraph{Attribution.}
The algorithmic threshold is \cite[Theorem~1 and Corollary~1.8]{hs2023};
for even $p$ it was proved earlier in \cite{hs2021}. In
\cite[Section~3]{sellke2024}, Sellke combined it with a Lipschitz
approximation of Langevin dynamics on a fixed time interval to show that
the energy of the dynamics stays below the threshold. We follow that
argument in \cref{sec:general-forks}. Trajectories with a common past and
independent future noise are a standard construction in the physics of
aging. Cugliandolo and Dean considered replicas that share their noise up
to a waiting time and receive independent noise afterward in
\cite[arXiv version~1, Section~II]{cugliandolodean1995}. Barrat, Burioni, and M\'ezard classified aging through the overlaps of such copies in
\cite[Sections~2 and~4]{bbm1996}. Franz and Virasoro
proposed averaging cloned futures to estimate a quasi-state magnetization in
\cite[arXiv version~3, Section~III]{franzvirasoro2000}, and noise-averaged
magnetizations are central to Biroli's dynamical TAP approach in
\cite[arXiv version~2, Section~III.B]{biroli1999}. Here we prove the overlap and energy identities of the forks from the
current identity of \cref{sec:entropy-clusters} and from Gaussian
integration by parts. The
disclosure argument of \cref{sec:general-disclosure} is a Gaussian
conditioning argument; we credit its antecedents there.

\paragraph{Order of limits.}
Every horizon, number of branches, field strength, approximation accuracy,
and Lipschitz constant below is fixed while the dimension tends to
infinity, and we take the long-time and outer limits only afterward.
We take dimension limits along the subsequence $(N_j)$ of
\cref{thm:fisher-clock}, and we often write $N\to\infty$ for $N=N_j$,
$j\to\infty$. The parameters and their order are as follows:
\begin{equation}\label{sel:aging-limit-order}
 N_j\to\infty,\qquad n\to\infty,\qquad \varepsilon\downarrow0,\qquad
 K\to\infty,\qquad L\to\infty,\qquad M\to\infty.
\end{equation}
Here $n$ indexes a sequence of late times ($T_n$ in
\cref{sec:general-forks} and $t_n$ in \cref{sec:general-transfer}),
$\varepsilon$ is the accuracy of the Lipschitz approximation in
\cref{sel:finite-forks}, $K$ is the number of branches of a fork, $L$ is
their length, and $M$ is the number of steps of the Euler--Maruyama scheme of
the proof of \cref{sel:general-fixed-wait}. Each parameter is fixed while the limits to its
left are taken. Part~\ref{item:threshold-ceiling} uses only the first four
limits, at a fixed $L$; part~\ref{item:threshold-forgetting} uses all six.
For the forgetting proof, it suffices to take the approximation accuracies
$\varepsilon=1/(m+1)$, $m\in\mathbb N$, so $\varepsilon\downarrow0$ there
means $m\to\infty$.
We remove the tolerance and the field strengths of the threshold theorem
inside the proof of \cref{sel:absolute-field-cost}, after the dimension
limit at fixed values of all other parameters. Constants of the Lipschitz
maps may depend on every fixed parameter, and they are used only before the
limits in which that parameter varies.

\section{The algorithmic threshold and its overlap penalty}
\label{sec:general-field-threshold}

Throughout this section, $\norm{\cdot}$ is the Euclidean norm,
\[
 B_N=\{x\in\R^N:\norm x\le\sqrt N\},
\]
and $\Pi_{B_N}$ is the metric projection onto $B_N$, which is $1$-Lipschitz.
The coefficient tensor $J$ of \eqref{eq:hamiltonian} lies in the
$N^p$-dimensional space $(\R^N)^{\otimes p}$, with its Euclidean norm, and
has the standard Gaussian law there. More explicitly, we keep the ordered
coefficients and use the Hilbert norm
\[
 \norm J_{\mathrm{coef}}^2
 =\sum_{i_1,\dots,i_p=1}^N J_{i_1\dots i_p}^2 .
\]
This norm is not rescaled with the dimension. When the dependence on $J$ matters, we
write $H_{N,J}$ for $H_N$. An algorithm is a measurable map
$\mathcal A_N=\mathcal A_N(J,\zeta)$ with values in $B_N$, where the seed
$\zeta$ is a random variable independent of $J$ (possibly absent). For
$0\le L_{\mathrm{Lip}}<\infty$, we say that $\mathcal A_N$ is $L_{\mathrm{Lip}}$-Lipschitz if, for
almost every value of $\zeta$,
\[
 \norm{\mathcal A_N(J,\zeta)-\mathcal A_N(J',\zeta)}\le L_{\mathrm{Lip}}\norm{J-J'}_{\mathrm{coef}}
 \qquad\text{for all }J,J'\in(\R^N)^{\otimes p}.
\]
The constant $L_{\mathrm{Lip}}$ is deterministic and does not depend on $N$ or on the
seed.

Recall the tensor norm $Q_N$ of \eqref{eq:tensor-norm} and the constants
$A_p,c_p>0$ of \cref{lem:landscape-moments}, for which
$\Prob(Q_N>A_p)\le e^{-c_pN}$. We write
\[
 Q_*=\sup_N(\E Q_N^2)^{1/2},
\]
which is finite by \cref{lem:landscape-moments}. By the definition of
$Q_N$, a contraction of $J$ against vectors $u_1,\dots,u_p$ is at most
$\sqrt N\,Q_N\norm{u_1}\cdots\norm{u_p}$ in absolute value. A derivative of
order $k\le2$ of $H_N$ at $x$, applied to $k$ vectors, is
$N^{-(p-1)/2}$ times a sum of $p(p-1)\cdots(p-k+1)$ such contractions in
which $p-k$ slots equal $x$. For $R\ge1$ and $\norm x\le R\sqrt N$, this
implies that
\begin{equation}\label{eq:general-ball-derivatives}
\begin{gathered}
 \frac{|H_N(x)|}N\le R^pQ_N,\qquad
 \norm{\nabla H_N(x)}\le pR^{p-1}Q_N\sqrt N,\\
 \norm{\nabla^2H_N(x)}_\op\le p(p-1)R^{p-2}Q_N.
\end{gathered}
\end{equation}
For $R=1$, the first two bounds read
\begin{equation}\label{sel:tensor-ui}
 \sup_{x\in B_N}\frac{|H_N(x)|}N\le Q_N,\qquad
 \sup_{x\in B_N}\norm{\nabla H_N(x)}\le pQ_N\sqrt N.
\end{equation}
In particular, the normalized energy of any algorithm is bounded by $Q_N$,
whose moments are bounded uniformly in $N$. Now let $x$ and $y$ be random
vectors of norm at most $R\sqrt N$. The segment between them lies in the
same ball. Using the mean value theorem, \eqref{eq:general-ball-derivatives},
and the Cauchy--Schwarz inequality, we find
\begin{equation}\label{eq:general-energy-transfer}
 \frac{\E|H_N(x)-H_N(y)|}N
 \le pR^{p-1}Q_*\left(\E\frac{\norm{x-y}^2}N\right)^{1/2}.
\end{equation}

For $\kappa\ge0$, let $q_1(\kappa)=(\kappa^2/(p(p-2)))^{1/(p-1)}$, and
define
\begin{equation}\label{eq:alg-field}
 \ALG_p(\kappa)=
 \begin{cases}
 \Eth+\dfrac{p-2}{p}\sqrt{p(p-1)}\,q_1(\kappa)^{p/2},
   &\kappa^2\le p(p-2),\\[3mm]
 \sqrt{p+\kappa^2},&\kappa^2\ge p(p-2).
 \end{cases}
\end{equation}
In terms of $\kappa$, the first line equals
$\Eth+\gamma_p\kappa^{p/(p-1)}$ with
$\gamma_p=((p-2)/p)\sqrt{p(p-1)}\,[p(p-2)]^{-p/(2(p-1))}$.

\begin{theorem}[External-field Lipschitz threshold of Huang and Sellke]
\label{sel:published-field-threshold}
Fix an integer $p\ge3$, $0\le L_{\mathrm{Lip}}<\infty$, $\kappa\ge0$, and $\eta>0$. There exist $c>0$
and $N_0$, depending only on $p$, $L_{\mathrm{Lip}}$, $\kappa$, and $\eta$, such that
the following holds for every $N\ge N_0$. Let $w\in\R^N$ be deterministic
with $\norm w=\sqrt N$, and let $\mathcal A_N:(\R^N)^{\otimes p}\to B_N$ be
$L_{\mathrm{Lip}}$-Lipschitz. Then
\begin{equation}\label{sel:field-threshold}
 \Prob\left(\frac{H_N(\mathcal A_N(J))+\kappa\ip{w}{\mathcal A_N(J)}}N
 \ge\ALG_p(\kappa)+\eta\right)\le e^{-cN}.
\end{equation}
\end{theorem}

\Cref{sel:published-field-threshold} is the case of one species and
$\xi(q)=q^p$ of \cite[Theorem~1 and Corollary~1.8]{hs2023}. We now check
that our setting satisfies the hypotheses there.
In \cite{hs2023}, the Hamiltonian is $\ip{h}{\sigma}$ plus a Gaussian part
on the ball of radius $\sqrt N$, the field is $h$ times the all-ones vector
$\mathbf 1$ with $h\ge0$, and an algorithm is a function of the disorder
coefficients, concatenated into one vector with the Euclidean distance.
Theorem~1 there states that, for fixed $\tau,\eta>0$ and $N$ large, every
$\tau$-Lipschitz algorithm reaches energy at least $\ALG+\eta$ with
probability at most $e^{-cN}$, with $c$ and the threshold for $N$ depending
on $\tau$, $\eta$, the covariance, and the field only. For one species with
$\xi(q)=q^p$, the Gaussian part is exactly $H_N$ of \eqref{eq:hamiltonian},
whose ordered coefficients are independent standard Gaussians; we take
$h=\kappa$. Our algorithms are functions of $J$ alone, the field being
deterministic. Corollary~1.8 of \cite{hs2023}, which covers odd and even
$p$, states that $\ALG=(\xi'(1)+\kappa^2)^{1/2}$ if
$\xi'(1)+\kappa^2\ge\xi''(1)$, and otherwise
\[
 \ALG=q_1\,\xi''(q_1)^{1/2}+\int_{q_1}^1\xi''(q)^{1/2}\,\dd q,
\]
where $q_1\in[0,1)$ solves $\xi'(q_1)+\kappa^2=q_1\xi''(q_1)$. For
$\xi(q)=q^p$, the first condition reads $\kappa^2\ge p(p-2)$, the equation
for $q_1$ reads $\kappa^2=p(p-2)q_1^{p-1}$, and
\[
\begin{aligned}
 q_1\xi''(q_1)^{1/2}+\int_{q_1}^1\xi''(q)^{1/2}\,\dd q
 &=\sqrt{p(p-1)}\left(q_1^{p/2}+\frac2p\big(1-q_1^{p/2}\big)\right)\\
 &=\Eth+\frac{p-2}p\sqrt{p(p-1)}\,q_1^{p/2},
\end{aligned}
\]
where we used $2\sqrt{p(p-1)}/p=\Eth$. This is \eqref{eq:alg-field}. At
$\kappa^2=p(p-2)$, we have $q_1=1$ and both lines of \eqref{eq:alg-field}
equal $\sqrt{p(p-1)}$.

It remains to replace $\mathbf 1$ by an arbitrary direction $w$. Let $O$ be
an orthogonal matrix such that $O\mathbf 1=w$. The map $J\mapsto O^{\otimes p}J$ is
a linear isometry of $(\R^N)^{\otimes p}$ that preserves the standard
Gaussian law, and $H_{N,O^{\otimes p}J}(Ox)=H_{N,J}(x)$ for all $x$. Given
$\mathcal A_N$, the map $\mathcal A'_N(J)=O^{\mathsf T}\mathcal A_N(O^{\otimes p}J)$
is $L_{\mathrm{Lip}}$-Lipschitz with values in $B_N$, and, with $J'=O^{\otimes p}J$,
\[
 H_{N,J}(\mathcal A'_N(J))+\kappa\ip{\mathbf 1}{\mathcal A'_N(J)}
 =H_{N,J'}(\mathcal A_N(J'))+\kappa\ip{w}{\mathcal A_N(J')}.
\]
The tensor $J'$ has the law of $J$, and \cite[Theorem~1]{hs2023} applied to
$\mathcal A'_N$ gives \eqref{sel:field-threshold}, with constants that do
not depend on $w$ or on $\mathcal A_N$. For $\kappa=0$, we have
$\ALG_p(0)=\Eth$, and \cref{sel:published-field-threshold} is the
pure-model threshold, recorded in this form in
\cite[Proposition~1.9]{sellke2024}.

In this argument, we may always use the larger positive bound
$L_{\mathrm{Lip}}+1$. This includes
$L_{\mathrm{Lip}}=0$ and gives the same uniform choice of $c$ and $N_0$
for every algorithm with the prescribed original bound. We need no separate
measurability assumption on a deterministic algorithm, since Lipschitz maps
are continuous.

\paragraph{Constructive finite-tree qualifications.}
The branching proof of the threshold uses finite trees whose depth and
branching degrees are fixed before the dimension tends to infinity. Two
auxiliary statements in \cite[Appendix~A]{hs2023} require qualifications
in a constructive version. In Lemma~A.8, if the absolute values of the leaf
energies are bounded by $C>0$ and the depth is $D\ge1$, the retained branching
degree $\lfloor K\varepsilon/(3CD)\rfloor$ is obtained in the range
$0<\varepsilon\le CD$, and $K$ is then taken large enough to make this
degree positive. This is the range used when $\varepsilon\downarrow0$;
an unrestricted error can make the displayed degree exceed $K$.

For the base-preserving fork correction of Lemma~A.9, we impose $k<N$
when $k$ orthogonal increments are retained. After dividing
the vertices by $\sqrt N$, suppose that the base $x$ has squared norm
$q\ge q_0>0$, the target increment has squared norm $q'\ge0$, and the
approximate increments $y_1,\dots,y_k$ satisfy, for $\varepsilon\ge0$,
\[
 |\norm{y_i}^2-q'|\le\varepsilon,\qquad
 |\ip x{y_i}|\le\varepsilon,\qquad
 |\ip{y_i}{y_j}|\le\varepsilon\quad(i\ne j).
\]
Under the sufficient small-error conditions
$\sqrt\varepsilon\le q_0$ and $4(k+1)\sqrt\varepsilon\le1$, the base can
be kept exactly and the increments corrected to mutually orthogonal
vectors $z_i\perp x$ with $\norm{z_i}^2=q'$ and
$\norm{z_i-y_i}\le e_k(\varepsilon)$, where
\[
 e_k(\varepsilon)
 =(4\sqrt{k+1}+1)\sqrt\varepsilon
   +2\sqrt{\sqrt\varepsilon}.
\]
This bound permits $q'=0$ and tends to zero for fixed $k$ as
$\varepsilon\downarrow0$. It supplies a vanishing modulus, independent of the dimension, in place of
the piecewise numerical modulus printed in that lemma.

We handle a schedule with zero initial overlap before this correction.
For fixed $\delta\in(0,1)$, we scale the normalized vertices, which have norm
at most one, by $\sqrt{1-\delta}$ and add a common orthogonal vector of norm
$\sqrt\delta$. The new schedule of target overlaps is
$q_i^{(\delta)}=(1-\delta)q_i+\delta$, so its root is at least $\delta$;
the terminal vectors change by at most $2\sqrt\delta$. We choose the dimension larger than the number of vertices of the retained
finite tree, as well as its branching degree, so that the common orthogonal
direction and all fork corrections exist. For fixed depth and retained
branching, the accumulated correction tends to zero, and $\delta$ can be
removed afterward. Repeated overlap levels are allowed: their target
increments have norm zero.

In the algorithmic application, we can retain a fixed deterministic set of children, because the success event imposes the desired energy bound
at every original terminal leaf. Internal vertices are first obtained
from averages of the original terminal outputs. Restricting the independent Gaussian increments of the tree to this fixed
set preserves their joint law. We therefore need no distributional assertion
about an adaptively selected subtree. The depths, branching degrees, error
tolerances, and dimension cutoffs are chosen before the algorithm and the
deterministic field direction. The calibrated overlap schedule may depend on the algorithm,
since the finite-tree ceiling is uniform over all admissible schedules.
This is the uniformity asserted in \cref{sel:published-field-threshold}.

The field version of the threshold penalizes the overlap of an algorithm
with a direction independent of $J$. To state the penalty, recall
$\qd=(p-2)/(p-1)$ and $\bd^2=(p-1)^{p-1}/(p(p-2)^{p-2})$. The constant $\bd$
is the dynamical transition point, and it enters here only as the
coefficient of the penalty. The inverse temperature $\beta$ of the dynamics
plays no role in \cref{sel:published-field-threshold,sel:absolute-field-cost}.

\begin{lemma}[Absolute-overlap energy cost]\label{sel:absolute-field-cost}
Fix an integer $p\ge3$ and $0\le L_{\mathrm{Lip}}<\infty$. There exist
finite real numbers $\phi_N=\phi_N(p,L_{\mathrm{Lip}})\ge0$, for every $N\ge1$,
with $\phi_N\to0$ as $N\to\infty$
such that the following holds. Let $\zeta$ be a random variable independent of $J$, let
$w=w(\zeta)\in\R^N$ satisfy $\norm w=\sqrt N$, let $\mathcal A_N$ be an
$L_{\mathrm{Lip}}$-Lipschitz algorithm with seed $\zeta$, and set
$a_N=\ip{w}{\mathcal A_N}/N$. Then, almost surely,
\begin{equation}\label{eq:conditional-field-cost}
 \E\left[\frac{H_N(\mathcal A_N)}N+\bd|a_N|^p\,\middle|\,\zeta\right]
 \le\Eth+\phi_N.
\end{equation}
In particular,
\begin{equation}\label{sel:absolute-field-cost-equation}
 \limsup_{N\to\infty}
 \left\{\E\frac{H_N(\mathcal A_N)}N+\bd\,\E|a_N|^p\right\}\le\Eth.
\end{equation}
\end{lemma}

The lemma makes no assumption on the concentration of $a_N$, and the seed
may be absent, in which case $w$ is deterministic.

\begin{proof}
The proof has three steps. We first compute the best ceiling that the field
thresholds give for a given overlap, then we use finitely many fields at
once, and finally we refine the set of fields after the dimension limit.

\paragraph{Step 1: the conjugate ceiling.}
For $0\le a\le1$, we define
\[
 \mathcal E_p(a)=\inf_{\kappa\ge0}\{\ALG_p(\kappa)-\kappa a\}.
\]
We claim that
\begin{equation}\label{eq:conjugate-ceiling}
 \mathcal E_p(a)=
 \begin{cases}
 \Eth-\bd a^p,&0\le a\le\sqrt{\qd},\\
 \sqrt{p(1-a^2)},&\sqrt{\qd}\le a\le1,
 \end{cases}
 \qquad\text{and}\qquad
 \Eth-\mathcal E_p(a)\ge\bd a^p\quad(0\le a\le1).
\end{equation}
The function $\ALG_p$ is convex on $[0,\infty)$, since $\kappa^{p/(p-1)}$ and
$\sqrt{p+\kappa^2}$ are convex. It is also continuously differentiable. Indeed, let $\kappa_*=\sqrt{p(p-2)}$. On $[0,\kappa_*]$, its derivative
is $p(p-1)^{-1}\gamma_p\kappa^{1/(p-1)}$, which increases from $0$ to
\[
 \frac p{p-1}\gamma_p\kappa_*^{1/(p-1)}
 =\frac{p-2}{p-1}\sqrt{\frac{p(p-1)}{p(p-2)}}=\sqrt{\qd}.
\]
On $[\kappa_*,\infty)$, its derivative is $\kappa/\sqrt{p+\kappa^2}$,
which increases from $\sqrt{\qd}$ to $1$. By convexity, for $0\le a<1$ the
infimum defining $\mathcal E_p(a)$ is attained where $\ALG_p'(\kappa)=a$.

If $a\le\sqrt{\qd}$, this point is
$\kappa=((p-1)a/(p\gamma_p))^{p-1}$. At this point, we have $\gamma_p\kappa^{1/(p-1)}=(p-1)a/p$, and therefore
\[
 \mathcal E_p(a)=\Eth+\kappa\big(\gamma_p\kappa^{1/(p-1)}-a\big)
 =\Eth-\frac{\kappa a}p
 =\Eth-\frac1p\left(\frac{p-1}{p\gamma_p}\right)^{p-1}a^p.
\]
Inserting the value of $\gamma_p$, we have
\[
\begin{aligned}
 \left(\frac{p-1}{p\gamma_p}\right)^{p-1}
 &=\left(\frac{p-1}{p-2}\right)^{p-1}\big(p(p-1)\big)^{-(p-1)/2}
  \big(p(p-2)\big)^{p/2}\\
 &=p^{1/2}(p-1)^{(p-1)/2}(p-2)^{-(p-2)/2},
\end{aligned}
\]
and dividing by $p$, we find that the coefficient of $a^p$ is
$((p-1)^{p-1}/(p(p-2)^{p-2}))^{1/2}=\bd$.

If $\sqrt{\qd}\le a<1$, the point is $\kappa=\sqrt p\,a/\sqrt{1-a^2}$, and
\[
 \mathcal E_p(a)=\sqrt{\frac p{1-a^2}}-\frac{\sqrt p\,a^2}{\sqrt{1-a^2}}
 =\sqrt{p(1-a^2)}.
\]
At $a=1$, we have $\mathcal E_p(1)=0$, because the function
$\sqrt{p+\kappa^2}-\kappa$ decreases to $0$ as $\kappa\to\infty$. The two formulas agree at
$a=\sqrt{\qd}$, where both equal $\sqrt{p/(p-1)}$. This proves the first
part of \eqref{eq:conjugate-ceiling}, and shows that $\mathcal E_p$ is
continuous on $[0,1]$.

Let $D_p(a)=\Eth-\mathcal E_p(a)$. For $a\le\sqrt{\qd}$, we have
$D_p(a)=\bd a^p$. For $\sqrt{\qd}<a<1$, we have
$D_p'(a)=\sqrt p\,a/\sqrt{1-a^2}$. The function $a^{p-2}\sqrt{1-a^2}$ on
$[0,1]$ is maximal at $a^2=\qd$, where it equals
\[
 \qd^{(p-2)/2}(1-\qd)^{1/2}=(p-2)^{(p-2)/2}(p-1)^{-(p-1)/2}
 =\frac{1}{\sqrt p\,\bd}.
\]
Since $D_p'(a)/(p\bd a^{p-1})=(\sqrt p\,\bd\,a^{p-2}\sqrt{1-a^2})^{-1}$, we conclude that $D_p'(a)\ge p\bd a^{p-1}$ on $(\sqrt{\qd},1)$. Integrating this
inequality from $\sqrt{\qd}$, where $D_p(\sqrt{\qd})=\bd\qd^{p/2}$, and
using the continuity of $D_p$ at $a=1$, we find that $D_p(a)\ge\bd a^p$ on
$[\sqrt{\qd},1]$. This proves \eqref{eq:conjugate-ceiling}.

\paragraph{Step 2: finitely many fields.}
Fix a finite set $\mathcal G\subset[0,\infty)$ with $0\in\mathcal G$, and
$\eta>0$. We define
\[
 V_{\mathcal G}(a)=\min_{\kappa\in\mathcal G}\{\ALG_p(\kappa)-\kappa a\},
 \qquad 0\le a\le1.
\]
Note that $|V_{\mathcal G}|\le M_{\mathcal G}$, where
$M_{\mathcal G}=\Eth+\max\mathcal G$, because $0\in\mathcal G$ and
$\ALG_p\ge0$. We condition on $\zeta$. Then $w$ is
fixed, $J$ is standard Gaussian, and $J\mapsto\mathcal A_N(J,\zeta)$ is
$L_{\mathrm{Lip}}$-Lipschitz. For each $\kappa\in\mathcal G$, we use
\cref{sel:published-field-threshold} with the field strength $\kappa$ and the
directions $w$ and $-w$. Let $c(\mathcal G,\eta)>0$ and $N_0(\mathcal G,\eta)$
be the smallest constant and the largest threshold for $N$ among these
$|\mathcal G|$ applications. For $N\ge N_0(\mathcal G,\eta)$, the event
\[
 \mathfrak B=\bigcup_{\kappa\in\mathcal G}\left\{
 \frac{H_N(\mathcal A_N)}N+\kappa|a_N|\ge\ALG_p(\kappa)+\eta\right\}
\]
has conditional probability at most $2|\mathcal G|e^{-c(\mathcal G,\eta)N}$,
uniformly in $\zeta$. Off $\mathfrak B$, we have
$H_N(\mathcal A_N)/N<V_{\mathcal G}(|a_N|)+\eta$; note that $|a_N|\le1$,
since $\norm w=\sqrt N$ and $\mathcal A_N\in B_N$. On $\mathfrak B$, we use
$|H_N(\mathcal A_N)|/N\le Q_N$ from \eqref{sel:tensor-ui}. Recall that $J$ is
independent of $\zeta$, which gives $\E[Q_N^2\mid\zeta]=\E Q_N^2\le Q_*^2$.
Let
\[
 r_N(\mathcal G,\eta)=(Q_*+M_{\mathcal G})
 \big(2|\mathcal G|e^{-c(\mathcal G,\eta)N}\big)^{1/2}.
\]
Using the Cauchy--Schwarz inequality, we obtain
\begin{equation}\label{eq:finite-grid-cost}
 \E\left[\frac{H_N(\mathcal A_N)}N-V_{\mathcal G}(|a_N|)\,\middle|\,\zeta\right]
 \le\eta+r_N(\mathcal G,\eta)
\end{equation}
for $N\ge N_0(\mathcal G,\eta)$. For $1\le N<N_0(\mathcal G,\eta)$, we set
$r_N(\mathcal G,\eta)=Q_*+M_{\mathcal G}$. Then
\eqref{eq:finite-grid-cost} still holds, directly from
$|H_N(\mathcal A_N)|/N\le Q_N$ and $|V_{\mathcal G}|\le M_{\mathcal G}$.
The finite quantity $r_N(\mathcal G,\eta)$ depends only on $p$, $L_{\mathrm{Lip}}$, $\mathcal G$, and $\eta$, and tends to zero as $N\to\infty$.

\paragraph{Step 3: refining the fields.}
Let $\mathcal G_k=\{i2^{-k}:0\le i\le4^k\}$. These sets increase with $k$,
and their union is dense in $[0,\infty)$. The functions
$V_{\mathcal G_k}$ are continuous on $[0,1]$ and decrease in $k$. For each
$a\in[0,1]$, the map $\kappa\mapsto\ALG_p(\kappa)-\kappa a$ is continuous,
and its infimum over the dense set $\bigcup_k\mathcal G_k$ equals
$\mathcal E_p(a)$. The limit $\mathcal E_p$ is continuous on $[0,1]$. By
Dini's theorem, the quantity
$s_k=\sup_{0\le a\le1}\{V_{\mathcal G_k}(a)-\mathcal E_p(a)\}$ tends to zero.
By \eqref{eq:conjugate-ceiling} and \eqref{eq:finite-grid-cost} with
$\mathcal G=\mathcal G_k$ and $\eta=1/k$, we have
\[
\begin{aligned}
 \E\left[\frac{H_N(\mathcal A_N)}N+\bd|a_N|^p\,\middle|\,\zeta\right]
 &\le\Eth+\E\left[\frac{H_N(\mathcal A_N)}N-\mathcal E_p(|a_N|)\,\middle|\,\zeta\right]\\
 &\le\Eth+s_k+\frac1k+r_N(\mathcal G_k,1/k)
\end{aligned}
\]
for every $k\ge1$. We define
$\phi_N=\inf_{k\ge1}\{s_k+1/k+r_N(\mathcal G_k,1/k)\}$, which depends only on
$p$ and $L_{\mathrm{Lip}}$. This infimum is nonnegative and finite for every
$N\ge1$, since already the term with $k=1$ is finite. Then
\eqref{eq:conditional-field-cost} holds. Moreover,
$\limsup_N\phi_N\le s_k+1/k$ for every $k$, and in particular $\phi_N\to0$. Taking
expectations in \eqref{eq:conditional-field-cost} and letting $N\to\infty$
gives \eqref{sel:absolute-field-cost-equation}.
\end{proof}

\section{Finite forks and the exact energy subtraction}\label{sec:general-forks}

In this section, we fix $\beta>0$ and a stationary cluster $(c,r,\ell)$ of
the translated CK kernels, as in \cref{prop:stationary-clusters}, realized
along a sequence $T_n\to\infty$. Its plateau is $q$ and its energy is
$e_\infty=(\ell-1)/(p\beta)$. Recall from \eqref{eq:cluster-fdt-properties}
that $0\le c\le1$, that $r=-c'\ge0$, so that $c$ is nonincreasing, and that
$\int_0^\infty r(\tau)\,\dd\tau=1-q$. Recall also that $e_\beta(T_n+a)\to e_\infty$
for every $a\ge0$, by \eqref{eq:energy-identity} and the convergence of
the translated multipliers.

\paragraph{Entropy notation.}
We use the notation of \cref{sec:entropy-clusters}. For fixed disorder,
$\rho_t$ is the density of $X_t$ with respect to $\omega_N$ and $g_N$ is the
Gibbs density. We write
\[
 v^{\mathrm{cur}}_t=\nabla_S\log(\rho_t/g_N),\qquad
 I_H(t)=\int\norm{v^{\mathrm{cur}}_t}^2\rho_t\,\dd\omega_N,\qquad
 a_N(t)=\frac1N\E I_H(t)
\]
for the relative score, the Fisher information, and its normalized mean.
This is the notation of \cref{sec:entropy-clusters}; the superscript
distinguishes the score from the tree vectors $v_t$. Let $\mu_N(\dd t)=a_N(t)\,\dd t$. \Cref{lem:entropy-budget} implies that
$\mu_N([0,\infty))\le\beta^2/2$. By \eqref{eq:fisher-density-budget},
$\mu_{N_j}(I)\to\mu(I)$ for every
bounded interval $I$, where $\mu(\dd t)=a_F(t)\,\dd t$ is the finite measure
of \eqref{eq:entropy-measure}. We write $P^H_t$ for the diffusion
semigroup at fixed disorder and $\mathcal L_N$ for its generator.

\paragraph{Forks.}
For each fixed integer $K\ge1$, we use the product probability space
carrying $J$, the uniform initial point $X_0$, a full Brownian path
$(B_u)_{u\ge0}$, and $K$ full Brownian paths $B^{(1)},\dots,B^{(K)}$
on $[0,\infty)$. These coordinate blocks are mutually independent. We call
\[
 \zeta_N=\big(X_0,(B^{(b)})_{1\le b\le K},B\big)
\]
the fork seed, with its uniform-initial and Wiener product law; the
disorder has its independent standard Gaussian law. The seed space and law
depend on $K$, which we suppress in the notation. Fix $T\ge0$ and $L\ge0$. The trajectory $X$
solves \eqref{eq:hard-sde} driven by $B$. For $b\le K$, let $Z_b$ be the
solution of \eqref{eq:hard-sde} on $[0,T+L]$ with $Z_b(0)=X_0$, driven by
the Brownian motion that equals $B_u$ for $u\le T$ and
$B_T+B^{(b)}_{u-T}$ for $u\ge T$. Then $Z_b(u)=X_u$ for $u\le T$, and after
time $T$ the paths $Z_1,\dots,Z_K$ are continuations of $X$ from $X_T$ by
independent noises. We write
\[
 Y_b=Z_b(T+L),\qquad m_{N,T}=\frac1K\sum_{b=1}^KY_b,
\]
and call $m_{N,T}$ the center of the fork. When needed, we indicate the
dependence on $K$ and $L$ in the notation. The Brownian motion that drives
$Z_b$ is a standard Brownian motion independent of $(J,X_0)$. Hence each
pair $(J,Z_b)$ has the same law as $(J,(X_u)_{u\le T+L})$. The branches
share their past before $T$ and are not independent. Note that $\norm{m_{N,T}}\le\sqrt N$, because
$\norm{Y_b}=\sqrt N$.

\paragraph{Response moments.}
Let $U_b(t,u)$ be the variational response matrix of the path $Z_b$, defined as in \eqref{eq:response-matrix}. Each driving process is Brownian for its own full filtration and is independent of $(J,X_0)$, even though the branches share a common past. Therefore the proof of \eqref{eq:finite-lag-flow-moments} applies to every branch. With the constants $K_S$ and $\Theta_S$ defined there, its coordinate columns and normalized trace have fourth moment at most $K_S$ whenever $0\le t-u\le S$. In particular,
\[
 \E\left|\frac{\tr U_b(t,u)}N\right|^2\le\Theta_S^2,
 \qquad \E\norm{U_b(t,u)}_{\op}^2\le N\Theta_S^2.
\]
We use these estimates on the whole interval $[0,T+L]$, with $S=T+L$, and on the final interval $[T,T+L]$, with $S=L$. The latter bounds are uniform in $T$.

\begin{lemma}[Entropy comparison for two branches]\label{sel:fork-moments}
Let $T,a,d\ge0$, and let $Y$ and $Z$ be the endpoints of two branches at
time $T$ of lengths $a$ and $d$, that is, of two continuations of $X$ from
$X_T$ driven by independent Brownian motions. Set
$Q_{N,T}(a,d)=\ip{Y}{Z}/N$. Then, for $k\in\{1,2\}$,
\begin{equation}\label{sel:fork-moment-comparison}
 \left|\E Q_{N,T}(a,d)^k-\E C_N(T+2a+d,T+a)^k\right|
 \le2k\,e^{\beta p^2A_p(a+d)}\sqrt{a\,\mu_N([T,T+a])}+2e^{-c_pN}.
\end{equation}
Moreover, if $T_n\to\infty$ is a sequence along which the translated kernels
converge to a stationary cluster with correlation profile $c$, then, for
every $L<\infty$,
\begin{equation}\label{sel:fork-overlap-concentration}
 \lim_{n\to\infty}\limsup_{j\to\infty}\sup_{0\le a,d\le L}
 \E\left|Q_{N_j,T_n}(a,d)-c(a+d)\right|^2=0.
\end{equation}
\end{lemma}

\begin{proof}
For $k\in\{1,2\}$, let $\Phi_k(x)=(x/\sqrt N)^{\otimes k}$, with values in
$(\R^N)^{\otimes k}$ and its Euclidean norm, so that
$\ip{\Phi_k(x)}{\Phi_k(y)}=(\ip xy/N)^k$. On $S_N$, we have
$\norm{\Phi_k}=1$ and $\norm{D_S\Phi_k}_\op\le k/\sqrt N$. Fix the disorder.
For $0\le z\le a$, we define
\[
 f_z=P^H_{a-z}\Phi_k,\qquad g_z=P^H_{d+z}\Phi_k,\qquad
 F_J(z)=\int_{S_N}\rho_{T+z}\ip{f_z}{g_z}\,\dd\omega_N,
\]
where the semigroup acts on each coordinate of $\Phi_k$.

We first identify the endpoint values. By the Markov property and the
conditional independence of the two branches given the past up to time $T$,
we have $\E F_J(0)=\E Q_{N,T}(a,d)^k$. By the semigroup property,
$F_J(a)$ is the conditional expectation of
$\ip{\Phi_k(X_{T+a})}{\Phi_k(X_{T+2a+d})}$ given $J$, so that
$\E F_J(a)=\E C_N(T+2a+d,T+a)^k$.

Next we compute $F_J'$. We have $\partial_zf_z=-\mathcal L_Nf_z$, $\partial_zg_z=\mathcal L_Ng_z$, and
$\partial_z\rho_{T+z}=\mathcal L_N^*\rho_{T+z}$. By the product rule for the
generator, applied to each pair of coordinates,
$\mathcal L_N\ip{f}{g}=\ip{\mathcal L_Nf}{g}+\ip{f}{\mathcal L_Ng}
+2\sum_i\ip{D_Sf[e_i]}{D_Sg[e_i]}$, where $(e_i)$ is an orthonormal frame of
the tangent space. Using these identities, we find
\[
\begin{aligned}
 F_J'(z)&=\int\rho_{T+z}\left\{\mathcal L_N\ip{f_z}{g_z}
 -\ip{\mathcal L_Nf_z}{g_z}+\ip{f_z}{\mathcal L_Ng_z}\right\}\dd\omega_N\\
 &=2\int\rho_{T+z}\left\{\ip{f_z}{\mathcal L_Ng_z}
 +\sum_i\ip{D_Sf_z[e_i]}{D_Sg_z[e_i]}\right\}\dd\omega_N.
\end{aligned}
\]
When we apply \eqref{eq:gibbs-current-integration-by-parts} to each
coordinate of $f_z$ and $g_z$, the gradient terms cancel, and we obtain the
exact current identity
\begin{equation}\label{sel:fork-current-identity}
 F_J'(z)=-2\int\rho_{T+z}
 \ip{P^H_{a-z}\Phi_k}{D_S\big(P^H_{d+z}\Phi_k\big)\big[v^{\mathrm{cur}}_{T+z}\big]}
 \,\dd\omega_N.
\end{equation}
Note that this identity does not use stationarity of $\rho_{T+z}$: the
score $v^{\mathrm{cur}}_{T+z}$ records how far $\rho_{T+z}$ is from the
Gibbs density.

We now bound the right side of \eqref{sel:fork-current-identity} on the
event $G_N=\{Q_N\le A_p\}$. By Jensen's inequality,
$\norm{P^H_t\Phi_k}\le1$. In the proof of \cref{lem:waiting-clock-current},
the Bochner formula, the nonnegative Ricci curvature of $S_N$, and the Hessian bound of \eqref{eq:landscape-derivative-bounds} give the gradient
bound $\norm{\nabla_SP^H_t\varphi}^2\le e^{2\beta p^2Q_Nt}P^H_t\norm{\nabla_S\varphi}^2$
for smooth scalar $\varphi$; this is the classical gradient bound under a
lower bound on the curvature. Applied to $\varphi=\ip{w}{\Phi_k}$ for a unit
vector $w$ in $(\R^N)^{\otimes k}$, and followed by the supremum over $w$, it
gives
\[
 \norm{D_SP^H_t\Phi_k}_\op\le\frac k{\sqrt N}e^{\beta p^2A_pt}
 \qquad\text{on }G_N,
\]
with no factor depending on the dimension of $(\R^N)^{\otimes k}$. Using the
Cauchy--Schwarz inequality, we find that, on $G_N$ and for $z\le a$,
\[
 |F_J'(z)|\le\frac{2k}{\sqrt N}e^{\beta p^2A_p(a+d)}
 \int\rho_{T+z}\norm{v^{\mathrm{cur}}_{T+z}}\,\dd\omega_N
 \le2k\,e^{\beta p^2A_p(a+d)}\left(\frac{I_H(T+z)}N\right)^{1/2}.
\]
Taking expectations and using Jensen's inequality, then integrating in $z$
and using the Cauchy--Schwarz inequality, we obtain
\[
\begin{aligned}
 \int_0^a\E\big[\1_{G_N}|F_J'(z)|\big]\,\dd z
 &\le2k\,e^{\beta p^2A_p(a+d)}\int_0^a\sqrt{a_N(T+z)}\,\dd z\\
 &\le2k\,e^{\beta p^2A_p(a+d)}\sqrt{a\,\mu_N([T,T+a])}.
\end{aligned}
\]
On $G_N^c$, we have $|F_J(0)|,|F_J(a)|\le1$, and $\Prob(G_N^c)\le e^{-c_pN}$.
Writing
$\E[F_J(0)-F_J(a)]=-\E[\1_{G_N}\int_0^aF_J'(z)\,\dd z]+\E[\1_{G_N^c}(F_J(0)-F_J(a))]$
and combining the last two bounds with the identification of the endpoint
values, we obtain \eqref{sel:fork-moment-comparison}.

For \eqref{sel:fork-overlap-concentration}, fix $L$ and let $a,d\le L$. In
\eqref{sel:fork-moment-comparison}, we may replace $a\,\mu_N([T,T+a])$ by
$L\,\mu_N([T,T+L])$ and the exponential factor by $e^{2\beta p^2A_pL}$. Then
\eqref{eq:fisher-density-budget} implies that
$\mu_{N_j}([T,T+L])\to\mu([T,T+L])$ for each fixed $T$. By
\cref{prop:hard-bridge} on the horizon $T+3L$ and $|C_N|\le1$, the moments
$\E C_N(T+2a+d,T+a)^k$ converge to $C(T+2a+d,T+a)^k$ uniformly in
$a,d\le L$. Combining these facts with \eqref{sel:fork-moment-comparison},
we find
\[
 \limsup_{j\to\infty}\sup_{0\le a,d\le L}
 \left|\E Q_{N_j,T}(a,d)^k-C(T+2a+d,T+a)^k\right|
 \le4e^{2\beta p^2A_pL}\sqrt{L\,\mu([T,T+L])}.
\]
Along $T=T_n$, the right side tends to zero because $\mu$ is a finite
measure, and $C(T_n+2a+d,T_n+a)\to c(a+d)$ uniformly in $a,d\le L$, by the
locally uniform convergence of the translated correlations on
$[0,3L]^2$. Finally, writing $Q=Q_{N,T}(a,d)$ and $c_0=c(a+d)\in[0,1]$, we
have
\[
 \E|Q-c_0|^2=(\E Q^2-c_0^2)-2c_0(\E Q-c_0),
\]
and the convergence of the first two moments, uniformly in $a,d\le L$,
proves \eqref{sel:fork-overlap-concentration}.
\end{proof}

We now state the main result of this section. For $K\ge1$ and $L\ge0$, we
define
\begin{align}
 Q_{K,L}&=\frac1K+\left(1-\frac1K\right)c(2L),
 \label{sel:fork-norm}\\
 k_{K,L}(\tau)&=\frac{c(\tau)+(K-1)c(2L-\tau)}K,\qquad 0\le\tau\le L,
 \label{sel:fork-mixed-overlap}\\
 h_{K,L}&=e_\infty-\beta p\int_0^Lr(\tau)
 \left\{c(\tau)^{p-1}-k_{K,L}(\tau)^{p-1}\right\}\dd\tau.
 \label{sel:fork-center-energy}
\end{align}
Note that $Q_{K,L}$ is a convex combination of $1$ and $c(2L)$, so that
$q\le c(2L)\le Q_{K,L}\le1$, and that $0\le k_{K,L}\le1$.

\begin{proposition}[Finite forks and their center]\label{sel:finite-forks}
Fix $\beta>0$ and let $(c,r,\ell)$ be a stationary cluster with plateau
$q>0$ and energy $e_\infty$, realized along $T_n\to\infty$. Fix $K\ge1$ and
$L\ge0$, and let $m_{N,T}$ be the center of the fork at time $T$ with $K$
branches of length $L$.
\begin{enumerate}[label=(\alph*)]
\item\label{item:fork-limits}
The squared norm and the energy of the center satisfy
\begin{align}
 &\lim_{n\to\infty}\limsup_{j\to\infty}
 \E\left|\frac{\norm{m_{N_j,T_n}}^2}{N_j}-Q_{K,L}\right|^2=0,
 \label{sel:fork-center-concentration}\\
 &\lim_{n\to\infty}\limsup_{j\to\infty}
 \left|\frac{\E H_{N_j}(m_{N_j,T_n})}{N_j}-h_{K,L}\right|=0.
 \label{sel:fork-center-energy-limit}
\end{align}
\item\label{item:fork-ceiling}
The limiting energy is below the threshold:
\begin{equation}\label{sel:fork-center-threshold}
 h_{K,L}\le\Eth\,Q_{K,L}^{p/2}.
\end{equation}
\item\label{item:fork-lipschitz}
For every $T\ge0$ and $\varepsilon>0$, there exist algorithms
$\mathcal A^\varepsilon_{N,T}=\mathcal A^\varepsilon_{N,T}(J,\zeta_N)$ with
values in $B_N$ and seed $\zeta_N$, which are $L_{\mathrm{Lip}}$-Lipschitz for a
constant $L_{\mathrm{Lip}}$ depending only on $p$, $\beta$, $T$, $K$, $L$, and
$\varepsilon$, such that
\begin{equation}\label{sel:fork-lipschitz}
 \limsup_{j\to\infty}\frac1{N_j}\E\norm{\mathcal A^\varepsilon_{N_j,T}
 -\Pi_{B_{N_j}}\big(m_{N_j,T}/\sqrt{Q_{K,L}}\big)}^2\le\varepsilon^2.
\end{equation}
Moreover, along $T=T_n$, the vector $\Pi_{B_N}(m_{N,T}/\sqrt{Q_{K,L}})$ may
be replaced by $m_{N,T}/\sqrt{Q_{K,L}}$, or by
$\widehat m_{N,T}=\sqrt N\,m_{N,T}/\norm{m_{N,T}}$ (with $\widehat m_{N,T}=0$
if $m_{N,T}=0$): the scaled mean-square distance between
$\Pi_{B_N}(m_{N,T}/\sqrt{Q_{K,L}})$ and either vector tends to zero as
$j\to\infty$ and then $n\to\infty$.
\end{enumerate}
\end{proposition}

By \eqref{sel:tensor-ui}, the normalized energies of the algorithms in
part~\ref{item:fork-lipschitz} are bounded by $Q_N$. Part~\ref{item:fork-lipschitz}
also transfers overlaps. If $V_N$ is a random vector on the same probability
space with $\norm{V_N}\le\sqrt N$ (for example $V_N=X_s$ with $s\le T$), then, by the Cauchy--Schwarz inequality, we have
\begin{equation}\label{eq:fork-overlap-transfer}
 \left|\E\frac{\ip{\mathcal A^\varepsilon_{N,T}}{V_N}}N
 -\E\frac{\ip{m_{N,T}}{V_N}}{N\sqrt{Q_{K,L}}}\right|
 \le\left(\frac1N\E\norm{\mathcal A^\varepsilon_{N,T}
 -m_{N,T}/\sqrt{Q_{K,L}}}^2\right)^{1/2}.
\end{equation}

\begin{proof}
The proof has five steps. In the first two, we compute the energy of the
center exactly by Gaussian integration by parts and subtract the energy of
a single branch. In the third, we show that the common past before $T$
contributes nothing in the dimension limit, and in the fourth we evaluate
the contribution of the last $L$ units of time; this proves
part~\ref{item:fork-limits}. In the fifth step, we construct the Lipschitz
algorithms of part~\ref{item:fork-lipschitz} and deduce
part~\ref{item:fork-ceiling} from the threshold. We write $t=T+L$.

\paragraph{Step 1: Gaussian integration by parts.}
For an ordered index $I=(i_1,\dots,i_p)$, let
$G_I(x)=N^{-(p-1)/2}x_{i_1}\cdots x_{i_p}$, so that
$H_N=\sum_IJ_IG_I$. Since $\sum_IG_I(x)G_I(z)=N(\ip xz/N)^p$,
differentiating first in $x_i$ and then in $z_l$ gives
\begin{equation}\label{sel:fork-covariance-derivative}
 \sum_I\partial_iG_I(x)\,\partial_lG_I(z)
 =p\left(\frac{\ip xz}N\right)^{p-1}\delta_{il}
 +\frac{p(p-1)}N\left(\frac{\ip xz}N\right)^{p-2}z_ix_l.
\end{equation}
Differentiating the actual terminal solution into $L^2$ with respect to the disorder, as in the proof of
\cref{prop:energy-bridge}, we have
\[
 \partial_{J_I}m_{N,T}
 =\frac\beta K\sum_{b=1}^K\int_0^tU_b(t,u)\nabla G_I(Z_b(u))\,\dd u.
\]
The seed is independent of $J$, and the response has the column and operator moments above. As in the proof of \cref{prop:energy-bridge}, apply Gaussian integration by parts to the expected coordinate observable in each disorder coefficient, using its genuine $L^2$ derivative:
\[
 \E\big[J_IG_I(m_{N,T})\big]=\E\big[\ip{\nabla G_I(m_{N,T})}{\partial_{J_I}m_{N,T}}\big],
\]
and $\E H_N(m_{N,T})$ is the sum of these terms over $I$. For $u\le t$, we
write
\[
 \gamma_b(u)=\frac{\ip{Y_b}{Z_b(u)}}N,\qquad
 \bar\gamma_b(u)=\frac{\ip{m_{N,T}}{Z_b(u)}}N
 =\frac1K\sum_{b'=1}^K\frac{\ip{Y_{b'}}{Z_b(u)}}N.
\]
Contracting \eqref{sel:fork-covariance-derivative} with $U_b(t,u)$ gives the
exact identity
\begin{align}
 \frac{\E H_N(m_{N,T})}N
 ={}&\frac{\beta p}K\sum_{b=1}^K\int_0^t
 \E\left[\frac{\tr U_b(t,u)}N\,\bar\gamma_b(u)^{p-1}\right]\dd u
 +\mathcal R_{N,T},\label{sel:fork-exact-ibp}\\
 \mathcal R_{N,T}={}&\frac{\beta p(p-1)}K\sum_{b=1}^K\int_0^t
 \E\left[\bar\gamma_b(u)^{p-2}
 \frac{\ip{Z_b(u)}{U_b(t,u)m_{N,T}}}{N^2}\right]\dd u.\notag
\end{align}
Using $\norm{m_{N,T}},\norm{Z_b(u)}\le\sqrt N$ and $|\bar\gamma_b|\le1$, we
find that $|\mathcal R_{N,T}|\le\beta p(p-1)t\,\Theta_t/\sqrt N$.

\paragraph{Step 2: the exact subtraction.}
The same computation with $m_{N,T}$ replaced by $Y_b$ gives
$\E H_N(Y_b)/N=\beta p\int_0^t\E[N^{-1}\tr U_b(t,u)\gamma_b(u)^{p-1}]\,\dd u
+O(t\,\Theta_t/\sqrt N)$. The pair $(J,Y_b)$ has the law of $(J,X_t)$, and
therefore $\E H_N(Y_b)/N=e_N(t)$, where $e_N(t)=\E H_N(X_t)/N$. Averaging over $b$ and
subtracting \eqref{sel:fork-exact-ibp}, before taking any limit, we obtain
\begin{equation}\label{sel:fork-exact-subtraction}
\begin{aligned}
 e_N(t)-\frac{\E H_N(m_{N,T})}N
 ={}&\frac{\beta p}K\sum_{b=1}^K\int_0^t
 \E\left[\frac{\tr U_b(t,u)}N
 \big\{\gamma_b(u)^{p-1}-\bar\gamma_b(u)^{p-1}\big\}\right]\dd u\\
 &+O\big(t\,\Theta_tN^{-1/2}\big),
\end{aligned}
\end{equation}
where the implicit constant depends only on $p$ and $\beta$.
Note that this identity retains the whole history of the trajectory.

\paragraph{Step 3: the common past.}
Let $u\le T$. Then $Z_b(u)=X_u$ for every $b$, and the pair $(Y_b,X_u)$ has
the law of $(X_t,X_u)$. By \cref{prop:hard-bridge} and $|C_N|\le1$, we
have $\E|\gamma_b(u)-C(t,u)|^2\to0$ as $N\to\infty$, uniformly in $u\le T$.
The same holds for $\bar\gamma_b(u)$, which is the average of
$\gamma_{b'}(u)$ over $b'$. The function $x\mapsto x^{p-1}$ is
$(p-1)$-Lipschitz on $[-1,1]$, and the normalized trace has second moment at most $\Theta_t^2$. Using the
Cauchy--Schwarz inequality, we find that the part of the integral in
\eqref{sel:fork-exact-subtraction} over $u\le T$ is at most
\[
 \beta p(p-1)\,T\,\Theta_t
 \sup_{b\le K,\,u\le T}\big(\E|\gamma_b(u)-\bar\gamma_b(u)|^2\big)^{1/2}
\]
in absolute value. This tends to zero as $N\to\infty$ with $T$ fixed. We
use no estimate on the decay of the memory here. The contribution of the
past is small because, at a fixed time horizon, both overlaps concentrate
around the same deterministic value $C(t,u)$.

\paragraph{Step 4: the last $L$ units of time.}
Let $u=T+v$ with $0\le v\le L$, and write $\tau=L-v$. The pair
$(Y_b,Z_b(T+v))$ has the law of $(X_t,X_{T+v})$. By \cref{prop:hard-bridge},
the overlap $\gamma_b(T+v)$ converges in $L^2$ to $C(T+L,T+v)$ as
$N\to\infty$, uniformly in $v$. Along $T=T_n$, this tends to $c(L-v)=c(\tau)$, uniformly in
$v\le L$. For $b'\ne b$, the vectors $Y_{b'}$ and $Z_b(T+v)$ are the
endpoints of two branches at $T$ of lengths $L$ and $v$. By
\eqref{sel:fork-overlap-concentration}, the overlap
$\ip{Y_{b'}}{Z_b(T+v)}/N$ converges in $L^2$ to $c(L+v)=c(2L-\tau)$, in the
double limit $j\to\infty$ and then $n\to\infty$, uniformly in $v\le L$.
Averaging over $b'$, we find
\[
 \lim_{n\to\infty}\limsup_{j\to\infty}\sup_{0\le v\le L}
 \Big(\E|\gamma_b(T_n+v)-c(L-v)|^2
 +\E|\bar\gamma_b(T_n+v)-k_{K,L}(L-v)|^2\Big)=0.
\]
On this interval, we have $t-u\le L$, and the normalized response trace has second moment at most $\Theta_L^2$, uniformly in $T$. As in Step~3, using the Cauchy--Schwarz inequality, we may replace $\gamma_b^{p-1}-\bar\gamma_b^{p-1}$ in
\eqref{sel:fork-exact-subtraction} by the deterministic weight
\[
 \omega(\tau)=c(\tau)^{p-1}-k_{K,L}(\tau)^{p-1},\qquad\tau=t-u,
\]
with an error whose $\limsup_j$ tends to zero as $n\to\infty$.

It remains to pass to the limit in
$\beta p\int_T^t\bar r_N(t,u)\,\omega(t-u)\,\dd u$, where
$\bar r_N(t,u)=\E\tr U_b(t,u)/N$ is the averaged response of
\cref{lem:brown-force-response}, which does not depend on $b$.
\Cref{lem:brown-force-response} and \cref{prop:hard-bridge} imply that
$\int_{s_1}^{s_2}\bar r_N(t,u)\,\dd u\to\int_{s_1}^{s_2}R(t,u)\,\dd u$
for all $T\le s_1\le s_2\le t$. The signed measure $\bar r_N(t,u)\,\dd u$ on
$[T,t]$ has total variation
\[
 \int_T^t|\bar r_N(t,u)|\,\dd u
 \le\int_T^t\E\left|\frac{\tr U_b(t,u)}N\right|\,\dd u
 \le L\,\Theta_L.
\]
As $R\ge0$, the integral of $R(t,\cdot)$ over $[T,t]$ is the limit of the
masses of these signed measures and obeys the same bound. The weight $\omega$ is continuous on $[0,L]$ and
is the uniform limit of step functions. For a step function, the integral converges as $N\to\infty$, and the total variation bound controls the error
of the approximation, uniformly in $N$. We conclude that, for fixed $T$,
\[
 \lim_{N\to\infty}\int_T^t\bar r_N(t,u)\,\omega(t-u)\,\dd u
 =\int_T^tR(t,u)\,\omega(t-u)\,\dd u.
\]
Along $T=T_n$, we have $R(T_n+L,T_n+v)\to r(L-v)$ uniformly in $v\le L$, by
the translated convergence on the closed causal triangle. Moreover,
$e_N(t)\to e_\beta(T+L)$ by \cref{prop:energy-bridge}, and
$e_\beta(T_n+L)\to e_\infty$. Inserting these limits into
\eqref{sel:fork-exact-subtraction}, we obtain
\[
 \lim_{n\to\infty}\limsup_{j\to\infty}
 \left|\frac{\E H_{N_j}(m_{N_j,T_n})}{N_j}
 -e_\infty+\beta p\int_0^Lr(\tau)\,\omega(\tau)\,\dd\tau\right|=0,
\]
which is \eqref{sel:fork-center-energy-limit} by the definition
\eqref{sel:fork-center-energy}.

For the norm, we have
\[
 \frac{\norm{m_{N,T}}^2}N=\frac1K+\frac1{K^2}\sum_{b\ne b'}
 \frac{\ip{Y_b}{Y_{b'}}}N.
\]
Each overlap $\ip{Y_b}{Y_{b'}}/N$ with $b\ne b'$ has the law of
$Q_{N,T}(L,L)$, which tends to $c(2L)$ in $L^2$ in the double limit, by
\eqref{sel:fork-overlap-concentration}. There are $K(K-1)$ such terms, and
$K(K-1)/K^2=1-1/K$. This proves \eqref{sel:fork-center-concentration} and
completes the proof of part~\ref{item:fork-limits}.

\paragraph{Step 5: Lipschitz algorithms and the threshold.}
Fix $T$ and $\varepsilon>0$. We may assume that $\varepsilon<1$. Let
$\varepsilon_1=\varepsilon/(2\sqrt K)$. We use the hard--soft approximation
of \cite[Lemmas~3.1--3.2]{sellke2024}. We first check its hypotheses.
In \cite{sellke2024}, the Hamiltonian is our $H_N$, with an ordered tensor of
independent standard Gaussian coefficients and the Euclidean distance between
coefficient tensors, and the dynamics is
$\dd x_v=(\beta_{\mathrm S}\nabla_{\mathrm{sp}}H_N(x_v)
-(N-1)(2N)^{-1}x_v)\,\dd v+P_{x_v}\,\dd W_v$, with noise coefficient one.
With $\beta_{\mathrm S}=\beta/2$, $X_u=x_{2u}$, and $B_u=W_{2u}/\sqrt2$, this
is \eqref{eq:hard-sde}, as in the proof of \cref{prop:hard-bridge}, and our
horizon $t$ becomes $2t$. Lemma~3.1 there states that, for a fixed initial
point on $S_N$ and a Hamiltonian satisfying their derivative bounds (which
hold for $H_N$ with probability at least $1-e^{-c_1N}$, by
\cite[Proposition~1.11]{sellke2024}), the hard dynamics stays within
$\varepsilon_1\sqrt N$ of a soft spherical dynamics driven by the same
Brownian motion, on the whole horizon, with probability at least
$1-e^{-c_1N}$. Here the confinement parameter of the soft dynamics is chosen large, depending on $p$, $\beta$, $t$, and $\varepsilon_1$, and $c_1>0$ and the
threshold for $N$ depend on the same parameters. Lemma~3.2 there states that, for a fixed initial point and
almost every Brownian path, there exists a map $\ell$ from coefficient tensors
to $B_N$, Lipschitz with a constant depending only on $p$, $\beta$, $t$, and
$\varepsilon_1$, such that $(1+\varepsilon_1)\ell(J)$ equals the soft
endpoint with probability at least $1-e^{-c_1N}$, where the probability is
over both the disorder and the Brownian motion. The coefficient-Lipschitz extension uses Kirszbraun's theorem in the proof of
Lemma~3.2 of \cite{sellke2024}; footnote~2 there gives a constructive choice
that depends measurably on the Brownian path.

The construction of \cite{sellke2024} is for a fixed initial point. To
include the uniform initial point, we choose measurably an orthogonal matrix
$O_{x_0}$ such that $O_{x_0}(\sqrt N e_1)=x_0$. Rotating the tensor by
$(O_{x_0}^{\mathsf T})^{\otimes p}$ and the Brownian path by
$O_{x_0}^{\mathsf T}$, solving from $\sqrt Ne_1$, and rotating the output back
by $O_{x_0}$ gives the solution from $x_0$, because $H_N$ and
\eqref{eq:hard-sde} are equivariant under rotations. Conditionally on
$X_0=x_0$, the rotated tensor and Brownian path have the same joint law as
the original ones, and the tensor rotation is an isometry of
$(\R^N)^{\otimes p}$, which does not change Lipschitz constants.

We apply this to each concatenated path $Z_b$, whose driving Brownian
motion is a function of the seed $\zeta_N$. We obtain maps
$\ell_b(J)=\ell_b(J,\zeta_N)\in B_N$, $b\le K$, which are Lipschitz in $J$
with a constant $L_{\mathrm{Lip},1}$ depending only on $p$, $\beta$, $t$, and
$\varepsilon_1$. Let $\mathcal G_N$ be the event on which $H_N$ satisfies the
derivative bounds of \cite{sellke2024} and, for every $b\le K$, the
conclusions of both lemmas hold for the branch $b$. On $\mathcal G_N$, we
have $\norm{(1+\varepsilon_1)\ell_b(J)-Y_b}\le\varepsilon_1\sqrt N$ for every
$b\le K$, and a union bound gives $\Prob(\mathcal G_N^c)\le3Ke^{-c_1N}$. The
union bound does not require the branches to be independent. We define
\begin{equation}\label{sel:fork-projected-algorithm}
 \mathcal A^\varepsilon_{N,T}(J)
 =\Pi_{B_N}\left(\frac{1+\varepsilon_1}{\sqrt{Q_{K,L}}}
 \frac1K\sum_{b=1}^K\ell_b(J)\right).
\end{equation}
Since $\Pi_{B_N}$ is $1$-Lipschitz and $Q_{K,L}\ge1/K$, this is an
$L_{\mathrm{Lip}}$-Lipschitz algorithm with seed $\zeta_N$ and
$L_{\mathrm{Lip}}=2L_{\mathrm{Lip},1}\sqrt K$. On $\mathcal G_N$, the distance between
$\mathcal A^\varepsilon_{N,T}$ and $\Pi_{B_N}(m_{N,T}/\sqrt{Q_{K,L}})$ is at
most $\varepsilon_1\sqrt N/\sqrt{Q_{K,L}}\le\varepsilon_1\sqrt{KN}
=\varepsilon\sqrt N/2$; on
$\mathcal G_N^c$, it is at most $2\sqrt N$. Using
$\Prob(\mathcal G_N^c)\le3Ke^{-c_1N}$, we find
\[
 \frac1N\E\norm{\mathcal A^\varepsilon_{N,T}
 -\Pi_{B_N}(m_{N,T}/\sqrt{Q_{K,L}})}^2
 \le\frac{\varepsilon^2}4+12Ke^{-c_1N},
\]
which proves \eqref{sel:fork-lipschitz}.

Next we compare the three normalizations. Let
$\nu_N=\norm{m_{N,T}}^2/N$. For $x\ge0$, we have
\begin{equation}\label{eq:sqrt-comparison}
 \left|\sqrt{x/Q_{K,L}}-1\right|
 =\frac{|x-Q_{K,L}|}{\sqrt{Q_{K,L}}\,(\sqrt x+\sqrt{Q_{K,L}})}
 \le\frac{|x-Q_{K,L}|}{Q_{K,L}}.
\end{equation}
The scaled distance between $m_{N,T}/\sqrt{Q_{K,L}}$ and its projection
onto $B_N$ is $(\sqrt{\nu_N/Q_{K,L}}-1)_+$. If $m_{N,T}\ne0$, the scaled
distance between the projection and $\widehat m_{N,T}$ is
$(1-\sqrt{\nu_N/Q_{K,L}})_+$, and both vectors vanish if $m_{N,T}=0$. By
\eqref{eq:sqrt-comparison} and $Q_{K,L}\ge q$, both distances are at most
$|\nu_N-Q_{K,L}|/q$. Then \eqref{sel:fork-center-concentration} implies that
their mean squares tend to zero as $j\to\infty$ and then $n\to\infty$. This completes
the proof of part~\ref{item:fork-lipschitz}.

Finally, we prove part~\ref{item:fork-ceiling}. Since
$\mathcal A^\varepsilon_{N,T}$ is an $L_{\mathrm{Lip}}$-Lipschitz algorithm with seed
$\zeta_N$, \cref{sel:absolute-field-cost}, applied with any
$\zeta_N$-measurable direction and with the nonnegative overlap term
dropped, gives
\[
 \limsup_{N\to\infty}\frac{\E H_N(\mathcal A^\varepsilon_{N,T})}N\le\Eth.
\]
This is the pure-model threshold, the case $\kappa=0$ of
\cref{sel:published-field-threshold}, averaged over the seed. Next we
compare energies. Note that $m_{N,T}/\sqrt{Q_{K,L}}$ has norm at most
$\sqrt{N/q}$. Using \eqref{eq:general-energy-transfer} with $R=q^{-1/2}$ and
the homogeneity $H_N(m_{N,T}/\sqrt{Q_{K,L}})=Q_{K,L}^{-p/2}H_N(m_{N,T})$, we
find
\[
 Q_{K,L}^{-p/2}\frac{\E H_N(m_{N,T})}N
 \le\frac{\E H_N(\mathcal A^\varepsilon_{N,T})}N
 +p\,q^{-(p-1)/2}Q_*\,\Delta_{N,T},
\]
where
\[
 \Delta_{N,T}=\left(\frac1N\E\norm{\mathcal A^\varepsilon_{N,T}
 -m_{N,T}/\sqrt{Q_{K,L}}}^2\right)^{1/2}.
\]
By \eqref{sel:fork-lipschitz} and the comparison of normalizations, we have
\[
 \limsup_{n\to\infty}\limsup_{j\to\infty}\Delta_{N_j,T_n}\le\varepsilon.
\]
We take
$j\to\infty$ with $T=T_n$ fixed, then $n\to\infty$, and finally
$\varepsilon\downarrow0$, in the order \eqref{sel:aging-limit-order}. By
\eqref{sel:fork-center-energy-limit}, we obtain
$Q_{K,L}^{-p/2}h_{K,L}\le\Eth$, which is \eqref{sel:fork-center-threshold}.
The constants of the algorithms depend on $T$, but we use them only at a
fixed $T$, before the limit $n\to\infty$.
\end{proof}

We next record the limits of the fork parameters as $K\to\infty$ and then
$L\to\infty$. Recall that $d_L=c(2L)$ and $h_L$ are given by
\eqref{sel:prefix-norm} and \eqref{sel:prefix-energy}, with the cluster
energy $e_\infty$, and that $h$ is the quantity in
\eqref{sel:limiting-center-energy}.

\begin{lemma}[The fork parameters as $K\to\infty$ and $L\to\infty$]
\label{lem:fork-parameter-limits}
Let $(c,r,\ell)$ be a stationary cluster with plateau $q\ge0$ and energy
$e_\infty$.
\begin{enumerate}[label=(\roman*)]
\item\label{item:fork-K-limit}
For every $L\ge0$ and $K\ge1$,
\begin{equation}\label{eq:fork-K-limit}
 0\le Q_{K,L}-d_L\le\frac1K,\qquad
 |h_{K,L}-h_L|\le\frac{\beta p(p-1)}K.
\end{equation}
\item\label{item:fork-L-limit}
As $L\to\infty$, we have $d_L\to q$ and
$h_L\to h=e_\infty-\beta\{1-q^p-pq^{p-1}(1-q)\}$.
\item\label{item:fork-saturation}
If $q>0$ and $\beta^2p(p-1)q^{p-2}(1-q)^2=1$, then
\begin{equation}\label{sel:fork-saturation}
 \lim_{L\to\infty}\lim_{K\to\infty}\frac{h_{K,L}}{Q_{K,L}^{p/2}}=\Eth.
\end{equation}
\end{enumerate}
\end{lemma}

\begin{proof}
For part~\ref{item:fork-K-limit}, we have
$Q_{K,L}-d_L=(1-c(2L))/K\in[0,1/K]$. By \eqref{sel:prefix-energy} and
\eqref{sel:fork-center-energy},
\[
 h_{K,L}-h_L=\beta p\int_0^Lr(\tau)
 \left\{k_{K,L}(\tau)^{p-1}-c(2L-\tau)^{p-1}\right\}\dd\tau.
\]
The integrand is dominated by $r(\tau)$ times a quantity that tends to zero
uniformly. Indeed, $k_{K,L}(\tau)-c(2L-\tau)=(c(\tau)-c(2L-\tau))/K$, and both
$k_{K,L}(\tau)$ and $c(2L-\tau)$ lie in $[0,1]$. By the mean value theorem,
\[
 \left|k_{K,L}(\tau)^{p-1}-c(2L-\tau)^{p-1}\right|\le\frac{p-1}K.
\]
Together with $\int_0^Lr\le1$, this proves \eqref{eq:fork-K-limit}.

Part~\ref{item:fork-L-limit} is part of \cref{lem:center-energy}. For
completeness, we recall the computation. Using $r=-c'$, we find that
$\int_0^Lr(\tau)c(\tau)^{p-1}\,\dd\tau=(1-c(L)^p)/p$. As $c$ is
nonincreasing, we have $c(2L-\tau)\in[q,c(L)]$ for $0\le\tau\le L$, and $\int_0^Lr=1-c(L)$. Thus the integral of
$r(\tau)c(2L-\tau)^{p-1}$ over $[0,L]$ lies between $q^{p-1}(1-c(L))$ and $c(L)^{p-1}(1-c(L))$. Both bounds
tend to $q^{p-1}(1-q)$. Inserting these two integrals into
\eqref{sel:prefix-energy}, we find that $h_L$ tends to
$e_\infty+\beta pq^{p-1}(1-q)-\beta(1-q^p)=h$.

For part~\ref{item:fork-saturation}, part~\ref{item:fork-K-limit} and
$d_L\ge q>0$ imply that $h_{K,L}/Q_{K,L}^{p/2}\to h_L/d_L^{p/2}$ as $K\to\infty$,
and part~\ref{item:fork-L-limit} gives $h_L/d_L^{p/2}\to h/q^{p/2}$ as
$L\to\infty$. By the equality case of \eqref{sel:center-amgm} in
\cref{lem:center-energy}, the marginality hypothesis implies that
$h/q^{p/2}=\Eth$.
\end{proof}

Part~\ref{item:fork-saturation} uses marginality, which is proved in
\cref{sel:positive-marginality} from part~\ref{item:threshold-ceiling} of
\cref{prop:threshold-inputs}. The proof of that part in
\cref{sec:general-proof-inputs} uses only \cref{sel:finite-forks} and
part~\ref{item:fork-K-limit}, and there is no circularity.

\section{Adaptive disclosure and orthogonality to early queries}
\label{sec:general-disclosure}

The overlap penalty of \cref{sel:absolute-field-cost} applies to directions
independent of $J$. At positive times, however, an earlier state $X_s$ depends on $J$. The goal of this section is to prove that an algorithm that saturates the threshold is almost
orthogonal to every subspace explored by a fixed finite number of gradient
queries. We condition $J$ on the information revealed by the queries. This is adaptive Gaussian conditioning, in the line of Bolthausen's
iterative construction of solutions of the TAP equations
\cite[Section~3]{bolthausen2014}. Bayati and Montanari used the same method
in their analysis of approximate message passing
\cite[arXiv version~4, Section~3.8, Lemma~11]{bayatimontanari2011}, and
they credit it to Bolthausen. The underlying facts on Gaussian vectors and
their orthogonal decompositions are classical; see \cite[Chapter~9, Theorem~9.1, pp.~127--128, and Theorems~9.8--9.9, pp.~131--132]{janson1997}. We use projection onto the span of the deterministic observations, which also permits linearly dependent observation families.
El Alaoui, Montanari, and Sellke conditioned the disorder tensor on adaptive
gradient observations in their analysis of message passing for spin
glasses \cite[arXiv version~1, Appendix~A, Lemma~A.2, and Theorem~7(a)]{eams2021}, and Huang and Sellke extended this conditioning to
multi-species spherical models \cite[arXiv version~1, Appendix~A, Lemma~A.2, and Theorem~3(a)]{hs2024amp}. The displayed inverse-Gram formulas in these AMP analyses use rank and conditioning hypotheses. Our projection-based disclosure below also covers dependent queries. Here we disclose all tensor blocks needed to
determine gradients throughout the explored span. We combine this disclosure with the external-field threshold to prove forgetting of a fixed
finite history.

\paragraph{Outer parameters.}
Let $\theta$ be a finite list of outer parameters with a prescribed order of
limits; in \cref{sec:general-transfer}, $\theta=(n,\varepsilon,K,L)$, in the
order \eqref{sel:aging-limit-order}. For nonnegative quantities
$F_{N,\theta}$, we write $\lim_\theta\limsup_NF_{N,\theta}=0$ if the upper
limit in $N$, followed by the upper limits in the outer parameters in their
prescribed order, vanishes. We do not assume that any intermediate limit
exists.

\paragraph{Standing hypotheses.}
For each $N$ and $\theta$, let $\zeta_{N,\theta}$ be a random variable independent of $J$, which contains all
auxiliary randomness used both by the algorithms and by the disclosure
procedure below. Its measurable space and law may depend on $\theta$;
below we write $\zeta_N$ with $\theta$ suppressed. For each $N$ and $\theta$, let
$\mathcal A_{N,\theta}=\mathcal A_{N,\theta}(J,\zeta_N)$ be an
$L_{\mathrm{Lip},\theta}$-Lipschitz algorithm with seed $\zeta_N$. The constant
$L_{\mathrm{Lip},\theta}$ is deterministic and independent of $N$, but it may be
unbounded in the outer limits. We define the energy deficit
\begin{equation}\label{sel:adaptive-deficit}
 \delta_\theta=\max\left\{0,\;\Eth-\liminf_{N\to\infty}
 \E\frac{H_N(\mathcal A_{N,\theta})}N\right\}.
\end{equation}
\Cref{sel:absolute-field-cost} implies that
$\limsup_N\E H_N(\mathcal A_{N,\theta})/N\le\Eth$. The hypothesis
$\lim_\theta\delta_\theta=0$ says exactly that the algorithms saturate the
threshold in mean.

For a subspace $W\subset\R^N$, we write
\[
 \mathcal T(W)=(W^\perp)^{\otimes p}\subset(\R^N)^{\otimes p},
\]
and $P_V$ for the orthogonal projection onto a subspace $V$ of $\R^N$ or of
$(\R^N)^{\otimes p}$. If $W\subset W'$, then $\mathcal T(W')\subset\mathcal T(W)$,
and we write $\mathcal T(W)\ominus\mathcal T(W')$ for the orthogonal
complement of $\mathcal T(W')$ in $\mathcal T(W)$. We choose orthonormal
bases of measurable subspaces measurably, by Gram--Schmidt in a fixed order
with deterministic fallbacks.

\begin{lemma}[Adaptive tensor disclosure]\label{sel:adaptive-disclosure}
Let $W_0\subset\R^N$ be a $\zeta_N$-measurable subspace, and let
$\mathcal F_0$ be the $\sigma$-algebra generated by $\zeta_N$ and
$P_{\mathcal T(W_0)^\perp}J$. For $j\ge0$, given $W_j$ and $\mathcal F_j$,
let $u_j$ be an $\mathcal F_j$-measurable vector that is either zero or a
unit vector in $W_j^\perp$, and set
\[
 W_{j+1}=W_j+\Span(u_j),\qquad
 \mathcal F_{j+1}=\sigma\Big(\mathcal F_j,\;
 P_{\mathcal T(W_j)\ominus\mathcal T(W_{j+1})}J\Big).
\]
Then, for every $j$:
\begin{enumerate}[label=(\roman*)]
\item\label{item:disclosure-gaussian}
$P_{\mathcal T(W_j)^\perp}J$ is $\mathcal F_j$-measurable, and conditionally
on $\mathcal F_j$, the tensor $P_{\mathcal T(W_j)}J$ is standard Gaussian on
$\mathcal T(W_j)$;
\item\label{item:disclosure-gradient}
if $x_j\in W_j$ is $\mathcal F_j$-measurable, then $\nabla H_N(x_j)$ is
$\mathcal F_j$-measurable. In particular, $u_j$ may be chosen as the
normalized projection of $\nabla H_N(x_j)$ onto $W_j^\perp$, and as zero if
that projection vanishes.
\end{enumerate}
\end{lemma}

\begin{proof}
The conditional Gaussian assertions are understood almost surely with respect to the disclosed history. A zero query discloses the zero tensor block and leaves the residual law unchanged.
We prove part~\ref{item:disclosure-gaussian} by induction on $j$. For
$j=0$, conditionally on $\zeta_N$, the tensor $J$ is standard Gaussian,
because it is independent of $\zeta_N$, and the subspaces $\mathcal T(W_0)$
and $\mathcal T(W_0)^\perp$ are fixed. The projections of a standard
Gaussian vector onto two complementary orthogonal subspaces are independent
standard Gaussian vectors on them. Revealing $P_{\mathcal T(W_0)^\perp}J$
leaves $P_{\mathcal T(W_0)}J$ standard Gaussian on
$\mathcal T(W_0)$.

Suppose that part~\ref{item:disclosure-gaussian} holds for $j$. Conditionally on
$\mathcal F_j$, the vector $u_j$ and the subspaces in the orthogonal
decomposition
\[
 \mathcal T(W_j)=\mathcal T(W_{j+1})\oplus
 \big(\mathcal T(W_j)\ominus\mathcal T(W_{j+1})\big)
\]
are fixed, and $P_{\mathcal T(W_j)}J$ is standard Gaussian on
$\mathcal T(W_j)$. Its projections onto the two summands are
independent standard Gaussian vectors. Revealing the second leaves the
first, which is $P_{\mathcal T(W_{j+1})}J$, standard Gaussian on
$\mathcal T(W_{j+1})$. Since
$\mathcal T(W_{j+1})^\perp=\mathcal T(W_j)^\perp\oplus
(\mathcal T(W_j)\ominus\mathcal T(W_{j+1}))$, the projection
$P_{\mathcal T(W_{j+1})^\perp}J$ is $\mathcal F_{j+1}$-measurable. This
completes the induction. At each step, we choose the new direction before we reveal the
corresponding part of $J$; this is what makes the induction work, and we
never condition on a subspace that depends on undisclosed randomness.

For part~\ref{item:disclosure-gradient}, let $x\in W_j$ and $y\in\R^N$. We
have
\[
 \ip{\nabla H_N(x)}{y}=N^{-(p-1)/2}\Big\langle J,\;\sum_{i=1}^p
 x^{\otimes(i-1)}\otimes y\otimes x^{\otimes(p-i)}\Big\rangle.
\]
Since $p\ge2$, each tensor in the sum has at least one factor $x\in W_j$.
It is orthogonal to every product tensor with all factors in $W_j^\perp$,
and these product tensors span $\mathcal T(W_j)$. The right side is equal to
$N^{-(p-1)/2}$ times the inner product of $P_{\mathcal T(W_j)^\perp}J$ with
that tensor. By part~\ref{item:disclosure-gaussian}, it is
$\mathcal F_j$-measurable when $x=x_j$ is.
\end{proof}

\begin{proposition}[Orthogonality to every fixed finite query span]
\label{sel:adaptive-orthogonality}
Assume the standing hypotheses and $\lim_\theta\delta_\theta=0$. Fix integers
$k_0\ge0$ and $M\ge0$. For each $N$ and $\theta$, let
$W_0\subset W_1\subset\dots\subset W_M$ be subspaces constructed as in
\cref{sel:adaptive-disclosure} from the seed $\zeta_N$, with
$\dim W_0\le k_0$. Then
\begin{equation}\label{sel:finite-span-orthogonality}
 \lim_\theta\limsup_{N\to\infty}
 \E\frac{\norm{P_{W_M}\mathcal A_{N,\theta}}^2}N=0.
\end{equation}
In particular, if $Y_N\in W_M\cap B_N$ is a random vector, then
$\lim_\theta\limsup_N\E|\ip{\mathcal A_{N,\theta}}{Y_N}/N|^2=0$.
\end{proposition}

\begin{proof}
For $0\le j\le M$, we define
\[
 \psi_{j,\theta}=\limsup_{N\to\infty}
 \E\frac{\norm{P_{W_j}\mathcal A_{N,\theta}}^2}N.
\]
We bound $\psi_{0,\theta}$ directly and then show by induction that
$\psi_{j+1,\theta}$ is controlled by $\psi_{j,\theta}$ and $\delta_\theta$.
We use the constants of the threshold only at fixed $\theta$, before any
outer limit, and the argument does not require $L_{\mathrm{Lip},\theta}$ to stay
bounded.

\paragraph{Step 1: the initial span.}
Let $e_1,\dots,e_{k_0}$ be a $\zeta_N$-measurable orthonormal basis of
$W_0$, completed by zero vectors if $\dim W_0<k_0$. We obtain such a frame by Gram--Schmidt on the projected coordinate vectors,
keeping nonzero columns in a fixed order and then padding by zero columns. For $N\ge1$, fix a
deterministic unit vector $e_{\mathrm{ref}}\in\R^N$, and for every $i$ set
$w_i=\sqrt N e_i$ on $\{e_i\ne0\}$ and
$w_i=\sqrt N e_{\mathrm{ref}}$ on $\{e_i=0\}$. This gives a
$\zeta_N$-measurable direction of norm $\sqrt N$ on every seed event.
We apply \cref{sel:absolute-field-cost} with
$L_{\mathrm{Lip}}=L_{\mathrm{Lip},\theta}$ and direction $w_i$, and we drop
its nonnegative overlap penalty on $\{e_i=0\}$. Since the pairing with
$e_i$ vanishes there, this yields
\[
 \bd\limsup_{N\to\infty}\E\left|\frac{\ip{\mathcal A_{N,\theta}}{e_i}}{\sqrt N}\right|^p
 \le\Eth-\liminf_{N\to\infty}\E\frac{H_N(\mathcal A_{N,\theta})}N
 \le\delta_\theta.
\]
For each $i$, by Jensen's inequality and $p\ge2$, we have
$\E|\ip{\mathcal A_{N,\theta}}{e_i}|^2/N
\le(\E|\ip{\mathcal A_{N,\theta}}{e_i}/\sqrt N|^p)^{2/p}$. Summing over $i$,
we obtain
\begin{equation}\label{sel:adaptive-initial-estimate}
 \psi_{0,\theta}\le k_0\left(\frac{\delta_\theta}{\bd}\right)^{2/p}.
\end{equation}

\paragraph{Step 2: the energy of the residual.}
Fix $j<M$, write $W=W_j$, and let
$\mathcal A^\perp=P_{W^\perp}\mathcal A_{N,\theta}$. Both
$\mathcal A_{N,\theta}$ and $\mathcal A^\perp$ lie in $B_N$. Using \eqref{sel:tensor-ui} and the mean value theorem, we have
\[
 \frac{|H_N(\mathcal A_{N,\theta})-H_N(\mathcal A^\perp)|}N
 \le pQ_N\frac{\norm{P_W\mathcal A_{N,\theta}}}{\sqrt N}.
\]
By the Cauchy--Schwarz inequality, the expectation of the right side has
upper limit at most $pQ_*\sqrt{\psi_{j,\theta}}$. Combining this with the
definition \eqref{sel:adaptive-deficit} of $\delta_\theta$, we find
\begin{equation}\label{eq:residual-energy-lower}
 \liminf_{N\to\infty}\E\frac{H_N(\mathcal A^\perp)}N
 \ge\Eth-\delta_\theta-pQ_*\sqrt{\psi_{j,\theta}}.
\end{equation}

\paragraph{Step 3: the residual model.}
Since $k_0$ and $M$ are fixed before the dimension limit, we may restrict
to $N>k_0+M$. Every residual dimension below is then positive.
We now show that $\mathcal A^\perp$ is the output of a Lipschitz algorithm
for a pure $p$-spin model in dimension $N'=N-\dim W$, and that its
energy has the same normalization up to a factor tending to one. We
condition on $\mathcal F_j$. Then $W$ and $u_j$ are fixed, and we choose an
$\mathcal F_j$-measurable orthonormal basis $f_1,\dots,f_{N'}$ of
$W^\perp$. Let $J'\in(\R^{N'})^{\otimes p}$ be the coordinates of
$P_{\mathcal T(W)}J$ in the basis
$(f_{i_1}\otimes\dots\otimes f_{i_p})$. By
\cref{sel:adaptive-disclosure}\ref{item:disclosure-gaussian}, $J'$ is a
standard Gaussian vector in the $N'^p$-dimensional space
$(\R^{N'})^{\otimes p}$, conditionally on $\mathcal F_j$. Let
$y'\in\R^{N'}$ be the coordinates of $\mathcal A^\perp$ in the basis
$(f_i)$, and set $x=\sqrt{N'/N}\,y'$. Then $\norm x\le\sqrt{N'}$, that is,
$x\in B_{N'}$. Since $(\mathcal A^\perp)^{\otimes p}\in\mathcal T(W)$, only
$P_{\mathcal T(W)}J$ contributes to $H_N(\mathcal A^\perp)$, and a change of
normalization yields
\begin{equation}\label{sel:residual-normalization}
 \frac{H_N(\mathcal A^\perp)}N
 =N^{-(p+1)/2}\ip{J'}{y'^{\otimes p}}
 =\sqrt{\frac{N'}N}\,\frac{H'_{N'}(x)}{N'},
\end{equation}
where $H'_{N'}(x)=N'^{-(p-1)/2}\ip{J'}{x^{\otimes p}}$ is the pure $p$-spin
Hamiltonian in dimension $N'$ with coefficient tensor $J'$. Indeed,
$\ip{J'}{y'^{\otimes p}}=(N/N')^{p/2}\ip{J'}{x^{\otimes p}}$, and
$N^{-(p+1)/2}(N/N')^{p/2}N'^{(p-1)/2}=(N'/N)^{1/2}N'^{-1}$. Moreover, $J$
is an affine isometric function of $J'$ given $\mathcal F_j$, and
$\mathcal A^\perp\mapsto x$ is a contraction. Therefore, conditionally on $\mathcal F_j$, the map $J'\mapsto x$ is $L_{\mathrm{Lip},\theta}$-Lipschitz with
values in $B_{N'}$. Let $u'\in\R^{N'}$ be the coordinates of $u_j$ when
$u_j\ne0$, and the first coordinate vector otherwise. The direction
$w'=\sqrt{N'}u'$ has norm $\sqrt{N'}$, and when $u_j\ne0$,
\[
 \frac{\ip{w'}{x}}{N'}=\frac{\ip{u'}{y'}}{\sqrt N}
 =\frac{\ip{u_j}{\mathcal A_{N,\theta}}}{\sqrt N},
\]
where we used $u_j\in W^\perp$ in the last equality.

\paragraph{Step 4: the penalty in the residual model.}
Let $\alpha_j=\ip{u_j}{\mathcal A_{N,\theta}}/\sqrt N$, which is zero if
$u_j=0$. The conditional law of $J'$ given $\mathcal F_j$ is standard
Gaussian, and on each event $\{N'=n'\}$ the tensor $J'$ is independent of
$\mathcal F_j$. We may therefore apply \eqref{eq:conditional-field-cost} in
dimension $N'$, with $\mathcal F_j$ in the role of the seed, to the
algorithm $J'\mapsto x$ and the direction $w'$. If $u_j\ne0$, its overlap
term is $\bd|\alpha_j|^p$; if $u_j=0$, we drop its nonnegative overlap
term. Since $N'\ge N-k_0-j$ and $\phi_N\ge0$, we obtain
\[
 \E\left[\frac{H'_{N'}(x)}{N'}+\bd|\alpha_j|^p\,\middle|\,\mathcal F_j\right]
 \le\Eth+\sup_{N''\ge N-k_0-j}\phi_{N''}(p,L_{\mathrm{Lip},\theta}).
\]
We use \eqref{sel:residual-normalization} to write
\[
 \frac{H_N(\mathcal A^\perp)}N+\bd|\alpha_j|^p
 =\sqrt{\frac{N'}N}\left(\frac{H'_{N'}(x)}{N'}+\bd|\alpha_j|^p\right)
 +\left(1-\sqrt{\frac{N'}N}\right)\bd|\alpha_j|^p.
\]
Here $\sqrt{N'/N}\le1$, $0\le1-\sqrt{N'/N}\le(k_0+j)/N$, and $|\alpha_j|\le1$.
Taking expectations and letting $N\to\infty$, we find
\[
 \limsup_{N\to\infty}\left\{\E\frac{H_N(\mathcal A^\perp)}N
 +\bd\,\E|\alpha_j|^p\right\}\le\Eth.
\]
Combining this with \eqref{eq:residual-energy-lower}, we obtain
$\bd\limsup_N\E|\alpha_j|^p\le\delta_\theta+pQ_*\sqrt{\psi_{j,\theta}}$. By
Jensen's inequality as in Step~1, and since
$\norm{P_{W_{j+1}}\mathcal A_{N,\theta}}^2=\norm{P_{W_j}\mathcal A_{N,\theta}}^2
+N|\alpha_j|^2$, we conclude that
\begin{equation}\label{sel:adaptive-recursion}
 \psi_{j+1,\theta}\le\psi_{j,\theta}
 +\left(\frac{\delta_\theta+pQ_*\sqrt{\psi_{j,\theta}}}{\bd}\right)^{2/p}.
\end{equation}

\paragraph{Step 5: the outer limits.}
The right sides of \eqref{sel:adaptive-initial-estimate} and
\eqref{sel:adaptive-recursion} are continuous nondecreasing functions of
$(\delta_\theta,\psi_{j,\theta})$ that vanish at $(0,0)$. Upper limits pass
through such functions. Taking the outer limits in their prescribed order
and inducting on $j\le M$, we find
$\lim_\theta\psi_{j,\theta}=0$ for every $j\le M$. This is
\eqref{sel:finite-span-orthogonality}. Finally, if
$Y_N\in W_M\cap B_N$, then
\[
 \left|\frac{\ip{\mathcal A_{N,\theta}}{Y_N}}N\right|^2
 =\left|\frac{\ip{P_{W_M}\mathcal A_{N,\theta}}{Y_N}}N\right|^2
 \le\frac{\norm{P_{W_M}\mathcal A_{N,\theta}}^2}N,
\]
which proves the last assertion.
\end{proof}

\section{Transfer to a fixed waiting-time state}\label{sec:general-transfer}

We now combine the forks with finite-query orthogonality on their fixed
product probability spaces. The hypotheses of the next proposition are those of
\cref{prop:threshold-inputs}\ref{item:threshold-forgetting}. Under them,
the translated family $(C_T,R_T,\lambda_T)$ converges along every sequence
$T_n\to\infty$ to one stationary cluster $(c,r,\ell)$. Indeed, by the
precompactness in \cref{prop:stationary-clusters}, every subsequence of
$(T_n)$ has a further subsequence along which the family converges, and the
limit is a stationary cluster whose kernels are $(c,r)$. Its multiplier is
determined by $r$: setting $\tau=0$ in \eqref{eq:stationary-cluster-response}
and using $r(0)=1$, we find $\ell=-r'(0)$. All subsequential limits are
therefore equal. In particular, \cref{sel:finite-forks} applies along every
sequence $T_n\to\infty$.

\begin{proposition}[Fixed-waiting-time decorrelation]
\label{sel:general-fixed-wait}
Fix $p\ge3$ and $\beta>0$. Suppose that the translated kernels converge
along every sequence $T\to\infty$ to one stationary profile $(c,r)$, and
that its plateau $q$ is positive and satisfies
$\beta^2p(p-1)q^{p-2}(1-q)^2=1$. Then
\[
 \lim_{t\to\infty}C(t,s)=0\qquad\text{for every fixed }s\ge0.
\]
\end{proposition}

\begin{proof}
The proof is by contradiction. Recall from \eqref{eq:ck-physical-bounds}
that $C\ge0$. Suppose that, for some $s\ge0$, there exist
$\gamma>0$ and $t_n\to\infty$ such that $C(t_n,s)\ge\gamma$ for all $n$. The
proof has three steps. We first express $C(t_n,s)$ through the center of a
fork that ends at $t_n$. Then we show that the Lipschitz algorithms
approximating the normalized center saturate the threshold in the order
\eqref{sel:aging-limit-order}. Finally, on the fixed fork product law, we approximate $X_s$ by a
finite-query Euler vector and use \cref{sel:adaptive-orthogonality} to force
their overlap with $X_s$ to vanish, which contradicts the first step.

\paragraph{Step 1: the overlap of the center.}
Fix $K\ge1$ and $L\ge0$, and set $T_n=\max\{s+1,t_n-L\}$ for every $n$.
Then $T_n\to\infty$, $T_n\ge s$, and for all sufficiently large $n$ we have
$T_n+L=t_n$. Let $m_{N,n}$ be the center of the fork at time $T_n$ with $K$
branches of length $L$. For these sufficiently large $n$, the branches
end at time $t_n$. The time $s$ lies in the common past,
and each pair $(Y_b,X_s)$ has the law of $(X_{t_n},X_s)$. This implies that
\begin{equation}\label{sel:fork-fixed-wait-overlap}
 \E\frac{\ip{m_{N,n}}{X_s}}N=\E C_N(t_n,s),
\end{equation}
and by \cref{prop:hard-bridge}, the right side tends to $C(t_n,s)$ as
$N\to\infty$. In the limit in $n$, the endpoint of the fork is the time $t_n$ itself,
so we do not need any regularity of $C(\cdot,s)$ in its first variable.

\paragraph{Step 2: saturation of the threshold.}
Let $\theta=(n,\varepsilon,K,L)$, where $\varepsilon=1/(m+1)$ and
$K=k+1$ for $m,k\in\mathbb N$, with the order of
\eqref{sel:aging-limit-order}, and let
$\mathcal A_{N,\theta}=\mathcal A^\varepsilon_{N,T_n}$ be the algorithms of
\cref{sel:finite-forks}\ref{item:fork-lipschitz} for this fork. They satisfy
the standing hypotheses of \cref{sec:general-disclosure} on the fixed fork
product law of \cref{sec:general-forks}, with seed $\zeta_N$ consisting
of $X_0$, the $K$ full future Wiener paths, and the full trunk Wiener path $B$,
jointly independent of $J$, and with a constant $L_{\mathrm{Lip},\theta}$
depending only on $p$, $\beta$, the fixed waiting time $s$, the horizon
$T_n+L$, $K$, and $\varepsilon$. We write
$Q=Q_{K,L}$, and recall that $q\le Q\le1$. Let
\[
 \Delta_{N,\theta}=\left(\frac1N\E\norm{\mathcal A_{N,\theta}
 -m_{N,n}/\sqrt Q}^2\right)^{1/2}.
\]
By \eqref{sel:fork-lipschitz} and the comparison of normalizations in
\cref{sel:finite-forks}\ref{item:fork-lipschitz}, we have
\begin{equation}\label{eq:transfer-approximation}
 \limsup_{n\to\infty}\limsup_{j\to\infty}\Delta_{N_j,\theta}\le\varepsilon.
\end{equation}
The vector $m_{N,n}/\sqrt Q$ has norm at most $\sqrt{N/q}$. Using
\eqref{eq:general-energy-transfer} with $R=q^{-1/2}$ and the homogeneity of
$H_N$, we find
\[
 \E\frac{H_N(\mathcal A_{N,\theta})}N
 \ge Q^{-p/2}\E\frac{H_N(m_{N,n})}N-C_q\Delta_{N,\theta},\qquad
 C_q=p\,q^{-(p-1)/2}Q_*.
\]
The constant $C_q$ does not depend on $n$, $\varepsilon$, $K$, or $L$; this
is where we use $q>0$. Then \eqref{sel:fork-center-energy-limit} implies that
$\lim_n\liminf_jQ^{-p/2}\E H_{N_j}(m_{N_j,n})/N_j=Q^{-p/2}h_{K,L}$.
Together with \eqref{eq:transfer-approximation}, this gives
\[
 \limsup_{n\to\infty}\delta_\theta
 \le\max\left\{0,\;\Eth-\frac{h_{K,L}}{Q_{K,L}^{p/2}}+C_q\varepsilon\right\}.
\]
We let $\varepsilon\downarrow0$, then $K\to\infty$, then $L\to\infty$. By
\eqref{sel:fork-saturation}, which applies because $q$ is marginal, we
conclude that $\lim_\theta\delta_\theta=0$ in the order
\eqref{sel:aging-limit-order}.

\paragraph{Step 3: orthogonality on the fork product law and the contradiction.}
We use the algorithms of Step~2 on precisely the fork seed and disorder
product space just described, and take the limits in precisely the order
\eqref{sel:aging-limit-order}, with $\varepsilon=1/(m+1)$. All disclosure
histories below include the full fork seed. If $s=0$, then $X_0$ lies in the
seed-measurable subspace $W_0=\Span(X_0)$, and
\cref{sel:adaptive-orthogonality} with $k_0=1$ and no disclosure steps gives
zero limiting squared overlap. For $s>0$, fix $M\ge1$ before any limit.
We construct the Euler--Maruyama approximation
\cite[Sections~4--5]{higham2001} of this fork's trunk state $X_s$ and prove
the required normalized error directly.

\subparagraph{The coefficients.}
Let $\Psi_N=\Pi_{B_N}$ be the metric projection onto the ball of radius $\sqrt N$.
We use $P_y=I-yy^{\mathsf T}/N$ for every $y\in\R^N$, as in \cref{sec:model}.
For $x\in\R^N$, we define
\[
 b_N(x)=\beta P_{\Psi_N(x)}\nabla H_N(\Psi_N(x))-\frac{N-1}N\Psi_N(x),\qquad
 \sigma_N(x)=\sqrt2\,P_{\Psi_N(x)}.
\]
Since $\Psi_N(x)=x$ on $S_N$, these are the coefficients of
\eqref{eq:hard-sde} on $S_N$, and
\[
 X_t=X_0+\int_0^tb_N(X_r)\,\dd r+\int_0^t\sigma_N(X_r)\,\dd B_r.
\]
Let $G_N=\{Q_N\le A_p\}$, and let $y=\Psi_N(x)$ and $z=\Psi_N(x')$, so that
$\norm y,\norm z\le\sqrt N$ and $\norm{y-z}\le\norm{x-x'}$. The matrix $P_y$
has eigenvalues $1$ and $1-\norm y^2/N\in[0,1]$, which gives
$\norm{P_y}_\op\le1$ and $\norm{P_y}_{\HS}^2\le N$. Moreover,
\[
 \norm{P_y-P_z}_{\HS}=\frac{\norm{yy^{\mathsf T}-zz^{\mathsf T}}_{\HS}}N
 \le\frac{(\norm y+\norm z)\norm{y-z}}N\le\frac{2\norm{x-x'}}{\sqrt N}.
\]
On $G_N$, \eqref{eq:general-ball-derivatives} with $R=2$ gives
$\norm{\nabla H_N(y)}\le p2^{p-1}A_p\sqrt N$ and
$\norm{\nabla H_N(y)-\nabla H_N(z)}\le p(p-1)2^{p-2}A_p\norm{y-z}$. Writing
$P_y\nabla H_N(y)-P_z\nabla H_N(z)
=P_y(\nabla H_N(y)-\nabla H_N(z))+(P_y-P_z)\nabla H_N(z)$, we obtain, on
$G_N$ and for all $x,x'\in\R^N$,
\begin{equation}\label{eq:euler-coefficient-bounds}
\begin{gathered}
 \norm{b_N(x)-b_N(x')}\le K_b\norm{x-x'},\qquad
 \norm{\sigma_N(x)-\sigma_N(x')}_{\HS}\le\frac{2\sqrt2}{\sqrt N}\norm{x-x'},\\
 \norm{b_N(x)}^2\le K_b'N,\qquad\norm{\sigma_N(x)}_{\HS}^2\le2N,
\end{gathered}
\end{equation}
with $K_b=1+\beta pA_p2^{p-2}(3p+5)$ and
$K_b'=(3\beta p2^{p-1}A_p+2)^2$. The Hilbert--Schmidt bound on
$\sigma_N(x)-\sigma_N(x')$ is the reason that the error below does not grow with $N$: a bound of order one for each matrix entry would lose a factor of
$N$.

\subparagraph{The Euler--Maruyama scheme.}
Let $t_i=is/M$ for $0\le i\le M$. We define $Z_0=X_0$ and
\[
 Z_{i+1}=Z_i+\frac sM\,b_N(Z_i)+\sigma_N(Z_i)\big(B_{t_{i+1}}-B_{t_i}\big),
 \qquad 0\le i<M.
\]
For $r\in[t_i,t_{i+1})$, let $r_-=t_i$, and let
$Z_r=Z_i+(r-t_i)b_N(Z_i)+\sigma_N(Z_i)(B_r-B_{t_i})$, so that
\[
 Z_t=X_0+\int_0^tb_N(Z_{r_-})\,\dd r+\int_0^t\sigma_N(Z_{r_-})\,\dd B_r
\]
and $Z_s=Z_M$. We set $Y_{N,M}=\Pi_{B_N}(Z_M)$.

\subparagraph{The error on $G_N$.}
Fix $J\in G_N$ and $X_0$, and let $\E_B$ denote expectation over the
Brownian motion. Let $u(t)=\E_B\norm{X_t-Z_t}^2/N$ and
$v(r)=\E_B\norm{Z_r-Z_{r_-}}^2/N$. Since $Z_{r_-}$ is independent of
$B_r-B_{r_-}$, the It\^o isometry and \eqref{eq:euler-coefficient-bounds}
give
\[
 v(r)\le2\Big(\frac sM\Big)^2K_b'+2\,\frac sM\,\frac{2N}N
 \le C_1\frac sM,\qquad C_1=2sK_b'+4.
\]
Subtracting the integral equations for $X_t$ and $Z_t$, and using the
Cauchy--Schwarz inequality for the drift and the It\^o isometry for the
noise, we find, for $t\le s$,
\[
 u(t)\le2sK_b^2\int_0^t\E_B\frac{\norm{X_r-Z_{r_-}}^2}N\,\dd r
 +\frac{16}N\int_0^t\E_B\frac{\norm{X_r-Z_{r_-}}^2}N\,\dd r.
\]
Inserting the bound $\E_B\norm{X_r-Z_{r_-}}^2/N\le2u(r)+2v(r)$, we obtain
$u(t)\le C_2\int_0^tu(r)\,\dd r+C_2C_1s^2/M$ with $C_2=4sK_b^2+32$. By
Gronwall's inequality,
\[
 u(s)\le C_1C_2s\,e^{C_2s}\,\frac sM.
\]
The projection $\Pi_{B_N}$ is $1$-Lipschitz and fixes $X_s\in S_N$, which
gives
$\E_B\norm{X_s-Y_{N,M}}^2/N\le u(s)$ on $G_N$. On $G_N^c$, we have
$\norm{X_s-Y_{N,M}}^2/N\le4$, and $\Prob(G_N^c)\le e^{-c_pN}$. Thus, with $C=C_1C_2se^{C_2s}$,
\begin{equation}\label{sel:euler-approximation}
 \limsup_{N\to\infty}\E\frac{\norm{X_s-Y_{N,M}}^2}N\le C\,\frac sM.
\end{equation}
The expectation here is on the fixed fork product law. Its projection
onto $(J,X_0,B)$ has exactly the original disorder, uniform initial and
Wiener product law, so the error bound is independent of $n$, $\varepsilon$,
$K$, and $L$.

\subparagraph{The query span.}
Let $W_0$ be the span of $X_0$ and the increments $B_{t_{i+1}}-B_{t_i}$,
$0\le i<M$. It is measurable with respect to the full fork seed $\zeta_N$ and has dimension at
most $M+1$. Suppose that $Z_i\in W_i$ and that $Z_i$ is
$\mathcal F_i$-measurable, where $\mathcal F_i$ is the disclosure history
generated using this full fork seed, which holds for $i=0$. The point $\Psi_N(Z_i)$
is a multiple of $Z_i$ and lies in $W_i$. By \cref{sel:adaptive-disclosure}\ref{item:disclosure-gradient}, the gradient
$g_i=\nabla H_N(\Psi_N(Z_i))$ is $\mathcal F_i$-measurable. We choose $u_i$
as the normalized projection of $g_i$ onto $W_i^\perp$ (zero if it
vanishes), so that $g_i\in W_{i+1}$. Since
$P_yw=w-y\ip yw/N$, the step from $Z_i$ to $Z_{i+1}$ is a linear
combination of $Z_i$, $\Psi_N(Z_i)$, $g_i$, and $B_{t_{i+1}}-B_{t_i}$, with
coefficients that are $\mathcal F_i$-measurable. This shows that
$Z_{i+1}\in W_{i+1}$ and that $Z_{i+1}$ is $\mathcal F_{i+1}$-measurable. After $M$ steps, $Z_M\in W_M$,
and $Y_{N,M}$, a multiple of $Z_M$, lies in $W_M$. Each step adds at most
one dimension, which gives $\dim W_M\le2M+1$.

\subparagraph{Passing from the Euler vector to the waiting state.}
By Step~2, \cref{sel:adaptive-orthogonality} applies to these particular
full-fork disclosure histories, with $k_0=M+1$ and $M$ queries. Hence the
squared overlap of $\mathcal A_{N,\theta}$ with $Y_{N,M}$ vanishes after the
dimension limit and then the limits in $n$, $\varepsilon=1/(m+1)$, $K$, and $L$.
Since $\norm{\mathcal A_{N,\theta}}\le\sqrt N$, we have
\[
 \E\left|\frac{\ip{\mathcal A_{N,\theta}}{X_s}}N\right|^2
 \le2\E\left|\frac{\ip{\mathcal A_{N,\theta}}{Y_{N,M}}}N\right|^2
 +2\E\frac{\norm{X_s-Y_{N,M}}^2}N.
\]
By \eqref{sel:euler-approximation}, the nested upper limit of the left side
is at most $2Cs/M$. It does not depend on $M$. Sending $M\to\infty$ last
(and using the $s=0$ argument above when needed), we obtain
\begin{equation}\label{eq:transfer-orthogonality}
 \lim_\theta\limsup_{N\to\infty}
 \E\left|\frac{\ip{\mathcal A_{N,\theta}}{X_s}}N\right|^2=0.
\end{equation}
On the other hand, by \eqref{eq:fork-overlap-transfer} with $V_N=X_s$ and by
\eqref{sel:fork-fixed-wait-overlap}, we have
\[
 \frac{\E C_N(t_n,s)}{\sqrt Q}
 \le\left(\E\left|\frac{\ip{\mathcal A_{N,\theta}}{X_s}}N\right|^2\right)^{1/2}
 +\Delta_{N,\theta}.
\]
As $N=N_j\to\infty$, the left side tends to $C(t_n,s)/\sqrt Q$, which is at
least $\gamma$ because $Q\le1$. We conclude that, for every fixed
$\varepsilon$, $K$, and $L$ and all sufficiently large $n$,
\[
 \gamma\le\limsup_{j\to\infty}
 \left(\E\left|\frac{\ip{\mathcal A_{N_j,\theta}}{X_s}}{N_j}\right|^2\right)^{1/2}
 +\limsup_{j\to\infty}\Delta_{N_j,\theta}.
\]
We take the upper limit as $n\to\infty$, then let $\varepsilon\downarrow0$, $K\to\infty$, and $L\to\infty$, in the order \eqref{sel:aging-limit-order}.
By \eqref{eq:transfer-orthogonality}, the first term tends to zero, and by
\eqref{eq:transfer-approximation}, the second is at most $\varepsilon$ after
the limit in $n$. Hence $\gamma\le0$, which contradicts $\gamma>0$. This completes the proof.
\end{proof}

\section{Proof of \texorpdfstring{\cref{prop:threshold-inputs}}{the threshold inputs}}
\label{sec:general-proof-inputs}

We now prove both parts of \cref{prop:threshold-inputs} for every $p\ge3$.

\begin{proof}[Proof of \cref{prop:threshold-inputs}\ref{item:threshold-ceiling}]
Let $T_n\to\infty$ be a sequence along which the translated kernels converge
to a stationary cluster $(c,r,\ell)$ with plateau $q>0$, and fix $L\ge0$.
For every $K\ge1$, \cref{sel:finite-forks} applies along $T_n$, and
\eqref{sel:fork-center-threshold} gives
\[
 h_{K,L}\le\Eth\,Q_{K,L}^{p/2}.
\]
By \cref{lem:fork-parameter-limits}\ref{item:fork-K-limit}, we have
$|Q_{K,L}-d_L|\le1/K$ and $|h_{K,L}-h_L|\le\beta p(p-1)/K$. The second bound
is the dominated convergence argument made explicit: the difference
$h_{K,L}-h_L$ is $\beta p$ times the integral of $r(\tau)$ times
$k_{K,L}(\tau)^{p-1}-c(2L-\tau)^{p-1}$, which is at most $(p-1)/K$ in absolute
value, and $\int_0^Lr\le1$. Letting $K\to\infty$ and using the continuity of
$x\mapsto x^{p/2}$, we obtain
\[
 h_L\le\Eth\,d_L^{p/2},
\]
which is \eqref{eq:threshold-input-ceiling}.
\end{proof}

Letting $L\to\infty$ in \eqref{eq:threshold-input-ceiling} and using
\cref{lem:fork-parameter-limits}\ref{item:fork-L-limit}, we obtain
\begin{equation}\label{sel:general-positive-upper-bound}
 h\le\Eth\,q^{p/2},
\end{equation}
where $h$ is given by \eqref{sel:limiting-center-energy}. This is the form
in which part~\ref{item:threshold-ceiling} enters the proof of
\cref{sel:positive-marginality}. It holds for every stationary cluster with
positive plateau, and its proof does not use marginality.

\begin{proof}[Proof of \cref{prop:threshold-inputs}\ref{item:threshold-forgetting}]
The hypotheses of part~\ref{item:threshold-forgetting} are those of
\cref{sel:general-fixed-wait}, which gives $C(t,s)\to0$ as $t\to\infty$ for
each fixed $s$. It remains to make the convergence uniform in
$s\in[0,S]$. Let $k\ge1$ be an integer and let $s_i=i/k$ for
$0\le i\le\lceil kS\rceil$. Every $s\in[0,S]$ is within $1/k$ of some
$s_i$, and by the correlation modulus \eqref{eq:global-modulus} in the
second variable,
\[
 |C(t,s)-C(t,s_i)|\le\frac{D_{p,\beta}}k+\sqrt{\frac2k}
\]
for all $t$. Since $C\ge0$, we obtain
\[
 \sup_{0\le s\le S}|C(t,s)|\le\max_{i}C(t,s_i)+\frac{D_{p,\beta}}k+\sqrt{\frac2k}.
\]
The maximum is over finitely many points, and it tends to zero as
$t\to\infty$ by \cref{sel:general-fixed-wait}. Letting $t\to\infty$ and then
$k\to\infty$ proves \eqref{eq:threshold-input-forgetting}.
\end{proof}


\part{Response and the fluctuation--dissipation relation}\label{part:response}
\chapter{Localization of the response and of correlation derivatives}\label{resp:chapter}

In this section, we study complete rows of the response and of both partial derivatives of the correlation and show that they converge to their stationary limits. Our main result, \cref{resp:strong-row}, states that for every $p\ge3$ and $\beta>0$, the lag row $v\mapsto R(t,t-v)$, extended by zero, converges to $r$ in $L^2(0,\infty)$. From it we deduce uniform convergence of response rows (\cref{resp:uniform-row}). We also show that the earlier derivative $\partial_sC(t,s)$ converges in the same $L^2$ sense, so that the squared backward FDT defect vanishes on whole rows (\cref{resp:backward-fdt}). Further, we bound the later derivative $\partial_tC(t,s)$ far from the diagonal for every $p\ge3$ (\cref{resp:forward-derivative}), prove a smoothing estimate in the earlier time (\cref{resp:remote-earlier}), and prove that the thermal part of the Green covariance identity converges uniformly on late triangles (\cref{resp:thermal}).

We use structural identities from \cref{sec:ck-structure} (\cref{prop:green-covariance,prop:uniform-bounds,prop:backward-response,prop:response-compactness}), the Fisher bound of \cref{thm:fisher-clock}, convergence and forgetting from \cref{sec:stationary-selection} (\cref{sel:all-temperature,sel:fixed-wait-forgetting}), and the causal tree realization (\cref{prop:geometry-tree-realization,lem:geometry-projected-drift,sel:prefix-reflection,sel:cluster-complete-monotonicity}). When the plateau is positive, the proof of \cref{resp:strong-row} also uses the cubic tree ceiling of \cref{sec:tree-ceilings} if $p=3$, and the finite forks of \cref{app:general-selection} for general $p$. We give the second argument at the end of this section, in \cref{resp:appendix-p}. The results of this section are used in \cref{resp:integrated-chapter}, which relies on \cref{resp:backward-fdt}, and in \cref{sec:three-time,resp:fisher-chapter}, which use \cref{resp:strong-row,resp:uniform-row,resp:forward-derivative}, together with the Fisher bound \eqref{resp:fisher-domination}. Combined with \cref{thm:fisher-clock}, they prove \eqref{eq:main-strong-rows}, \eqref{eq:main-uniform-response}, and \eqref{eq:main-defect-l2} in \cref{thm:main-response}.

All convergence statements in this section are in the $L^2$ norm or in the uniform norm on a row. They do not control the integral of a row over its remote part. In \cref{resp:integrated-chapter}, we prove that the backward FDT defect, although small in $L^2$, keeps a positive integrated mass there.

\paragraph{Standing assumptions and notation.}
Throughout this section, an integer $p\ge3$ and $\beta>0$ are fixed, except where a statement restricts $p$. Let $(c,r)$ be the stationary profile of \cref{sel:all-temperature}, and let $q=\lim_{v\to\infty}c(v)\in[0,1)$ be its plateau. Recall from \cref{prop:stationary-clusters,sel:cluster-complete-monotonicity} that
\[
 c(0)=r(0)=1,\qquad r=-c',\qquad \int_0^\infty r(v)\,\dd v=1-q,
\]
and that $c$ and $r$ are completely monotone on $(0,\infty)$. In particular, $r$ is nonnegative and nonincreasing, and $r(v)\to0$ as $v\to\infty$ because $r$ is integrable.

By \cref{sel:all-temperature}, the translated kernels converge along every sequence of times tending to infinity. Taking the translation time equal to $t-A$, we find that, for every fixed $0\le A<\infty$,
\begin{equation}\label{eq:resp-recent-strip}
 \lim_{t\to\infty}\sup_{0\le v\le A}\bigl\{|C(t,t-v)-c(v)|+|R(t,t-v)-r(v)|\bigr\}=0.
\end{equation}
We call $\{(t,s):t-A\le s\le t\}$ the recent strip of length $A$. Recall also the forgetting statement \eqref{sel:fixed-wait-decay}: for every fixed $S\ge0$, we have $\sup_{0\le s\le S}\{C(t,s)+R(t,s)\}\to0$ as $t\to\infty$.

We use the bounds $0\le R(t,s)\le C(t,s)\le1$ for $t\ge s$ from \cref{prop:response-compactness}, and $1\le\lambda\le\Lambda$ from \eqref{eq:uniform-multiplier}. Since the memory integral in the response equation \eqref{eq:ck-response} is nonnegative, we have $\partial_tR(t,s)\ge-\lambda(t)R(t,s)$ for $t>s$. Integrating this inequality from the diagonal, where $R(s,s)=1$, we find
\begin{equation}\label{eq:resp-response-lower}
 R(t,s)\ge\exp\Bigl(-\int_s^t\lambda(u)\,\dd u\Bigr)\ge e^{-\Lambda(t-s)},\qquad t\ge s\ge0.
\end{equation}

For $t\ge0$, we extend the lag row of the response by zero:
\begin{equation}
 f_t(v)=\mathbf1_{[0,t]}(v)R(t,t-v),\qquad v\ge0.
 \label{resp:lag-row}
\end{equation}
We write $\norm{\cdot}_2$ for the norm of $L^2(0,\infty)$. By \eqref{eq:initial-response-row-bound}, we have $\norm{f_t}_2^2=\int_0^tR(t,u)^2\,\dd u\le1/2$. For $t\ge L\ge0$, we write
\[
 S_2(t,L)=\int_0^{t-L}R(t,u)^2\,\dd u=\int_L^\infty f_t(v)^2\,\dd v
\]
for the squared response in the remote part of the row. Here and below, an integral with upper endpoint $t-L$ is considered only for $t\ge L$.

Let $v_t$ be the tree vectors of \cref{prop:geometry-tree-realization}, so that $\ip{v_t}{v_s}=C(t,s)$ and $\norm{v_t}=1$. Let $P_T$ be the prefix projections, let $\Pi_j$ be the projections onto trees of height $j$, and let $h(v)=\sqrt p\ip{v}{\mathcal I(v^{\otimes(p-1)})}$ be the tree energy. Recall that $v_s$ belongs to the range of $P_T$ when $s\le T$. For $T,a\ge0$, we set $M_T(a)=P_Tv_{T+a}$. As the whole translated family converges, \cref{sel:prefix-reflection} applies along every sequence $T\to\infty$. For every fixed $L\ge0$, it implies that
\begin{equation}
 \norm{M_T(L)}^2\longrightarrow d_L=c(2L),\qquad
 h(M_T(L))\longrightarrow h_L\qquad(T\to\infty),
 \label{resp:prefix-norm}
\end{equation}
where $h_L$ is given by \eqref{sel:prefix-energy} with $e=e_\beta(\infty)$. If $q>0$, then $d_L\ge q>0$, and we define
\[
 \delta_L=\Eth-\frac{h_L}{d_L^{p/2}}.
\]
By \cref{prop:threshold-inputs}\ref{item:threshold-ceiling}, we have $\delta_L\ge0$ for every $L$. Recall also that the plateau is then marginal, by \cref{sel:positive-marginality}. We prove in \eqref{resp:p-center-deficit} below that $\delta_L\to0$ as $L\to\infty$ for every $p\ge3$. For $p\in\{3,4\}$, both properties are part of \cref{sel:threshold-centers}, whose hypotheses hold by \cref{sel:all-temperature,sel:positive-marginality}. In words, the tree energy of the normalized prefix centers $M_T(L)/\norm{M_T(L)}$ approaches the threshold energy $\Eth$ when $T\to\infty$ first and then $L\to\infty$.

\section{The full response row}\label{sec:resp-full-row}

\begin{theorem}[Strong localization of the response]
\label{resp:strong-row}
Fix $p\ge3$ and $\beta>0$. Then
\begin{equation}
 \lim_{t\to\infty}\norm{f_t-r}_{L^2(0,\infty)}=0.
 \label{resp:strong-row-limit}
\end{equation}
Further, for each fixed $L\ge0$,
\begin{equation}
 \lim_{t\to\infty}\int_0^{t-L}R(t,u)^2\,\dd u=\int_L^\infty r(v)^2\,\dd v,
 \qquad \lim_{t\to\infty}\frac{\chi(t,0)}{\sqrt t}=0.
 \label{resp:square-tail}
\end{equation}
\end{theorem}

The main step is to show that the response keeps no squared mass in the remote part of the row:
\begin{equation}\label{eq:resp-remote-vanishing}
 \lim_{L\to\infty}\limsup_{t\to\infty}S_2(t,L)=0.
\end{equation}
Once \eqref{eq:resp-remote-vanishing} is known, the theorem follows from stationary convergence on the recent strip. We prove \eqref{eq:resp-remote-vanishing} in three cases, which use different inputs.
\begin{itemize}
\item \emph{$q=0$ and every $p\ge3$.} The white-leaf component of the prefix center $M_T(L)$ carries the squared response $2S_2(T+L,L)$, and its norm is at most $\norm{M_T(L)}$. Here we use only the tree realization (\cref{prop:geometry-tree-realization}) and reflection (\cref{sel:prefix-reflection}).
\item \emph{$q>0$ and $p=3$.} We normalize the prefix center and control its white-leaf component by the height inequality of the cubic tree ceiling (\cref{thm:geometry-tree-ceiling}). Besides the inputs of the first case, we need that $\delta_L\to0$, which is part of \cref{sel:threshold-centers} in \cref{sec:tree-ceilings}.
\item \emph{$q>0$ and every $p\ge3$.} We test the mean of finitely many independent continuations against a Gaussian direction built from the Brownian path, and use the overlap penalty of \cref{sel:absolute-field-cost} for an independent direction. This case relies on the finite forks of \cref{app:general-selection} and is carried out in \cref{resp:appendix-p} (\cref{prop:remote-response-general}).
\end{itemize}
The third case contains the second. We include the second because it is short and uses only the deterministic tree geometry.

\begin{proof}[Proof of \cref{resp:strong-row}]
We prove \eqref{eq:resp-remote-vanishing} in Steps~1--3 and deduce the theorem from it in Step~4.

\paragraph{Step 1: the white-leaf bound.}
Fix $T,L\ge0$. By \eqref{eq:geometry-tree-realization}, the component of $v_{T+L}$ in the one-leaf summand $\mathbb Re_0\oplus L^2(\R_+)$ is $R(T+L,0)e_0+\sqrt2\,\mathbf1_{[0,T+L]}R(T+L,\cdot)$. The remaining term of \eqref{eq:geometry-tree-realization} lies in the range of the grafting map, which is orthogonal to that summand. Recall that $P_T$ commutes with $\Pi_0$, and that on the one-leaf summand it restricts the noise-time function to $[0,T]$. We conclude that
\begin{equation}
 \norm{\Pi_0M_T(L)}^2=R(T+L,0)^2+2\int_0^TR(T+L,u)^2\,\dd u.
 \label{resp:white-prefix-norm}
\end{equation}
As $\Pi_0$ is an orthogonal projection, \eqref{resp:white-prefix-norm} implies that $2S_2(T+L,L)\le\norm{\Pi_0M_T(L)}^2\le\norm{M_T(L)}^2$. Letting $T\to\infty$ and using \eqref{resp:prefix-norm}, we find that, with $t=T+L$,
\begin{equation}
 \limsup_{t\to\infty}S_2(t,L)\le\frac{d_L}2=\frac{c(2L)}2.
 \label{resp:basic-tail-bound}
\end{equation}
This bound holds for every $p\ge3$ and every plateau. If $q=0$, then $c(2L)\to0$ as $L\to\infty$, and \eqref{resp:basic-tail-bound} proves \eqref{eq:resp-remote-vanishing}.

\paragraph{Step 2: a positive plateau in the cubic model.}
Suppose that $q>0$ and $p=3$, and fix $L\ge0$. For all large $T$, we have $\norm{M_T(L)}^2\ge q/2$ by \eqref{resp:prefix-norm}; for such $T$, let $z_T=M_T(L)/\norm{M_T(L)}$. For $p=3$, the tree energy is $h=\sqrt3f$, where $f(v)=\ip{v}{\mathcal I(v\otimes v)}$ is the functional of \cref{thm:geometry-tree-ceiling}, and $\Eth=2\sqrt{2/3}=\sqrt3F_*$ with $F_*=2\sqrt2/3$. Keeping only the term $j=0$ in the first inequality of \eqref{eq:geometry-height-defect}, we have
\[
 \norm{\Pi_0z_T}^4\le4\sqrt2\,\bigl[F_*-f(z_T)\bigr]=4\sqrt{2/3}\,\bigl[\Eth-h(z_T)\bigr].
\]
On the other hand, \eqref{resp:white-prefix-norm} implies that $\norm{\Pi_0z_T}^2\ge2S_2(T+L,L)/\norm{M_T(L)}^2$. Combining the two bounds, we find
\[
 S_2(T+L,L)\le\frac{\norm{M_T(L)}^2}2\Bigl(4\sqrt{2/3}\,\bigl[\Eth-h(z_T)\bigr]\Bigr)^{1/2}.
\]
The tree energy is homogeneous of degree three, which gives $h(z_T)=h(M_T(L))/\norm{M_T(L)}^3$, and \eqref{resp:prefix-norm} implies that $h(z_T)\to h_L/d_L^{3/2}=\Eth-\delta_L$. Letting $T\to\infty$ with $L$ fixed, we obtain
\begin{equation}
 \limsup_{t\to\infty}S_2(t,L)
 \le\frac{d_L}2\sqrt{4\sqrt{2/3}\,\delta_L}.
 \label{resp:threshold-tail-bound}
\end{equation}
\Cref{sel:threshold-centers} states that $\delta_L\to0$ as $L\to\infty$. As $d_L\le1$, the right side of \eqref{resp:threshold-tail-bound} tends to zero, and \eqref{eq:resp-remote-vanishing} follows.

\paragraph{Step 3: a positive plateau for general $p$.}
Suppose that $q>0$ and $p\ge3$. \Cref{prop:remote-response-general} states that, for every $L>0$,
\[
 \limsup_{t\to\infty}S_2(t,L)\le\frac{d_L}2\,\bd^{-2/p}\,\delta_L^{2/p},
\]
and the right side tends to zero as $L\to\infty$ by \eqref{resp:p-center-deficit}. The proof of \cref{prop:remote-response-general}, in \cref{resp:appendix-p}, is independent of \cref{resp:strong-row} and its consequences. This completes the proof of \eqref{eq:resp-remote-vanishing} in all cases.

\paragraph{Step 4: from the remote bound to strong convergence.}
We first check that $r\in L^2(0,\infty)$. For fixed $V$, we have $f_t\to r$ uniformly on $[0,V]$ by \eqref{eq:resp-recent-strip}, so $\int_0^Vr^2=\lim_t\int_0^Vf_t^2\le1/2$. Letting $V\to\infty$, we find $\norm{r}_2^2\le1/2$.

Next, fix $L>0$ and $t\ge L$. Splitting $\norm{f_t-r}_2^2$ at $L$, and using $(x-y)^2\le2x^2+2y^2$ on $[L,\infty)$ together with $\int_L^\infty f_t^2=S_2(t,L)$, we have
\[
 \norm{f_t-r}_2^2
 \le L\sup_{0\le v\le L}|R(t,t-v)-r(v)|^2+2S_2(t,L)+2\int_L^\infty r(v)^2\,\dd v.
\]
The first term tends to zero by \eqref{eq:resp-recent-strip}, and taking the upper limit in $t$, we find
\[
 \limsup_{t\to\infty}\norm{f_t-r}_2^2\le2\limsup_{t\to\infty}S_2(t,L)+2\int_L^\infty r(v)^2\,\dd v.
\]
Letting $L\to\infty$, and using \eqref{eq:resp-remote-vanishing} and $r\in L^2$, we obtain \eqref{resp:strong-row-limit}.

For the first limit in \eqref{resp:square-tail}, fix $L\ge0$. Using the triangle inequality, and the fact that restriction to $[L,\infty)$ is a contraction of $L^2(0,\infty)$, we have
\[
 \Bigl|S_2(t,L)^{1/2}-\Bigl(\int_L^\infty r(v)^2\,\dd v\Bigr)^{1/2}\Bigr|
 \le\norm{\mathbf1_{[L,\infty)}(f_t-r)}_2\le\norm{f_t-r}_2,
\]
which tends to zero by \eqref{resp:strong-row-limit}. Finally, note that $\chi(t,0)=\int_0^tR(t,u)\,\dd u=\int_0^tf_t(v)\,\dd v$. Using the Cauchy--Schwarz inequality on $[0,t]$, we obtain
\[
 0\le\chi(t,0)\le\int_0^tr(v)\,\dd v+\int_0^t|f_t(v)-r(v)|\,\dd v\le1-q+\sqrt t\,\norm{f_t-r}_2.
\]
After division by $\sqrt t$, the second limit in \eqref{resp:square-tail} follows from \eqref{resp:strong-row-limit}.
\end{proof}

We next upgrade \cref{resp:strong-row} to uniform convergence. The backward response equation prevents $s\mapsto R(t,s)$ from increasing quickly. A large remote value $R(t,s)$ would then force squared mass on the interval $[s-1,s]$, which \eqref{resp:square-tail} excludes.

\begin{corollary}[Uniform absolute localization]
\label{resp:uniform-row}
Fix $p\ge3$ and $\beta>0$. Then
\begin{equation}
 \lim_{t\to\infty}\sup_{0\le s\le t}|R(t,s)-r(t-s)|=0.
 \label{resp:uniform-row-limit}
\end{equation}
In particular, if $s_n\to\infty$ and $t_n-s_n\to\infty$, then $R(t_n,s_n)\to0$.
\end{corollary}

\begin{proof}
Recall that the kernel $\Sigma$ in \eqref{eq:self-energy-kernel} is nonnegative. Since $R\ge0$, the backward response equation \eqref{eq:backward-response} implies that $\partial_sR(t,s)\le\lambda(s)R(t,s)\le\Lambda R(t,s)$ for $t>s$. Then $w\mapsto e^{-\Lambda w}R(t,w)$ is nonincreasing on $[0,t]$, that is, $R(t,u)\ge e^{-\Lambda(s-u)}R(t,s)$ for $u\le s\le t$. Squaring this inequality and integrating it over $u\in[s-1,s]$, we have, for $1\le s\le t$,
\[
 R(t,s)^2\,\frac{1-e^{-2\Lambda}}{2\Lambda}\le\int_{s-1}^sR(t,u)^2\,\dd u.
\]
If moreover $s\le t-L$, the integral on the right is at most $S_2(t,L)$. Using \eqref{resp:square-tail}, we find that, for every fixed $L\ge0$,
\[
 \limsup_{t\to\infty}\sup_{1\le s\le t-L}R(t,s)
 \le\Bigl(\frac{2\Lambda}{1-e^{-2\Lambda}}\int_L^\infty r(v)^2\,\dd v\Bigr)^{1/2}.
\]
Recall from \eqref{sel:fixed-wait-decay} that $\sup_{0\le s\le1}R(t,s)\to0$. On the recent strip $t-L\le s\le t$, the difference $|R(t,s)-r(t-s)|$ tends to zero uniformly by \eqref{eq:resp-recent-strip}. On the remaining part $0\le s\le t-L$, we have $|R(t,s)-r(t-s)|\le R(t,s)+r(L)$, since $r$ is nonincreasing. Combining these three estimates, we have
\[
 \limsup_{t\to\infty}\sup_{0\le s\le t}|R(t,s)-r(t-s)|
 \le\Bigl(\frac{2\Lambda}{1-e^{-2\Lambda}}\int_L^\infty r(v)^2\,\dd v\Bigr)^{1/2}+r(L).
\]
The right side tends to zero as $L\to\infty$, which proves \eqref{resp:uniform-row-limit}. For the last assertion, note that $R(t_n,s_n)\le|R(t_n,s_n)-r(t_n-s_n)|+r(t_n-s_n)$, and both terms tend to zero.
\end{proof}

\section{The later correlation derivative}\label{sec:resp-later-derivative}

We now bound $\partial_tC(t,s)$ when $t-s$ is large. By \cref{lem:geometry-projected-drift}, the projected curve $a\mapsto M_T(a)$ is continuously differentiable for $a>0$, with
\[
 \frac{\dd}{\dd a}M_T(a)=P_T\mathcal B(v_{T+a}),\qquad
 \mathcal B(v)=-v+\beta\nabla h(v)-p\beta h(v)v,
\]
and $\mathcal B$ is Lipschitz on the closed unit ball of $\mathcal H$. Reflection identifies limits of inner products of prefix centers, not of their derivatives. For this reason, we pass through difference quotients. This requires a modulus of continuity for the projected derivative that is uniform in $T$, which we now record.

\begin{lemma}[Hölder continuity of the projected drift]\label{lem:resp-drift-holder}
Fix $p\ge3$ and $\beta>0$. Let $L_{\mathcal B}$ be a Lipschitz constant of $\mathcal B$ on the closed unit ball of $\mathcal H$, and set $\kappa_{\mathcal B}=L_{\mathcal B}\sqrt{2\Lambda}$. Then, for all $t\ge s\ge0$,
\[
 0\le1-C(t,s)\le\Lambda(t-s),\qquad \norm{v_t-v_s}\le\sqrt{2\Lambda(t-s)},
\]
and, for all $t,t'\ge0$ and $T\ge0$,
\[
 \norm{\mathcal B(v_t)-\mathcal B(v_{t'})}\le\kappa_{\mathcal B}|t-t'|^{1/2},
 \qquad
 \norm{P_T\mathcal B(v_t)-P_T\mathcal B(v_{t'})}\le\kappa_{\mathcal B}|t-t'|^{1/2}.
\]
\end{lemma}

\begin{proof}
The first bound is \eqref{eq:diagonal-increment} in \cref{prop:global-regularity}. The second bound follows from the identity $\norm{v_t-v_s}^2=2-2C(t,s)$. The vectors $v_t$ lie in the unit ball, and the third bound follows from the second and the Lipschitz property of $\mathcal B$. The last bound holds because $P_T$ is a contraction.
\end{proof}

\begin{theorem}[Remote derivative bound and forward FDT]
\label{resp:forward-derivative}
Fix an integer $p\ge3$ and $\beta>0$. For every fixed $L>0$,
\begin{equation}
 \limsup_{t\to\infty}\sup_{0\le s\le t-L}
 |\partial_tC(t,s)|\le\sqrt{c''(2L)}
 \le\frac{\sqrt{1-q}}{eL}.
 \label{resp:remote-forward}
\end{equation}
Further,
\begin{equation}
 \lim_{t\to\infty}\sup_{0\le s\le t}|\partial_tC(t,s)-c'(t-s)|=0,
 \qquad
 \lim_{t\to\infty}\sup_{0\le s\le t}|\partial_tC(t,s)+R(t,s)|=0.
 \label{resp:uniform-forward-fdt}
\end{equation}
Here $\partial_tC(t,t)$ is the right derivative of $u\mapsto C(u,t)$ at $u=t$; it equals $-1$ for every $t\ge0$. The value $c'(0)=-1$ is also the right derivative.
\end{theorem}

\begin{proof}
All inputs to this proof hold for every integer $p\ge3$: the projected drift of \cref{lem:geometry-projected-drift}, its uniform modulus in \cref{lem:resp-drift-holder}, prefix reflection in \cref{sel:prefix-reflection}, complete monotonicity in \cref{sel:cluster-complete-monotonicity}, and uniform response localization in \cref{resp:uniform-row}.

The proof has three steps. First, we bound the norm of the projected derivative $M_T'(L)$ through difference quotients and reflection. Since $\partial_tC(T+L,s)=\ip{M_T'(L)}{v_s}$ for every $s\le T$, we deduce \eqref{resp:remote-forward} from this bound. Second, we show that $\partial_tC(t,t-v)\to c'(v)$ uniformly on each recent strip, by differentiating stationary convergence through forward difference quotients. Third, we combine the two estimates with \cref{resp:uniform-row}.

\paragraph{Step 1: the remote bound.}
Fix $L>0$ and $\eta\in(0,1]$. For every $T\ge0$, we have $M_T(L+\eta)-M_T(L)=\int_0^\eta M_T'(L+w)\,\dd w$. \Cref{lem:resp-drift-holder} then implies that
\[
 \Bigl\|M_T'(L)-\frac{M_T(L+\eta)-M_T(L)}\eta\Bigr\|
 \le\frac1\eta\int_0^\eta\norm{M_T'(L+w)-M_T'(L)}\,\dd w
 \le\frac23\kappa_{\mathcal B}\sqrt\eta.
\]
Applying \eqref{sel:prefix-gram} to the three inner products in the expansion of the squared norm, we have
\[
 \lim_{T\to\infty}\Bigl\|\frac{M_T(L+\eta)-M_T(L)}\eta\Bigr\|^2
 =\frac{c(2L+2\eta)-2c(2L+\eta)+c(2L)}{\eta^2}.
\]
Combining the last two displays, we find
\[
 \limsup_{T\to\infty}\norm{M_T'(L)}
 \le\frac23\kappa_{\mathcal B}\sqrt\eta
 +\Bigl(\frac{c(2L+2\eta)-2c(2L+\eta)+c(2L)}{\eta^2}\Bigr)^{1/2}.
\]
The left side does not depend on $\eta$. The function $c$ is smooth on $(0,\infty)$, and the second difference quotient tends to $c''(2L)$ as $\eta\downarrow0$. Taking $T\to\infty$ first and then $\eta\downarrow0$, we have
\[
 \limsup_{T\to\infty}\norm{M_T'(L)}\le\sqrt{c''(2L)}.
\]
If $s\le T$, then $v_s$ lies in the range of $P_T$, so $C(T+a,s)=\ip{M_T(a)}{v_s}$ for all $a\ge0$, and $\partial_tC(T+L,s)=\ip{M_T'(L)}{v_s}$. As $\norm{v_s}=1$, we have $|\partial_tC(T+L,s)|\le\norm{M_T'(L)}$ for every $s\le T$. Setting $T=t-L$ proves the first inequality in \eqref{resp:remote-forward}.

For the second inequality, we write the representation \eqref{eq:stationary-spectral-measure} as $c(v)=q+\int_{(0,\infty)}e^{-xv}\,\nu(\dd x)$, where $\nu((0,\infty))=1-q$. Differentiating under the integral for $v>0$, we have
\[
 c''(2L)=\int_{(0,\infty)}x^2e^{-2Lx}\,\nu(\dd x)\le\frac{1-q}{e^2L^2},
\]
where we used that $\sup_{x>0}x^2e^{-2Lx}=(eL)^{-2}$, attained at $x=1/L$.

\paragraph{Step 2: the recent strip.}
For $t>s$, we apply the projected drift with $T=s$ to $C(t,s)=\ip{P_sv_t}{v_s}$. This gives
\[
 \partial_tC(t,s)=\ip{P_s\mathcal B(v_t)}{v_s}=\ip{\mathcal B(v_t)}{v_s}.
\]
The right side is continuous in $t\ge s$. At $t=s$ it equals $-1$, because $\ip{\nabla h(v)}{v}=ph(v)$ by homogeneity and $\norm{v_s}=1$; this is the derivative from the causal side at the diagonal. \Cref{lem:resp-drift-holder} also implies that $|\partial_tC(t,s)-\partial_tC(t',s)|\le\kappa_{\mathcal B}|t-t'|^{1/2}$ for all $t,t'\ge s$.

Fix $L>0$ and $\eta\in(0,1]$. For $0\le v\le L$ and $t\ge L$, we hold the earlier time $t-v$ fixed in a forward difference quotient. The Hölder bound just stated gives
\[
 \Bigl|\partial_tC(t,t-v)-\frac{C(t+\eta,t-v)-C(t,t-v)}\eta\Bigr|
 \le\frac1\eta\int_0^\eta\kappa_{\mathcal B}\sqrt w\,\dd w=\frac23\kappa_{\mathcal B}\sqrt\eta.
\]
By \eqref{eq:resp-recent-strip} with $A=L+\eta$, the difference quotient converges to $[c(v+\eta)-c(v)]/\eta$ uniformly in $v\in[0,L]$. We obtain
\[
 \limsup_{t\to\infty}\sup_{0\le v\le L}|\partial_tC(t,t-v)-c'(v)|
 \le\frac23\kappa_{\mathcal B}\sqrt\eta
 +\sup_{0\le v\le L}\Bigl|\frac{c(v+\eta)-c(v)}\eta-c'(v)\Bigr|.
\]
The function $c'=-r$ is uniformly continuous on $[0,L+1]$, including the right derivative $c'(0)=-1$, and the last supremum tends to zero as $\eta\downarrow0$. We conclude that $\sup_{t-L\le s\le t}|\partial_tC(t,s)-c'(t-s)|\to0$ for every fixed $L$.

\paragraph{Step 3: the whole row.}
On the remote part $0\le s\le t-L$, we have $|\partial_tC(t,s)-c'(t-s)|\le|\partial_tC(t,s)|+r(L)$. Using Step 2 and \eqref{resp:remote-forward}, we find
\[
 \limsup_{t\to\infty}\sup_{0\le s\le t}|\partial_tC(t,s)-c'(t-s)|\le\sqrt{c''(2L)}+r(L).
\]
Letting $L\to\infty$ proves the first limit in \eqref{resp:uniform-forward-fdt}. As $c'=-r$, the second limit follows from the first and \cref{resp:uniform-row}.
\end{proof}

\section{The earlier derivative and the squared FDT defect}\label{sec:resp-earlier-derivative}

For $t>s>0$, let $E(t,s)=\partial_sC(t,s)-R(t,s)$ be the backward FDT defect. By \cref{thm:fisher-clock}, the derivative $\partial_sC$ is continuous on the open causal triangle, and there exists a nonnegative integrable function $a_F$, obtained on one dimension subsequence, such that
\begin{equation}
 |E(t,s)|^2\le a_F(s)\quad(t>s>0),\qquad
 \int_0^\infty a_F(s)\,\dd s<\infty,\qquad
 \lim_{S\to\infty}\sup_{s\ge S}a_F(s)=0.
 \label{resp:fisher-domination}
\end{equation}
We do not need $a_F$ to be unique. The bound \eqref{resp:fisher-domination} is the spherical form used here of the entropy-production bounds on FDT violation of Cugliandolo, Dean, and Kurchan \cite[arXiv version~2, equations~(12)--(13), p.~3, and~(18), p.~4]{cdk1997}. The common majorant in \eqref{resp:fisher-domination} is proved in \cref{thm:fisher-clock}, and the whole-row conclusions below are proved here. It shows that the defect is uniformly small when the earlier time is late. The row $[0,t]$ grows with $t$, however, and this alone does not control the defect on the whole row. The next theorem shows that the integral of the squared defect over the row nevertheless tends to zero.

\begin{theorem}[Strong backward FDT]
\label{resp:backward-fdt}
Fix $p\ge3$ and $\beta>0$. Then
\begin{equation}
 \lim_{t\to\infty}\sup_{0\le s\le S}|\partial_sC(t,s)|=0
 \quad\text{for each fixed }0\le S<\infty,
 \qquad
 \lim_{t\to\infty}\int_0^t|E(t,s)|^2\,\dd s=0.
 \label{resp:backward-fdt-limit}
\end{equation}
Further, if $g_t(v)=\mathbf1_{[0,t]}(v)\partial_sC(t,t-v)$, then
\begin{equation}
 \lim_{t\to\infty}\norm{g_t-r}_{L^2(0,\infty)}=0.
 \label{resp:earlier-row-limit}
\end{equation}
For $t>0$, the earlier derivative at $s=0$ is the right derivative. At $s=t>0$, it is the left derivative of $u\mapsto C(t,u)$ and equals $1$; we set $E(t,t)=0$.
\end{theorem}

\begin{proof}
We proceed in three steps. First, we represent the colored memory term of the Green identity through feature vectors. A positivity argument, together with forgetting, shows that this term is negligible on every bounded interval of earlier times. Second, we differentiate the Green identity in its earlier time. This expresses $\partial_sC(t,s)$ through $C$, $R$, and the colored term, and it proves the first claim in \eqref{resp:backward-fdt-limit}. Third, the Fisher majorant allows dominated convergence over the whole row, and \cref{resp:strong-row} implies \eqref{resp:earlier-row-limit}.

\paragraph{Step 1: the colored term forgets bounded earlier intervals.}
For $s\ge0$, let $\phi_s=v_s^{\otimes(p-1)}\in\mathcal H^{\otimes(p-1)}$. These are unit vectors with $\ip{\phi_u}{\phi_s}=C(u,s)^{p-1}$ and $\norm{\phi_u-\phi_s}^2=2[1-C(u,s)^{p-1}]$. In particular, $s\mapsto\phi_s$ is norm continuous. We define
\[
 F_t=\int_0^tR(t,u)\phi_u\,\dd u,\qquad
 \mathcal K(t,s)=\ip{F_t}{\phi_s}=\int_0^tR(t,u)C(u,s)^{p-1}\,\dd u,
\]
where the first integral is a Bochner integral of a continuous function on $[0,t]$. Note that $\mathcal K(t,s)\ge0$. When $s=t$, the colored memory term in the Green identity \eqref{eq:green-covariance} equals $\beta^2p\norm{F_t}^2$. Since all three terms of \eqref{eq:green-covariance} are nonnegative, we obtain $\beta^2p\norm{F_t}^2\le C(t,t)=1$. By the Cauchy--Schwarz inequality, we have, for all $t,s,v\ge0$,
\begin{equation}
 |\mathcal K(t,s)-\mathcal K(t,v)|\le(\beta^2p)^{-1/2}\sqrt{2[1-C(s,v)^{p-1}]}.
 \label{resp:colored-modulus}
\end{equation}
This modulus does not depend on $t$. For $s\le t$, Fubini's theorem expresses the colored memory term in \eqref{eq:green-covariance} as $\beta^2p\int_0^sR(s,v)\mathcal K(t,v)\,\dd v$. The other two terms are nonnegative, and therefore
\begin{equation}
 C(t,s)\ge\beta^2p\int_0^sR(s,v)\mathcal K(t,v)\,\dd v.
 \label{resp:colored-test}
\end{equation}

We claim that $\mathcal K(t,s)\to0$ as $t\to\infty$, for every fixed $s>0$. Suppose instead that $\mathcal K(t_n,s)\ge\varepsilon$ along a sequence $t_n\to\infty$, for some $\varepsilon>0$. Since $C$ is continuous and $C(s,s)=1$, there exists $h\in(0,s]$ such that the right side of \eqref{resp:colored-modulus} is at most $\varepsilon/2$ for $v\in[s-h,s]$. Then $\mathcal K(t_n,v)\ge\varepsilon/2$ on $[s-h,s]$. Inserting this and the lower bound \eqref{eq:resp-response-lower} into \eqref{resp:colored-test}, we obtain
\[
 C(t_n,s)\ge\frac{\beta^2p\,\varepsilon}2\int_{s-h}^se^{-\Lambda(s-v)}\,\dd v
 =\frac{\beta^2p\,\varepsilon}{2\Lambda}\bigl(1-e^{-\Lambda h}\bigr)>0,
\]
which contradicts $C(t_n,s)\to0$ from \eqref{sel:fixed-wait-decay}. For $s=0$ and any $s'>0$, \eqref{resp:colored-modulus} implies that $\mathcal K(t,0)\le\mathcal K(t,s')+(\beta^2p)^{-1/2}\sqrt{2[1-C(s',0)^{p-1}]}$. Letting $t\to\infty$ and then $s'\downarrow0$, we find $\mathcal K(t,0)\to0$.

Finally, fix $S$ and $\varepsilon>0$. The correlation is uniformly continuous on $[0,S]^2$ with $C(s,s)=1$. Hence there exists a finite set $\{s_1,\dots,s_k\}\subset[0,S]$ such that every $s\in[0,S]$ has some $s_i$ with $(\beta^2p)^{-1/2}\sqrt{2[1-C(s,s_i)^{p-1}]}\le\varepsilon$. Then \eqref{resp:colored-modulus} implies that $\sup_{0\le s\le S}\mathcal K(t,s)\le\max_i\mathcal K(t,s_i)+\varepsilon$. Letting $t\to\infty$ and then $\varepsilon\downarrow0$, we conclude that
\begin{equation}\label{eq:resp-colored-forgetting}
 \lim_{t\to\infty}\sup_{0\le s\le S}\mathcal K(t,s)=0.
\end{equation}

\paragraph{Step 2: the backward correlation equation.}
Fix $t>0$. By the Green identity \eqref{eq:green-covariance} and Fubini's theorem, we have, for $0\le s\le t$,
\[
 C(t,s)=R(t,0)R(s,0)+2\int_0^sR(t,u)R(s,u)\,\dd u+\beta^2p\int_0^sR(s,v)\mathcal K(t,v)\,\dd v.
\]
By \cref{prop:backward-response}, the derivative $\partial_1R$ is continuous on each compact causal triangle, and $\mathcal K(t,\cdot)$ is continuous. We may therefore differentiate in $s\in(0,t)$ under the integral signs. Using $R(s,s)=1$, we obtain
\begin{align*}
 \partial_sC(t,s)={}&2R(t,s)+\beta^2p\,\mathcal K(t,s)+R(t,0)\,\partial_1R(s,0)\\
 &+2\int_0^sR(t,u)\,\partial_1R(s,u)\,\dd u+\beta^2p\int_0^s\partial_1R(s,v)\mathcal K(t,v)\,\dd v.
\end{align*}
The first two terms are the boundary terms of the thermal and colored integrals. We now insert the response equation \eqref{eq:ck-response}, which reads $\partial_1R(s,v)=-\lambda(s)R(s,v)+\int_v^s\Sigma(s,w)R(w,v)\,\dd w$ with $\Sigma(s,w)=\beta^2p(p-1)C(s,w)^{p-2}R(s,w)$. The terms with $-\lambda(s)$ reproduce $-\lambda(s)$ times the Green representation of $C(t,s)$. In the memory terms, we exchange the order of integration on the triangle $\{0\le u\le w\le s\}$, and they become
\[
 \int_0^s\Sigma(s,w)\Bigl[R(t,0)R(w,0)+2\int_0^wR(t,u)R(w,u)\,\dd u+\beta^2p\int_0^wR(w,v)\mathcal K(t,v)\,\dd v\Bigr]\dd w.
\]
By the Green identity at $(t,w)$, the bracket equals $C(t,w)$. We have shown that, for $0<s<t$,
\begin{align}
 \partial_sC(t,s)={}&-\lambda(s)C(t,s)+2R(t,s)
       +\beta^2p(p-1)\int_0^s
          C(s,u)^{p-2}R(s,u)C(t,u)\,\dd u\notag\\
       &+\beta^2p\,\mathcal K(t,s).
 \label{resp:backward-correlation}
\end{align}
All integrations and differentiations here take place on finite rectangles, where the integrands are continuous and bounded. The right side of \eqref{resp:backward-correlation} is continuous on $[0,t]$, so it also gives the right derivative at $s=0$ and the left derivative at $s=t$.

As a check, we let $s\uparrow t$. Using the definition of $\mathcal K$ and \eqref{eq:ck-multiplier}, we find
\[
\begin{aligned}
 \beta^2p\,\mathcal K(t,t)+\beta^2p(p-1)\int_0^tC(t,u)^{p-1}R(t,u)\,\dd u
 &=\beta^2p^2\int_0^tC(t,u)^{p-1}R(t,u)\,\dd u\\
 &=\lambda(t)-1.
\end{aligned}
\]
Inserting this into \eqref{resp:backward-correlation}, we obtain $\partial_sC(t,t-)=-\lambda(t)+2+\lambda(t)-1=1$, and $E(t,t-)=0$, as the unit diagonal response requires.

Now fix $S$ and let $t>S$. Bounding $C(s,u)^{p-2}\le1$ and $R(s,u)\le1$ in \eqref{resp:backward-correlation}, we have, for $0\le s\le S$,
\[
 |\partial_sC(t,s)|\le\bigl(\Lambda+\beta^2p(p-1)S\bigr)\sup_{0\le u\le S}C(t,u)
 +2\sup_{0\le u\le S}R(t,u)+\beta^2p\sup_{0\le u\le S}\mathcal K(t,u).
\]
Each supremum tends to zero, by \eqref{sel:fixed-wait-decay} and \eqref{eq:resp-colored-forgetting}. This proves the first claim in \eqref{resp:backward-fdt-limit}.

\paragraph{Step 3: the whole row.}
Let $\widetilde E_t(s)=\mathbf1_{(0,t)}(s)E(t,s)$ for $s>0$. For fixed $s>0$, the first claim in \eqref{resp:backward-fdt-limit} and \eqref{sel:fixed-wait-decay} imply that $\widetilde E_t(s)\to0$ as $t\to\infty$. Further, $|\widetilde E_t(s)|^2\le a_F(s)$ for all $t$ by \eqref{resp:fisher-domination}, and $a_F$ is integrable on $(0,\infty)$. By the dominated convergence theorem, we conclude that
\[
 \lim_{t\to\infty}\int_0^t|E(t,s)|^2\,\dd s=\lim_{t\to\infty}\int_0^\infty|\widetilde E_t(s)|^2\,\dd s=0.
\]
This is the second claim in \eqref{resp:backward-fdt-limit}. After the change of variables $s=t-v$, we have $\norm{g_t-f_t}_2^2=\int_0^t|E(t,s)|^2\,\dd s$. Then \cref{resp:strong-row} implies that $\norm{g_t-r}_2\le\norm{g_t-f_t}_2+\norm{f_t-r}_2\to0$, which is \eqref{resp:earlier-row-limit}.
\end{proof}

The proof does not require a unique limit of the Fisher information, nor a majorant chosen anew for each terminal time: one function $a_F$ dominates every row. The positivity of $C$ and $R$ enters through \eqref{resp:colored-test}.

The next corollary states that remote values of the response and of the earlier derivative are small, uniformly on the late triangle, and that remote correlations change little across a window of fixed length in the earlier time. It is used in \cref{sec:three-time}.

\begin{corollary}[Uniform remote decay and earlier smoothing]
\label{resp:remote-earlier}
Fix $p\ge3$ and $\beta>0$. Then
\begin{equation}
 \lim_{D\to\infty}\limsup_{S\to\infty}
 \sup_{s\ge S,\ t\ge s+D}
 \bigl(R(t,s)+|\partial_sC(t,s)|\bigr)=0.
 \label{resp:remote-earlier-limit}
\end{equation}
More precisely, for each fixed $D>0$,
\begin{equation}
 \limsup_{S\to\infty}
 \sup_{s\ge S+D,\ t\ge s+D}
 \int_{s-D}^s|\partial_uC(t,u)|\,\dd u
 \le\sqrt D\left(\int_D^\infty r(v)^2\,\dd v\right)^{1/2}.
 \label{resp:earlier-smoothing}
\end{equation}
If $r(v)=O(v^{-1-a})$ as $v\to\infty$ for some $a>0$, the right side of \eqref{resp:earlier-smoothing} is $O(D^{-a})$ as $D\to\infty$.
\end{corollary}

When $\beta\ge\bd$, \cref{thm:stationary-critical-amplitude} gives $r(v)\sim aA_\beta v^{-1-a}$, with the exponent $a\in(0,1/2)$ defined there, and the last hypothesis holds.

\begin{proof}
Let
\[
 \epsilon(U)=\sup_{t\ge U}\norm{f_t-r}_2,\qquad
 \epsilon_\infty(U)=\sup_{t\ge U}\sup_{0\le s\le t}|R(t,s)-r(t-s)|.
\]
Both tend to zero as $U\to\infty$, by \cref{resp:strong-row,resp:uniform-row}. For $t>s>0$, we have $|\partial_sC(t,s)|\le R(t,s)+\sqrt{a_F(s)}$ by \eqref{resp:fisher-domination}.

For the pointwise statement, let $s\ge S$ and $t\ge s+D$. As $r$ is nonincreasing, $R(t,s)\le r(D)+\epsilon_\infty(S+D)$, which implies
\[
 R(t,s)+|\partial_sC(t,s)|\le2r(D)+2\epsilon_\infty(S+D)+\sup_{u\ge S}\sqrt{a_F(u)}.
\]
We take $S\to\infty$, using \eqref{resp:fisher-domination}, and then $D\to\infty$. This proves \eqref{resp:remote-earlier-limit}.

For the integrated statement, fix $D>0$, and let $s\ge S+D$ and $t\ge s+D$. Then $[s-D,s]\subset[S,t-D]$ and $t\ge S+2D$. Recall that $\mu([S,\infty))=\int_S^\infty a_F(u)\,\dd u$, where $\mu$ is the measure of \eqref{eq:entropy-measure}. By the Cauchy--Schwarz inequality,
\[
\begin{aligned}
 \int_{s-D}^s|\partial_uC(t,u)|\,\dd u
 &\le\sqrt D\Bigl(\int_{s-D}^sR(t,u)^2\,\dd u\Bigr)^{1/2}+\sqrt{D\,\mu([S,\infty))}\\
 &\le\sqrt D\,S_2(t,D)^{1/2}+\sqrt{D\,\mu([S,\infty))}.
\end{aligned}
\]
As in Step 4 of the proof of \cref{resp:strong-row}, we have $S_2(t,D)^{1/2}\le(\int_D^\infty r^2)^{1/2}+\norm{f_t-r}_2$. The supremum in \eqref{resp:earlier-smoothing} is therefore at most
\[
 \sqrt D\Bigl[\Bigl(\int_D^\infty r(v)^2\,\dd v\Bigr)^{1/2}+\epsilon(S+2D)\Bigr]+\sqrt{D\,\mu([S,\infty))}.
\]
Since $a_F$ is integrable, $\mu([S,\infty))\to0$, and letting $S\to\infty$ proves \eqref{resp:earlier-smoothing}. Finally, if $r(v)\le A_0v^{-1-a}$ for $v\ge v_0$, then for $D\ge v_0$ we have
\[
 \sqrt D\Bigl(\int_D^\infty r(v)^2\,\dd v\Bigr)^{1/2}\le\sqrt D\Bigl(\frac{A_0^2D^{-1-2a}}{1+2a}\Bigr)^{1/2}=\frac{A_0}{\sqrt{1+2a}}D^{-a}.
\]
This completes the proof.
\end{proof}

\section{The direct thermal covariance}\label{sec:resp-thermal}

The Green identity \eqref{eq:green-covariance} splits $C(t,s)$ into an initial term, a thermal term, which comes directly from the white noise, and a colored memory term. We define the thermal term and its stationary counterpart by
\[
 B(t,s)=2\int_0^{t\wedge s}R(t,u)R(s,u)\,\dd u,
 \qquad b_{\rm th}(\tau)=2\int_0^\infty r(\tau+v)r(v)\,\dd v.
\]
We say that a sequence of pairs $(t_n,s_n)$ with $t_n\ge s_n$ is a late pair with diverging lag if $s_n\to\infty$ and $t_n-s_n\to\infty$.

\begin{corollary}[Thermal covariance on the full late triangle]
\label{resp:thermal}
Fix $p\ge3$ and $\beta>0$. Then
\begin{equation}
 \lim_{S\to\infty}\sup_{t\ge s\ge S}|B(t,s)-b_{\rm th}(t-s)|=0.
 \label{resp:thermal-uniform}
\end{equation}
The function $b_{\rm th}$ is completely monotone on $(0,\infty)$ and tends to zero at infinity. In particular, $B(t_n,s_n)\to0$ along every late pair with diverging lag, and $B(t,t)\to2\int_0^\infty r(v)^2\,\dd v$. If $C(t_n,s_n)\to c_0$ along a late pair with diverging lag, then
\begin{equation}
 \lim_{n\to\infty}\beta^2p\int_0^{t_n}\int_0^{s_n}
 R(t_n,u)R(s_n,v)C(u,v)^{p-1}\,\dd v\,\dd u=c_0.
 \label{resp:colored-aging}
\end{equation}
If $r(\tau)\sim aA\tau^{-1-a}$ as $\tau\to\infty$ for some constants $a,A>0$, then
\begin{equation}
 b_{\rm th}(\tau)\sim2(1-q)\,r(\tau)
 \sim2(1-q)aA\tau^{-1-a}.
 \label{resp:thermal-tail}
\end{equation}
\end{corollary}

By \cref{thm:stationary-critical-amplitude}, the hypothesis of the last assertion holds when $\beta\ge\bd$, with $A=A_\beta$. The identity \eqref{resp:colored-aging} says that any correlation retained along a late pair with diverging lag comes from the colored memory term, with its whole integration rectangle.

\begin{proof}
Let $\tau=t-s\ge0$. After the substitution $u=s-v$, the zero extension \eqref{resp:lag-row} gives
\[
 B(t,s)=2\int_0^\infty f_t(\tau+v)f_s(v)\,\dd v,
\]
because $f_s(v)=0$ for $v>s$. Let $(\mathcal T_\tau g)(v)=g(\tau+v)$; this translation is a contraction of $L^2(0,\infty)$. Then $B(t,s)=2\ip{\mathcal T_\tau f_t}{f_s}$ and $b_{\rm th}(\tau)=2\ip{\mathcal T_\tau r}{r}$. Using $\norm{f_s}_2,\norm{r}_2\le1/\sqrt2$, we have
\[
 |B(t,s)-b_{\rm th}(\tau)|
 \le2\bigl|\ip{\mathcal T_\tau(f_t-r)}{f_s}\bigr|+2\bigl|\ip{\mathcal T_\tau r}{f_s-r}\bigr|
 \le\sqrt2\,\bigl(\norm{f_t-r}_2+\norm{f_s-r}_2\bigr),
\]
uniformly in $\tau\ge0$. With $\epsilon(S)=\sup_{t\ge S}\norm{f_t-r}_2$, the left side of \eqref{resp:thermal-uniform} is at most $2\sqrt2\,\epsilon(S)$, which tends to zero by \cref{resp:strong-row}.

Next, we write $c(v)=q+\int_{(0,\infty)}e^{-xv}\nu(\dd x)$ as in \eqref{eq:stationary-spectral-measure}. Let $\sigma(\dd x)=x\,\nu(\dd x)$ on $(0,\infty)$. Then $r(v)=\int_{(0,\infty)}e^{-xv}\sigma(\dd x)$ for $v>0$, and $\sigma((0,\infty))=\lim_{v\downarrow0}r(v)=1$ by monotone convergence. By Tonelli's theorem,
\[
 b_{\rm th}(\tau)
 =2\int_{(0,\infty)^2}\frac{e^{-x\tau}}{x+y}\,\sigma(\dd x)\,\sigma(\dd y)
 =\int_{(0,\infty)}e^{-x\tau}\rho_{\rm th}(\dd x),
\]
where $\rho_{\rm th}(\dd x)=2\sigma(\dd x)\int_{(0,\infty)}(x+y)^{-1}\sigma(\dd y)$. The measure $\rho_{\rm th}$ is positive with total mass $b_{\rm th}(0)=2\norm{r}_2^2\le1$. Hence $b_{\rm th}$ is the Laplace transform of a finite positive measure on $(0,\infty)$, and it is completely monotone on $(0,\infty)$. Also $b_{\rm th}(\tau)\le2\norm{r}_2\,\norm{\mathbf1_{[\tau,\infty)}r}_2\to0$.

Along a late pair with diverging lag, we have $|B(t_n,s_n)|\le|B(t_n,s_n)-b_{\rm th}(t_n-s_n)|+b_{\rm th}(t_n-s_n)\to0$. Taking $\tau=0$ in \eqref{resp:thermal-uniform}, we find $B(t,t)\to b_{\rm th}(0)=2\norm{r}_2^2$. The initial term of the Green identity satisfies $0\le R(t_n,0)R(s_n,0)\le R(t_n,0)\to0$, by \eqref{sel:fixed-wait-decay}. Subtracting the initial and thermal terms from \eqref{eq:green-covariance} at $(t_n,s_n)$ proves \eqref{resp:colored-aging}.

For \eqref{resp:thermal-tail}, note that $r>0$ by \cref{prop:scalar-clusters}, and that
\[
 \frac{b_{\rm th}(\tau)}{r(\tau)}=2\int_0^\infty\frac{r(\tau+v)}{r(\tau)}\,r(v)\,\dd v.
\]
As $r$ is nonincreasing, the ratio $r(\tau+v)/r(\tau)$ lies in $[0,1]$. It tends to one for each fixed $v$, by the assumed asymptotics. The measure $r(v)\,\dd v$ has finite mass $1-q$. Using dominated convergence, we find $b_{\rm th}(\tau)/r(\tau)\to2(1-q)$.
\end{proof}

\begin{remark}[Square integrable versus integrated response]\label{rem:resp-l2-versus-l1}
The convergence statements of this section concern $L^2$ norms and uniform norms of rows. A slowly varying contribution spread over a growing part of the row can have vanishing squared norm and a positive integral. This happens for the response in the cubic model. For $p=3$ and $\beta\ge\sqrt{4/3}$, \cref{thm:geometry-susceptibility-lower} gives $\liminf_t\chi(t,0)\ge2(1-q)$, while $\int_0^Lf_t(v)\,\dd v\to\int_0^Lr(v)\,\dd v\le1-q$ for each fixed $L$. We conclude that, for every fixed $L$,
\[
 \liminf_{t\to\infty}\int_0^{t-L}R(t,u)\,\dd u\ge1-q>0,
\]
although $\int_0^{t-L}R(t,u)^2\,\dd u\to\int_L^\infty r^2$, which tends to zero as $L\to\infty$. In \cref{resp:integrated-chapter}, we prove a statement of this kind for the backward FDT defect.
\end{remark}

\section{The remote response for general pure models}
\label{resp:appendix-p}

In this section, we prove \cref{prop:remote-response-general}. It is the case $q>0$ of the remote estimate \eqref{eq:resp-remote-vanishing} for general $p\ge3$, which we used in Step~3 of the proof of \cref{resp:strong-row}. For $p=3$, we controlled the white-leaf component of a normalized prefix center directly, by the height inequality of the cubic tree ceiling. For general $p$, we have no such inequality. Instead, we test the mean of finitely many independent continuations of the dynamics against one Gaussian direction built from the Brownian path. The overlap penalty for a direction independent of the disorder (\cref{sel:absolute-field-cost}) then plays the role of the height inequality. That lemma rests on the external-field algorithmic threshold of Huang and Sellke \cite[arXiv version~2, Theorem~1, p.~4, and Corollary~1.8, p.~8]{hs2023}; see \cref{sel:published-field-threshold}. As the test direction is independent of the disorder, we do not need the adaptive disclosure argument that \cref{app:general-selection} uses for forgetting.

\paragraph{Setting.}
We work in the setting of \cref{app:general-selection}. We take dimension limits along the subsequence $(N_j)$ of \cref{thm:fisher-clock} and write $N\to\infty$ for $N=N_j\to\infty$. We keep every time horizon, branch number, and approximation parameter fixed while $N\to\infty$, and take long-time limits afterward.

Fix an integer $K\ge1$ and times $T>0$ and $L\ge0$. We run the diffusion \eqref{eq:hard-sde} up to time $T$, driven by the Brownian motion $(B_u)_{u\le T}$, which we call the trunk. From $X_T$, we run $K$ continuations of duration $L$, driven by Brownian motions independent of each other and of $(J,X_0,(B_u)_{u\le T})$. Let $Y_1,\dots,Y_K$ be their endpoints at time $T+L$, and let $m_{N,T}=K^{-1}\sum_{b=1}^KY_b$ be their mean, which also depends on $K$ and $L$. For each $b$, the Brownian path driving the $b$-th continuation, extended by the trunk on $[0,T]$, is a standard Brownian motion on $[0,T+L]$ independent of $(J,X_0)$. Since \eqref{eq:hard-sde} has a unique strong solution, the vector $(J,X_0,(B_u)_{u\le T},Y_b)$ has the same law as $(J,X_0,(B_u)_{u\le T},X_{T+L})$.

Let $Q_{K,L}$ and $h_{K,L}$ be given by \eqref{sel:fork-norm} and \eqref{sel:fork-center-energy}, with $e=e_\beta(\infty)$. When $q>0$, every sequence $T\to\infty$ realizes the stationary profile $(c,r)$ by \cref{sel:all-temperature}, and we may apply \cref{sel:finite-forks} along every such sequence. By \cref{lem:fork-parameter-limits}\ref{item:fork-K-limit}, we have, for every fixed $L\ge0$,
\[
 \lim_{K\to\infty}Q_{K,L}=c(2L)=d_L,\qquad \lim_{K\to\infty}h_{K,L}=h_L.
\]
We now show that the quantity $\delta_L$ defined at the start of \cref{resp:chapter} satisfies
\begin{equation}
 \delta_L=\Eth-\frac{h_L}{d_L^{p/2}}\ge0,
 \qquad \lim_{L\to\infty}\delta_L=0.
 \label{resp:p-center-deficit}
\end{equation}
The inequality is \eqref{eq:threshold-input-ceiling}; it also follows by letting $K\to\infty$ in \eqref{sel:fork-center-threshold}. The plateau is marginal by \cref{sel:positive-marginality}, so \eqref{sel:fork-saturation} applies, and it gives the limit, because $h_{K,L}/Q_{K,L}^{p/2}\to h_L/d_L^{p/2}$ as $K\to\infty$. For $p\in\{3,4\}$, the quantity $\delta_L$ is the same as in \eqref{sel:center-deficit}.

\paragraph{Two lemmas.}
The first lemma converts the absolute-overlap cost of \cref{sel:absolute-field-cost} into a bound on the overlap with a Gaussian vector from an independent seed. We write $B_N=\{x\in\R^N:\norm x\le\sqrt N\}$. Recall that $\bd=\bd(p)$ is a constant of the model, defined in \cref{sec:intro-picture}; it does not depend on the inverse temperature $\beta$.

\begin{lemma}[Gaussian seed overlap has a threshold cost]
\label{resp:seed-cost}
Fix $p\ge3$ and $0\le L_0<\infty$. For each $N$, let $\zeta_N$ be a random element independent of the coefficient tensor $J$, let $G_N$ be a standard Gaussian vector in $\R^N$ that is a measurable function of $\zeta_N$, and let $A_N=A_N(J,\zeta_N)\in B_N$ be measurable. Suppose that, for almost every value of $\zeta_N$, the map $J\mapsto A_N(J,\zeta_N)$ is globally $L_0$-Lipschitz, for the Euclidean distances on the ordered coefficient tensors and on $\R^N$. Let $\delta\ge0$, and suppose that
\[
 \liminf_{N\to\infty}\frac{\E H_N(A_N)}N\ge\Eth-\delta.
\]
Then
\begin{equation}
 \limsup_{N\to\infty}\E\Bigl|\frac{\ip{A_N}{G_N}}N\Bigr|\le\bd^{-1/p}\,\delta^{1/p}.
 \label{resp:seed-cost-bound}
\end{equation}
The constant $\bd^{-1/p}$ depends neither on $L_0$ nor on the law of $\zeta_N$. The same holds when $N$ runs through a subsequence.
\end{lemma}

The algorithm $A_N$ may depend on $G_N$; only $J$ must be independent of the seed. If $A_N$ also depends on parameters held fixed while $N\to\infty$, the lemma applies at each of their values, and its conclusion passes to limits in these parameters taken afterward.

\begin{proof}
Let $w_N=\sqrt N\,G_N/\norm{G_N}$ if $G_N\ne0$, and $w_N=\sqrt N\,e_1$ otherwise. Then $w_N$ is a measurable function of $\zeta_N$ with $\norm{w_N}=\sqrt N$. Since $\norm{A_N}\le\sqrt N$, we have
\[
 \frac{|\ip{A_N}{G_N}-\ip{A_N}{w_N}|}N\le\frac{\norm{G_N-w_N}}{\sqrt N}=\Bigl|\frac{\norm{G_N}}{\sqrt N}-1\Bigr|.
\]
For $x\ge0$, we have $|x-1|\le|x^2-1|$. The variable $\norm{G_N}^2/N$ has mean one and variance $2/N$, and the expectation of the right side is at most $\E|\norm{G_N}^2/N-1|\le\sqrt{2/N}$, which tends to zero.

We use \cref{sel:absolute-field-cost} conditionally on $\zeta_N$, with direction $w_N$, which is $\zeta_N$-measurable, and Lipschitz constant $L_0$. Writing $a_N=\ip{A_N}{w_N}/N$, we have
\[
 \limsup_{N\to\infty}\Bigl\{\frac{\E H_N(A_N)}N+\bd\,\E|a_N|^p\Bigr\}\le\Eth.
\]
Since $\limsup_N(x_N+y_N)\ge\liminf_Nx_N+\limsup_Ny_N$, the hypothesis on the energy implies that $\bd\limsup_N\E|a_N|^p\le\delta$. Hölder's inequality gives $\E|a_N|\le(\E|a_N|^p)^{1/p}$, and hence $\limsup_N\E|a_N|\le(\delta/\bd)^{1/p}$. Combining this with the first paragraph proves \eqref{resp:seed-cost-bound}. The constants in \cref{sel:absolute-field-cost} depend neither on the seed nor on the direction.
\end{proof}

The second lemma identifies the overlap of the mean $m_{N,T}$ with a Gaussian vector built from the trunk. It is a weighted form of the Brownian response identity of \cref{lem:brown-force-response} and \cref{prop:hard-bridge}.

\begin{lemma}[A weighted Brownian response identity]
\label{resp:weighted-brownian}
Fix $K\ge1$, $T>0$, and $L\ge0$, and let $k:[0,T]\to\R$ be a continuous deterministic function with $\int_0^Tk(u)^2\,\dd u=1$. Let $G_N=\int_0^Tk(u)\,\dd B_u$, where $B$ is the trunk Brownian motion. Then $G_N$ is a standard Gaussian vector in $\R^N$, independent of the disorder, and
\begin{equation}
 \lim_{N\to\infty}\E\Bigl[\frac{\ip{m_{N,T}}{G_N}}N\Bigr]
 =\sqrt2\int_0^Tk(u)R(T+L,u)\,\dd u.
 \label{resp:weighted-brownian-identity}
\end{equation}
\end{lemma}

\begin{proof}
The coordinates of $G_N$ are Wiener integrals of the independent coordinates of $B$. They are independent centered Gaussian variables with variance $\int_0^Tk^2=1$, and they are independent of $(J,X_0)$ because $B$ is.

Let $t=T+L$. Since $G_N$ is a function of the trunk, the equality of laws noted above gives $\E\ip{Y_b}{G_N}=\E\ip{X_t}{G_N}$ for each $b$. Averaging over $b$, we have $\E\ip{m_{N,T}}{G_N}=\E\ip{X_t}{G_N}$. Suppose first that $k=\sum_{i=1}^nk_i\mathbf1_{(u_{i-1},u_i]}$ is a step function, with $0=u_0<\dots<u_n=T$. Then
\[
 \frac{\E\ip{X_t}{G_N}}N=\sum_{i=1}^nk_i\,\frac{\E\ip{X_t}{B_{u_i}-B_{u_{i-1}}}}N
 =\sqrt2\sum_{i=1}^nk_i\bigl[\E\Gamma_N(t,u_i)-\E\Gamma_N(t,u_{i-1})\bigr],
\]
with $\Gamma_N$ as in \cref{prop:hard-bridge}. That proposition states that $\E\Gamma_N(t,u)\to\int_0^uR(t,w)\,\dd w$, and the right side converges to $\sqrt2\int_0^Tk(u)R(t,u)\,\dd u$.

For a general continuous $k$ and $\varepsilon>0$, we choose a step function $k_\varepsilon$ with $\sup_{[0,T]}|k-k_\varepsilon|\le\varepsilon$. We have $\norm{X_t}^2=N$ and, by the Itô isometry, $\E\norm{\int_0^T(k-k_\varepsilon)\,\dd B}^2=N\int_0^T(k-k_\varepsilon)^2$. Using the Cauchy--Schwarz inequality, we find
\[
 \frac1N\Bigl|\E\Bigl\langle X_t,\int_0^T(k-k_\varepsilon)\,\dd B\Bigr\rangle\Bigr|
 \le\Bigl(\int_0^T(k-k_\varepsilon)^2\Bigr)^{1/2}\le\sqrt T\,\varepsilon,
\]
uniformly in $N$. Also $|\sqrt2\int_0^T(k-k_\varepsilon)R(t,\cdot)|\le\sqrt2\,T\varepsilon$, since $0\le R\le1$. Using the step-function case for $k_\varepsilon$, we obtain
\[
 \limsup_{N\to\infty}\Bigl|\frac{\E\ip{X_t}{G_N}}N-\sqrt2\int_0^Tk(u)R(t,u)\,\dd u\Bigr|\le(\sqrt T+\sqrt2\,T)\varepsilon.
\]
Letting $\varepsilon\downarrow0$ proves \eqref{resp:weighted-brownian-identity}.
\end{proof}

The Brownian observable goes back to Ben Arous, Dembo, and Guionnet \cite[arXiv version~1, equation~(1.9), p.~4, and Theorem~1.2, p.~5]{bag2006}, and its hard-sphere limit is \cite[arXiv version~5, Definition~3, p.~22, and Theorem~4.3, p.~23]{sellke2024}, which underlies \cref{prop:hard-bridge}. By \cref{lem:brown-force-response}, the left side of \eqref{resp:weighted-brownian-identity} at finite $N$ equals $\sqrt2\int_0^Tk(u)\,\bar r_N(T+L,u)\,\dd u$ when $k$ is a step function; this is the Cameron--Martin identity for a deterministic force with profile $k$. We need only its limit, which \cref{prop:hard-bridge} provides without a pointwise limit of the finite-dimensional response.

\paragraph{The remote response.}
We now prove the remaining case of \eqref{eq:resp-remote-vanishing}.

\begin{proposition}[Remote squared response for general pure models]\label{prop:remote-response-general}
Fix $p\ge3$ and $\beta>0$, and suppose that the plateau $q$ is positive. Then, for every $L>0$,
\begin{equation}\label{eq:resp-general-remote-bound}
 \limsup_{t\to\infty}\int_0^{t-L}R(t,u)^2\,\dd u\le\frac{d_L}2\,\bd^{-2/p}\,\delta_L^{2/p}.
\end{equation}
In particular, \eqref{eq:resp-remote-vanishing} holds.
\end{proposition}

The idea is as follows. We choose the weight $k$ in \cref{resp:weighted-brownian} proportional to the remote response row $u\mapsto R(T+L,u)$ on $[0,T]$. Then the overlap of $m_{N,T}$ with $G_N$ is $\sqrt{2S_2(T+L,L)}$ in the limit. After normalization by $\sqrt{Q_{K,L}}$, the mean $m_{N,T}$ has energy close to $\Eth$, and \cref{resp:seed-cost} forces this overlap to be small. The coefficient-Lipschitz approximations of \cref{sel:finite-forks} make \cref{resp:seed-cost} applicable.

\begin{proof}
Fix $L>0$ and $K\ge1$, and abbreviate $Q=Q_{K,L}$ and $\delta_{K,L}=\Eth-h_{K,L}/Q^{p/2}$. Recall from \eqref{sel:fork-center-threshold} that $\delta_{K,L}\ge0$, and note that $Q\ge c(2L)\ge q>0$. The proof has four steps. We take the limits in the following order:
\begin{enumerate}[label=(\roman*)]
\item $N=N_j\to\infty$, with $T$, $K$, $L$, and the approximation parameter $\varepsilon$ fixed;
\item $T\to\infty$, along an arbitrary sequence;
\item $\varepsilon\downarrow0$;
\item $K\to\infty$.
\end{enumerate}

\paragraph{Step 1: the test direction.}
For $T>0$, the response is positive by \eqref{eq:resp-response-lower}, and $S_2(T+L,L)=\int_0^TR(T+L,u)^2\,\dd u>0$. We define
\[
 k_{T,L}(u)=\frac{R(T+L,u)}{S_2(T+L,L)^{1/2}},\qquad 0\le u\le T,
 \qquad G_N=\int_0^Tk_{T,L}(u)\,\dd B_u.
\]
The weight $k_{T,L}$ is continuous and deterministic: it depends on the CK solution and on the fixed times $T$ and $L$, not on the disorder or the dimension. It satisfies $\int_0^Tk_{T,L}^2=1$. \Cref{resp:weighted-brownian} implies that
\begin{equation}\label{eq:resp-test-overlap}
 \lim_{N\to\infty}\E\Bigl[\frac{\ip{m_{N,T}}{G_N}}N\Bigr]
 =\sqrt2\,\frac{\int_0^TR(T+L,u)^2\,\dd u}{S_2(T+L,L)^{1/2}}
 =\sqrt{2S_2(T+L,L)}.
\end{equation}

\paragraph{Step 2: the normalized mean and its Lipschitz approximation.}
Let $\Pi_{B_N}$ be the metric projection onto $B_N$, and let $\bar m_{N,T}=\Pi_{B_N}(m_{N,T}/\sqrt Q)$. For $\varepsilon>0$, let $\mathcal A^\varepsilon_{N,T}$ be the maps of \cref{sel:finite-forks}, so that
\[
 \limsup_{N\to\infty}\frac1N\E\norm{\mathcal A^\varepsilon_{N,T}-\bar m_{N,T}}^2\le\varepsilon^2,
\]
as in \eqref{sel:fork-lipschitz}. These maps take values in $B_N$. Conditionally on the initial point and on all Brownian paths, they are globally Lipschitz in $J$, with a constant that depends neither on $N$ nor on the seed. The vector $G_N$ is a function of the trunk Brownian path. It is therefore measurable with respect to this seed, which is independent of $J$.

We next compare energies. The vectors $m_{N,T}/\sqrt Q$, $\bar m_{N,T}$, and $\mathcal A^\varepsilon_{N,T}$ have norm at most $\sqrt{N/q}$. By homogeneity and \eqref{sel:tensor-ui}, we have $\norm{\nabla H_N(x)}\le pQ_Nq^{-(p-1)/2}\sqrt N$ when $\norm x\le\sqrt{N/q}$, where $Q_N$ is the tensor norm \eqref{eq:tensor-norm}. Integrating along segments and applying the Cauchy--Schwarz inequality, we find that, for random vectors $x,y$ of norm at most $\sqrt{N/q}$,
\[
 \Bigl|\frac{\E H_N(x)-\E H_N(y)}N\Bigr|\le D_q\Bigl(\frac{\E\norm{x-y}^2}N\Bigr)^{1/2},
 \qquad D_q=pq^{-(p-1)/2}\sup_N(\E Q_N^2)^{1/2},
\]
and $D_q<\infty$ by \cref{lem:landscape-moments}. We set
\[
 \gamma_T=\limsup_{N\to\infty}\Bigl|\frac{\E H_N(m_{N,T})}N-h_{K,L}\Bigr|,
 \qquad
 \kappa_T=\limsup_{N\to\infty}\Bigl(\frac{\E\norm{\bar m_{N,T}-m_{N,T}/\sqrt Q}^2}N\Bigr)^{1/2}.
\]
Note that $H_N(m_{N,T}/\sqrt Q)=Q^{-p/2}H_N(m_{N,T})$ by homogeneity. Comparing energies through the intermediate vector $\bar m_{N,T}$, we obtain
\begin{equation}\label{eq:resp-algorithm-energy}
\begin{aligned}
 \liminf_{N\to\infty}\frac{\E H_N(\mathcal A^\varepsilon_{N,T})}N
 &\ge Q^{-p/2}(h_{K,L}-\gamma_T)-D_q(\varepsilon+\kappa_T)\\
 &=\Eth-\delta_{K,L}-Q^{-p/2}\gamma_T-D_q(\varepsilon+\kappa_T).
\end{aligned}
\end{equation}

\paragraph{Step 3: the seed cost.}
By Step 2, the maps $\mathcal A^\varepsilon_{N,T}$ satisfy the hypotheses of \cref{resp:seed-cost}. The seed is formed by the initial point and all Brownian paths, the Gaussian vector is the vector $G_N$ of Step 1, and $\delta=\delta_{K,L}+Q^{-p/2}\gamma_T+D_q(\varepsilon+\kappa_T)$ by \eqref{eq:resp-algorithm-energy}. \Cref{resp:seed-cost} implies that
\[
 \limsup_{N\to\infty}\E\Bigl|\frac{\ip{\mathcal A^\varepsilon_{N,T}}{G_N}}N\Bigr|
 \le\bd^{-1/p}\bigl(\delta_{K,L}+Q^{-p/2}\gamma_T+D_q(\varepsilon+\kappa_T)\bigr)^{1/p}.
\]
To return to $m_{N,T}$, we use $\E\norm{G_N}^2=N$ and the Cauchy--Schwarz inequality; the fact that $G_N$ and the approximation share the trunk seed plays no role here. We have
\[
 \Bigl|\E\frac{\ip{m_{N,T}/\sqrt Q-\mathcal A^\varepsilon_{N,T}}{G_N}}N\Bigr|
 \le\Bigl(\frac{\E\norm{m_{N,T}/\sqrt Q-\bar m_{N,T}}^2}N\Bigr)^{1/2}
 +\Bigl(\frac{\E\norm{\bar m_{N,T}-\mathcal A^\varepsilon_{N,T}}^2}N\Bigr)^{1/2}.
\]
Combining the last two displays with \eqref{eq:resp-test-overlap}, we find that, for every $T>0$,
\begin{equation}\label{eq:resp-fixed-T-bound}
 \sqrt{\frac{2S_2(T+L,L)}Q}
 \le\bd^{-1/p}\bigl(\delta_{K,L}+Q^{-p/2}\gamma_T+D_q(\varepsilon+\kappa_T)\bigr)^{1/p}+\varepsilon+\kappa_T.
\end{equation}

\paragraph{Step 4: the long-time limits.}
Let $T_n\to\infty$ be an arbitrary sequence. Using \eqref{sel:fork-center-energy-limit}, we have $\gamma_{T_n}\to0$. For $\kappa_T$, note that $\norm{\Pi_{B_N}(x)-x}=(\norm x-\sqrt N)_+$, and that $|\sqrt{y/Q}-1|\le|y-Q|/Q$ for $y\ge0$. With $y=\norm{m_{N,T}}^2/N$, we find $\kappa_T^2\le\limsup_N\E|\norm{m_{N,T}}^2/N-Q|^2/Q^2$, and \eqref{sel:fork-center-concentration} implies that $\kappa_{T_n}\to0$. The left side of \eqref{eq:resp-fixed-T-bound} does not depend on $\varepsilon$. Taking $T=T_n\to\infty$ in \eqref{eq:resp-fixed-T-bound} and then $\varepsilon\downarrow0$, we obtain
\[
 \limsup_{n\to\infty}S_2(T_n+L,L)\le\frac Q2\,\bd^{-2/p}\,\delta_{K,L}^{2/p}.
\]
Since the sequence $T_n$ was arbitrary, the same bound holds for $\limsup_{t\to\infty}S_2(t,L)$. This upper limit does not depend on $K$. As $K\to\infty$, we have $Q_{K,L}\to d_L$ and $\delta_{K,L}\to\delta_L$, by \cref{lem:fork-parameter-limits}\ref{item:fork-K-limit}. This proves \eqref{eq:resp-general-remote-bound}. As $d_L\le1$ and $\delta_L\to0$ by \eqref{resp:p-center-deficit}, the right side of \eqref{eq:resp-general-remote-bound} tends to zero as $L\to\infty$, which is \eqref{eq:resp-remote-vanishing}.
\end{proof}

Together with Steps~1, 2, and~4 of the proof of \cref{resp:strong-row}, \cref{prop:remote-response-general} proves strong $L^2$ convergence of the response rows, the exact limits of the squared tails, and $\chi(t,0)=o(\sqrt t)$ for every pure model with $p\ge3$.

\chapter{Persistent integrated violation of the fluctuation--dissipation relation}\label{resp:integrated-chapter}

By \cref{resp:backward-fdt}, the backward FDT defect $E(t,s)=\partial_sC(t,s)-R(t,s)$ tends to zero in $L^2(0,t)$ as $t\to\infty$. In this section, we prove that, above the dynamical transition, the defect nevertheless does not disappear. Its positive part keeps an integrated mass bounded below by a positive constant on the part of the row away from any fixed initial interval $[0,S]$ and any fixed recent strip $[t-L,t]$. This is what we mean by a persistent integrated violation: the defect is small in $L^2$, but it has positive integrated mass, spread over a history whose length grows with $t$. The two statements are compatible. By the Cauchy--Schwarz inequality, a mass $m$ spread over an interval of length $\ell$ forces a squared norm of at least $m^2/\ell$, and here $\ell$ grows like $t$.

The main result is \cref{resp:persistent-defect}, and its proof has two ingredients. The first is \cref{resp:weighted-ratio}. The energy identity determines the response integrated against the weight $pC(t,u)^{p-1}$ over the whole row, and stationary convergence determines the contribution of a recent strip. The remainder identifies one weighted ratio of the slow response with the number $x=(p-2)(1-q)/q$. The second ingredient is \cref{resp:transport}, an elementary inequality for one absolutely continuous function, which converts a deficit in this weighted response budget into a lower bound on the unweighted positive defect. This inequality is sharp (\cref{rem:transport-sharp}).

The proofs use the energy identity \eqref{eq:energy-identity}, the selected plateau and energy of \cref{sel:selected-sector} together with stationary convergence \eqref{sel:stationary-window}, forgetting \eqref{sel:fixed-wait-decay}, the absolute continuity of correlation rows from \cref{thm:fisher-clock}, and \cref{resp:backward-fdt}. The results of this section prove \eqref{eq:main-weighted-ratio} and \eqref{eq:main-defect-budget} in \cref{thm:main-response}. We return to them in \cref{sec:discussion}.

The physical mechanism is the one described by Cugliandolo, Dean, and Kurchan in \cite[arXiv version~2, equations~(12)--(14), p.~3, equation~(18), p.~4, and the integrated-bounds discussion on p.~3]{cdk1997}. Entropy production bounds the violation of the fluctuation--dissipation relation, and the violation is therefore small at late times; its integral over a time interval that grows with the age of the system may nevertheless remain of order one. In our setting, the smallness at late times is \cref{resp:backward-fdt}, which rests on the entropy bound of \cref{thm:fisher-clock}. The lower bound on the integrated violation for the CK solution is proved here, from the energy identity.

\paragraph{Standing assumptions and notation.}
Throughout this section, we fix an integer $p\ge3$ and $\beta>\bd$. By \cref{sel:selected-sector}, the plateau $q=q_\beta$ is the larger root of
\begin{equation}\label{eq:viol-marginality}
 \beta^2p(p-1)q^{p-2}(1-q)^2=1,
\end{equation}
and the limiting energy is
\begin{equation}\label{eq:viol-energy}
 e_\infty=e_\beta(\infty)=\lim_{t\to\infty}e_\beta(t)=\beta\bigl[1+(p-2)q^{p-1}-(p-1)q^p\bigr].
\end{equation}
Since $\beta>\bd$, we have $q>\qd=(p-2)/(p-1)$, as noted before \cref{sel:selected-sector}. Indeed, the function $g(y)=y^{p-2}(1-y)^2$ is decreasing on $[q_{\mathrm m},1]$, where $q_{\mathrm m}=(p-2)/p<\qd$, and $\beta^2p(p-1)g(\qd)=\beta^2/\bd^2>1$. The larger root of \eqref{eq:viol-marginality} therefore lies in $(\qd,1)$. We set
\begin{equation}\label{eq:viol-x-dstar}
 x=\frac{(p-2)(1-q)}q,\qquad d_*=1-\frac{e_\infty}\beta.
\end{equation}
The inequality $q>\qd$ is equivalent to $(p-2)(1-q)<q$, which means that $x\in(0,1)$. Using \eqref{eq:viol-energy}, we find
\[
 d_*=(p-1)q^p-(p-2)q^{p-1}=q^p\Bigl[(p-1)-\frac{p-2}q\Bigr]=(1-x)q^p>0,
\]
where we used $(p-1)-(p-2)/q=1-(p-2)(1-q)/q$. Equivalently,
\begin{equation}\label{eq:viol-old-budget}
 \frac{e_\infty}\beta-(1-q^p)=q^p-d_*=xq^p.
\end{equation}

We write $E(t,s)=\partial_sC(t,s)-R(t,s)$ for $t>s>0$, as in \cref{sec:resp-earlier-derivative}, and $E_+=\max\{E,0\}$. For $t>0$, we use the right earlier derivative at $s=0$ and set $E(t,t)=0$, consistently with the left earlier derivative at the diagonal. Recall from \cref{thm:fisher-clock} that, for each $t$, the row $u\mapsto C(t,u)$ is absolutely continuous on $[0,t]$, and that its derivative on $(0,t)$ is the continuous function $\partial_uC(t,u)$. In particular, $E(t,\cdot)$ is integrable on $(0,t)$. Let $(c,r)$ be the stationary profile, with $r=-c'$, $c(0)=1$, and $c$ decreasing to $q$. Stationary convergence \eqref{sel:stationary-window} implies that, for every fixed $A\ge0$,
\begin{equation}\label{eq:viol-recent-strip}
 \lim_{t\to\infty}\sup_{0\le v\le A}\bigl\{|C(t,t-v)-c(v)|+|R(t,t-v)-r(v)|\bigr\}=0.
\end{equation}

\section{A weighted slow ratio}\label{sec:viol-weighted-ratio}

\begin{proposition}[An exact weighted slow ratio]
\label{resp:weighted-ratio}
For each fixed $A>0$, the denominator below is positive for all large $t$, the ratio has a limit as $t\to\infty$, and the subsequent limit as $A\to\infty$ is
\begin{equation}
 \lim_{A\to\infty}\lim_{t\to\infty}
 \frac{p\int_0^{t-A}C(t,u)^{p-1}R(t,u)\,\dd u}
      {C(t,t-A)^p-C(t,0)^p}=x.
 \label{resp:weighted-ratio-limit}
\end{equation}
Further, the weighted signed defect has the full-row limit
\begin{equation}
 \lim_{t\to\infty}p\int_0^tC(t,u)^{p-1}E(t,u)\,\dd u=d_*.
 \label{resp:weighted-defect}
\end{equation}
\end{proposition}

The idea is as follows. By the energy identity, the total weighted response $p\int_0^tC^{p-1}R$ tends to $e_\infty/\beta$. A recent strip of length $A$ contributes $1-c(A)^p$, as for an equilibrium profile, and its contribution tends to $1-q^p$ as $A\to\infty$. By \eqref{eq:viol-old-budget}, the older part of the row carries the remaining weighted response $xq^p$, while the correlation raised to the power $p$ falls from about $q^p$ to zero across it.

\begin{proof}
By the energy identity \eqref{eq:energy-identity}, we have $p\int_0^tC(t,u)^{p-1}R(t,u)\,\dd u=e_\beta(t)/\beta$, which tends to $e_\infty/\beta=1-d_*$. Fix $A>0$. Using the change of variables $u=t-v$ and \eqref{eq:viol-recent-strip}, we find
\[
\begin{aligned}
 p\int_{t-A}^tC(t,u)^{p-1}R(t,u)\,\dd u&=p\int_0^AC(t,t-v)^{p-1}R(t,t-v)\,\dd v\\
 &\longrightarrow p\int_0^Ac(v)^{p-1}r(v)\,\dd v=1-c(A)^p,
\end{aligned}
\]
where we used $r=-c'$ and $c(0)=1$ in the last step. Subtracting, we see that the numerator in \eqref{resp:weighted-ratio-limit} converges to $c(A)^p-d_*$. We have $C(t,0)\to0$ by \eqref{sel:fixed-wait-decay} and $C(t,t-A)\to c(A)$ by \eqref{eq:viol-recent-strip}. The denominator therefore converges to $c(A)^p\ge q^p>0$; in particular, it is positive for all large $t$. The ratio converges to
\[
 \frac{c(A)^p-d_*}{c(A)^p}=1-\frac{d_*}{c(A)^p}.
\]
As $A\to\infty$, we have $c(A)\to q$, and this tends to $1-d_*/q^p=x$. The two limits are taken in this order, and no interchange of limits occurs.

For \eqref{resp:weighted-defect}, fix $t>0$. The row $u\mapsto C(t,u)$ is absolutely continuous on $[0,t]$ with values in $[0,1]$. It follows that $u\mapsto C(t,u)^p$ is absolutely continuous, with derivative $pC(t,u)^{p-1}\partial_uC(t,u)$ almost everywhere. Integrating over $[0,t]$ and using $C(t,t)=1$, we have
\[
 p\int_0^tC(t,u)^{p-1}\partial_uC(t,u)\,\dd u=1-C(t,0)^p.
\]
Subtracting the energy identity, we arrive at the exact identity
\[
 p\int_0^tC(t,u)^{p-1}E(t,u)\,\dd u=1-C(t,0)^p-\frac{e_\beta(t)}\beta,
\]
whose right side tends to $1-(1-d_*)=d_*$. This argument does not require the row to be monotone.
\end{proof}

\begin{remark}[Interpretation of $x$]\label{rem:viol-x-interpretation}
In the asymptotic solution of Cugliandolo and Kurchan, the number $x=(p-2)(1-q)/q$ is the value of the fluctuation--dissipation ratio in the aging regime; see the paragraph after \cite[arXiv version~1, Equation~(12), p.~6]{ck1993}. There, for late pairs of times with correlation below the plateau, the response is predicted to be $x$ times the earlier derivative of the correlation, in the normalization of this paper. \Cref{resp:weighted-ratio} identifies $x$ with one ratio of weighted integrals over the older part of the row. It does not show that the ratio $R(t,s)/\partial_sC(t,s)$ converges to $x$ along late pairs; we discuss that statement, and its integrated form, in \cref{sec:discussion}.
\end{remark}

\section{Transport of the response budget}\label{sec:viol-transport}

We now remove the weight $C^{p-1}$. On the interval $[S,t-L]$, the correlation $u\mapsto C(t,u)$ rises from about zero to about $c(L)$. The weighted response budget $p\int_S^{t-L}C^{p-1}R$ on this interval tends to $c(L)^p-d_*$, which is smaller than $c(L)^p$. If the defect were nonpositive on the interval, the correlation could not rise faster than the response allows, and this budget would control the increase of $C(t,\cdot)^p$. The following lemma quantifies the failure of this control in terms of the positive part of the defect. No monotonicity is required.

To see the idea, let $z\ge0$ be absolutely continuous and $r_0\ge0$. If $z'\le r_0$ almost everywhere, then $(z^p)'=pz^{p-1}z'\le pz^{p-1}r_0$, which implies that $z(v)^p\le z(u)^p+p\int_u^vz^{p-1}r_0$. The lemma bounds the positive part of $z'-r_0$ from below when this inequality fails.

\begin{lemma}[Response-budget transport]
\label{resp:transport}
Let $u<v$, let $z:[u,v]\to[0,\infty)$ be absolutely continuous, and let $r_0\ge0$ be integrable on $[u,v]$. Set
\[
 J=\int_u^vz(s)^{p-1}r_0(s)\,\dd s,
 \qquad e_0=z'-r_0,
 \qquad D=\Bigl[z(v)-\bigl(z(u)^p+pJ\bigr)^{1/p}\Bigr]_+.
\]
Then
\begin{equation}
 \int_u^v(e_0)_+\,\dd s\ge D,
 \qquad (v-u)\int_u^v(e_0)_+^2\,\dd s\ge D^2.
 \label{resp:transport-bound}
\end{equation}
\end{lemma}

Here $J$ is finite because $z$ is bounded on $[u,v]$, and $e_0$ is integrable because $z'$ and $r_0$ are. If $(e_0)_+$ is not square integrable, the left side of the second inequality is infinite.

\begin{proof}
Let $\bar z=(z(u)^p+pJ)^{1/p}$. If $\bar z\ge z(v)$, then $D=0$, and both claims hold because $(e_0)_+\ge0$. Suppose that $\bar z<z(v)$, so that $D=z(v)-\bar z>0$.

We first treat the case $\bar z>0$. For $a>0$, we define the clipped power weight and its primitive by
\[
 w_a(y)=\min\Bigl\{\Bigl(\frac ya\Bigr)^{p-1},1\Bigr\},
 \qquad
 V_a(y)=\begin{cases}
 \dfrac{y^p}{pa^{p-1}},&0\le y\le a,\\[2mm]
 y-\dfrac{p-1}pa,&y\ge a.
 \end{cases}
\]
The function $V_a$ is continuously differentiable on $[0,\infty)$, and its derivative $V_a'=w_a$ takes values in $[0,1]$. Then $V_a\circ z$ is absolutely continuous, with derivative $w_a(z)z'$ almost everywhere. We use two pointwise bounds. As $0\le w_a\le1$, we have $(e_0)_+\ge w_a(z)e_0$. As $r_0\ge0$ and $w_a(y)\le(y/a)^{p-1}$, we have $w_a(z)r_0\le z^{p-1}r_0/a^{p-1}$. Integrating, we obtain
\[
 \int_u^v(e_0)_+\,\dd s\ge\int_u^vw_a(z)z'\,\dd s-\int_u^vw_a(z)r_0\,\dd s
 \ge V_a(z(v))-V_a(z(u))-\frac J{a^{p-1}}.
\]
We now take $a=\bar z$ and compute the right side. From $z(u)^p\le\bar z^p$, we have $z(u)\le\bar z$ and $V_{\bar z}(z(u))=z(u)^p/(p\bar z^{p-1})$. From $z(v)>\bar z$, we have $V_{\bar z}(z(v))=z(v)-(p-1)\bar z/p$. Using $z(u)^p+pJ=\bar z^p$, we find
\[
\begin{aligned}
 V_{\bar z}(z(v))-V_{\bar z}(z(u))-\frac J{\bar z^{p-1}}
 &=z(v)-\frac{p-1}p\,\bar z-\frac{z(u)^p+pJ}{p\,\bar z^{p-1}}\\
 &=z(v)-\frac{p-1}p\,\bar z-\frac{\bar z}p=z(v)-\bar z=D.
\end{aligned}
\]
This is the first inequality in \eqref{resp:transport-bound}. The primitive $V_{\bar z}$ is evaluated only at the endpoints, and excursions and returns of $z$ inside $[u,v]$ do not affect the bound.

If $\bar z=0$, then $z(u)=0$ and $J=0$. For $0<\varepsilon<z(v)$, the computation above with $a=\varepsilon$ gives
\[
 \int_u^v(e_0)_+\,\dd s\ge V_\varepsilon(z(v))-V_\varepsilon(0)-0=z(v)-\frac{p-1}p\,\varepsilon.
\]
Letting $\varepsilon\downarrow0$, we find $\int_u^v(e_0)_+\ge z(v)=D$. Finally, the Cauchy--Schwarz inequality gives $D^2\le(\int_u^v(e_0)_+)^2\le(v-u)\int_u^v(e_0)_+^2$, which is the second inequality.
\end{proof}

\section{The persistent defect}\label{sec:viol-persistent}

\begin{theorem}[A positive slow defect and a physical rate constraint]
\label{resp:persistent-defect}
Set
\[
 m_L=c(L)-\bigl[c(L)^p-d_*\bigr]^{1/p},
 \qquad m_\infty=q\bigl(1-x^{1/p}\bigr)>0.
\]
For every fixed $S\ge0$ and $L>0$,
\begin{align}
 \liminf_{t\to\infty}\int_S^{t-L}E_+(t,u)\,\dd u
       &\ge m_L,\label{resp:persistent-l1}\\
 \liminf_{t\to\infty}(t-S-L)
                \int_S^{t-L}E_+(t,u)^2\,\dd u
       &\ge m_L^2.\label{resp:persistent-l2}
\end{align}
Further, $m_L\uparrow m_\infty$ as $L\to\infty$. In particular,
\begin{equation}
 \lim_{t\to\infty}\int_0^tE_+(t,u)^2\,\dd u=0,
 \qquad
 \liminf_{t\to\infty}t\int_0^tE_+(t,u)^2\,\dd u
       \ge q^2\bigl(1-x^{1/p}\bigr)^2>0.
 \label{resp:physical-rate}
\end{equation}
\end{theorem}

The backward FDT defect tends to zero in $L^2(0,t)$, but not at a rate $o(t^{-1/2})$: by \eqref{resp:physical-rate}, we have $\liminf_t\sqrt t\,\norm{E(t,\cdot)}_{L^2(0,t)}\ge m_\infty$. The constant $m_L$ is the limit of the transport distance $D$ of \cref{resp:transport} on the interval $[S,t-L]$; we show in the proof that its radicand $c(L)^p-d_*$ is positive.

\begin{proof}
We first compute the limits of the endpoint values and of the weighted response budget on $[S,t-L]$. We then apply \cref{resp:transport} on that interval. Finally, we study $m_L$ and deduce \eqref{resp:physical-rate}.

\paragraph{Step 1: the budget on the remote part of the row.}
Fix $S\ge0$ and $L>0$, and let $t>S+L$. We set
\[
 J_t=\int_S^{t-L}C(t,u)^{p-1}R(t,u)\,\dd u.
\]
As in the proof of \cref{resp:weighted-ratio}, with $A=L$, we have $p\int_0^{t-L}C(t,u)^{p-1}R(t,u)\,\dd u\to c(L)^p-d_*$. The omitted initial interval is negligible. Indeed, since $0\le C\le1$, we have
\[
 0\le p\int_0^SC(t,u)^{p-1}R(t,u)\,\dd u\le pS\sup_{0\le u\le S}R(t,u)\longrightarrow0,
\]
by \eqref{sel:fixed-wait-decay}. We conclude that $pJ_t\to c(L)^p-d_*$. Further, $C(t,S)\to0$ by \eqref{sel:fixed-wait-decay}, and $C(t,t-L)\to c(L)$ by \eqref{eq:viol-recent-strip}.

\paragraph{Step 2: transport.}
We use \cref{resp:transport} on $[S,t-L]$ with $z(u)=C(t,u)$ and $r_0(u)=R(t,u)$. The function $z$ is absolutely continuous with values in $[0,1]$, $r_0$ is nonnegative and bounded, and $e_0=z'-r_0=E(t,\cdot)$ almost everywhere on $(S,t-L)$. The transport distance is
\[
 D_t=\Bigl[C(t,t-L)-\bigl(C(t,S)^p+pJ_t\bigr)^{1/p}\Bigr]_+.
\]
By Step 1 and the continuity of $y\mapsto y^{1/p}$ and of $y\mapsto y_+$, we have
\[
 \lim_{t\to\infty}D_t=\Bigl[c(L)-\bigl(c(L)^p-d_*\bigr)^{1/p}\Bigr]_+=m_L.
\]
The radicand is positive: $c(L)\ge q$ and \eqref{eq:viol-old-budget} give $c(L)^p-d_*\ge q^p-d_*=xq^p>0$. It is also smaller than $c(L)^p$, because $d_*>0$. Hence $m_L>0$. Using \cref{resp:transport}, we obtain
\[
 \int_S^{t-L}E_+(t,u)\,\dd u\ge D_t,
 \qquad
 (t-L-S)\int_S^{t-L}E_+(t,u)^2\,\dd u\ge D_t^2.
\]
Taking the lower limit as $t\to\infty$ proves \eqref{resp:persistent-l1} and \eqref{resp:persistent-l2}.

\paragraph{Step 3: the constants and the rate.}
Let $\varphi(y)=y-(y^p-d_*)^{1/p}$ for $y>d_*^{1/p}$. Then
\[
 \varphi'(y)=1-y^{p-1}\bigl(y^p-d_*\bigr)^{(1-p)/p}
 =1-\Bigl(\frac{y^p}{y^p-d_*}\Bigr)^{(p-1)/p}<0,
\]
and $\varphi$ is strictly decreasing. We have $q^p>d_*$, and $c(L)$ decreases to $q$ as $L\to\infty$. It follows that $m_L=\varphi(c(L))$ increases to
\[
 \varphi(q)=q-\bigl(q^p-d_*\bigr)^{1/p}=q-\bigl(xq^p\bigr)^{1/p}=q\bigl(1-x^{1/p}\bigr)=m_\infty,
\]
where we used \eqref{eq:viol-old-budget}. As $x<1$, we have $m_\infty>0$.

We have $0\le E_+\le|E|$, and \cref{resp:backward-fdt} implies that $\int_0^tE_+(t,u)^2\,\dd u\to0$. For the lower bound in \eqref{resp:physical-rate}, fix $L>0$ and take $S=0$. For $t>L$, we have
\[
 t\int_0^tE_+(t,u)^2\,\dd u\ge(t-L)\int_0^{t-L}E_+(t,u)^2\,\dd u.
\]
The lower limit of the right side is at least $m_L^2$, by \eqref{resp:persistent-l2}. Letting $L\to\infty$, we find $\liminf_tt\int_0^tE_+^2\ge m_\infty^2=q^2(1-x^{1/p})^2$. This completes the proof.
\end{proof}

The lower limit in \eqref{resp:persistent-l1} is nonincreasing in $L$, because the integrand is nonnegative and the interval shrinks as $L$ grows. It therefore has a limit in $[0,\infty]$ as $L\to\infty$. Since $m_{L'}\ge m_L$ for $L'\ge L$, we obtain, for every fixed $S\ge0$,
\begin{equation}\label{eq:viol-iterated-l1}
 \lim_{L\to\infty}\liminf_{t\to\infty}\int_S^{t-L}E_+(t,u)\,\dd u\ge m_\infty.
\end{equation}
Together with \eqref{resp:physical-rate}, this is \eqref{eq:main-defect-budget}.

\begin{remark}[Sharpness of the transport inequality]\label{rem:transport-sharp}
The constant in \cref{resp:transport} is sharp for the problem with prescribed endpoint values, interval length, and weighted response budget. Let nonnegative endpoint values $z_0,z_1$ and a budget $J\ge0$ be given, with $\bar z=(z_0^p+pJ)^{1/p}<z_1$, and let $0<\varepsilon<v-u$. On $[u,u+\varepsilon]$, we let $z$ increase linearly from $z_0$ to $\bar z$ and set $r_0=z'$. Then $e_0=0$ there, and
\[
 \int_u^{u+\varepsilon}z^{p-1}r_0\,\dd s=\int_u^{u+\varepsilon}z^{p-1}z'\,\dd s=\frac{\bar z^p-z_0^p}p=J.
\]
On $[u+\varepsilon,v]$, we let $z$ increase linearly from $\bar z$ to $z_1$ and set $r_0=0$. Then $e_0=(z_1-\bar z)/(v-u-\varepsilon)$ there, and
\[
 \int_u^v(e_0)_+\,\dd s=z_1-\bar z=D,
 \qquad
 (v-u)\int_u^v(e_0)_+^2\,\dd s=\frac{v-u}{v-u-\varepsilon}\,D^2.
\]
The first inequality in \eqref{resp:transport-bound} is an equality for each $\varepsilon$, and the second becomes an equality as $\varepsilon\downarrow0$. This optimality concerns the inequality of \cref{resp:transport} for the stated class of functions. We do not know whether the constants $m_L$ and $m_\infty$ are optimal for the CK solution.
\end{remark}


\part{Three-time geometry in the cubic model}\label{part:geometry}
\chapter{Saturated tree vectors and the product inequality}\label{sec:geometry-results}

This section and the next concern the cubic model. In this section, we study unit vectors of the tree Hilbert space whose cubic tree energy is close to the ceiling $\Eth$ of \cref{thm:geometry-tree-ceiling}. We prove two results about such vectors. The first, \cref{thm:geometry-universal-root}, states that the root operator $\mathcal I^*v$ of a vector close to the ceiling is close, in suitable coordinates, to one universal integral operator on $L^2(0,1)$. The second, \cref{thm:geometry-threshold-product}, is a product inequality for three such vectors whose overlaps are nonnegative at every height. \Cref{lem:geometry-coherent-root,lem:geometry-pair-deficit} are the intermediate steps.

The section is deterministic. It uses the tree Hilbert space of \cref{sec:geometry-core} and the cubic ceiling of \cref{thm:geometry-tree-ceiling}, with the retained defects \eqref{eq:geometry-height-defect} and \eqref{eq:geometry-angular-young-defects} proved in \cref{sec:tree-ceilings}. Nothing about the CK equations is used. In \cref{sec:three-time}, we apply both theorems to the normalized prefix centers of the cubic CK solution: \cref{thm:geometry-threshold-product} gives the three-time product inequality of \cref{thm:geometry-actual-product}, and \cref{thm:geometry-universal-root} is the main input of the susceptibility bound in \cref{thm:geometry-susceptibility-lower}.

\paragraph{Notation.}
Throughout this section, $p=3$, so that the grafting map $\mathcal I\colon\mathcal H\otimes\mathcal H\to\mathcal H$ is binary. Recall from \cref{sec:geometry-core} that $\mathcal H$ is a real separable Hilbert space, that $\mathcal I$ is an isometry, and that its range is the orthogonal complement of the one-leaf summand. The one-leaf summand is the subspace of height zero, and we have
\begin{equation}\label{eq:geometry-graft-range}
 \mathcal I\mathcal I^*=\Id-\Pi_0 .
\end{equation}
We write $\Pi_j$ for the projection onto trees of height exactly $j$ and $\Pi_{\le H}=\sum_{j=0}^H\Pi_j$; the projections $\Pi_j$ are orthogonal and sum to the identity. As in \cref{sec:geometry-core}, we identify a tensor $x\otimes y$ with the Hilbert--Schmidt operator $|x\rangle\langle y|$, which maps $z$ to $\ip{y}{z} x$. Then $\mathcal I^*v$ is a Hilbert--Schmidt operator on $\mathcal H$, and $\norm{\mathcal I^*v}_{\HS}\le\norm{v}$. We write $\ip{A}{B}_{\HS}=\tr(A^*B)$, so that
\[
 \bigl\langle |a\rangle\langle b|,|c\rangle\langle d|\bigr\rangle_{\HS}
 =\ip{a}{c}\ip{b}{d} .
\]
For $v\in\mathcal H$, let
\[
 f(v)=\ip{v}{\mathcal I(v\otimes v)}
 =\bigl\langle\mathcal I^*v,|v\rangle\langle v|\bigr\rangle_{\HS},
 \qquad F_*=\frac{2\sqrt2}{3},
\]
and, for a unit vector $v$, let $\delta(v)=F_*-f(v)$ be its \emph{deficit}. By \cref{thm:geometry-tree-ceiling}, we have $\delta(v)>0$ for every unit vector $v$. The cubic tree energy of \eqref{eq:geometry-tree-energy} is $h(v)=\sqrt3 f(v)$, and $\Eth=\sqrt3F_*$. Hence the deficit $\delta(v)=(\Eth-h(v))/\sqrt3$ is the energy deficit divided by $\sqrt3$.

For a unit vector $v$, we use the height masses of \cref{thm:geometry-tree-ceiling}:
\[
 v_j=\Pi_jv,\qquad b_j=\norm{v_j}^2,\qquad B_j=\sum_{i\le j}b_i,\qquad
 a_j=\sqrt{b_j},\qquad \bar b=\sup_{j\ge0}b_j,
\]
with $B_{-1}=B_{-2}=0$. Note that $B_j\uparrow1$ as $j\to\infty$.

\paragraph{The height rule.}
A graft has height one plus the larger of the heights of its two children. This means that $\mathcal I$ maps $\Pi_i\mathcal H\otimes\Pi_k\mathcal H$ into $\Pi_{\max(i,k)+1}\mathcal H$, and
\begin{equation}\label{eq:geometry-height-rule}
 \Pi_j\mathcal I(x\otimes y)=\sum_{\max(i,k)=j-1}\mathcal I(\Pi_ix\otimes\Pi_ky)\quad(j\ge1),
 \qquad \Pi_0\mathcal I=0 .
\end{equation}
Taking adjoints, we find that $\mathcal I^*\Pi_0=0$ and, for $j\ge1$,
\begin{equation}\label{eq:geometry-root-blocks}
 \mathcal I^*\Pi_jv=\sum_{\max(i,k)=j-1}\Pi_i(\mathcal I^*v)\Pi_k .
\end{equation}
In words, the height-$j$ component of $v$ is carried by the blocks $\Pi_i(\mathcal I^*v)\Pi_k$ of the root operator with $\max(i,k)=j-1$. We call $j$ the \emph{root height} of these blocks. Blocks with different root heights are orthogonal in the Hilbert--Schmidt inner product, and the restriction of $\mathcal I^*$ to $\Pi_j\mathcal H$ is an isometry for $j\ge1$, by \eqref{eq:geometry-graft-range}.

\paragraph{The retained defects.}
Let $v$ be a unit vector with deficit $\delta$. By \eqref{eq:geometry-height-rule} and the isometry of $\mathcal I$ (see also \cref{lem:tree-graft-heights}), the projected graft $g_j=\Pi_j\mathcal I(v\otimes v)$ satisfies
\begin{equation}\label{eq:geometry-graft-norm}
 \norm{g_j}^2=\sum_{\max(i,k)=j-1}b_ib_k=B_{j-1}^2-B_{j-2}^2=b_{j-1}s_j^2,
 \qquad s_j=\sqrt{B_{j-1}+B_{j-2}} .
\end{equation}
Let $\theta_j$ be the angle between $v_j$ and $g_j$, with the convention $\cos\theta_j=1$ when $v_j=0$ or $g_j=0$. We use three consequences of \cref{thm:geometry-tree-ceiling}. The first inequality in \eqref{eq:geometry-height-defect} implies that
\begin{equation}\label{eq:geometry-mesh-bound}
 \bar b\le\Bigl(\sum_{j\ge0}b_j^2\Bigr)^{1/2}\le(4\sqrt2\,\delta)^{1/2}=2^{5/4}\delta^{1/2},
\end{equation}
and \eqref{eq:geometry-angular-young-defects} states that
\begin{equation}\label{eq:geometry-retained-recalled}
 \sum_{j\ge1}a_ja_{j-1}s_j(1-\cos\theta_j)\le\delta,
 \qquad
 \frac12\sum_{j\ge1}s_j(a_j-a_{j-1})^2\le\delta .
\end{equation}

\section{The universal operator at a saturated root}

By \eqref{eq:geometry-mesh-bound}, a vector close to the ceiling has small mass at every height. The retained defects \eqref{eq:geometry-retained-recalled} give more information: they bound weighted sums of the variations $a_j-a_{j-1}$ of the height amplitudes and of the angular errors $1-\cos\theta_j$ between $v_j$ and the projected graft $g_j$ of the lower components. The next lemma shows that, as a consequence, the root operator of $v$ is close to an operator whose coefficients in the normalized height directions $v_j/\sqrt{b_j}$, for $b_j>0$, depend only on the height masses.

\begin{lemma}[Coherent approximation of a root]
\label{lem:geometry-coherent-root}
Let $v$ be a unit tree vector with $\delta=\delta(v)\le1$. There exists a symmetric Hilbert--Schmidt operator $\widetilde A_v$ with $\norm{\widetilde A_v}_{\HS}=1$ and height blocks
\begin{equation}
 \Pi_i\widetilde A_v\Pi_k
 =\frac{|v_i\rangle\langle v_k|}{\sqrt{B_h+B_{h-1}}},\qquad h=\max(i,k),
 \label{eq:geometry-coherent-root-blocks}
\end{equation}
where a zero denominator means a zero block, such that
\begin{equation}
 \norm{\mathcal I^*v-\widetilde A_v}_{\HS}\le K\delta^{1/4}.
 \label{eq:geometry-coherent-root-error}
\end{equation}
Here $K$ is an absolute constant; one can take $K=5$.
\end{lemma}

\begin{proof}
The proof has three steps. First, we define $\widetilde A_v$ one root height at a time and compute its norm. Second, we compute the error at each root height. Third, we sum the errors: at root heights above a small cumulative mass $\varepsilon$, we use the retained defects; the remaining root heights form an initial segment of small mass.

\paragraph{Step 1: definition and normalization.}
For $j\ge1$, let $g_j$ and $s_j$ be as in \eqref{eq:geometry-graft-norm}. We define the component of $\widetilde A_v$ at root height $j$ by
\[
 \widetilde A^{(j)}=\frac{\mathcal I^*g_j}{s_j}\quad\text{if } s_j>0,
 \qquad \widetilde A^{(j)}=0\quad\text{if } s_j=0,
\]
and we set $\widetilde A_v=\sum_{j\ge1}\widetilde A^{(j)}$. Using \eqref{eq:geometry-height-rule} and $\mathcal I^*\mathcal I=\Id$, we have
\[
 \mathcal I^*g_j=\sum_{\max(i,k)=j-1}|v_i\rangle\langle v_k| .
\]
With $h=j-1$, the denominator $s_j$ equals $\sqrt{B_h+B_{h-1}}$, and the blocks of $\widetilde A_v$ are those in \eqref{eq:geometry-coherent-root-blocks}. The block with indices $(k,i)$ is the adjoint of the block with indices $(i,k)$, and therefore $\widetilde A_v$ is symmetric. By \eqref{eq:geometry-graft-norm} and the isometry of $\mathcal I$, we have $\norm{\widetilde A^{(j)}}_{\HS}^2=\norm{g_j}^2/s_j^2=b_{j-1}$ when $s_j>0$. When $s_j=0$, we have $B_{j-1}=0$, and in particular $b_{j-1}=0$. The components with different root heights are orthogonal. Summing over $j$, we find
\[
 \norm{\widetilde A_v}_{\HS}^2=\sum_{j\ge1}b_{j-1}=\sum_{h\ge0}b_h=1 .
\]

\paragraph{Step 2: the error at one root height.}
By \eqref{eq:geometry-root-blocks}, we have $\mathcal I^*v=\sum_{j\ge1}\mathcal I^*v_j$, and $\mathcal I^*v_j$ lies in the same root-height component as $\widetilde A^{(j)}$. By orthogonality,
\[
 \norm{\mathcal I^*v-\widetilde A_v}_{\HS}^2=\sum_{j\ge1}\norm{\mathcal I^*v_j-\widetilde A^{(j)}}_{\HS}^2 .
\]
For $j\ge1$, the vector $v_j$ lies in the range of $\mathcal I$, and \eqref{eq:geometry-graft-range} gives $\norm{\mathcal I^*v_j}_{\HS}^2=b_j$. If $s_j>0$, then, since $g_j$ also lies in the range of $\mathcal I$,
\[
 \bigl\langle\mathcal I^*v_j,\widetilde A^{(j)}\bigr\rangle_{\HS}
 =\frac{\ip{v_j}{\mathcal I\mathcal I^*g_j}}{s_j}
 =\frac{\ip{v_j}{g_j}}{s_j}
 =a_ja_{j-1}\cos\theta_j,
\]
where we used $\norm{g_j}=a_{j-1}s_j$ from \eqref{eq:geometry-graft-norm}. The error at root height $j$ is
\begin{equation}\label{eq:geometry-root-height-error}
 \norm{\mathcal I^*v_j-\widetilde A^{(j)}}_{\HS}^2
 =b_j+b_{j-1}-2a_ja_{j-1}\cos\theta_j
 =(a_j-a_{j-1})^2+2a_ja_{j-1}(1-\cos\theta_j).
\end{equation}
If $s_j=0$, then $b_{j-1}=0$ and the error is $b_j$, in agreement with \eqref{eq:geometry-root-height-error}.

\paragraph{Step 3: summation.}
If $\delta=1$, then \eqref{eq:geometry-coherent-root-error} holds with $K=2$, because both operators have Hilbert--Schmidt norm at most one. We assume from now on that $\delta<1$. Fix $\varepsilon\in(0,1)$.

We first consider the root heights $j\ge1$ with $B_{j-2}\ge\varepsilon$. For these, $s_j^2=B_{j-1}+B_{j-2}\ge2\varepsilon$. Multiplying each term of \eqref{eq:geometry-root-height-error} by $s_j/\sqrt{2\varepsilon}\ge1$ and using \eqref{eq:geometry-retained-recalled}, we obtain
\[
 \sum_{j:\,B_{j-2}\ge\varepsilon}\norm{\mathcal I^*v_j-\widetilde A^{(j)}}_{\HS}^2
 \le\frac1{\sqrt{2\varepsilon}}\sum_{j\ge1}s_j\Bigl[(a_j-a_{j-1})^2+2a_ja_{j-1}(1-\cos\theta_j)\Bigr]
 \le\frac{4\delta}{\sqrt{2\varepsilon}} .
\]

The remaining root heights form an initial segment. Let $m_\varepsilon=\min\{k\ge0:B_k\ge\varepsilon\}$, which is finite because $B_k\uparrow1$. The condition $B_{j-2}<\varepsilon$ holds exactly when $1\le j\le m_\varepsilon+1$, because $B$ is nondecreasing. On this segment,
\[
 \sum_{j=1}^{m_\varepsilon+1}b_j\le B_{m_\varepsilon-1}+b_{m_\varepsilon}+b_{m_\varepsilon+1}\le\varepsilon+2\bar b,
 \qquad
 \sum_{j=1}^{m_\varepsilon+1}b_{j-1}=B_{m_\varepsilon}\le\varepsilon+\bar b .
\]
Using $\norm{x-y}^2\le2\norm{x}^2+2\norm{y}^2$ together with $\norm{\mathcal I^*v_j}_{\HS}^2=b_j$ and $\norm{\widetilde A^{(j)}}_{\HS}^2=b_{j-1}$, we find that the errors on this segment sum to at most $4\varepsilon+6\bar b$.

Adding the two bounds and using \eqref{eq:geometry-mesh-bound}, we obtain
\[
 \norm{\mathcal I^*v-\widetilde A_v}_{\HS}^2
 \le4\varepsilon+6\cdot2^{5/4}\delta^{1/2}+\frac{4\delta}{\sqrt{2\varepsilon}} .
\]
We choose $\varepsilon=\delta^{1/2}$. Then $\delta^{3/4}\le\delta^{1/2}$, because $\delta<1$, and the right side is at most $(4+6\cdot2^{5/4}+2\sqrt2)\delta^{1/2}\le25\delta^{1/2}$. This proves \eqref{eq:geometry-coherent-root-error} with $K=5$.
\end{proof}

When the height mesh $\bar b$ is small, the denominator in \eqref{eq:geometry-coherent-root-blocks} is close to $\sqrt{2B_h}$, and $B_h$ varies almost continuously with the height. In the coordinate $x=B_h$ on $(0,1)$, the blocks then become the values of one function of two variables. This is the content of the next theorem.

\begin{theorem}[Universal threshold root]
\label{thm:geometry-universal-root}
Let $\mathcal K$ be a fixed separable infinite-dimensional Hilbert space, and let $\mathbf1$ denote the constant function one in $L^2(0,1)$. If $\norm{v_n}=1$ and $f(v_n)\to F_*$, then there exist isometric embeddings $V_n\colon\mathcal H\to L^2(0,1)\oplus\mathcal K$ satisfying $V_nv_n=\mathbf1$, for which
\begin{equation}
 V_n\mathcal I^*v_nV_n^*\longrightarrow\mathcal T\oplus0
 \quad\text{in Hilbert--Schmidt norm},
 \qquad
 T(x,y)=\frac1{\sqrt{2\max(x,y)}} .
 \label{eq:geometry-universal-root}
\end{equation}
Here $\mathcal T$ is the integral operator on $L^2(0,1)$ with kernel $T$. The operator $\mathcal T$ is positive and trace class, with the positive frame
\begin{equation}
 \mathcal T=\frac1{\sqrt2}|\mathbf1\rangle\langle\mathbf1|
 +\frac1{2\sqrt2}\int_0^1z^{-3/2}
 |\mathbf1_{[0,z]}\rangle\langle\mathbf1_{[0,z]}|\,\dd z .
 \label{eq:geometry-universal-root-frame}
\end{equation}
In particular, $\tr\mathcal T=\sqrt2$ and $\norm{\mathcal T}_{\HS}=1$.
\end{theorem}

\begin{remark}[Interpretation of the universal kernel]\label{rem:geometry-universal-kernel}
In the coordinates of \cref{thm:geometry-universal-root}, a point $x\in(0,1)$ stands for the height at which the cumulative squared mass of $v_n$ reaches $x$. The kernel says that, at the threshold, the root of $v_n$ couples two children at cumulative masses $x$ and $y$ with a weight that depends only on the larger of the two, that is, on the taller child. In this continuum description, the weight $(2\max(x,y))^{-1/2}$ is, up to null sets, the unique nonnegative function of $\max(x,y)$ that gives squared root mass $\dd x$ to each band of cumulative mass $\dd x$: the pairs $(x',y')$ with $\max(x',y')\in[x,x+\dd x]$ have area $2x\,\dd x$, and $T^2\cdot2x\,\dd x=\dd x$. (At a finite root height $j$, the coherent approximation of \cref{lem:geometry-coherent-root} carries the squared mass $b_{j-1}$ of the height below.) The threshold value itself is the quadratic form of $\mathcal T$ on the constant function:
\[
 \ip{\mathbf1}{\mathcal T\mathbf1}=2\int_0^1\frac{x}{\sqrt{2x}}\,\dd x=\frac{2\sqrt2}3=F_* .
\]
This is consistent with $f(v_n)=\langle\mathcal I^*v_n,|v_n\rangle\langle v_n|\rangle_{\HS}\to F_*$. The convergence \eqref{eq:geometry-universal-root} holds in Hilbert--Schmidt norm only, and we do not claim that the traces of the operators $\mathcal I^*v_n$ converge to $\tr\mathcal T$. For this reason, \cref{sec:geometry-susceptibility} uses the frame \eqref{eq:geometry-universal-root-frame} only through the lower semicontinuity of the trace.
\end{remark}

\begin{proof}[Proof of \cref{thm:geometry-universal-root}]
The proof has six steps. After a bound on the height mesh, we embed each $v_n$ so that its heights become consecutive intervals of $(0,1)$. In these coordinates, the coherent approximation of \cref{lem:geometry-coherent-root} becomes an integral operator with a step kernel. We then replace the step kernel by a simpler one and show that it converges to $T$. In the last step, we prove the properties of $\mathcal T$.

We write $\delta_n=\delta(v_n)$, and $b_{n,j}$, $B_{n,j}$, and $\bar b_n$ for the height masses of $v_n$. When no confusion can arise, we drop the index $n$.

\paragraph{Step 1: the height mesh.}
By \eqref{eq:geometry-mesh-bound}, we have $\bar b_n\le2^{5/4}\delta_n^{1/2}$, which tends to zero because $\delta_n\to0$.

\paragraph{Step 2: the embedding.}
For each $j$ with $b_j>0$, let $\hat v_j=v_j/\sqrt{b_j}$ and $J_j=(B_{j-1},B_j)$. The vectors $\hat v_j$ are orthonormal, and so are the functions $b_j^{-1/2}\mathbf1_{J_j}$, because the intervals $J_j$ are disjoint and have lengths $b_j$. Let $E_n$ be the closed span of the vectors $\hat v_j$, and let $W_n\colon E_n\to L^2(0,1)$ be the isometry with $W_n\hat v_j=b_j^{-1/2}\mathbf1_{J_j}$. The orthogonal complement $E_n^\perp$ is separable, because $\mathcal H$ is separable (it is a countable direct sum of tensor powers of a separable space). It is also infinite-dimensional, because it contains every vector of the one-leaf summand that is orthogonal to $\hat v_0$ (or the whole summand if $b_0=0$), and the one-leaf summand is infinite-dimensional by \cref{sec:tree-space}. We may therefore choose a unitary map $U_n\colon E_n^\perp\to\mathcal K$. We define $V_n(x+y)=W_nx\oplus U_ny$ for $x\in E_n$ and $y\in E_n^\perp$. Then $V_n$ is an isometry, and
\[
 V_nv_n=\sum_{j:\,b_j>0}\sqrt{b_j}\,W_n\hat v_j=\sum_{j:\,b_j>0}\mathbf1_{J_j}=\mathbf1,
\]
since the intervals $J_j$ cover $(0,1)$ up to a countable set. Conjugation by $V_n$ preserves Hilbert--Schmidt norms, positivity, and traces, because $V_n^*V_n$ is the identity of $\mathcal H$. Further, $V_n|x\rangle\langle y|V_n^*=|V_nx\rangle\langle V_ny|$.

\paragraph{Step 3: the coherent approximation as a step kernel.}
Let $\widetilde A_n=\widetilde A_{v_n}$ be the operator of \cref{lem:geometry-coherent-root}, which applies once $\delta_n\le1$. Since $V_nv_{n,i}=\mathbf1_{J_i}$, the blocks \eqref{eq:geometry-coherent-root-blocks} are transported to
\[
 V_n\frac{|v_{n,i}\rangle\langle v_{n,k}|}{\sqrt{B_h+B_{h-1}}}V_n^*
 =\frac{|\mathbf1_{J_i}\rangle\langle\mathbf1_{J_k}|}{\sqrt{B_h+B_{h-1}}},
 \qquad h=\max(i,k).
\]
In other words, $V_n\widetilde A_nV_n^*$ is the integral operator on $L^2(0,1)$, extended by zero on $\mathcal K$, with the step kernel
\[
 \widetilde T_n(x,y)=\frac1{\sqrt{B_h+B_{h-1}}},\qquad x\in J_i,\ y\in J_k,\ h=\max(i,k).
\]
Recall that the Hilbert--Schmidt norm of an integral operator is the $L^2$ norm of its kernel. Then \cref{lem:geometry-coherent-root} implies that
\[
 \norm{V_n\mathcal I^*v_nV_n^*-\widetilde T_n}_{\HS}=\norm{\mathcal I^*v_n-\widetilde A_n}_{\HS}\le5\delta_n^{1/4}\to0 .
\]
Here and below, we identify a kernel with its integral operator.

\paragraph{Step 4: replacement of the denominator.}
Let $T_n$ be the step kernel equal to $(2B_h)^{-1/2}$ on $J_i\times J_k$, with $h=\max(i,k)$. The set of pairs $(x,y)$ with $\max(i,k)=h$ has area $\sum_{\max(i,k)=h}b_ib_k=B_h^2-B_{h-1}^2=b_h(B_h+B_{h-1})$. Using $B_h+B_{h-1}=2B_h-b_h$, we find
\[
 \norm{\widetilde T_n-T_n}_{L^2}^2
 =\sum_{h:\,b_h>0}b_h\Bigl(1-\sqrt{1-\frac{b_h}{2B_h}}\Bigr)^2
 \le\sum_{h:\,b_h>0}\min\Bigl\{b_h,\frac{b_h^3}{4B_h^2}\Bigr\},
\]
where we used $0\le1-\sqrt{1-x}\le x$ for $x\in[0,1]$. Fix $\varepsilon\in(0,1)$. The heights with $B_h<\varepsilon$ form an initial segment, on which $\sum b_h<\varepsilon$. On the heights with $B_h\ge\varepsilon$, we use the second term in the minimum and $\sum_hb_h=1$. We obtain
\[
 \norm{\widetilde T_n-T_n}_{L^2}^2\le\varepsilon+\frac{\bar b_n^2}{4\varepsilon^2} .
\]
By Step 1, the upper limit of the right side as $n\to\infty$ is $\varepsilon$, and letting $\varepsilon\downarrow0$ gives $\norm{\widetilde T_n-T_n}_{L^2}\to0$.

\paragraph{Step 5: convergence to the universal kernel.}
Let $x\in J_i$ and $y\in J_k$, with $h=\max(i,k)$. Then $\max(x,y)<\max(B_i,B_k)=B_h$, because $x<B_i$, $y<B_k$, and $B$ is nondecreasing. If $h=i$, then $x>B_{h-1}\ge B_h-\bar b_n$, and if $h=k$, the same holds for $y$. This gives
\[
 \max(x,y)<B_h<\max(x,y)+\bar b_n .
\]
Recalling that $T_n(x,y)=(2B_h)^{-1/2}$, we find
\[
 T(x,y)\Bigl(\frac{\max(x,y)}{\max(x,y)+\bar b_n}\Bigr)^{1/2}\le T_n(x,y)\le T(x,y).
\]
The union of the rectangles $J_i\times J_k$ has full measure, and $\bar b_n\to0$ by Step 1. Hence $T_n\to T$ almost everywhere, with $0\le T_n\le T$. The kernel $T$ is square integrable:
\[
 \int_0^1\!\!\int_0^1T(x,y)^2\,\dd y\,\dd x=2\int_0^1\!\!\int_0^x\frac1{2x}\,\dd y\,\dd x=1 .
\]
By dominated convergence with the dominating function $T^2$, we have $\norm{T_n-T}_{L^2}\to0$. Combining Steps 3--5, we obtain
\[
 \norm{V_n\mathcal I^*v_nV_n^*-\mathcal T\oplus0}_{\HS}
 \le\norm{\mathcal I^*v_n-\widetilde A_n}_{\HS}+\norm{\widetilde T_n-T_n}_{L^2}+\norm{T_n-T}_{L^2}\longrightarrow0,
\]
which is \eqref{eq:geometry-universal-root}. The computation above also shows that $\norm{\mathcal T}_{\HS}=\norm{T}_{L^2}=1$.

\paragraph{Step 6: the frame.}
For $z\in(0,1]$, let $P_z=|\mathbf1_{[0,z]}\rangle\langle\mathbf1_{[0,z]}|$. It is positive, of rank one, and has trace $z$. The map $z\mapsto P_z$ is continuous in trace norm, and
\[
 \int_0^1z^{-3/2}\norm{P_z}_1\,\dd z=\int_0^1z^{-1/2}\,\dd z=2 .
\]
Thus the integral in \eqref{eq:geometry-universal-root-frame} converges in trace norm, and the right side of \eqref{eq:geometry-universal-root-frame} defines a positive trace-class operator $\mathcal T'$ with
\[
 \tr\mathcal T'=\frac1{\sqrt2}+\frac1{2\sqrt2}\int_0^1z^{-3/2}\cdot z\,\dd z=\frac1{\sqrt2}+\frac1{\sqrt2}=\sqrt2 .
\]
The kernel of $P_z$ is $\mathbf1\{\max(x,y)\le z\}$. For $f,g\in L^2(0,1)$, Fubini's theorem gives
\begin{align*}
 \ip{f}{\mathcal T'g}
 &=\int_0^1\!\!\int_0^1f(x)g(y)\Bigl[\frac1{\sqrt2}+\frac1{2\sqrt2}\int_{\max(x,y)}^1z^{-3/2}\,\dd z\Bigr]\dd x\,\dd y\\
 &=\int_0^1\!\!\int_0^1f(x)T(x,y)g(y)\,\dd x\,\dd y .
\end{align*}
Fubini's theorem applies because $|f(x)g(y)|$ times the bracket is integrable: the bracket equals $T(x,y)$, since $\int_m^1z^{-3/2}\,\dd z=2(m^{-1/2}-1)$, and $T$ is square integrable. We have shown that $\mathcal T'=\mathcal T$, which proves \eqref{eq:geometry-universal-root-frame}, the positivity of $\mathcal T$, and the value of its trace.
\end{proof}

\section{Positive saturated triples}

We next compare different vectors close to the ceiling. When two vectors have nonnegative inner products at every height, the overlaps at different heights cannot cancel. The pair-deficit lemma below shows that, for two vectors close to the ceiling, little cross mass is carried by heights at which a smoothed cosine between their prefixes is far from one. For three vectors, we choose for each of two pairs a prefix that carries almost all of the overlap of the pair and on which the two vectors are nearly parallel. These two prefixes are nested, and this gives the product inequality.

For two unit vectors $v,w$, we write $B^v_j$ and $B^w_j$ for their cumulative height masses.

\begin{lemma}[Prefix alignment measured by a pair deficit]
\label{lem:geometry-pair-deficit}
Let $v,w$ be unit tree vectors with $d_j=\ip{\Pi_jv}{\Pi_jw}\ge0$ for all $j\ge0$. Let $D_j=\sum_{i\le j}d_i$, with $D_{-1}=0$, and set
\[
 \rho_j=\frac{D_j+D_{j-1}}{\sqrt{(B_j^v+B_{j-1}^v)(B_j^w+B_{j-1}^w)}},\qquad
 \Delta(v,w)=\sum_{j\ge0}d_j(1-\rho_j),
\]
with $\rho_j=0$ when the denominator vanishes. Then $0\le\rho_j\le1$ for all $j$. If moreover $\delta(v)\le1$ and $\delta(w)\le1$, then
\begin{equation}
 0\le\Delta(v,w)\le K\bigl(\delta(v)^{1/4}+\delta(w)^{1/4}\bigr),
 \label{eq:geometry-pair-deficit}
\end{equation}
where $K$ is an absolute constant.
\end{lemma}

The quantity $\Delta(v,w)$ is an average of $1-\rho_j$ with respect to the cross masses $d_j$, whose total is $\ip{v}{w}$. The ratio $\rho_j$ is a smoothed version of $D_j/\sqrt{B^v_jB^w_j}$, the cosine of the angle between the prefixes $\Pi_{\le j}v$ and $\Pi_{\le j}w$. When the overlap of $v$ and $w$ is bounded away from zero and both deficits are small, the bound \eqref{eq:geometry-pair-deficit} says that most of the overlap of $v$ and $w$, measured by the cross masses $d_j$, is collected at heights where $\rho_j$ is close to one.

\begin{proof}
The ratio $\rho_j$ is nonnegative, because all $d_j$ are nonnegative. By the Cauchy--Schwarz inequality on prefixes, $D_j=\ip{\Pi_{\le j}v}{\Pi_{\le j}w}\le(B_j^vB_j^w)^{1/2}$, and the same holds with $j-1$ in place of $j$. The two-term Cauchy--Schwarz inequality $\sqrt{\alpha_1\beta_1}+\sqrt{\alpha_2\beta_2}\le\sqrt{(\alpha_1+\alpha_2)(\beta_1+\beta_2)}$ then gives
\[
 D_j+D_{j-1}\le\sqrt{(B_j^v+B_{j-1}^v)(B_j^w+B_{j-1}^w)} .
\]
This shows that $\rho_j\le1$, and $\Delta(v,w)\ge0$ follows.

Next, we compute the Hilbert--Schmidt inner product of the coherent approximations of \cref{lem:geometry-coherent-root}. Blocks with different index pairs $(i,k)$ are orthogonal, and the block $(i,k)$ contributes $\ip{v_i}{w_i}\ip{v_k}{w_k}=d_id_k$ divided by the two denominators in \eqref{eq:geometry-coherent-root-blocks}. For each $h$, we have $\sum_{\max(i,k)=h}d_id_k=D_h^2-D_{h-1}^2=d_h(D_h+D_{h-1})$. Therefore
\begin{equation}\label{eq:geometry-coherent-pair}
 \ip{\widetilde A_v}{\widetilde A_w}_{\HS}
 =\sum_{h\ge0}\frac{D_h^2-D_{h-1}^2}{\sqrt{(B_h^v+B_{h-1}^v)(B_h^w+B_{h-1}^w)}}
 =\sum_{h\ge0}d_h\rho_h
 =\ip{v}{w}-\Delta(v,w).
\end{equation}
If a denominator vanishes, then the corresponding components of $v$ or $w$ vanish at all heights up to $h$, and the terms on both sides are zero. The series converge absolutely, since the sum of the absolute values of the block inner products is at most $\norm{\widetilde A_v}_{\HS}\norm{\widetilde A_w}_{\HS}=1$ by the Cauchy--Schwarz inequality. In the last equality, we used $\sum_hd_h=\ip{v}{w}$.

On the other hand, by \eqref{eq:geometry-graft-range}, we have
\[
 \ip{\mathcal I^*v}{\mathcal I^*w}_{\HS}=\ip{v}{\mathcal I\mathcal I^*w}=\ip{v}{w}-d_0 .
\]
Subtracting this from \eqref{eq:geometry-coherent-pair}, we find
\[
 \Delta(v,w)=d_0+\ip{\mathcal I^*v}{\mathcal I^*w}_{\HS}-\ip{\widetilde A_v}{\widetilde A_w}_{\HS}.
\]
Using $\norm{\mathcal I^*w}_{\HS}\le1$, $\norm{\widetilde A_v}_{\HS}=1$, and \cref{lem:geometry-coherent-root}, we bound the difference of the two inner products by
\[
 \norm{\mathcal I^*v-\widetilde A_v}_{\HS}\norm{\mathcal I^*w}_{\HS}+\norm{\widetilde A_v}_{\HS}\norm{\mathcal I^*w-\widetilde A_w}_{\HS}
 \le5\delta(v)^{1/4}+5\delta(w)^{1/4} .
\]
The leaf overlap satisfies $d_0\le(b_0^v+b_0^w)/2$, and \eqref{eq:geometry-mesh-bound} implies that $b_0^v\le2^{5/4}\delta(v)^{1/2}\le2^{5/4}\delta(v)^{1/4}$, and similarly for $w$. Combining these bounds, we obtain \eqref{eq:geometry-pair-deficit} with $K=7$.
\end{proof}

We now state the product inequality. Note that if $\ip{\Pi_hv_i}{\Pi_hv_j}\ge0$ for all $h$, then $\ip{v_i}{v_j}=\sum_h\ip{\Pi_hv_i}{\Pi_hv_j}$ lies in $[0,1]$.

\begin{theorem}[Product inequality for positive threshold states]
\label{thm:geometry-threshold-product}
For $\delta\in(0,1]$, let $\omega(\delta)$ be the supremum of
\[
 \bigl[\ip{v_1}{v_3}\ip{v_2}{v_3}-\ip{v_1}{v_2}\bigr]_+
\]
over all triples of unit tree vectors $v_1,v_2,v_3$ such that $\ip{\Pi_hv_i}{\Pi_hv_j}\ge0$ for all $h\ge0$ and $i,j\in\{1,2,3\}$, and such that $\max_i\delta(v_i)\le\delta$. Then $\omega(\delta)\to0$ as $\delta\downarrow0$. In other words, every such triple satisfies
\begin{equation}
 \ip{v_1}{v_2}\ge
 \ip{v_1}{v_3}\ip{v_2}{v_3}-\omega\Bigl(\max_i\delta(v_i)\Bigr),
 \label{eq:geometry-threshold-product}
\end{equation}
with an error that tends to zero uniformly over such triples as their maximal deficit tends to zero.
\end{theorem}

\begin{proof}
The proof has four steps. We fix $\kappa\in(0,1]$ and first treat triples with $\ip{v_1}{v_3}\ge\kappa$ and $\ip{v_2}{v_3}\ge\kappa$. For each $i\in\{1,2\}$, we choose a height $H_i$ such that the prefix $\Pi_{\le H_i}$ contains almost all of the overlap between $v_i$ and $v_3$, and such that the prefixes of $v_i$ and $v_3$ up to $H_i$ are nearly parallel. The prefixes up to $H_1$ and $H_2$ are nested, which lets us compare $v_1$ and $v_2$ on the shorter one. The last step treats triples with a small overlap.

Throughout, we write
\[
 c_{ij}=\ip{v_i}{v_j},\qquad d_{ij,h}=\ip{\Pi_hv_i}{\Pi_hv_j},\qquad
 D_{ij,h}=\sum_{k\le h}d_{ij,k},\qquad B_{i,h}=\sum_{k\le h}\norm{\Pi_kv_i}^2 .
\]
We let $\rho_{i3,h}$ be the ratio of \cref{lem:geometry-pair-deficit} for the pair $(v_i,v_3)$, and set
\[
 \bar b=\sup_{i,h}\norm{\Pi_hv_i}^2,\qquad
 \varepsilon=\bar b+\Delta(v_1,v_3)+\Delta(v_2,v_3),\qquad \eta=\sqrt\varepsilon .
\]
Let $\delta=\max_i\delta(v_i)\le1$. By \eqref{eq:geometry-mesh-bound} and \cref{lem:geometry-pair-deficit}, we have
\begin{equation}\label{eq:geometry-triple-epsilon}
 \varepsilon\le2^{5/4}\delta^{1/2}+28\delta^{1/4}\le K_0\delta^{1/4}
\end{equation}
with the absolute constant $K_0=31$. In Steps 1--3, we assume that $c_{13},c_{23}\ge\kappa$ and that $\eta\le\kappa/4$; in particular, $\varepsilon\le1$, so that $\varepsilon\le\eta\le\varepsilon^{1/4}$.

\paragraph{Step 1: the choice of the grades $H_i$.}
Fix $i\in\{1,2\}$. As $h\to\infty$, the partial sums $D_{i3,h}$ increase to $c_{i3}$. Let $h_0$ be the first height with $D_{i3,h_0}\ge c_{i3}-\eta$. The heights $h\ge h_0$ carry the cross mass
\[
 \sum_{h\ge h_0}d_{i3,h}=c_{i3}-D_{i3,h_0-1}>\eta,
\]
since either $h_0=0$ and $c_{i3}\ge\kappa>\eta$, or $D_{i3,h_0-1}<c_{i3}-\eta$. We claim that there exists $H_i\ge h_0$ with
\begin{equation}\label{eq:geometry-grade-choice}
 d_{i3,H_i}>0\qquad\text{and}\qquad 1-\rho_{i3,H_i}\le\frac{2\Delta(v_i,v_3)}\eta\le2\eta .
\end{equation}
Indeed, if $1-\rho_{i3,h}>2\Delta(v_i,v_3)/\eta$ for every $h\ge h_0$ with $d_{i3,h}>0$, then
\[
 \Delta(v_i,v_3)\ge\sum_{h\ge h_0}d_{i3,h}(1-\rho_{i3,h})>\frac{2\Delta(v_i,v_3)}\eta\cdot\eta=2\Delta(v_i,v_3),
\]
which is impossible, since $\Delta(v_i,v_3)\ge0$. (The strict inequality holds because some $d_{i3,h}$ with $h\ge h_0$ is positive.) The second inequality in \eqref{eq:geometry-grade-choice} follows from $\Delta(v_i,v_3)\le\varepsilon=\eta^2$. This is the reason for the choice of $H_i$: the pair deficit is an average of $1-\rho_{i3,h}$ with weights $d_{i3,h}$, the alignment can fail only at heights where the cross mass is small, and the heights beyond $h_0$ carry cross mass at least $\eta$. Because $H_i\ge h_0$, we also have
\begin{equation}\label{eq:geometry-grade-mass}
 D_{i3,H_i}\ge c_{i3}-\eta .
\end{equation}

\paragraph{Step 2: alignment on the prefix up to $H_i$.}
Let $x_i=\Pi_{\le H_i}v_i$ and $z_i=\Pi_{\le H_i}v_3$, and write
\[
 \Gamma=D_{i3,H_i}=\ip{x_i}{z_i},\qquad X=B_{i,H_i}=\norm{x_i}^2,\qquad Z=B_{3,H_i}=\norm{z_i}^2 .
\]
By \eqref{eq:geometry-grade-mass}, we have $\Gamma\ge\kappa-\kappa/4=3\kappa/4$. Using $\Gamma\le\sqrt{XZ}$ and $X,Z\le1$, we find $X,Z\ge\Gamma^2\ge9\kappa^2/16$. The last increments $\xi=\norm{\Pi_{H_i}v_i}^2$, $\zeta=\norm{\Pi_{H_i}v_3}^2$, and $\gamma=d_{i3,H_i}$ satisfy $\xi,\zeta\le\bar b\le\varepsilon$ and $\gamma\ge0$, and
\[
 \rho_{i3,H_i}=\frac{\Gamma-\gamma/2}{\sqrt{(X-\xi/2)(Z-\zeta/2)}}\le\frac{\Gamma}{\sqrt{XZ}}\cdot\frac1{1-\sigma},
 \qquad \sigma=\frac{8\varepsilon}{9\kappa^2}\le\frac1{18},
\]
where we used $(X-\xi/2)/X\ge1-\sigma$, $(Z-\zeta/2)/Z\ge1-\sigma$, and $\varepsilon\le\kappa^2/16$. Using $\Gamma/\sqrt{XZ}\le1$ and $1/(1-\sigma)\le1+2\sigma$, we have $\rho_{i3,H_i}\le\Gamma/\sqrt{XZ}+2\sigma$. Combining this with \eqref{eq:geometry-grade-choice}, we obtain
\[
 1-\frac{\Gamma}{\sqrt{XZ}}\le(1-\rho_{i3,H_i})+2\sigma\le2\eta+\frac{16\varepsilon}{9\kappa^2}\le\frac{4\eta}{\kappa^2} .
\]
Let $\alpha_i=\Gamma/Z$. Then
\[
 \norm{x_i-\alpha_iz_i}^2=X-\frac{\Gamma^2}Z=X\Bigl(1-\frac{\Gamma^2}{XZ}\Bigr)\le2\Bigl(1-\frac{\Gamma}{\sqrt{XZ}}\Bigr)\le\frac{8\eta}{\kappa^2},
\]
and $\norm{\alpha_iz_i}=\Gamma/\sqrt Z\le\sqrt X\le1$. We have established that
\begin{equation}\label{eq:geometry-prefix-alignment}
 x_i=\alpha_iz_i+e_i,\qquad \norm{e_i}\le\frac{2\sqrt2}{\kappa}\varepsilon^{1/4},\qquad\norm{\alpha_iz_i}\le1\qquad(i=1,2).
\end{equation}

\paragraph{Step 3: comparison on the shorter prefix.}
By symmetry in the indices $1$ and $2$, we may assume that $H_1\le H_2$. Applying $\Pi_{\le H_1}$ to \eqref{eq:geometry-prefix-alignment} with $i=2$, and using $\Pi_{\le H_1}z_2=\Pi_{\le H_1}v_3=z_1$, we find $\Pi_{\le H_1}v_2=\alpha_2z_1+e_2'$ with $\norm{e_2'}\le\norm{e_2}$. Because the overlaps $d_{12,h}$ are all nonnegative, we have
\[
 c_{12}=\sum_hd_{12,h}\ge\sum_{h\le H_1}d_{12,h}=\ip{x_1}{\Pi_{\le H_1}v_2}
 =\ip{\alpha_1z_1+e_1}{\alpha_2z_1+e_2'} .
\]
We expand the last inner product. We have $\norm{\alpha_1z_1}\le1$, $\norm{\alpha_2z_1}\le\norm{\alpha_2z_2}\le1$, and $\norm{e_1}\le\norm{x_1}+\norm{\alpha_1z_1}\le2$. The three cross terms are therefore bounded in absolute value by $\norm{e_2'}+\norm{e_1}+2\norm{e_2'}\le(8\sqrt2/\kappa)\varepsilon^{1/4}$, and
\[
 c_{12}\ge\alpha_1\alpha_2B_{3,H_1}-\frac{8\sqrt2}{\kappa}\varepsilon^{1/4} .
\]
For the main term, we use \eqref{eq:geometry-grade-mass}, $B_{3,H_2}\le1$, and $c_{13},c_{23}\le1$:
\[
 \alpha_1\alpha_2B_{3,H_1}=D_{13,H_1}\frac{D_{23,H_2}}{B_{3,H_2}}
 \ge(c_{13}-\eta)(c_{23}-\eta)\ge c_{13}c_{23}-2\eta\ge c_{13}c_{23}-2\varepsilon^{1/4}.
\]
Combining the last two displays, we find
\begin{equation}\label{eq:geometry-product-quantitative}
 c_{12}\ge c_{13}c_{23}-\frac{14}{\kappa}\varepsilon^{1/4}
 \qquad\text{whenever } c_{13},c_{23}\ge\kappa\text{ and }\varepsilon\le\frac{\kappa^2}{16} .
\end{equation}

\paragraph{Step 4: uniformity.}
Let $\kappa\in(0,1]$, and let $\delta\le1$ be so small that $K_0\delta^{1/4}\le\kappa^2/16$. Let $v_1,v_2,v_3$ be an admissible triple with $\max_i\delta(v_i)\le\delta$. If $c_{13}<\kappa$ or $c_{23}<\kappa$, then $c_{13}c_{23}-c_{12}\le c_{13}c_{23}<\kappa$, because all overlaps lie in $[0,1]$. Otherwise, \eqref{eq:geometry-triple-epsilon} and \eqref{eq:geometry-product-quantitative} imply that $c_{13}c_{23}-c_{12}\le14\kappa^{-1}K_0^{1/4}\delta^{1/16}$. This shows that
\[
 \omega(\delta)\le\max\bigl\{\kappa,\,14\kappa^{-1}K_0^{1/4}\delta^{1/16}\bigr\},
\]
and $\limsup_{\delta\downarrow0}\omega(\delta)\le\kappa$. Since $\kappa\in(0,1]$ is arbitrary, we conclude that $\omega(\delta)\to0$, as desired.
\end{proof}

\chapter{Three-time correlations and the susceptibility}\label{sec:three-time}

This section applies the deterministic results of \cref{sec:geometry-results} to the cubic CK solution. In \cref{sec:geometry-center-correlations}, we establish that the normalized prefix centers have asymptotically the same overlaps as the correlations divided by the plateau (\cref{prop:geometry-center-gram}). Combined with \cref{thm:geometry-threshold-product}, this gives the three-time product inequality of \cref{thm:geometry-actual-product} and composition bounds for last-crossing times (\cref{cor:geometry-last-crossing}). In \cref{sec:geometry-susceptibility}, we prove the lower bound $\liminf_{t\to\infty}\chi(t,0)\ge2(1-q)$ (\cref{thm:geometry-susceptibility-lower}). These results prove \cref{thm:main-geometry}. The product inequality is used again in \cref{resp:fisher-chapter}, where it bounds the memory term of the response equation, and the last-crossing bounds are discussed in \cref{sec:discussion}.

\paragraph{Standing assumptions and notation.}
Throughout this section, $p=3$ and $\beta\ge\bd=\sqrt{4/3}$, including the critical endpoint. We write $q$ for the selected plateau and $(c,r)$ for the stationary profile of \cref{sel:selected-sector}, with $r=-c'$. Recall that $q\ge1/2$, and that the bounded-lag translates of $C$ and $R$ converge locally uniformly, as in \eqref{sel:stationary-window}. We use the tree vectors $v_t$ and the prefix projections $P_T$ of \cref{prop:geometry-tree-realization}, and the notation $f$, $F_*$, and $\delta(\cdot)$ of \cref{sec:geometry-results}.

It is convenient to index a prefix center by its terminal time and its age. For $s\ge D\ge0$, we define
\[
 M_{s,D}=P_{s-D}v_s,\qquad U_{s,D}=\frac{M_{s,D}}{\norm{M_{s,D}}},
\]
the latter whenever $M_{s,D}\ne0$. In the notation of \cref{sec:selection}, $M_{s,D}=M_{s-D}(D)$: it is the part of the state at time $s$ that is determined by the noise up to time $s-D$. We call $D$ the smoothing age of the center.

Fix $D\ge0$. Recall from \cref{sel:selected-sector} that the translated kernels converge along every sequence of times tending to infinity. Hence \cref{sel:prefix-reflection} applies along every such sequence, and its conclusions hold as $T\to\infty$. In particular, \eqref{sel:prefix-norm} gives $\norm{M_{s,D}}^2\to c(2D)\ge q$ as $s\to\infty$, and $U_{s,D}$ is defined for all large $s$. The plateau is positive and marginal (\cref{sel:selected-sector,sel:positive-marginality}), which means that \cref{sel:threshold-centers} applies. Let $h_D$ and $\delta_D$ be as in \eqref{sel:prefix-energy} and \eqref{sel:center-deficit}. The cubic tree energy \eqref{eq:geometry-tree-energy} is $h=\sqrt3f$, the function $f$ is homogeneous of degree three, and $\Eth=\sqrt3F_*$. By \eqref{sel:prefix-norm} and \eqref{sel:prefix-energy}, we have
\[
 f(U_{s,D})=\frac{f(M_{s,D})}{\norm{M_{s,D}}^3}=\frac{h(M_{s,D})}{\sqrt3\norm{M_{s,D}}^3}\longrightarrow\frac{h_D}{\sqrt3\,c(2D)^{3/2}}=\frac{\Eth-\delta_D}{\sqrt3}=F_*-\frac{\delta_D}{\sqrt3}
\]
as $s\to\infty$. Together with \eqref{sel:center-deficit}, this gives
\begin{equation}
 \lim_{s\to\infty}\bigl[F_*-f(U_{s,D})\bigr]=\frac{\delta_D}{\sqrt3},
 \qquad\lim_{D\to\infty}\delta_D=0 .
 \label{eq:geometry-selected-center-deficit}
\end{equation}
In all limits in this section, we let the waiting time tend to infinity with the smoothing age fixed, and increase the smoothing age only afterward.

\section{From causal centers to actual correlations}\label{sec:geometry-center-correlations}

In this section, we assume that $\beta\ge\bd$. We use the following inputs:
\begin{itemize}
\item from \cref{resp:chapter}, the remote bound \eqref{resp:remote-forward} on the later derivative (\cref{resp:forward-derivative}), and the bound \eqref{resp:earlier-smoothing} on the variation of a correlation row across a window (\cref{resp:remote-earlier}), which rests on the strong localization of the response (\cref{resp:strong-row}) and on the Fisher control of the backward defect (\cref{thm:fisher-clock});
\item from \cref{sec:tree-ceilings}, \cref{sel:threshold-centers}, through \eqref{eq:geometry-selected-center-deficit};
\item from \cref{sec:geometry-results}, \cref{thm:geometry-threshold-product};
\item from \cref{sec:selection}, the complete monotonicity of $c$ (\cref{sel:cluster-complete-monotonicity}).
\end{itemize}

At $\beta=\bd$, we have $q=1/2$ by \cref{sel:selected-sector}, and $6\beta^2q(1-q)^2=1$ by \cref{sel:positive-marginality}. Then the plateau is positive and marginal, so \cref{sel:threshold-centers} applies at the endpoint. The response and Fisher inputs above hold for every $\beta>0$. The argument below uses $q>0$ and these inputs, and therefore includes $\beta=\bd$ directly. It does not use a strictly positive Fisher margin; the positive relaxation cost of \cref{resp:fisher-chapter} retains the strict hypothesis $\beta>\bd$.

These inputs imply that late, remote correlations change little across a window of length $D$, uniformly after the waiting-time limit. For the later coordinate, let $s\ge S$ and $t\ge s+D$. Then $s\le t-D$ and $t\ge S+D$, so \eqref{resp:remote-forward} with $L=D$ gives
\begin{equation}
 \limsup_{S\to\infty}\sup_{t\ge s+D,\ s\ge S}|\partial_tC(t,s)|\le\sqrt{c''(2D)} .
 \label{eq:geometry-remote-derivative-input}
\end{equation}
The earlier coordinate is controlled by \eqref{resp:earlier-smoothing}. Both bounds become small after multiplication by the window length:
\begin{equation}
 D^2c''(2D)\longrightarrow0,\qquad D\int_D^\infty r(u)^2\,\dd u\longrightarrow0\qquad(D\to\infty).
 \label{eq:geometry-smoothing-moduli}
\end{equation}
To prove the first limit, we use the spectral representation \eqref{eq:stationary-spectral-measure}. Differentiating under the integral for $t>0$, we have
\[
 D^2c''(2D)=\int_{[0,\infty)}(Dx)^2e^{-2Dx}\,\nu(\dd x).
\]
The integrand is bounded by $\sup_{y\ge0}y^2e^{-2y}=e^{-2}$ and tends to zero as $D\to\infty$ for every $x\ge0$, including $x=0$. Since $\nu$ is a probability measure, the first limit follows from dominated convergence. For the second, recall that $r$ is nonincreasing and $\int_0^\infty r=1-q$. This implies that $Dr(D)\le\int_0^Dr\le1-q$, and
\[
 D\int_D^\infty r(u)^2\,\dd u\le Dr(D)\int_D^\infty r(u)\,\dd u\le(1-q)\int_D^\infty r(u)\,\dd u\longrightarrow0 .
\]
We use only the complete monotonicity of $c$ here, and not the power law of \cref{sec:stationary-relaxation}.

We place the causal centers and the three-time inequality below in the line of the normalized time-averaged magnetizations and the spherical product relation discussed by Cugliandolo and Kurchan \cite[arXiv version~1, Section~5.1, Equations~(5.1.8)--(5.1.9)]{ck1994}. In that work, the product relation is an asymptotic relation of the physical theory. Here we prove a uniform one-sided inequality for the CK solution, without assuming a slow clock.

The next lemma compares a center with the average of the trajectory over the window that precedes it. For $s\ge D>0$, we define the Bochner integral
\[
 Z_{s,D}=\frac1D\int_{s-D}^sv_u\,\dd u .
\]
It is well defined because $u\mapsto v_u$ is continuous, and $\norm{Z_{s,D}}\le1$.

\begin{lemma}[Causal centers are approximated by local averages]
\label{lem:geometry-center-average}
For every fixed $D>0$,
\begin{equation}
 \lim_{s\to\infty}\norm{M_{s,D}-Z_{s,D}}^2=\eta_D,
 \qquad
 \eta_D=c(2D)+\frac2{D^2}\int_0^D(D-\tau)c(\tau)\,\dd\tau-\frac2D\int_D^{2D}c(\tau)\,\dd\tau .
 \label{eq:geometry-center-average}
\end{equation}
Moreover, $\eta_D\ge0$ and $\eta_D\to0$ as $D\to\infty$.
\end{lemma}

\begin{proof}
We expand $\norm{M_{s,D}-Z_{s,D}}^2=\norm{M_{s,D}}^2+\norm{Z_{s,D}}^2-2\ip{M_{s,D}}{Z_{s,D}}$ and compute the limit of each term. By \eqref{sel:prefix-norm}, the first term tends to $c(2D)$. Next,
\begin{align*}
 \norm{Z_{s,D}}^2&=\frac1{D^2}\int_{s-D}^s\!\!\int_{s-D}^sC(u,w)\,\dd u\,\dd w\\
 &\longrightarrow\frac1{D^2}\int_0^D\!\!\int_0^Dc(|u-w|)\,\dd u\,\dd w=\frac2{D^2}\int_0^D(D-\tau)c(\tau)\,\dd\tau,
\end{align*}
where we used the uniform convergence \eqref{sel:stationary-window} on the window $[s-D,s]$. For the cross term, let $T=s-D$, and write $Q_T(a,b)=\ip{M_T(a)}{M_T(b)}$ for the Gram kernel of the prefix centers of \cref{sec:selection}. For $u=T+a$ with $a\in[0,D]$, we have $\ip{M_{s,D}}{v_u}=\ip{P_Tv_{T+D}}{P_Tv_{T+a}}=Q_T(D,a)$. As noted at the start of this section, \eqref{sel:prefix-gram} holds as $T\to\infty$. It implies that $Q_T(D,a)$ tends to $c(D+a)$ uniformly in $a\in[0,D]$, and
\[
 \ip{M_{s,D}}{Z_{s,D}}=\frac1D\int_0^DQ_T(D,a)\,\dd a\longrightarrow\frac1D\int_0^Dc(D+a)\,\dd a=\frac1D\int_D^{2D}c(\tau)\,\dd\tau .
\]
This proves \eqref{eq:geometry-center-average}. The limit $\eta_D$ is nonnegative, as a limit of squared norms.

To show that $\eta_D\to0$, let $g=c-q$, which is nonnegative, nonincreasing, and tends to zero. The constant $q$ contributes $q+q-2q=0$ to $\eta_D$, so
\[
 \eta_D=g(2D)+\frac2{D^2}\int_0^D(D-\tau)g(\tau)\,\dd\tau-\frac2D\int_D^{2D}g(\tau)\,\dd\tau .
\]
The first term tends to zero. The second is at most $2D^{-1}\int_0^Dg$, which tends to zero because $g(\tau)\to0$. The absolute value of the third is at most $2g(D)$, which also tends to zero.
\end{proof}

\begin{proposition}[Uniform aging Gram of normalized centers]
\label{prop:geometry-center-gram}
We have
\begin{equation}
 \lim_{D\to\infty}\limsup_{S\to\infty}
 \sup_{s\ge S+D,\ t\ge s+2D}
 \left|\ip{U_{t,D}}{U_{s,D}}-\frac{C(t,s)}q\right|=0 .
 \label{eq:geometry-center-gram}
\end{equation}
In particular, the correlation of two late times with a diverging separation does not exceed the plateau in the limit: if $s_n\to\infty$ and $t_n-s_n\to\infty$, then
\begin{equation}\label{eq:geometry-remote-envelope}
 \limsup_{n\to\infty}C(t_n,s_n)\le q .
\end{equation}
\end{proposition}

In \eqref{eq:geometry-center-gram}, the supremum is over pairs for which both centers are defined. For fixed $D$, this includes all pairs in the range once $S$ is large.

\begin{proof}
The proof has four steps. We first replace each center by the local average of \cref{lem:geometry-center-average}. We then compare the overlap of two local averages with the correlation of their endpoints. Combining the two steps gives the claim for the unnormalized centers, and in the last step we divide by their norms.

\paragraph{Step 1: centers and local averages.}
Since $\norm{M_{s,D}}\le\norm{v_s}=1$ and $\norm{Z_{t,D}}\le1$, we have
\[
 \bigl|\ip{M_{t,D}}{M_{s,D}}-\ip{Z_{t,D}}{Z_{s,D}}\bigr|
 \le\norm{M_{t,D}-Z_{t,D}}+\norm{M_{s,D}-Z_{s,D}} .
\]
The limit in \eqref{eq:geometry-center-average} is a limit as the single variable $s$ tends to infinity, so that $\limsup_{S\to\infty}\sup_{s\ge S}\norm{M_{s,D}-Z_{s,D}}=\eta_D^{1/2}$.

\paragraph{Step 2: local averages and correlations.}
Let $s\ge S+D$ and $t\ge s+2D$. Then
\[
 \ip{Z_{t,D}}{Z_{s,D}}=\frac1{D^2}\int_{t-D}^t\!\!\int_{s-D}^sC(u,w)\,\dd w\,\dd u .
\]
For $u\in[t-D,t]$ and $w\in[s-D,s]$, we compare $C(u,w)$ with $C(t,s)$ by changing the two coordinates in succession:
\[
 |C(u,w)-C(t,s)|\le\int_u^t|\partial_xC(x,w)|\,\dd x+\int_{s-D}^s|\partial_yC(t,y)|\,\dd y .
\]
In the first integral, $x\ge u\ge t-D\ge s+D\ge w+D$ and $w\ge s-D\ge S$, and the integrand is bounded by the supremum in \eqref{eq:geometry-remote-derivative-input}. The second integral is the quantity bounded in \eqref{resp:earlier-smoothing}, since $s\ge S+D$ and $t\ge s+D$. Averaging over $u$ and $w$ and using these two bounds, we find
\[
 \limsup_{S\to\infty}\sup_{s\ge S+D,\ t\ge s+2D}\bigl|\ip{Z_{t,D}}{Z_{s,D}}-C(t,s)\bigr|
 \le D\sqrt{c''(2D)}+\Bigl(D\int_D^\infty r(u)^2\,\dd u\Bigr)^{1/2}.
\]

\paragraph{Step 3: unnormalized centers.}
Combining Steps 1 and 2, we obtain
\begin{equation}\label{eq:geometry-unnormalized-gram}
 \limsup_{S\to\infty}\sup_{s\ge S+D,\ t\ge s+2D}\bigl|\ip{M_{t,D}}{M_{s,D}}-C(t,s)\bigr|\le\mathcal E_D,
\end{equation}
where
\[
 \mathcal E_D=2\eta_D^{1/2}+D\sqrt{c''(2D)}+\Bigl(D\int_D^\infty r(u)^2\,\dd u\Bigr)^{1/2}.
\]
\Cref{lem:geometry-center-average} and \eqref{eq:geometry-smoothing-moduli} imply that $\mathcal E_D\to0$ as $D\to\infty$.

\paragraph{Step 4: normalization.}
For centers that are defined,
\[
 \Bigl|\ip{U_{t,D}}{U_{s,D}}-\frac{C(t,s)}q\Bigr|
 \le\frac{|\ip{M_{t,D}}{M_{s,D}}-C(t,s)|}{\norm{M_{t,D}}\norm{M_{s,D}}}
 +C(t,s)\Bigl|\frac1{\norm{M_{t,D}}\norm{M_{s,D}}}-\frac1q\Bigr| .
\]
As $s\to\infty$, the squared norm $\norm{M_{s,D}}^2$ tends to $c(2D)\ge q$, and this limit is again uniform over $s\ge S$ as $S\to\infty$. Using \eqref{eq:geometry-unnormalized-gram} and $C\le1$, we find that the upper limit as $S\to\infty$ of the supremum in \eqref{eq:geometry-center-gram} is at most
\[
 \frac{\mathcal E_D}{c(2D)}+\Bigl|\frac1{c(2D)}-\frac1q\Bigr|\le\frac{\mathcal E_D}q+\frac{c(2D)-q}{q^2} .
\]
This tends to zero as $D\to\infty$, because $\mathcal E_D\to0$ and $c(2D)\to q$. This proves \eqref{eq:geometry-center-gram}.

For \eqref{eq:geometry-remote-envelope}, let $e(S,D)$ denote the supremum in \eqref{eq:geometry-center-gram}. By the Cauchy--Schwarz inequality, $\ip{U_{t,D}}{U_{s,D}}\le1$, and we have $C(t,s)\le q+q\,e(S,D)$ for all pairs in its range. For fixed $S$ and $D$, the pair $(s_n,t_n)$ lies in this range for all large $n$, and therefore $\limsup_nC(t_n,s_n)\le q+q\,e(S,D)$. Letting $S\to\infty$ and then $D\to\infty$ proves \eqref{eq:geometry-remote-envelope}.
\end{proof}

We now transfer the product inequality of \cref{thm:geometry-threshold-product} to the correlations. For $S\ge0$ and $D>0$, let
\[
 \mathcal A_{S,D}=\{(u,s,t):u\ge S+D,\ s\ge u+2D,\ t\ge s+2D\}.
\]

\begin{theorem}[Actual three-time product inequality]
\label{thm:geometry-actual-product}
For the cubic CK solution at $\beta\ge\bd$, including the critical endpoint,
\begin{equation}
 \lim_{D\to\infty}\limsup_{S\to\infty}
 \sup_{(u,s,t)\in\mathcal A_{S,D}}
 \bigl[C(t,s)C(s,u)-qC(t,u)\bigr]_+=0 .
 \label{eq:geometry-actual-product}
\end{equation}
The same holds with the bracket replaced by $[C(t,s)C(t,u)-qC(s,u)]_+$ or by $[C(t,u)C(s,u)-qC(t,s)]_+$. In particular, for every sequence of triples $u_n<s_n<t_n$ with $u_n\to\infty$, $s_n-u_n\to\infty$, and $t_n-s_n\to\infty$,
\begin{equation}\label{eq:geometry-actual-product-sequences}
 \liminf_{n\to\infty}\bigl\{qC(t_n,u_n)-C(t_n,s_n)C(s_n,u_n)\bigr\}\ge0,
\end{equation}
and the same holds for the other two inequalities.
\end{theorem}

\begin{proof}
We first check that the centers satisfy the positivity hypothesis of \cref{thm:geometry-threshold-product}. Recall from \cref{prop:geometry-tree-realization} that every component of $v_t$ is a nonnegative function of the leaf labels. The projections $P_T$ and $\Pi_h$ preserve this nonnegativity, since they multiply every component by an indicator function (\cref{sec:tree-space}). The inner product of two nonnegative vectors is nonnegative, and therefore $\ip{\Pi_hM_{t,D}}{\Pi_hM_{s,D}}\ge0$ for all $h,s,t$, and $D$, and the same holds for the normalized centers.

Let $\omega$ be the function of \cref{thm:geometry-threshold-product}. Fix $D$ so large that $\delta_D/\sqrt3+1/D\le1$, and write $\delta_D'=\delta_D/\sqrt3+1/D$. Using \eqref{eq:geometry-selected-center-deficit}, we choose $S_D$ such that $U_{x,D}$ is defined and $F_*-f(U_{x,D})\le\delta_D'$ for all $x\ge S_D$. Let $S\ge S_D$ and $(u,s,t)\in\mathcal A_{S,D}$. All three times are at least $S$, so the triple
\[
 v_1=U_{t,D},\qquad v_2=U_{u,D},\qquad v_3=U_{s,D}
\]
satisfies the hypotheses of \cref{thm:geometry-threshold-product} with maximal deficit at most $\delta_D'$. Writing $g_{xy}=\ip{U_{x,D}}{U_{y,D}}$, we deduce from \eqref{eq:geometry-threshold-product} that
\[
 g_{tu}\ge g_{ts}g_{su}-\omega(\delta_D') .
\]

We now replace the overlaps by correlations. Let $e(S,D)$ be the supremum in \eqref{eq:geometry-center-gram}. The three pairs $(s,t)$, $(u,s)$, and $(u,t)$ all lie in the range of that supremum, since $s\ge u+2D\ge S+D$, $t\ge s+2D$, and $t\ge u+2D$. We obtain $|C(x,y)-q\,g_{xy}|\le q\,e$ for each pair, with $e=e(S,D)$. All correlations and normalized overlaps lie in $[0,1]$. Using $C(t,s)\le q(g_{ts}+e)$, $C(s,u)\le q(g_{su}+e)$, and $C(t,u)\ge q(g_{tu}-e)$, we find
\[
 C(t,s)C(s,u)-qC(t,u)
 \le q^2(g_{ts}g_{su}-g_{tu})+q^2(3e+e^2)
 \le q^2\omega(\delta_D')+q^2(3e+e^2) .
\]
We take the supremum over $\mathcal A_{S,D}$ and let $S\to\infty$. \Cref{prop:geometry-center-gram} implies that $\bar e_D=\limsup_{S\to\infty}e(S,D)$ tends to zero as $D\to\infty$, and $\omega(\delta_D')\to0$ by \cref{thm:geometry-threshold-product}, since $\delta_D'\to0$. We conclude that
\[
 \limsup_{S\to\infty}\sup_{\mathcal A_{S,D}}\bigl[C(t,s)C(s,u)-qC(t,u)\bigr]_+\le q^2\omega(\delta_D')+q^2(3\bar e_D+\bar e_D^2)\longrightarrow0
\]
as $D\to\infty$, which is \eqref{eq:geometry-actual-product}. The other two inequalities follow in the same way, by choosing $U_{t,D}$ or $U_{u,D}$ as the vector $v_3$ in \cref{thm:geometry-threshold-product}.

Finally, let $(u_n,s_n,t_n)$ be a sequence as in the statement. For fixed $S$ and $D$, the triple $(u_n,s_n,t_n)$ lies in $\mathcal A_{S,D}$ for all large $n$. Then the upper limit of $[C(t_n,s_n)C(s_n,u_n)-qC(t_n,u_n)]_+$ as $n\to\infty$ is at most the supremum over $\mathcal A_{S,D}$, for every $S$ and $D$. Letting $S\to\infty$ and then $D\to\infty$, we find that this upper limit is zero, which implies \eqref{eq:geometry-actual-product-sequences}. The other two inequalities are handled in the same way.
\end{proof}

The product inequality implies composition bounds for relaxation times. For a level $c\in(0,q)$, we consider the last time at which the correlation with a given earlier time is at least $c$. Here and in \cref{resp:fisher-chapter}, the letter $c$ without an argument, possibly with an index, denotes such a level; the stationary profile is always written with its argument, as in $c(W)$.

\begin{corollary}[Composition of last-crossing relaxation times]
\label{cor:geometry-last-crossing}
For a level $c\in(0,q)$ and $s\ge0$, define
\[
 \Theta_c(s)=\sup\{t\ge s:C(t,s)\ge c\}.
\]
For levels $c_1,c_2\in(0,q)$ and $0<c_0<c_1c_2/q$, we have, for every sufficiently large $s$,
\begin{equation}
 \Theta_{c_0}(s)\ge\Theta_{c_2}(\Theta_{c_1}(s)).
 \label{eq:geometry-last-crossing}
\end{equation}
More generally, let $c_1,\dots,c_k\in(0,q)$ be a fixed finite list of levels, and let $s_0=s$ and $s_j=\Theta_{c_j}(s_{j-1})$ for $1\le j\le k$. Then
\[
 \liminf_{s\to\infty}C(s_k,s)\ge q^{1-k}\prod_{j=1}^kc_j .
\]
\end{corollary}

\begin{proof}
We first show that $\Theta_c(s)$ is finite and that $C(\Theta_c(s),s)=c$. By \cref{sel:fixed-wait-forgetting}, $C(t,s)\to0$ as $t\to\infty$, and the set $\{t\ge s:C(t,s)\ge c\}$ is bounded. It contains $s$, because $C(s,s)=1$, and it is closed, because $C$ is continuous. Its supremum $\Theta_c(s)$ therefore belongs to it, that is, $C(\Theta_c(s),s)\ge c$. If this inequality were strict, then by continuity some $t>\Theta_c(s)$ would also satisfy $C(t,s)>c$, which is impossible.

Next, we show that $\Theta_c(s)-s\to\infty$ as $s\to\infty$. Fix $W>0$. By the stationary window \eqref{sel:stationary-window}, we have $C(s+W,s)\to c(W)$, and the profile satisfies $c(W)\ge q>c$. Then $C(s+W,s)>c$ for all large $s$, which gives $\Theta_c(s)\ge s+W$.

Let $s_1=\Theta_{c_1}(s)$ and $s_2=\Theta_{c_2}(s_1)$. As $s\to\infty$, we have $s\to\infty$, $s_1-s\to\infty$, and, since $s_1\ge s$, also $s_2-s_1\to\infty$. Applying \eqref{eq:geometry-actual-product-sequences} to the triple $s<s_1<s_2$ along any sequence of values of $s$ tending to infinity, and using $C(s_2,s_1)=c_2$ and $C(s_1,s)=c_1$, we obtain
\[
 \liminf_{s\to\infty}C(s_2,s)\ge\frac{c_1c_2}q>c_0 .
\]
Then $C(s_2,s)\ge c_0$ for all large $s$, and $s_2$ belongs to the set that defines $\Theta_{c_0}(s)$. This proves \eqref{eq:geometry-last-crossing}.

For the finite list, we prove by induction on $j$ that $\liminf_{s\to\infty}C(s_j,s)\ge q^{1-j}\prod_{i\le j}c_i$. For $j=1$, we have $C(s_1,s)=c_1$. Let $j\ge2$. The times $s_j$ are nondecreasing in $j$, which gives $s_{j-1}-s\ge s_1-s\to\infty$, and $s_j-s_{j-1}\to\infty$ because $s_{j-1}\to\infty$. Applying \eqref{eq:geometry-actual-product-sequences} to the triple $s<s_{j-1}<s_j$ and using $C(s_j,s_{j-1})=c_j$, we find
\[
 q\liminf_{s\to\infty}C(s_j,s)\ge c_j\liminf_{s\to\infty}C(s_{j-1},s)\ge c_jq^{2-j}\prod_{i<j}c_i ,
\]
where the last inequality is the induction hypothesis. This completes the induction.
\end{proof}

\begin{remark}\label{rem:geometry-last-crossing-scope}
The definition of $\Theta_c(s)$ allows returns: the correlation $C(t,s)$ may fall below $c$ and rise above it again before the last crossing. \Cref{thm:geometry-actual-product,cor:geometry-last-crossing} are inequalities. They do not assert a multiplicative composition law for correlations, and they do not identify an aging clock; see \cref{sec:discussion}.
\end{remark}

\section{The susceptibility lower bound}
\label{sec:geometry-susceptibility}

Throughout this section, $\beta\ge\bd$; the endpoint $\beta=\bd$, at which $q=1/2$, is included. We have $\int_0^\infty r=1-q$, because $r=-c'$, $c(0)=1$, and $c(\tau)\to q$. For $p=3$, the marginality equation \eqref{sel:positive-plateau-equation} reads $6\beta^2q(1-q)^2=1$, that is,
\begin{equation}\label{eq:geometry-marginal-normalization}
 \beta\sqrt3\sqrt q\,(1-q)=\frac1{\sqrt2} .
\end{equation}

In this section, we use \cref{thm:geometry-universal-root} from \cref{sec:geometry-results}, \cref{sel:threshold-centers} from \cref{sec:tree-ceilings} (through \eqref{eq:geometry-selected-center-deficit}), and the strong localization of the response row, \cref{resp:strong-row}, from \cref{resp:chapter}. We also use the root frame \eqref{eq:geometry-prefix-root-frame} and \cref{sel:prefix-reflection}. For $p=3$, the root frame reads
\begin{equation}\label{eq:geometry-cubic-root-frame}
 \mathcal I^*P_Tv_t=\beta\sqrt3\int_0^tR(t,u)\,|P_Tv_u\rangle\langle P_Tv_u|\,\dd u .
\end{equation}
The root operator of a center is a positive integral of the rank-one operators $|P_Tv_u\rangle\langle P_Tv_u|$, weighted by the response. \Cref{thm:geometry-universal-root} determines the limit of the normalized root in Hilbert--Schmidt norm, but this alone does not determine the limit of its trace, which is where the total response mass appears. We obtain a lower bound on the trace by first replacing each old rank-one term $|v_u\rangle\langle v_u|$ by the smoothed term $|M_{u,D}\rangle\langle M_{u,D}|$, whose norm converges. Because of causality, the errors made on time strips at distance at least $D$ from each other are orthogonal. Their total is therefore controlled by the squared response row on the old times, which tends to zero by \cref{resp:strong-row}. We need no bound on the integral $\chi(t,0)$ of the response row; by \eqref{resp:square-tail}, it is only known to be $o(\sqrt t)$.

\begin{lemma}[Orthogonality of disjoint causal strips]
\label{lem:geometry-disjoint-strips}
For $D>0$ and $u\ge D$, let
\[
 F_{u,D}=|v_u\rangle\langle v_u|-|M_{u,D}\rangle\langle M_{u,D}| .
\]
If $u,w\ge D$ and $|u-w|\ge D$, then $\ip{F_{u,D}}{F_{w,D}}_{\HS}=0$. For $D\le T\le t$, we have
\begin{equation}
 \Bigl\|\int_D^TR(t,u)F_{u,D}\,\dd u\Bigr\|_{\HS}^2\le8D\int_D^TR(t,u)^2\,\dd u .
 \label{eq:geometry-disjoint-strips}
\end{equation}
If $t_n\to\infty$, $L_n\to\infty$, and $t_n-L_n\ge D$, then the Hilbert--Schmidt distance between
\[
 \int_0^{t_n-L_n}R(t_n,u)\,|v_u\rangle\langle v_u|\,\dd u
 \quad\text{and}\quad
 \int_D^{t_n-L_n}R(t_n,u)\,|M_{u,D}\rangle\langle M_{u,D}|\,\dd u
\]
tends to zero.
\end{lemma}

\begin{proof}
We first note that the integrals are Bochner integrals in trace norm. For each fixed $x\in\mathcal H$, the map $T\mapsto P_Tx$ is continuous: on each tree summand, $P_T$ multiplies the leaf kernel by indicator functions of leaf times at most $T$, and the continuity follows from dominated convergence, first on each summand and then on their direct sum. Together with the continuity of $u\mapsto v_u$, this shows that $u\mapsto M_{u,D}$ is continuous on $[D,\infty)$. Then the rank-one operators above are continuous in trace norm, with trace norms at most one.

Suppose that $u\ge w+D$; the opposite ordering follows by symmetry. Then $w\le u-D$, and we have $v_w\in\mathcal H_w\subset\mathcal H_{u-D}$ and $M_{w,D}\in\mathcal H_{w-D}\subset\mathcal H_{u-D}$. For both vectors $x\in\{v_w,M_{w,D}\}$, this gives $\ip{v_u}{x}=\ip{P_{u-D}v_u}{x}=\ip{M_{u,D}}{x}$. Expanding the four rank-one pairings and using $\langle|a\rangle\langle a|,|b\rangle\langle b|\rangle_{\HS}=\ip{a}{b}^2$, we find
\begin{align*}
 \ip{F_{u,D}}{F_{w,D}}_{\HS}
 &=\ip{v_u}{v_w}^2-\ip{v_u}{M_{w,D}}^2-\ip{M_{u,D}}{v_w}^2+\ip{M_{u,D}}{M_{w,D}}^2\\
 &=\ip{M_{u,D}}{v_w}^2-\ip{M_{u,D}}{M_{w,D}}^2-\ip{M_{u,D}}{v_w}^2+\ip{M_{u,D}}{M_{w,D}}^2=0 .
\end{align*}

For \eqref{eq:geometry-disjoint-strips}, note that $\norm{F_{u,D}}_{\HS}\le\norm{v_u}^2+\norm{M_{u,D}}^2\le2$, and that $R\ge0$. Let $g=\mathbf1_{[D,T]}R(t,\cdot)$, extended by zero to $\R$, and $k=\mathbf1_{[-D,D]}$. By the orthogonality just proved,
\begin{align*}
 \Bigl\|\int_D^TR(t,u)F_{u,D}\,\dd u\Bigr\|_{\HS}^2
 &=\int_D^T\!\!\int_D^TR(t,u)R(t,w)\ip{F_{u,D}}{F_{w,D}}_{\HS}\,\dd u\,\dd w\\
 &\le4\int_\R g(u)(k*g)(u)\,\dd u .
\end{align*}
By the Cauchy--Schwarz inequality and Young's convolution inequality, the right side is at most $4\norm{g}_2\norm{k*g}_2\le4\norm{k}_1\norm{g}_2^2=8D\int_D^TR(t,u)^2\,\dd u$.

For the last assertion, let $T_n=t_n-L_n$, and let $f_t$ be the lag row \eqref{resp:lag-row}. Substituting $\tau=t_n-u$, we have
\begin{equation}\label{eq:geometry-old-row-l2}
 \int_0^{T_n}R(t_n,u)^2\,\dd u=\int_{L_n}^{t_n}f_{t_n}(\tau)^2\,\dd\tau
 \le2\norm{f_{t_n}-r}_{L^2(0,\infty)}^2+2\int_{L_n}^\infty r(\tau)^2\,\dd\tau\longrightarrow0,
\end{equation}
by \eqref{resp:strong-row-limit} and $r\in L^2(0,\infty)$. The difference of the two operators in the statement is
\[
 \int_0^DR(t_n,u)\,|v_u\rangle\langle v_u|\,\dd u+\int_D^{T_n}R(t_n,u)F_{u,D}\,\dd u .
\]
The first term has Hilbert--Schmidt norm at most $\int_0^DR(t_n,u)\,\dd u\le\sqrt D\,(\int_0^{T_n}R(t_n,u)^2\,\dd u)^{1/2}$, by the Cauchy--Schwarz inequality. The second term is bounded by \eqref{eq:geometry-disjoint-strips} with $T=T_n$. Both bounds tend to zero by \eqref{eq:geometry-old-row-l2}.
\end{proof}

\begin{theorem}[Susceptibility lower bound]
\label{thm:geometry-susceptibility-lower}
For the uniform-start cubic CK solution at $\beta\ge\bd$,
\begin{equation}
 \liminf_{t\to\infty}\chi(t,0)\ge2(1-q) .
 \label{eq:geometry-susceptibility-lower}
\end{equation}
More precisely, for every fixed $D>0$,
\begin{equation}
 \liminf_{t\to\infty}\chi(t,0)\ge(1-q)+\frac{q(1-q)}{c(2D)} .
 \label{eq:geometry-susceptibility-smoothed}
\end{equation}
\end{theorem}

\begin{proof}
Let $\gamma=\liminf_{t\to\infty}\chi(t,0)\in[0,\infty]$. If $\gamma=\infty$, there is nothing to prove. Otherwise, we choose $t_n\to\infty$ with $\chi(t_n,0)\to\gamma$. All later bounds on the total response mass are along this sequence. Because $c(2D)\downarrow q$ as $D\to\infty$, the bound \eqref{eq:geometry-susceptibility-smoothed} implies \eqref{eq:geometry-susceptibility-lower}. We therefore fix $D>0$ and prove \eqref{eq:geometry-susceptibility-smoothed}.

The proof has six steps.
\begin{enumerate}
\item We choose ages $L_n\to\infty$ such that the centers $U_n=U_{t_n,L_n}$ approach the ceiling and have the limits \eqref{eq:geometry-susceptibility-diagonal} below.
\item By \eqref{eq:geometry-cubic-root-frame}, the root operator $\mathcal I^*U_n$ splits into an old part $A_n$, which is an integral over the times before $T_n=t_n-L_n$, and a recent part $B_n$.
\item In the coordinates of \cref{thm:geometry-universal-root}, the recent part converges to the rank-one operator $2^{-1/2}|\mathbf1\rangle\langle\mathbf1|$.
\item It follows that the old part converges in Hilbert--Schmidt norm to a positive operator $\mathcal S$ with trace $2^{-1/2}$.
\item We smooth the old part with \cref{lem:geometry-disjoint-strips}. The smoothed operator is positive, and the lower semicontinuity of the trace shows that its trace is at least $2^{-1/2}$ in the limit.
\item The trace of the smoothed operator is explicit in the limit, and it is proportional to the old response mass $\gamma-(1-q)$. Comparing with Step 5 gives \eqref{eq:geometry-susceptibility-smoothed}.
\end{enumerate}

\paragraph{Step 1: the diagonal choice.}
Recall the notation $M_T(a)=P_Tv_{T+a}$ and $Q_T(a,b)=\ip{M_T(a)}{M_T(b)}$ from the proof of \cref{lem:geometry-center-average}. Fix an integer $k\ge1$. By \eqref{sel:prefix-gram}, by \eqref{eq:geometry-selected-center-deficit} with $D=k$ and $s=T+k$, and by \eqref{sel:stationary-window}, the following three statements hold for all large $T$:
\begin{enumerate}[label=(\roman*)]
\item $\sup_{0\le a,b\le k}|Q_T(a,b)-c(a+b)|\le1/k$;
\item $M_T(k)\ne0$, and $|f(M_T(k)/\norm{M_T(k)})-F_*+\delta_k/\sqrt3|\le1/k$;
\item $|\int_0^kR(T+k,T+k-\tau)\,\dd\tau-(1-c(k))|\le1/k$.
\end{enumerate}
For (iii), we used that $\int_0^kR(T+k,T+k-\tau)\,\dd\tau\to\int_0^kr=1-c(k)$. Let $\theta_k$ be a threshold such that (i)--(iii) hold for all $T\ge\theta_k$, chosen nondecreasing in $k$ and satisfying $\theta_k\ge k$. Discarding finitely many $n$, we may assume that $t_n\ge1+\theta_1$. Let $L_n$ be the largest integer $k\ge1$ with $t_n\ge k+\theta_k$. For each fixed $K$, we have $t_n\ge K+\theta_K$ for all large $n$, because $t_n\to\infty$; therefore $L_n\to\infty$. Let $T_n=t_n-L_n$. Then $T_n\ge\theta_{L_n}\ge L_n\to\infty$, and (i)--(iii) hold with $k=L_n$ and $T=T_n$. We define
\[
 d_n=\norm{M_{t_n,L_n}}^2=Q_{T_n}(L_n,L_n),\qquad U_n=U_{t_n,L_n}=\frac{M_{T_n}(L_n)}{\sqrt{d_n}} .
\]
Note that $\chi(t_n,T_n)=\int_{T_n}^{t_n}R(t_n,u)\,\dd u=\int_0^{L_n}R(t_n,t_n-\tau)\,\dd\tau$. Since $c(L_n)\to q$, $c(2L_n)\to q$, and $\delta_{L_n}\to0$, the statements (i)--(iii) give
\begin{equation}
 d_n\to q,\qquad
 \chi(t_n,T_n)\to1-q,\qquad
 f(U_n)\to F_* .
 \label{eq:geometry-susceptibility-diagonal}
\end{equation}
We also record a consequence of (i) for the recent states. Fix $W>0$, and let $n$ be so large that $L_n\ge W$. For $0\le\tau\le W$, we have $P_{T_n}v_{t_n-\tau}=M_{T_n}(L_n-\tau)$, and (i) gives
\begin{align*}
 \bigl\|P_{T_n}v_{t_n-\tau}-\sqrt q\,U_n\bigr\|^2
 &=Q_{T_n}(L_n-\tau,L_n-\tau)+q-\frac{2\sqrt q\,Q_{T_n}(L_n-\tau,L_n)}{\sqrt{d_n}}\\
 &=c(2L_n-2\tau)+q-\frac{2\sqrt q\,c(2L_n-\tau)}{\sqrt{d_n}}+o(1),
\end{align*}
uniformly in $\tau\in[0,W]$. The function $c$ is nonincreasing with limit $q$, which gives $c(2L_n-2\tau)\to q$ and $c(2L_n-\tau)\to q$ uniformly in $\tau\in[0,W]$. Together with $d_n\to q$, this shows that
\begin{equation}\label{eq:geometry-recent-centers}
 \sup_{0\le\tau\le W}\bigl\|P_{T_n}v_{t_n-\tau}-\sqrt q\,U_n\bigr\|\longrightarrow0 .
\end{equation}

\paragraph{Step 2: the old and recent parts of the root.}
For $u\le T_n$, we have $v_u\in\mathcal H_u\subset\mathcal H_{T_n}$ and $P_{T_n}v_u=v_u$. Applying \eqref{eq:geometry-cubic-root-frame} with $T=T_n$ and $t=t_n$, and dividing by $\sqrt{d_n}$, we obtain $\mathcal I^*U_n=A_n+B_n$, where
\begin{align*}
 A_n&=\frac{\beta\sqrt3}{\sqrt{d_n}}\int_0^{T_n}R(t_n,u)\,|v_u\rangle\langle v_u|\,\dd u,\\
 B_n&=\frac{\beta\sqrt3}{\sqrt{d_n}}\int_{T_n}^{t_n}R(t_n,u)\,|P_{T_n}v_u\rangle\langle P_{T_n}v_u|\,\dd u .
\end{align*}
Both operators are positive and trace class, as Bochner integrals in trace norm of positive rank-one operators with trace norms at most one.

\paragraph{Step 3: the recent part.}
By \eqref{eq:geometry-susceptibility-diagonal}, \cref{thm:geometry-universal-root} applies to the unit vectors $U_n$. Let $V_n$ be the embeddings it provides, so that $V_nU_n=\mathbf1$ and $V_n\mathcal I^*U_nV_n^*\to\mathcal T\oplus0$ in Hilbert--Schmidt norm. We transport every operator by $V_n$ and write $\widehat X=V_nXV_n^*$. This preserves Hilbert--Schmidt distances, positivity, and traces. We write $|\mathbf1\rangle\langle\mathbf1|$ also for the operator $|\mathbf1\rangle\langle\mathbf1|\oplus0$. We claim that
\begin{equation}
 \widehat B_n\longrightarrow\frac1{\sqrt2}|\mathbf1\rangle\langle\mathbf1|\quad\text{in Hilbert--Schmidt norm.}
 \label{eq:geometry-recent-root-atom}
\end{equation}

To prove the claim, fix $W>0$ and let $n$ be so large that $L_n\ge W$. We split $B_n=B_n'+B_n''$ at the time $t_n-W$, where $B_n'$ is the integral over $[t_n-W,t_n]$. We write $w_{n,\tau}=P_{T_n}v_{t_n-\tau}$ and $\chi_n(W)=\int_0^WR(t_n,t_n-\tau)\,\dd\tau$. Using the rank-one estimate
\[
 \bigl\||x\rangle\langle x|-|y\rangle\langle y|\bigr\|_{\HS}\le(\norm{x}+\norm{y})\norm{x-y}
\]
with $x=w_{n,\tau}$ and $y=\sqrt q\,U_n$, both of norm at most one, we find
\[
 \Bigl\|B_n'-\frac{\beta\sqrt3\,q\,\chi_n(W)}{\sqrt{d_n}}|U_n\rangle\langle U_n|\Bigr\|_{\HS}
 \le\frac{2\beta\sqrt3}{\sqrt{d_n}}\,\chi_n(W)\sup_{0\le\tau\le W}\norm{w_{n,\tau}-\sqrt q\,U_n}\longrightarrow0,
\]
by \eqref{eq:geometry-recent-centers}. Note that $\chi_n(W)\to\int_0^Wr$, by \eqref{sel:stationary-window}. Since $V_n|U_n\rangle\langle U_n|V_n^*=|\mathbf1\rangle\langle\mathbf1|$ and $d_n\to q$, it follows that
\[
 \widehat B_n'\longrightarrow\beta\sqrt3\sqrt q\Bigl(\int_0^Wr(\tau)\,\dd\tau\Bigr)|\mathbf1\rangle\langle\mathbf1| .
\]
For the remaining part, since $\norm{P_{T_n}v_u}\le1$, we have
\begin{align*}
 \norm{\widehat B_n''}_{\HS}&\le\frac{\beta\sqrt3}{\sqrt{d_n}}\int_{T_n}^{t_n-W}R(t_n,u)\,\dd u
 =\frac{\beta\sqrt3}{\sqrt{d_n}}\bigl[\chi(t_n,T_n)-\chi_n(W)\bigr]\\
 &\longrightarrow\frac{\beta\sqrt3}{\sqrt q}\Bigl[(1-q)-\int_0^Wr(\tau)\,\dd\tau\Bigr],
\end{align*}
by \eqref{eq:geometry-susceptibility-diagonal}. Combining the last two displays, we obtain
\[
 \limsup_{n\to\infty}\Bigl\|\widehat B_n-\frac1{\sqrt2}|\mathbf1\rangle\langle\mathbf1|\Bigr\|_{\HS}
 \le\Bigl|\beta\sqrt3\sqrt q\int_0^Wr-\frac1{\sqrt2}\Bigr|+\frac{\beta\sqrt3}{\sqrt q}\Bigl[(1-q)-\int_0^Wr\Bigr].
\]
The left side does not depend on $W$. As $W\to\infty$, we have $\int_0^Wr\to1-q$, and the right side tends to zero by \eqref{eq:geometry-marginal-normalization}. This proves \eqref{eq:geometry-recent-root-atom}.

\paragraph{Step 4: the old part.}
By \cref{thm:geometry-universal-root} and \eqref{eq:geometry-recent-root-atom}, $\widehat A_n=V_n\mathcal I^*U_nV_n^*-\widehat B_n$ converges in Hilbert--Schmidt norm to $\mathcal S\oplus0$, where
\[
 \mathcal S=\mathcal T-\frac1{\sqrt2}|\mathbf1\rangle\langle\mathbf1|
 =\frac1{2\sqrt2}\int_0^1z^{-3/2}|\mathbf1_{[0,z]}\rangle\langle\mathbf1_{[0,z]}|\,\dd z .
\]
The second expression comes from the frame \eqref{eq:geometry-universal-root-frame}. It shows that $\mathcal S$ is positive and that $\tr\mathcal S=\tr\mathcal T-2^{-1/2}=2^{-1/2}$.

\paragraph{Step 5: smoothing and lower semicontinuity of the trace.}
Let $n$ be so large that $T_n\ge D$, and define the smoothed old part
\[
 A_{n,D}=\frac{\beta\sqrt3}{\sqrt{d_n}}\int_D^{T_n}R(t_n,u)\,|M_{u,D}\rangle\langle M_{u,D}|\,\dd u .
\]
It is a positive trace-class operator. By the last assertion of \cref{lem:geometry-disjoint-strips}, applied with the sequences $t_n$ and $L_n$, and since $d_n\to q$, we have $\norm{A_{n,D}-A_n}_{\HS}\to0$. By Step 4, $\widehat A_{n,D}\to\mathcal S\oplus0$ in Hilbert--Schmidt norm. We claim that
\begin{equation}
 \liminf_{n\to\infty}\tr A_{n,D}\ge\tr\mathcal S=\frac1{\sqrt2} .
 \label{eq:geometry-trace-fatou}
\end{equation}
This is the classical lower semicontinuity of the trace on positive operators, and not continuity, since the convergence holds only in Hilbert--Schmidt norm. We include the short proof. Let $Q$ be an orthogonal projection of finite rank on $L^2(0,1)\oplus\mathcal K$. For Hilbert--Schmidt operators $X$ and $Y$, we have $\tr(QXQ)=\tr(QX)=\ip{Q}{X}_{\HS}$, so
\[
 |\tr(QXQ)-\tr(QYQ)|\le\norm{Q}_{\HS}\norm{X-Y}_{\HS}=(\operatorname{rank}Q)^{1/2}\norm{X-Y}_{\HS} .
\]
In particular, $\tr(Q\widehat A_{n,D}Q)\to\tr(Q(\mathcal S\oplus0)Q)$. For a positive trace-class operator $X$, we have $\tr(QXQ)\le\tr X$, as one sees by computing the trace in an orthonormal basis that contains a basis of the range of $Q$. Since $\tr\widehat A_{n,D}=\tr A_{n,D}$, we obtain
\[
 \liminf_{n\to\infty}\tr A_{n,D}\ge\tr\bigl(Q(\mathcal S\oplus0)Q\bigr).
\]
Let $(e_i)$ be an orthonormal basis of $L^2(0,1)$, and let $Q_k$ be the projection onto the span of $e_1,\dots,e_k$. Then $\tr(Q_k(\mathcal S\oplus0)Q_k)=\sum_{i\le k}\ip{e_i}{\mathcal Se_i}$ increases to $\tr\mathcal S$ as $k\to\infty$. This proves \eqref{eq:geometry-trace-fatou}.

\paragraph{Step 6: the trace of the smoothed old part.}
The trace of a rank-one operator $|x\rangle\langle x|$ is $\norm{x}^2$, so
\[
 \tr A_{n,D}=\frac{\beta\sqrt3}{\sqrt{d_n}}\int_D^{T_n}R(t_n,u)\,z_D(u)\,\dd u,
 \qquad z_D(u)=\norm{M_{u,D}}^2 .
\]
Recall from \eqref{sel:prefix-norm} that $z_D(u)\to c(2D)$ as $u\to\infty$. By \eqref{eq:geometry-susceptibility-diagonal}, the old response mass satisfies
\[
 \int_0^{T_n}R(t_n,u)\,\dd u=\chi(t_n,0)-\chi(t_n,T_n)\longrightarrow\gamma-(1-q) .
\]
Fix $u_0\ge D$. For $n$ so large that $T_n\ge u_0$, the Cauchy--Schwarz inequality and \eqref{eq:geometry-old-row-l2} give $\int_0^{u_0}R(t_n,u)\,\dd u\le\sqrt{u_0}\,(\int_0^{T_n}R(t_n,u)^2\,\dd u)^{1/2}\to0$. Splitting the integrals at $u_0$ and using $0\le z_D\le1$ and $c(2D)\le1$, we find
\begin{align*}
 &\Bigl|\int_D^{T_n}R(t_n,u)z_D(u)\,\dd u-c(2D)\int_0^{T_n}R(t_n,u)\,\dd u\Bigr|\\
 &\qquad\le2\int_0^{u_0}R(t_n,u)\,\dd u+\sup_{u\ge u_0}|z_D(u)-c(2D)|\int_0^{T_n}R(t_n,u)\,\dd u .
\end{align*}
We recall that $\gamma<\infty$. The upper limit of the right side as $n\to\infty$ is at most $\sup_{u\ge u_0}|z_D(u)-c(2D)|\,(\gamma-(1-q))$, which tends to zero as $u_0\to\infty$. Together with $d_n\to q$, this gives
\[
 \tr A_{n,D}\longrightarrow\frac{\beta\sqrt3}{\sqrt q}\,c(2D)\bigl(\gamma-(1-q)\bigr).
\]
Comparing with \eqref{eq:geometry-trace-fatou} and using \eqref{eq:geometry-marginal-normalization}, we obtain
\[
 \frac{\beta\sqrt3}{\sqrt q}\,c(2D)\bigl(\gamma-(1-q)\bigr)\ge\frac1{\sqrt2}=\beta\sqrt3\sqrt q\,(1-q) .
\]
Multiplying by $\sqrt q/(\beta\sqrt3)$, we find $c(2D)(\gamma-(1-q))\ge q(1-q)$, which is \eqref{eq:geometry-susceptibility-smoothed}. This completes the proof.
\end{proof}

\begin{remark}\label{rem:geometry-susceptibility-prediction}
In the asymptotic solution of Cugliandolo and Kurchan \cite[arXiv version~1, paragraph after Equation~(12) and Equation~(13), p.~6]{ck1993}, the slow sector has the constant fluctuation--dissipation ratio $x=(p-2)(1-q)/q$. Under this prediction, the total susceptibility tends to $(1-q)+xq=(p-1)(1-q)$, which is $2(1-q)$ for $p=3$; formally, this is the relation \eqref{eq:discussion-slow-fdr} with $c=0$. \Cref{thm:geometry-susceptibility-lower} proves the corresponding lower bound. The matching upper bound \eqref{eq:discussion-susceptibility-upper} is open; see \cref{sec:discussion}.
\end{remark}

\chapter{A Fisher cost for relaxation below the plateau}\label{resp:fisher-chapter}

In this section, we prove \cref{resp:fisher-cost}, which is \cref{thm:main-fisher-cost}. It concerns the relaxation of the cubic CK solution below the plateau, measured in physical time. We show that a growing time window $[s,s+\tau]$ on which $\tau\int_s^{s+\tau}a_F$ tends to zero ends at the plateau, and that every first passage of the correlation below a level $c\in(0,q)$ costs a positive amount of this quantity.

In the proof, we compare the backward response equation of \cref{prop:backward-response}, integrated over the window, with a lower bound on its memory term. The Fisher control of the backward defect (\cref{thm:fisher-clock}), which is the form taken here by the entropy-production mechanism of Cugliandolo, Dean, and Kurchan \cite[arXiv version~2, Equations~(12)--(13), p.~3, and Equation~(18), p.~4]{cdk1997}, makes the integrated defect on the window small. The additional input is the uniform three-time inequality of \cref{thm:geometry-actual-product}: for late times $v<u<w$ with large gaps, it gives $C(w,u)C(u,v)\le qC(w,v)+o(1)$, and we use it to bound the memory term from below. Its product form has a physical predecessor in \cite[arXiv version~1, Section~5.1, Equation~(5.1.8)]{ck1994}. After the calibration of the selected sector, the two bounds differ by $3\beta^2(2q-1)(q-c)^2$, which is positive because $q>1/2$. We also use the remote envelope \eqref{eq:geometry-remote-envelope} of \cref{prop:geometry-center-gram}, the response localization of \cref{resp:uniform-row,resp:forward-derivative}, the stationary window \eqref{sel:stationary-window}, and the forgetting of fixed waiting times (\cref{sel:fixed-wait-forgetting}). These results prove \cref{thm:main-fisher-cost}, and their implications for relaxation rates are discussed in \cref{sec:discussion}.

\paragraph{Standing assumptions and notation.}
Throughout this section, $p=3$ and $\beta>\bd=\sqrt{4/3}$. By \cref{sel:selected-sector} and the discussion preceding it, the selected plateau satisfies $q>\qd=1/2$. We write $\ell=\ell_\beta=\lim_{t\to\infty}\lambda(t)$. For $p=3$, the marginality equation \eqref{sel:positive-plateau-equation} and the multiplier formula \eqref{sel:selected-multiplier} read
\begin{equation}
 6\beta^2q(1-q)^2=1,
 \qquad \ell=\frac1{1-q}+3\beta^2(1-q^2).
 \label{resp:fisher-calibration}
\end{equation}
Let $a_F$ be the Fisher majorant fixed in \cref{thm:fisher-clock}, from the same dimension extraction used in the physical-fork comparisons. Let $E(t,s)=\partial_sC(t,s)-R(t,s)$ be the backward FDT defect. Recall from \eqref{resp:fisher-domination} that $|E(t,s)|^2\le a_F(s)$ for $t>s>0$, and that $\int_0^\infty a_F<\infty$. Every row $C(t,\cdot)$ is absolutely continuous on $[0,t]$, with derivative $R(t,\cdot)+E(t,\cdot)$ almost everywhere (\cref{thm:fisher-clock}). We also use that $0\le R\le C\le1$ (\cref{prop:response-compactness}), that $\lambda\ge1$ by \eqref{eq:ck-multiplier}, which gives $\ell\ge1$, and that for $p=3$ the kernel \eqref{eq:self-energy-kernel} is $\Sigma(t,u)=6\beta^2C(t,u)R(t,u)$. As in \cref{sec:geometry-center-correlations}, the letter $c$ without an argument denotes a level in $(0,q)$, and the stationary profile is always written with its argument, as in $c(2D)$.

\begin{theorem}[Fisher-small windows and all-level passage costs]
\label{resp:fisher-cost}
Suppose that $s_n,\tau_n\ge0$, $s_n\to\infty$, $\tau_n\to\infty$, and
\begin{equation}
 \tau_n\int_{s_n}^{s_n+\tau_n}a_F(u)\,\dd u\longrightarrow0 .
 \label{resp:fisher-small-window}
\end{equation}
Then
\begin{equation}
 C(s_n+\tau_n,s_n)\to q,
 \qquad \chi(s_n+\tau_n,s_n)\to1-q .
 \label{resp:fisher-window-limit}
\end{equation}
For a level $c\in(0,q)$ and $s\ge0$, define the first passage lag
\[
 \tau_c(s)=\inf\{\tau\ge0:C(s+\tau,s)\le c\}.
\]
It is positive and finite, satisfies $C(s+\tau_c(s),s)=c$, tends to infinity as $s\to\infty$, and satisfies
\begin{align}
 \liminf_{s\to\infty}\tau_c(s)\int_s^{s+\tau_c(s)}a_F(u)\,\dd u
 &\ge\mathfrak c_\beta(c),\label{resp:local-fisher-cost}\\
 \liminf_{s\to\infty}\tau_c(s)\int_s^\infty a_F(u)\,\dd u
 &\ge\mathfrak c_\beta(c),\label{resp:tail-fisher-cost}
\end{align}
where the explicit positive constant is
\begin{equation}
 \mathfrak c_\beta(c)=\min\left\{1,
 \left[\frac{3\beta^2(2q-1)(q-c)^2}{\ell+15\beta^2}\right]^2\right\}.
 \label{resp:fisher-constant}
\end{equation}
\end{theorem}

\begin{remark}[The constant and the endpoint]\label{rem:fisher-constant}
By \eqref{resp:fisher-calibration}, we have $3\beta^2=(2q(1-q)^2)^{-1}$ and $\ell=(1+3q)/(2q(1-q))$, so that
\[
 \ell+15\beta^2=\frac{6+2q-3q^2}{2q(1-q)^2},
 \qquad
 \frac{3\beta^2(2q-1)(q-c)^2}{\ell+15\beta^2}=\frac{(2q-1)(q-c)^2}{6+2q-3q^2} .
\]
For $1/2<q<1$ and a level $c\in(0,q)$, we have $6+2q-3q^2=5+(1-q)(1+3q)>5$ and $0<(2q-1)(q-c)^2<1$. The ratio is therefore less than $1/5$, and the minimum in \eqref{resp:fisher-constant} is always attained by its second term. At the endpoint $\beta=\bd$, the plateau is $q=1/2$ and the formula for the constant vanishes. Although the cubic product inequality of \cref{thm:geometry-actual-product} includes this endpoint, the positive Fisher cost and the Fisher-small-window conclusion above are asserted only for $\beta>\bd$. The strict inequality is essential to the argument below, in which the positive margin $3\beta^2(2q-1)(q-c)^2$ is the source of the cost.
\end{remark}

The theorem follows from a quantitative statement about a single sequence of windows, which we prove first.

\begin{lemma}[A sequential physical-window obstruction]
\label{resp:fisher-obstruction}
Let $0\le s_n\le t_n$, $s_n\to\infty$, and $t_n-s_n\to\infty$, and suppose that $C(t_n,s_n)\to z\in(0,q)$. Suppose further that
\[
 \varepsilon_n=\left[(t_n-s_n)\int_{s_n}^{t_n}a_F(u)\,\dd u\right]^{1/2}\longrightarrow\varepsilon\le1 .
\]
Then
\begin{equation}
 3\beta^2(2q-1)(q-z)^2\le(\ell+9\beta^2)\varepsilon+6\beta^2\varepsilon^2 .
 \label{resp:sequential-obstruction}
\end{equation}
\end{lemma}

\begin{proof}
The proof has six steps.
\begin{enumerate}
\item The Fisher control bounds the integrated defect on every subinterval of the window by $\varepsilon_n$.
\item We integrate the backward response equation over the window. The result relates the response mass of the window to the memory integral $I_n$ defined in \eqref{eq:fisher-cost-memory-integral} below, with an error of order $\varepsilon_n$.
\item We split $I_n$ into an early strip, a bulk, and a recent strip of length $2D$, and compute the recent strip and the values of the correlation row at the ends of the bulk.
\item In the bulk, the product inequality bounds $1-C(u,s)^2$ from below by a function of the correlation row $C(t,u)$, which we integrate exactly.
\item We take the limits $n\to\infty$ and then $D\to\infty$, and obtain a lower bound for $I_n$.
\item The calibration \eqref{resp:fisher-calibration} turns the two estimates into \eqref{resp:sequential-obstruction}.
\end{enumerate}
We temporarily suppress the index $n$ and write $s=s_n$, $t=t_n$, $c_n=C(t,s)$, and $Y(u)=C(t,u)$ for $u\in[s,t]$. We assume throughout that $n$ is large; in particular, $s>0$.

\paragraph{Step 1: Fisher control on the window.}
For all $s\le v\le w\le t$, the bound $|E(w,u)|\le\sqrt{a_F(u)}$, the Cauchy--Schwarz inequality, and $a_F\ge0$ give
\begin{equation}
 \int_v^w|E(w,u)|\,\dd u\le\Bigl((w-v)\int_v^wa_F(u)\,\dd u\Bigr)^{1/2}\le\varepsilon_n .
 \label{resp:fisher-entire-triangle}
\end{equation}
The function $Y$ is absolutely continuous with derivative $R(t,\cdot)+E(t,\cdot)$, and $Y(t)=1$. Integrating over $[s,t]$, we find $1-c_n=\chi(t,s)+\int_s^tE(t,u)\,\dd u$. By \eqref{resp:fisher-entire-triangle} with $v=s$ and $w=t$, and since $c_n\ge0$, we have
\begin{equation}
 |\chi(t,s)-(1-c_n)|\le\varepsilon_n,
 \qquad \chi(t,s)\le1+\varepsilon_n .
 \label{resp:fisher-chi-bound}
\end{equation}
This bound on the response mass of the window comes from the Fisher control alone. We have no bound on $\chi(t,0)$, and none is needed.

\paragraph{Step 2: the integrated response equation.}
Recall from \cref{prop:backward-response} that the derivative $\partial_vR(t,v)$ is continuous on the closed triangle $s\le v\le t$, and note that $R(t,t)=1$. Integrating \eqref{eq:backward-response} in its second variable over $[s,t]$, and exchanging the order of integration on the finite triangle $\{s\le v\le u\le t\}$, where the integrand is continuous and bounded, we obtain
\begin{equation}
 1-R(t,s)=\int_s^t\lambda(u)R(t,u)\,\dd u
 -6\beta^2\int_s^tR(t,u)\int_s^uC(u,v)R(u,v)\,\dd v\,\dd u .
 \label{resp:fisher-response-primitive}
\end{equation}
We rewrite the inner integral with the chain rule. For $s\le u\le t$, the function $v\mapsto C(u,v)^2$ is absolutely continuous on $[s,u]$, with derivative $2C(u,v)(R(u,v)+E(u,v))$, and $C(u,u)=1$. Hence
\[
 \int_s^uC(u,v)R(u,v)\,\dd v=\frac{1-C(u,s)^2}2-J_n(u),
 \qquad J_n(u)=\int_s^uC(u,v)E(u,v)\,\dd v .
\]
Using $0\le C\le1$ and \eqref{resp:fisher-entire-triangle} with $v=s$ and $w=u$, we find $|J_n(u)|\le\varepsilon_n$. We define the memory integral
\begin{equation}\label{eq:fisher-cost-memory-integral}
 I_n=\int_s^t\bigl[1-C(u,s)^2\bigr]R(t,u)\,\dd u .
\end{equation}
Inserting the expression for the inner integral into \eqref{resp:fisher-response-primitive}, writing $\lambda=\ell+(\lambda-\ell)$, and rearranging, we find the exact identity
\begin{equation}\label{eq:fisher-cost-exact-rearrangement}
 \begin{split}
 1-\ell(1-c_n)+3\beta^2I_n
 ={}&R(t,s)+\int_s^t(\lambda(u)-\ell)R(t,u)\,\dd u\\
 &+\ell\bigl[\chi(t,s)-(1-c_n)\bigr]+6\beta^2\int_s^tR(t,u)J_n(u)\,\dd u .
 \end{split}
\end{equation}
We bound the absolute values of the last three terms. The first is at most $\sup_{u\ge s_n}|\lambda(u)-\ell|\,(1+\varepsilon_n)$, which tends to zero because $\lambda(u)\to\ell$ (\cref{sel:selected-sector}) and $\varepsilon_n$ is bounded. By \eqref{resp:fisher-chi-bound}, the second is at most $\ell\varepsilon_n$, and the third is at most $6\beta^2\varepsilon_n\chi(t,s)\le6\beta^2(1+\varepsilon_n)\varepsilon_n$. We conclude that
\begin{equation}
 \bigl|1-\ell(1-c_n)+3\beta^2I_n\bigr|
 \le R(t,s)+o(1)+\bigl[\ell+6\beta^2(1+\varepsilon_n)\bigr]\varepsilon_n,
 \label{resp:fisher-primitive-error}
\end{equation}
where $o(1)$ tends to zero as $n\to\infty$.

\paragraph{Step 3: the recent strip and the ends of the bulk.}
Fix $D>0$, and let $n$ be so large that $t-s>4D$. Let $s'=s+2D$ and $t'=t-2D$, and split $I_n$ into the integrals over $[s,s']$, $[s',t']$, and $[t',t]$. The first is nonnegative, since $R\ge0$ and $C\le1$.

On the recent strip, we first show that $C(u,s)$ is close to $c_n$. For $u\in[t',t]$,
\[
 |C(u,s)-c_n|\le\int_u^t|\partial_xC(x,s)|\,\dd x\le2D\sup_{x\in[t',t]}|\partial_xC(x,s)| .
\]
Fix $L>0$. For large $n$, every $x\in[t',t]$ satisfies $x-s\ge t-s-2D\ge L$, and $x\ge t'\to\infty$. Then \eqref{resp:remote-forward} implies that $\limsup_n\sup_{x\in[t',t]}|\partial_xC(x,s)|\le\sqrt{c''(2L)}$. The right side tends to zero as $L\to\infty$, by the second inequality in \eqref{resp:remote-forward}. Together with $c_n\to z$, this shows that $\sup_{u\in[t',t]}|C(u,s)-z|\to0$. Using \eqref{sel:stationary-window}, we find that the response mass of the strip satisfies $\int_{t'}^tR(t,u)\,\dd u=\int_0^{2D}R(t,t-\tau)\,\dd\tau\to\int_0^{2D}r=1-c(2D)$. Therefore
\begin{equation}
 \int_{t'}^t\bigl[1-C(u,s)^2\bigr]R(t,u)\,\dd u\longrightarrow\bigl[1-c(2D)\bigr](1-z^2) .
 \label{resp:fisher-upper-strip}
\end{equation}

We next consider the values of $Y$ at the ends of the bulk. The stationary window \eqref{sel:stationary-window} also gives $Y(t')=C(t,t-2D)\to c(2D)$. At the other end,
\begin{align*}
 |Y(s')-c_n|&\le\int_s^{s+2D}R(t,u)\,\dd u+\int_s^{s+2D}|E(t,u)|\,\dd u\\
 &\le2D\sup_{s\le u\le s+2D}R(t,u)+\Bigl(2D\int_{s_n}^\infty a_F(u)\,\dd u\Bigr)^{1/2}.
\end{align*}
The first term tends to zero by \cref{resp:uniform-row}, because $t-u\ge t-s-2D\to\infty$ uniformly for $u\le s+2D$ and $r(\tau)\to0$. The second tends to zero because $a_F$ is integrable and $s_n\to\infty$. Therefore $Y(s')\to z$.

\paragraph{Step 4: the bulk.}
Let
\[
 \pi_{n,D}=\sup_{s'\le u\le t'}\bigl[Y(u)C(u,s)-q\,c_n\bigr]_+ .
\]
For $u\in[s',t']$, the times $s<u<t$ have gaps $u-s\ge2D$ and $t-u\ge2D$. For fixed $S$, we have $s_n\ge S+D$ for large $n$, and then the triple $(s,u,t)$ belongs to the set $\mathcal A_{S,D}$ of \cref{thm:geometry-actual-product}, and $[Y(u)C(u,s)-qc_n]_+=[C(t,u)C(u,s)-qC(t,s)]_+$ is the bracket in \eqref{eq:geometry-actual-product}. Letting $n\to\infty$ and then $S\to\infty$, we find
\begin{equation}\label{eq:fisher-cost-product-error}
 \lim_{D\to\infty}\limsup_{n\to\infty}\pi_{n,D}=0 .
\end{equation}

To use the product bound, we introduce a clipped function and its primitive. For $\kappa>0$, let
\[
 \varphi_\kappa(y)=\Bigl[1-\frac{\kappa^2}{y^2}\Bigr]_+\quad(y>0),\qquad\varphi_\kappa(0)=0,
\]
and
\[
 \Phi_\kappa(y)=\int_0^y\varphi_\kappa(x)\,\dd x=\begin{cases}0,&0\le y\le\kappa,\\ y+\kappa^2/y-2\kappa,&y>\kappa.\end{cases}
\]
The function $\Phi_\kappa$ is continuously differentiable on $[0,\infty)$, with $\Phi_\kappa'=\varphi_\kappa\in[0,1]$, and $(\kappa,y)\mapsto\Phi_\kappa(y)$ is continuous on $(0,\infty)\times[0,\infty)$.

Let $\kappa_n=qc_n$ and $\kappa_n'=\kappa_n+\pi_{n,D}$. For large $n$, we have $c_n\ge z/2>0$ and $\kappa_n>0$. We claim that
\[
 1-C(u,s)^2\ge\varphi_{\kappa_n'}(Y(u)),\qquad s'\le u\le t'.
\]
If $Y(u)\le\kappa_n'$, the right side is zero. Otherwise, by the definition of $\pi_{n,D}$, we have $Y(u)C(u,s)\le\kappa_n'$, that is, $0\le C(u,s)\le\kappa_n'/Y(u)$, and the claim follows. Next, for $0<\kappa\le\kappa'$ and $y\ge0$, we have
\[
 0\le\varphi_\kappa(y)-\varphi_{\kappa'}(y)\le1-\frac{\kappa^2}{\kappa'^2}\le\frac{2(\kappa'-\kappa)}{\kappa},
\]
as one checks separately for $y\le\kappa$, $\kappa<y\le\kappa'$, and $y>\kappa'$. With $\kappa=\kappa_n$ and $\kappa'=\kappa_n'$, and using \eqref{resp:fisher-chi-bound}, we obtain
\[
 \int_{s'}^{t'}\bigl[1-C(u,s)^2\bigr]R(t,u)\,\dd u
 \ge\int_{s'}^{t'}\varphi_{\kappa_n}(Y(u))R(t,u)\,\dd u-\frac{2(1+\varepsilon_n)}{qc_n}\pi_{n,D} .
\]
The function $Y$ is absolutely continuous on $[s',t']$ with $Y'=R(t,\cdot)+E(t,\cdot)$ almost everywhere, and $\Phi_{\kappa_n}$ is continuously differentiable with a bounded derivative. Hence $\Phi_{\kappa_n}\circ Y$ is absolutely continuous, with derivative $\varphi_{\kappa_n}(Y)Y'$, and
\[
 \int_{s'}^{t'}\varphi_{\kappa_n}(Y(u))R(t,u)\,\dd u
 =\Phi_{\kappa_n}(Y(t'))-\Phi_{\kappa_n}(Y(s'))-\int_{s'}^{t'}\varphi_{\kappa_n}(Y(u))E(t,u)\,\dd u .
\]
This identity is oriented: it does not require $Y$ to be monotone, and it holds when $Y$ returns to earlier values or crosses the level $\kappa_n$. By \eqref{resp:fisher-entire-triangle} and $0\le\varphi_{\kappa_n}\le1$, the last integral is at most $\varepsilon_n$ in absolute value. We conclude that
\begin{equation}
 \int_{s'}^{t'}\bigl[1-C(u,s)^2\bigr]R(t,u)\,\dd u
 \ge\Phi_{qc_n}(Y(t'))-\Phi_{qc_n}(Y(s'))-\varepsilon_n-\frac{2(1+\varepsilon_n)}{qc_n}\pi_{n,D} .
 \label{resp:fisher-bulk-calibration}
\end{equation}

\paragraph{Step 5: the limits.}
We first take the lower limit as $n\to\infty$ with $D$ fixed. We have $qc_n\to qz$, and by Step 3, $Y(t')\to c(2D)$ and $Y(s')\to z$. By the continuity of $(\kappa,y)\mapsto\Phi_\kappa(y)$,
\[
 \Phi_{qc_n}(Y(t'))-\Phi_{qc_n}(Y(s'))\longrightarrow\Phi_{qz}(c(2D))-\Phi_{qz}(z)=\int_z^{c(2D)}\Bigl(1-\frac{q^2z^2}{y^2}\Bigr)\dd y,
\]
where we used that $y\ge z>qz$ on $[z,c(2D)]$, so that no clipping occurs. Combining the nonnegativity of the early strip, \eqref{resp:fisher-upper-strip}, and \eqref{resp:fisher-bulk-calibration}, we obtain
\[
 \liminf_{n\to\infty}I_n\ge\bigl[1-c(2D)\bigr](1-z^2)+\int_z^{c(2D)}\Bigl(1-\frac{q^2z^2}{y^2}\Bigr)\dd y-\varepsilon-\frac{2(1+\varepsilon)}{qz}\limsup_{n\to\infty}\pi_{n,D} .
\]
The left side does not depend on $D$. We let $D\to\infty$, using $c(2D)\downarrow q$ and \eqref{eq:fisher-cost-product-error}. Since $\int_z^q(1-q^2z^2/y^2)\,\dd y=q-z-q^2z+qz^2$, this gives
\begin{equation}
 \liminf_{n\to\infty}I_n\ge(1-q)(1-z^2)+q-z-q^2z+qz^2-\varepsilon .
 \label{resp:fisher-memory-lower}
\end{equation}

\paragraph{Step 6: the calibrated margin.}
Let $\Psi(z)=(1-q)(1-z^2)+q-z-q^2z+qz^2$ be the main term in \eqref{resp:fisher-memory-lower}. Expanding, we find $\Psi(z)=1-(1+q^2)z+(2q-1)z^2$, and hence
\[
 \Psi(z)-(1-q^2)(1-z)=q^2-2q^2z+(2q-1)z^2 .
\]
By the second identity in \eqref{resp:fisher-calibration}, we have $1-\ell(1-z)=(z-q)/(1-q)-3\beta^2(1-q^2)(1-z)$, and by the first, $1/(1-q)=6\beta^2q(1-q)$. Therefore
\begin{equation}\label{eq:fisher-cost-calibrated-margin}
 \begin{split}
 1-\ell(1-z)+3\beta^2\Psi(z)
 &=3\beta^2\bigl[2q(1-q)(z-q)+q^2-2q^2z+(2q-1)z^2\bigr]\\
 &=3\beta^2(2q-1)(q-z)^2 .
 \end{split}
\end{equation}
In the second line, we expanded both sides: each equals $3\beta^2[(2q-1)z^2-4q^2z+2qz+2q^3-q^2]$.

We now combine the estimates. Since $s_n\to\infty$ and $t_n-s_n\to\infty$, \cref{resp:uniform-row} implies that $R(t_n,s_n)\to0$. Taking the lower limit of the left side of \eqref{resp:fisher-primitive-error} and using $c_n\to z$, \eqref{resp:fisher-memory-lower}, and \eqref{eq:fisher-cost-calibrated-margin}, we obtain
\[
 3\beta^2(2q-1)(q-z)^2-3\beta^2\varepsilon
 \le\liminf_{n\to\infty}\bigl[1-\ell(1-c_n)+3\beta^2I_n\bigr]
 \le\bigl[\ell+6\beta^2(1+\varepsilon)\bigr]\varepsilon .
\]
The memory bound \eqref{resp:fisher-memory-lower} loses $3\beta^2\varepsilon$, and the response equation \eqref{resp:fisher-primitive-error} loses $\ell\varepsilon+6\beta^2\varepsilon+6\beta^2\varepsilon^2$. Adding these losses, we arrive at
\[
 3\beta^2(2q-1)(q-z)^2\le(\ell+9\beta^2)\varepsilon+6\beta^2\varepsilon^2,
\]
which is \eqref{resp:sequential-obstruction}.
\end{proof}

\begin{proof}[Proof of \cref{resp:fisher-cost}]
We first prove \eqref{resp:fisher-window-limit}. Let $t_n=s_n+\tau_n$. Then $\varepsilon_n=(\tau_n\int_{s_n}^{t_n}a_F)^{1/2}\to0$ by \eqref{resp:fisher-small-window}, and \eqref{eq:geometry-remote-envelope} implies that $\limsup_nC(t_n,s_n)\le q$. Since $C$ takes values in $[0,1]$, it suffices to show that every subsequential limit $z$ of $C(t_n,s_n)$ equals $q$; we already know that $z\le q$. Passing to a subsequence, we may assume that $C(t_n,s_n)\to z$.

If $z\in(0,q)$, then \cref{resp:fisher-obstruction} with $\varepsilon=0$ implies that $3\beta^2(2q-1)(q-z)^2\le0$. This is impossible, because $q>1/2$ and $z<q$.

If $z=0$, we fix $z'\in(0,q)$. For large $n$, we have $C(t_n,s_n)<z'$, while $C(s_n,s_n)=1>z'$. The correlation is continuous, and the first crossing lag $\sigma_n=\inf\{\sigma\ge0:C(s_n+\sigma,s_n)\le z'\}$ satisfies $\sigma_n\le\tau_n$ and $C(s_n+\sigma_n,s_n)=z'$. The lags $\sigma_n$ tend to infinity. Indeed, for fixed $W>0$, \eqref{sel:stationary-window} shows that $C(s_n+\sigma,s_n)\to c(\sigma)\ge q>z'$ uniformly in $\sigma\in[0,W]$, and this forces $\sigma_n>W$ for large $n$. Since $a_F\ge0$ and $\sigma_n\le\tau_n$,
\[
 \sigma_n\int_{s_n}^{s_n+\sigma_n}a_F(u)\,\dd u\le\tau_n\int_{s_n}^{s_n+\tau_n}a_F(u)\,\dd u\longrightarrow0 .
\]
\Cref{resp:fisher-obstruction}, applied to the windows $[s_n,s_n+\sigma_n]$ with $z'$ in place of $z$ and with $\varepsilon=0$, again gives a contradiction. We conclude that $z=q$, and that $C(t_n,s_n)\to q$. Finally, by \eqref{resp:fisher-chi-bound}, we have $|\chi(t_n,s_n)-(1-C(t_n,s_n))|\le\varepsilon_n\to0$, and therefore $\chi(t_n,s_n)\to1-q$.

We now turn to the passage lag. Fix a level $c\in(0,q)$ and $s\ge0$. By \cref{sel:fixed-wait-forgetting}, $C(s+\tau,s)\to0$ as $\tau\to\infty$, and the set $\{\tau\ge0:C(s+\tau,s)\le c\}$ is nonempty. It is closed, because $C$ is continuous. Then $\tau_c(s)$ is finite and $C(s+\tau_c(s),s)\le c$. We have $C(s,s)=1$ and $C(s+\tau,s)>c$ for $0\le\tau<\tau_c(s)$, and by continuity, $C(s+\tau_c(s),s)=c$. As in the case $z=0$ above, \eqref{sel:stationary-window} and the bound $c(\tau)\ge q>c$ on the profile show that, for each fixed $W$, we have $C(s+\tau,s)>c$ for all $\tau\in[0,W]$ once $s$ is large. We conclude that $\tau_c(s)\to\infty$.

Let $A=\liminf_{s\to\infty}\tau_c(s)\int_s^{s+\tau_c(s)}a_F(u)\,\dd u\in[0,\infty]$. If $A\ge1$, then \eqref{resp:local-fisher-cost} holds, since $\mathfrak c_\beta(c)\le1$. Suppose that $A<1$, and choose $s_n\to\infty$ along which the quantity in the definition of $A$ converges to $A$. Let $t_n=s_n+\tau_c(s_n)$. Then $t_n-s_n\to\infty$, $C(t_n,s_n)=c\in(0,q)$, and $\varepsilon_n\to\sqrt A<1$. By \cref{resp:fisher-obstruction} with $z=c$ and $\varepsilon=\sqrt A$, and since $A\le\sqrt A$, we have
\[
 3\beta^2(2q-1)(q-c)^2
 \le(\ell+9\beta^2)\sqrt A+6\beta^2A
 \le(\ell+9\beta^2)\sqrt A+6\beta^2\sqrt A
 =(\ell+15\beta^2)\sqrt A .
\]
Squaring, we find $A\ge[3\beta^2(2q-1)(q-c)^2/(\ell+15\beta^2)]^2\ge\mathfrak c_\beta(c)$. This proves \eqref{resp:local-fisher-cost}. Because $a_F\ge0$, enlarging the domain of integration from $[s,s+\tau_c(s)]$ to $[s,\infty)$ can only increase the left side, and \eqref{resp:tail-fisher-cost} follows.
\end{proof}

\Cref{resp:fisher-cost} bounds the physical time needed to relax below the plateau in terms of the Fisher majorant. It does not determine the growth of $\tau_c(s)$, and it does not identify an aging clock; see \cref{sec:discussion}.


\part{The stationary approach to the plateau}\label{part:relaxation}
\chapter{Algebraic relaxation in the stationary sector}\label{sec:stationary-relaxation}

In this section, we prove \cref{thm:stationary-critical-amplitude}. Together with the discussion of the exponent in \cref{sec:sc-statement}, it proves \cref{thm:main-stationary-tail}. For $\beta\ge\bd$, it describes how the selected stationary profile $c$ approaches its plateau $q$. The difference $c(t)-q$ decays like $A_\beta t^{-a}$, where the amplitude $A_\beta$ is finite and positive, and the exponent $a\in(0,1/2)$ is determined by $q$ through the relation \eqref{eq:stationary-critical-exponent}. Every fixed derivative of $c$ has the corresponding power asymptotic.

The argument is an analysis of one scalar equation. It uses four results proved earlier: the regularity of a stationary cluster (\cref{prop:stationary-clusters}), the scalar memory equation and the plateau identity (\cref{prop:scalar-clusters}), the selection of the plateau together with marginality (\cref{sel:selected-sector}), and the complete monotonicity of the profile (\cref{sel:cluster-complete-monotonicity}). The precise properties that we use are listed at the start of \cref{sec:sc-statement}. The results of this section prove \cref{thm:main-stationary-tail} and verify the power-law hypotheses in \cref{resp:remote-earlier,resp:thermal}. The section can be read once \cref{part:selection} is available.

The relation \eqref{eq:stationary-critical-exponent} between the exponent and the plateau is the classical exponent relation of mode-coupling theory. G\"otze analyzed the leading critical and beta-scaling asymptotics of schematic and full mode-coupling equations in \cite[p.~144, Equation~(11)]{goetze1984} and \cite[p.~195, Equation~(1b), and Sections~3 and~5]{goetze1985}, the latter under a generic instability with an isolated critical stability eigenvalue. The cited critical power-law expansions concern an exponent parameter strictly between zero and one. An explicit statement of the relation is given in \cite[Section~4, p.~881, paragraph following Equation~(6)]{goetze1998}. The liquid-theory review gives the exponent relation for $\lambda_{\mathrm{MCT}}\in[1/2,1)$ \cite[Section~3.3, Equation~(3.19), p.~306]{goetzesjogren1992}. The general polynomial discussion in \cite[Section~6.1.2, Equations~(6.11d) and~(6.17b), pp.~444--447]{goetze2009} includes $0<\lambda_{\mathrm{MCT}}<1$, but assumes that the constructed low-frequency solution determines the asymptote of the actual critical correlator; it explicitly states on p.~447 that a proof of this matching assumption is unavailable.
An aging analogue, relating the approach and departure exponents to the fluctuation--dissipation ratio, was analyzed in \cite[arXiv version~1, Sections~III and~VI.D--E]{cugliandololedoussal1996} and \cite[Section~IV, Equation~(IV.6)]{bckm1996}; see also \cite[arXiv version~2, Section~12.3.8, Equation~(12.36)]{cugliandolo2004}. For the pure spherical model, its parameter is $\vartheta=(p-2)(1-q)/(2q)$, with $0<\vartheta\le1/2$ for the selected aging profile, and it is half of the fluctuation--dissipation ratio $(p-2)(1-q)/q$ of the aging regime predicted in \cite[arXiv version~1, paragraph after Equation~(12), p.~6]{ck1993}, and the relation is written in this form in \cite[arXiv version~1, p.~4, paragraph following Equation~(12)]{kimlatz2001}. The structural theory of scalar mode-coupling equations in \cite[Section~2, Theorem~1 and Corollaries~1.1--1.2, pp.~232--233]{gs1995}, its extension to Brownian mixtures with completely monotone solutions \cite[arXiv version~1, Section~III]{fv2002}, and the regenerative-process viewpoint of \cite[Sections~1 and~5.2 and Appendix~C]{sjogren2003} are mathematical predecessors of the argument below. The critical power law, with an amplitude that depends on the transient dynamics, is a classical asymptotic result of mode-coupling theory \cite[Section~4, Equation~(6), pp.~880--881]{goetze1998}. We give a self-contained proof for the stationary profile selected by the CK solution, including a finite positive amplitude and the asymptotics of all fixed derivatives. The proof does not assume a power expansion of the profile and does not differentiate an asymptotic equivalence. The tools are classical: comparison principles, beta integrals, the Schur test with its exact remainder, a factorization of gamma variables, the key renewal theorem \cite{feller1971}, and the monotone density theorem \cite{bgt1987}.

\section{Inputs and statement of the result}\label{sec:sc-statement}

Throughout this section, we fix $p\ge3$ and $\beta\ge\bd$. Let $(c,r)$ be the stationary profile and $\ell=\ell_\beta$ the limiting multiplier given by \cref{sel:selected-sector}, and let $q=q_\beta$ be the plateau of $c$. We use the following properties of $(c,r,q,\ell)$, and no other information about the CK solution.
\begin{enumerate}[label=(S\arabic*)]
\item\label{item:sc-input-regularity} (\Cref{prop:stationary-clusters}.) The function $r$ is continuously differentiable on $[0,\infty)$, and $c'=-r$ and $c(0)=r(0)=1$.
\item\label{item:sc-input-scalar} (\Cref{prop:scalar-clusters}.) The function $r$ is strictly positive on $[0,\infty)$. With
\[
 \phi(y)=\gamma+\beta^2py^{p-1},\qquad \gamma=\ell-\beta^2p,
\]
we have the scalar memory equation \eqref{eq:scalar-memory},
\[
 c'(t)=-1+\int_0^t\phi(c(t-s))\,r(s)\,\dd s,\qquad t\ge0,
\]
and the plateau identity $\phi(q)(1-q)=1$ from \eqref{eq:plateau-multiplier}.
\item\label{item:sc-input-marginal} (\Cref{sel:selected-sector}.) The plateau satisfies $\qd\le q<1$ and the marginality condition
\begin{equation}\label{eq:sc-marginality}
 \beta^2p(p-1)q^{p-2}(1-q)^2=1.
\end{equation}
Moreover, $q=\qd$ if and only if $\beta=\bd$.
\item\label{item:sc-input-cm} (\Cref{sel:cluster-complete-monotonicity}.) There exists a probability measure $\nu$ on $[0,\infty)$ such that
\[
 c(t)=\int_{[0,\infty)}e^{-xt}\,\nu(\dd x),\qquad t\ge0,\qquad\nu(\{0\})=q.
\]
\end{enumerate}
The last sentence of \ref{item:sc-input-marginal} follows from \cref{sel:selected-sector}, which identifies $q$ as the larger root of \eqref{eq:sc-marginality}, and this root lies in $[(p-2)/p,1)$. The function $y\mapsto y^{p-2}(1-y)^2$ is strictly decreasing on this interval, which implies that the larger root is strictly increasing in $\beta$. The point $\qd$ lies in the same interval, and a direct substitution shows that it is a root when $\beta=\bd$.

\paragraph{The exponent.} We define
\begin{equation}\label{eq:stationary-critical-exponent}
 \vartheta=\frac{(p-2)(1-q)}{2q},\qquad
 \frac{\Gamma(1-a)^2}{\Gamma(1-2a)}=\vartheta,
 \qquad 0<a<\frac12,
\end{equation}
where $\Gamma$ is Euler's gamma function. We now check that the second relation determines $a$ uniquely. Recall that $q\ge\qd=(p-2)/(p-1)$. Then $(1-q)/q\le1/(p-2)$, which gives $0<\vartheta\le1/2$. Let $\psi=\Gamma'/\Gamma$ be the digamma function. For $0<\alpha<1/2$, we have
\[
 \frac{\dd}{\dd\alpha}\log\frac{\Gamma(1-\alpha)^2}{\Gamma(1-2\alpha)}
 =2\bigl[\psi(1-2\alpha)-\psi(1-\alpha)\bigr]<0,
\]
since $\psi$ is strictly increasing on $(0,\infty)$. The quotient equals $1$ at $\alpha=0$ and tends to $0$ as $\alpha\uparrow1/2$, because $\Gamma(1-2\alpha)\to\infty$. We conclude that it is a strictly decreasing bijection from $[0,1/2)$ onto $(0,1]$. In particular, there exists exactly one $a\in(0,1/2)$ satisfying \eqref{eq:stationary-critical-exponent}.

At $\beta=\bd$, we have $q=\qd$, and then $\vartheta=1/2$ for all $p$. Solving $\Gamma(1-a)^2/\Gamma(1-2a)=1/2$ numerically, by bisection in double-precision arithmetic with the logarithm of the gamma function, yields $a=0.3952630630\ldots$ at $\beta=\bd$. For $\beta>\bd$, we have $q>\qd$ by \ref{item:sc-input-marginal}. Then $\vartheta<1/2$, and $a$ lies strictly between this value and $1/2$. In the cubic model, $\vartheta=(1-q)/(2q)$. These numerical values are not used in the proofs.

The relation \eqref{eq:stationary-critical-exponent} has a simple meaning. We write $(f*g)(t)=\int_0^tf(t-s)g(s)\,\dd s$ for the convolution of two functions on $[0,\infty)$; we use this notation throughout the section. For the pure power $w(t)=t^{-a}$, the beta integral gives $(w*w)(t)=\mathrm B(1-a,1-a)\,t^{1-2a}$, where $\mathrm B$ is Euler's beta function. Using $\Gamma(2-2a)=(1-2a)\Gamma(1-2a)$, we find
\begin{equation}\label{eq:sc-power-balance}
 (w*w)'(t)=(1-2a)\mathrm B(1-a,1-a)\,t^{-2a}
 =\frac{\Gamma(1-a)^2}{\Gamma(1-2a)}\,w(t)^2 .
\end{equation}
In other words, \eqref{eq:stationary-critical-exponent} says that $(w*w)'=\vartheta w^2$. This is the balance of the two quadratic terms in the normalized equation \eqref{eq:sc-normalized} below; see \eqref{eq:sc-quadratic-balance}. Because both sides of this balance are quadratic, it fixes the exponent but not the amplitude.

\paragraph{The result.}

\begin{theorem}[Stationary critical relaxation]
\label{thm:stationary-critical-amplitude}
Fix $p\ge3$ and $\beta\ge\bd$. Let $c$ be the stationary profile of \cref{sel:selected-sector}, let $q$ be its plateau, and let $a\in(0,1/2)$ be defined by \eqref{eq:stationary-critical-exponent}. Then there exists $A_\beta\in(0,\infty)$ such that, as $t\to\infty$,
\begin{equation}
 c(t)-q\sim A_\beta t^{-a},\qquad
 (-1)^n c^{(n)}(t)\sim A_\beta(a)_n t^{-a-n}
 \quad\text{for every fixed integer }n\ge1,
 \label{eq:stationary-critical-amplitude}
\end{equation}
where $(a)_n=a(a+1)\cdots(a+n-1)$. In particular, $r(t)\sim aA_\beta t^{-1-a}$.
\end{theorem}

\begin{remark}\label{rem:sc-amplitude}
The exponent $a$ depends only on $\vartheta$, that is, on $p$ and $q$. The proof determines the amplitude as a limit, $A_\beta=(1-q)^{1+a}A_0$ with $A_0$ given in \eqref{eq:sc-amplitude-limit}, and gives no closed formula for it. In line with the remark after \eqref{eq:sc-power-balance}, marginality determines the exponent, while the amplitude is determined by the whole profile $c$.
\end{remark}

\paragraph{Outline of the proof.} The proof occupies the rest of this section. We first rescale time and the profile so that marginality says that the coefficient of the linear memory term equals one (\cref{lem:sc-normalized-equation}). The normalized profile $u$ then solves $u'+u+M(u)*u'=0$, where $M$ is a polynomial with $M(y)=y+\vartheta y^2+O(y^3)$. The proof then has three steps.
\begin{enumerate}[label=\arabic*.]
\item \emph{Comparison profiles} (\cref{subsec:sc-normalization,subsec:sc-endpoint}). A comparison principle (\cref{lem:sc-comparison}) turns the sign of the residual of an explicit profile into an inequality between that profile and $u$. For the shifted power $g(t)=(1+t)^{-a}$, the choice of $a$ makes the quadratic part of the residual positive but of smaller order than $g^3$, and the sign of the residual is decided by cubic terms. We fix this sign by correcting $g$ with a small multiple of $g^2$ (\cref{lem:sc-shifted,lem:sc-critical-residuals}). The comparison profiles keep $u$ itself up to a fixed time and continue it by a rescaled corrected power. This yields $A_-(1+t)^{-a}\le u(t)\le A_+(1+t)^{-a}$ (\cref{thm:stationary-critical-endpoint}). These bounds identify the exponent, but they still allow $t^au(t)$ to oscillate between $A_-$ and $A_+$ on a logarithmic time scale.
\item \emph{The logarithmic Dirichlet form} (\cref{subsec:sc-dirichlet}). The function $Q=u*u$ satisfies $Q'=\vartheta u^2+\mathcal E$, where the error $\mathcal E$ is integrable against $t^{-(1-2a)}$ (\cref{lem:stationary-critical-balance-error}). Integrating this balance against the same weight and comparing the result with the Schur test for a homogeneous kernel, we deduce that $D(x)=e^{ax}u(e^x)$ has a finite Dirichlet energy for an explicit kernel on the line (\cref{lem:stationary-critical-dirichlet}). This implies that $D(x+h)-D(x)\to0$ as $x\to\infty$ for every fixed $h$. This still allows $D$ to vary slowly without converging.
\item \emph{The renewal equation} (\cref{subsec:sc-renewal}). We take the Laplace transform of the balance at the point $e^{-x}$. A factorization of gamma random variables turns the result into a linear renewal equation $f-\mathcal Bf=g_0$ for a smoothed version $f$ of $D^2$, with an integrable source $g_0$. We control the variance part of $g_0$ by the Dirichlet energy of the second step, and the rest by the error $\mathcal E$. Using the key renewal theorem, we show that $f$ converges. By the second step, $f-D^2\to0$, and by the first step the limit of $D$ lies in $(0,\infty)$. Finally, complete monotonicity and the monotone density theorem give the asymptotics of every derivative.
\end{enumerate}

\section{Normalization and comparison}
\label{subsec:sc-normalization}

We first rescale the profile so that marginality becomes a statement about a single coefficient. Recall the function $\phi$ and the plateau identity $\phi(q)(1-q)=1$ from \ref{item:sc-input-scalar}. We define the polynomial
\begin{equation}\label{eq:sc-memory-polynomial}
 M(y)=(1-q)\bigl[\phi(q+(1-q)y)-\phi(q)\bigr],\qquad y\in\R .
\end{equation}

\begin{lemma}[The normalized equation]\label{lem:sc-normalized-equation}
Define $u\colon[0,\infty)\to\R$ by
\[
 u(z)=\frac{c((1-q)z)-q}{1-q}.
\]
Then $u$ is continuously differentiable, and
\begin{equation}
 u'+u+M(u)*u'=0,\qquad u(0)=1,\qquad
 M(y)=y+\sum_{j=2}^{p-1}\lambda_jy^j,
 \quad\lambda_2=\vartheta,
 \label{eq:sc-normalized}
\end{equation}
where
\begin{equation}\label{eq:sc-memory-coefficients}
 \lambda_j=\frac1{p-1}\binom{p-1}{j}\Bigl(\frac{1-q}{q}\Bigr)^{j-1}\ge0,\qquad 2\le j\le p-1.
\end{equation}
Further, there exists a probability measure $\mu$ on $(0,\infty)$ such that
\begin{equation}\label{eq:sc-u-spectral}
 u(z)=\int_{(0,\infty)}e^{-\eta z}\,\mu(\dd\eta),\qquad z\ge0.
\end{equation}
In particular, $u$ is strictly positive and completely monotone on $(0,\infty)$, and $u(z)\to0$ as $z\to\infty$. The function $v=-u'$ is continuous, strictly positive, and nonincreasing on $[0,\infty)$, with $v(0)=1$ and $\int_0^\infty v(z)\,\dd z=1$.
\end{lemma}

\begin{proof}
Recall from \ref{item:sc-input-regularity} that $c$ is continuously differentiable with $c'=-r$. By the chain rule, $u'(z)=c'((1-q)z)=-r((1-q)z)$; the factor $1-q$ produced by the chain rule cancels the normalization. In particular, $u$ is continuously differentiable, $u(0)=1$, and $v(z)=r((1-q)z)$; by \ref{item:sc-input-scalar}, $v$ is strictly positive, and $v(0)=r(0)=1$.

Fix $z\ge0$. We use \eqref{eq:scalar-memory} at the time $(1-q)z$ and substitute $s=(1-q)w$ in its integral. Using $r((1-q)w)=v(w)$ and $c((1-q)(z-w))=q+(1-q)u(z-w)$, we find
\[
 u'(z)=-1+(1-q)\int_0^z\phi\bigl(q+(1-q)u(z-w)\bigr)\,v(w)\,\dd w .
\]
By \eqref{eq:sc-memory-polynomial}, we have $(1-q)\phi(q+(1-q)y)=(1-q)\phi(q)+M(y)=1+M(y)$, where we used the plateau identity. Inserting this and $\int_0^zv(w)\,\dd w=1-u(z)$ into the previous display, we conclude that
\[
 u'(z)=-1+\bigl(1-u(z)\bigr)+\int_0^zM(u(z-w))\,v(w)\,\dd w
 =-u(z)-(M(u)*u')(z),
\]
which is the first relation in \eqref{eq:sc-normalized}.

Next, we compute $M$. By the binomial theorem,
\[
 M(y)=(1-q)\beta^2p\bigl[(q+(1-q)y)^{p-1}-q^{p-1}\bigr]
 =\beta^2p\sum_{j=1}^{p-1}\binom{p-1}{j}q^{p-1-j}(1-q)^{j+1}y^j .
\]
The coefficient of $y$ is $\beta^2p(p-1)q^{p-2}(1-q)^2$, which equals $1$ by the marginality condition \eqref{eq:sc-marginality}. Dividing the coefficient of $y^j$ by this quantity, we find that it equals $\lambda_j$ in \eqref{eq:sc-memory-coefficients}. For $j=2$, this gives $\lambda_2=(p-2)(1-q)/(2q)=\vartheta$. All coefficients are nonnegative because $0<q<1$. For $p=3$, we have simply $M(y)=y+\vartheta y^2$.

Finally, by \ref{item:sc-input-cm}, we have $c(t)-q=\int_{(0,\infty)}e^{-xt}\,\nu(\dd x)$ and $\nu((0,\infty))=1-q>0$. Let $\mu$ be the image of the probability measure $(1-q)^{-1}\nu(\,\cdot\,\cap(0,\infty))$ under the map $x\mapsto(1-q)x$. Then \eqref{eq:sc-u-spectral} holds. By the representation, we have $u>0$, since $\mu$ is a nonzero measure on $(0,\infty)$. By dominated convergence, $u(z)\to0$ as $z\to\infty$. For $n\ge1$ and $z_0>0$, the function $\eta^ne^{-\eta z'}$ is bounded uniformly in $\eta>0$ and $z'\ge z_0$. By induction on $n$, this justifies differentiating $n$ times under the integral sign on $(z_0,\infty)$, and we have
\[
 (-1)^nu^{(n)}(z)=\int_{(0,\infty)}\eta^ne^{-\eta z}\,\mu(\dd\eta)>0,\qquad z>0,\ n\ge1 .
\]
For $n=2$, this shows that $v$ is nonincreasing on $(0,\infty)$, and by continuity also on $[0,\infty)$. Lastly, $\int_0^\infty v=u(0)-\lim_{z\to\infty}u(z)=1$.
\end{proof}

From now on, we write $t$ for the normalized time variable. The physical lag reappears only at the end of the proof of \cref{thm:stationary-critical-amplitude}. We record two properties of $M$. Its coefficients are nonnegative, and in particular $M$ is increasing and convex on $[0,\infty)$, with $M(0)=0$. The number $L_M=M'(1)$ is a Lipschitz constant for $M$ on $[0,1]$.

For a locally absolutely continuous, nonincreasing function $P\colon[0,\infty)\to[0,1]$ with $P(0)=1$, we define its residual
\[
 \mathcal R_P=P'+P+M(P)*P' .
\]
It is defined almost everywhere, and $\mathcal R_u=0$ by \eqref{eq:sc-normalized}. By the following lemma, a nonnegative residual gives an upper bound for $u$, and a nonpositive residual gives a lower bound.

\begin{lemma}[Comparison]
\label{lem:sc-comparison}
Let $P\colon[0,\infty)\to[0,1]$ be locally absolutely continuous and nonincreasing, with $P(0)=1$. If $\mathcal R_P\ge0$ almost everywhere, then $u\le P$. If $\mathcal R_P\le0$ almost everywhere, then $P\le u$.
\end{lemma}

\begin{proof}
We begin with the first assertion. Set $w=P-u$, and recall that $v=-u'\ge0$. We have
\[
 M(P)*P'-M(u)*u'=M(P)*w'-[M(P)-M(u)]*v .
\]
Subtracting \eqref{eq:sc-normalized} from the definition of $\mathcal R_P$ and using this identity, we find
\begin{equation}\label{eq:sc-comparison-difference}
 w'+w+M(P)*w'=\mathcal R_P+[M(P)-M(u)]*v .
\end{equation}
The functions $M\circ P$ and $w$ are locally absolutely continuous, and $w(0)=0$ and $P(0)=1$. Integrating by parts in the convolution, we have, for all $t\ge0$,
\[
 (M(P)*w')(t)=\int_0^tM(P(t-s))\,w'(s)\,\dd s=M(1)\,w(t)-(k_P*w)(t),
\]
where $k_P=-(M\circ P)'$.
We have $k_P\ge0$ almost everywhere and $\int_0^tk_P=M(1)-M(P(t))\le M(1)$, because $M$ is increasing on $[0,1]$ and $P$ is nonincreasing with values in $[0,1]$. Inserting this into \eqref{eq:sc-comparison-difference}, we obtain, with $b_0=1+M(1)$,
\begin{equation}\label{eq:sc-comparison-ode}
 w'+b_0w=\mathcal R_P+[M(P)-M(u)]*v+k_P*w\qquad\text{almost everywhere.}
\end{equation}

We now estimate the negative part $w_-=\max\{-w,0\}$. Recall that $M$ is increasing on $[0,1]$, with Lipschitz constant $L_M$ there. Since $u$ and $P$ take values in $[0,1]$, we have $M(P)-M(u)\ge-L_Mw_-$. We also have $w\ge-w_-$, and $v$ and $k_P$ are nonnegative. Using also $\mathcal R_P\ge0$, we deduce from \eqref{eq:sc-comparison-ode} that
\[
 \bigl(e^{b_0t}w(t)\bigr)'\ge-e^{b_0t}\bigl[L_M(w_-*v)(t)+(k_P*w_-)(t)\bigr]
 \qquad\text{almost everywhere.}
\]
The function $e^{b_0t}w(t)$ is locally absolutely continuous and vanishes at $t=0$. Integrating, and using $e^{-b_0(t-s)}\le1$, we find
\[
 w(t)\ge-\int_0^t\bigl[L_M(w_-*v)(s)+(k_P*w_-)(s)\bigr]\,\dd s .
\]
Since the integral on the right is nonnegative, $w_-(t)$ is bounded by the same integral. For a nonnegative locally integrable function $k$, we have, by Tonelli's theorem,
\[
 \int_0^t(w_-*k)(s)\,\dd s=\int_0^tk(\sigma)\int_0^{t-\sigma}w_-(s')\,\dd s'\,\dd\sigma
 \le\|k\|_{L^1(0,t)}\int_0^tw_-(s)\,\dd s .
\]
We have $\|v\|_{L^1(0,t)}=1-u(t)\le1$ and $\|k_P\|_{L^1(0,t)}\le M(1)$. Combining these bounds, we find
\[
 w_-(t)\le\bigl(L_M+M(1)\bigr)\int_0^tw_-(s)\,\dd s,\qquad t\ge0 .
\]
The function $w_-$ is continuous and nonnegative, and Gronwall's inequality implies that $w_-=0$, that is, $u\le P$.

For the second assertion, we set $w=u-P$ and $v_P=-P'\ge0$. Repeating the computation above, with the roles of $u$ and $P$ exchanged and with $u(0)=1$ in the integration by parts, we have
\[
 w'+b_0w=-\mathcal R_P+[M(u)-M(P)]*v_P+k_u*w,\qquad k_u=-(M\circ u)'\ge0,
\]
almost everywhere. Now $-\mathcal R_P\ge0$, $\|v_P\|_{L^1(0,t)}=1-P(t)\le1$, and $\|k_u\|_{L^1(0,t)}\le M(1)$. Arguing as in the previous paragraph, we conclude that $w_-=0$, that is, $P\le u$.
\end{proof}

\section{Comparison profiles at the critical exponent}
\label{subsec:sc-endpoint}

To apply \cref{lem:sc-comparison}, we need explicit profiles whose residuals have a definite sign. The natural candidates are multiples of the shifted power
\[
 g(t)=(1+t)^{-a}.
\]
We define
\[
 D_g=(g*g)'-\vartheta g^2,\qquad
 H=(g*g^2)'-\vartheta g^3,\qquad
 D_2=(g^2*g^2)'-\vartheta g^4,
\]
and, for $|\varepsilon|\le1/4$,
\[
 f_\varepsilon=(1+\varepsilon)g-\varepsilon g^2 .
\]
The next lemma explains why an uncorrected power does not suffice. By the choice of $a$, the quadratic part $D_g$ of the residual of $g$ is positive but of smaller order than $g^3$, and the sign of the residual is decided by cubic terms. The cross term $H$ that appears when $g$ is corrected by a multiple of $g^2$ is of order $g^3$ and has a definite sign. Unless stated otherwise, constants denoted by $K$, possibly with indices, depend only on $p$ and $q$, through $a$, $\vartheta$, and the coefficients $\lambda_j$, and on the indicated indices; their values may change from line to line. Constants that also depend on the profile $u$ are indicated as such.

\begin{lemma}[Shifted-power estimates]
\label{lem:sc-shifted}
For every $t\ge0$, we have $D_g(t)>0$. Moreover, as $t\to\infty$,
\begin{equation}
 D_g(t)=O((1+t)^{-a-1})=o(g(t)^3),\qquad
 \frac{H(t)}{g(t)^3}\longrightarrow
 (1-3a)\,\mathrm B(1-a,1-2a)-\vartheta<0.
 \label{eq:sc-shifted-limits}
\end{equation}
When $3a\ne1$, the limit equals $\Gamma(1-a)\Gamma(1-2a)/\Gamma(1-3a)-\vartheta$. There exists $K$ such that $|D_2|\le Kg^3$ on $[0,\infty)$. For every $\varepsilon$ with $|\varepsilon|\le1/4$, the profile $f_\varepsilon$ is positive and strictly decreasing, $f_\varepsilon(0)=1$, $f_\varepsilon\le Kg$, and $|f_\varepsilon'(t)|\le K(1+t)^{-a-1}$. Finally, for every integer $j\ge2$ there exists $K_j$ such that, for all $|\varepsilon|\le1/4$ and $t\ge0$,
\begin{equation}
 |(f_\varepsilon^j*f_\varepsilon)'(t)|\le K_jg(t)^3 .
 \label{eq:sc-shifted-derivative}
\end{equation}
\end{lemma}

\begin{proof}
The proof has four steps. We first express the derivative of a convolution of two shifted powers through a truncated beta integral. We then use this formula to prove that $D_g>0$, and next to prove \eqref{eq:sc-shifted-limits}. Finally, we bound derivatives of convolutions by splitting them at the midpoint, without differentiating an asymptotic formula.

\paragraph{Step 1: a beta formula.} Let $\alpha,\zeta\in(0,1)$ and $g_\alpha(t)=(1+t)^{-\alpha}$. For $t\ge0$, set $T=t+2$ and $x=1/T\in(0,1/2]$. Substituting $y=(1+s)/T$ in the convolution integral, we find
\[
 (g_\alpha*g_\zeta)(t)=\int_0^t(1+t-s)^{-\alpha}(1+s)^{-\zeta}\,\dd s
 =T^{1-\alpha-\zeta}\int_x^{1-x}y^{-\alpha}(1-y)^{-\zeta}\,\dd y,
\]
where we used the symmetry $y\mapsto1-y$ of the interval $[x,1-x]$. We differentiate in $t$, using $\dd T/\dd t=1$ and $\dd x/\dd t=-x^2$. The two endpoints of the integral contribute boundary terms, multiplied by $x^2T^{1-\alpha-\zeta}=xT^{-\alpha-\zeta}$, and we arrive at
\begin{align}
 (g_\alpha*g_\zeta)'(t)
 = T^{-\alpha-\zeta}\bigg[
 &(1-\alpha-\zeta)\int_x^{1-x}
             y^{-\alpha}(1-y)^{-\zeta}\,\dd y\nonumber\\
 &+x^{1-\alpha}(1-x)^{-\zeta}
  +x^{1-\zeta}(1-x)^{-\alpha}\bigg].
 \label{eq:sc-beta-derivative}
\end{align}
Both boundary terms are positive.

\paragraph{Step 2: positivity of $D_g$.} We take $\alpha=\zeta=a$ in \eqref{eq:sc-beta-derivative}. With $\omega(x)=[x(1-x)]^{-a}$, the bracket in \eqref{eq:sc-beta-derivative} is
\[
 J(x)=(1-2a)\int_x^{1-x}\omega(y)\,\dd y+2x\,\omega(x).
\]
Since $g(t)^2=(T-1)^{-2a}=T^{-2a}(1-x)^{-2a}$, we have
\begin{equation}\label{eq:sc-quotient-J}
 \frac{(g*g)'(t)}{g(t)^2}=(1-x)^{2a}J(x).
\end{equation}
Using $\omega(1-x)=\omega(x)$ and $\omega'(x)=-a(1-2x)\,\omega(x)/[x(1-x)]$, a direct computation shows that $J'(x)=2a\,\omega(x)/(1-x)$, and then
\[
 \frac{\dd}{\dd x}\bigl[(1-x)^{2a}J(x)\bigr]=2a(1-x)^{2a-1}L(x),
\]
where
\[
 L(x)=(1-2x)\,\omega(x)-(1-2a)\int_x^{1-x}\omega(y)\,\dd y .
\]
We have $L(1/2)=0$, and
\[
 L'(x)=-\frac{a\,\omega(x)}{x(1-x)}<0,\qquad 0<x<\frac12,
\]
where we used $(1-2x)^2+4x(1-x)=1$. It follows that $L>0$ on $(0,1/2)$, and that $(1-x)^{2a}J(x)$ is strictly increasing on $(0,1/2]$. As $x\downarrow0$, we have $2x\,\omega(x)\to0$ and $J(x)\to(1-2a)\mathrm B(1-a,1-a)$. As in \eqref{eq:sc-power-balance}, this limit equals $\Gamma(1-a)^2/\Gamma(1-2a)=\vartheta$. Every $t\ge0$ corresponds to some $x\in(0,1/2]$, and \eqref{eq:sc-quotient-J} implies that $(g*g)'(t)>\vartheta g(t)^2$, that is, $D_g(t)>0$, for all $t\ge0$, including $t=0$.

\paragraph{Step 3: the limits.} By \eqref{eq:sc-quotient-J}, we have $D_g(t)=T^{-2a}[J(x)-\vartheta(1-x)^{-2a}]$. Using $\vartheta=(1-2a)\int_0^1\omega$, we can write
\[
 J(x)-\vartheta(1-x)^{-2a}
 =-2(1-2a)\int_0^x\omega(y)\,\dd y+2x\,\omega(x)
 -\vartheta\bigl[(1-x)^{-2a}-1\bigr].
\]
For $0<y\le x\le1/2$, we have $\omega(y)\le2^ay^{-a}$. Then the first two terms are bounded by $Kx^{1-a}$, and the last term is bounded by $Kx\le Kx^{1-a}$. With $x=1/T$ and $1+t\le T\le2(1+t)$, this gives
\[
 |D_g(t)|\le KT^{-2a}x^{1-a}=KT^{-1-a}\le K(1+t)^{-1-a}.
\]
This is the first estimate in \eqref{eq:sc-shifted-limits}. We also have $D_g=o(g^3)$, because $(1+t)^{-1-a}=g(t)^3(1+t)^{-(1-2a)}$ and $1-2a>0$.

For the second limit, we take $(\alpha,\zeta)=(a,2a)$ in \eqref{eq:sc-beta-derivative}, which is allowed since $2a<1$. Using $g(t)^3=T^{-3a}(1-x)^{-3a}$, we find
\begin{multline*}
 \frac{(g*g^2)'(t)}{g(t)^3}=(1-x)^{3a}\bigg[(1-3a)\int_x^{1-x}y^{-a}(1-y)^{-2a}\,\dd y\\
 +x^{1-a}(1-x)^{-2a}+x^{1-2a}(1-x)^{-a}\bigg].
\end{multline*}
As $t\to\infty$, we have $x\to0$, the boundary terms tend to zero, and the integral tends to $\mathrm B(1-a,1-2a)$. This proves the convergence in \eqref{eq:sc-shifted-limits}. When $3a\ne1$, the identity $\Gamma(2-3a)=(1-3a)\Gamma(1-3a)$ yields $(1-3a)\mathrm B(1-a,1-2a)=\Gamma(1-a)\Gamma(1-2a)/\Gamma(1-3a)$.

It remains to show that the limit is negative. If $3a\ge1$, then $(1-3a)\mathrm B(1-a,1-2a)\le0<\vartheta$. If $3a<1$, consider
\[
 \Upsilon(\sigma)=\frac{\Gamma(1-a)\Gamma(1-\sigma)}{\Gamma(1-a-\sigma)},\qquad 0<\sigma<1-a .
\]
Its logarithmic derivative is $\psi(1-a-\sigma)-\psi(1-\sigma)<0$, and hence $\Upsilon$ is strictly decreasing. We have $\Upsilon(a)=\vartheta$ by \eqref{eq:stationary-critical-exponent}, and $2a<1-a$. We conclude that $(1-3a)\mathrm B(1-a,1-2a)=\Upsilon(2a)<\Upsilon(a)=\vartheta$.

\paragraph{Step 4: derivative bounds.} Let $F$ and $G$ be continuously differentiable on $[0,\infty)$. Splitting the convolution at $t/2$ and substituting $s\mapsto t-s$ on $[t/2,t]$, we have
\[
 (F*G)(t)=\int_0^{t/2}F(s)G(t-s)\,\dd s+\int_0^{t/2}G(s)F(t-s)\,\dd s .
\]
Differentiating both integrals, we find the exact identity
\begin{align}
 (F*G)'(t)={}&F(t/2)G(t/2)
 +\int_0^{t/2}F(s)G'(t-s)\,\dd s\nonumber\\
 &+\int_0^{t/2}G(s)F'(t-s)\,\dd s,
 \label{eq:sc-midpoint-derivative}
\end{align}
in which each of the two boundary terms contributed $F(t/2)G(t/2)/2$.

We next consider $f_\varepsilon$. We can write
\[
 f_\varepsilon=g\bigl[1+\varepsilon(1-g)\bigr],\qquad
 f_\varepsilon'=g'\bigl[1+\varepsilon(1-2g)\bigr],
\]
and both brackets lie in $[3/4,5/4]$ when $|\varepsilon|\le1/4$, because $0<g\le1$. Recalling that $g'(t)=-a(1+t)^{-a-1}<0$, we deduce that $f_\varepsilon$ is positive and strictly decreasing, and that
\begin{equation}\label{eq:sc-f-eps-bounds}
 \tfrac34 g\le f_\varepsilon\le\tfrac54 g,\qquad
 \tfrac34a(1+t)^{-a-1}\le|f_\varepsilon'(t)|\le\tfrac54a(1+t)^{-a-1}.
\end{equation}
Also $f_\varepsilon(0)=(1+\varepsilon)-\varepsilon=1$.

Fix $j\ge2$, and apply \eqref{eq:sc-midpoint-derivative} with $F=f_\varepsilon^j$ and $G=f_\varepsilon$. Using \eqref{eq:sc-f-eps-bounds}, we have $|F(t)|\le K(1+t)^{-ja}$, $|F'(t)|=jf_\varepsilon^{j-1}|f_\varepsilon'|\le K_j(1+t)^{-ja-1}$, $|G(t)|\le K(1+t)^{-a}$, and $|G'(t)|\le K(1+t)^{-a-1}$. For $0\le s\le t/2$, we have $1+t-s\ge(1+t)/2$. Inserting these bounds into the three terms of \eqref{eq:sc-midpoint-derivative}, we obtain
\begin{multline*}
 |(f_\varepsilon^j*f_\varepsilon)'(t)|\le K_j\bigg[(1+t)^{-(j+1)a}
 +(1+t)^{-a-1}\int_0^{t/2}(1+s)^{-ja}\,\dd s\\
 +(1+t)^{-ja-1}\int_0^{t/2}(1+s)^{-a}\,\dd s\bigg].
\end{multline*}
The last term is at most $K_j(1+t)^{-(j+1)a}$, since $\int_0^{t/2}(1+s)^{-a}\,\dd s\le(1+t)^{1-a}/(1-a)$. The integral in the middle term is at most $(1+t)^{1-ja}/(1-ja)$ if $ja<1$, at most $\log(2+t)$ if $ja=1$, and at most $1/(ja-1)$ if $ja>1$. Accordingly, the middle term is at most $K_j(1+t)^{-(j+1)a}$, $K(1+t)^{-a-1}\log(2+t)$, or $K_j(1+t)^{-a-1}$. Each of these bounds is at most $K_j(1+t)^{-3a}=K_jg(t)^3$: the first because $j\ge2$, and the other two because $a+1>3a$ and $(1+t)^{-(1-2a)}\log(2+t)$ is bounded. This proves \eqref{eq:sc-shifted-derivative}.

Finally, we take $(\alpha,\zeta)=(2a,2a)$ in \eqref{eq:sc-beta-derivative}, which is allowed since $2a<1$. The integral in the bracket is at most $\mathrm B(1-2a,1-2a)$, and the boundary terms are bounded on $(0,1/2]$. Then $|(g^2*g^2)'(t)|\le KT^{-4a}\le K(1+t)^{-3a}$. Together with $g^4\le g^3$, this shows that $|D_2|\le Kg^3$.
\end{proof}

We now correct $g$ by a multiple of $g^2$. The following quantity plays the role of the residual of a profile $hf(\cdot/\tau)$ attached at a time where $u$ has fallen to the value $h$; its precise relation to the residual $\mathcal R_P$ appears in \eqref{eq:sc-upper-decomposition} below. For a continuously differentiable function $f$ on $[0,\infty)$ and $h\in(0,1]$, we write $j_f=-f'$ and define the ideal residual
\begin{equation}
 I_h(f)=hf-(1-h)M(hf)-h[M(hf)*j_f].
 \label{eq:sc-ideal-residual}
\end{equation}

\begin{lemma}[Corrected critical profiles]
\label{lem:sc-critical-residuals}
There exist constants $\kappa,\eta,h_0>0$, depending only on $p$ and $q$, with $h_0\le1$ and $\kappa h_0\le1/4$, such that the following holds for every $h\in(0,h_0]$. The profiles
\[
 f_+=(1+\kappa h)g-\kappa hg^2,\qquad f_-=(1-\kappa h)g+\kappa hg^2
\]
are positive and strictly decreasing, with $f_\pm(0)=1$. Further, $I_h(f_+)\ge\eta h^3g^3$ on $[0,\infty)$, and there exists $X_h<\infty$ such that $I_h(f_-)\le-\eta h^3g^3$ on $[X_h,\infty)$.
\end{lemma}

\begin{proof}
The profiles are $f_\pm=f_{\pm\kappa h}$. Once $\kappa h_0\le1/4$, the first assertion follows from \cref{lem:sc-shifted}. We begin by expanding $I_h(f)$ for a continuously differentiable $f$ with $f(0)=1$. Differentiating $(f^j*f)(t)=\int_0^tf(s)^jf(t-s)\,\dd s$ and using $f(0)=1$, we have
\[
 (f^j*f)'=f^j-f^j*j_f,\qquad j\ge1 .
\]
Inserting $M(y)=\sum_{j=1}^{p-1}\lambda_jy^j$, with $\lambda_1=1$, into \eqref{eq:sc-ideal-residual} and using this identity, we obtain
\[
 I_h(f)=hf-\sum_{j=1}^{p-1}\lambda_jh^jf^j+\sum_{j=1}^{p-1}\lambda_jh^{j+1}(f^j*f)' .
\]
The term $hf$ cancels against the term $j=1$ of the first sum. Separating the terms with $j=1$ and $j=2$ from the rest, we have
\begin{align}
 I_h(f)={}&h^2[(f*f)'-\vartheta f^2]+N_h(f),\nonumber\\
 N_h(f)={}&\vartheta h^3(f^2*f)'
 +\sum_{j=3}^{p-1}\lambda_j
       [-h^jf^j+h^{j+1}(f^j*f)'].
 \label{eq:sc-ideal-expansion}
\end{align}

We now bound $N_h(f_\varepsilon)$ for $|\varepsilon|\le1/4$ and $0<h\le1$. By \eqref{eq:sc-shifted-derivative}, we have $h^{j+1}|(f_\varepsilon^j*f_\varepsilon)'|\le K_jh^3g^3$ for $j\ge2$. Using \eqref{eq:sc-f-eps-bounds}, we have $h^jf_\varepsilon^j\le(5/4)^jh^3g^3$ for $j\ge3$, where we used $h\le1$ and $g\le1$. Then there exists $K_0$ such that
\begin{equation}\label{eq:sc-N-bound}
 |N_h(f_\varepsilon)|\le K_0h^3g^3,\qquad |\varepsilon|\le1/4,\quad 0<h\le1 .
\end{equation}
Next, expanding $f_\varepsilon*f_\varepsilon$ and $f_\varepsilon^2$ by bilinearity, we find the exact quadratic expansion
\begin{equation}
 (f_\varepsilon*f_\varepsilon)'-\vartheta f_\varepsilon^2
 =(1+\varepsilon)^2D_g
 -2\varepsilon(1+\varepsilon)H+\varepsilon^2D_2.
 \label{eq:sc-quadratic-expansion}
\end{equation}

We now fix the constants. Let
\[
 \Delta=\vartheta-(1-3a)\mathrm B(1-a,1-2a),
\]
which is positive by \cref{lem:sc-shifted}. By \eqref{eq:sc-shifted-limits}, there exists $t_*<\infty$ such that $H(t)\le-(\Delta/2)g(t)^3$ for $t\ge t_*$. There also exist $K_D$ and $K_H$ such that $|D_2|\le K_Dg^3$ and $|H|\le K_Hg^3$ on $[0,\infty)$, by \cref{lem:sc-shifted}; for $H$ we used that $H/g^3$ is continuous and has a finite limit. Let $\mu_*=\min_{[0,t_*]}D_g/g^3$, which is positive by \cref{lem:sc-shifted}. We set
\[
 \kappa=1+\frac{4K_0}{\Delta},
\]
so that $K_0\le\kappa\Delta/4$. Let $h_0\in(0,1]$ be so small that
\[
 \kappa h_0\le\frac14,\qquad \kappa^2K_Dh_0\le\frac{\kappa\Delta}4,\qquad
 \Bigl(\tfrac52\kappa K_H+\kappa^2K_D+K_0\Bigr)h_0\le\frac{\mu_*}2 ,
\]
and we set $\eta=\min\{\kappa\Delta/8,\mu_*/2\}$. Fix $h\in(0,h_0]$.

For the upper profile, we take $\varepsilon=\kappa h\in(0,1/4]$ in \eqref{eq:sc-quadratic-expansion}. On $[t_*,\infty)$, we have $(1+\varepsilon)^2D_g\ge0$, $-2\varepsilon(1+\varepsilon)H\ge\varepsilon\Delta g^3$, and $\varepsilon^2D_2\ge-\kappa^2h^2K_Dg^3$. Together with \eqref{eq:sc-ideal-expansion} and \eqref{eq:sc-N-bound}, these bounds imply that
\[
 I_h(f_+)\ge h^3\bigl[\kappa\Delta-\kappa^2hK_D-K_0\bigr]g^3\ge\frac{\kappa\Delta}2h^3g^3\ge\eta h^3g^3
 \qquad\text{on }[t_*,\infty),
\]
where we used $\kappa^2hK_D\le\kappa\Delta/4$ and $K_0\le\kappa\Delta/4$. On $[0,t_*]$, we divide \eqref{eq:sc-ideal-expansion} by $h^2g^3$ and use $(1+\varepsilon)^2\ge1$ and $|2\varepsilon(1+\varepsilon)|\le5\kappa h/2$. This yields
\[
 \frac{I_h(f_+)}{h^2g^3}\ge\mu_*-\Bigl(\tfrac52\kappa K_H+\kappa^2K_D+K_0\Bigr)h\ge\frac{\mu_*}2,
\]
that is, $I_h(f_+)\ge(\mu_*/2)h^2g^3\ge\eta h^3g^3$ on $[0,t_*]$, since $h\le1$.

For the lower profile, we take $\varepsilon=-\kappa h$. Then $1-\kappa h\ge3/4$, and on $[t_*,\infty)$ the cross term satisfies $2\kappa h(1-\kappa h)H\le-(3/4)\kappa h\Delta g^3$. Using also $(1-\kappa h)^2D_g\le D_g$, $\varepsilon^2D_2\le\kappa^2h^2K_Dg^3$, and \eqref{eq:sc-N-bound}, we find
\[
 I_h(f_-)\le h^2D_g+h^3\Bigl[-\frac34\kappa\Delta+\kappa^2hK_D+K_0\Bigr]g^3
 \le h^2D_g-\frac{\kappa\Delta}4h^3g^3
 \qquad\text{on }[t_*,\infty).
\]
The ratio $D_g/g^3$ tends to zero by \eqref{eq:sc-shifted-limits}. Then there exists $X_h\ge t_*$ such that $D_g\le(\kappa\Delta/8)hg^3$ on $[X_h,\infty)$, and there $I_h(f_-)\le-(\kappa\Delta/8)h^3g^3\le-\eta h^3g^3$. This completes the proof.
\end{proof}

The cutoff $X_h$ depends on $h$; since $D_g/g^3\le K(1+t)^{-(1-2a)}$, it can be taken of order $h^{-1/(1-2a)}$. In the next proof, $h$ is fixed before $X_h$ is used.

\begin{proposition}[Critical endpoint bounds]
\label{thm:stationary-critical-endpoint}
There exist constants $0<A_-\le A_+<\infty$ such that
\begin{equation}
 A_-(1+t)^{-a}\le u(t)\le A_+(1+t)^{-a},\qquad t\ge0.
 \label{eq:sc-endpoint-bounds}
\end{equation}
For every integer $n\ge1$, there exist constants $0<\underline K_n\le\overline K_n<\infty$ and $t_n<\infty$ such that
\[
 \underline K_n(1+t)^{-a-n}\le(-1)^nu^{(n)}(t)\le\overline K_n(1+t)^{-a-n},\qquad t\ge t_n .
\]
\end{proposition}

\begin{proof}
The proof has three steps. First, we construct a profile $P_+$ with nonnegative residual. It equals $u$ up to a time $T_0$, is constant up to $T_1=2T_0$, and then follows a copy of $f_+$ on a long time scale $\tau$; we choose the time scale large enough to absorb the error caused by the memory of the initial segment of $u$. Second, we construct a profile $P_-$ with nonpositive residual by attaching a copy of $f_-$ directly at $T_0$, on a short time scale. In both cases, we then obtain a bound for $u$ from \cref{lem:sc-comparison}. Third, we deduce the derivative bounds from complete monotonicity.

We will use the following consequence of \cref{lem:sc-shifted}. Let $f=f_{\pm}$ be one of the profiles of \cref{lem:sc-critical-residuals}. For $0\le\delta\le1$ and $x\ge0$, every $y\ge(x-\delta)_+$ satisfies $1+y\ge(1+x)/2$. By the mean value theorem and the bound on $|f_\varepsilon'|$ in \cref{lem:sc-shifted}, we have
\begin{equation}\label{eq:sc-shift-cost}
 0\le f\bigl((x-\delta)_+\bigr)-f(x)\le K\delta(1+x)^{-a-1},\qquad 0\le\delta\le1,\ x\ge0 .
\end{equation}
We also use the convolution bound
\begin{equation}\label{eq:sc-power-convolution}
 \int_0^x(1+z)^{-a-1}(1+x-z)^{-a-1}\,\dd z\le2\Bigl(1+\frac x2\Bigr)^{-a-1}\int_0^\infty(1+z)^{-a-1}\,\dd z\le K(1+x)^{-a-1},
\end{equation}
which follows from the symmetry $z\mapsto x-z$ and splitting the integral at $x/2$.

\paragraph{Step 1: the upper bound.} Let $\kappa,\eta,h_0$ be as in \cref{lem:sc-critical-residuals}. Since $u$ decreases to zero, there exists $T_0\ge1$ with $h=u(T_0)\le h_0$. We fix such a $T_0$; note that $h>0$. Let $T_1=2T_0$, and let $\tau\ge T_1$ be a parameter, to be fixed below. We write $f=f_+$ and $j=-f'$, and define
\[
 P_+(t)=
 \begin{cases}
 u(t),&0\le t\le T_0,\\
 h,&T_0\le t\le T_1,\\
 hf((t-T_1)/\tau),&t\ge T_1.
 \end{cases}
\]
This function is continuous, since $u(T_0)=h=hf(0)$. It is locally absolutely continuous, nonincreasing, takes values in $[0,1]$, and satisfies $P_+(0)=1$. We check that $\mathcal R_{P_+}\ge0$ on each of the three intervals.

On $(0,T_0)$, we have $\mathcal R_{P_+}=\mathcal R_u=0$, since the residual at time $t$ involves only the values of $P_+$ on $[0,t]$.

On $(T_0,T_1)$, we have $P_+'(t)=0$ and $P_+(t)=h$, and $-P_+'$ vanishes on $(T_0,t)$. For $0\le s\le T_0$, we have $t-s\ge T_0-s$, which implies $P_+(t-s)\le P_+(T_0-s)=u(T_0-s)$. Since $M$ is increasing on $[0,1]$, we have
\[
 \mathcal R_{P_+}(t)=h-\int_0^{T_0}M(P_+(t-s))\,v(s)\,\dd s\ge h-(M(u)*v)(T_0)=v(T_0)>0 .
\]
In the equality on the right, we used \eqref{eq:sc-normalized} at time $T_0$, in the form $(M(u)*v)(T_0)=u(T_0)+u'(T_0)=h-v(T_0)$.

On $(T_1,\infty)$, we write $x=(t-T_1)/\tau$. Then $P_+(t)=hf(x)$ and $P_+'(t)=(h/\tau)f'(x)$. The measure $-P_+'(s)\,\dd s$ equals $v(s)\,\dd s$ on $[0,T_0]$, vanishes on $(T_0,T_1)$, and equals $(h/\tau)j((s-T_1)/\tau)\,\dd s$ on $(T_1,\infty)$. For $0\le s\le T_0$, we have $t-s\ge T_1-T_0=T_0$ and $P_+(t-s)=hf((x-s/\tau)_+)$: this is the definition of $P_+$ if $t-s\ge T_1$, and if $T_0\le t-s<T_1$ both sides equal $h$. In the integral over $(T_1,t)$, we substitute $s=T_1+\tau(x-z)$, for which $t-s=\tau z$. Recalling that $\int_0^{T_0}v=1-u(T_0)=1-h$ and comparing with \eqref{eq:sc-ideal-residual}, we obtain the exact decomposition
\begin{equation}\label{eq:sc-upper-decomposition}
 \mathcal R_{P_+}(t)=\frac h\tau f'(x)+I_h(f)(x)-E_1(x)-E_2(x),
\end{equation}
where
\begin{align*}
 E_1(x)&=\int_0^{T_0}\Bigl[M\bigl(hf((x-s/\tau)_+)\bigr)-M\bigl(hf(x)\bigr)\Bigr]v(s)\,\dd s,\\
 E_2(x)&=h\int_0^x\bigl[M(P_+(\tau z))-M(hf(z))\bigr]j(x-z)\,\dd z .
\end{align*}
The terms $E_1$ and $E_2$ measure the difference between the actual memory of $P_+$, which contains the initial segment of $u$, and the ideal memory in $I_h(f)$. Both are nonnegative. For $E_1$, this holds because $f$ is decreasing and $M$ is increasing. For $E_2$, we have $P_+(\tau z)\ge hf(z)$ for all $z\ge0$: if $\tau z\le T_1$, then $P_+(\tau z)\ge h\ge hf(z)$, and if $\tau z\ge T_1$, then $P_+(\tau z)=hf(z-T_1/\tau)\ge hf(z)$.

We now bound $E_1$ and $E_2$ from above, keeping the whole initial segment of $u$. For $s\le T_0$, we have $s/\tau\le T_0/\tau\le1$. By \eqref{eq:sc-shift-cost} with $\delta=s/\tau$, the Lipschitz bound for $M$ on $[0,1]$, and $\int_0^{T_0}v\le1$, we have
\[
 0\le E_1(x)\le L_MKh\,\frac{T_0}\tau\,(1+x)^{-a-1}.
\]
For $E_2$, we split the integral at $z=T_0/\tau$. On $0\le z\le\min\{x,T_0/\tau\}$, the difference in brackets is at most $M(1)$, and $j(x-z)\le K(1+x-z)^{-a-1}\le K(1+x)^{-a-1}$, since $z\le1$. This part of $E_2$ is at most $KhM(1)(T_0/\tau)(1+x)^{-a-1}$. On $T_0/\tau\le z\le x$, we have $\tau z\ge T_0$ and $P_+(\tau z)=hf((z-T_1/\tau)_+)$. By \eqref{eq:sc-shift-cost} with $\delta=T_1/\tau\le1$, the difference in brackets is at most $L_MKh(T_1/\tau)(1+z)^{-a-1}$. Using also $j(x-z)\le K(1+x-z)^{-a-1}$ and \eqref{eq:sc-power-convolution}, we find that this part of $E_2$ is at most $Kh^2(T_1/\tau)(1+x)^{-a-1}$. Finally, $(h/\tau)f'(x)\ge-K(h/\tau)(1+x)^{-a-1}$.

We insert these bounds and $I_h(f_+)\ge\eta h^3g^3$ into \eqref{eq:sc-upper-decomposition}. Using also $(1+x)^{-a-1}\le g(x)^3$, which holds because $a+1>3a$, we conclude that
\[
 \mathcal R_{P_+}(t)\ge\Bigl[\eta h^3-\frac{K_1h}\tau\Bigr]g(x)^3,\qquad t>T_1,
\]
where $K_1$ depends on $p$, $q$, and $T_1$, but not on $\tau\ge T_1$. We now fix
\[
 \tau=\max\Bigl\{T_1,\frac{2K_1}{\eta h^2}\Bigr\}.
\]
Then $\mathcal R_{P_+}>0$ on $(T_1,\infty)$. We have shown that $\mathcal R_{P_+}\ge0$ almost everywhere, and \cref{lem:sc-comparison} implies that $u\le P_+$. For $t\ge T_1$, we have $\tau+t-T_1\ge t\ge(1+t)/2$, since $\tau\ge T_1$ and $t\ge1$. Using \eqref{eq:sc-f-eps-bounds}, we find
\[
 u(t)\le hf_+\Bigl(\frac{t-T_1}\tau\Bigr)\le\frac54h\Bigl(\frac{\tau+t-T_1}\tau\Bigr)^{-a}
 \le\frac54h(2\tau)^a(1+t)^{-a},\qquad t\ge T_1 .
\]
For $t\le T_1$, we have $u(t)\le1\le(1+T_1)^a(1+t)^{-a}$. This proves the upper bound in \eqref{eq:sc-endpoint-bounds}.

\paragraph{Step 2: the lower bound.} We now write $f=f_-$ and $j=-f'$. We fix $T_0$ with $h=u(T_0)\le h_0$ and the cutoff $X_h$ of \cref{lem:sc-critical-residuals}. For a parameter $\tau>0$, to be fixed below, we define
\[
 P_-(t)=
 \begin{cases}
 u(t),&0\le t\le T_0,\\
 hf((t-T_0)/\tau),&t\ge T_0.
 \end{cases}
\]
This function is continuous, locally absolutely continuous, nonincreasing, takes values in $[0,1]$, and satisfies $P_-(0)=1$. As before, $\mathcal R_{P_-}=0$ on $(0,T_0)$.

For $t>T_0$, let $x=(t-T_0)/\tau$. Splitting the memory integral of $\mathcal R_{P_-}(t)$ at $T_0$ and substituting $s=T_0+\tau(x-z)$ in the second part, we have
\[
 \mathcal R_{P_-}(t)=\frac h\tau f'(x)+hf(x)-\int_0^{T_0}M(P_-(t-s))\,v(s)\,\dd s
 -h\int_0^xM(P_-(\tau z))\,j(x-z)\,\dd z .
\]
The actual values of $P_-$ in the memory integrals are at least their ideal replacements. For $0\le s\le T_0$, we have $P_-(t-s)\ge P_-(t)=hf(x)$, since $P_-$ is nonincreasing. For $z\ge0$, we have $P_-(\tau z)\ge hf(z)$: if $\tau z\le T_0$, then $P_-(\tau z)=u(\tau z)\ge u(T_0)=h\ge hf(z)$, and if $\tau z\ge T_0$, then $P_-(\tau z)=hf(z-T_0/\tau)\ge hf(z)$. The function $M$ is increasing, and the memory terms enter with a negative sign. Recalling also that $\int_0^{T_0}v=1-h$, we conclude that
\[
 \mathcal R_{P_-}(t)\le\frac h\tau f'(x)+I_h(f)(x),\qquad t>T_0 .
\]
In contrast with the upper bound, this inequality has no error term.

It remains to fix $\tau$ so that the right side is nonpositive for all $x\ge0$. For $x\ge X_h$, both terms are nonpositive by \cref{lem:sc-critical-residuals}. On $[0,X_h]$, the function $I_h(f)$ is continuous, and $L_h=\max_{[0,X_h]}I_h(f)_+$ is finite. By \eqref{eq:sc-f-eps-bounds}, we have $-f'(x)\ge\delta_h=(3/4)a(1+X_h)^{-a-1}>0$ there. Let
\[
 \tau=\frac{h\delta_h}{L_h+1}.
\]
Then $(h/\tau)f'(x)\le-(L_h+1)<-I_h(f)(x)$ for $0\le x\le X_h$. We have shown that $\mathcal R_{P_-}\le0$ almost everywhere; the single jump of $P_-'$ at $T_0$ does not affect the argument. By \cref{lem:sc-comparison}, we have $P_-\le u$. For $t\ge T_0$, using \eqref{eq:sc-f-eps-bounds} and $1+(t-T_0)/\tau\le1+t/\tau\le\max\{1,1/\tau\}(1+t)$, we find
\[
 u(t)\ge hf_-\Bigl(\frac{t-T_0}\tau\Bigr)\ge\frac34h\Bigl(1+\frac{t-T_0}\tau\Bigr)^{-a}
 \ge\frac34h\min\{1,\tau\}^a(1+t)^{-a}.
\]
For $t\le T_0$, we have $u(t)\ge u(T_0)=h\ge h(1+t)^{-a}$. This proves the lower bound in \eqref{eq:sc-endpoint-bounds}.

\paragraph{Step 3: the derivatives.} For $n\ge0$ and $t>0$, let $F_n(t)=(-1)^nu^{(n)}(t)$. By \cref{lem:sc-normalized-equation}, each $F_n$ is positive. Since $F_n'=-F_{n+1}\le0$, each $F_n$ is nonincreasing. Using \eqref{eq:sc-endpoint-bounds}, we have $\underline k_0t^{-a}\le F_0(t)\le\overline k_0t^{-a}$ for $t\ge1$, with $\underline k_0=2^{-a}A_-$ and $\overline k_0=A_+$.

Suppose that, for some $n\ge0$ and $t_n'\ge1$,
\[
 \underline k_nt^{-\alpha}\le F_n(t)\le\overline k_nt^{-\alpha},\qquad t\ge t_n',\qquad \alpha=a+n .
\]
Using that $F_{n+1}$ is nonincreasing, we have, for $t\ge2t_n'$,
\[
 \frac t2F_{n+1}(t)\le\int_{t/2}^tF_{n+1}(s)\,\dd s=F_n(t/2)-F_n(t)\le F_n(t/2)\le2^\alpha\overline k_nt^{-\alpha},
\]
that is, $F_{n+1}(t)\le2^{1+\alpha}\overline k_nt^{-\alpha-1}$. For the lower bound, let $\Lambda>1$ be such that $\overline k_n\Lambda^{-\alpha}\le\underline k_n/2$. For $t\ge t_n'$, we have
\[
 \frac{\underline k_n}2t^{-\alpha}\le F_n(t)-F_n(\Lambda t)=\int_t^{\Lambda t}F_{n+1}(s)\,\dd s\le(\Lambda-1)\,t\,F_{n+1}(t),
\]
that is, $F_{n+1}(t)\ge\underline k_nt^{-\alpha-1}/(2(\Lambda-1))$. The same two-sided bound therefore holds for $F_{n+1}$, with $\alpha$ replaced by $\alpha+1$ and $t_{n+1}'=2t_n'$. By induction, it holds for all $n\ge0$. This proves the derivative bounds, because $(1+t)/2\le t\le1+t$ for $t\ge1$.
\end{proof}

The function $v$ is nonincreasing with $v(0)=1$. Together with the case $n=1$ of \cref{thm:stationary-critical-endpoint}, this implies that, for a constant $K$ depending on $u$,
\begin{equation}\label{eq:sc-v-bound}
 0<v(t)\le K(1+t)^{-a-1},\qquad t\ge0 .
\end{equation}

\section{The quadratic balance and its logarithmic Dirichlet form}
\label{subsec:sc-dirichlet}

The bounds \eqref{eq:sc-endpoint-bounds} identify the exponent, but they do not show that $t^au(t)$ has a limit. The goal of this section is to prove that $t^au(t)$ varies more and more slowly on a logarithmic time scale. We first isolate the quadratic balance behind \eqref{eq:sc-power-balance}.

\begin{lemma}[Integrable critical balance error]
\label{lem:stationary-critical-balance-error}
Let $Q=u*u$ and $\rho=1-2a$. Then the function $\mathcal E=Q'-\vartheta u^2$ satisfies
\begin{equation}
 Q'=\vartheta u^2+\mathcal E,
 \qquad |\mathcal E(t)|\le K(1+t)^{-3a},\qquad
 \int_0^\infty|\mathcal E(t)|t^{-\rho}\,\dd t<\infty,
 \label{eq:sc-quadratic-balance}
\end{equation}
where $K$ depends on $u$.
\end{lemma}

\begin{proof}
Recall that $v=-u'$. Differentiating $Q=u*u$ and using $u(0)=1$, we have $Q'=u-u*v$. By \eqref{eq:sc-normalized}, we have $u=v+M(u)*v=v+u*v+\sum_{j=2}^{p-1}\lambda_ju^j*v$. Inserting this into the expression for $Q'$, we find
\[
 Q'=v+\sum_{j=2}^{p-1}\lambda_j\,u^j*v .
\]
In the same way, we have $(u^j*u)'=u^j-u^j*v$. Inserting this identity and separating the term $\lambda_2u^2=\vartheta u^2$, we obtain $Q'=\vartheta u^2+\mathcal E$ with the exact error
\begin{equation}
 \mathcal E=v+\sum_{j=3}^{p-1}\lambda_ju^j
               -\sum_{j=2}^{p-1}\lambda_j(u^j*u)'.
 \label{eq:sc-error-exact}
\end{equation}
For $p=3$, the first sum is empty.

We bound each term. Using \eqref{eq:sc-endpoint-bounds} and \eqref{eq:sc-v-bound}, we have $u(t)\le K(1+t)^{-a}$ and $v(t)\le K(1+t)^{-a-1}$. The term $v$ is at most $K(1+t)^{-3a}$, since $a+1>3a$, and so is $u^j$ for $j\ge3$. For $j\ge2$, we apply \eqref{eq:sc-midpoint-derivative} with $F=u^j$ and $G=u$, which are continuously differentiable. We have $|F|\le K(1+t)^{-ja}$, $|F'|=ju^{j-1}v\le K_j(1+t)^{-ja-1}$, $|G|\le K(1+t)^{-a}$, and $|G'|=v\le K(1+t)^{-a-1}$. These are the bounds used in Step~4 of the proof of \cref{lem:sc-shifted}, and the computation there yields $|(u^j*u)'(t)|\le K_j(1+t)^{-3a}$, including the case $ja=1$. This proves the pointwise bound in \eqref{eq:sc-quadratic-balance}.

The error is bounded near zero and $\rho<1$, which implies that $\int_0^1|\mathcal E(t)|t^{-\rho}\,\dd t<\infty$. On $[1,\infty)$, we use $3a+\rho=1+a$ and find $\int_1^\infty|\mathcal E(t)|t^{-\rho}\,\dd t\le K\int_1^\infty t^{-1-a}\,\dd t<\infty$.
\end{proof}

We now pass to logarithmic time. In the following lemma, we convert the balance \eqref{eq:sc-quadratic-balance} into a finite Dirichlet energy for the function $e^{ax}u(e^x)$, which equals $t^au(t)$ at $t=e^x$.

\begin{lemma}[Logarithmic Dirichlet bound]
\label{lem:stationary-critical-dirichlet}
The function $D(x)=e^{ax}u(e^x)$, $x\in\R$, is bounded and Lipschitz on $\R$, it is bounded below by a positive constant on $[0,\infty)$, and
\begin{equation}
 \int_{\R}\int_{\R}
 \frac{[D(x)-D(y)]^2}{[2\cosh((x-y)/2)]^{2-2a}}
       \,\dd x\,\dd y<\infty.
 \label{eq:stationary-log-dirichlet-finite}
\end{equation}
In particular, for every fixed $h\in\R$,
\begin{equation}
 D(x+h)-D(x)\longrightarrow0\qquad(x\to\infty),
 \label{eq:stationary-log-increment-vanishing}
\end{equation}
and the convergence is uniform when $h$ ranges over a compact interval.
\end{lemma}

\begin{proof}
The proof has four steps. We first record the elementary properties of $D$. We then integrate the balance \eqref{eq:sc-quadratic-balance} against $t^{-\rho}$ up to a finite time $T$ and obtain a quantity $J_T$ that is bounded above uniformly in $T$. Next, we use the identity behind the Schur test to bound a truncated version of the energy in \eqref{eq:stationary-log-dirichlet-finite} by $J_T$. Finally, we deduce \eqref{eq:stationary-log-increment-vanishing}. We keep the cutoff $T$ finite until the last moment, because the separate integrals that appear in the second and third steps diverge as $T\to\infty$.

\paragraph{Step 1: properties of $D$.} By \eqref{eq:sc-endpoint-bounds}, we have
\[
 A_-\Bigl(\frac{e^x}{1+e^x}\Bigr)^a\le D(x)\le A_+\Bigl(\frac{e^x}{1+e^x}\Bigr)^a,\qquad x\in\R .
\]
In particular, $0<D\le A_+$, $D(x)\ge2^{-a}A_-$ for $x\ge0$, and $D(x)\to0$ as $x\to-\infty$. The derivative of $D$ is
\[
 D'(x)=ae^{ax}u(e^x)-e^{(a+1)x}v(e^x).
\]
For $x\ge0$, the bounds \eqref{eq:sc-endpoint-bounds} and \eqref{eq:sc-v-bound} show that both terms are bounded. For $x\le0$, both terms are bounded by $1$, since $u\le1$ and $v\le1$. We conclude that $D$ is Lipschitz on $\R$. Let $L_D$ denote a Lipschitz constant.

\paragraph{Step 2: the balance in integrated form.} Let
\[
 C_\rho=\mathrm B\Bigl(\frac{1+\rho}2,\frac{1+\rho}2\Bigr)=\mathrm B(1-a,1-a),
\]
so that $\vartheta=\rho C_\rho$ by \eqref{eq:sc-power-balance}. For $T>0$, we define
\[
 I_T=\int_0^TQ(t)t^{-1-\rho}\,\dd t,
 \qquad
 J_T=C_\rho\int_0^T u(t)^2t^{-\rho}\,\dd t-I_T.
\]
Note that $0\le Q(t)\le t$, because $0\le u\le1$. Then both integrals are finite, and $Q(t)t^{-\rho}\to0$ as $t\to0$. The function $Q$ is continuously differentiable. Integrating by parts, we find
\[
 \int_0^TQ'(t)t^{-\rho}\,\dd t=Q(T)T^{-\rho}+\rho I_T .
\]
We insert \eqref{eq:sc-quadratic-balance} on the left side, use $\vartheta=\rho C_\rho$, and divide by $\rho$. We arrive at the exact identity
\begin{equation}
 J_T=\frac{Q(T)T^{-\rho}}\rho
       -\frac1\rho\int_0^T\mathcal E(t)t^{-\rho}\,\dd t.
 \label{eq:sc-Schur-deficit}
\end{equation}
Using \eqref{eq:sc-endpoint-bounds}, we have
\[
 Q(T)\le A_+^2\int_0^T(1+T-s)^{-a}(1+s)^{-a}\,\dd s\le A_+^2\int_0^T(T-s)^{-a}s^{-a}\,\dd s=A_+^2C_\rho T^{\rho}.
\]
Together with \cref{lem:stationary-critical-balance-error}, this shows that
\[
 \sup_{T>0}J_T\le\frac{A_+^2C_\rho}\rho+\frac1\rho\int_0^\infty|\mathcal E(t)|t^{-\rho}\,\dd t<\infty .
\]

\paragraph{Step 3: the weighted Schur identity.} We define the symmetric kernel
\[
 \Xi(s,t)=\frac{(st)^{-a}}{(s+t)^{2-2a}},\qquad s,t>0 .
\]
Substituting $t=s\varsigma$ and using the beta integral $\int_0^\infty\varsigma^{-a}(1+\varsigma)^{2a-2}\,\dd\varsigma=\mathrm B(1-a,1-a)$, we find, for $s>0$,
\begin{equation}\label{eq:sc-Schur-row}
 \int_0^\infty\Xi(s,t)\,\dd t=\frac{C_\rho}s .
\end{equation}
Let $u_T=u\mathbf 1_{[0,T]}$ and $\Phi_T(s)=s^au_T(s)$. Using \eqref{eq:sc-Schur-row}, $2a-1=-\rho$, and $2-2a=1+\rho$, we have
\[
 \int_0^\infty\!\!\int_0^\infty\Xi(s,t)\Phi_T(s)^2\,\dd s\,\dd t=C_\rho\int_0^Tu(s)^2s^{-\rho}\,\dd s,
 \qquad
 \Xi(s,t)\Phi_T(s)\Phi_T(t)=\frac{u_T(s)u_T(t)}{(s+t)^{1+\rho}} .
\]
The first quantity is finite, since $u\le1$ and $\rho<1$. By the Cauchy--Schwarz inequality with respect to the measure $\Xi(s,t)\,\dd s\,\dd t$, the function $\Xi(s,t)|\Phi_T(s)\Phi_T(t)|$ is integrable as well. Expanding the square and using the symmetry of $\Xi$, we obtain
\begin{align}
 &\frac12\int_0^\infty\!\!\int_0^\infty
 \Xi(s,t)\bigl[\Phi_T(s)-\Phi_T(t)\bigr]^2\,\dd s\,\dd t\nonumber\\
 &\quad=C_\rho\int_0^T u(s)^2s^{-\rho}\,\dd s
     -\int_0^T\!\!\int_0^T\frac{u(s)u(t)}{(s+t)^{1+\rho}}
                         \,\dd s\,\dd t
 \le J_T.
 \label{eq:sc-Schur-square}
\end{align}
For the inequality, we used the representation of $I_T$ obtained by writing out $Q$ and applying Tonelli's theorem,
\[
 I_T=\iint_{\{s,t>0,\ s+t\le T\}}\frac{u(s)u(t)}{(s+t)^{1+\rho}}\,\dd s\,\dd t,
\]
together with the facts that the triangle $\{s+t\le T\}$ is contained in the square $[0,T]^2$ and that the integrand is nonnegative. The identity in \eqref{eq:sc-Schur-square} is the sum-of-squares identity behind the Schur test \cite[Section~7, Equation~(1), pp.~23--24]{schur1911} for the kernel $(st)^{\rho/2}(s+t)^{-1-\rho}$ on $L^2(0,\infty)$ with the test function $s^{-1/2}$, written in the variables $\Phi_T$.

As $T\to\infty$, the integrand on the left of \eqref{eq:sc-Schur-square} converges pointwise to $\Xi(s,t)[s^au(s)-t^au(t)]^2$; indeed, it equals this limit as soon as $T>\max\{s,t\}$. Fatou's lemma and Step~2 then imply that
\[
 \int_0^\infty\!\!\int_0^\infty
 \frac{(st)^{-a}}{(s+t)^{2-2a}}
 [s^au(s)-t^au(t)]^2\,\dd s\,\dd t\le2\sup_{T>0}J_T<\infty.
\]
We now substitute $s=e^x$ and $t=e^y$. Using $e^x+e^y=e^{(x+y)/2}\,2\cosh((x-y)/2)$, we have
\begin{align*}
 \Xi(e^x,e^y)\,e^{x+y}&=e^{-a(x+y)}\,e^{(a-1)(x+y)}\,[2\cosh((x-y)/2)]^{2a-2}\,e^{x+y}\\
 &=[2\cosh((x-y)/2)]^{-(2-2a)},
\end{align*}
and $s^au(s)=D(x)$. This proves \eqref{eq:stationary-log-dirichlet-finite}.

\paragraph{Step 4: vanishing increments.} Suppose that \eqref{eq:stationary-log-increment-vanishing} fails for some $h\in\R$. Then there exist $\epsilon>0$ and a sequence $x_k\to\infty$ with $|D(x_k+h)-D(x_k)|\ge\epsilon$. Let $\delta=\epsilon/(6L_D)$, and consider the regions
\[
 \mathcal S_k=\bigl\{(x,y):|x-x_k|\le\delta,\ |y-x-h|\le\delta\bigr\}.
\]
For $(x,y)\in\mathcal S_k$, we have $|x-x_k|\le\delta$ and $|y-(x_k+h)|\le2\delta$. Using the Lipschitz bound, we find
\[
 |D(y)-D(x)|\ge|D(x_k+h)-D(x_k)|-3L_D\delta\ge\epsilon/2 .
\]
Further, since $|x-y|\le|h|+\delta$ on $\mathcal S_k$, the kernel in \eqref{eq:stationary-log-dirichlet-finite} is at least $[2\cosh((|h|+\delta)/2)]^{-(2-2a)}>0$ there. Passing to a subsequence, we may assume that the intervals $[x_k-\delta,x_k+\delta]$ are disjoint; then the regions $\mathcal S_k$ are disjoint. Each of them has area $4\delta^2$ and contributes at least a fixed positive amount to the integral in \eqref{eq:stationary-log-dirichlet-finite}. This contradicts the finiteness of the integral, and proves \eqref{eq:stationary-log-increment-vanishing}.

For the uniformity, let $I$ be a compact interval and $\epsilon>0$. There exist finitely many points $h_1,\dots,h_N\in I$ such that every $h\in I$ satisfies $|h-h_i|\le\epsilon/(2L_D)$ for some $i$. Then $|D(x+h)-D(x)|\le|D(x+h_i)-D(x)|+\epsilon/2$, and the claim follows from \eqref{eq:stationary-log-increment-vanishing} for the finitely many $h_i$.
\end{proof}

\section{An integrable-source renewal equation}
\label{subsec:sc-renewal}

The vanishing of the increments in \eqref{eq:stationary-log-increment-vanishing} does not by itself imply that $D$ converges. In this section, we show that a smoothed version of $D^2$ satisfies a renewal equation with an integrable source and use the renewal theorem to prove that it converges.

\paragraph{The renewal theorem.} We use the renewal theorem in the following form. A probability measure $F$ on $(0,\infty)$ is called arithmetic if it is concentrated on a set of the form $\{h,2h,3h,\dots\}$ for some $h>0$. A function $z\colon[0,\infty)\to[0,\infty)$ is called directly Riemann integrable if, for every $h>0$, the upper and lower sums
\[
 \overline\sigma(h)=h\sum_{n\ge1}\sup_{[(n-1)h,nh]}z,\qquad
 \underline\sigma(h)=h\sum_{n\ge1}\inf_{[(n-1)h,nh]}z
\]
are finite, and they tend to the same limit as $h\to0$. This limit is then $\int_0^\infty z$.

\begin{theorem}[Key renewal theorem; {\cite[Chapter~XI, Section~1, pp.~362--363]{feller1971}}]\label{thm:sc-key-renewal}
Let $F$ be a probability measure on $(0,\infty)$ which is not arithmetic and has finite mean $m_F$, and let $F^{*n}$ denote its $n$-fold convolution, with $F^{*0}=\delta_0$. Let $z\colon[0,\infty)\to[0,\infty)$ be Borel measurable and directly Riemann integrable, and define
\[
 Z(x)=\sum_{n\ge0}\int_{[0,x]}z(x-y)\,F^{*n}(\dd y),\qquad x\ge0 .
\]
Then $Z(x)\to m_F^{-1}\int_0^\infty z(y)\,\dd y$ as $x\to\infty$.
\end{theorem}

We extend the test $z$ by zero to negative arguments. Borel measurability on $[0,\infty)$ is equivalent to Borel measurability of this extension and makes the integrals well-defined even for singular increment laws. The Darboux-sum condition alone does not imply measurability for such laws. The application below uses a continuous nonnegative source.

In \cite[Chapter~XI, Section~1, pp.~362--363]{feller1971}, the function $Z$ is the solution of the renewal equation $Z(x)=z(x)+\int_{[0,x]}Z(x-y)\,F(\dd y)$, written as $Z=U\star z$ with $U=\sum_{n\ge0}F^{*n}$, and the theorem is stated there as the alternative form of the renewal theorem. We check direct Riemann integrability through the following elementary criterion, which we prove for completeness.

\begin{lemma}[A criterion for direct Riemann integrability]\label{lem:sc-dri}
Let $z\colon[0,\infty)\to[0,\infty)$ be continuous, and suppose that $\sum_{k\ge0}s_k<\infty$, where $s_k=\sup_{[k,k+1]}z$. Then $z$ is directly Riemann integrable.
\end{lemma}

\begin{proof}
Let $S=\sum_ks_k$. For $0<h\le1$ and $n\ge1$, the interval $[(n-1)h,nh]$ lies in $[k,k+2]$ with $k=\lfloor(n-1)h\rfloor$, and its supremum is at most $s_k+s_{k+1}$. At most $2/h$ indices $n$ share the same value of $k$, which gives $\overline\sigma(h)\le4S$. For $h>1$, each interval $[(n-1)h,nh]$ is covered by the intervals $[k,k+1]$ that meet it, and each $[k,k+1]$ meets at most two of them. Then $\overline\sigma(h)\le2hS$. This shows that all upper and lower sums are finite. Now let $X\ge2$ be an integer. The function $z$ is continuous on $[0,X]$, and the parts of $\overline\sigma(h)$ and $\underline\sigma(h)$ coming from intervals contained in $[0,X]$ tend to $\int_0^Xz$ as $h\to0$. For $h\le1$, the remaining parts of both sums are at most $4\sum_{k\ge X-2}s_k$, by the argument above. This tail and $\int_X^\infty z$ tend to zero as $X\to\infty$, and we conclude that both sums converge to $\int_0^\infty z$.
\end{proof}

\paragraph{A renewal density.} The renewal equation that we derive is driven by a logarithmic beta variable, whose density is unbounded near zero. The following lemma records the behavior of its renewal density.

\begin{lemma}[Renewal density for a logarithmic beta variable]
\label{lem:log-beta-renewal-density}
Let $0<a<1$ and $\rho>0$, and let $Y=-\log B$, where $B$ has the beta distribution with parameters $(\rho,a)$. The density of $Y$ is
\[
 k(y)=\frac{e^{-\rho y}(1-e^{-y})^{a-1}}{\mathrm B(\rho,a)},
 \qquad y>0,
\]
and its mean is $m=\psi(\rho+a)-\psi(\rho)>0$.
The renewal density $U(y)=\sum_{n\ge1}k^{*n}(y)$ is locally integrable on $[0,\infty)$, bounded on $[1,\infty)$, and satisfies
\begin{equation}
 U(y)\longrightarrow\frac1m\qquad(y\to\infty).
 \label{eq:log-beta-renewal-limit}
\end{equation}
\end{lemma}

Here we set $k=0$ on $(-\infty,0]$, $k^{*1}=k$, and $k^{*(n+1)}=k^{*n}*k$. Each convolution power is a nonnegative Borel probability density supported in $(0,\infty)$. The parameter $a$ in this lemma is arbitrary in $(0,1)$; we apply it with the exponent $a$ of \eqref{eq:stationary-critical-exponent} and $\rho=1-2a$.

\begin{proof}
The density of $B$ is $b^{\rho-1}(1-b)^{a-1}/\mathrm B(\rho,a)$ on $(0,1)$. The change of variables $b=e^{-y}$ gives the formula for $k$. Differentiating $\log\mathrm B(\rho,a)=\log\int_0^1b^{\rho-1}(1-b)^{a-1}\,\dd b$ in $\rho$, we find $\E\log B=\psi(\rho)-\psi(\rho+a)$; the differentiation under the integral sign is justified because $|\log b|\,b^{\rho'-1}(1-b)^{a-1}$ is integrable on $(0,1)$, uniformly for $\rho'$ near $\rho$. We conclude that $m=\E Y=\psi(\rho+a)-\psi(\rho)$, which is positive since $\psi$ is strictly increasing.

We first bound the convolution powers. The function $(1-e^{-y})/y$ decreases from $1$ to $1-e^{-1}$ on $(0,1]$, and $y^{1-a}e^{-\rho y/2}$ is bounded on $[1,\infty)$. It follows that there exists $K$, depending only on $a$ and $\rho$, such that $k(y)\le Ky^{a-1}e^{-\rho y/2}$ for all $y>0$. The beta integral
\[
 \int_0^ys^{na-1}(y-s)^{a-1}\,\dd s=\frac{\Gamma(na)\Gamma(a)}{\Gamma((n+1)a)}\,y^{(n+1)a-1}
\]
and $e^{-\rho s/2}e^{-\rho(y-s)/2}=e^{-\rho y/2}$ imply, by induction on $n$, that
\begin{equation}
 k^{*n}(y)\le
 \frac{K^n\Gamma(a)^n}{\Gamma(na)}y^{na-1}e^{-\rho y/2},\qquad y>0,\ n\ge1 .
 \label{eq:sc-beta-renewal-envelope}
\end{equation}
The series $\sum_n(K\Gamma(a))^nY_0^{na}/\Gamma(na)$ converges for all $Y_0>0$, because $\Gamma(na)$ grows faster than any exponential in $n$. For $0<y\le1$, we have $y^{na-1}\le y^{a-1}$, so $U(y)\le K'y^{a-1}$. In particular, $U$ is finite on $(0,\infty)$ and locally integrable on $[0,\infty)$.

We next show that each $k^{*n}$ is continuous on $(0,\infty)$. Substituting $s_i=y\varrho_i$ in the $(n-1)$-fold integral that defines $k^{*n}(y)$, we have
\[
 k^{*n}(y)=y^{n-1}\int_{\Sigma_{n-1}}\prod_{i=1}^nk(y\varrho_i)\,\dd\varrho_1\cdots\dd\varrho_{n-1},
 \qquad\varrho_n=1-\sum_{i<n}\varrho_i,
\]
where $\Sigma_{n-1}=\{\varrho_i>0,\ \sum_{i<n}\varrho_i<1\}$. For $y$ in a compact subset of $(0,\infty)$, the integrand is continuous in $y$ and dominated by a constant times $\prod_i\varrho_i^{a-1}$, which is integrable on $\Sigma_{n-1}$. The continuity follows by dominated convergence.

Now fix an integer $j$ with $ja>1$, and let $k_j=k^{*j}$. By \eqref{eq:sc-beta-renewal-envelope}, we have $k_j(y)\to0$ as $y\downarrow0$, and $k_j$ extends to a continuous function on $[0,\infty)$ with $k_j(0)=0$. Its zero extension is therefore Borel measurable. The same bound shows that $\sum_{l\ge0}\sup_{[l,l+1]}k_j<\infty$. Then \cref{lem:sc-dri} implies that $k_j$ is directly Riemann integrable. It is the density of the sum of $j$ independent copies of $Y$. The corresponding probability measure is supported in $(0,\infty)$, it is not arithmetic because it has a density, and its finite positive mean is $jm$. Thus this law and the nonnegative test $z=k_j$ satisfy every hypothesis of \cref{thm:sc-key-renewal}. In this case, $Z=k_j+\sum_{n\ge1}k_j^{*n}*k_j=\sum_{n\ge1}k_j^{*n}$, where the interchange of sum and integral is justified by Tonelli's theorem. The theorem gives
\[
 U_j(y)=\sum_{n\ge1}k_j^{*n}(y)\longrightarrow\frac1{jm}\qquad(y\to\infty).
\]
Note that $k_j^{*n}=k^{*jn}$. Using \eqref{eq:sc-beta-renewal-envelope} with $jna-1>0$, we see that the series defining $U_j$ converges uniformly on every interval $[0,Y_0]$. Each term of this series is continuous on $[0,\infty)$, and so is $U_j$. Since $U_j$ is continuous and converges at infinity, it is bounded on $[0,\infty)$.

Finally, we decompose the indices $n\ge1$ according to their residues modulo $j$, writing $n=i+jl$ with $1\le i\le j$ and $l\ge0$. Using Tonelli's theorem, we obtain
\[
 U=\sum_{i=1}^j k^{*i}+\sum_{i=1}^j k^{*i}*U_j .
\]
By \eqref{eq:sc-beta-renewal-envelope}, each $k^{*i}$ is bounded on $[1,\infty)$ and tends to zero at infinity. For each $i$, we have $(k^{*i}*U_j)(y)=\int_0^\infty k^{*i}(s)U_j(y-s)\mathbf 1_{\{s\le y\}}\,\dd s$. The integrand is bounded by $\|U_j\|_\infty k^{*i}(s)$, which is integrable, and it converges to $k^{*i}(s)/(jm)$ for every $s$. Then $(k^{*i}*U_j)(y)\to1/(jm)$ by dominated convergence, and also $(k^{*i}*U_j)\le\|U_j\|_\infty$. Summing over $i$, we find that $U$ is bounded on $[1,\infty)$ and that $U(y)\to j/(jm)=1/m$. This proves \eqref{eq:log-beta-renewal-limit}.
\end{proof}

\paragraph{The gamma--beta factorization.} The second ingredient is a classical identity between gamma and beta variables. For a shape parameter $\alpha>0$, we say that a random variable has the gamma distribution with shape $\alpha$ if its density is $y^{\alpha-1}e^{-y}/\Gamma(\alpha)$ on $(0,\infty)$.

\begin{lemma}[Gamma--beta factorization]\label{lem:sc-gamma-beta}
Let $\rho,a>0$ and $b=\rho+a$. Let $G_b$ and $B$ be independent random variables, where $G_b$ has the gamma distribution with shape $b$ and $B$ has the beta distribution with parameters $(\rho,a)$. Then $BG_b$ has the gamma distribution with shape $\rho$.
\end{lemma}

\begin{proof}
Let $X$ and $W$ be independent, with gamma distributions of shapes $\rho$ and $a$. The map $(x,w)\mapsto(g,y)=(x+w,x/(x+w))$ is a bijection from $(0,\infty)^2$ onto $(0,\infty)\times(0,1)$, with inverse $(g,y)\mapsto(gy,g(1-y))$ and Jacobian $g$. The joint density of $(X+W,X/(X+W))$ is therefore
\[
 \frac{(gy)^{\rho-1}(g(1-y))^{a-1}e^{-g}}{\Gamma(\rho)\Gamma(a)}\,g
 =\frac{g^{b-1}e^{-g}}{\Gamma(b)}\cdot\frac{y^{\rho-1}(1-y)^{a-1}}{\mathrm B(\rho,a)},
\]
where we used $\Gamma(\rho)\Gamma(a)=\Gamma(b)\mathrm B(\rho,a)$. This shows that $X+W$ and $X/(X+W)$ are independent, with the laws of $G_b$ and $B$. Then $BG_b$ has the law of $X=[X/(X+W)]\cdot(X+W)$, which is the gamma distribution with shape $\rho$.
\end{proof}

\paragraph{The proof of the theorem.} We now combine the three steps.

\begin{proof}[Proof of \cref{thm:stationary-critical-amplitude}]
We keep the normalized time variable until the last step. Recall that $D(x)=e^{ax}u(e^x)$ and $\rho=1-2a$, and set $b=1-a=\rho+a$. Let $G_b$ have the gamma distribution with shape $b$, and let $B$ be an independent random variable with the beta distribution with parameters $(\rho,a)$. For bounded measurable functions $\zeta$ on $\R$, we define the averaging operators
\[
 (\mathcal K\zeta)(x)=\E\,\zeta(x+\log G_b),\qquad
 (\mathcal B\zeta)(x)=\E\,\zeta(x+\log B)=\int_0^\infty k(y)\,\zeta(x-y)\,\dd y,
\]
where $k$ is the density of $-\log B$ from \cref{lem:log-beta-renewal-density}. By the independence of $B$ and $G_b$ and Fubini's theorem, $\mathcal B\mathcal K\zeta(x)=\E\,\zeta(x+\log(BG_b))$. Recall from \cref{lem:sc-gamma-beta} that $BG_b$ has the gamma distribution with shape $\rho$.

The proof has five steps. We first transform the balance \eqref{eq:sc-quadratic-balance} into an exact equation in the logarithmic variable. We then show that it is a renewal equation with an integrable source, solve this equation by iteration, and use \cref{lem:log-beta-renewal-density} to find the limit of its solution. Next, we deduce that $D$ converges. Finally, we treat the derivatives and return to the physical lag.

\paragraph{Step 1: the balance in logarithmic time.} For $z>0$, we write $\widehat\zeta(z)=\int_0^\infty e^{-zt}\zeta(t)\,\dd t$ for the Laplace transform of a bounded function $\zeta$ on $[0,\infty)$. By \cref{lem:stationary-critical-balance-error}, the functions $u$, $u^2$, and $\mathcal E$ are bounded. The function $Q$ is continuously differentiable with $Q(0)=0$ and $0\le Q(t)\le t$. An integration by parts gives $\widehat{Q'}(z)=z\widehat Q(z)$, and $\widehat Q=\widehat u^{\,2}$ by Fubini's theorem. Taking the Laplace transform of \eqref{eq:sc-quadratic-balance}, we obtain
\begin{equation}\label{eq:sc-laplace-balance}
 z\,\widehat u(z)^2=\vartheta\,\widehat{u^2}(z)+\widehat{\mathcal E}(z),\qquad z>0 .
\end{equation}
We evaluate the transforms at $z=e^{-x}$, $x\in\R$. We have $u(t)=t^{-a}D(\log t)$, and the substitution $t=e^xw$ gives
\[
 \widehat u(e^{-x})=\int_0^\infty e^{-w}\,(e^xw)^{-a}D(x+\log w)\,e^x\,\dd w
 =e^{bx}\Gamma(b)\,(\mathcal KD)(x),
\]
where we used that $w^{b-1}e^{-w}/\Gamma(b)$ is the density of $G_b$. In the same way, since $BG_b$ has the gamma distribution with shape $\rho=1-2a$,
\[
 \widehat{u^2}(e^{-x})=\int_0^\infty e^{-w}(e^xw)^{-2a}D(x+\log w)^2\,e^x\,\dd w
 =e^{\rho x}\Gamma(\rho)\,\mathcal B\mathcal K(D^2)(x).
\]
Further, $e^{-x}\widehat u(e^{-x})^2=e^{(2b-1)x}\Gamma(b)^2(\mathcal KD)(x)^2=e^{\rho x}\Gamma(b)^2(\mathcal KD)(x)^2$. Inserting these three identities into \eqref{eq:sc-laplace-balance} at $z=e^{-x}$ and dividing by $e^{\rho x}\Gamma(b)^2$, we obtain the exact equation
\begin{equation}
 [\mathcal KD(x)]^2=\mathcal B\mathcal K(D^2)(x)+\sigma(x),
 \qquad
 \sigma(x)=\frac{e^{-\rho x}\widehat{\mathcal E}(e^{-x})}
                   {\Gamma(b)^2},
 \label{eq:sc-log-renewal-balance}
\end{equation}
where we used the exponent relation in the form $\vartheta\Gamma(\rho)=\Gamma(b)^2$, which is \eqref{eq:stationary-critical-exponent}.

The source $\sigma$ is integrable on $\R$. Indeed, for $t>0$, the substitution $w=te^{-x}$ yields $\int_\R e^{-\rho x}e^{-te^{-x}}\,\dd x=t^{-\rho}\Gamma(\rho)$. By Tonelli's theorem and \cref{lem:stationary-critical-balance-error}, we have
\begin{equation}
 \int_{\R}|\sigma(x)|\,\dd x
 \le\frac{\Gamma(\rho)}{\Gamma(b)^2}
       \int_0^\infty|\mathcal E(t)|t^{-\rho}\,\dd t<\infty.
 \label{eq:sc-log-source-integrable}
\end{equation}

\paragraph{Step 2: the linear renewal equation.} We set
\[
 f=\mathcal K(D^2),\qquad V=\mathcal K(D^2)-[\mathcal KD]^2 .
\]
For each $x$, $V(x)$ is the variance of $D(x+\log G_b)$; in particular, $V\ge0$. We have $\mathcal B\mathcal K(D^2)=\mathcal Bf$ and $[\mathcal KD]^2=f-V$, and \eqref{eq:sc-log-renewal-balance} becomes the linear renewal equation
\begin{equation}
 f-\mathcal Bf=g_0,\qquad g_0=V+\sigma .
 \label{eq:sc-linear-renewal}
\end{equation}

We now verify that $V$ is integrable. Let $G^{(1)}$ and $G^{(2)}$ be independent copies of $G_b$. Then
\[
 V(x)=\frac12\E\bigl[D(x+\log G^{(1)})-D(x+\log G^{(2)})\bigr]^2 .
\]
Integrating out $G^{(1)}$, we find that the ratio $G^{(2)}/G^{(1)}$ has the density $\varsigma^{b-1}(1+\varsigma)^{-2b}/\mathrm B(b,b)$ on $(0,\infty)$. Taking logarithms, we see that the difference $\log G^{(2)}-\log G^{(1)}$ has the density
\[
 \frac{e^{bh}}{\mathrm B(b,b)(1+e^h)^{2b}}=\frac{[2\cosh(h/2)]^{-2b}}{\mathrm B(b,b)},\qquad h\in\R .
\]
By Tonelli's theorem and the translation $x\mapsto x-\log G^{(1)}$, and since $2b=2-2a$, we obtain
\[
 \int_{\R}V(x)\,\dd x
 =\frac1{2\mathrm B(b,b)}\int_{\R}\int_{\R}
 \frac{[D(x)-D(y)]^2}{[2\cosh((x-y)/2)]^{2-2a}}
                 \,\dd x\,\dd y<\infty,
\]
where the finiteness is \eqref{eq:stationary-log-dirichlet-finite}. Together with \eqref{eq:sc-log-source-integrable}, this shows that $g_0\in L^1(\R)$.

We next record the regularity of the terms in \eqref{eq:sc-linear-renewal}. Averages of translates preserve bounds and Lipschitz constants. By \cref{lem:stationary-critical-dirichlet}, $D$ is bounded and Lipschitz. The functions $D^2$, $\mathcal KD$, $f$, $[\mathcal KD]^2$, $V$, and $\mathcal Bf$ inherit these properties. By \eqref{eq:sc-log-renewal-balance} and \eqref{eq:sc-linear-renewal}, the same holds for $\sigma$ and $g_0$. A uniformly continuous integrable function on $\R$ tends to zero at $\pm\infty$: otherwise there would be infinitely many disjoint intervals of a fixed length on which its absolute value is bounded below by a fixed positive number. In particular, $g_0(x)\to0$ as $x\to\pm\infty$. Further, $f(x)\to0$ as $x\to-\infty$, by dominated convergence, since $D$ is bounded and $D(y)\to0$ as $y\to-\infty$.

\paragraph{Step 3: the renewal representation and the limit of $f$.} The source $g_0$ is bounded, continuous, and integrable, but it may be signed. We use the local integrability and bounded tail of the renewal density from \cref{lem:log-beta-renewal-density} to handle this source. Let $Y_1,Y_2,\dots$ be independent copies of $Y=-\log B$, and let $S_n=Y_1+\cdots+Y_n$, with $S_0=0$. Then $\mathcal B^n\zeta(x)=\E\,\zeta(x-S_n)$, and for $n\ge1$ the variable $S_n$ has the density $k^{*n}$. Applying \eqref{eq:sc-linear-renewal} repeatedly, we have, for all $n\ge1$,
\[
 f(x)=\sum_{l=0}^{n-1}\E\,g_0(x-S_l)+\E\,f(x-S_n).
\]
By the strong law of large numbers, $S_n/n\to m>0$ almost surely, and in particular $S_n\to\infty$ almost surely. Then $\E f(x-S_n)\to0$ by bounded convergence, because $f$ is bounded and $f(-\infty)=0$. For every fixed $x$, we have
\begin{align*}
 \sum_{l\ge1}\E\,|g_0(x-S_l)|&=\int_0^\infty U(y)|g_0(x-y)|\,\dd y\\
 &\le\|g_0\|_\infty\int_0^1U(y)\,\dd y+\sup_{y\ge1}U(y)\,\|g_0\|_{L^1(\R)}<\infty,
\end{align*}
by \cref{lem:log-beta-renewal-density} and Tonelli's theorem. This justifies summing the series term by term, and we obtain
\begin{equation}
 f(x)=g_0(x)+\int_0^\infty U(y)g_0(x-y)\,\dd y,\qquad x\in\R .
 \label{eq:sc-renewal-representation}
\end{equation}

We now let $x\to\infty$ in \eqref{eq:sc-renewal-representation}. The term $g_0(x)$ tends to zero. The part of the integral with $0<y<1$ is at most $\sup_{[x-1,x]}|g_0|\int_0^1U$, which tends to zero. In the remaining part, we substitute $s=x-y$:
\[
 \int_1^\infty U(y)g_0(x-y)\,\dd y=\int_\R U(x-s)\,g_0(s)\,\mathbf 1_{\{s\le x-1\}}\,\dd s .
\]
The integrand is dominated by $\sup_{y\ge1}U(y)\,|g_0(s)|$, which is integrable, and by \eqref{eq:log-beta-renewal-limit} it converges to $g_0(s)/m$ for every $s$. We conclude by dominated convergence that
\begin{equation}\label{eq:sc-f-limit}
 f(x)\longrightarrow f_\infty=\frac1m\int_{\R}g_0(s)\,\dd s\qquad(x\to\infty).
\end{equation}

\paragraph{Step 4: the limit of $D$.} For every fixed $h\in\R$, we have
\[
 D(x+h)^2-D(x)^2=\bigl[D(x+h)-D(x)\bigr]\bigl[D(x+h)+D(x)\bigr]\longrightarrow0
\]
as $x\to\infty$, by \eqref{eq:stationary-log-increment-vanishing} and the boundedness of $D$. These differences are bounded by $2\|D\|_\infty^2$. By dominated convergence with respect to the law of $\log G_b$, we find
\[
 f(x)-D(x)^2=\E\bigl[D(x+\log G_b)^2-D(x)^2\bigr]\longrightarrow0 .
\]
Combined with \eqref{eq:sc-f-limit}, this shows that $D(x)^2\to f_\infty$. Recall from \cref{lem:stationary-critical-dirichlet} that $D$ is bounded, and bounded below by a positive constant on $[0,\infty)$. This implies that $f_\infty\in(0,\infty)$. Since $D>0$, we conclude that
\begin{equation}\label{eq:sc-amplitude-limit}
 D(x)\longrightarrow A_0=\sqrt{f_\infty}\in(0,\infty),\qquad\text{that is,}\qquad u(t)\sim A_0t^{-a}\quad(t\to\infty).
\end{equation}
We did not need to know the sign of $\int g_0$ in advance: its positivity follows from the lower bound in \eqref{eq:sc-endpoint-bounds}.

\paragraph{Step 5: the derivatives and the physical lag.} As in the proof of \cref{thm:stationary-critical-endpoint}, let $F_n=(-1)^nu^{(n)}$ on $(0,\infty)$ for $n\ge0$. Recall from \cref{lem:sc-normalized-equation} that each $F_{n+1}=-F_n'$ is positive and nonincreasing. We show by induction that
\begin{equation}\label{eq:sc-derivative-asymptotics}
 F_n(t)\sim A_0(a)_nt^{-a-n}\qquad(t\to\infty),\qquad n\ge0,
\end{equation}
where $(a)_0=1$. For $n=0$, this is \eqref{eq:sc-amplitude-limit}. Suppose that \eqref{eq:sc-derivative-asymptotics} holds for some $n\ge0$, and write $\alpha=a+n$ and $A_n=A_0(a)_n$. For $\Lambda>1$, since $F_{n+1}$ is nonincreasing, we have
\[
 \frac{F_n(t)-F_n(\Lambda t)}{(\Lambda-1)t}\le F_{n+1}(t)\le\frac{F_n(t/\Lambda)-F_n(t)}{(1-1/\Lambda)t}.
\]
Multiplying by $t^{\alpha+1}$ and letting $t\to\infty$, we obtain
\[
 A_n\,\frac{1-\Lambda^{-\alpha}}{\Lambda-1}\le\liminf_{t\to\infty}t^{\alpha+1}F_{n+1}(t)
 \le\limsup_{t\to\infty}t^{\alpha+1}F_{n+1}(t)\le A_n\,\frac{\Lambda^\alpha-1}{1-1/\Lambda}.
\]
Both bounds tend to $\alpha A_n$ as $\Lambda\downarrow1$, and we conclude that $F_{n+1}(t)\sim(a+n)A_0(a)_nt^{-a-n-1}=A_0(a)_{n+1}t^{-a-n-1}$, which is \eqref{eq:sc-derivative-asymptotics} for $n+1$. This argument is the monotone density theorem; see \cite[Section~1.7.3, Theorem~1.7.2 and its following tail-integral variant, p.~39]{bgt1987}.

Finally, we return to the physical lag, which we again denote by $t$. By the definition of $u$, we have $c(t)=q+(1-q)u(t/(1-q))$ and $c^{(n)}(t)=(1-q)^{1-n}u^{(n)}(t/(1-q))$ for $n\ge1$. Using \eqref{eq:sc-derivative-asymptotics}, we obtain
\[
 c(t)-q\sim(1-q)A_0\Bigl(\frac t{1-q}\Bigr)^{-a}=A_0(1-q)^{1+a}t^{-a},
\]
and, for all $n\ge1$,
\[
 (-1)^nc^{(n)}(t)\sim(1-q)^{1-n}A_0(a)_n\Bigl(\frac t{1-q}\Bigr)^{-a-n}=A_0(1-q)^{1+a}(a)_nt^{-a-n}.
\]
This is \eqref{eq:stationary-critical-amplitude} with $A_\beta=A_0(1-q)^{1+a}\in(0,\infty)$. The case $n=1$ gives $r(t)\sim aA_\beta t^{-1-a}$, because $r=-c'$. This completes the proof.
\end{proof}

\addtocontents{toc}{\protect\addvspace{16pt plus 2pt}}
\chapter{Further questions}\label{sec:discussion}

We have proved correlation aging for the uniformly initialized CK solution and identified its stationary endpoint. For every $p\ge3$ and $\beta\ge\bd$, the solution satisfies the stationary fluctuation--dissipation theorem at fixed lag, selects a positive marginal plateau and an explicit limiting energy, and forgets every fixed earlier state (\cref{thm:main-aging}). For $\beta>\bd$, we have also proved an integrated violation of the fluctuation--dissipation theorem (\cref{thm:main-response}). In the cubic model, we obtained quantitative restrictions on correlations at three late times for $\beta\ge\bd$ and on relaxation below the plateau for $\beta>\bd$ (\cref{thm:main-geometry,thm:main-fisher-cost}). Several central parts of the physical picture of Cugliandolo and Kurchan \cite{ck1993} remain to be proved. We now formulate them in terms of observables of the CK solution itself, without an asymptotic ansatz.

\section{The complete slow fluctuation--dissipation relation}\label{sec:discussion-fdr}

Fix $\beta>\bd$, write $q=q_\beta$, and let $x=(p-2)(1-q)/q$. For a level $c\in(0,q)$, the predicted integrated relation in the aging regime is
\begin{equation}
 \chi(t_n,s_n)\longrightarrow 1-q+x(q-c)
 \quad\text{if }s_n\to\infty,\ t_n-s_n\to\infty,
 \ C(t_n,s_n)\to c.
 \label{eq:discussion-slow-fdr}
\end{equation}
This is the slow fluctuation--dissipation relation of Cugliandolo and Kurchan, with the ratio $x$ given after \cite[Equation~(12)]{ck1993}. We write it in the conventions of \cref{sec:model}, in which the fluctuation--dissipation theorem in the stationary regime reads $R(t,s)=\partial_sC(t,s)$. We refer to \cite[Sections~3.3--3.4]{bckm1998} for a review of the two time sectors and their response, and to \cite{kimlatz2001} for a numerical test of the relation on the CK equations. The weighted ratio in \eqref{eq:main-weighted-ratio} identifies one moment of the slow response. Proving \eqref{eq:discussion-slow-fdr} requires more information than this moment.

A natural intermediate target is the weak convergence, as finite measures on $[0,1]$, of the response-weighted correlation measures
\[
 \nu_t=\int_0^t R(t,u)\,\delta_{C(t,u)}\,\dd u
\]
to the predicted measure
\[
 x\mathbf1_{(0,q)}(y)\,\dd y
       +\mathbf1_{(q,1)}(y)\,\dd y.
\]
This target includes the statement that no mass accumulates near correlation zero,
\[
 \lim_{\varepsilon\downarrow0}\limsup_{t\to\infty}\nu_t([0,\varepsilon])=0,
\]
and, since $\nu_t([0,1])=\chi(t,0)$, it also includes the convergence of $\chi(t,0)$ to the total mass $1-q+xq$ of the predicted measure. Such a statement concerns the amount of response at each correlation level. To deduce \eqref{eq:discussion-slow-fdr} from it, we must also link the correlation levels to chronological order. Specifically, we would need the response masses
\[
 \int_0^{s_n}R(t_n,u)\,
       \mathbf1_{\{C(t_n,u)>C(t_n,s_n)\}}\,\dd u
 \quad\text{and}\quad
 \int_{s_n}^{t_n}R(t_n,u)\,
       \mathbf1_{\{C(t_n,u)<C(t_n,s_n)\}}\,\dd u
\]
to vanish, together with control at the boundary level $C(t_n,s_n)$. These masses measure returns of the row $u\mapsto C(t_n,u)$ across the level $C(t_n,s_n)$. The product inequality \eqref{eq:main-product} constrains the correlations at three times, but it does not exclude such returns. Note also that the vanishing of the squared defect in \eqref{eq:main-defect-l2} does not control the integrated negative part of the defect over a growing interval.

\section{The sharp susceptibility upper bound}\label{sec:discussion-susceptibility}

In the cubic model with $\beta\ge\bd$, we proved in \eqref{eq:main-susceptibility} that $\liminf_{t\to\infty}\chi(t,0)\ge2(1-q)$. The expected matching estimate is
\begin{equation}
 \limsup_{t\to\infty}\chi(t,0)\le2(1-q).
 \label{eq:discussion-susceptibility-upper}
\end{equation}
For $p=3$, we have $xq=1-q$, so $2(1-q)$ is the total mass $1-q+xq$ of the predicted measure in \cref{sec:discussion-fdr}. Together with the lower bound \eqref{eq:main-susceptibility}, the estimate \eqref{eq:discussion-susceptibility-upper} would determine the limit of the total response mass. This is much more precise than the bound $\chi(t,0)=o(\sqrt t)$ in \eqref{eq:main-uniform-response}. A related scalar target is
\[
 \int_0^t C(t,u)R(t,u)\,\dd u
          \longrightarrow\frac{(1-q)(1+2q)}2,
\]
which is the first moment of the predicted measure when $p=3$. The energy identity fixes the moment with $C^2$ in place of $C$, since $e_\beta(t)=3\beta\int_0^tC(t,u)^2R(t,u)\,\dd u$ converges to the value $e_\beta(\infty)$ given by \eqref{eq:main-energy}. It does not determine the first moment, nor the unweighted mass near correlation zero. In the proof of \cref{thm:geometry-susceptibility-lower}, we obtained the lower bound by positivity, through a trace argument. An upper bound requires a separate estimate on the response-weighted frame used in that proof. Even equality in \eqref{eq:discussion-susceptibility-upper} would leave open the question of chronological returns raised in \cref{sec:discussion-fdr}.

\section{A nondegenerate clock and physical relaxation rates}\label{sec:discussion-clock}

The noncommutation of limits in \eqref{eq:main-aging} does not specify the time scale on which correlations move below $q$. A stronger description would exhibit an increasing function $\mathfrak h\colon[0,\infty)\to(0,\infty)$ and a nonconstant profile $\mathcal C\colon(0,1)\to[0,q]$ such that, for every $a\in(0,1)$,
\[
 C(t_n,s_n)\longrightarrow\mathcal C(a)
 \quad\text{whenever }s_n\to\infty,\ t_n-s_n\to\infty
 \text{ and }\frac{\mathfrak h(s_n)}{\mathfrak h(t_n)}\to a,
\]
and such that sequences with these properties exist for every $a\in(0,1)$. The last requirement excludes clocks for which the statement is vacuous, such as $\mathfrak h(t)=e^t$, for which $\mathfrak h(s_n)/\mathfrak h(t_n)\to0$ whenever $t_n-s_n\to\infty$. Our results fix neither the shape of $\mathcal C$ nor the growth of $\mathfrak h$. In the physical theory, the asymptotic equations for the aging regime admit reparametrizations of time, as observed by Cugliandolo and Kurchan \cite{ck1994}, so they do not select $\mathfrak h$ either. The clock has been studied by numerical integration of the CK equations \cite{kimlatz2001}, and by a matching calculation \cite{andreanovlefevre2006} that proposes a clock within a scaling expansion whose remaining exponent is estimated numerically. For $p=3$ and $\beta\ge\bd$, the composition inequalities for last-crossing times in \cref{cor:geometry-last-crossing} give one-sided restrictions on passage times. A composition law in the limit that is compatible with these inequalities, such as the triangle relations of \cite{ck1994}, would give much more information on the time scale.

For $p=3$, $\beta>\bd$, and a level $c\in(0,q)$, the local Fisher cost \eqref{eq:main-passage-cost} is a constraint in physical time: relaxation to the level $c$ costs a positive amount of $\tau_c(s)\int_s^{s+\tau_c(s)}a_F(u)\,\dd u$. It does not determine the growth of $\tau_c(s)$. We expect that quantitative estimates on the energy current or on the tail of the Fisher density, together with matching dynamical bounds on passage times, would help to determine this growth. The majorant $a_F$ need not be unique, and we do not claim that it defines a physical aging clock. By contrast, the exponent $a$ in \eqref{eq:main-tail-exponent} describes stationary relaxation in the lag variable after the waiting-time limit. It is not an exponent of the aging dynamics.

\section{Selection below the dynamical transition}\label{sec:discussion-below}

In the interval $\beta_{\mathrm m}\le\beta<\bd$, we proved convergence in \cref{sel:all-temperature}, but we did not decide between the zero plateau and the upper marginal root $q_+(\beta)$. The physical picture of the dynamical transition \cite{ck1993,chs1993} predicts the zero plateau throughout the interval, and proving this requires excluding the positive alternative. One sufficient estimate would be $e_\beta(\infty)\le\beta$ throughout the interval, since by \cref{sel:intermediate-range} the positive branch has $e_\beta(\infty)>\beta$. Such an estimate would complete the classification of the stationary plateau from below.

\section{Other orders of limits}\label{sec:discussion-limits}

All of these questions concern a fixed positive temperature and the order of limits in which the dimension tends to infinity before the times do. The bridge of \cref{prop:hard-bridge} transfers the conclusions about correlations and about primitives of the response to the Langevin dynamics in this order, and \cref{prop:energy-bridge} gives the limit of the annealed energy. The annealed response means also converge uniformly on fixed causal time windows. The concentration statement in the bridge concerns the Brownian observables representing their primitives; see \cref{rem:no-pointwise-response}. These fixed-window results give no information on observation times that grow with the dimension, and the low-temperature limit in \cref{rem:low-temperature} is taken after all other limits. We regard a description of the dynamics on time scales that grow with the dimension, and a theorem at zero temperature, as separate problems.


\backmatter
\bibliographystyle{plainnat}
\bibliography{refs,refs-add-A,refs-add-B,refs-add-C,refs-add-D,refs-add-E,refs-add-H,refs-add-I,refs-add-O,refs-add-J,refs-add-K}
\end{document}
